\section{Canonical resolution of marked ideals (Hironaka)}
\label{sec:hironaka}

This section develops the resolution of marked ideals on complex manifolds from which the two resolution
statements of the development, Lemmas~\ref{ana:relative-principalization} and~\ref{ana:embedded-resolution}, are
derived (Lemmas~\ref{hir:factorization-from-resolution} and~\ref{hir:embedded-from-resolution}). The route
is Włodarczyk's algorithm for marked ideals, in the formalism of marked ideals \((M,N,E,\mathcal I,d)\) of
Bierstone and Milman, with Kollár's lectures as a cross-check. The blowup calculus it uses (exceptional divisors,
controlled transforms, strict transforms of submanifolds through the centre, restriction to open sets) is in
Section~\ref{sec:blowup-calculus}.

The route differs from the sources as follows.
\begin{enumerate}
\item The category is that of complex manifolds and holomorphic maps, not smooth varieties. Étale morphisms are
replaced by local biholomorphisms (open embeddings composed with biholomorphisms between open sets); no fibre
product is formed. In the glueing lemma the formal automorphism of Włodarczyk's Lemma ``homo0'' becomes the
\emph{convergent} coordinate change between the coordinate systems \((u,u_2,\dots,u_n)\) and
\((v,u_2,\dots,u_n)\), and the inclusion of its pullback of the ideal in the homogenized ideal is proved from the
convergent Taylor expansion in the first variable.
\item Global finiteness of the algorithm on a compact manifold comes from local finiteness of the invariant and
compactness.
\item The submanifold \(N\) of a marked ideal is a closed submanifold of the fixed ambient manifold \(M\), and the
ideal of \(\mathcal O_N\) is represented by an ideal of \(\mathcal O_M\) of which only the restrictions to \(N\)
matter. All blowups are blowups of the ambient manifold; the ambient manifold never changes along the induction
on \(\dim N\).
\item The order of an ideal at a point is defined with powers of the maximal ideal of germs of holomorphic
functions (no completion); its characterisation by derivatives is a lemma (Lemma~\ref{hir:order-taylor}).
\item Steps of a sequence of blowups whose centre is empty are allowed; they are biholomorphisms. They encode the
isomorphisms in Włodarczyk's multiple test blow-ups and the ``trivial blowings-up'' that Bierstone and Milman insert
when a resolution is restricted to an open set.
\end{enumerate}

\subsection{Order of an ideal along a submanifold}

\paragraph{Notation.} For a complex manifold \(M\) and \(x\in M\), \(\mathfrak m_x=\mathcal I_{\{x\}}\) is the ideal
sheaf of germs vanishing at \(x\), and \(\mathfrak m_x^q\) its \(q\)-th power (\(\mathfrak m_x^0=\mathcal O_M\)).
Ideal sheaves are recorded by their stalks as in Definition~\ref{ana:ideal-sheaves}; products and powers are
generated stalkwise by products of germs.

\begin{definition}[Order along a submanifold]
\label{hir:order}
\uses{ana:ideal-sheaves}
\source{wlodarczyk05-hironaka,bierstone-milman08-functoriality}
\lean{TranscendentalNakamaye.Pending.GermIdealOrder.pointIdeal,TranscendentalNakamaye.Pending.GermIdealOrder.VanishesToOrderOn,TranscendentalNakamaye.Pending.GermIdealOrder.OrderGE,TranscendentalNakamaye.Pending.GermIdealOrder.order}
\leanok
Let \(N\subset M\) be a subset and \(x\in M\). A function \(g\) holomorphic near \(x\) \emph{vanishes to order
\(q\) along \(N\) at \(x\)} if there is \(h\in\mathfrak m_x^q\) with \(g=h\) on \(N\) near \(x\). For an ideal
\(\mathfrak a\subset\mathcal O_M\), \(\operatorname{ord}_x(\mathfrak a|_N)\ge q\) means that every germ of
\(\mathfrak a_x\) vanishes to order \(q\) along \(N\) at \(x\), and
\(\operatorname{ord}_x(\mathfrak a|_N)\in\mathbb N\cup\{\infty\}\) is the supremum of such \(q\).
\end{definition}
When \(N\) is a closed submanifold through \(x\), \(\mathcal O_{N,x}=\mathcal O_{M,x}/\mathcal I_{N,x}\) and the
maximal ideal of \(\mathcal O_{N,x}\) is the image of \(\mathfrak m_x\); so \(g\) vanishes to order \(q\) along
\(N\) exactly when \(g|_N\in\mathfrak m_{N,x}^q\), and \(\operatorname{ord}_x(\mathfrak a|_N)\) is the order of the
ideal \(\mathfrak a|_N\) of \(\mathcal O_N\) used by Bierstone--Milman. For \(N=M\) it is Włodarczyk's
\(\operatorname{ord}_x\).

\begin{lemma}[Elementary properties of the order]
\label{hir:order-basic}
\uses{hir:order}
\lean{TranscendentalNakamaye.Pending.GermIdealOrder.idealPow_pointIdeal_succ_le,TranscendentalNakamaye.Pending.GermIdealOrder.idealPow_le_of_le,TranscendentalNakamaye.Pending.GermIdealOrder.mul_mem_idealPow_add,TranscendentalNakamaye.Pending.GermIdealOrder.VanishesToOrderOn.mono,TranscendentalNakamaye.Pending.GermIdealOrder.VanishesToOrderOn.mul,TranscendentalNakamaye.Pending.GermIdealOrder.VanishesToOrderOn.mul_left,TranscendentalNakamaye.Pending.GermIdealOrder.vanishesToOrderOn_one_iff,TranscendentalNakamaye.Pending.GermIdealOrder.le_order_iff}
\leanok
Let \(N\subset M\), \(x\in M\), and let \(f,g\) be holomorphic near \(x\).
\begin{enumerate}
\item \(\mathfrak m_x^{p}\subset\mathfrak m_x^{q}\) for \(q\le p\), and \(\mathfrak m_x^p\,\mathfrak m_x^q\subset\mathfrak m_x^{p+q}\).
\item If \(f\) vanishes to order \(p\) and \(g\) to order \(q\) along \(N\) at \(x\), then \(fg\) vanishes to order
\(p+q\); a function vanishing to order \(p\) vanishes to every order \(q\le p\).
\item If \(x\in N\), then \(g\) vanishes to order \(1\) along \(N\) at \(x\) if and only if \(g(x)=0\).
\item For an ideal \(\mathfrak a\) of holomorphic germs, \(q\le\operatorname{ord}_x(\mathfrak a|_N)\) if and only if
\(\operatorname{ord}_x(\mathfrak a|_N)\ge q\) in the sense of Definition~\ref{hir:order}.
\end{enumerate}
\end{lemma}
\begin{proof}
\leanok
(1) A generator of \(\mathfrak m_x^{q+1}=\mathfrak m_x^q\mathfrak m_x\) is a product \(ab\) with \(a\in\mathfrak m_x^q\) and \(b\)
holomorphic, and \(\mathfrak m_x^q\) is an ideal; induction on \(p-q\). For the product, induct on \(q\): a generator of
\(\mathfrak m_x^{q+1}\) is \(bc\) with \(b\in\mathfrak m_x^q\), \(c\in\mathfrak m_x\), and \(a(bc)=(ab)c\) with
\(ab\in\mathfrak m_x^{p+q}\). (2) If \(f=h\) and \(g=k\) on \(N\) near \(x\) with \(h\in\mathfrak m_x^p\), \(k\in\mathfrak m_x^q\),
then \(fg=hk\) on \(N\) near \(x\) and \(hk\in\mathfrak m_x^{p+q}\) by (1); the second claim is (1). (3) Every germ of
\(\mathfrak m_x=\mathfrak m_x^0\mathfrak m_x\) vanishes at \(x\), so \(g(x)=h(x)=0\) when \(x\in N\). Conversely if \(g(x)=0\) then
\(g=1\cdot g\in\mathfrak m_x\). (4) If the order is \(\ge q\) the supremum is \(\ge q\). If it is not \(\ge q\), then by (2) no
\(k\ge q\) occurs, so every \(k\) that occurs is \(\le q-1\) (and \(q\ge1\), the order being always \(\ge0\) for an ideal
of holomorphic germs), hence the supremum is \(\le q-1<q\).
\end{proof}

\begin{definition}[Cosupport]
\label{hir:cosupport}
\uses{hir:order}
\source{wlodarczyk05-hironaka,bierstone-milman08-functoriality,kollar07-resolution}
\lean{TranscendentalNakamaye.Pending.GermIdealOrder.cosupport}
\leanok
For \(N\subset M\), an ideal \(\mathfrak a\subset\mathcal O_M\) and \(d\in\mathbb N\),
\(\operatorname{cosupp}(\mathfrak a|_N,d)=\{x\in N:\operatorname{ord}_x(\mathfrak a|_N)\ge d\}\).
\end{definition}

\begin{lemma}[Elementary properties of the cosupport]
\label{hir:cosupport-basic}
\uses{hir:cosupport,hir:order-basic}
\source{kollar07-resolution}
\lean{TranscendentalNakamaye.Pending.GermIdealOrder.cosupport_subset,TranscendentalNakamaye.Pending.GermIdealOrder.cosupport_anti,TranscendentalNakamaye.Pending.GermIdealOrder.cosupport_zero,TranscendentalNakamaye.Pending.GermIdealOrder.cosupport_one,TranscendentalNakamaye.Pending.GermIdealOrder.cosupport_unitIdeal}
\leanok
Let \(\mathfrak a\) be an ideal of holomorphic germs and \(N\subset M\). Then
\(\operatorname{cosupp}(\mathfrak a|_N,d')\subset\operatorname{cosupp}(\mathfrak a|_N,d)\subset N\) for \(d\le d'\);
\(\operatorname{cosupp}(\mathfrak a|_N,0)=N\); \(\operatorname{cosupp}(\mathfrak a|_N,1)=N\cap V(\mathfrak a)\), where
\(V(\mathfrak a)\) is the zero locus; and \(\operatorname{cosupp}(\mathcal O_M|_N,d)=\emptyset\) for \(d\ge1\).
\end{lemma}
\begin{proof}
\leanok
The inclusions are Lemma~\ref{hir:order-basic}(2). Every holomorphic germ vanishes to order \(0\). Order \(\ge1\) at
\(x\in N\) means that every germ of \(\mathfrak a_x\) vanishes at \(x\) (Lemma~\ref{hir:order-basic}(3)), i.e.
\(x\in V(\mathfrak a)\). The constant \(1\in\mathcal O_{M,x}\) does not vanish at \(x\).
\end{proof}

\subsection{Simple normal crossings}

\begin{definition}[Adapted coordinates, simple normal crossings, transversality]
\label{hir:snc}
\uses{ana:blowups}
\source{wlodarczyk05-hironaka,bierstone-milman08-functoriality}
\lean{TranscendentalNakamaye.Pending.NormalCrossings.coordSubspace,TranscendentalNakamaye.Pending.NormalCrossings.IsCoordinateFamilyAt,TranscendentalNakamaye.Pending.NormalCrossings.IsSNCDivisor,TranscendentalNakamaye.Pending.NormalCrossings.HasNormalCrossings,TranscendentalNakamaye.Pending.NormalCrossings.IsTransverse}
\leanok
For \(S\subset\{1,\dots,n\}\) let \(L_S=\{z\in\mathbf C^n:z_j=0\text{ for }j\in S\}\). A finite family \((A_i)\) of subsets of
\(M\) is \emph{adapted at \(x\)} to \((S_i)\) if there are holomorphic coordinates \(\phi\) at \(x\) with \(\phi(x)=0\) in which every
\(A_i\) with \(x\in A_i\) is \(\phi^{-1}(L_{S_i})\) near \(x\).
\begin{enumerate}
\item An \emph{ordered simple normal crossings divisor} is a finite sequence \(E=(H_1,\dots,H_s)\) of closed subsets such that
at every \(x\in M\) the family is adapted to some \((S_i)\) with \(|S_i|=1\) for \(x\in H_i\) and \(S_i\ne S_j\) for distinct \(i,j\)
with \(x\in H_i\cap H_j\).
\item A subset \(C\) has \emph{only normal crossings with} \(E\) if at every \(x\in C\) the family \((C,H_1,\dots,H_s)\) is adapted to
some \((S_0,S_1,\dots,S_s)\) with \(|S_i|=1\) for \(x\in H_i\).
\item \(E\) is \emph{transverse to} \(N\) if moreover, at every \(x\in N\), \(S_i\cap S_0=\emptyset\) for \(x\in H_i\) (with \(C=N\)).
\end{enumerate}
\end{definition}
The members of \(E\) are smooth hypersurfaces, not necessarily connected, and may be empty; the order of \(E\) is the
order of the sequence, a new exceptional divisor being appended last.

\begin{lemma}[Normal crossings with the empty divisor]
\label{hir:snc-basic}
\uses{hir:snc}
\lean{TranscendentalNakamaye.Pending.NormalCrossings.isSNCDivisor_nil,TranscendentalNakamaye.Pending.NormalCrossings.hasNormalCrossings_nil_of_isClosedSubmanifold,TranscendentalNakamaye.Pending.NormalCrossings.isTransverse_nil_of_isClosedSubmanifold,TranscendentalNakamaye.Pending.NormalCrossings.IsTransverse.hasNormalCrossings,TranscendentalNakamaye.Pending.NormalCrossings.coordSubspace_eq_complexCoordinatePlane}
\leanok
The empty sequence is a simple normal crossings divisor. A closed submanifold has only normal crossings with the empty
divisor, and the empty divisor is transverse to it. Transversality implies normal crossings.
\end{lemma}
\begin{proof}
\leanok
Centred coordinates exist at every point (a chart followed by a translation), and the empty family is adapted to them.
At a point of a closed submanifold \(C\), flattening coordinates make \(C\) the coordinate plane
\(\mathbf C^k\times0=L_{\{k+1,\dots,n\}}\). The last claim forgets the disjointness condition.
\end{proof}

\subsection{Towers of blowups}

\begin{definition}[Tower of blowups]
\label{hir:blowup-tower}
\uses{ana:blowups}
\source{wlodarczyk05-hironaka,bierstone-milman08-functoriality}
\lean{TranscendentalNakamaye.Pending.BlowupTowers.CMfd,TranscendentalNakamaye.Pending.BlowupTowers.BlowupTower,TranscendentalNakamaye.Pending.BlowupTowers.BlowupTower.proj,TranscendentalNakamaye.Pending.BlowupTowers.BlowupTower.length,TranscendentalNakamaye.Pending.BlowupTowers.BlowupTower.append}
\leanok
A \emph{tower of blowups} over a complex manifold \(X\) is a finite sequence
\(X=X_0\xleftarrow{\sigma_1}X_1\leftarrow\cdots\xleftarrow{\sigma_r}X_r\) together with closed submanifolds \(C_{i-1}\subset X_{i-1}\)
such that \(\sigma_i\) is a blowup of \(X_{i-1}\) along \(C_{i-1}\) (Definition~\ref{ana:blowups}; \(C_{i-1}=\emptyset\) is allowed,
and then \(\sigma_i\) is a biholomorphism). Its \emph{length} is \(r\), its \emph{projection} is
\(\sigma_1\circ\cdots\circ\sigma_r\), and the \emph{concatenation} of a tower over \(X\) ending at \(X_r\) with a tower over \(X_r\) is
the tower obtained by continuing the first with the second.
\end{definition}

\begin{lemma}[Towers and modifications]
\label{hir:tower-sequence}
\uses{hir:blowup-tower}
\lean{TranscendentalNakamaye.Pending.BlowupTowers.BlowupTower.isBlowupSequence,TranscendentalNakamaye.Pending.BlowupTowers.BlowupTower.exists_proj_eq_of_isBlowupSequence,TranscendentalNakamaye.Pending.BlowupTowers.BlowupTower.append_assoc,TranscendentalNakamaye.Pending.BlowupTowers.BlowupTower.proj_append,TranscendentalNakamaye.Pending.BlowupTowers.BlowupTower.length_append}
\leanok
The projection of a tower is a finite composition of blowups with smooth centres, and every such composition is the
projection of a tower. Concatenation is associative; the projection of a concatenation is the composite of the
projections, and its length is the sum of the lengths.
\end{lemma}
\begin{proof}
\leanok
Induction on the length; for the converse, induction on the number of blowups of the composition.
\end{proof}

\subsection{Marked ideals}

\begin{definition}[Marked ideal]
\label{hir:marked-ideal}
\uses{hir:cosupport,hir:snc}
\source{bierstone-milman08-functoriality,wlodarczyk05-hironaka,kollar07-resolution}
\lean{TranscendentalNakamaye.Pending.MarkedIdeals.MarkedIdeal,TranscendentalNakamaye.Pending.MarkedIdeals.MarkedIdeal.cosupp,TranscendentalNakamaye.Pending.MarkedIdeals.MarkedIdeal.IsWellFormed}
\leanok
A \emph{marked ideal} on a complex manifold \(M\) is \(\underline{\mathcal I}=(M,N,E,\mathfrak a,d)\): a subset \(N\subset M\), a finite
sequence \(E\) of subsets of \(M\), an ideal \(\mathfrak a\subset\mathcal O_M\) (representing the ideal \(\mathfrak a|_N\) of \(\mathcal O_N\)) and
\(d\in\mathbb N\). Its \emph{cosupport} is \(\operatorname{cosupp}\underline{\mathcal I}=\operatorname{cosupp}(\mathfrak a|_N,d)\). It is
\emph{well formed} if \(N\) is a closed submanifold, \(\mathfrak a\) is coherent, and \(E\) is an ordered simple normal crossings
divisor transverse to \(N\).
\end{definition}

\begin{definition}[Admissible centre and transform]
\label{hir:transform}
\uses{hir:marked-ideal,ana:blowups}
\source{bierstone-milman08-functoriality,wlodarczyk05-hironaka,kollar07-resolution}
\lean{TranscendentalNakamaye.Pending.MarkedIdeals.MarkedIdeal.IsAdmissibleCentre,TranscendentalNakamaye.Pending.MarkedIdeals.MarkedIdeal.transform,TranscendentalNakamaye.Pending.MarkedIdeals.MarkedIdeal.isAdmissibleCentre_empty,TranscendentalNakamaye.Pending.MarkedIdeals.MarkedIdeal.length_E_transform,TranscendentalNakamaye.Pending.MarkedIdeals.MarkedIdeal.getLast_E_transform}
\leanok
A closed submanifold \(C\subset M\) is an \emph{admissible centre} for \(\underline{\mathcal I}=(M,N,E,\mathfrak a,d)\) if
\(C\subset\operatorname{cosupp}\underline{\mathcal I}\) and \(C\) has only normal crossings with \(E\); the empty set is admissible. Let
\(\sigma\colon M'\to M\) be a blowup along \(C\). The \emph{transform} of \(\underline{\mathcal I}\) is
\((M',N',E',\mathfrak a',d)\), where \(N'=\overline{\sigma^{-1}(N\setminus C)}\) is the strict transform of \(N\) relative to the centre, \(E'\) is
the sequence of the strict transforms \(\overline{\sigma^{-1}(H\setminus C)}\) of the members \(H\) of \(E\) followed by \(\sigma^{-1}(C)\), and
\[
 \mathfrak a'=\bigl\{g:\ g\,\ell\in\sigma^*\mathfrak a+\mathcal I_{N'}\text{ for every }\ell\in(\sigma^*\mathcal I_C)^d\bigr\},
\]
the quotient along \(N'\) of \(\sigma^*\mathfrak a\) by the \(d\)-th power of the ideal of the exceptional divisor. On \(N'\) this is
\(\mathcal I_{\sigma^{-1}(C)}^{-d}\,\sigma^*(\mathfrak a|_N)\).
\end{definition}
The transform depends on \(C\), not only on \(\sigma\): along a centre of codimension one \(\sigma\) is a biholomorphism, but the
transform still divides by \(\mathcal I_C^d\) and appends \(C\) to \(E\).

\begin{definition}[Multiple test blow-up, resolution]
\label{hir:test-tower}
\uses{hir:transform,hir:blowup-tower}
\source{wlodarczyk05-hironaka,bierstone-milman08-functoriality}
\lean{TranscendentalNakamaye.Pending.MarkedIdeals.MarkedIdeal.transformAlong,TranscendentalNakamaye.Pending.MarkedIdeals.MarkedIdeal.IsTestTower,TranscendentalNakamaye.Pending.MarkedIdeals.MarkedIdeal.IsResolution}
\leanok
The transform of \(\underline{\mathcal I}\) along a tower is obtained by transforming stage by stage. A tower is a \emph{multiple test
blow-up} of \(\underline{\mathcal I}\) if each centre \(C_{i-1}\) is admissible for the transform of \(\underline{\mathcal I}\) on
\(X_{i-1}\), and a \emph{resolution} of \(\underline{\mathcal I}\) if moreover the transform on \(X_r\) has empty cosupport.
\end{definition}

\begin{lemma}[Bookkeeping along towers]
\label{hir:test-tower-basic}
\uses{hir:test-tower,hir:tower-sequence}
\lean{TranscendentalNakamaye.Pending.MarkedIdeals.MarkedIdeal.transformAlong_append,TranscendentalNakamaye.Pending.MarkedIdeals.MarkedIdeal.isTestTower_append_iff,TranscendentalNakamaye.Pending.MarkedIdeals.MarkedIdeal.IsResolution.append,TranscendentalNakamaye.Pending.MarkedIdeals.MarkedIdeal.isResolution_refl_iff,TranscendentalNakamaye.Pending.MarkedIdeals.MarkedIdeal.length_E_transformAlong,TranscendentalNakamaye.Pending.MarkedIdeals.MarkedIdeal.d_transformAlong}
\leanok
The transform along a concatenation is the transform along the second tower of the transform along the first. A
concatenation is a multiple test blow-up of \(\underline{\mathcal I}\) if and only if the first tower is one and the second is one
of the transform; a multiple test blow-up followed by a resolution of the transform is a resolution. The empty tower is a
resolution exactly when the cosupport is empty. Along a tower of length \(r\) the mark is unchanged and \(E\) gains exactly
\(r\) members.
\end{lemma}
\begin{proof}
\leanok
Induction on the length of the second tower, using that the transform along a tower is defined stage by stage.
\end{proof}

\subsection{Strict transforms relative to the centre}

\begin{definition}[Strict transform relative to the centre; corrected transform]
\label{hir:geo-transform}
\uses{hir:transform,hir:test-tower,bcalc:strict-submanifold}
\source{bierstone-milman08-functoriality,wlodarczyk05-hironaka}
\lean{TranscendentalNakamaye.Pending.MarkedIdeals.MarkedIdeal.geoTransform,TranscendentalNakamaye.Pending.MarkedIdeals.MarkedIdeal.geoTransformAlong,TranscendentalNakamaye.Pending.MarkedIdeals.MarkedIdeal.IsGeoTestTower,TranscendentalNakamaye.Pending.MarkedIdeals.MarkedIdeal.IsGeoResolution}
\leanok
Let \(\sigma\colon M'\to M\) be a blowup along \(C\). The \emph{strict transform} of a subset \(A\subset M\) is
\(A'=\overline{\sigma^{-1}(A\setminus C)}\). The \emph{transform} of \(\underline{\mathcal I}=(M,N,E,\mathfrak a,d)\) is defined as in
Definition~\ref{hir:transform}, with \(N'=\overline{\sigma^{-1}(N\setminus C)}\) and the members \(H'=\overline{\sigma^{-1}(H\setminus C)}\) of \(E'\)
(followed by \(\sigma^{-1}(C)\)); transforms along towers, multiple test blow-ups and resolutions are defined from it as in
Definition~\ref{hir:test-tower}. This is the transform of Definition~\ref{hir:transform}.
\end{definition}
The set \(\overline{\sigma^{-1}(A)\setminus\operatorname{Exc}\sigma}\) agrees with \(A'\) when \(\operatorname{Exc}\sigma=\sigma^{-1}(C)\), e.g.\ when \(C\) has
codimension at least two everywhere or is empty, but not along a centre of codimension one, where \(\sigma\) is biholomorphic and
\(\operatorname{Exc}\sigma=\emptyset\): there it keeps \(A\cap C\). This matters for the algorithm: if a member \(H\) of \(E\) is blown up again as a
centre \(C=H\) of codimension one (Włodarczyk's treatment of codimension-one components of the cosupport, e.g.\ \(\mathfrak a=\mathfrak m^{2\mu}\) on
a surface after one point blowup), then \(H'=\emptyset\) and \(E'=(\dots,\emptyset,\dots,H)\) is again a simple normal crossings divisor, whereas
keeping \(H\) would list it twice.

\begin{lemma}[Bookkeeping for the corrected transform]
\label{hir:geo-transform-basic}
\uses{hir:geo-transform,hir:tower-sequence,bcalc:strict-geo}
\lean{TranscendentalNakamaye.Pending.MarkedIdeals.MarkedIdeal.geoTransform_ideal,TranscendentalNakamaye.Pending.MarkedIdeals.MarkedIdeal.length_E_geoTransform,TranscendentalNakamaye.Pending.MarkedIdeals.MarkedIdeal.getLast_E_geoTransform,TranscendentalNakamaye.Pending.MarkedIdeals.MarkedIdeal.geoStrictTransform_eq_empty_of_subset,TranscendentalNakamaye.Pending.MarkedIdeals.MarkedIdeal.geoStrictTransform_eq_strictTransform_of_exceptionalLocus_eq,TranscendentalNakamaye.Pending.MarkedIdeals.MarkedIdeal.geoTransform_eq_transform,TranscendentalNakamaye.Pending.MarkedIdeals.MarkedIdeal.geoTransform_empty_eq_transform,TranscendentalNakamaye.Pending.MarkedIdeals.MarkedIdeal.geoTransformAlong_append,TranscendentalNakamaye.Pending.MarkedIdeals.MarkedIdeal.isGeoTestTower_append_iff,TranscendentalNakamaye.Pending.MarkedIdeals.MarkedIdeal.IsGeoResolution.append,TranscendentalNakamaye.Pending.MarkedIdeals.MarkedIdeal.isGeoResolution_refl_iff,TranscendentalNakamaye.Pending.MarkedIdeals.MarkedIdeal.d_geoTransformAlong,TranscendentalNakamaye.Pending.MarkedIdeals.MarkedIdeal.length_E_geoTransformAlong,TranscendentalNakamaye.Pending.MarkedIdeals.MarkedIdeal.N_geoTransformAlong_step}
\leanok
A subset of the centre has empty strict transform. The strict transform of Definition~\ref{hir:geo-transform} is contained in
\(\overline{\sigma^{-1}(A)\setminus\operatorname{Exc}\sigma}\), with equality when \(\operatorname{Exc}\sigma=\sigma^{-1}(C)\); then the two transforms of a marked
ideal coincide, in particular for the empty centre. Along towers the statements of Lemma~\ref{hir:test-tower-basic} hold for the
corrected transform: transforms compose along concatenations, concatenations of multiple test blow-ups are characterised stagewise, a
multiple test blow-up followed by a resolution of the transform is a resolution, the empty tower resolves exactly the marked ideals with
empty cosupport, the mark is unchanged and \(E\) gains one member per blowup.
\end{lemma}
\begin{proof}
\leanok
\(\overline{\sigma^{-1}(A\setminus C)}=\overline{\emptyset}\) if \(A\subset C\); \(\operatorname{Exc}\sigma\subset\sigma^{-1}(C)\) because \(\sigma\) is biholomorphic over \(M\setminus C\)
(the blowup is an isomorphism over the complement of its centre), which gives the inclusion and, under \(\operatorname{Exc}\sigma=\sigma^{-1}(C)\), the equality; for \(C=\emptyset\) both sides of
\(\operatorname{Exc}\sigma\subset\sigma^{-1}(C)\) are empty. The tower statements are proved by induction on the second tower exactly as in
Lemma~\ref{hir:test-tower-basic}.
\end{proof}

\begin{lemma}[Well-formedness of the corrected transform: submanifold and ideal]
\label{hir:geo-transform-wellformed}
\uses{hir:geo-transform,bcalc:strict-submanifold,bcalc:strict-ideal,bcalc:controlled-transform}
\lean{TranscendentalNakamaye.Pending.MarkedIdeals.MarkedIdeal.isClosedSubmanifold_geoTransform_N,TranscendentalNakamaye.Pending.MarkedIdeals.MarkedIdeal.one_mem_idealPow,TranscendentalNakamaye.Pending.MarkedIdeals.MarkedIdeal.one_mem_pullback_vanishingIdeal,TranscendentalNakamaye.Pending.MarkedIdeals.MarkedIdeal.isInvertibleAlong_geoStrictTransform,TranscendentalNakamaye.Pending.MarkedIdeals.MarkedIdeal.restrictLe_pullback,TranscendentalNakamaye.Pending.MarkedIdeals.MarkedIdeal.isCoherent_geoTransform_ideal}
\leanok
Let \(\underline{\mathcal I}=(M,N,E,\mathfrak a,d)\) with \(N\) a closed submanifold, \(C\subset N\) a closed submanifold and \(\sigma\) a blowup along \(C\). Then
\(N'\) is a closed submanifold of \(M'\), and the ideal \((\sigma^*\mathcal I_C)^d\) is invertible along \(N'\): near every point of \(N'\) it is generated by
one holomorphic function vanishing identically on no nonempty open subset of \(N'\). If \(\mathfrak a\) is coherent and
\(\mathfrak a|_N\subset\mathcal I_C^d|_N\), the transformed ideal \(\mathfrak a'\) is coherent.
\end{lemma}
\begin{proof}
\leanok
At a point \(z\in N'\) the image \(\sigma(z)\) lies in \(N\), where \(N\) is regular of some dimension \(k_N\); if \(\sigma(z)\in C\), Lemma~\ref{bcalc:strict-submanifold}
applied near \(\sigma(z)\) shows that \(N'\) is regular of dimension \(k_N\) at \(z\), and otherwise \(\sigma\) is a local biholomorphism at \(z\) and \(N'=\sigma^{-1}(N)\)
near \(z\). For invertibility along \(N'\), let \(h\) generate \((\sigma^*\mathcal I_C)^d\) on an open \(U\) and let \(V\subset U\) be open with \(V\cap N'\ne\emptyset\).
Since \(N'\) is the closure of \(\sigma^{-1}(N\setminus C)\), \(V\) contains a point \(y\in N'\) with \(\sigma(y)\notin C\); there \(1\in\mathcal I_C\), so \(1\in(\sigma^*\mathcal I_C)^d\)
at \(y\), and \(1=\sum a_i h\) near \(y\) forces \(h(y)\ne0\). The inclusion \(\mathfrak a|_N\subset\mathcal I_C^d|_N\) pulls back to
\(\sigma^*\mathfrak a|_{\sigma^{-1}N}\subset\sigma^*(\mathcal I_C^d)|_{\sigma^{-1}N}=(\sigma^*\mathcal I_C)^d|_{\sigma^{-1}N}\), and \(N'\subset\sigma^{-1}(N)\). The quotient along a closed set \(W\) of a
coherent ideal by a coherent ideal invertible along \(W\) that contains it along \(W\) is coherent when \(\mathcal I_W\) is coherent
(the restricted quotient of Copy-4 Step~1), and \(\mathcal I_{N'}\) is coherent since \(N'\) is a closed submanifold.
\end{proof}
For \(N=M\) the hypothesis \(\mathfrak a\subset\mathcal I_C^d\) is the Hadamard bridge (Lemma~\ref{hir:hadamard-bridge}) for an admissible centre. With the
strict transform \(\overline{\sigma^{-1}(N)\setminus\operatorname{Exc}\sigma}\) the invertibility along \(N'\) fails when a codimension-one centre contains a
component of \(N\).

\begin{lemma}[The corrected transform for \(N=M\)]
\label{hir:geo-transform-ambient}
\uses{hir:geo-transform-basic,hir:geo-transform-wellformed,hir:bridge-consumers,bcalc:tower-exceptional,bcalc:principalization}
\source{wlodarczyk05-hironaka}
\lean{TranscendentalNakamaye.Pending.MarkedIdeals.MarkedIdeal.geoStrictTransform_univ,TranscendentalNakamaye.Pending.MarkedIdeals.MarkedIdeal.geoTransform_N_eq_univ,TranscendentalNakamaye.Pending.MarkedIdeals.MarkedIdeal.geoTransform_ideal_eq_controlledTransform,TranscendentalNakamaye.Pending.MarkedIdeals.MarkedIdeal.pullback_eq_mul_geoTransformAlong,TranscendentalNakamaye.Pending.MarkedIdeals.MarkedIdeal.pullback_eq_exceptionalDivisor_of_isGeoResolution,TranscendentalNakamaye.Pending.MarkedIdeals.MarkedIdeal.isInvertible_pullback_of_isGeoResolution}
\leanok
For \(N=M\) the strict transform of \(M\) is \(M'\) and the transformed ideal is the controlled transform \(\sigma^{\mathrm c}(\mathfrak a,d)\). Along a
multiple test blow-up \(T\) of \((M,M,E,\mathfrak a,d)\) with \(\mathfrak a\) coherent, every transform \(\mathfrak a_T\) is coherent and
\(T^*\mathfrak a=\mathcal O(-dE_T)\,\mathfrak a_T\), where \(E_T\) is the exceptional divisor of the tower. If \(T\) is a resolution and \(d=1\), then
\(T^*\mathfrak a=\mathcal O(-E_T)\) is invertible.
\end{lemma}
\begin{proof}
\leanok
\(\sigma^{-1}(C)\) has empty interior, so \(\overline{\sigma^{-1}(M\setminus C)}=M'\), and the quotient along \(M'\) is the colon ideal
\((\sigma^*\mathfrak a:\mathcal O(-dD))=\sigma^{\mathrm c}(\mathfrak a,d)\). Induction on \(T\): at a step along an admissible centre \(C\), the Hadamard bridge gives
\(\mathfrak a_T\subset\mathcal I_C^d\), so \(\sigma^*\mathfrak a_T=\mathcal O(-dD)\,\sigma^{\mathrm c}(\mathfrak a_T,d)\) with a coherent controlled transform
(Lemma~\ref{bcalc:controlled-transform}); then \((T\sigma)^*\mathfrak a=\sigma^*\mathcal O(-dE_T)\,\mathcal O(-dD)\,\mathfrak a_{T\sigma}=\mathcal O(-d(\sigma^*E_T+D))\,\mathfrak a_{T\sigma}\).
With \(d=1\) the cosupport of \(\mathfrak a_T\) is its zero locus; if it is empty, \(\mathfrak a_T=\mathcal O\).
\end{proof}

\subsection{Derivative ideals, order and the Hadamard bridge}

\begin{definition}[Derivative ideals]
\label{hir:derivative-ideal}
\uses{ana:ideal-sheaves}
\source{wlodarczyk05-hironaka,bierstone-milman08-functoriality}
\lean{TranscendentalNakamaye.Pending.DerivativeIdeal.chartDeriv,TranscendentalNakamaye.Pending.DerivativeIdeal.DerivGen,TranscendentalNakamaye.Pending.DerivativeIdeal.derivIdeal,TranscendentalNakamaye.Pending.DerivativeIdeal.iterDerivIdeal}
\leanok
For a chart \(e\) of \(M\), a direction \(v\in\mathbf C^n\) and \(f\) holomorphic near a point of the domain of \(e\), let
\(\partial^e_vf=D(f\circ e^{-1})(e(\cdot))\,v\). For an ideal \(\mathfrak a\subset\mathcal O_M\), \(\mathcal D(\mathfrak a)_x\) is the ideal generated by the
germs of \(\mathfrak a_x\) and their derivatives \(\partial^e_vf\) for every chart \(e\) at \(x\) and every \(v\); \(\mathcal D^0(\mathfrak a)=\mathfrak a\) and
\(\mathcal D^{j+1}(\mathfrak a)=\mathcal D(\mathcal D^j(\mathfrak a))\).
\end{definition}
These are the derivative ideals of marked ideals with \(N=M\), and they are used only for \(N=M\) in this section. For
\(N\subsetneq M\) the derivatives of a marked ideal are taken along \(N\) (Bierstone and Milman's \(\mathrm{Der}_N\)); the ambient
\(\mathcal D(\mathfrak a)\) restricted to \(N\) can be larger: for \(N=\{u=0\}\) and \(\mathfrak a=(u)\), \(\mathfrak a|_N=0\) but
\(1=\partial_uu\in\mathcal D(\mathfrak a)\).

\begin{lemma}[Derivative ideals in one chart]
\label{hir:derivative-ideal-generators}
\uses{hir:derivative-ideal}
\source{bierstone-milman08-functoriality}
\lean{TranscendentalNakamaye.Pending.DerivativeIdeal.chartDeriv_congr,TranscendentalNakamaye.Pending.DerivativeIdeal.chartDeriv_eventuallyEq_sum_mul,TranscendentalNakamaye.Pending.DerivativeIdeal.chartDeriv_eventuallyEq_sum,TranscendentalNakamaye.Pending.DerivativeIdeal.holomorphicAt_chartDeriv,TranscendentalNakamaye.Pending.DerivativeIdeal.holomorphicAt_transitionCoeff,TranscendentalNakamaye.Pending.DerivativeIdeal.derivFamily,TranscendentalNakamaye.Pending.DerivativeIdeal.derivIdeal_mem_iff_of_generators,TranscendentalNakamaye.Pending.DerivativeIdeal.isCoherent_derivIdeal,TranscendentalNakamaye.Pending.DerivativeIdeal.isCoherent_iterDerivIdeal,TranscendentalNakamaye.Pending.DerivativeIdeal.le_derivIdeal,TranscendentalNakamaye.Pending.DerivativeIdeal.derivIdeal_mono,TranscendentalNakamaye.Pending.DerivativeIdeal.le_iterDerivIdeal}
\leanok
If \(f_1,\dots,f_m\), holomorphic on an open set \(U\) contained in the domain of a chart \(e\), generate \(\mathfrak a\) at every point of \(U\),
then at every point of \(U\) the ideal \(\mathcal D(\mathfrak a)\) is generated by the \(f_i\) and the \(\partial^e_{e_j}f_i\), \(1\le j\le n\). In particular
\(\mathcal D^j(\mathfrak a)\) is coherent for every \(j\) when \(\mathfrak a\) is, and \(\mathfrak a\subset\mathcal D(\mathfrak a)\).
\end{lemma}
\begin{proof}
\leanok
The listed functions are generators of \(\mathcal D(\mathfrak a)\) at each point of \(U\). Conversely, let \(y\in U\), \(e'\) a chart at \(y\), \(v\in\mathbf C^n\) and
\(f\in\mathfrak a_y\), so \(f=\sum_k a_kf_{i_k}\) near \(y\) with \(a_k\) holomorphic. The derivative \(\partial^{e'}_v\) depends only on the germ, and by the
product rule \(\partial^{e'}_vf=\sum_k(\partial^{e'}_va_k)f_{i_k}+a_k\,\partial^{e'}_vf_{i_k}\) near \(y\). By the chain rule applied to
\(f\circ e'^{-1}=(f\circ e^{-1})\circ(e\circ e'^{-1})\),
\(\partial^{e'}_vg=\sum_j c_j\,\partial^e_{e_j}g\) with \(c_j=\bigl(D(e\circ e'^{-1})(e'(\cdot))\,v\bigr)_j\) holomorphic. Hence \(\partial^{e'}_vf\) is a holomorphic
combination of the listed functions. Coherence of \(\mathcal D^j(\mathfrak a)\) follows by induction on \(j\), taking for \(e\) the chart at the point.
\end{proof}

\begin{lemma}[Derivatives lower the order by at most one]
\label{hir:derivative-order}
\uses{hir:derivative-ideal-generators,hir:order-basic,hir:cosupport}
\source{wlodarczyk05-hironaka}
\lean{TranscendentalNakamaye.Pending.DerivativeIdeal.chartDeriv_mem_idealPow,TranscendentalNakamaye.Pending.DerivativeIdeal.cosupport_succ_subset_cosupport_derivIdeal}
\leanok
If \(g\in\mathfrak m_x^{q+1}\) and \(e\) is a chart at \(x\), then \(\partial^e_vg\in\mathfrak m_x^q\). Consequently, for \(N=M\),
\(\operatorname{cosupp}(\mathfrak a,\mu+1)\subset\operatorname{cosupp}(\mathcal D(\mathfrak a),\mu)\).
\end{lemma}
\begin{proof}
\leanok
Induction on \(q\); for \(q=0\), \(\partial^e_vg\) is holomorphic. A germ of \(\mathfrak m_x^{q+2}\) is \(\sum_kc_ka_kb_k\) with \(a_k\in\mathfrak m_x^{q+1}\),
\(b_k\in\mathfrak m_x\), and by the product rule its derivative is \(\sum_k(\partial c_k)a_kb_k+c_k(\partial a_k)b_k+c_ka_k\,\partial b_k\); the first and
last terms lie in \(\mathfrak m_x^{q+1}\) because \(a_k\) does, and the middle one because \(\partial a_k\in\mathfrak m_x^q\) by induction and \(b_k\in\mathfrak m_x\)
(Lemma~\ref{hir:order-basic}(1)). For the cosupport: if every germ of \(\mathfrak a_x\) lies in \(\mathfrak m_x^{\mu+1}\), each generator of
\(\mathcal D(\mathfrak a)_x\) agrees near \(x\) with a germ of \(\mathfrak m_x^\mu\) (the germs of \(\mathfrak a_x\) because
\(\mathfrak m_x^{\mu+1}\subset\mathfrak m_x^\mu\), their derivatives by the first part and locality), and \(\mathfrak m_x^\mu\) is an ideal.
\end{proof}

\begin{lemma}[Pulling back powers of the ideal of normal coordinates]
\label{hir:normal-power-pullback}
\uses{mc:normal-power,hir:order-basic}
\lean{TranscendentalNakamaye.Pending.OrderTaylor.mem_idealPow_congr,TranscendentalNakamaye.Pending.OrderTaylor.mem_idealPow_add,TranscendentalNakamaye.Pending.OrderTaylor.mem_idealPow_sum,TranscendentalNakamaye.Pending.OrderTaylor.mul_mem_idealPow_succ,TranscendentalNakamaye.Pending.OrderTaylor.holomorphicAt_comp,TranscendentalNakamaye.Pending.OrderTaylor.mem_idealPow_of_inNormalPow}
\leanok
Let \(\mathcal I\subset\mathcal O_M\) be an ideal of holomorphic germs, \(x\in M\), \(0\le k\le n\), and let \(\psi\colon M\to\mathbf C^n\) be holomorphic
near \(x\), with \(q=\psi(x)\) and with coordinates \(\psi_i\in\mathcal I_x\) for \(i\ge k\) (coordinates indexed from \(0\), as in
Definition~\ref{mc:normal-power}). If \(F\) lies near \(q\) in the \(\mu\)-th power of the ideal of the normal coordinates of \(P_k\), then
\(F\circ\psi\in\mathcal I^\mu_x\).
\end{lemma}
\begin{proof}
\leanok
Induction on \(\mu\). For \(\mu=0\), \(F\) is holomorphic near \(q\), so \(F\circ\psi\) is holomorphic near \(x\). For \(\mu+1\), \(F=\sum_{i\ge k}u_iH_i\)
near \(q\) with each \(H_i\) in the \(\mu\)-th power, so \(F\circ\psi=\sum_{i\ge k}(H_i\circ\psi)\,\psi_i\) near \(x\). Each summand is the product of a
germ of \(\mathcal I^\mu_x\) (induction) and a germ of \(\mathcal I_x\), hence a germ of \(\mathcal I^{\mu+1}_x\), and \(\mathcal I^{\mu+1}_x\) is closed
under sums and only depends on germs.
\end{proof}

\begin{lemma}[Order and derivatives]
\label{hir:order-taylor}
\uses{hir:order-basic,hir:derivative-order,hir:normal-power-pullback,mc:normal-vanishing}
\source{wlodarczyk05-hironaka}
\lean{TranscendentalNakamaye.Pending.OrderTaylor.iterChartDeriv,TranscendentalNakamaye.Pending.OrderTaylor.iterChartDeriv_append,TranscendentalNakamaye.Pending.OrderTaylor.iterChartDeriv_eq_iterDeriv,TranscendentalNakamaye.Pending.OrderTaylor.iterDeriv_comp_add_right,TranscendentalNakamaye.Pending.OrderTaylor.iterChartDeriv_mem_idealPow,TranscendentalNakamaye.Pending.OrderTaylor.iterChartDeriv_eq_zero_of_mem_idealPow,TranscendentalNakamaye.Pending.OrderTaylor.mem_idealPow_of_iterChartDeriv_eq_zero,TranscendentalNakamaye.Pending.OrderTaylor.mem_idealPow_pointIdeal_iff}
\leanok
Let \(e\) be a chart of \(M\) at \(x\), let \(g\) be holomorphic near \(x\), and let \(q\in\mathbb N\). Then \(g\in\mathfrak m_x^q\) if and only if
\((\partial^e_{v_1}\cdots\partial^e_{v_m}g)(x)=0\) for every \(m<q\) and all \(v_1,\dots,v_m\in\mathbf C^n\) (iterated partial derivatives in the chart
\(e\), Definition~\ref{hir:derivative-ideal}).
\end{lemma}
\begin{proof}
\leanok
On the domain of \(e\), \(\partial^e_{v_1}\cdots\partial^e_{v_m}g=(\partial_{v_1}\cdots\partial_{v_m}G)\circ e\), where \(G=g\circ e^{-1}\) and \(\partial_v\) is the
directional derivative on \(\mathbf C^n\) (induction on \(m\); \(e\circ e^{-1}\) is the identity near \(e(p)\)).

If \(g\in\mathfrak m_x^q\) and \(m<q\), then \(m\) applications of Lemma~\ref{hir:derivative-order} give
\(\partial^e_{v_1}\cdots\partial^e_{v_m}g\in\mathfrak m_x^{q-m}\) with \(q-m\ge1\), so it vanishes at \(x\) (Lemma~\ref{hir:order-basic}(3)).

Conversely, let \(c=e(x)\) and \(F(u)=G(u+c)\), holomorphic near \(0\). Directional derivatives commute with translations, so every iterated
directional derivative of order \(<q\) of \(F\) vanishes at \(0\), which is the only point of \(P_0=\{0\}\). By Lemma~\ref{mc:normal-vanishing} with
\(k=0\), \(F\) lies near \(0\) in the \(q\)-th power of the ideal generated by all coordinates. The coordinates of \(\psi=e-c\) are holomorphic near \(x\)
and vanish at \(x\), so they lie in \(\mathfrak m_x\), and \(g=F\circ\psi\) near \(x\); Lemma~\ref{hir:normal-power-pullback} gives \(g\in\mathfrak m_x^q\).
\end{proof}

\begin{lemma}[Giraud--Villamayor: the cosupport through derivatives]
\label{hir:giraud-support}
\uses{hir:derivative-ideal,hir:derivative-order,hir:order-taylor,hir:cosupport-basic,hir:derivative-ideal-generators}
\source{wlodarczyk05-hironaka,bierstone-milman08-functoriality}
\lean{TranscendentalNakamaye.Pending.OrderTaylor.mem_idealPow_of_vanishesToOrderOn_univ,TranscendentalNakamaye.Pending.OrderTaylor.cosupport_derivIdeal_subset_cosupport_succ,TranscendentalNakamaye.Pending.OrderTaylor.cosupport_succ_eq_cosupport_derivIdeal,TranscendentalNakamaye.Pending.OrderTaylor.cosupport_add_eq_cosupport_iterDerivIdeal,TranscendentalNakamaye.Pending.OrderTaylor.cosupport_succ_eq_zeroLocus_iterDerivIdeal,TranscendentalNakamaye.Pending.OrderTaylor.isClosed_zeroLocus_of_isCoherent,TranscendentalNakamaye.Pending.OrderTaylor.isClosed_cosupport}
\leanok
Let \(N=M\) and let \(\mathfrak a\) be an ideal of holomorphic germs. For \(\mu\ge1\) and \(i\ge0\),
\(\operatorname{cosupp}(\mathfrak a,\mu+i)=\operatorname{cosupp}(\mathcal D^i(\mathfrak a),\mu)\). In particular
\(\operatorname{cosupp}(\mathfrak a,j+1)=V(\mathcal D^j(\mathfrak a))\) for every \(j\ge0\), and if \(\mathfrak a\) is coherent, every
\(\operatorname{cosupp}(\mathfrak a,\mu)\) is closed.
\end{lemma}
In Włodarczyk's notation: \(\operatorname{supp}(\mathfrak a,\mu)=\operatorname{supp}(\mathcal D^i(\mathfrak a),\mu-i)\) for \(i<\mu\). For \(i=\mu\) the right side
would be \(\operatorname{cosupp}(\mathcal D^\mu(\mathfrak a),0)=M\).
\begin{proof}
\leanok
First let \(i=1\): \(\operatorname{cosupp}(\mathfrak a,\mu+1)=\operatorname{cosupp}(\mathcal D(\mathfrak a),\mu)\) for \(\mu\ge1\). The inclusion \(\subset\) is
Lemma~\ref{hir:derivative-order}. Conversely let \(x\in\operatorname{cosupp}(\mathcal D(\mathfrak a),\mu)\), \(f\in\mathfrak a_x\), and \(e\) a chart at \(x\). Both
\(f\) and its partial derivatives \(\partial^e_vf\) are germs of \(\mathcal D(\mathfrak a)_x\), hence of \(\mathfrak m_x^\mu\). By Lemma~\ref{hir:order-taylor},
\(f(x)=0\) (as \(\mu\ge1\)) and every iterated partial derivative of order \(<\mu\) of each \(\partial^e_vf\) vanishes at \(x\), i.e.\ every iterated partial
derivative of \(f\) of order \(m\) with \(1\le m\le\mu\). So all iterated partial derivatives of \(f\) of order \(<\mu+1\) vanish at \(x\), and
\(f\in\mathfrak m_x^{\mu+1}\) by Lemma~\ref{hir:order-taylor}.

For general \(i\), induct on \(i\) for all \(\mu\ge1\): \(\operatorname{cosupp}(\mathfrak a,\mu+i+1)=\operatorname{cosupp}(\mathcal D^i(\mathfrak a),\mu+1)
=\operatorname{cosupp}(\mathcal D(\mathcal D^i(\mathfrak a)),\mu)=\operatorname{cosupp}(\mathcal D^{i+1}(\mathfrak a),\mu)\). With \(\mu=1\) and
Lemma~\ref{hir:cosupport-basic}, \(\operatorname{cosupp}(\mathfrak a,1+j)=\operatorname{cosupp}(\mathcal D^j(\mathfrak a),1)=V(\mathcal D^j(\mathfrak a))\).
If \(\mathfrak a\) is coherent, so is \(\mathcal D^j(\mathfrak a)\) (Lemma~\ref{hir:derivative-ideal-generators}), and the zero locus \(V(\mathfrak b)\) of a
coherent ideal is closed: if \(x\notin V(\mathfrak b)\), some local generator \(s\) of \(\mathfrak b\) near \(x\) has \(s(x)\ne0\); it is a germ of \(\mathfrak b\)
at every point near \(x\) and does not vanish there. Finally \(\operatorname{cosupp}(\mathfrak a,0)=M\) (Lemma~\ref{hir:cosupport-basic}).
\end{proof}

\begin{lemma}[Pulling back powers of a maximal ideal]
\label{hir:point-power-pullback}
\uses{hir:order-basic}
\lean{TranscendentalNakamaye.Pending.HadamardBridge.holomorphicAt_comp_map,TranscendentalNakamaye.Pending.HadamardBridge.comp_mem_idealPow_pointIdeal}
\leanok
Let \(\varphi\colon Y\to M\) be a map of complex manifolds which is holomorphic near \(u\in Y\). If \(f\in\mathfrak m_{\varphi(u)}^\mu\), then
\(f\circ\varphi\in\mathfrak m_u^\mu\).
\end{lemma}
\begin{proof}
\leanok
Induction on \(\mu\). A germ of \(\mathfrak m_{\varphi(u)}^{\mu+1}\) is \(\sum_kc_ka_kb_k\) with \(c_k\) holomorphic, \(a_k\in\mathfrak m^\mu_{\varphi(u)}\) and
\(b_k\in\mathfrak m_{\varphi(u)}\); then \(f\circ\varphi=\sum_k(c_k\circ\varphi)(a_k\circ\varphi)(b_k\circ\varphi)\) near \(u\), with \(a_k\circ\varphi\in\mathfrak m_u^\mu\)
by induction and \(b_k\circ\varphi\) holomorphic near \(u\) and vanishing at \(u\).
\end{proof}

\begin{lemma}[Hadamard bridge: an ideal of order \(\ge\mu\) along a smooth centre lies in \(\mathcal I_C^\mu\)]
\label{hir:hadamard-bridge}
\uses{hir:cosupport,hir:order-taylor,hir:point-power-pullback,hir:normal-power-pullback,mc:normal-vanishing,ana:blowups}
\source{wlodarczyk05-hironaka}
\lean{TranscendentalNakamaye.Pending.HadamardBridge.eventually_mem_of_isCoherent,TranscendentalNakamaye.Pending.HadamardBridge.mem_idealPow_of_one_mem,TranscendentalNakamaye.Pending.HadamardBridge.le_idealPow_vanishingIdeal_of_subset_cosupport}
\leanok
Let \(\mathfrak a\subset\mathcal O_M\) be a coherent ideal, \(\mu\in\mathbb N\), and \(C\) a closed submanifold of \(M\) with
\(C\subset\operatorname{cosupp}(\mathfrak a|_M,\mu)\). Then \(\mathfrak a\subset\mathcal I_C^\mu\).
\end{lemma}
\begin{proof}
\leanok
Let \(y\in M\) and \(f\in\mathfrak a_y\). If \(y\notin C\), then \(1\in\mathcal I_{C,y}\) because \(C\) is closed, and every holomorphic germ at \(y\) lies in
\(\mathcal I_{C,y}^\mu\). Let \(y\in C\), where \(C\) has dimension \(k\), and let \(\psi=\Phi\circ e\) be flattening coordinates at \(y\)
(Definition~\ref{ana:blowups}): \(\psi(y)=0\) and, near \(y\), \(C=\psi^{-1}(P_k)\). Put \(F=f\circ\psi^{-1}\), holomorphic near \(0\). By coherence, \(f\) is
a germ of \(\mathfrak a\) at every point \(z\) near \(y\) (local generators of \(\mathfrak a\) near \(y\) generate it at \(z\)). For \(w\in P_k\) near \(0\), the point
\(z=\psi^{-1}(w)\) lies in \(C\subset\operatorname{cosupp}(\mathfrak a,\mu)\), so \(f\in\mathfrak m_z^\mu\); hence \(F\in\mathfrak m_w^\mu\) on \(\mathbf C^n\)
(Lemma~\ref{hir:point-power-pullback}), and every iterated directional derivative of order \(<\mu\) of \(F\) vanishes at \(w\)
(Lemma~\ref{hir:order-taylor} on \(\mathbf C^n\) with its identity chart). By Lemma~\ref{mc:normal-vanishing}, \(F\) lies near \(0\) in the \(\mu\)-th power
of the ideal of the normal coordinates of \(P_k\). The coordinates \(\psi_i\), \(i\ge k\), vanish on \(C\) near \(y\), so they lie in \(\mathcal I_{C,y}\), and
\(f=F\circ\psi\) near \(y\); Lemma~\ref{hir:normal-power-pullback} gives \(f\in\mathcal I_{C,y}^\mu\).
\end{proof}

\begin{corollary}[Admissible centres and the controlled transform]
\label{hir:bridge-consumers}
\uses{hir:hadamard-bridge,hir:transform,bcalc:controlled-transform}
\lean{TranscendentalNakamaye.Pending.HadamardBridge.ideal_le_idealPow_of_isAdmissibleCentre,TranscendentalNakamaye.Pending.HadamardBridge.pullback_eq_mul_controlledTransform_of_subset_cosupport,TranscendentalNakamaye.Pending.HadamardBridge.isCoherent_controlledTransform_of_subset_cosupport}
\leanok
Let \(\mathfrak a\) be coherent. If \(C\) is an admissible centre of a marked ideal \((M,M,E,\mathfrak a,d)\), then \(\mathfrak a\subset\mathcal I_C^d\). If
\(\pi\colon Y\to M\) is a blowup along \(C\subset\operatorname{cosupp}(\mathfrak a,\mu)\) with exceptional divisor \(D\), then
\(\pi^*\mathfrak a=\mathcal O_Y(-\mu D)\cdot\sigma^c(\mathfrak a,\mu)\) and the controlled transform \(\sigma^c(\mathfrak a,\mu)\) is coherent.
\end{corollary}
\begin{proof}
\leanok
An admissible centre, and the centre of a blowup, is a closed submanifold; Lemma~\ref{hir:hadamard-bridge} supplies the hypothesis
\(\mathfrak a\subset\mathcal I_C^\mu\) of Lemma~\ref{bcalc:controlled-transform}.
\end{proof}

\begin{lemma}[Hadamard bridge along a submanifold]
\label{hir:hadamard-bridge-along}
\uses{hir:hadamard-bridge,bcalc:nested-flattening}
\source{wlodarczyk05-hironaka,bierstone-milman08-functoriality}
\lean{TranscendentalNakamaye.Pending.HadamardBridgeAlong.planeRetract,TranscendentalNakamaye.Pending.HadamardBridgeAlong.planeRetract_apply,TranscendentalNakamaye.Pending.HadamardBridgeAlong.planeRetract_mem,TranscendentalNakamaye.Pending.HadamardBridgeAlong.planeRetract_of_mem,TranscendentalNakamaye.Pending.HadamardBridgeAlong.planeRetract_zero,TranscendentalNakamaye.Pending.HadamardBridgeAlong.restrictLe_idealPow_vanishingIdeal_of_subset_cosupport}
\leanok
Let \(\mathfrak a\) be a coherent ideal on \(M\), \(N\) a closed submanifold and \(C\subset N\) a closed submanifold with
\(C\subset\operatorname{cosupp}(\mathfrak a|_N,\mu)\). Then \(\mathfrak a|_N\subset\mathcal I_C^\mu|_N\): every germ of \(\mathfrak a\) at a point of \(N\) agrees on \(N\)
with a germ of \(\mathcal I_C^\mu\).
\end{lemma}
\begin{proof}
\leanok
Off \(C\) the ideal \(\mathcal I_C\) contains \(1\). At \(x\in C\) choose coordinates \(\psi\) centred at \(x\) with \(C=\psi^{-1}(\mathbb C^{k_C}\times0)\) and
\(N=\psi^{-1}(\mathbb C^{k_N}\times0)\) (Lemma~\ref{bcalc:nested-flattening}), and let \(P\) be the linear retraction of \(\mathbb C^n\) onto \(\mathbb C^{k_N}\times0\).
For \(g\in\mathfrak a_x\) put \(F=g\circ\psi^{-1}\circ P\); then \(\hat g=F\circ\psi\) agrees with \(g\) on \(N\) near \(x\). Let \(w\in\mathbb C^{k_C}\times0\) be near
\(0\) and \(z=\psi^{-1}(w)\in C\). By coherence \(g\in\mathfrak a_z\), so \(g\) agrees on \(N\) near \(z\) with some \(h\in\mathfrak m_z^\mu\). The map \(\psi^{-1}\circ P\)
takes a neighbourhood of \(w\) into \(N\), near \(z\), and sends \(w\) to \(z\); hence \(F=h\circ\psi^{-1}\circ P\) near \(w\), which lies in \(\mathfrak m_w^\mu\)
(Lemma~\ref{hir:point-power-pullback}). So every derivative of order \(<\mu\) of \(F\) vanishes on \(\mathbb C^{k_C}\times0\) near \(0\)
(Lemma~\ref{hir:order-taylor}), and the proof of Lemma~\ref{hir:hadamard-bridge} gives \(\hat g\in(\mathcal I_C^\mu)_x\).
\end{proof}

\begin{corollary}[Coherence of the corrected transform]
\label{hir:geo-transform-coherent}
\uses{hir:hadamard-bridge-along,hir:geo-transform-wellformed}
\lean{TranscendentalNakamaye.Pending.HadamardBridgeAlong.isCoherent_geoTransform_ideal_of_isAdmissibleCentre}
\leanok
If \(N\) is a closed submanifold, \(\mathfrak a\) is coherent and \(C\) is an admissible centre of \((M,N,E,\mathfrak a,d)\), the ideal of the transform
of Definition~\ref{hir:geo-transform} is coherent.
\end{corollary}
\begin{proof}
\leanok
Lemma~\ref{hir:hadamard-bridge-along} gives \(\mathfrak a|_N\subset\mathcal I_C^d|_N\), the hypothesis of Lemma~\ref{hir:geo-transform-wellformed}.
\end{proof}

\begin{lemma}[Derivative ideals of products and sums]
\label{hir:derivative-leibniz}
\uses{hir:derivative-ideal,hir:derivative-ideal-generators}
\source{wlodarczyk05-hironaka,bierstone-milman08-functoriality}
\lean{TranscendentalNakamaye.Pending.IdealOrderArithmetic.mem_mul_iff,TranscendentalNakamaye.Pending.IdealOrderArithmetic.mem_sumIdeal_iff,TranscendentalNakamaye.Pending.IdealOrderArithmetic.mem_derivIdeal_iff,TranscendentalNakamaye.Pending.IdealOrderArithmetic.isInSpanAt_mul_of_mem,TranscendentalNakamaye.Pending.IdealOrderArithmetic.mul_mem_mul,TranscendentalNakamaye.Pending.IdealOrderArithmetic.le_sumIdeal_left',TranscendentalNakamaye.Pending.IdealOrderArithmetic.le_sumIdeal_right',TranscendentalNakamaye.Pending.IdealOrderArithmetic.derivIdeal_mul_le,TranscendentalNakamaye.Pending.IdealOrderArithmetic.derivIdeal_sumIdeal_le,TranscendentalNakamaye.Pending.IdealOrderArithmetic.sumIdeal_le_derivIdeal_sumIdeal}
\leanok
Let \(\mathfrak a,\mathfrak b\) be ideals of holomorphic germs, \(\mathfrak a\mathfrak b\) their product and \(\mathfrak a+\mathfrak b\) their sum (the ideal generated
stalkwise by the union of the stalks). Then \(\mathcal D(\mathfrak a\mathfrak b)\subset\mathcal D(\mathfrak a)\,\mathfrak b+\mathfrak a\,\mathcal D(\mathfrak b)\) and
\(\mathcal D(\mathfrak a+\mathfrak b)=\mathcal D(\mathfrak a)+\mathcal D(\mathfrak b)\).
\end{lemma}
\begin{proof}
\leanok
It suffices to check the generators of the left sides. A germ of \(\mathfrak a\mathfrak b\) lies in \(\mathcal D(\mathfrak a)\,\mathfrak b\) since \(\mathfrak a\subset\mathcal D(\mathfrak a)\). If
\(f=\sum_kc_ka_kb_k\) near \(x\) with \(a_k\in\mathfrak a_x\), \(b_k\in\mathfrak b_x\), then by the product rule
\(\partial^e_vf=\sum_k(\partial^e_vc_k)a_kb_k+c_k(\partial^e_va_k)b_k+c_ka_k(\partial^e_vb_k)\) near \(x\), and the three summands lie in
\(\mathcal D(\mathfrak a)\mathfrak b\), \(\mathcal D(\mathfrak a)\mathfrak b\) and \(\mathfrak a\mathcal D(\mathfrak b)\). For sums, a germ of \(\mathfrak a+\mathfrak b\) lies in
\(\mathcal D(\mathfrak a)+\mathcal D(\mathfrak b)\), and if \(f=\sum_kc_kr_k\) with each \(r_k\) in \(\mathfrak a_x\) or \(\mathfrak b_x\), then
\(\partial^e_vf=\sum_k(\partial^e_vc_k)r_k+c_k\,\partial^e_vr_k\) with \(\partial^e_vr_k\) in \(\mathcal D(\mathfrak a)_x\) or \(\mathcal D(\mathfrak b)_x\). Conversely
\(\mathcal D\) is monotone and \(\mathfrak a,\mathfrak b\subset\mathfrak a+\mathfrak b\).
\end{proof}

\begin{lemma}[Polarization]
\label{hir:polarization}
\lean{TranscendentalNakamaye.Pending.OrderMultiplicativity.update_const_comp_swap,TranscendentalNakamaye.Pending.OrderMultiplicativity.eq_zero_of_symmetric_of_apply_diag}
\leanok
A symmetric continuous \(k\)-linear map \(\Phi\) between complex normed spaces with \(\Phi(y,\dots,y)=0\) for every \(y\) is zero.
\end{lemma}
\begin{proof}
\leanok
Induction on \(k\); \(k=0\) is the hypothesis at \(y=0\). For \(k+1\), the function \(t\mapsto\Phi(y+tw,\dots,y+tw)\) vanishes identically, and
its derivative at \(t=0\) is \(\sum_i\Phi(y,\dots,w,\dots,y)\) (\(w\) in the \(i\)-th slot) \(=(k+1)\,\Phi(w,y,\dots,y)\) by symmetry; so
\(\Phi(w,y,\dots,y)=0\). For fixed \(w\) the \(k\)-linear map \(\Phi(w,\cdot)\) is symmetric and vanishes on the diagonal, hence is zero by
induction.
\end{proof}

\begin{lemma}[The order of a power]
\label{hir:order-power}
\uses{hir:order-taylor,hir:polarization,mc:taylor-tail}
\source{wlodarczyk05-hironaka}
\lean{TranscendentalNakamaye.Pending.OrderMultiplicativity.iterDeriv_ofFn_eq_iteratedFDeriv,TranscendentalNakamaye.Pending.OrderMultiplicativity.iterDeriv_eq_iteratedFDeriv,TranscendentalNakamaye.Pending.OrderMultiplicativity.mem_idealPow_pointIdeal_iff_iteratedFDeriv,TranscendentalNakamaye.Pending.OrderMultiplicativity.pow_mem_idealPow,TranscendentalNakamaye.Pending.OrderMultiplicativity.mem_idealPow_of_pow_mem}
\leanok
Let \(g\) be holomorphic near \(x\), \(\nu\ge1\) and \(\mu\in\mathbb N\). If \(g^\nu\in\mathfrak m_x^{\nu\mu}\), then \(g\in\mathfrak m_x^\mu\). In other words
\(\operatorname{ord}_x(g^\nu)=\nu\operatorname{ord}_x(g)\).
\end{lemma}
\begin{proof}
\leanok
In a chart \(e\) at \(x\) with \(p=e(x)\) and \(G=g\circ e^{-1}\), Lemma~\ref{hir:order-taylor} says \(g\in\mathfrak m_x^q\) if and only if the
iterated derivatives \(D^jG(p)\) vanish for \(j<q\) (iterated directional derivatives are the values of these symmetric multilinear maps).
Suppose \(g\notin\mathfrak m_x^\mu\) and let \(a<\mu\) be least with \(D^aG(p)\ne0\). By Lemma~\ref{hir:polarization} there is \(y\) with
\(D^aG(p)(y,\dots,y)\ne0\). The holomorphic function \(\varphi(t)=G(p+ty)\) of one variable has \(\varphi^{(j)}(0)=D^jG(p)(y,\dots,y)\), so it vanishes
to order exactly \(a\) at \(0\), and \(\varphi^\nu\) to order exactly \(\nu a<\nu\mu\). But \(\varphi^\nu(t)=G^\nu(p+ty)\) and \(g^\nu\in\mathfrak m_x^{\nu\mu}\)
gives \((\varphi^\nu)^{(j)}(0)=D^j(G^\nu)(p)(y,\dots,y)=0\) for \(j<\nu\mu\), a contradiction.
\end{proof}

\begin{lemma}[Cosupports of sums, products and powers]
\label{hir:marked-ideal-operations}
\uses{hir:cosupport-basic,hir:order-basic,hir:order-power,hir:derivative-leibniz}
\source{wlodarczyk05-hironaka}
\lean{TranscendentalNakamaye.Pending.IdealOrderArithmetic.vanishesToOrderOn_sum_mul,TranscendentalNakamaye.Pending.IdealOrderArithmetic.IsInSpanAt.vanishesToOrderOn,TranscendentalNakamaye.Pending.IdealOrderArithmetic.cosupport_sumIdeal,TranscendentalNakamaye.Pending.IdealOrderArithmetic.inter_cosupport_subset_cosupport_mul,TranscendentalNakamaye.Pending.IdealOrderArithmetic.cosupport_subset_cosupport_idealPow,TranscendentalNakamaye.Pending.OrderMultiplicativity.cosupport_idealPow_mul_eq,TranscendentalNakamaye.Pending.OrderMultiplicativity.cosupport_sumIdeal_idealPow}
\leanok
Let \(\mathfrak a,\mathfrak b\) be ideals of holomorphic germs and \(N\subset M\). Then
\(\operatorname{cosupp}(\mathfrak a+\mathfrak b|_N,\mu)=\operatorname{cosupp}(\mathfrak a|_N,\mu)\cap\operatorname{cosupp}(\mathfrak b|_N,\mu)\),
\(\operatorname{cosupp}(\mathfrak a|_N,\mu)\cap\operatorname{cosupp}(\mathfrak b|_N,\nu)\subset\operatorname{cosupp}(\mathfrak a\mathfrak b|_N,\mu+\nu)\) and
\(\operatorname{cosupp}(\mathfrak a|_N,\mu)\subset\operatorname{cosupp}(\mathfrak a^\nu|_N,\mu\nu)\). For \(N=M\) and \(\nu\ge1\),
\(\operatorname{cosupp}(\mathfrak a^\nu,\mu\nu)=\operatorname{cosupp}(\mathfrak a,\mu)\); hence for \(\mu,\nu\ge1\) the sum of marked ideals
\((\mathfrak a,\mu)+(\mathfrak b,\nu)=(\mathfrak a^\nu+\mathfrak b^\mu,\mu\nu)\) has cosupport
\(\operatorname{cosupp}(\mathfrak a,\mu)\cap\operatorname{cosupp}(\mathfrak b,\nu)\).
\end{lemma}
\begin{proof}
\leanok
Germs vanishing to order \(\mu\) along \(N\) at \(x\) are closed under holomorphic combinations, and products add orders
(Lemma~\ref{hir:order-basic}); this gives the first three statements (the first in both directions because
\(\mathfrak a,\mathfrak b\subset\mathfrak a+\mathfrak b\)), the third by induction on \(\nu\). For the reverse inclusion of the fourth, if
\(x\in\operatorname{cosupp}(\mathfrak a^\nu,\mu\nu)\) and \(g\in\mathfrak a_x\), then \(g^\nu\in(\mathfrak a^\nu)_x\subset\mathfrak m_x^{\nu\mu}\), so \(g\in\mathfrak m_x^\mu\) by
Lemma~\ref{hir:order-power}. The last statement combines the first and the fourth.
\end{proof}

\begin{lemma}[Invariance and semicontinuity of the order]
\label{hir:order-invariance}
\uses{hir:cosupport,hir:point-power-pullback,hir:giraud-support,hir:marked-ideal-operations}
\source{wlodarczyk05-hironaka}
\lean{TranscendentalNakamaye.Pending.OrderInvariance.cosupport_pullback_eq,TranscendentalNakamaye.Pending.OrderInvariance.order_pullback_eq,TranscendentalNakamaye.Pending.OrderInvariance.eventually_order_le}
\leanok
Let \(\varphi\colon M'\to M\) be a biholomorphism and \(\mathfrak a\) an ideal of holomorphic germs on \(M\) (\(N=M\)). Then
\(\operatorname{cosupp}(\varphi^*\mathfrak a,\mu)=\varphi^{-1}(\operatorname{cosupp}(\mathfrak a,\mu))\) for every \(\mu\), i.e.
\(\operatorname{ord}_{x'}(\varphi^*\mathfrak a)=\operatorname{ord}_{\varphi(x')}(\mathfrak a)\). If \(\mathfrak a\) is coherent, \(x\mapsto\operatorname{ord}_x\mathfrak a\) is upper
semicontinuous.
\end{lemma}
\begin{proof}
\leanok
If every germ of \(\mathfrak a_{\varphi(x')}\) lies in \(\mathfrak m_{\varphi(x')}^\mu\), their pullbacks lie in \(\mathfrak m_{x'}^\mu\)
(Lemma~\ref{hir:point-power-pullback}), and so does every germ of \(\varphi^*\mathfrak a\) at \(x'\), a holomorphic combination of them. Conversely,
if \(a\in\mathfrak a_{\varphi(x')}\) and \(a\circ\varphi\in\mathfrak m_{x'}^\mu\), then \(a=(a\circ\varphi)\circ\varphi^{-1}\in\mathfrak m_{\varphi(x')}^\mu\) by the same lemma for
\(\varphi^{-1}\). For semicontinuity: if \(\operatorname{ord}_x\mathfrak a=m<\infty\), then \(x\) lies in the open complement of the closed set
\(\operatorname{cosupp}(\mathfrak a,m+1)\) (Lemma~\ref{hir:giraud-support}), so \(\operatorname{ord}_y\mathfrak a\le m\) for \(y\) near \(x\).
\end{proof}

\begin{lemma}[Chain rule in charts]
\label{hir:derivative-chain-rule}
\uses{hir:derivative-ideal-generators}
\lean{TranscendentalNakamaye.Pending.DerivativeIdealInvariance.eventually_differentiableAt_chartMap,TranscendentalNakamaye.Pending.DerivativeIdealInvariance.chainCoeff,TranscendentalNakamaye.Pending.DerivativeIdealInvariance.holomorphicAt_chainCoeff,TranscendentalNakamaye.Pending.DerivativeIdealInvariance.chartDeriv_comp_eventuallyEq_sum}
\leanok
Let \(\varphi\colon M'\to M\) be holomorphic, \(e'\) a chart of \(M'\) at \(y\) and \(e\) a chart of \(M\) at \(\varphi(y)\), and \(f\) holomorphic near
\(\varphi(y)\). Near \(y\),
\[
 \partial_v^{e'}(f\circ\varphi)=\sum_j c_j\,(\partial_j^{e}f)\circ\varphi,\qquad c_j(p)=\bigl(D(e\circ\varphi\circ e'^{-1})(e'(p))\,v\bigr)_j,
\]
and the coefficients \(c_j\) are holomorphic near \(y\).
\end{lemma}
\begin{proof}
\leanok
Near \(e'(p)\), \(f\circ\varphi\circ e'^{-1}=(f\circ e^{-1})\circ(e\circ\varphi\circ e'^{-1})\) with both factors holomorphic; differentiate the composite
and expand \(D(e\circ\varphi\circ e'^{-1})(e'(p))v\) in the standard basis. The coefficients are coordinates of the derivative of a holomorphic map
in a direction, composed with \(e'\).
\end{proof}

\begin{lemma}[Derivative ideals commute with biholomorphisms]
\label{hir:derivative-invariance}
\uses{hir:derivative-chain-rule,hir:derivative-ideal,hir:point-power-pullback}
\source{wlodarczyk05-hironaka,bierstone-milman08-functoriality}
\lean{TranscendentalNakamaye.Pending.DerivativeIdealInvariance.isInSpanAt_comp_of_mem_derivIdeal,TranscendentalNakamaye.Pending.DerivativeIdealInvariance.derivIdeal_pullback_le,TranscendentalNakamaye.Pending.DerivativeIdealInvariance.pullback_derivIdeal_le,TranscendentalNakamaye.Pending.DerivativeIdealInvariance.pullback_derivIdeal,TranscendentalNakamaye.Pending.DerivativeIdealInvariance.pullback_iterDerivIdeal}
\leanok
For a holomorphic map \(\varphi\colon M'\to M\) and an ideal \(\mathfrak a\) of holomorphic germs on \(M\),
\(\mathcal D(\varphi^*\mathfrak a)\subset\varphi^*\mathcal D(\mathfrak a)\). If \(\varphi\) is a biholomorphism, then
\(\varphi^*\mathcal D(\mathfrak a)=\mathcal D(\varphi^*\mathfrak a)\) and \(\varphi^*\mathcal D^j(\mathfrak a)=\mathcal D^j(\varphi^*\mathfrak a)\) for all \(j\).
\end{lemma}
\begin{proof}
\leanok
A germ of \(\varphi^*\mathfrak a\) at \(y\) is \(f=\sum_k b_k\,(a_k\circ\varphi)\) with \(a_k\in\mathfrak a_{\varphi(y)}\). By the product rule and
Lemma~\ref{hir:derivative-chain-rule}, \(\partial_v^{e'}f=\sum_k(\partial_v^{e'}b_k)(a_k\circ\varphi)+\sum_{k,j}b_kc_j\,(\partial_j^ea_k)\circ\varphi\), a combination of
pullbacks of germs of \(\mathcal D(\mathfrak a)\); with \(\mathfrak a\subset\mathcal D(\mathfrak a)\) this gives the first inclusion. For a biholomorphism with inverse
\(\psi\), a generator of \(\mathcal D(\mathfrak a)\) at \(x=\varphi(y)\) is a germ of \(\mathfrak a\) or a derivative \(\partial_w^ef\), \(f\in\mathfrak a_x\). The chain rule for
\(\psi\) applied to \(f\circ\varphi\) (note \((f\circ\varphi)\circ\psi=f\)) gives \(\partial_w^ef=\sum_jc'_j\,(\partial_j^{e'}(f\circ\varphi))\circ\psi\) near \(x\), so
\((\partial_w^ef)\circ\varphi=\sum_j(c'_j\circ\varphi)\,\partial_j^{e'}(f\circ\varphi)\) near \(y\), a combination of generators of \(\mathcal D(\varphi^*\mathfrak a)\). The
iterated statement follows by induction, \(\mathcal D^j(\mathfrak a)\) being an ideal of holomorphic germs.
\end{proof}

\begin{lemma}[Controlled transforms of products, powers and sums]
\label{hir:controlled-operations}
\uses{bcalc:controlled-transform,bcalc:exceptional-divisor}
\source{wlodarczyk05-hironaka}
\lean{TranscendentalNakamaye.Pending.MarkedIdealOperationsTransform.mul_le_mul',TranscendentalNakamaye.Pending.MarkedIdealOperationsTransform.idealPow_mono,TranscendentalNakamaye.Pending.MarkedIdealOperationsTransform.idealPow_idealPow,TranscendentalNakamaye.Pending.MarkedIdealOperationsTransform.mul_sumIdeal,TranscendentalNakamaye.Pending.MarkedIdealOperationsTransform.pullback_sumIdeal,TranscendentalNakamaye.Pending.MarkedIdealOperationsTransform.isCoherent_D,TranscendentalNakamaye.Pending.MarkedIdealOperationsTransform.controlledTransform_mul,TranscendentalNakamaye.Pending.MarkedIdealOperationsTransform.controlledTransform_idealPow,TranscendentalNakamaye.Pending.MarkedIdealOperationsTransform.controlledTransform_sumIdeal,TranscendentalNakamaye.Pending.MarkedIdealOperationsTransform.controlledTransform_sumIdeal_idealPow}
\leanok
Let \(\sigma\) be a blowup along \(C\) and \(\mathfrak a\subset\mathcal I_C^\mu\), \(\mathfrak b\subset\mathcal I_C^\nu\) coherent. Then
\(\sigma^{\mathrm c}(\mathfrak a\mathfrak b,\mu+\nu)=\sigma^{\mathrm c}(\mathfrak a,\mu)\,\sigma^{\mathrm c}(\mathfrak b,\nu)\), \(\sigma^{\mathrm c}(\mathfrak a^k,\mu k)=\sigma^{\mathrm c}(\mathfrak a,\mu)^k\),
\(\sigma^{\mathrm c}(\mathfrak a+\mathfrak b,\mu)=\sigma^{\mathrm c}(\mathfrak a,\mu)+\sigma^{\mathrm c}(\mathfrak b,\mu)\) if also \(\mathfrak b\subset\mathcal I_C^\mu\), and
\(\sigma^{\mathrm c}(\mathfrak a^\nu+\mathfrak b^\mu,\mu\nu)=\sigma^{\mathrm c}(\mathfrak a,\mu)^\nu+\sigma^{\mathrm c}(\mathfrak b,\nu)^\mu\).
\end{lemma}
\begin{proof}
\leanok
Pullback commutes with products and sums of ideals and \(\mathfrak l(\mathfrak a+\mathfrak b)=\mathfrak l\mathfrak a+\mathfrak l\mathfrak b\); so, with \(\mathcal O(-\mu D)\) the
\(\mu\)-th power of the exceptional ideal, \(\sigma^*(\mathfrak a\mathfrak b)=\mathcal O(-(\mu+\nu)D)\,\sigma^{\mathrm c}(\mathfrak a,\mu)\sigma^{\mathrm c}(\mathfrak b,\nu)\),
\(\sigma^*(\mathfrak a^k)=\mathcal O(-\mu kD)\,\sigma^{\mathrm c}(\mathfrak a,\mu)^k\) and
\(\sigma^*(\mathfrak a+\mathfrak b)=\mathcal O(-\mu D)(\sigma^{\mathrm c}(\mathfrak a,\mu)+\sigma^{\mathrm c}(\mathfrak b,\mu))\), and the controlled transform is the unique coherent
ideal with this property (Lemma~\ref{bcalc:controlled-transform}). For the last statement, \(\mathfrak a^\nu\subset\mathcal I_C^{\mu\nu}\) and
\(\mathfrak b^\mu\subset\mathcal I_C^{\mu\nu}\); apply the sum rule and then the power rule.
\end{proof}

\begin{definition}[Sum and product of marked ideals]
\label{hir:marked-sum-product}
\uses{hir:marked-ideal}
\source{wlodarczyk05-hironaka}
\lean{TranscendentalNakamaye.Pending.MarkedIdealOperationsTransform.MarkedIdeal.add,TranscendentalNakamaye.Pending.MarkedIdealOperationsTransform.MarkedIdeal.mul}
\leanok
For marked ideals \(\underline{\mathcal I}_1=(M,N,E,\mathfrak a,\mu)\) and \(\underline{\mathcal I}_2=(M,N,E,\mathfrak b,\nu)\),
\(\underline{\mathcal I}_1+\underline{\mathcal I}_2=(M,N,E,\mathfrak a^\nu+\mathfrak b^\mu,\mu\nu)\) and \(\underline{\mathcal I}_1\cdot\underline{\mathcal I}_2=(M,N,E,\mathfrak a\mathfrak b,\mu+\nu)\).
\end{definition}

\begin{lemma}[Sums and products of marked ideals under multiple test blow-ups]
\label{hir:marked-operations-transform}
\uses{hir:marked-sum-product,hir:controlled-operations,hir:marked-ideal-operations,hir:geo-transform-ambient}
\source{wlodarczyk05-hironaka}
\lean{TranscendentalNakamaye.Pending.MarkedIdealOperationsTransform.MarkedIdeal.ext_of,TranscendentalNakamaye.Pending.MarkedIdealOperationsTransform.MarkedIdeal.cosupp_add,TranscendentalNakamaye.Pending.MarkedIdealOperationsTransform.MarkedIdeal.inter_cosupp_subset_cosupp_mul,TranscendentalNakamaye.Pending.MarkedIdealOperationsTransform.MarkedIdeal.IsAmbientPair,TranscendentalNakamaye.Pending.MarkedIdealOperationsTransform.MarkedIdeal.IsAmbientPair.geoTransform,TranscendentalNakamaye.Pending.MarkedIdealOperationsTransform.MarkedIdeal.isAdmissibleCentre_add_iff,TranscendentalNakamaye.Pending.MarkedIdealOperationsTransform.MarkedIdeal.IsAdmissibleCentre.mul,TranscendentalNakamaye.Pending.MarkedIdealOperationsTransform.MarkedIdeal.geoTransform_add,TranscendentalNakamaye.Pending.MarkedIdealOperationsTransform.MarkedIdeal.geoTransform_mul,TranscendentalNakamaye.Pending.MarkedIdealOperationsTransform.MarkedIdeal.geoTransformAlong_add,TranscendentalNakamaye.Pending.MarkedIdealOperationsTransform.MarkedIdeal.isGeoResolution_add_iff,TranscendentalNakamaye.Pending.MarkedIdealOperationsTransform.MarkedIdeal.geoTransformAlong_mul}
\leanok
Let \(\underline{\mathcal I}_1=(M,M,E,\mathfrak a,\mu)\) and \(\underline{\mathcal I}_2=(M,M,E,\mathfrak b,\nu)\) with \(\mathfrak a,\mathfrak b\) coherent, and use the transform of
Definition~\ref{hir:geo-transform}.
\begin{enumerate}
\item If \(\mu,\nu\ge1\), the multiple test blow-ups of \(\underline{\mathcal I}_1+\underline{\mathcal I}_2\) are exactly the simultaneous multiple test blow-ups of
\(\underline{\mathcal I}_1\) and \(\underline{\mathcal I}_2\); along them the transform of the sum is the sum of the transforms. Hence a tower resolves the sum if
and only if it is a simultaneous multiple test blow-up after which the two cosupports are disjoint.
\item A simultaneous multiple test blow-up of \(\underline{\mathcal I}_1\) and \(\underline{\mathcal I}_2\) is a multiple test blow-up of
\(\underline{\mathcal I}_1\cdot\underline{\mathcal I}_2\), and along it the transform of the product is the product of the transforms.
\end{enumerate}
\end{lemma}
\begin{proof}
\leanok
Induction on the tower. For \(N=M\) the transformed ideals are controlled transforms (Lemma~\ref{hir:geo-transform-ambient}), and an admissible centre
\(C\) of \((M,M,E,\mathfrak a,\mu)\) satisfies \(\mathfrak a\subset\mathcal I_C^\mu\) (Lemma~\ref{hir:bridge-consumers}). A centre is admissible for the sum if and only
if it is admissible for both, since \(\operatorname{cosupp}(\mathfrak a^\nu+\mathfrak b^\mu,\mu\nu)=\operatorname{cosupp}(\mathfrak a,\mu)\cap\operatorname{cosupp}(\mathfrak b,\nu)\)
(Lemma~\ref{hir:marked-ideal-operations}) and both have the same \(E\); a centre admissible for both is admissible for the product, since
\(\operatorname{cosupp}(\mathfrak a,\mu)\cap\operatorname{cosupp}(\mathfrak b,\nu)\subset\operatorname{cosupp}(\mathfrak a\mathfrak b,\mu+\nu)\). At each step
Lemma~\ref{hir:controlled-operations} identifies the transformed ideal; \(N\) stays \(M\), the \(E\) are transformed alike, the marks do not change, and
coherence is preserved.
\end{proof}

\begin{lemma}[Blowups under biholomorphisms of the base]
\label{hir:biholomorphism-blowup}
\uses{ana:blowups}
\lean{TranscendentalNakamaye.Pending.MarkedIdealBiholomorphism.notMem_exceptionalLocus_of_inverse,TranscendentalNakamaye.Pending.MarkedIdealBiholomorphism.isClosedSubmanifold_preimage,TranscendentalNakamaye.Pending.MarkedIdealBiholomorphism.preimage_preimage_eq,TranscendentalNakamaye.Pending.MarkedIdealBiholomorphism.pullback_vanishingIdeal_preimage,TranscendentalNakamaye.Pending.MarkedIdealBiholomorphism.isBlowupAlong_comp_inverse,TranscendentalNakamaye.Pending.MarkedIdealBiholomorphism.isBlowupAlong_empty_of_inverse}
\leanok
Let \(\varphi\colon M'\to M\) be a biholomorphism with inverse \(\psi\). Then \(\varphi\) has no exceptional points, \(\varphi^{-1}(C)\) is a closed
submanifold for every closed submanifold \(C\), and \(g^*\mathcal I_{\varphi^{-1}A}=(\varphi\circ g)^*\mathcal I_A\) for every holomorphic map \(g\) into \(M'\) and every
\(A\subset M\). If \(\sigma\colon Y\to M\) is a blowup along \(C\), then \(\psi\circ\sigma\colon Y\to M'\) is a blowup along \(\varphi^{-1}(C)\); in particular every
biholomorphism is a blowup along the empty set.
\end{lemma}
\begin{proof}
\leanok
The chain rule applied to \(\psi\circ\varphi=\mathrm{id}\) and \(\varphi\circ\psi=\mathrm{id}\) shows that the differentials are invertible; regularity of \(C\) pulls back
along the local biholomorphism \(\varphi\). A generator \(t\circ g\) of \(g^*\mathcal I_{\varphi^{-1}A}\) equals \((t\circ\psi)\circ(\varphi\circ g)\) with
\(t\circ\psi\in\mathcal I_A\), and a generator \(s\circ\varphi\circ g\) of \((\varphi\circ g)^*\mathcal I_A\) is \((s\circ\varphi)\circ g\) with \(s\circ\varphi\in\mathcal I_{\varphi^{-1}A}\). Hence
\((\psi\circ\sigma)^*\mathcal I_{\varphi^{-1}C}=\sigma^*\mathcal I_C\) is invertible, \(\varphi^{-1}(C)\) has empty interior (\(\varphi\) is a homeomorphism), and the universal
property of \(\sigma\) applied to \(\varphi\circ g\) gives that of \(\psi\circ\sigma\): \(g^*\mathcal I_{\varphi^{-1}C}=(\varphi\circ g)^*\mathcal I_C\), and \(\sigma\circ h=\varphi\circ g\) if and only if
\(\psi\circ\sigma\circ h=g\). For the last statement take \(\sigma=\mathrm{id}_M\), a blowup along \(\emptyset\).
\end{proof}

\begin{lemma}[Naturality of the transform under biholomorphisms]
\label{hir:biholomorphism-transform}
\uses{hir:biholomorphism-blowup,hir:geo-transform,hir:order-invariance}
\source{wlodarczyk05-hironaka,bierstone-milman08-functoriality}
\lean{TranscendentalNakamaye.Pending.MarkedIdealBiholomorphism.MarkedIdeal.pullback,TranscendentalNakamaye.Pending.MarkedIdealBiholomorphism.MarkedIdeal.cosupp_pullback,TranscendentalNakamaye.Pending.MarkedIdealBiholomorphism.MarkedIdeal.geoTransform_pullback}
\leanok
Let \(\varphi\colon M'\to M\) be a biholomorphism with inverse \(\psi\) and \(\underline{\mathcal I}=(M,N,E,\mathfrak a,d)\); put
\(\varphi^*\underline{\mathcal I}=(M',\varphi^{-1}N,\varphi^{-1}E,\varphi^*\mathfrak a,d)\). Then \(\operatorname{cosupp}\varphi^*\underline{\mathcal I}=\varphi^{-1}\operatorname{cosupp}\underline{\mathcal I}\), and if
\(\mathfrak a\) is coherent and \(\sigma\colon Y\to M\) is holomorphic, the transform of \(\varphi^*\underline{\mathcal I}\) along \(\psi\circ\sigma\) with centre \(\varphi^{-1}(C)\) equals
the transform of \(\underline{\mathcal I}\) along \(\sigma\) with centre \(C\) (Definition~\ref{hir:geo-transform}).
\end{lemma}
\begin{proof}
\leanok
Orders along \(N\) pull back along holomorphic maps (Lemma~\ref{hir:point-power-pullback} applied to the germ that \(g\) agrees with on \(N\)); apply this to
\(\varphi\) and to \(\psi\), as in Lemma~\ref{hir:order-invariance}. For the transform: \((\psi\circ\sigma)^{-1}(\varphi^{-1}A\setminus\varphi^{-1}C)=\sigma^{-1}(A\setminus C)\) gives the same
\(N'\) and \(E'\); \((\psi\circ\sigma)^*\varphi^*\mathfrak a=\sigma^*(\varphi\circ\psi)^*\mathfrak a=\sigma^*\mathfrak a\) and \((\psi\circ\sigma)^*\mathcal I_{\varphi^{-1}C}=\sigma^*\mathcal I_C\)
(Lemma~\ref{hir:biholomorphism-blowup}) give the same ideal.
\end{proof}

\begin{lemma}[Normal crossings pull back along local biholomorphisms and open inclusions]
\label{hir:snc-invariance}
\uses{hir:snc,ana:local-holomorphic-inverse,rem:open-subset-analytic}
\source{wlodarczyk05-hironaka,bierstone-milman08-functoriality}
\lean{TranscendentalNakamaye.Pending.NormalCrossingsInvariance.IsCoordinateFamilyAt.reindex,TranscendentalNakamaye.Pending.NormalCrossingsInvariance.isCoordinateFamilyAt_of_preimage,TranscendentalNakamaye.Pending.NormalCrossingsInvariance.isCoordinateFamilyAt_preimage_val_iff,TranscendentalNakamaye.Pending.NormalCrossingsInvariance.PullsBackCoordinates,TranscendentalNakamaye.Pending.NormalCrossingsInvariance.pullsBackCoordinates_of_forall_notMem_exceptionalLocus,TranscendentalNakamaye.Pending.NormalCrossingsInvariance.pullsBackCoordinates_val,TranscendentalNakamaye.Pending.NormalCrossingsInvariance.get_map_preimage,TranscendentalNakamaye.Pending.NormalCrossingsInvariance.isSNCDivisor_map_preimage,TranscendentalNakamaye.Pending.NormalCrossingsInvariance.hasNormalCrossings_preimage,TranscendentalNakamaye.Pending.NormalCrossingsInvariance.isTransverse_preimage}
\leanok
Let \(g\colon Z\to Y\) be a holomorphic map between complex manifolds of the same dimension and \(z\in Z\) a point where the
differential of \(g\) is bijective. If \(A'_i\cap O=g^{-1}(A_i)\cap O\) for an open \(O\ni z\) and every \(i\), and the family
\((A_i)\) is adapted at \(g(z)\) to \((S_i)\) (Definition~\ref{hir:snc}), then \((A'_i)\) is adapted at \(z\) to \((S_i)\). For an
open subset \(U\subset Y\) and \(x\in U\), the family \((A_i\cap U)\) is adapted at \(x\) in \(U\) to \((S_i)\) if and only if \((A_i)\) is
adapted at \(x\) in \(Y\). Consequently, if \(g\) is continuous and pulls back adapted families at every point (for instance if its
differential is bijective everywhere, or if \(g\) is the inclusion of an open subset), then for an ordered simple normal crossings
divisor \(E=(H_1,\dots,H_s)\) on \(Y\) the sequence \(g^{-1}E=(g^{-1}H_1,\dots,g^{-1}H_s)\) is an ordered simple normal crossings
divisor; if \(C\) has only normal crossings with \(E\), then \(g^{-1}C\) has only normal crossings with \(g^{-1}E\); and if \(E\) is
transverse to \(N\), then \(g^{-1}E\) is transverse to \(g^{-1}N\).
\end{lemma}
\begin{proof}
\leanok
In charts \(e\) at \(z\) and \(e'\) at \(g(z)\), the map \(F=e'\circ g\circ e^{-1}\) has invertible differential at \(e(z)\), so it restricts
to a biholomorphism \(\Psi\) near \(e(z)\) (Lemma~\ref{ana:local-holomorphic-inverse}), which we shrink so that its source lies in the
chart image of \(O\cap g^{-1}(\operatorname{dom}e')\). If \(\phi\) are the adapted coordinates at \(g(z)\), then \(\phi\circ\Psi\) are holomorphic
coordinates at \(z\) vanishing at \(z\), and for \(y\) near \(z\): \(y\in A'_i\iff g(y)\in A_i\iff\phi(F(e(y)))\in L_{S_i}\); a member through \(z\)
has its image through \(g(z)\). On \(U\) the preferred chart at \(x\) is that of \(Y\) restricted to \(U\)
(Lemma~\ref{rem:open-subset-analytic}): coordinates of \(U\) at \(x\) are coordinates of \(Y\) at \(x\) whose source lies in \(U\), and coordinates
of \(Y\) restrict to coordinates of \(U\) after shrinking their source to the chart image of \(U\). For the families take at \(z\) the index
sets given at \(g(z)\): the conditions \(|S_i|=1\), \(S_i\ne S_j\) and \(S_i\cap S_0=\emptyset\) only concern members through \(z\), whose
images pass through \(g(z)\); the preimages of closed sets are closed.
\end{proof}

\begin{lemma}[Invariance under biholomorphisms]
\label{hir:biholomorphism-invariance}
\uses{hir:snc-invariance,hir:biholomorphism-blowup,hir:biholomorphism-transform,hir:marked-ideal,hir:transform}
\source{wlodarczyk05-hironaka,bierstone-milman08-functoriality}
\lean{TranscendentalNakamaye.Pending.MarkedIdealTowerTransport.notMem_exceptionalLocus,TranscendentalNakamaye.Pending.MarkedIdealTowerTransport.pullsBackCoordinates,TranscendentalNakamaye.Pending.MarkedIdealTowerTransport.symm_preimage_preimage,TranscendentalNakamaye.Pending.MarkedIdealTowerTransport.map_symm_preimage_map_preimage,TranscendentalNakamaye.Pending.MarkedIdealTowerTransport.isClosedSubmanifold_preimage_iff,TranscendentalNakamaye.Pending.MarkedIdealTowerTransport.isSNCDivisor_map_preimage_iff,TranscendentalNakamaye.Pending.MarkedIdealTowerTransport.hasNormalCrossings_preimage_iff,TranscendentalNakamaye.Pending.MarkedIdealTowerTransport.isTransverse_preimage_iff,TranscendentalNakamaye.Pending.MarkedIdealTowerTransport.cosupp_pullback_eq,TranscendentalNakamaye.Pending.MarkedIdealTowerTransport.isAdmissibleCentre_pullback_iff,TranscendentalNakamaye.Pending.MarkedIdealTowerTransport.IsWellFormed.pullback}
\leanok
Let \(\varphi\colon M'\to M\) be a biholomorphism. Then \(\varphi^{-1}C\) is a closed submanifold if and only if \(C\) is; \(\varphi^{-1}E\) is an
ordered simple normal crossings divisor if and only if \(E\) is; \(\varphi^{-1}C\) has only normal crossings with \(\varphi^{-1}E\) if and only if
\(C\) has with \(E\); \(\varphi^{-1}E\) is transverse to \(\varphi^{-1}N\) if and only if \(E\) is transverse to \(N\). For a marked ideal
\(\underline{\mathcal I}\) with coherent ideal, \(\varphi^{-1}C\) is an admissible centre of \(\varphi^*\underline{\mathcal I}\) (Lemma~\ref{hir:biholomorphism-transform})
if and only if \(C\) is an admissible centre of \(\underline{\mathcal I}\), and \(\varphi^*\underline{\mathcal I}\) is well formed if \(\underline{\mathcal I}\) is.
\end{lemma}
\begin{proof}
\leanok
\(\varphi\) and \(\varphi^{-1}\) have bijective differentials (Lemma~\ref{hir:biholomorphism-blowup}), so Lemma~\ref{hir:snc-invariance} applies to both;
the converse directions follow by applying the forward ones to \(\varphi^{-1}\), since \((\varphi^{-1})^{-1}(\varphi^{-1}A)=A\). For admissibility,
\(\operatorname{cosupp}\varphi^*\underline{\mathcal I}=\varphi^{-1}\operatorname{cosupp}\underline{\mathcal I}\) (Lemma~\ref{hir:biholomorphism-transform}) and \(\varphi\) is surjective,
so \(\varphi^{-1}C\subset\varphi^{-1}\operatorname{cosupp}\underline{\mathcal I}\) if and only if \(C\subset\operatorname{cosupp}\underline{\mathcal I}\). The pullback of a coherent ideal along a
holomorphic map is coherent.
\end{proof}

\begin{definition}[Transport of a tower along a biholomorphism]
\label{hir:tower-transport}
\uses{hir:blowup-tower,hir:biholomorphism-blowup}
\source{wlodarczyk05-hironaka}
\lean{TranscendentalNakamaye.Pending.MarkedIdealTowerTransport.TowerTransport,TranscendentalNakamaye.Pending.MarkedIdealTowerTransport.transport,TranscendentalNakamaye.Pending.MarkedIdealTowerTransport.transport_refl_target,TranscendentalNakamaye.Pending.MarkedIdealTowerTransport.transport_refl_tower,TranscendentalNakamaye.Pending.MarkedIdealTowerTransport.transport_refl_back,TranscendentalNakamaye.Pending.MarkedIdealTowerTransport.transport_step_tower,TranscendentalNakamaye.Pending.MarkedIdealTowerTransport.transport_step_back}
\leanok
Let \(\varphi\colon X'\to X\) be a biholomorphism and \(T\colon X=X_0\xleftarrow{\sigma_1}X_1\leftarrow\cdots\xleftarrow{\sigma_r}X_r\) a tower with centres
\(C_0,\dots,C_{r-1}\). Its \emph{transport} \(\varphi^*T\) is the tower over \(X'\) with the same stages \(X_1,\dots,X_r\), the first blowup replaced by
\(\varphi^{-1}\circ\sigma_1\colon X_1\to X'\) with centre \(\varphi^{-1}C_0\), and the other blowups and centres unchanged; for \(r=0\) it is the empty
tower over \(X'\). The stages are related by the biholomorphisms \(\chi_0=\varphi\colon X'\to X_0\) and \(\chi_i=\mathrm{id}_{X_i}\) for \(i\ge1\).
\end{definition}
By Lemma~\ref{hir:biholomorphism-blowup}, \(\varphi^{-1}\circ\sigma_1\) is a blowup along \(\varphi^{-1}C_0\), so \(\varphi^*T\) is a tower.

\begin{lemma}[Multiple test blow-ups and resolutions under biholomorphisms]
\label{hir:tower-transport-resolution}
\uses{hir:tower-transport,hir:biholomorphism-invariance,hir:biholomorphism-transform,hir:geo-transform}
\source{wlodarczyk05-hironaka,bierstone-milman08-functoriality}
\lean{TranscendentalNakamaye.Pending.MarkedIdealTowerTransport.length_transport,TranscendentalNakamaye.Pending.MarkedIdealTowerTransport.proj_transport,TranscendentalNakamaye.Pending.MarkedIdealTowerTransport.back_comp_step,TranscendentalNakamaye.Pending.MarkedIdealTowerTransport.geoTransform_transport,TranscendentalNakamaye.Pending.MarkedIdealTowerTransport.geoTransformAlong_transport_step,TranscendentalNakamaye.Pending.MarkedIdealTowerTransport.cosupp_geoTransformAlong_transport,TranscendentalNakamaye.Pending.MarkedIdealTowerTransport.isAdmissibleCentre_transport_iff,TranscendentalNakamaye.Pending.MarkedIdealTowerTransport.isWellFormed_geoTransformAlong_transport,TranscendentalNakamaye.Pending.MarkedIdealTowerTransport.isGeoTestTower_transport_iff,TranscendentalNakamaye.Pending.MarkedIdealTowerTransport.isGeoResolution_transport_iff}
\leanok
Let \(\varphi\colon X'\to X\) be a biholomorphism, \(T\) a tower over \(X\) and \(\underline{\mathcal I}\) a marked ideal on \(X\) with coherent ideal. The
transport \(\varphi^*T\) has the length of \(T\), and \(\varphi\circ p_{\varphi^*T}=p_T\circ\chi_r\) for the projections. Along \(\varphi^*T\), with the
transform of Definition~\ref{hir:geo-transform}, the transform of \(\varphi^*\underline{\mathcal I}\) at stage \(0\) is \(\varphi^*\underline{\mathcal I}\) and at every stage
\(i\ge1\) it equals the transform of \(\underline{\mathcal I}\) along \(T\). At every stage the cosupports correspond under \(\chi_i\), a centre \(C\) is
admissible for the transform of \(\underline{\mathcal I}\) if and only if \(\chi_i^{-1}C\) is admissible for that of \(\varphi^*\underline{\mathcal I}\), and well-formedness of
the transform of \(\underline{\mathcal I}\) implies that of the transform of \(\varphi^*\underline{\mathcal I}\). In particular \(\varphi^*T\) is a multiple test blow-up
(a resolution) of \(\varphi^*\underline{\mathcal I}\) if and only if \(T\) is one of \(\underline{\mathcal I}\).
\end{lemma}
\begin{proof}
\leanok
Induction on \(T\). The one-step statement is: blowing up \(\chi_{i}^{-1}C\) by \(\chi_i^{-1}\circ\sigma\) after \(\varphi^*T\) gives the transform of \(\underline{\mathcal I}\)
along \(T\) followed by \(\sigma\) along \(C\). For \(T\) empty this is the naturality of the transform (Lemma~\ref{hir:biholomorphism-transform}); for \(i\ge1\),
\(\chi_i=\mathrm{id}\) and it follows from the induction hypothesis, the transform being taken stage by stage. The statements on cosupports, centres
and well-formedness hold at stage \(0\) by Lemmas~\ref{hir:biholomorphism-transform} and~\ref{hir:biholomorphism-invariance} and at later stages because
the transforms are equal; the multiple test blow-up statement follows by induction, and the resolution statement because \(\chi_r\) is bijective,
so \(\chi_r^{-1}\) of the final cosupport is empty if and only if the final cosupport is. For the projections, let \(T'\) be \(T\) without its
last blowup \(\sigma\); the last blowup of \(\varphi^*T\) is \(\chi_{r-1}^{-1}\circ\sigma\), and by induction
\(\varphi\circ p_{\varphi^*T}=\varphi\circ p_{\varphi^*T'}\circ\chi_{r-1}^{-1}\circ\sigma=p_{T'}\circ\chi_{r-1}\circ\chi_{r-1}^{-1}\circ\sigma=p_T\).
\end{proof}

\begin{lemma}[Ideals restricted to an open subset]
\label{hir:open-restriction-ideal}
\uses{ana:restricted-ideals,ana:restricted-colon,bcalc:restrict-open,hir:order}
\source{wlodarczyk05-hironaka,bierstone-milman08-functoriality}
\lean{TranscendentalNakamaye.Pending.MarkedIdealRestrictOpen.mem_pullback_val_iff,TranscendentalNakamaye.Pending.MarkedIdealRestrictOpen.comp_val_mem_pullback_val,TranscendentalNakamaye.Pending.MarkedIdealRestrictOpen.restrictOn_pullback_val_iff,TranscendentalNakamaye.Pending.MarkedIdealRestrictOpen.extendOpens_comp_val_mul,TranscendentalNakamaye.Pending.MarkedIdealRestrictOpen.extendOpens_mul,TranscendentalNakamaye.Pending.MarkedIdealRestrictOpen.colonAlong_pullback_val,TranscendentalNakamaye.Pending.MarkedIdealRestrictOpen.mem_idealPow_pointIdeal_extendOpens}
\leanok
Let \(O\subset Y\) be open with inclusion \(\iota\), and for a function \(f\) on \(O\) let \(\tilde f\) be its extension by zero to \(Y\). For a coherent
ideal \(\mathfrak c\) on \(Y\) and \(x\in O\), \(f\in(\iota^*\mathfrak c)_x\) if and only if \(\tilde f\in\mathfrak c_x\). For a coherent \(\mathfrak b\) and
\(W\subset Y\), \(F\in\widetilde{(\iota^*\mathfrak b)}^{W\cap O}_x\) if and only if \(\tilde F\in\widetilde{\mathfrak b}^W_x\)
(Definition~\ref{ana:restricted-ideals}). If \(\mathfrak b\), \(\mathfrak l\) and the quotient \((\operatorname{Res}_W\mathfrak b:\mathfrak l)\) of
Lemma~\ref{ana:restricted-colon} are coherent, then
\((\operatorname{Res}_{W\cap O}\iota^*\mathfrak b:\iota^*\mathfrak l)=\iota^*(\operatorname{Res}_W\mathfrak b:\mathfrak l)\). If \(h\in\mathfrak m_x^q\) on \(O\), then
\(\tilde h\in\mathfrak m_x^q\) on \(Y\).
\end{lemma}
\begin{proof}
\leanok
A germ of \(\iota^*\mathfrak c\) at \(x\) is a holomorphic combination of restrictions \(g\circ\iota\), \(g\in\mathfrak c_x\); near \(x\) such a combination is
the restriction of the same combination of the \(g\) with the extended coefficients, so \(f\in(\iota^*\mathfrak c)_x\) if and only if \(\tilde f\) is a
holomorphic combination of germs of \(\mathfrak c_x\), which for a coherent ideal means \(\tilde f\in\mathfrak c_x\). The statement for
\(\widetilde{\mathfrak b}^W\) is the same argument with equalities on \(W\) near \(x\). For the quotient: \(f\) lies in the left-hand side if it is
holomorphic and \(\ell f\in\widetilde{(\iota^*\mathfrak b)}^{W\cap O}\) for every \(\ell\in(\iota^*\mathfrak l)_x\); testing with \(\ell=g\circ\iota\),
\(g\in\mathfrak l_x\), and using \(\widetilde{(g\circ\iota)f}=g\tilde f\) gives \(g\tilde f\in\widetilde{\mathfrak b}^W_x\), i.e.\ \(\tilde f\) lies in the quotient;
conversely \(\widetilde{\ell f}=\tilde\ell\tilde f\) with \(\tilde\ell\in\mathfrak l_x\). Then apply the first statement to the coherent quotient. The last
statement is by induction on \(q\): a germ of \(\mathfrak m_x^{q+1}\) is a combination of products \(ab\), \(a\in\mathfrak m_x^q\), \(b\in\mathfrak m_x\), and
\(\tilde a\tilde b\) restricts to \(ab\), with \(\tilde a\in\mathfrak m_x^q\) by induction and \(\tilde b\in\mathfrak m_x\) since \(\tilde b\) is holomorphic near \(x\)
and vanishes at \(x\).
\end{proof}

\begin{lemma}[Marked ideals and transforms restricted to an open subset]
\label{hir:open-restriction-marked}
\uses{hir:open-restriction-ideal,hir:snc-invariance,hir:biholomorphism-transform,hir:geo-transform,hir:point-power-pullback,bcalc:restrict-open}
\source{wlodarczyk05-hironaka,bierstone-milman08-functoriality}
\lean{TranscendentalNakamaye.Pending.MarkedIdealRestrictOpen.cosupp_pullback_val,TranscendentalNakamaye.Pending.MarkedIdealRestrictOpen.IsAdmissibleCentre.pullback_val,TranscendentalNakamaye.Pending.MarkedIdealRestrictOpen.IsWellFormed.pullback_val,TranscendentalNakamaye.Pending.MarkedIdealRestrictOpen.geoStrictTransform_restrict,TranscendentalNakamaye.Pending.MarkedIdealRestrictOpen.geoTransform_pullback_val}
\leanok
Let \(U\subset X\) be open with inclusion \(\iota\), and let \(\underline{\mathcal I}|_U=\iota^*\underline{\mathcal I}=(U,N\cap U,E\cap U,\iota^*\mathfrak a,d)\) be the
restriction of \(\underline{\mathcal I}=(X,N,E,\mathfrak a,d)\) (Lemma~\ref{hir:biholomorphism-transform} for the notation). If \(\mathfrak a\) consists of holomorphic
germs, \(\operatorname{cosupp}\underline{\mathcal I}|_U=U\cap\operatorname{cosupp}\underline{\mathcal I}\). If \(C\) is an admissible centre of \(\underline{\mathcal I}\), then
\(C\cap U\) is one of \(\underline{\mathcal I}|_U\); if \(\underline{\mathcal I}\) is well formed, so is \(\underline{\mathcal I}|_U\). Let \(\sigma\colon Y\to X\) be a blowup along
\(C\), \(V=\sigma^{-1}(U)\) and \(\sigma_U\colon V\to U\) the restriction, a blowup along \(C\cap U\) (Lemma~\ref{bcalc:restrict-open}). Strict transforms
commute with restriction: \(\overline{\sigma_U^{-1}((A\cap U)\setminus C)}=V\cap\overline{\sigma^{-1}(A\setminus C)}\). If \(\mathfrak a\) and the ideal of the
transform of \(\underline{\mathcal I}\) are coherent, the transform of \(\underline{\mathcal I}|_U\) along \(\sigma_U\) with centre \(C\cap U\) is the restriction to
\(V\) of the transform of \(\underline{\mathcal I}\) along \(\sigma\) (Definition~\ref{hir:geo-transform}).
\end{lemma}
\begin{proof}
\leanok
Cosupport: if \(x\in U\) lies in \(\operatorname{cosupp}\underline{\mathcal I}\), every generator \(a\circ\iota\) of \((\iota^*\mathfrak a)_x\) agrees on \(N\cap U\) with
\(h\circ\iota\), where \(a=h\) on \(N\) near \(x\) and \(h\in\mathfrak m_x^d\), and \(h\circ\iota\in\mathfrak m_x^d\) on \(U\) (Lemma~\ref{hir:point-power-pullback}); combinations
of such germs vanish to order \(d\) along \(N\cap U\). Conversely, if \(x\in\operatorname{cosupp}\underline{\mathcal I}|_U\) and \(a\in\mathfrak a_x\), then \(a\circ\iota\) agrees on
\(N\cap U\) near \(x\) with some \(h\in\mathfrak m_x^d\) on \(U\), so \(a\) agrees on \(N\) near \(x\) with \(\tilde h\in\mathfrak m_x^d\)
(Lemma~\ref{hir:open-restriction-ideal}). Admissibility and well-formedness: \(C\cap U\) is a closed submanifold of \(U\) (Lemma~\ref{bcalc:restrict-open}),
lies in \(U\cap\operatorname{cosupp}\underline{\mathcal I}\), and normal crossings, simple normal crossings and transversality restrict
(Lemma~\ref{hir:snc-invariance}); \(\iota^*\mathfrak a\) is coherent. Strict transforms: \(\sigma_U^{-1}((A\cap U)\setminus(C\cap U))=V\cap\sigma^{-1}(A\setminus C)\) and the
closure in the open set \(V\) of \(V\cap S\) is \(V\cap\overline S\). Transform: \(N'\) and the members of \(E'\) restrict by the previous statement, and
\(\sigma_U^{-1}(C\cap U)=V\cap\sigma^{-1}(C)\). For the ideal, \(\iota\circ\sigma_U=\sigma\circ\iota_V\) gives \(\sigma_U^*\iota^*\mathfrak a=\iota_V^*\sigma^*\mathfrak a\),
\(\sigma_U^*\mathcal I_{C\cap U}=\iota_V^*\sigma^*\mathcal I_C\) (Lemma~\ref{bcalc:restrict-open}) and \((\iota_V^*\mathfrak l)^d=\iota_V^*(\mathfrak l^d)\) for coherent ideals; the
quotient along \(V\cap N'\) of the restricted ideals is the restriction of the quotient along \(N'\) by Lemma~\ref{hir:open-restriction-ideal}.
\end{proof}

\begin{lemma}[Multiple test blow-ups and resolutions restricted to an open subset]
\label{hir:open-restriction-tower}
\uses{hir:open-restriction-marked,bcalc:tower-restrict,hir:geo-transform-coherent,hir:geo-transform-wellformed,hir:test-tower}
\source{wlodarczyk05-hironaka,bierstone-milman08-functoriality}
\lean{TranscendentalNakamaye.Pending.MarkedIdealRestrictOpen.isClosedSubmanifold_N_and_isCoherent_geoTransformAlong,TranscendentalNakamaye.Pending.MarkedIdealRestrictOpen.geoTransformAlong_restrict,TranscendentalNakamaye.Pending.MarkedIdealRestrictOpen.isGeoTestTower_restrict,TranscendentalNakamaye.Pending.MarkedIdealRestrictOpen.isGeoResolution_restrict,TranscendentalNakamaye.Pending.MarkedIdealRestrictOpen.isWellFormed_geoTransformAlong_restrict}
\leanok
Let \(\underline{\mathcal I}=(X,N,E,\mathfrak a,d)\) with \(N\) a closed submanifold and \(\mathfrak a\) coherent, \(T\) a multiple test blow-up of
\(\underline{\mathcal I}\) (Definition~\ref{hir:geo-transform}) and \(U\subset X\) open. Along \(T\) every transform has \(N\) a closed submanifold and coherent ideal.
Along the restriction \(T|_U\) of \(T\) (Lemma~\ref{bcalc:tower-restrict}), the transform of \(\underline{\mathcal I}|_U\) at each stage is the restriction of the
transform of \(\underline{\mathcal I}\) to the open preimage of \(U\); hence \(T|_U\) is a multiple test blow-up of \(\underline{\mathcal I}|_U\), it is a resolution if
\(T\) is one, and well-formed stages restrict to well-formed stages.
\end{lemma}
\begin{proof}
\leanok
Induction on \(T\). At a step along an admissible centre, \(N'\) is a closed submanifold (Lemma~\ref{hir:geo-transform-wellformed}) and the transformed ideal
is coherent (Corollary~\ref{hir:geo-transform-coherent}); so Lemma~\ref{hir:open-restriction-marked} applies at every stage and identifies the transform
of the restriction with the restriction of the transform. The restriction of an admissible centre is admissible, and the cosupport of the restriction is
the trace of the cosupport, so a resolution restricts to a resolution; well-formedness restricts by Lemma~\ref{hir:open-restriction-marked}.
\end{proof}

\begin{lemma}[Local retraction onto a submanifold]
\label{hir:local-retraction}
\uses{fpgeom:complex-analytic-subvariety,hir:hadamard-bridge-along}
\lean{TranscendentalNakamaye.Pending.MarkedIdealOperationsAlong.exists_localRetraction}
\leanok
Let \(N\subset M\) be regular of dimension \(k\) at \(x\in N\). There is a map \(r\colon M\to M\), holomorphic near \(x\), with \(r(x)=x\),
\(r(y)\in N\) for \(y\) near \(x\), and \(r(y)=y\) for \(y\in N\) near \(x\).
\end{lemma}
\begin{proof}
\leanok
Let \(e\) be the chart at \(x\) and \(\phi\) flattening coordinates, \(\phi(e(N))=\mathbf C^k\times0\) near \(\phi(e(x))=0\), and let \(P\) be the linear
retraction of \(\mathbf C^n\) onto \(\mathbf C^k\times0\). Put \(r=e^{-1}\circ\phi^{-1}\circ P\circ\phi\circ e\) near \(x\); since \(P(0)=0\) and \(P\) is continuous,
\(P\phi e(y)\) stays in the target of \(\phi\) for \(y\) near \(x\). Then \(\phi e(r(y))=P\phi e(y)\in\mathbf C^k\times0\), so \(r(y)\in N\); on \(N\),
\(P\) fixes \(\phi e(y)\) and \(r(y)=y\); \(r(x)=x\) because \(P(0)=0\).
\end{proof}

\begin{lemma}[Order of powers along a submanifold]
\label{hir:order-power-along}
\uses{hir:local-retraction,hir:order-power,hir:point-power-pullback,hir:marked-ideal-operations}
\source{wlodarczyk05-hironaka}
\lean{TranscendentalNakamaye.Pending.MarkedIdealOperationsAlong.vanishesToOrderOn_of_pow,TranscendentalNakamaye.Pending.MarkedIdealOperationsAlong.cosupport_idealPow_mul_eq_along,TranscendentalNakamaye.Pending.MarkedIdealOperationsAlong.cosupport_sumIdeal_idealPow_along}
\leanok
Let \(N\) be regular at \(x\in N\), \(g\) holomorphic near \(x\) and \(\nu\ge1\). If \(g^\nu|_N\) vanishes to order \(\nu\mu\) at \(x\), then \(g|_N\)
vanishes to order \(\mu\) at \(x\). Hence, if \(N\) is a submanifold and \(\mathfrak a,\mathfrak b\) are ideals of holomorphic germs,
\(\operatorname{cosupp}(\mathfrak a^\nu|_N,\mu\nu)=\operatorname{cosupp}(\mathfrak a|_N,\mu)\), and for \(\mu,\nu\ge1\)
\(\operatorname{cosupp}((\mathfrak a^\nu+\mathfrak b^\mu)|_N,\mu\nu)=\operatorname{cosupp}(\mathfrak a|_N,\mu)\cap\operatorname{cosupp}(\mathfrak b|_N,\nu)\).
\end{lemma}
\begin{proof}
\leanok
Let \(r\) be a local retraction onto \(N\) (Lemma~\ref{hir:local-retraction}) and \(\hat g=g\circ r\), which agrees with \(g\) on \(N\) near \(x\). If
\(g^\nu=h\) on \(N\) near \(x\) with \(h\in\mathfrak m_x^{\nu\mu}\), then \(\hat g^\nu=h\circ r\) near \(x\) (since \(r\) maps into \(N\)), and
\(h\circ r\in\mathfrak m_x^{\nu\mu}\) (Lemma~\ref{hir:point-power-pullback}, \(r(x)=x\)). So \(\hat g\in\mathfrak m_x^\mu\) by Lemma~\ref{hir:order-power}, and
\(g|_N\) vanishes to order \(\mu\). For the cosupports: \(\supseteq\) is Lemma~\ref{hir:marked-ideal-operations}; conversely, at a point of
\(\operatorname{cosupp}(\mathfrak a^\nu|_N,\mu\nu)\) and for \(g\in\mathfrak a_x\), \(g^\nu\in(\mathfrak a^\nu)_x\) vanishes to order \(\mu\nu\) along \(N\). The last statement
follows with \(\operatorname{cosupp}(\mathfrak a+\mathfrak b|_N,\mu)=\operatorname{cosupp}(\mathfrak a|_N,\mu)\cap\operatorname{cosupp}(\mathfrak b|_N,\mu)\).
\end{proof}

\begin{lemma}[Calculus of restricted inclusions]
\label{hir:restricted-calculus}
\uses{ana:restricted-ideals,hir:cosupport,hir:cosupport-basic}
\lean{TranscendentalNakamaye.Pending.MarkedIdealOperationsAlong.IsInSpanOnAt.mul,TranscendentalNakamaye.Pending.MarkedIdealOperationsAlong.restrictLe_mul,TranscendentalNakamaye.Pending.MarkedIdealOperationsAlong.restrictLe_idealPow,TranscendentalNakamaye.Pending.MarkedIdealOperationsAlong.restrictLe_sumIdeal,TranscendentalNakamaye.Pending.MarkedIdealOperationsAlong.restrictLe_mono,TranscendentalNakamaye.Pending.MarkedIdealOperationsAlong.isInSpanOnAt_of_notMem,TranscendentalNakamaye.Pending.MarkedIdealOperationsAlong.restrictOn_eq_of_restrictLe,TranscendentalNakamaye.Pending.MarkedIdealOperationsAlong.cosupport_subset_of_restrictLe}
\leanok
Let \(W\subset M\). If \(\operatorname{Res}_W\mathfrak a\subset\operatorname{Res}_W\mathfrak b\) and \(\operatorname{Res}_W\mathfrak a'\subset\operatorname{Res}_W\mathfrak b'\)
(Definition~\ref{ana:restricted-ideals}), with \(\mathfrak b,\mathfrak b'\) ideals of holomorphic germs, then
\(\operatorname{Res}_W\mathfrak a\mathfrak a'\subset\operatorname{Res}_W\mathfrak b\mathfrak b'\), \(\operatorname{Res}_W\mathfrak a^k\subset\operatorname{Res}_W\mathfrak b^k\) and
\(\operatorname{Res}_W(\mathfrak a+\mathfrak a')\subset\operatorname{Res}_W(\mathfrak b+\mathfrak b')\); restricted inclusions pass to subsets of \(W\). If \(W\) is closed and
\(\operatorname{Res}_W\mathfrak a=\operatorname{Res}_W\mathfrak b\), then \(\widetilde{\mathfrak a}^W=\widetilde{\mathfrak b}^W\). If
\(\operatorname{Res}_N\mathfrak a\subset\operatorname{Res}_N\mathfrak b\), then \(\operatorname{cosupp}(\mathfrak b|_N,d)\subset\operatorname{cosupp}(\mathfrak a|_N,d)\).
\end{lemma}
\begin{proof}
\leanok
Products: if \(a=\sum c_is_i\) and \(a'=\sum c'_js'_j\) on \(W\) near \(x\), then \(aa'=\sum_{i,j}c_ic'_j\,s_is'_j\) there; a germ of \(\mathfrak a\mathfrak a'\) is a
combination of such products. Powers follow by induction, sums since a germ of \(\mathfrak a+\mathfrak a'\) is a combination of germs of \(\mathfrak a\) and of
\(\mathfrak a'\). Off the closed set \(W\) every germ is the empty combination on \(W\); at points of \(W\) the two extensions agree by the two restricted
inclusions. For the cosupport, a germ of \(\mathfrak a_x\) agrees on \(N\) with a combination of germs of \(\mathfrak b_x\), each vanishing to order \(d\) along \(N\).
\end{proof}

\begin{definition}[Restricted equivalence of marked ideals]
\label{hir:restricted-equivalence}
\uses{hir:marked-ideal,ana:restricted-ideals}
\source{bierstone-milman08-functoriality}
\lean{TranscendentalNakamaye.Pending.MarkedIdealOperationsAlong.MarkedIdeal.RestrictEquiv}
\leanok
Two marked ideals \((M,N,E,\mathfrak a,d)\) and \((M,N,E,\mathfrak b,d)\) with the same \(N\), \(E\) and \(d\) are \emph{restricted-equivalent} if
\(\operatorname{Res}_N\mathfrak a=\operatorname{Res}_N\mathfrak b\) (Definition~\ref{ana:restricted-ideals}), i.e.\ they represent the same ideal of \(\mathcal O_N\).
\end{definition}

\begin{lemma}[Restricted-equivalent marked ideals]
\label{hir:restricted-equivalence-basic}
\uses{hir:restricted-equivalence,hir:restricted-calculus,hir:geo-transform}
\lean{TranscendentalNakamaye.Pending.MarkedIdealOperationsAlong.MarkedIdeal.RestrictEquiv.refl,TranscendentalNakamaye.Pending.MarkedIdealOperationsAlong.MarkedIdeal.RestrictEquiv.cosupp_eq,TranscendentalNakamaye.Pending.MarkedIdealOperationsAlong.MarkedIdeal.RestrictEquiv.isAdmissibleCentre_iff,TranscendentalNakamaye.Pending.MarkedIdealOperationsAlong.MarkedIdeal.RestrictEquiv.geoTransform_eq}
\leanok
A marked ideal with coherent ideal is restricted-equivalent to itself. Restricted-equivalent marked ideals have the same cosupport and the same
admissible centres; if \(N\) is closed and the ideals are coherent, they have the same transform (Definition~\ref{hir:geo-transform}) along every
holomorphic map.
\end{lemma}
\begin{proof}
\leanok
A coherent ideal is restrictedly contained in itself. The cosupport statement is Lemma~\ref{hir:restricted-calculus} in both directions, and admissibility
depends only on the cosupport and \(E\). For the transform, \(N'\subset\sigma^{-1}(N)\) is closed, the restricted equality pulls back to
\(\operatorname{Res}_{N'}\sigma^*\mathfrak a=\operatorname{Res}_{N'}\sigma^*\mathfrak b\), so the extensions along \(N'\) of \(\sigma^*\mathfrak a\) and \(\sigma^*\mathfrak b\) are equal
(Lemma~\ref{hir:restricted-calculus}), and so are their quotients by \((\sigma^*\mathcal I_C)^d\).
\end{proof}

\begin{lemma}[Sums and products of marked ideals on a submanifold: one blowup]
\label{hir:marked-operations-along-step}
\uses{hir:restricted-equivalence,hir:restricted-calculus,hir:order-power-along,hir:marked-sum-product,hir:hadamard-bridge-along,hir:geo-transform-coherent,hir:geo-transform-wellformed,ana:restricted-colon}
\source{wlodarczyk05-hironaka,bierstone-milman08-functoriality}
\lean{TranscendentalNakamaye.Pending.MarkedIdealOperationsAlong.MarkedIdeal.IsSubmanifoldPair,TranscendentalNakamaye.Pending.MarkedIdealOperationsAlong.MarkedIdeal.hol_of_isCoherent,TranscendentalNakamaye.Pending.MarkedIdealOperationsAlong.MarkedIdeal.restrictLe_pullback_geoTransform,TranscendentalNakamaye.Pending.MarkedIdealOperationsAlong.MarkedIdeal.IsAdmissibleCentre.add,TranscendentalNakamaye.Pending.MarkedIdealOperationsAlong.MarkedIdeal.geoTransform_add_restrictEquiv,TranscendentalNakamaye.Pending.MarkedIdealOperationsAlong.MarkedIdeal.geoTransform_mul_restrictEquiv,TranscendentalNakamaye.Pending.MarkedIdealOperationsAlong.MarkedIdeal.cosupp_add_along,TranscendentalNakamaye.Pending.MarkedIdealOperationsAlong.MarkedIdeal.isAdmissibleCentre_add_iff_along,TranscendentalNakamaye.Pending.MarkedIdealOperationsAlong.MarkedIdeal.IsSubmanifoldPair.geoTransform}
\leanok
Let \(\underline{\mathcal I}_1=(M,N,E,\mathfrak a,\mu)\) and \(\underline{\mathcal I}_2=(M,N,E,\mathfrak b,\nu)\) with \(N\) a closed submanifold and \(\mathfrak a,\mathfrak b\)
coherent, and let \(\sigma\) be a blowup along a centre \(C\) admissible for both. Then \(C\) is admissible for \(\underline{\mathcal I}_1+\underline{\mathcal I}_2\) and
for \(\underline{\mathcal I}_1\cdot\underline{\mathcal I}_2\) (Definition~\ref{hir:marked-sum-product}), and the transform of the sum (product) is restricted-equivalent to
the sum (product) of the transforms. If \(\mu,\nu\ge1\), \(\operatorname{cosupp}(\underline{\mathcal I}_1+\underline{\mathcal I}_2)=\operatorname{cosupp}\underline{\mathcal I}_1\cap\operatorname{cosupp}\underline{\mathcal I}_2\),
so a centre is admissible for the sum if and only if it is admissible for both. The transforms are again marked ideals on a common closed
submanifold \(N'\) with the same \(E'\) and coherent ideals.
\end{lemma}
\begin{proof}
\leanok
Write \(\mathfrak l=\sigma^*\mathcal I_C\), invertible along \(N'\) (Lemma~\ref{hir:geo-transform-wellformed}). For an admissible centre the Hadamard bridge along \(N\)
(Lemma~\ref{hir:hadamard-bridge-along}) gives \(\operatorname{Res}_N\mathfrak a\subset\operatorname{Res}_N\mathcal I_C^\mu\), hence
\(\operatorname{Res}_{N'}\sigma^*\mathfrak a\subset\operatorname{Res}_{N'}\mathfrak l^\mu\), and the quotient \(\mathfrak a'\) of Lemma~\ref{ana:restricted-colon} satisfies
\(\operatorname{Res}_{N'}\sigma^*\mathfrak a=\operatorname{Res}_{N'}\mathfrak l^\mu\mathfrak a'\); likewise for \(\mathfrak b\), the sum \(\mathfrak s=\mathfrak a^\nu+\mathfrak b^\mu\) with mark \(\mu\nu\) and the
product with mark \(\mu+\nu\). Then, along \(N'\) (Lemma~\ref{hir:restricted-calculus}),
\(\mathfrak l^{\mu\nu}\mathfrak s'\equiv\sigma^*\mathfrak s=(\sigma^*\mathfrak a)^\nu+(\sigma^*\mathfrak b)^\mu\equiv(\mathfrak l^\mu\mathfrak a')^\nu+(\mathfrak l^\nu\mathfrak b')^\mu=\mathfrak l^{\mu\nu}(\mathfrak a'^\nu+\mathfrak b'^\mu)\),
and \(\mathfrak l^{\mu\nu}\) can be cancelled along \(N'\) since it is invertible along \(N'\); the product is the same computation with
\(\mathfrak l^\mu\mathfrak a'\,\mathfrak l^\nu\mathfrak b'=\mathfrak l^{\mu+\nu}\mathfrak a'\mathfrak b'\). Admissibility for the sum and the product: the cosupports of the summands lie in that of the
sum and of the product (Lemma~\ref{hir:marked-ideal-operations}); the reverse inclusion for the sum is Lemma~\ref{hir:order-power-along}. The last statement is
Lemma~\ref{hir:geo-transform-wellformed} and Corollary~\ref{hir:geo-transform-coherent}.
\end{proof}

\begin{lemma}[Sums and products of marked ideals on a submanifold along multiple test blow-ups]
\label{hir:marked-operations-along}
\uses{hir:marked-operations-along-step,hir:restricted-equivalence-basic,hir:test-tower}
\source{wlodarczyk05-hironaka}
\lean{TranscendentalNakamaye.Pending.MarkedIdealOperationsAlong.MarkedIdeal.geoTransformAlong_add_along,TranscendentalNakamaye.Pending.MarkedIdealOperationsAlong.MarkedIdeal.isGeoResolution_add_iff_along,TranscendentalNakamaye.Pending.MarkedIdealOperationsAlong.MarkedIdeal.geoTransformAlong_mul_along}
\leanok
Let \(\underline{\mathcal I}_1=(M,N,E,\mathfrak a,\mu)\) and \(\underline{\mathcal I}_2=(M,N,E,\mathfrak b,\nu)\) with \(N\) a closed submanifold and \(\mathfrak a,\mathfrak b\)
coherent. If \(\mu,\nu\ge1\), the multiple test blow-ups of \(\underline{\mathcal I}_1+\underline{\mathcal I}_2\) are exactly the simultaneous multiple test blow-ups of
\(\underline{\mathcal I}_1\) and \(\underline{\mathcal I}_2\), and along them the transform of the sum is restricted-equivalent to the sum of the transforms; a tower
resolves the sum if and only if it is a simultaneous multiple test blow-up after which the two cosupports are disjoint. A simultaneous multiple
test blow-up of \(\underline{\mathcal I}_1\) and \(\underline{\mathcal I}_2\) is a multiple test blow-up of \(\underline{\mathcal I}_1\cdot\underline{\mathcal I}_2\), and along it the
transform of the product is restricted-equivalent to the product of the transforms.
\end{lemma}
\begin{proof}
\leanok
Induction on the tower, carrying the restricted equivalence. At a step, restricted-equivalent marked ideals have the same admissible centres and the
same transform (Lemma~\ref{hir:restricted-equivalence-basic}); so the transform of the transformed sum is the transform of the sum of the transforms,
which is restricted-equivalent to the sum of the next transforms by Lemma~\ref{hir:marked-operations-along-step}; the transformed pair is again a pair on a
closed submanifold with coherent ideals. The resolution statement follows since restricted-equivalent marked ideals have the same cosupport.
\end{proof}

\begin{lemma}[Biholomorphisms between open subsets from local inverses]
\label{hir:open-biholomorphism}
\uses{rem:open-subset-analytic,bcalc:restrict-open}
\lean{TranscendentalNakamaye.Pending.MarkedIdealOpenBiholomorphism.contMDiffAt_of_eventually_mdifferentiableAt',TranscendentalNakamaye.Pending.MarkedIdealOpenBiholomorphism.exists_diffeomorph_opens_of_localInverse}
\leanok
A map between complex manifolds that is holomorphic near a point is smooth there. If \(\varphi,\psi\colon M\to M\) are holomorphic near \(x\), \(\varphi(x)=x\),
and \(\psi\circ\varphi=\mathrm{id}\) and \(\varphi\circ\psi=\mathrm{id}\) near \(x\), then there are open neighbourhoods \(U,V\) of \(x\) such that \(\varphi\) restricts to a
biholomorphism \(U\to V\) with inverse \(\psi\).
\end{lemma}
\begin{proof}
\leanok
Smoothness: in a chart of the target the map is a holomorphic map into \(\mathbf C^n\), hence smooth. Let \(O\ni x\) be an open set on which \(\varphi,\psi\) are
holomorphic and mutually inverse, \(U=O\cap\varphi^{-1}(O)\) and \(V=O\cap\psi^{-1}(U)\), both open since \(\varphi,\psi\) are continuous on \(O\), and both contain \(x\)
since \(\varphi(x)=x=\psi(x)\). For \(p\in U\), \(\varphi(p)\in O\) and \(\psi(\varphi(p))=p\in U\), so \(\varphi(U)\subset V\); \(\psi(V)\subset U\) by definition; the
compositions are the identity because both sets lie in \(O\). Maps of open submanifolds with values in an open set are holomorphic as maps into it.
\end{proof}

\begin{lemma}[Stages of multiple test blow-ups are well formed]
\label{hir:wellformed-tower}
\uses{bcalc:transform-wellformed,hir:geo-transform}
\source{bierstone-milman08-functoriality}
\lean{TranscendentalNakamaye.Pending.MarkedIdealOpenBiholomorphism.isWellFormed_geoTransformAlong}
\leanok
Along a multiple test blow-up (Definition~\ref{hir:geo-transform}) of a well-formed marked ideal, every transform is well formed.
\end{lemma}
\begin{proof}
\leanok
Induction on the tower with Lemma~\ref{bcalc:transform-wellformed} at each admissible step.
\end{proof}

\begin{lemma}[Towers along a biholomorphism onto an open subset]
\label{hir:tower-transport-open}
\uses{hir:tower-transport-resolution,hir:open-restriction-tower,hir:wellformed-tower,hir:biholomorphism-invariance,hir:open-restriction-marked}
\source{wlodarczyk05-hironaka}
\lean{TranscendentalNakamaye.Pending.MarkedIdealOpenBiholomorphism.geoTransformAlong_transport_restrict_step,TranscendentalNakamaye.Pending.MarkedIdealOpenBiholomorphism.isGeoTestTower_transport_restrict,TranscendentalNakamaye.Pending.MarkedIdealOpenBiholomorphism.isGeoResolution_transport_restrict,TranscendentalNakamaye.Pending.MarkedIdealOpenBiholomorphism.isWellFormed_geoTransformAlong_transport_restrict}
\leanok
Let \(\underline{\mathcal I}=(X,N,E,\mathfrak a,d)\) with \(N\) a closed submanifold and \(\mathfrak a\) coherent, \(T\) a multiple test blow-up of \(\underline{\mathcal I}\),
\(V\subset X\) open and \(\Phi\colon U'\to V\) a biholomorphism. The tower \(\Phi^*(T|_V)\) over \(U'\) (Definition~\ref{hir:tower-transport} applied to
the restriction of Lemma~\ref{hir:open-restriction-tower}) is a multiple test blow-up of \(\Phi^*(\underline{\mathcal I}|_V)\), a resolution if \(T\) is one,
and from its first blowup on the transforms are the restrictions of those of \(\underline{\mathcal I}\) along \(T\). If \(\underline{\mathcal I}\) is well formed, so are all
the transforms along \(\Phi^*(T|_V)\).
\end{lemma}
\begin{proof}
\leanok
Combine Lemma~\ref{hir:open-restriction-tower} (restriction to \(V\)) with Lemma~\ref{hir:tower-transport-resolution} (transport along \(\Phi\)); the ideal of
\(\underline{\mathcal I}|_V\) is coherent. Well-formedness: \(\Phi^*(\underline{\mathcal I}|_V)\) is well formed (Lemmas~\ref{hir:open-restriction-marked}
and~\ref{hir:biholomorphism-invariance}), and Lemma~\ref{hir:wellformed-tower} applies.
\end{proof}

\begin{definition}[Equivalence of marked ideals]
\label{hir:marked-equivalence}
\uses{hir:geo-transform,hir:marked-ideal}
\source{wlodarczyk05-hironaka}
\lean{TranscendentalNakamaye.Pending.MarkedIdealEquivalence.MarkedIdeal.WEquiv}
\leanok
Two marked ideals \(\underline{\mathcal I},\underline{\mathcal J}\) on \(X\) are \emph{equivalent}, \(\underline{\mathcal I}\simeq\underline{\mathcal J}\), if they have the same ordered
\(E\), the same multiple test blow-ups (Definition~\ref{hir:geo-transform}), and along every multiple test blow-up the cosupports of their transforms
coincide.
\end{definition}

\begin{lemma}[Basic properties of equivalence]
\label{hir:marked-equivalence-basic}
\uses{hir:marked-equivalence,hir:restricted-equivalence-basic,hir:geo-transform-basic}
\source{wlodarczyk05-hironaka}
\lean{TranscendentalNakamaye.Pending.MarkedIdealEquivalence.MarkedIdeal.WEquiv.refl,TranscendentalNakamaye.Pending.MarkedIdealEquivalence.MarkedIdeal.WEquiv.symm,TranscendentalNakamaye.Pending.MarkedIdealEquivalence.MarkedIdeal.WEquiv.trans,TranscendentalNakamaye.Pending.MarkedIdealEquivalence.MarkedIdeal.WEquiv.cosupp_eq,TranscendentalNakamaye.Pending.MarkedIdealEquivalence.MarkedIdeal.WEquiv.isGeoResolution_iff,TranscendentalNakamaye.Pending.MarkedIdealEquivalence.MarkedIdeal.E_geoTransformAlong_eq,TranscendentalNakamaye.Pending.MarkedIdealEquivalence.MarkedIdeal.WEquiv.geoTransformAlong,TranscendentalNakamaye.Pending.MarkedIdealEquivalence.MarkedIdeal.RestrictEquiv.geoTransformAlong_step_eq,TranscendentalNakamaye.Pending.MarkedIdealEquivalence.MarkedIdeal.RestrictEquiv.cosupp_admissible_geoTransformAlong,TranscendentalNakamaye.Pending.MarkedIdealEquivalence.MarkedIdeal.RestrictEquiv.wEquiv}
\leanok
Equivalence is an equivalence relation. Equivalent marked ideals have the same cosupport and the same resolutions, and their transforms along a common
multiple test blow-up are equivalent. Restricted-equivalent marked ideals (Definition~\ref{hir:restricted-equivalence}) with closed \(N\) and coherent ideals
are equivalent; from the first blowup on, their transforms along any tower are equal.
\end{lemma}
\begin{proof}
\leanok
Reflexivity, symmetry and transitivity are immediate. The cosupport is the case of the empty tower, and a resolution is a multiple test blow-up with
empty final cosupport. For the transforms along \(T\): a tower \(S\) over the last stage is a multiple test blow-up of the transform if and only if the
concatenation \(T\cdot S\) is one of the original marked ideal, and the transform along \(S\) of the transform along \(T\) is the transform along
\(T\cdot S\) (Lemma~\ref{hir:geo-transform-basic}); the \(E\) of a transform only depends on \(E\) and the tower. For restricted-equivalent marked ideals,
the first transforms are equal (Lemma~\ref{hir:restricted-equivalence-basic}), hence so are all later ones; at stage \(0\) the cosupports and admissible
centres coincide, so the multiple test blow-ups coincide by induction.
\end{proof}

\begin{lemma}[Sums and equivalence]
\label{hir:marked-sum-associative}
\uses{hir:marked-equivalence-basic,hir:marked-operations-transform,hir:marked-sum-product}
\source{wlodarczyk05-hironaka}
\lean{TranscendentalNakamaye.Pending.MarkedIdealEquivalence.MarkedIdeal.WEquiv.add_left,TranscendentalNakamaye.Pending.MarkedIdealEquivalence.MarkedIdeal.isCoherent_add_of_isAmbientPair,TranscendentalNakamaye.Pending.MarkedIdealEquivalence.MarkedIdeal.add_assoc_wEquiv}
\leanok
Let the marked ideals below have \(N=M\), a common \(E\), coherent ideals and marks \(\ge1\). If \(\underline{\mathcal I}_1\simeq\underline{\mathcal I}'_1\), then
\(\underline{\mathcal I}_1+\underline{\mathcal I}_2\simeq\underline{\mathcal I}'_1+\underline{\mathcal I}_2\); and
\((\underline{\mathcal I}_1+\underline{\mathcal I}_2)+\underline{\mathcal I}_3\simeq\underline{\mathcal I}_1+(\underline{\mathcal I}_2+\underline{\mathcal I}_3)\).
\end{lemma}
\begin{proof}
\leanok
By Lemma~\ref{hir:marked-operations-transform} the multiple test blow-ups of a sum are the simultaneous ones of the summands, and along them the transform
of the sum is the sum of the transforms, whose cosupport is the intersection of the cosupports. Both statements follow: in the first the simultaneous
multiple test blow-ups and the intersections agree by the equivalence, in the second both sides have the simultaneous multiple test blow-ups of the three
summands and the cosupports \((c_1\cap c_2)\cap c_3=c_1\cap(c_2\cap c_3)\).
\end{proof}

\subsection{Orders of products}

\begin{lemma}[The order of a product]
\label{hir:order-product}
\uses{hir:order-basic,hir:order-taylor,hir:polarization,hir:point-power-pullback}
\source{wlodarczyk05-hironaka,bierstone-milman08-functoriality}
\lean{TranscendentalNakamaye.Pending.GermOrderProduct.analyticOnNhd_diag,TranscendentalNakamaye.Pending.GermOrderProduct.exists_apply_diag_ne_zero,TranscendentalNakamaye.Pending.GermOrderProduct.exists_apply_diag_ne_zero_both,TranscendentalNakamaye.Pending.GermOrderProduct.natCast_le_analyticOrderAt_line,TranscendentalNakamaye.Pending.GermOrderProduct.analyticOrderAt_line_eq,TranscendentalNakamaye.Pending.GermOrderProduct.exists_least_ne_zero,TranscendentalNakamaye.Pending.GermOrderProduct.mem_idealPow_of_mul_mem,TranscendentalNakamaye.Pending.GermOrderProduct.mul_notMem_idealPow,TranscendentalNakamaye.Pending.GermOrderProduct.vanishesToOrderOn_of_mem,TranscendentalNakamaye.Pending.GermOrderProduct.mem_cosupport_of_mem_cosupport_mul}
\leanok
Let \(f,g\) be holomorphic near \(x\in M\). Then \(\operatorname{ord}_x(fg)=\operatorname{ord}_xf+\operatorname{ord}_xg\). Consequently:
\begin{enumerate}
\item if \(fg\in\mathfrak m_x^{a+b}\) and \(g\notin\mathfrak m_x^{b+1}\), then \(f\in\mathfrak m_x^{a}\);
\item if \(\mathfrak b,\mathfrak c\) are ideals of holomorphic germs, \(o\le d\), \(\operatorname{ord}_x\mathfrak b\le o\) and
\(\operatorname{ord}_x(\mathfrak b\mathfrak c)\ge d\), then \(\operatorname{ord}_x\mathfrak c\ge d-o\).
\end{enumerate}
\end{lemma}
\begin{proof}
\leanok
The inequality \(\ge\) is Lemma~\ref{hir:order-basic}(1). For \(\le\) we may assume \(a=\operatorname{ord}_xf\) and \(b=\operatorname{ord}_xg\) finite. Let \(e\) be a
chart at \(x\), \(p=e(x)\), \(F=f\circ e^{-1}\) and \(G=g\circ e^{-1}\). By Lemma~\ref{hir:order-taylor}, \(D^jF(p)=0\) for \(j<a\) and \(D^aF(p)\ne0\), and likewise
for \(G\) and \(b\). The symmetric multilinear maps \(D^aF(p)\) and \(D^bG(p)\) are nonzero, so by Lemma~\ref{hir:polarization} the homogeneous polynomials
\(P(y)=D^aF(p)(y,\dots,y)\) and \(Q(y)=D^bG(p)(y,\dots,y)\) do not vanish identically on \(\mathbf C^n\). The set \(\{P\ne0\}\) is open and nonempty, and the
holomorphic function \(Q\) on the connected space \(\mathbf C^n\) does not vanish identically on it (identity theorem); choose \(y\) with \(P(y)Q(y)\ne0\). The
one-variable functions \(\varphi(t)=F(p+ty)\) and \(\psi(t)=G(p+ty)\) satisfy \(\varphi^{(j)}(0)=D^jF(p)(y,\dots,y)\) and \(\psi^{(j)}(0)=D^jG(p)(y,\dots,y)\), so
they have orders exactly \(a\) and \(b\) at \(0\), and \(\varphi\psi\) has order exactly \(a+b\). If \(fg\in\mathfrak m_x^{a+b+1}\), then by
Lemma~\ref{hir:point-power-pullback}, applied to the holomorphic map \(t\mapsto e^{-1}(p+ty)\), the function \(\varphi\psi\) vanishes to order \(a+b+1\) at \(0\), a
contradiction.

(1) If \(f\notin\mathfrak m_x^a\), then \(\operatorname{ord}_xf<a\) and \(\operatorname{ord}_xg\le b\), so \(\operatorname{ord}_x(fg)<a+b\), i.e.\ \(fg\notin\mathfrak m_x^{a+b}\).
(2) Choose \(g\in\mathfrak b_x\) with \(g\notin\mathfrak m_x^{o+1}\). For \(f\in\mathfrak c_x\), \(fg\in(\mathfrak b\mathfrak c)_x\subset\mathfrak m_x^{(d-o)+o}\), and (1)
gives \(f\in\mathfrak m_x^{d-o}\).
\end{proof}

\subsection{The resolution algorithm: monomial part and companion ideal}

The resolution of marked ideals is constructed by induction on \(\dim N\), following the proof of Proposition ``pr:resol'' of Włodarczyk in the
arrangement of Bierstone and Milman (their Sections~5 and~7). In the ambient formalism of Definition~\ref{hir:marked-ideal} the ambient manifold never
changes: the induction on dimension replaces the submanifold \(N\) by a hypersurface of \(N\) (a hypersurface of maximal contact, or an intersection
with members of \(E\)), and the statements of Section~\ref{sec:maximal-contact} are used for marked ideals with \(N\ne M\) in their form along \(N\)
(Lemma~\ref{hir:along-calculus}).
The centres are constructed as maximum loci of an upper semicontinuous invariant (the second of the two constructions of Bierstone and Milman,
their Section~7). Then the constructions made on small open sets (hypersurfaces of maximal contact, intersections with members of \(E\)) are
compared only through restriction to open sets and biholomorphisms between open sets, so no fibre product of towers and no descent of towers
along local biholomorphisms from disjoint unions is needed; finiteness of the resolution over compact sets takes the place of resolving germs
of manifolds at compact sets.

\begin{lemma}[Order along a submanifold in flat coordinates]
\label{hir:order-along-flat}
\uses{hir:order-taylor,hir:point-power-pullback,hir:normal-power-pullback,mc:normal-vanishing,mc:taylor-tail,hir:hadamard-bridge-along,mc:derivative-along,hir:derivative-ideal}
\source{bierstone-milman08-functoriality}
\lean{TranscendentalNakamaye.Pending.GiraudAlong.iterDeriv_eventuallyEq,TranscendentalNakamaye.Pending.GiraudAlong.flatMap,TranscendentalNakamaye.Pending.GiraudAlong.flatInv,TranscendentalNakamaye.Pending.GiraudAlong.flatten,TranscendentalNakamaye.Pending.GiraudAlong.flatField,TranscendentalNakamaye.Pending.GiraudAlong.flatMap_flatInv,TranscendentalNakamaye.Pending.GiraudAlong.flatInv_mem_flatDomain,TranscendentalNakamaye.Pending.GiraudAlong.flatInv_flatMap,TranscendentalNakamaye.Pending.GiraudAlong.flatMap_mem_target,TranscendentalNakamaye.Pending.GiraudAlong.flatInv_zero,TranscendentalNakamaye.Pending.GiraudAlong.zero_mem_target,TranscendentalNakamaye.Pending.GiraudAlong.continuousAt_flatInv,TranscendentalNakamaye.Pending.GiraudAlong.differentiableAt_comp_flatInv,TranscendentalNakamaye.Pending.GiraudAlong.eventually_planeRetract_mem_target,TranscendentalNakamaye.Pending.GiraudAlong.tendsto_flatInv_planeRetract,TranscendentalNakamaye.Pending.GiraudAlong.eventually_differentiableAt_flatten,TranscendentalNakamaye.Pending.GiraudAlong.mem_iff_flatMap_mem,TranscendentalNakamaye.Pending.GiraudAlong.flatInv_mem,TranscendentalNakamaye.Pending.GiraudAlong.iterDeriv_flatten_flatMap_eq_zero,TranscendentalNakamaye.Pending.GiraudAlong.mem_of_flat,TranscendentalNakamaye.Pending.GiraudAlong.iterDeriv_flatten_eq_zero_of_vanishesToOrderOn,TranscendentalNakamaye.Pending.GiraudAlong.vanishesToOrderOn_of_iterDeriv_flatten_eq_zero,TranscendentalNakamaye.Pending.GiraudAlong.vanishesToOrderOn_iff_iterDeriv_flatten,TranscendentalNakamaye.Pending.GiraudAlong.holomorphicAt_flatField,TranscendentalNakamaye.Pending.GiraudAlong.fieldDeriv_flatField,TranscendentalNakamaye.Pending.GiraudAlong.isTangentField_flatField,TranscendentalNakamaye.Pending.GiraudAlong.dirDeriv_flatten_eventuallyEq}
\leanok
Let \(N\subset M\) be regular of dimension \(k\) at \(x\in N\): with \(e\) the chart at \(x\) there is a biholomorphism \(\Phi\) between open subsets of
\(\mathbf C^n\) with \(\Phi(e(x))=0\) such that, near \(x\), \(y\in N\) if and only if \(\psi(y)=\Phi(e(y))\in\mathbf C^k\times0\). Put
\(\varphi=e^{-1}\circ\Phi^{-1}\), let \(P\) be the linear projection of \(\mathbf C^n\) onto \(\mathbf C^k\times0\) that sets the coordinates of index \(\ge k\) to
zero, and for \(f\) holomorphic near \(x\) put \(f^\flat=f\circ\varphi\circ P\), holomorphic near \(0\). For \(u\in\mathbf C^n\) let \(\xi_u\) be the
holomorphic vector field near \(x\) whose coefficients in the chart \(e\) are the coordinates of \(D\Phi^{-1}(\psi(p))\,u\), so that
\(\xi_u(f)(p)=D(f\circ\varphi)(\psi(p))\,u\).
\begin{enumerate}
\item If \(y\in N\) is near \(x\) and \(f|_N\in\mathfrak m_{N,y}^q\), then every iterated derivative of order \(<q\) of \(f^\flat\) vanishes at \(\psi(y)\).
Conversely, if every iterated derivative of order \(<q\) of \(f^\flat\) vanishes at \(0\), then \(f|_N\in\mathfrak m_{N,x}^q\).
\item For \(u\in\mathbf C^k\times0\) the field \(\xi_u\) is tangent to \(N\) (Definition~\ref{mc:derivative-along}) at every point near \(x\).
\item Near \(0\), \(\partial_v(f^\flat)=(\xi_{Pv}f)^\flat\) for every \(v\in\mathbf C^n\).
\end{enumerate}
\end{lemma}
\begin{proof}
\leanok
(1) If \(f=h\) on \(N\) near \(y\) with \(h\in\mathfrak m_y^q\), then \(f^\flat=h\circ\varphi\circ P\) near \(\psi(y)\): since \(\psi(y)\in\mathbf C^k\times0\),
\(P(\psi(y))=\psi(y)\), and \(\varphi\circ P\) maps a neighbourhood of \(\psi(y)\) into \(N\), near \(y\). By Lemma~\ref{hir:point-power-pullback},
\(h\circ\varphi\circ P\in\mathfrak m_{\psi(y)}^q\), and its derivatives of order \(<q\) vanish at \(\psi(y)\) (Lemma~\ref{hir:order-taylor} in the identity chart
of \(\mathbf C^n\)). Conversely, if the derivatives of order \(<q\) of \(f^\flat\) vanish at \(0\), then \(f^\flat\in\mathfrak m_0^q\)
(Lemma~\ref{mc:normal-vanishing} with the plane \(\{0\}\)), hence \(f^\flat\circ\psi\in\mathfrak m_x^q\) (Lemma~\ref{hir:normal-power-pullback}); and
\(f^\flat\circ\psi=f\) on \(N\) near \(x\), because \(P\circ\psi=\psi\) and \(\varphi\circ\psi=\mathrm{id}\) there.
(2) The coefficients of \(\xi_u\) are holomorphic. Let \(g\) vanish on \(N\) near a point \(p\in N\) near \(x\). As in (1), \(g\circ\varphi\circ P\) vanishes near
\(\psi(p)\), so by the chain rule \(D(g\circ\varphi)(\psi(p))\,Pu=0\); since \(Pu=u\), \(\xi_u(g)(p)=0\). So \(\xi_u(g)\) vanishes on \(N\) near \(p\).
(3) By the chain rule \(D(f\circ\varphi\circ P)(w)\,v=D(f\circ\varphi)(Pw)\,Pv=\xi_{Pv}(f)(\varphi(Pw))\), because \(\psi(\varphi(Pw))=Pw\).
\end{proof}

\begin{lemma}[Giraud--Villamayor along a submanifold]
\label{hir:giraud-along}
\uses{hir:order-along-flat,mc:derivative-along,mc:derivative-along-basic,hir:giraud-support,hir:derivative-order,hir:cosupport,hir:cosupport-basic,hir:marked-ideal-operations,hir:hadamard-bridge}
\source{bierstone-milman08-functoriality,wlodarczyk05-hironaka}
\lean{TranscendentalNakamaye.Pending.GiraudAlong.holomorphicAt_sub,TranscendentalNakamaye.Pending.GiraudAlong.zero_mem_idealPow,TranscendentalNakamaye.Pending.GiraudAlong.vanishesToOrderOn_of_mem_vanishingIdeal,TranscendentalNakamaye.Pending.GiraudAlong.fieldDeriv_add_eventuallyEq,TranscendentalNakamaye.Pending.GiraudAlong.orderGE_derivIdealAlong_of_orderGE_succ,TranscendentalNakamaye.Pending.GiraudAlong.cosupport_succ_subset_cosupport_derivIdealAlong,TranscendentalNakamaye.Pending.GiraudAlong.orderGE_succ_of_orderGE_derivIdealAlong,TranscendentalNakamaye.Pending.GiraudAlong.cosupport_derivIdealAlong_subset_cosupport_succ,TranscendentalNakamaye.Pending.GiraudAlong.cosupport_succ_eq_cosupport_derivIdealAlong,TranscendentalNakamaye.Pending.GiraudAlong.holomorphicAt_of_mem_iterDerivIdealAlong,TranscendentalNakamaye.Pending.GiraudAlong.cosupport_add_eq_cosupport_iterDerivIdealAlong,TranscendentalNakamaye.Pending.GiraudAlong.cosupport_succ_eq_inter_zeroLocus_iterDerivIdealAlong,TranscendentalNakamaye.Pending.GiraudAlong.isClosed_cosupport_along}
\leanok
Let \(\mathfrak a\) be an ideal of holomorphic germs and \(N\subset M\).
\begin{enumerate}
\item For every \(x\) and \(\mu\), \(\operatorname{ord}_x(\mathfrak a|_N)\ge\mu+1\) implies \(\operatorname{ord}_x((\mathcal D_N\mathfrak a)|_N)\ge\mu\); hence
\(\operatorname{cosupp}(\mathfrak a|_N,\mu+1)\subset\operatorname{cosupp}((\mathcal D_N\mathfrak a)|_N,\mu)\).
\item If \(N\) is a submanifold at each of its points and \(\mu\ge1\), then
\(\operatorname{cosupp}(\mathfrak a|_N,\mu+i)=\operatorname{cosupp}((\mathcal D_N^i\mathfrak a)|_N,\mu)\) for every \(i\ge0\); in particular
\(\operatorname{cosupp}(\mathfrak a|_N,j+1)=N\cap V(\mathcal D_N^j\mathfrak a)\) for every \(j\ge0\).
\item If \(\mathfrak a\) is coherent and \(N\) is a closed submanifold, every \(\operatorname{cosupp}(\mathfrak a|_N,\mu)\) is closed.
\end{enumerate}
\end{lemma}
\begin{proof}
\leanok
(1) Consider the generators of \(\mathcal D_N(\mathfrak a)_x\) (Definition~\ref{mc:derivative-along}). Germs of \(\mathfrak a_x\) have order \(\ge\mu+1\) along \(N\), and germs
of \((\mathcal I_N)_x\) vanish on \(N\). For \(f\in\mathfrak a_x\) and \(\xi\) tangent to \(N\) at \(x\), write \(f=h\) on \(N\) near \(x\) with
\(h\in\mathfrak m_x^{\mu+1}\); then \(f-h\in(\mathcal I_N)_x\), so \(\xi(f-h)\in(\mathcal I_N)_x\), and \(\xi(f)\) agrees on \(N\) near \(x\) with
\(\xi(h)=\sum_m\xi_m\partial^e_mh\in\mathfrak m_x^\mu\) (Lemma~\ref{hir:derivative-order}). Holomorphic combinations of germs of order \(\ge\mu\) along \(N\) have order
\(\ge\mu\) (Lemma~\ref{hir:marked-ideal-operations}).
(2) It suffices to prove the converse of (1) at points \(x\in N\) for \(\mu\ge1\); the rest follows by induction on \(i\) as for \(N=M\)
(Lemma~\ref{hir:giraud-support}), and the last statement is the case \(\mu=1\) with \(\operatorname{cosupp}(\mathfrak b|_N,1)=N\cap V(\mathfrak b)\)
(Lemma~\ref{hir:cosupport-basic}). Use the notation of Lemma~\ref{hir:order-along-flat} and let \(f\in\mathfrak a_x\); by
Lemma~\ref{hir:order-along-flat}(1) it suffices that the iterated derivatives of order \(\le\mu\) of \(f^\flat\) vanish at \(0\). Those of order \(<\mu\) do,
because \(f\in\mathcal D_N(\mathfrak a)_x\) has order \(\ge\mu\) along \(N\) (for order \(0\) this uses \(\mu\ge1\)). A derivative of order \(\mu\) is
\(\partial_{v_1}\cdots\partial_{v_{\mu-1}}\partial_vf^\flat(0)=\partial_{v_1}\cdots\partial_{v_{\mu-1}}(\xi_{Pv}f)^\flat(0)\) by
Lemma~\ref{hir:order-along-flat}(3); \(\xi_{Pv}\) is tangent to \(N\) at \(x\) (Lemma~\ref{hir:order-along-flat}(2)), so \(\xi_{Pv}f\in\mathcal D_N(\mathfrak a)_x\)
has order \(\ge\mu\) along \(N\), and its derivatives of order \(\mu-1\) vanish at \(0\) by Lemma~\ref{hir:order-along-flat}(1).
(3) Let \(x\notin\operatorname{cosupp}(\mathfrak a|_N,\mu)\). If \(x\notin N\), a neighbourhood of \(x\) misses the closed set \(N\). If \(x\in N\), some
\(f\in\mathfrak a_x\) has order \(<\mu\) along \(N\), so by Lemma~\ref{hir:order-along-flat}(1) some derivative of order \(<\mu\) of \(f^\flat\) is nonzero at
\(0\), hence near \(0\). For \(y\in N\) near \(x\), \(f\in\mathfrak a_y\) by coherence (Lemma~\ref{hir:hadamard-bridge}), \(\psi(y)\) is near \(0\), and
Lemma~\ref{hir:order-along-flat}(1) at \(y\) shows that \(f\) has order \(<\mu\) along \(N\) at \(y\). So a neighbourhood of \(x\) misses the cosupport.
\end{proof}

\begin{lemma}[Derivative ideals along \(N\) under a blowup]
\label{hir:along-transform}
\uses{mc:inclusions,mc:exceptional-derivation,mc:derivative-along,mc:derivative-along-coherent,mc:homogenized-along,mc:maximal-order-along,bcalc:strict-ideal,bcalc:strict-geo,hir:hadamard-bridge-along,hir:geo-transform}
\source{wlodarczyk05-hironaka,bierstone-milman08-functoriality}
\lean{TranscendentalNakamaye.Pending.DerivativeIdealAlongTransform.chartDeriv_eventuallyEq_sum_flat_uniform,TranscendentalNakamaye.Pending.DerivativeIdealAlongTransform.derivationData_of_coordinates_uniform,TranscendentalNakamaye.Pending.DerivativeIdealAlongTransform.IsUniformDerivationGenerator,TranscendentalNakamaye.Pending.DerivativeIdealAlongTransform.isUniformDerivationGenerator_of_notMem,TranscendentalNakamaye.Pending.DerivativeIdealAlongTransform.exists_isUniformDerivationGenerator_of_mem,TranscendentalNakamaye.Pending.DerivativeIdealAlongTransform.exists_isUniformDerivationGenerator,TranscendentalNakamaye.Pending.DerivativeIdealAlongTransform.IsUniformDerivationGenerator.ne_zero,TranscendentalNakamaye.Pending.DerivativeIdealAlongTransform.eventually_eq_on_of_pow_mul,TranscendentalNakamaye.Pending.DerivativeIdealAlongTransform.exists_tangent_lift,TranscendentalNakamaye.Pending.DerivativeIdealAlongTransform.fieldDeriv_pow,TranscendentalNakamaye.Pending.DerivativeIdealAlongTransform.fieldDeriv_const,TranscendentalNakamaye.Pending.DerivativeIdealAlongTransform.fieldDeriv_eventuallyEq_add_sum,TranscendentalNakamaye.Pending.DerivativeIdealAlongTransform.fieldDeriv_mem_idealPow,TranscendentalNakamaye.Pending.DerivativeIdealAlongTransform.restrictLe_refl,TranscendentalNakamaye.Pending.DerivativeIdealAlongTransform.restrictLe_trans,TranscendentalNakamaye.Pending.DerivativeIdealAlongTransform.restrictLe_derivIdealAlong_idealPow,TranscendentalNakamaye.Pending.DerivativeIdealAlongTransform.restrictLe_iterDerivIdealAlong_idealPow,TranscendentalNakamaye.Pending.DerivativeIdealAlongTransform.restrictLe_derivIdealAlong,TranscendentalNakamaye.Pending.DerivativeIdealAlongTransform.restrictLe_colonAlong_derivIdealAlong,TranscendentalNakamaye.Pending.DerivativeIdealAlongTransform.restrictLe_colonAlong_iterDerivIdealAlong,TranscendentalNakamaye.Pending.DerivativeIdealAlongTransform.restrictLe_geoTransform_derivMarkedAlong,TranscendentalNakamaye.Pending.DerivativeIdealAlongTransform.isMaximalOrderAlong_geoTransform}
\leanok
Let \(\underline{\mathcal I}=(M,N,E,\mathfrak a,d)\) be a marked ideal with \(N\) a closed submanifold and \(\mathfrak a\) coherent, \(C\) an admissible centre, \(\sigma\)
the blowup along \(C\), \(N'\) the strict transform of \(N\) and \(\underline{\mathcal I}'=(M',N',E',\mathfrak a',d)\) the transform. For \(j\le d\) let
\(\mathcal D^j_N(\underline{\mathcal I})=(M,N,E,\mathcal D_N^j\mathfrak a,d-j)\). Then
\[
 [\mathcal D^j_N(\underline{\mathcal I})]'\subset\mathcal D^j_{N'}(\underline{\mathcal I}')\quad\text{on }N',
\]
i.e.\ every germ of the ideal of the left side is, on \(N'\), a combination of germs of \(\mathcal D^j_{N'}(\mathfrak a')\). In particular, if
\(\underline{\mathcal I}\) has maximal order along \(N\) (no point of \(N\) has order \(\ge d+1\) along \(N\)), so has \(\underline{\mathcal I}'\). More generally, for
\(N\supset C\) a closed submanifold and \(\mathfrak a|_N\subset\mathcal I_C^{\mu+r}|_N\),
\((\sigma^*\mathcal D^r_N\mathfrak a:\mathcal O(-\mu D))|_{N'}\subset\mathcal D^r_{N'}\bigl((\sigma^*\mathfrak a:\mathcal O(-(\mu+r)D))|_{N'}\bigr)\), the colons taken along
\(N'\).
\end{lemma}
\begin{proof}
\leanok
The ambient proof of Lemma~\ref{mc:inclusions} rests on Lemma~\ref{mc:exceptional-derivation}: near \(z\) there are a generator \(y\) of \(\mathcal O(-D)\) and, for
every vector field \(\xi\) at \(x=\sigma(z)\), a derivation \(D'\) with \(y\,(\xi f)\circ\sigma=D'(f\circ\sigma)\) and \(D'(y)\in(y)\); its construction (flat
coordinates at \(x\), a monomial chart at \(z\)) gives one \(D'\) for all \(f\) at once. Fix \(z\in N'\).

\emph{Density.} \(N'=\overline{\sigma^{-1}(N\setminus C)}\) and \(y\ne0\) off \(D\); so if \(y^q g=y^qG\) on \(N'\) near \(z\) with \(g,G\) continuous, then \(g=G\) on \(N'\)
near \(z\).

\emph{Tangency.} If \(\xi\) is tangent to \(N\) at \(x\), then \(D'\) is tangent to \(N'\) at \(z\). Let \(w\in\mathcal I_{N'}\). By Lemma~\ref{bcalc:strict-ideal},
\(\sigma^*\mathcal I_N=\mathcal O(-D)\mathcal I_{N'}\), so \(yw=\sum c_i(t_i\circ\sigma)\) with \(t_i\in\mathcal I_N\). Then
\(\rho yw+yD'(w)=D'(yw)=\sum D'(c_i)(t_i\circ\sigma)+c_i\,y\,(\xi t_i)\circ\sigma\) vanishes on \(N'\) (\(\sigma(N')\subset N\), \(\xi t_i\in\mathcal I_N\)), so
\(yD'(w)=0\) on \(N'\) and \(D'(w)\in\mathcal I_{N'}\) by density.

\emph{One derivative.} Let \(\mathfrak a|_N\subset\mathcal I_C^{\mu+1}|_N\) and \(\mathfrak c=(\sigma^*\mathfrak a:\mathcal O(-(\mu+1)D))\) along \(N'\). For \(s\in\mathfrak a_x\),
\(s=\sum a_kh_k+t\) with \(h_k\in\mathcal I_C^{\mu+1}\), \(t\in\mathcal I_N\), so \(s\circ\sigma=y^{\mu+1}H+t\circ\sigma\) with \(H\in\mathfrak c_z\). For the generators of
\(\mathcal D_N\mathfrak a\) at \(x\): \(y(s\circ\sigma)=y^{\mu+1}(yH)\) on \(N'\); \(y(t\circ\sigma)=0\) on \(N'\) for \(t\in\mathcal I_N\); and for \(\xi\) tangent to \(N\),
\(y(\xi f)\circ\sigma=D'(y^{\mu+1}H+t\circ\sigma)=y^{\mu+1}((\mu+1)\rho H+D'(H))+y(\xi t)\circ\sigma\), whose last term vanishes on \(N'\); \(D'(H)\) is a generator
of \(\mathcal D_{N'}(\mathfrak c)\) by tangency. Combinations follow. If \(g\) lies in the left side, \(y^\mu g\) is on \(N'\) a combination of pullbacks of
elements of \(\mathcal D_N\mathfrak a\), so \(y^{\mu+1}g=y^{\mu+1}K\) on \(N'\) with \(K\) a combination of germs of \(\mathcal D_{N'}(\mathfrak c)\), and \(g=K\) on \(N'\) by
density.

\emph{Iteration.} A field tangent to \(N\) lowers \(\mathcal I_C^{q+1}|_N\) to \(\mathcal I_C^q|_N\) (\(\xi t\in\mathcal I_N\) for \(t\in\mathcal I_N\)), so
\(\mathcal D^r_N\mathfrak a|_N\subset\mathcal I_C^{\mu+1}|_N\); and \(\mathcal D_{N'}\) respects inclusions on \(N'\) (for \(f=\sum a_kg_k+t\), \(t\in\mathcal I_{N'}\),
\(\xi f=\sum\xi(a_k)g_k+a_k\xi(g_k)+\xi(t)\)). Induction on \(r\) as in Lemma~\ref{mc:inclusions}. For the marked ideals, \(C\subset\operatorname{cosupp}\underline{\mathcal I}\)
gives \(\mathfrak a|_N\subset\mathcal I_C^d|_N\) (Lemma~\ref{hir:hadamard-bridge-along}); take \(r=j\), \(\mu=d-j\).

\emph{Maximal order.} Maximal order along \(N\) means \(N\cap V(\mathcal D^d_N\mathfrak a)=\emptyset\) (Lemma~\ref{mc:maximal-order-along}). For \(z\in N'\) pick
\(u\in\mathcal D^d_N\mathfrak a\) with \(u(\sigma(z))\ne0\); \(u\circ\sigma\) lies in \([\mathcal D^d_N(\underline{\mathcal I})]'\) (mark \(0\)), hence on \(N'\) in
\(\mathcal D^d_{N'}(\mathfrak a')\), so \(z\notin V(\mathcal D^d_{N'}\mathfrak a')\).
\end{proof}

\begin{theorem}[Hypersurfaces of maximal contact inside \(N\) along multiple test blow-ups]
\label{hir:giraud-tower-along}
\uses{hir:along-transform,mc:tangent-direction-along,mc:tangent-direction-along-basic,mc:derivative-along,bcalc:giraud-local-along,bcalc:function-transform,mc:inclusions-tower-along,mc:strict-transform-local,mc:giraud-tower,hir:test-tower,bcalc:strict-geo}
\source{wlodarczyk05-hironaka,bierstone-milman08-functoriality}
\lean{TranscendentalNakamaye.Pending.MaximalContactAlongTower.apply_eq_zero_of_isInSpanOnAt,TranscendentalNakamaye.Pending.MaximalContactAlongTower.strictAlong_subset_N,TranscendentalNakamaye.Pending.MaximalContactAlongTower.exists_contact_along,TranscendentalNakamaye.Pending.MaximalContactAlongTower.cosupp_subset_strictAlong_along,TranscendentalNakamaye.Pending.MaximalContactAlongTower.centre_subset_strictAlong_along}
\leanok
Let \(\underline{\mathcal I}=(M,N,E,\mathfrak a,\mu)\) with \(N\) a closed submanifold, \(\mathfrak a\) coherent and \(\mu\ge1\), and let \(u\) be a tangent direction of
\((\mathfrak a|_N,\mu)\) along \(N\) on an open set \(U\) (Definition~\ref{mc:tangent-direction-along}). Then for every multiple test blow-up \(T\) of
\(\underline{\mathcal I}\),
\[
 \operatorname{cosupp}\underline{\mathcal I}_T\cap T^{-1}(U)\subset(N\cap V(u)\cap U)_T ,
\]
the strict transform along \(T\); in particular every admissible centre of \(\underline{\mathcal I}_T\) lies, over \(U\), in \((N\cap V(u)\cap U)_T\). More
precisely, every \(z\in N_T\) over \(U\) has an open neighbourhood \(O\) over \(U\) and a function \(w\) holomorphic on \(O\) which, at the points of
\(N_T\cap O\), agrees on \(N_T\) with germs of \(\mathcal D^{\mu-1}_{N_T}(\mathfrak a_T)\), has multiplicity one along \(N_T\) at its zeros there, and whose zero
set in \(N_T\cap O\) is \((N\cap V(u)\cap U)_T\cap O\).
\end{theorem}
\begin{proof}
\leanok
The last sentence, by induction on the length of \(T\). For \(T\) empty take \(O=U\), \(w=u\). Let \(C\) be the next centre (admissible for \(\underline{\mathcal I}_T\),
so \(C\subset\operatorname{cosupp}\underline{\mathcal I}_T\subset N_T\)), \(\sigma\) the blowup along \(C\) and \(z'\in N'_T\) over \(U\). Then \(\sigma(z')\in N_T\) (strict transforms
lie in preimages, Lemma~\ref{bcalc:strict-geo}); take \(O\), \(w\) for \(\sigma(z')\). \(N_T\) is a closed submanifold and \(\mathfrak a_T\) is coherent
(Lemma~\ref{mc:inclusions-tower-along}). \(w\) vanishes on \(C\cap O\): at \(x\in C\), \(w\) agrees on \(N_T\) with a combination of germs of
\(\mathcal D^{\mu-1}_{N_T}(\mathfrak a_T)\), and these vanish on \(\operatorname{cosupp}(\mathfrak a_T|_{N_T},\mu)\) (Lemma~\ref{mc:tangent-direction-along-basic}). Hence
\(w\circ\sigma=hw'\) near \(z'\), \(h\) a local equation of the exceptional divisor \(D\) (Lemma~\ref{bcalc:function-transform}). On a small open \(O'\ni z'\), by
Lemma~\ref{bcalc:giraud-local-along} the zeros of \(w'\) on \(N_T'\cap O'\) are the points of the strict transform of \(N_T\cap V(w)\cap O\), where \(w'\) has
multiplicity one along \(N_T'\); since \(N_T\cap V(w)\cap O=(N\cap V(u)\cap U)_T\cap O\) and strict transforms are local over the base
(Lemma~\ref{mc:strict-transform-local}), that strict transform is \((N\cap V(u)\cap U)_{T'}\) over \(O\). Finally let \(x\in N_T'\cap O'\) and \(p\) a germ of
\(\mathcal O(-D)\) at \(x\), so \(p=ch\) near \(x\); then \(pw'=c\,(w\circ\sigma)\) agrees on \(N_T'\subset\sigma^{-1}(N_T)\) with \(c\sum_i(a_i\circ\sigma)(s_i\circ\sigma)\), where
\(w=\sum_ia_is_i\) on \(N_T\) near \(\sigma(x)\) with \(s_i\in\mathcal D^{\mu-1}_{N_T}(\mathfrak a_T)_{\sigma(x)}\). So \(w'\) lies in the transform of
\((M,N_T,E_T,\mathcal D^{\mu-1}_{N_T}\mathfrak a_T,1)\), hence, on \(N_T'\), in \(\mathcal D^{\mu-1}_{N_T'}(\mathfrak a_{T'})\) (Lemma~\ref{hir:along-transform}).
The inclusion: a point \(z\) of \(\operatorname{cosupp}\underline{\mathcal I}_T\) over \(U\) lies in \(N_T\), and \(w(z)\) is a combination of values at \(z\) of germs of
\(\mathcal D^{\mu-1}_{N_T}(\mathfrak a_T)\), all \(0\); so \(z\) is a zero of \(w\) on \(N_T\). This is the analogue along \(N\) of Lemma~\ref{mc:giraud-tower}.
\end{proof}

\begin{lemma}[Calculus along a submanifold]
\label{hir:along-calculus}
\uses{hir:along-transform,hir:giraud-tower-along,mc:inclusions-tower-along,mc:homogenized-tower-along,mc:homogenized-equivalent-along,mc:coefficient-equivalent-along,mc:glueing-tower-pullback-along,mc:restrict-tower-along,mc:coefficient-restriction-tower-along,hir:giraud-along,hir:order-along-flat,mc:derivative-along,mc:derivative-along-basic,hir:local-retraction,hir:giraud-support,mc:inclusions,mc:maximal-order-transform,mc:tangent-direction-exists,mc:giraud-tower,mc:glueing-lemma,mc:coefficient-equivalent,hir:restricted-equivalence-basic,mc:derivative-along-generators,mc:derivative-along-coherent,mc:derivative-along-etale,mc:tangent-direction-along-exists,mc:tangent-direction-along-basic,mc:contact-hypersurface-along,bcalc:chart-retraction,mc:homogenized-along,mc:homogenized-along-basic,mc:homogenized-along-maximal-order,mc:homogenized-along-etale,mc:glueing-lemma-along,mc:coefficient-restriction-along,mc:coefficient-restriction-full-along,mc:glueing-lift-along}
\source{bierstone-milman08-functoriality,wlodarczyk05-hironaka}
\lean{TranscendentalNakamaye.Pending.GiraudAlong.cosupport_add_eq_cosupport_iterDerivIdealAlong,TranscendentalNakamaye.Pending.GiraudAlong.cosupport_succ_eq_inter_zeroLocus_iterDerivIdealAlong,TranscendentalNakamaye.Pending.GiraudAlong.isClosed_cosupport_along,TranscendentalNakamaye.Pending.DerivativeIdealAlongTransform.restrictLe_geoTransform_derivMarkedAlong,TranscendentalNakamaye.Pending.DerivativeIdealAlongTransform.isMaximalOrderAlong_geoTransform,TranscendentalNakamaye.Pending.TangentDirectionAlong.exists_isTangentDirectionAlong,TranscendentalNakamaye.Pending.TangentDirectionAlong.IsTangentDirectionAlong.cosupport_inter_subset,TranscendentalNakamaye.Pending.MaximalContactAlongTower.cosupp_subset_strictAlong_along,TranscendentalNakamaye.Pending.HomogenizedIdealAlongTower.homAlong_wEquiv,TranscendentalNakamaye.Pending.HomogenizedIdealAlongTower.coeffAlong_wEquiv,TranscendentalNakamaye.Pending.HomogenizedIdealAlongTower.coeff_homAlong_wEquiv,TranscendentalNakamaye.Pending.HomogenizedIdealAlong.IsMaximalOrderAlong.coeffMarkedAlong,TranscendentalNakamaye.Pending.HomogenizedIdealAlong.pullback_homMarkedAlong_of_inverse,TranscendentalNakamaye.Pending.HomogenizedIdealAlong.pullback_coeffMarkedAlong_of_inverse,TranscendentalNakamaye.Pending.HomogenizedGlueingAlong.glueing_lemma_along,TranscendentalNakamaye.Pending.HomogenizedGlueingLiftAlong.glueing_lemma_lift_along}
\leanok
Let \(\underline{\mathcal I}=(M,N,E,\mathfrak a,d)\) be well formed, \(d\ge1\). With the derivative ideals \(\mathcal D_N\) along \(N\)
(Definition~\ref{mc:derivative-along}) in place of \(\mathcal D\), the following hold.
\begin{enumerate}
\item (Giraud--Villamayor along \(N\)) \(\operatorname{cosupp}(\mathfrak a|_N,d+i)=\operatorname{cosupp}((\mathcal D_N^i\mathfrak a)|_N,d)\) for \(d\ge1\); in particular
\(\operatorname{cosupp}(\mathfrak a|_N,d)=N\cap V(\mathcal D_N^{d-1}\mathfrak a)\), and it is closed.
\item (Transforms) For an admissible centre \(C\) and the blowup \(\sigma\) along \(C\),
\([\mathcal D^j_N(\underline{\mathcal I})]'\subset\mathcal D^j_{N'}(\underline{\mathcal I}')\) on \(N'\), \(j\le d\); in particular maximal order along \(N\) is preserved.
\item (Tangent directions along \(N\)) If \(\underline{\mathcal I}\) is of maximal order along \(N\) and \(x\in\operatorname{cosupp}\underline{\mathcal I}\), there are an open
\(U\ni x\) and \(u\in\mathcal D_N^{d-1}(\mathfrak a)\) on \(U\) whose restriction to \(N\cap U\) has nonvanishing differential at its zeros; then
\(\operatorname{cosupp}\underline{\mathcal I}\cap U\subset V(u)\), and along every multiple test blow-up the cosupports stay in the strict transforms of
\(V(u)\cap N\cap U\).
\item (Homogenized and coefficient ideals, glueing) The homogenized ideal, the coefficient ideal and the glueing lemma
(Lemmas~\ref{mc:glueing-lemma} and~\ref{hir:homogenized-coefficient}) hold along \(N\), with \(\mathcal D_N\), tangent directions along \(N\), and
biholomorphisms of open sets of \(M\) preserving \(N\).
\end{enumerate}
\end{lemma}
\begin{proof}
\leanok
Item (1) is Lemma~\ref{hir:giraud-along}, proved with the flat coordinates of Lemma~\ref{hir:order-along-flat}. The single-stage part of (3) is
Lemmas~\ref{mc:tangent-direction-along-exists} and~\ref{mc:tangent-direction-along-basic} (a tangent direction along \(N\) exists at a point of
maximal order along \(N\), and the cosupport lies in its zero set) and Lemma~\ref{mc:contact-hypersurface-along} (\(N\cap V(u)\) is a submanifold of
codimension one in \(N\), with vanishing ideal \(\mathcal I_N+(u)\)). Proved inputs for (2) and (4): the derivative ideal along \(N\) is coherent, with explicit
generators in one flat chart (Lemmas~\ref{mc:derivative-along-generators} and~\ref{mc:derivative-along-coherent}), and it commutes with restriction to open sets and
biholomorphisms preserving \(N\) (Lemma~\ref{mc:derivative-along-etale}); the chart retraction \(r'\) of the blowup used below is
Lemma~\ref{bcalc:chart-retraction}. Item (2) is Lemma~\ref{hir:along-transform}. Item (4) at one stage: \(\mathcal D_N\) and \(\mathcal H_N\), \(\mathcal C_N\) under
restriction and biholomorphisms (Lemmas~\ref{mc:derivative-along-etale}, \ref{mc:homogenized-along-etale}), coherence, cosupports and maximal order
(Lemmas~\ref{mc:homogenized-along}, \ref{mc:homogenized-along-basic}, \ref{mc:homogenized-along-maximal-order}), the glueing lemma along \(N\)
(Lemma~\ref{mc:glueing-lemma-along}) and the restricted coefficient ideal (Lemma~\ref{mc:coefficient-restriction-along}). The tower half of (3) is
Theorem~\ref{hir:giraud-tower-along}. The tower statements of (4), on Lemma~\ref{hir:along-transform}: the inclusions along towers
(Lemma~\ref{mc:inclusions-tower-along}), the homogenized and coefficient ideals along towers and their equivalences (Lemmas~\ref{mc:homogenized-tower-along},
\ref{mc:homogenized-equivalent-along}, \ref{mc:coefficient-equivalent-along}), the glueing lemma along towers in pullback form
(Lemma~\ref{mc:glueing-tower-pullback-along}) and the restriction to \(S\subset N\) along towers with the inclusion of le:S (Lemmas~\ref{mc:restrict-tower-along},
\ref{mc:coefficient-restriction-tower-along}); the reverse inclusion of the restricted coefficient ideal along towers (le:S (2) in full along \(N\),
Lemma~\ref{mc:coefficient-restriction-full-along}) and the stage-by-stage glueing lifts along \(N\) (Lemma~\ref{mc:glueing-lift-along}). This proves the lemma.
Alternatively, all statements are local at points of \(N\) and depend only on \(\mathfrak a|_N\) (for (2), on the restricted ideals; restricted-equivalent marked
ideals have the same transforms, Lemma~\ref{hir:restricted-equivalence-basic}). Near \(x\in N\) choose coordinates \((z',z'')\) with
\(N=\{z''=0\}\) and the holomorphic retraction \(r(z',z'')=(z',0)\) (Lemma~\ref{hir:local-retraction}), and put \(\mathfrak a^\flat=r^*(\mathfrak a|_N)\). Then
\((\mathfrak a^\flat)|_N=\mathfrak a|_N\); for \(y\in N\), \(\operatorname{ord}_y(\mathfrak a|_N)=\operatorname{ord}_y\mathfrak a^\flat\) (pullback along the submersion \(r\)); and
\((\mathcal D_N\mathfrak a)|_N=(\mathcal D\mathfrak a^\flat)|_N\), because a vector field tangent to \(N\) acts on \(N\) through its \(z'\)-components and the
\(z''\)-derivatives of \(\mathfrak a^\flat\) vanish. Hence the statements at a single stage in (1), (3) and (4) follow from their versions for \(N=M\) (Lemmas~\ref{hir:giraud-support},
\ref{mc:tangent-direction-exists}, \ref{mc:glueing-lemma}, \ref{mc:coefficient-equivalent}) applied to \(\mathfrak a^\flat\) on the cylinder over \(N\cap U\).
For one blowup, choose the coordinates adapted to the pair \(C\subset N\) (\(C=\{z'_1=\dots=z'_m=0,\ z''=0\}\)); in each monomial chart of the blowup the
strict transform \(N'\) is \(\{w=0\}\) for the transformed normal coordinates \(w\), the pullback of \(\mathfrak a^\flat\) depends only on the remaining chart
coordinates, and so its controlled transform is again of the form \(r'^*(\,\cdot\,|_{N'})\) for the chart retraction \(r'\), and restricts on \(N'\) to the
transform of \(\underline{\mathcal I}\) (the exceptional ideal is invertible on \(N'\)). This gives the one-blowup forms of (3) and (4) from their
versions for \(N=M\) (Lemmas~\ref{mc:giraud-tower} and~\ref{mc:glueing-lemma}); the statements along multiple test blow-ups follow by induction on the
length, renormalizing at each stage. Item (2) and the preservation of maximal order are proved without the retraction, directly on \(N'\): the derivation
\(D'\) with \(y\,(\xi f)\circ\sigma=D'(f\circ\sigma)\) of Lemma~\ref{mc:inclusions} is tangent to \(N'\) when \(\xi\) is tangent to \(N\), and \(y^{\mu+1}\) is cancelled
on \(N'\) by density of \(\sigma^{-1}(N\setminus C)\) (Lemma~\ref{hir:along-transform}).
\end{proof}

\begin{lemma}[Exponents along a submanifold are locally constant]
\label{hir:exponent-locally-constant}
\uses{hir:order-along-flat,hir:hadamard-bridge,hir:order-taylor,mc:normal-vanishing,hir:normal-power-pullback,hir:point-power-pullback,hir:cosupport,hir:cosupport-basic,hir:marked-ideal-operations}
\source{bierstone-milman08-functoriality,wlodarczyk05-hironaka}
\lean{TranscendentalNakamaye.Pending.DivisorExponent.eventually_eq_zero_of_forall_mem_idealPow,TranscendentalNakamaye.Pending.DivisorExponent.exists_notMem_idealPow_of_frequently,TranscendentalNakamaye.Pending.DivisorExponent.eventually_mem_idealPow_vanishingIdeal,TranscendentalNakamaye.Pending.DivisorExponent.idealPow_le_idealPow,TranscendentalNakamaye.Pending.DivisorExponent.mem_idealPow_pointIdeal_of_mem_idealPow_vanishingIdeal,TranscendentalNakamaye.Pending.DivisorExponent.flatten_univ,TranscendentalNakamaye.Pending.DivisorExponent.flat_univ,TranscendentalNakamaye.Pending.DivisorExponent.iterDeriv_flatten_univ_eq_zero,TranscendentalNakamaye.Pending.DivisorExponent.mem_idealPow_vanishingIdeal_of_iterDeriv,TranscendentalNakamaye.Pending.DivisorExponent.mem_idealPow_vanishingIdeal_of_eventually,TranscendentalNakamaye.Pending.DivisorExponent.holomorphicAt_of_mdifferentiableOn_open,TranscendentalNakamaye.Pending.DivisorExponent.IsDivNear,TranscendentalNakamaye.Pending.DivisorExponent.isDivNear_zero,TranscendentalNakamaye.Pending.DivisorExponent.IsDivNear.mono,TranscendentalNakamaye.Pending.DivisorExponent.IsDivNear.eventually,TranscendentalNakamaye.Pending.DivisorExponent.isDivNear_of_notMem,TranscendentalNakamaye.Pending.DivisorExponent.isDivNear_of_generators,TranscendentalNakamaye.Pending.DivisorExponent.isDivNear_iff,TranscendentalNakamaye.Pending.DivisorExponent.mem_idealPow_vanishingIdeal_of_frequently,TranscendentalNakamaye.Pending.DivisorExponent.isDivNear_of_frequently,TranscendentalNakamaye.Pending.DivisorExponent.divExp,TranscendentalNakamaye.Pending.DivisorExponent.IsNonzeroAt,TranscendentalNakamaye.Pending.DivisorExponent.exists_not_isDivNear,TranscendentalNakamaye.Pending.DivisorExponent.isDivNear_divExp,TranscendentalNakamaye.Pending.DivisorExponent.not_isDivNear_divExp_succ,TranscendentalNakamaye.Pending.DivisorExponent.isDivNear_iff_le_divExp,TranscendentalNakamaye.Pending.DivisorExponent.eventually_divExp_eq}
\leanok
Let \(\mathfrak a\) be a coherent ideal of \(\mathcal O_M\) and \(H\subset M\) a submanifold at \(x\in H\). Say that \(\mathfrak a\) is \emph{divisible by
\(\mathcal I_H^e\) near \(x\)} if the points of \(H\) near \(x\) lie in \(\operatorname{cosupp}(\mathfrak a,e)\).
\begin{enumerate}
\item This holds if and only if \(\mathfrak a_x\subset(\mathcal I_H^e)_x\). More precisely, a germ \(f\) at \(x\) lies in \((\mathcal I_H^e)_x\) as soon as
\(f\in\mathfrak m_z^e\) at the points \(z\in H\) near \(x\).
\item A germ lying in \(\mathfrak m_x^q\) for every \(q\) vanishes near \(x\). Hence, if \(\mathfrak a_x\) contains a germ that does not vanish near \(x\),
there is a largest \(e=e_H(x)\) such that \(\mathfrak a\) is divisible by \(\mathcal I_H^e\) near \(x\).
\item (Identity principle along \(H\).) If the points \(y\in H\) near which \(\mathfrak a\) is divisible by \(\mathcal I_H^e\) accumulate at \(x\), then
\(\mathfrak a\) is divisible by \(\mathcal I_H^e\) near \(x\). The same holds for a single germ \(f\) holomorphic at \(x\): if the points \(y\in H\) near
which \(f\in\mathfrak m_z^e\) at every \(z\in H\) accumulate at \(x\), then \(f\in(\mathcal I_H^e)_x\).
\item If \(\mathfrak a_y\) contains a germ not vanishing near \(y\) for every \(y\) near \(x\), then \(e_H\) is constant on \(H\) near \(x\).
\end{enumerate}
\end{lemma}
\begin{proof}
\leanok
Take the flat coordinates \(\psi\), \(\varphi=\psi^{-1}\) of Lemma~\ref{hir:order-along-flat} for \(H\) at \(x\), with \(H=\psi^{-1}(\mathbf C^k\times0)\) near \(x\).
(1) If \(f\in\mathfrak m_z^e\) at the points \(z\in H\) near \(x\), the derivatives of order \(<e\) of \(f\circ\varphi\) vanish at the points of
\(\mathbf C^k\times0\) near \(0\) (Lemmas~\ref{hir:point-power-pullback} and~\ref{hir:order-taylor}), so \(f\circ\varphi\) lies in the \(e\)-th power of the
ideal of the normal coordinates near \(0\) (Lemma~\ref{mc:normal-vanishing}) and \(f\in(\mathcal I_H^e)_x\) (Lemma~\ref{hir:normal-power-pullback}). Applied
to the germs of \(\mathfrak a_x\), which are germs of \(\mathfrak a_z\) for \(z\) near \(x\) (coherence, Lemma~\ref{hir:hadamard-bridge}), this gives one
direction. Conversely, generators of \(\mathfrak a\) near \(x\) lying in \((\mathcal I_H^e)_x\) lie in \((\mathcal I_H^e)_z\subset\mathfrak m_z^e\) at the points
\(z\in H\) near \(x\), and holomorphic combinations keep the order (Lemma~\ref{hir:marked-ideal-operations}).
(2) If \(g\in\mathfrak m_x^q\) for every \(q\), all derivatives of \(g\) in a chart vanish at \(x\) (Lemma~\ref{hir:order-taylor}), so the power series of \(g\) at
\(x\) is zero and \(g\) vanishes near \(x\). If \(g\in\mathfrak a_x\) does not vanish near \(x\), choose \(q\) with \(g\notin\mathfrak m_x^q\); divisibility by
\(\mathcal I_H^{q+1}\) near \(x\) would give \(x\in\operatorname{cosupp}(\mathfrak a,q+1)\), so \(g\in\mathfrak m_x^{q+1}\subset\mathfrak m_x^q\).
(3) Let \(s_1,\dots,s_K\) generate \(\mathfrak a\) near \(x\); by (1) it suffices that each \(s_i\in(\mathcal I_H^e)_x\), which is the statement for a single
germ. Let \(B\) be a ball around \(0\) on which \(S=s_i\circ\varphi\) is holomorphic, \(L\) a list of fewer than \(e\) directions, \(D=\partial_LS\), and \(P\) the projection onto \(\mathbf C^k\times0\). The function
\(D\circ P\) is holomorphic on the convex, hence connected, open set \(P^{-1}(B)\). Choose \(y\in H\) near \(x\), with \(\psi(y)\in B\), near which \(\mathfrak a\) is
divisible by \(\mathcal I_H^e\). For \(u\) near \(\psi(y)\), \(\varphi(Pu)\) is a point of \(H\) near \(y\), so \(s_i\in\mathfrak m^e\) there and
\(D(Pu)=0\) (Lemma~\ref{hir:order-along-flat}(1) with \(N=M\)). By the identity theorem \(D\circ P=0\) on \(P^{-1}(B)\), so \(D=0\) on \((\mathbf C^k\times0)\cap B\), and
\(s_i\in(\mathcal I_H^e)_x\) as in (1).
(4) Let \(e=e_H(x)\). Divisibility by \(\mathcal I_H^e\) near \(x\) is divisibility near every point close to \(x\). Non-divisibility by
\(\mathcal I_H^{e+1}\) near \(x\) persists at the points of \(H\) near \(x\) by (3). So \(e_H(y)=e\) for \(y\in H\) near \(x\).
\end{proof}

\begin{lemma}[Simultaneous divisibility by coordinate hyperplanes]
\label{hir:monomial-factorization}
\uses{hir:exponent-locally-constant,hir:order-along-flat,hir:order-product,hir:order-basic,bcalc:hyperplane-ideal}
\source{bierstone-milman08-functoriality,wlodarczyk05-hironaka}
\lean{TranscendentalNakamaye.Pending.MonomialPart.exists_eq_mul_pow_of_mem_idealPow,TranscendentalNakamaye.Pending.MonomialPart.holomorphicAt_finset_prod,TranscendentalNakamaye.Pending.MonomialPart.exists_eq_mul_pow_of_ne_zero,TranscendentalNakamaye.Pending.MonomialPart.pow_mem_idealPow,TranscendentalNakamaye.Pending.MonomialPart.frequently_flatMap_ne_zero,TranscendentalNakamaye.Pending.MonomialPart.exists_eq_mul_prod_of_forall,TranscendentalNakamaye.Pending.MonomialPart.eventually_eq_zero_of_mul_eventually_eq_zero}
\leanok
Let \(\psi\) be holomorphic coordinates on an open set of \(M\) and \(A_i=\{\psi_{\sigma(i)}=0\}\) (\(i\in T\), \(T\) finite, \(\sigma\) injective on \(T\)) coordinate
hyperplanes which are submanifolds at their points; put \(u_i=\psi_{\sigma(i)}\).
\begin{enumerate}
\item If \(\mathcal I=(u)\) at \(y\), the germs of \(\mathcal I^a\) at \(y\) are the multiples of \(u^a\); if \(u(y)\ne0\), every germ at \(y\) is a multiple of \(u^a\).
\item Points at which the coordinates \(\psi_l\), \(l\in L\), are nonzero and the other coordinates are those of \(y\) accumulate at \(y\).
\item If a germ \(g\) at a point \(y\) of the chart is a multiple of \(u_i^{a_i}\) for each \(i\in T\), then it is a multiple of \(\prod_{i\in T}u_i^{a_i}\).
\item If \(m\) is holomorphic at \(y\) and does not vanish identically near \(y\), and \(fm=0\) near \(y\), then \(f=0\) near \(y\).
\end{enumerate}
\end{lemma}
\begin{proof}
\leanok
(1) Induction on \(a\): a germ of \(\mathcal I^{a+1}\) is a combination of products \(pq\) with \(p\) a multiple of \(u^a\) and \(q\) one of \(u\). If \(u(y)\neq0\),
\(g=(gu^{-a})u^a\). (2) Move along the line \(t\mapsto\psi(y)+t\sum_{l\in L}e_l\) for small \(t\ne0\).
(3) Induction on \(T\). Let \(T=T'\cup\{j\}\), \(g=c'p\) with \(p=\prod_{T'}u_i^{a_i}\), and \(g=du_j^{a_j}\). If \(u_j(y)\ne0\), use (1). Otherwise \(y\in A_j\). By (2) the
points \(y'\in A_j\) with \(p(y')\neq0\) (the coordinates \(u_i\), \(i\in T'\), are different from \(u_j\)) accumulate at \(y\); near such \(y'\),
\(c'=(d/p)\,u_j^{a_j}\) lies in \(\mathfrak m_z^{a_j}\) at every \(z\in A_j\). By the identity principle along \(A_j\) for one germ
(Lemma~\ref{hir:exponent-locally-constant}(3)), \(c'\in(\mathcal I_{A_j}^{a_j})_y\), and \(\mathcal I_{A_j}=(u_j)\) near \(y\) (Hadamard for a coordinate hyperplane), so
\(c'\) is a multiple of \(u_j^{a_j}\) by (1).
(4) \(m\notin\mathfrak m_y^{b+1}\) for some \(b\) (Lemma~\ref{hir:exponent-locally-constant}(2)); since \(fm=0\in\mathfrak m_y^{q+b}\), \(f\in\mathfrak m_y^q\) for every \(q\)
(Lemma~\ref{hir:order-product}), so \(f=0\) near \(y\).
\end{proof}

\begin{proposition}[Monomial and nonmonomial parts, \(N=M\)]
\label{hir:monomial-part-ambient}
\uses{hir:monomial-factorization,hir:exponent-locally-constant,hir:snc,bcalc:member-intersection,hir:hadamard-bridge}
\source{bierstone-milman08-functoriality,wlodarczyk05-hironaka}
\lean{TranscendentalNakamaye.Pending.MonomialPart.isClosedSubmanifold_get,TranscendentalNakamaye.Pending.MonomialPart.monomialPart,TranscendentalNakamaye.Pending.MonomialPart.nonmonomialPart,TranscendentalNakamaye.Pending.MonomialPart.le_monomialPart,TranscendentalNakamaye.Pending.MonomialPart.exists_generator_monomialPart,TranscendentalNakamaye.Pending.MonomialPart.mem_congr_of_isCoherent,TranscendentalNakamaye.Pending.MonomialPart.isCoherent_monomialPart,TranscendentalNakamaye.Pending.MonomialPart.exists_mem_nonmonomialPart_mul,TranscendentalNakamaye.Pending.MonomialPart.eq_mul_monomialPart,TranscendentalNakamaye.Pending.MonomialPart.exists_mem_nonmonomialPart_notMem,TranscendentalNakamaye.Pending.MonomialPart.isCoherent_nonmonomialPart}
\leanok
Let \(E\) be a simple normal crossings divisor on \(M\) and \(\mathfrak a\) a coherent ideal such that \(\mathfrak a_y\) contains a germ not vanishing near \(y\), for every
\(y\in M\). For a member \(H\) let \(e_H\) be the exponent of Lemma~\ref{hir:exponent-locally-constant}. Let \(\mathcal M(\mathfrak a)_y\) consist of the germs lying in
\((\mathcal I_H^{e_H(y)})_y\) for every member \(H\ni y\), and \(\mathcal N(\mathfrak a)=(\mathfrak a:\mathcal M(\mathfrak a))\). Then:
\begin{enumerate}
\item \(\mathfrak a\subset\mathcal M(\mathfrak a)\); near every \(x\), \(\mathcal M(\mathfrak a)\) is generated by one function \(m=\prod_{H\ni x}u_H^{e_H(x)}\) (\(u_H\) a
coordinate cutting out \(H\)), which vanishes identically near no point; \(\mathcal M(\mathfrak a)\) is coherent;
\item every germ of \(\mathfrak a_y\) is \(cm\) with \(c\in\mathcal N(\mathfrak a)_y\), \(m\in\mathcal M(\mathfrak a)_y\), and \(\mathfrak a=\mathcal N(\mathfrak a)\,\mathcal M(\mathfrak a)\);
\item at every point \(y\) of a member \(H\), \(\mathcal N(\mathfrak a)_y\not\subset(\mathcal I_H)_y\);
\item \(\mathcal N(\mathfrak a)\) is coherent.
\end{enumerate}
\end{proposition}
\begin{proof}
\leanok
The members are closed submanifolds (Lemma~\ref{bcalc:member-intersection}). (1) \(\mathfrak a\subset\mathcal M(\mathfrak a)\) is the definition of \(e_H\) with
Lemma~\ref{hir:exponent-locally-constant}(1). Take coordinates at \(x\) in which the members through \(x\) are the coordinate hyperplanes \(\{u_H=0\}\)
(Definition~\ref{hir:snc}), and a neighbourhood \(U\) of \(x\) meeting no other member, on which \(e_H\) is constant on \(H\)
(Lemma~\ref{hir:exponent-locally-constant}(4)). For \(y\in U\) and \(g\in\mathcal M(\mathfrak a)_y\), \(g\) is a multiple of \(u_H^{e_H(x)}\) for each \(H\ni x\) (by
Lemma~\ref{hir:monomial-factorization}(1), with \(\mathcal I_H=(u_H)\) if \(y\in H\), and because \(u_H\) is a unit at \(y\) otherwise), hence of \(m\)
(Lemma~\ref{hir:monomial-factorization}(3)); conversely \(m\in(\mathcal I_H^{e_H(y)})_y\) for \(H\ni y\). The points off all \(\{u_H=0\}\) accumulate everywhere
(Lemma~\ref{hir:monomial-factorization}(2)). (2) Write \(g\in\mathfrak a_y\subset\mathcal M(\mathfrak a)_y\) as \(cm\); for \(m'=c'm\in\mathcal M(\mathfrak a)_y\),
\(cm'=c'g\in\mathfrak a_y\), so \(c\in\mathcal N(\mathfrak a)_y\). Conversely products of \(\mathcal N(\mathfrak a)\) and \(\mathcal M(\mathfrak a)\) lie in \(\mathfrak a\) by definition, and
\(\mathfrak a\) is closed under holomorphic combinations. (3) Let \(e=e_H(y)\). By maximality some \(g\in\mathfrak a_y\) is not in \((\mathcal I_H^{e+1})_y\); write
\(g=cm\) as in (2); \(m\in(\mathcal I_H^e)_y\), so \(c\in(\mathcal I_H)_y\) would give \(g\in(\mathcal I_H^{e+1})_y\). (4) Let \(s_k\) generate \(\mathfrak a\) near \(x\);
\(s_k=t_km\) near \(x\). Near \(x\), \(\mathcal N(\mathfrak a)\) is generated by the \(t_k\): they lie in it since \(t_km'=c's_k\); and if \(h\in\mathcal N(\mathfrak a)_y\),
then \(hm=\sum a_ks_k=(\sum a_kt_k)m\), so \(h=\sum a_kt_k\) by Lemma~\ref{hir:monomial-factorization}(4).
\end{proof}

\begin{lemma}[Exponents along a submanifold]
\label{hir:exponent-along}
\uses{hir:exponent-locally-constant,hir:local-retraction,hir:point-power-pullback,hir:hadamard-bridge,hir:marked-ideal-operations,hir:cosupport,hir:cosupport-basic,hir:snc,bcalc:member-intersection-along}
\source{bierstone-milman08-functoriality,wlodarczyk05-hironaka}
\lean{TranscendentalNakamaye.Pending.DivisorExponentAlong.IsLocalRetraction,TranscendentalNakamaye.Pending.DivisorExponentAlong.exists_isLocalRetraction,TranscendentalNakamaye.Pending.DivisorExponentAlong.IsLocalRetraction.mem_self,TranscendentalNakamaye.Pending.DivisorExponentAlong.IsLocalRetraction.tendsto,TranscendentalNakamaye.Pending.DivisorExponentAlong.IsLocalRetraction.eventually,TranscendentalNakamaye.Pending.DivisorExponentAlong.IsLocalRetraction.eventually_comp_self,TranscendentalNakamaye.Pending.DivisorExponentAlong.IsLocalRetraction.holomorphicAt_comp,TranscendentalNakamaye.Pending.DivisorExponentAlong.IsLocalRetraction.eventually_eq_comp,TranscendentalNakamaye.Pending.DivisorExponentAlong.IsLocalRetraction.vanishesToOrderOn_iff,TranscendentalNakamaye.Pending.DivisorExponentAlong.IsLocalRetraction.eventually_eq_zero_of_forall,TranscendentalNakamaye.Pending.DivisorExponentAlong.IsLocalRetraction.exists_not_vanishesToOrderOn,TranscendentalNakamaye.Pending.DivisorExponentAlong.IsDivNearAlong,TranscendentalNakamaye.Pending.DivisorExponentAlong.isDivNearAlong_univ_iff,TranscendentalNakamaye.Pending.DivisorExponentAlong.isDivNearAlong_zero,TranscendentalNakamaye.Pending.DivisorExponentAlong.IsDivNearAlong.mono,TranscendentalNakamaye.Pending.DivisorExponentAlong.IsDivNearAlong.eventually,TranscendentalNakamaye.Pending.DivisorExponentAlong.isDivNearAlong_of_notMem,TranscendentalNakamaye.Pending.DivisorExponentAlong.isDivNearAlong_of_generators,TranscendentalNakamaye.Pending.DivisorExponentAlong.isDivNearAlong_of_forall_mem,TranscendentalNakamaye.Pending.DivisorExponentAlong.IsLocalRetraction.mem_idealPow_vanishingIdeal_comp,TranscendentalNakamaye.Pending.DivisorExponentAlong.isDivNearAlong_iff,TranscendentalNakamaye.Pending.DivisorExponentAlong.isDivNearAlong_of_frequently,TranscendentalNakamaye.Pending.DivisorExponentAlong.divExpAlong,TranscendentalNakamaye.Pending.DivisorExponentAlong.divExpAlong_univ,TranscendentalNakamaye.Pending.DivisorExponentAlong.IsNonzeroAlongAt,TranscendentalNakamaye.Pending.DivisorExponentAlong.isNonzeroAlongAt_univ_iff,TranscendentalNakamaye.Pending.DivisorExponentAlong.exists_not_isDivNearAlong,TranscendentalNakamaye.Pending.DivisorExponentAlong.isDivNearAlong_divExpAlong,TranscendentalNakamaye.Pending.DivisorExponentAlong.not_isDivNearAlong_divExpAlong_succ,TranscendentalNakamaye.Pending.DivisorExponentAlong.isDivNearAlong_iff_le_divExpAlong,TranscendentalNakamaye.Pending.DivisorExponentAlong.eventually_divExpAlong_eq,TranscendentalNakamaye.Pending.DivisorExponentAlong.memberInter_singleton,TranscendentalNakamaye.Pending.DivisorExponentAlong.regularAtDim_inter_get,TranscendentalNakamaye.Pending.DivisorExponentAlong.eventually_divExpAlong_get_eq,TranscendentalNakamaye.Pending.DivisorExponentAlong.isDivNearAlong_get_iff_le}
\leanok
Let \(\mathfrak a\) be a coherent ideal of \(\mathcal O_M\), \(N\subset M\) regular at \(x\in N\), \(r\) a local holomorphic retraction onto \(N\) at \(x\)
(Lemma~\ref{hir:local-retraction}), and \(H\subset M\) a set. Say that \(\mathfrak a\) is \emph{divisible by \(\mathcal I_{N\cap H}^e\) along \(N\) near \(x\)} if
the points of \(N\cap H\) near \(x\) lie in \(\operatorname{cosupp}(\mathfrak a|_N,e)\).
\begin{enumerate}
\item For \(y\in N\) near \(x\) and \(f\) holomorphic at \(y\): \(f|_N\in\mathfrak m_{N,y}^q\) if and only if \(f\circ r\in\mathfrak m_y^q\). Hence a germ lying in
\(\mathfrak m_{N,x}^q\) for every \(q\) vanishes on \(N\) near \(x\); and if \(\mathfrak a_x\) contains a germ which does not vanish identically on \(N\) near
\(x\), there is a largest \(e=e^N_H(x)\) such that \(\mathfrak a\) is divisible by \(\mathcal I_{N\cap H}^e\) along \(N\) near \(x\).
\item Let \(A\) be a submanifold at \(x\) which \(r\) maps into \(H\) near \(x\). If \(\mathfrak a\) is divisible by \(\mathcal I_{N\cap H}^e\) along \(N\) near \(x\), then
\(g\circ r\in(\mathcal I_A^e)_x\) for every \(g\in\mathfrak a_x\); and \(g=g\circ r\) on \(N\) near \(x\). Conversely, if \(A\supset N\cap H\) and every germ of
\(\mathfrak a_x\) agrees on \(N\) near \(x\) with a germ of \((\mathcal I_A^e)_x\), then \(\mathfrak a\) is divisible by \(\mathcal I_{N\cap H}^e\) along \(N\) near \(x\).
In particular, if \(N\cap H\) is a submanifold at \(x\), divisibility along \(N\) near \(x\) means that every germ of \(\mathfrak a_x\) agrees on \(N\) near \(x\) with a
germ of \((\mathcal I_{N\cap H}^e)_x\).
\item (Identity principle along \(N\cap H\).) Let \(N\cap H\) be a submanifold at \(x\). If the points \(y\in N\cap H\) near which \(\mathfrak a\) is divisible by
\(\mathcal I_{N\cap H}^e\) along \(N\) accumulate at \(x\), then \(\mathfrak a\) is divisible by \(\mathcal I_{N\cap H}^e\) along \(N\) near \(x\).
\item If \(N\cap H\) is a submanifold at \(x\) and, for every \(y\in N\) near \(x\), \(\mathfrak a_y\) contains a germ not vanishing identically on \(N\) near \(y\),
then \(e^N_H\) is constant on \(N\cap H\) near \(x\).
\item For \(N=M\) these are the notions of Lemma~\ref{hir:exponent-locally-constant}. If \(n\ge1\), \(N\) is closed and \(E\) is a simple normal crossings
divisor transverse to \(N\), then \(N\cap H\) is a submanifold at each of its points for every member \(H\) of \(E\), so (3) and (4) apply on \(N\cap H\).
\end{enumerate}
\end{lemma}
\begin{proof}
\leanok
The retraction data at \(x\) are retraction data at every point of \(N\) near \(x\), and \(r\circ r=r\) near \(x\).
(1) If \(f=h\) on \(N\) near \(y\) with \(h\in\mathfrak m_y^q\), then \(f\circ r=h\circ r\) near \(y\), because \(r\) maps a neighbourhood of \(y\) into \(N\), near
\(y\); and \(h\circ r\in\mathfrak m_y^q\) (Lemma~\ref{hir:point-power-pullback}, \(r(y)=y\)). Conversely \(f=f\circ r\) on \(N\) near \(y\). If
\(f|_N\in\mathfrak m_{N,x}^q\) for all \(q\), then \(f\circ r\in\mathfrak m_x^q\) for all \(q\), so \(f\circ r\) vanishes near \(x\)
(Lemma~\ref{hir:exponent-locally-constant}(2)) and \(f=f\circ r\) vanishes on \(N\) near \(x\). If \(g\in\mathfrak a_x\) does not vanish identically on \(N\) near
\(x\), choose \(q\) with \(g|_N\notin\mathfrak m_{N,x}^q\): divisibility by \(\mathcal I_{N\cap H}^{q+1}\) along \(N\) near \(x\) would put \(x\) into
\(\operatorname{cosupp}(\mathfrak a|_N,q+1)\).
(2) Let \(g\in\mathfrak a_x\) and \(z\in A\) near \(x\). Then \(w=r(z)\in N\cap H\) is near \(x\), so \(w\in\operatorname{cosupp}(\mathfrak a|_N,e)\) and
\(g\in\mathfrak a_w\) (coherence, Lemma~\ref{hir:hadamard-bridge}); by (1) at \(w\), \(g\circ r\in\mathfrak m_w^e\); pulling back by \(r\)
(Lemma~\ref{hir:point-power-pullback}), \(g\circ r\circ r\in\mathfrak m_z^e\), and \(g\circ r\circ r=g\circ r\) near \(z\). So \(g\circ r\in\mathfrak m_z^e\) at the
points \(z\in A\) near \(x\), and \(g\circ r\in(\mathcal I_A^e)_x\) by Lemma~\ref{hir:exponent-locally-constant}(1). Conversely, let \(s_1,\dots,s_K\) generate
\(\mathfrak a\) near \(x\) and \(s_i=h_i\) on \(N\) near \(x\) with \(h_i\in(\mathcal I_A^e)_x\); at the points \(z\in N\cap H\subset A\) near \(x\),
\(h_i\in(\mathcal I_A^e)_z\subset\mathfrak m_z^e\), so \(s_i|_N\in\mathfrak m_{N,z}^e\), and holomorphic combinations keep the order along \(N\)
(Lemma~\ref{hir:marked-ideal-operations}); so \(z\in\operatorname{cosupp}(\mathfrak a|_N,e)\). For the last sentence take \(A=N\cap H\): \(r\) maps the points of
\(N\cap H\) near \(x\) to themselves.
(3) Let \(s_1,\dots,s_K\) generate \(\mathfrak a\) near \(x\). Near an accumulating point \(y\), \(s_i|_N\in\mathfrak m_{N,z}^e\) at the points \(z\in N\cap H\), i.e.
\(s_i\circ r\in\mathfrak m_z^e\) by (1). The identity principle along the submanifold \(N\cap H\) for one germ (Lemma~\ref{hir:exponent-locally-constant}(3)),
applied to \(s_i\circ r\), gives \(s_i\circ r\in(\mathcal I_{N\cap H}^e)_x\); since \(s_i=s_i\circ r\) on \(N\), the converse half of (2) applies.
(4) Let \(e=e^N_H(x)\). Divisibility by \(\mathcal I_{N\cap H}^e\) along \(N\) near \(x\) is divisibility near every point close to \(x\); non-divisibility by
\(\mathcal I_{N\cap H}^{e+1}\) persists at the points of \(N\cap H\) near \(x\) by (3). By (1) at those points (which have retractions),
\(e^N_H(y)=e\) for \(y\in N\cap H\) near \(x\).
(5) For \(N=M\) the cosupport along \(M\) is the ambient cosupport. The trace \(N\cap H\) of a member is a closed submanifold by
Lemma~\ref{bcalc:member-intersection-along} (with the set of members \(\{H\}\)), and \(N\) has a local retraction at each of its points.
\end{proof}

\begin{lemma}[Simultaneous divisibility along a submanifold]
\label{hir:monomial-factorization-along}
\uses{hir:monomial-factorization,hir:exponent-along,hir:local-retraction,hir:point-power-pullback,bcalc:snc-simultaneous,hir:snc}
\source{bierstone-milman08-functoriality,wlodarczyk05-hironaka}
\lean{TranscendentalNakamaye.Pending.MonomialPartAlong.flatRetract,TranscendentalNakamaye.Pending.MonomialPartAlong.mdifferentiableAt_flatMap,TranscendentalNakamaye.Pending.MonomialPartAlong.eventually_planeRetract_flatMap_mem_target,TranscendentalNakamaye.Pending.MonomialPartAlong.isLocalRetraction_flatRetract,TranscendentalNakamaye.Pending.MonomialPartAlong.eventually_flatMap_flatRetract,TranscendentalNakamaye.Pending.MonomialPartAlong.eventually_flatRetract_mem,TranscendentalNakamaye.Pending.MonomialPartAlong.exists_eq_mul_prod_of_forall_along,TranscendentalNakamaye.Pending.MonomialPartAlong.exists_isLocalRetraction_preserving,TranscendentalNakamaye.Pending.MonomialPartAlong.isInSpanOnAt_eq_iff}
\leanok
Let \(\Phi\) be holomorphic coordinates near \(x\in M\) (after the chart at \(x\)), \(\psi\) the composite coordinate map, \(\varphi=\psi^{-1}\), and suppose
that near \(x\) a set \(N\) is \(\psi^{-1}(\mathbf C^{k_N}\times0)\). Let \(P\) be the linear projection onto \(\mathbf C^{k_N}\times0\) and \(r=\varphi\circ P\circ\psi\).
\begin{enumerate}
\item At every point \(y\in N\) of the coordinate domain, \(r\) is a local holomorphic retraction onto \(N\) (Lemma~\ref{hir:local-retraction}), and near \(y\) it
keeps the coordinates \(w_t\), \(t<k_N\); in particular it maps the coordinate hyperplane \(\{w_t=0\}\), \(t<k_N\), into itself near \(y\).
\item Let \(A_i=\{w_{\sigma(i)}=0\}\) (\(i\in T\), \(T\) finite, \(\sigma\) injective on \(T\), \(\sigma(i)<k_N\)) be submanifolds at their points and \(u_i=w_{\sigma(i)}\). If
\(g\) is holomorphic at a point \(y\in N\) of the coordinate domain and, for each \(i\in T\), \(g=c_iu_i^{a_i}\) on \(N\) near \(y\) with \(c_i\) holomorphic
at \(y\), then \(g=c\prod_{i\in T}u_i^{a_i}\) on \(N\) near \(y\) with \(c\) holomorphic at \(y\).
\item If \(n\ge1\) and \(E\) is a simple normal crossings divisor transverse to \(N\), then at every \(y\in N\) there is a local holomorphic retraction onto
\(N\) which maps every member of \(E\) through \(y\) into itself near \(y\).
\end{enumerate}
\end{lemma}
\begin{proof}
\leanok
(1) \(r(y)=\varphi(P\psi(y))=\varphi(\psi(y))=y\), and \(P\psi(p)\) stays in the target of \(\Phi\) for \(p\) near \(y\) (continuity, \(P\psi(y)=\psi(y)\)); so \(r\)
is holomorphic near \(y\), \(\psi(r(p))=P\psi(p)\in\mathbf C^{k_N}\times0\), i.e.\ \(r(p)\in N\), and \(r(p)=p\) for \(p\in N\) (then \(P\psi(p)=\psi(p)\)). The \(t\)-th
coordinate of \(P\psi(p)\) is that of \(\psi(p)\) for \(t<k_N\).
(2) By (1), \(g\circ r=(c_i\circ r)\,(u_i\circ r)^{a_i}=(c_i\circ r)\,u_i^{a_i}\) near \(y\) (ambiently), because \(r\) maps a neighbourhood of \(y\) into \(N\)
near \(y\), and \(u_i\circ r=u_i\) near \(y\). Lemma~\ref{hir:monomial-factorization}(3) applied to \(g\circ r\) gives \(g\circ r=c\prod_{i\in T}u_i^{a_i}\) near
\(y\), and \(g=g\circ r\) on \(N\) near \(y\).
(3) Take coordinates at \(y\) adapted to \(N\) and \(E\) (Lemma~\ref{bcalc:snc-simultaneous} with the centre \(N\)): \(N=\mathbf C^{k_N}\times0\) and the members
through \(y\) are hyperplanes \(\{w_{\sigma(i)}=0\}\) with \(\sigma(i)<k_N\) (Definition~\ref{hir:snc}); then (1).
\end{proof}

\begin{proposition}[Monomial and nonmonomial parts along a submanifold]
\label{hir:monomial-part-along}
\uses{hir:monomial-factorization-along,hir:exponent-along,hir:monomial-part-ambient,hir:monomial-factorization,hir:exponent-locally-constant,hir:snc,bcalc:snc-simultaneous,bcalc:member-intersection-along,ana:restricted-ideals,ana:restricted-calculus,ana:restricted-colon}
\source{bierstone-milman08-functoriality,wlodarczyk05-hironaka}
\lean{TranscendentalNakamaye.Pending.MonomialPartAlong.monomialPartAlong,TranscendentalNakamaye.Pending.MonomialPartAlong.nonmonomialPartAlong,TranscendentalNakamaye.Pending.MonomialPartAlong.le_monomialPartAlong,TranscendentalNakamaye.Pending.MonomialPartAlong.holomorphicAt_generator,TranscendentalNakamaye.Pending.MonomialPartAlong.mem_idealPow_generator,TranscendentalNakamaye.Pending.MonomialPartAlong.frequently_generator_ne_zero,TranscendentalNakamaye.Pending.MonomialPartAlong.monomialPartAlong_mem_iff,TranscendentalNakamaye.Pending.MonomialPartAlong.exists_generator_monomialPartAlong,TranscendentalNakamaye.Pending.MonomialPartAlong.IsMonomialAlong,TranscendentalNakamaye.Pending.MonomialPartAlong.isInSpanOnAt_iff_generator_of_isMonomialAlong,TranscendentalNakamaye.Pending.MonomialPartAlong.exists_mem_nonmonomialPartAlong_mul,TranscendentalNakamaye.Pending.MonomialPartAlong.restrictLe_mul_along,TranscendentalNakamaye.Pending.MonomialPartAlong.restrictLe_mul_along',TranscendentalNakamaye.Pending.MonomialPartAlong.exists_mem_nonmonomialPartAlong_not,TranscendentalNakamaye.Pending.MonomialPartAlong.monomialPartAlong_mem_of_notMem,TranscendentalNakamaye.Pending.MonomialPartAlong.isInSpanAt_range_one_of_holomorphicAt,TranscendentalNakamaye.Pending.MonomialPartAlong.isCoherent_monomialPartAlong,TranscendentalNakamaye.Pending.MonomialPartAlong.isCoherent_nonmonomialPartAlong}
\leanok
Let \(n\ge1\), \(N\subset M\) a closed submanifold, \(E\) a simple normal crossings divisor transverse to \(N\), \(\mathfrak a\) a coherent ideal such that for
every \(y\in N\) the stalk \(\mathfrak a_y\) contains a germ not vanishing identically on \(N\) near \(y\), and \(e_i=e^N_{E_i}\) the exponents of
Lemma~\ref{hir:exponent-along}. Let \(\mathcal M_N(\mathfrak a)_y\), for \(y\in N\), consist of the holomorphic germs whose restriction to \(N\) lies in
\((\mathcal I_{E_i}^{e_i(y)})|_N\) for every member \(E_i\ni y\) (every holomorphic germ if \(y\notin N\)), and
\(\mathcal N_N(\mathfrak a)=(\widetilde{\mathfrak a}^N:\mathcal M_N(\mathfrak a))\) (Lemma~\ref{ana:restricted-colon}, \(\widetilde{\mathfrak a}^N\) the germs agreeing
on \(N\) with germs of \(\mathfrak a\)). Then:
\begin{enumerate}
\item \(\mathfrak a\subset\mathcal M_N(\mathfrak a)\).
\item Near \(x\in N\), in coordinates adapted to \(N\) and \(E\) (\(N=\mathbf C^{k_N}\times0\), the members through \(x\) the hyperplanes \(\{w_{\sigma(i)}=0\}\),
\(\sigma(i)<k_N\)), let \(m=\prod_{x\in E_i}w_{\sigma(i)}^{e_i(x)}\). For \(y\) near \(x\), a germ lies in \(\mathcal M_N(\mathfrak a)_y\) if and only if it is holomorphic
and its restriction to \(N\) is a multiple of \(m|_N\); \(m\in(\mathcal I_{E_i}^{e_i(x)})_x\) for \(E_i\ni x\); and \(m\) vanishes identically on no relatively
open subset of \(N\) near \(x\).
\item \(\operatorname{Res}_N\mathfrak a=\operatorname{Res}_N(\mathcal M_N(\mathfrak a)\,\mathcal N_N(\mathfrak a))\); more precisely every \(g\in\mathfrak a_y\), \(y\in N\), is
\(m\,c\) on \(N\) near \(y\) with \(m\) as in (2) at \(y\) and \(c\in\mathcal N_N(\mathfrak a)_y\).
\item At every point \(y\) of \(N\cap E_i\) some germ of \(\mathcal N_N(\mathfrak a)_y\) does not restrict on \(N\) into \((\mathcal I_{E_i})|_N\).
\item \(\mathcal M_N(\mathfrak a)\) and \(\mathcal N_N(\mathfrak a)\) are coherent.
\item If \(\mathfrak a\) is \emph{monomial along \(N\)}, i.e.\ \(\operatorname{Res}_N\mathcal M_N(\mathfrak a)\subset\operatorname{Res}_N\mathfrak a\), then in the coordinates of
(2) the restriction of \(\mathfrak a_x\) to \(N\) is generated by \(m|_N\).
\end{enumerate}
\end{proposition}
\begin{proof}
\leanok
(1) Let \(g\in\mathfrak a_y\), \(y\in N\cap E_i\), and \(e=e_i(y)\). Take a retraction \(r\) onto \(N\) at \(y\) mapping \(E_i\) into itself
(Lemma~\ref{hir:monomial-factorization-along}(3)); \(\mathfrak a\) is divisible by \(\mathcal I_{N\cap E_i}^e\) along \(N\) near \(y\), so
\(g\circ r\in(\mathcal I_{E_i}^e)_y\) by Lemma~\ref{hir:exponent-along}(2) with \(A=E_i\) (a submanifold, Definition~\ref{hir:snc}), and \(g=g\circ r\) on \(N\).
(2) Near \(x\) the members through a point are among those through \(x\) (they are closed), the coordinate domain contains the point, and
\(e_i(y)=e_i(x)\) for \(y\in N\cap E_i\) (Lemma~\ref{hir:exponent-along}(4), (5)). Let \(g\in\mathcal M_N(\mathfrak a)_y\), \(y\in N\). For \(E_i\ni y\), \(g|_N\) is a
combination of germs of \(\mathcal I_{E_i}^{e_i(x)}\), which are multiples of \(w_{\sigma(i)}^{e_i(x)}\) (Lemma~\ref{hir:monomial-factorization}(1), \(\mathcal I_{E_i}\) being
generated by \(w_{\sigma(i)}\)); for \(E_i\ni x\) not through \(y\), \(w_{\sigma(i)}\) is a unit at \(y\). Lemma~\ref{hir:monomial-factorization-along}(2) gives
\(g=c\,m\) on \(N\) near \(y\). Conversely \(m=\bigl(\prod_{j\ne i}w_{\sigma(j)}^{e_j(x)}\bigr)w_{\sigma(i)}^{e_i(x)}\in\mathcal I_{E_i}^{e_i(x)}\) at the points of the coordinate
domain. Off \(N\) both sides contain every germ (\(N\) is closed). Points of \(N\) where the coordinates \(w_{\sigma(i)}\) are nonzero and the others are unchanged
accumulate at every point of \(N\) (Lemma~\ref{hir:monomial-factorization}(2); the coordinates of index \(\ge k_N\), which cut out \(N\), are unchanged), and \(m\)
does not vanish there.
(3) By (1) and (2) at \(y\), \(g=c\,m\) on \(N\) near \(y\); \(m\in\mathcal M_N(\mathfrak a)_y\), and \(c\in\mathcal N_N(\mathfrak a)_y\): for \(g'\in\mathcal M_N(\mathfrak a)_y\), \(g'=b\,m\)
on \(N\), so \(g'c=b\,g\) on \(N\), a germ of \(\widetilde{\mathfrak a}^N\). The converse inclusion \(\mathcal M_N\mathcal N_N\subset\widetilde{\mathfrak a}^N\) is the
definition of the colon (Lemma~\ref{ana:restricted-colon}).
(4) Let \(e=e_i(y)\). By Lemma~\ref{hir:exponent-along}(1) \(\mathfrak a\) is not divisible by \(\mathcal I_{N\cap E_i}^{e+1}\) along \(N\) near \(y\), so by the converse half
of Lemma~\ref{hir:exponent-along}(2) (with \(A=E_i\)) some \(g\in\mathfrak a_y\) agrees on \(N\) with no germ of \(\mathcal I_{E_i}^{e+1}\). Write \(g=c\,m\) on \(N\) as in (3).
If \(c|_N=\sum a_kt_k|_N\) with \(t_k\in(\mathcal I_{E_i})_y\), then \(g=\sum a_k\,(m\,t_k)\) on \(N\) with \(m\,t_k\in\mathcal I_{E_i}^{e+1}\) by (2), a contradiction.
(5) Near \(x\notin N\) both ideals contain every germ (\(1\in\widetilde{\mathfrak a}^N\) off the closed \(N\)) and are generated by \(1\). Near \(x\in N\),
\(\mathcal M_N(\mathfrak a)=(m)+\mathcal I_N\) by (2), and \(\mathcal I_N\) is coherent (Lemma~\ref{ana:restricted-calculus}). For \(\mathcal N_N(\mathfrak a)\), let \(s_k\) generate
\(\mathfrak a\) near \(x\) and \(s_k=t_k\,m\) on \(N\) near \(x\) (by (1), (2)). Then \(\mathcal N_N(\mathfrak a)\) is generated near \(x\) by the \(t_k\) and generators of
\(\mathcal I_N\): they lie in it (as in (3); a germ of \(\mathcal I_N\) times anything restricts to \(0\)); and if \(h\in\mathcal N_N(\mathfrak a)_y\), then \(m\,h|_N=\sum a_js_{\rho(j)}|_N=
m\sum a_jt_{\rho(j)}|_N\), and \(m\) can be cancelled on \(N\) because it vanishes identically on no relatively open subset of \(N\) (Lemma~\ref{ana:restricted-colon});
so \(h-\sum a_jt_{\rho(j)}\in(\mathcal I_N)_y\).
(6) A germ of \(\mathfrak a_x\) lies in \(\mathcal M_N(\mathfrak a)_x\) by (1), hence restricts to a multiple of \(m\) by (2); conversely \(m\in\mathcal M_N(\mathfrak a)_x\) by
(2), so \(m|_N\) is a combination of germs of \(\mathfrak a_x\) on \(N\) by the hypothesis.
\end{proof}

\begin{lemma}[Exponents along the members of \(E\)]
\label{hir:divisor-exponent}
\uses{hir:exponent-along,hir:monomial-part-along,hir:monomial-factorization-along,hir:monomial-part-ambient,hir:marked-ideal,hir:snc}
\source{bierstone-milman08-functoriality,wlodarczyk05-hironaka}
\lean{TranscendentalNakamaye.Pending.DivisorExponentAlong.IsDivNearAlong,TranscendentalNakamaye.Pending.DivisorExponentAlong.divExpAlong,TranscendentalNakamaye.Pending.DivisorExponentAlong.exists_not_isDivNearAlong,TranscendentalNakamaye.Pending.DivisorExponentAlong.isDivNearAlong_get_iff_le,TranscendentalNakamaye.Pending.DivisorExponentAlong.isDivNearAlong_iff,TranscendentalNakamaye.Pending.DivisorExponentAlong.eventually_divExpAlong_get_eq,TranscendentalNakamaye.Pending.MonomialPartAlong.monomialPartAlong_mem_iff,TranscendentalNakamaye.Pending.MonomialPartAlong.restrictLe_mul_along,TranscendentalNakamaye.Pending.MonomialPartAlong.restrictLe_mul_along',TranscendentalNakamaye.Pending.MonomialPartAlong.exists_mem_nonmonomialPartAlong_not}
\leanok
Let \(n\ge1\) and \(\underline{\mathcal I}=(M,N,E,\mathfrak a,d)\) be well formed, with \(\mathfrak a_y\) containing, for every \(y\in N\), a germ which does not vanish
identically on \(N\) near \(y\) (for \(\mathfrak a\) coherent: \(\mathfrak a|_N\) vanishes identically on no nonempty open subset of \(N\)), and let \(H\) be a member
of \(E\). For \(x\in H\cap N\) let \(e_H(x)\) be the largest \(e\) such that \(\mathfrak a|_N\subset\mathcal I_{H\cap N}^e\) near \(x\), i.e.\ such that the points of
\(H\cap N\) near \(x\) lie in \(\operatorname{cosupp}(\mathfrak a|_N,e)\). Then \(e_H(x)\) is finite, \(e_H\) is locally constant on \(H\cap N\), and near every
\(x\in N\)
\[
 \mathfrak a|_N=\Bigl(\prod_{H\ni x}\mathcal I_H^{e_H(x)}\Bigr)\Big|_N\cdot\mathfrak a^\natural,
\]
where \(\mathfrak a^\natural\) is contained in \((\mathcal I_H)|_N\) at no point of \(H\cap N\), for every member \(H\) through \(x\).
\end{lemma}
\begin{proof}
\leanok
The two descriptions of divisibility agree by Lemma~\ref{hir:exponent-along}(2) (the trace \(H\cap N\) is a submanifold, Lemma~\ref{hir:exponent-along}(5)); the
finiteness and the local constancy are Lemma~\ref{hir:exponent-along}(1), (4), (5): a local retraction onto \(N\) turns orders along \(N\) into ambient orders, and
the identity principle along the submanifold \(H\cap N\) of Lemma~\ref{hir:exponent-locally-constant} applies to the retracted generators. In coordinates
adapted to \(N\) and \(E\) the factor \(\bigl(\prod_{H\ni x}\mathcal I_H^{e_H(x)}\bigr)|_N\) is generated by the restriction of the monomial \(m\) of
Proposition~\ref{hir:monomial-part-along}(2) (simultaneous divisibility along \(N\), Lemma~\ref{hir:monomial-factorization-along}), which is the restriction of
\(\mathcal M_N(\mathfrak a)\); the factorization with \(\mathfrak a^\natural=\mathcal N_N(\mathfrak a)|_N\) and the non-containment are
Proposition~\ref{hir:monomial-part-along}(3), (4). For \(N=M\) this is Proposition~\ref{hir:monomial-part-ambient}. (As \(H\) is transverse to \(N\),
\((\mathcal I_H)|_N\) is the ideal of \(H\cap N\) in \(\mathcal O_N\): in adapted coordinates both are generated by the restriction of one coordinate of \(N\);
this identification is not used.)
\end{proof}

\begin{definition}[Monomial and nonmonomial part]
\label{hir:monomial-part}
\uses{hir:divisor-exponent,hir:monomial-part-ambient,hir:monomial-part-along}
\source{bierstone-milman08-functoriality,wlodarczyk05-hironaka}
\lean{TranscendentalNakamaye.Pending.MonomialPartAlong.monomialPartAlong,TranscendentalNakamaye.Pending.MonomialPartAlong.nonmonomialPartAlong,TranscendentalNakamaye.Pending.MonomialPartAlong.IsMonomialAlong}
\leanok
For \(\underline{\mathcal I}\) as in Lemma~\ref{hir:divisor-exponent}, the \emph{monomial part} \(\mathcal M(\underline{\mathcal I})\) is the ideal of \(\mathcal O_N\) whose
stalk at \(x\in N\) consists of the restrictions to \(N\) of the germs lying in \(\mathcal I_H^{e_H(x)}\) on \(N\) for every member \(H\ni x\), i.e.\
\(\bigl(\prod_{H\ni x}\mathcal I_H^{e_H(x)}\bigr)|_N\) (Proposition~\ref{hir:monomial-part-along}(2); in \(\mathcal O_N\) this is \(\prod_{H\ni x}\mathcal I_{H\cap N,x}^{e_H(x)}\)), and
the \emph{nonmonomial part} is \(\mathcal N(\underline{\mathcal I})=(\mathfrak a|_N:\mathcal M(\underline{\mathcal I}))\), so that
\(\mathfrak a|_N=\mathcal M(\underline{\mathcal I})\,\mathcal N(\underline{\mathcal I})\) and \(\mathcal N(\underline{\mathcal I})\) is contained in \((\mathcal I_H)|_N\) at no point of \(H\cap N\). The marked ideal is \emph{monomial} if
\(\mathcal M(\underline{\mathcal I})\subset\mathfrak a|_N\), i.e.\ \(\mathcal N(\underline{\mathcal I})=\mathcal O_N\). Both are coherent; their ambient representatives
are \(\mathcal M_N(\mathfrak a)\) and \(\mathcal N_N(\mathfrak a)\) of Proposition~\ref{hir:monomial-part-along}, coherent by its (5) (as in
Definition~\ref{hir:marked-ideal}, only their restrictions to \(N\) matter). For \(N=M\) this is Proposition~\ref{hir:monomial-part-ambient}.
\end{definition}

\begin{definition}[Companion ideal]
\label{hir:companion-ideal}
\uses{hir:monomial-part,hir:marked-sum-product,hir:cosupport,hir:invariant-order}
\source{bierstone-milman08-functoriality,wlodarczyk05-hironaka}
\lean{TranscendentalNakamaye.Pending.CompanionReductionAlong.companionAlong,TranscendentalNakamaye.Pending.CompanionRebase.companionOf,TranscendentalNakamaye.Pending.CompanionRebase.companionAlong_ideal,TranscendentalNakamaye.Pending.CompanionRebase.companionAlong_d,TranscendentalNakamaye.Pending.CompanionRebase.pullback_val_companionOf}
\leanok
Let \(\underline{\mathcal I}\) be as in Lemma~\ref{hir:divisor-exponent}, \(d\ge1\), and \(W\subset M\) open such that
\(o=\max\{\operatorname{ord}_x\mathcal N(\underline{\mathcal I}):x\in\operatorname{cosupp}\underline{\mathcal I}\cap W\}\) exists and is \(\ge1\) (for instance \(W=M\)
when the cosupport is compact and the maximum is \(\ge1\); \(x\mapsto\operatorname{ord}_x\mathcal N(\underline{\mathcal I})\) is upper semicontinuous with finite values, and
Lemma~\ref{hir:invariant-order} applies). The \emph{companion ideal} of \(\underline{\mathcal I}\) on \(W\) is the marked ideal on \(W\)
\[
 \mathcal O(\underline{\mathcal I})=
 \begin{cases}
 (W,N,E,\mathcal N(\underline{\mathcal I}),o)+(W,N,E,\mathcal M(\underline{\mathcal I}),d-o) & \text{if } o<d,\\
 (W,N,E,\mathcal N(\underline{\mathcal I}),o) & \text{if } o\ge d,
 \end{cases}
\]
with the sum of Definition~\ref{hir:marked-sum-product} (here \(N\) and \(E\) stand for their traces on \(W\)).
\end{definition}

\begin{lemma}[The companion ideal is of maximal order]
\label{hir:companion-basic}
\uses{hir:companion-ideal,hir:companion-maximal-order,hir:companion-maximal-order-along,hir:companion-reduction-along,hir:order-product,hir:order-power,hir:marked-ideal-operations,hir:marked-operations-along,hir:monomial-part-along,bcalc:companion-open,bcalc:companion-rebase,bcalc:monomial-part-etale}
\source{bierstone-milman08-functoriality}
\lean{TranscendentalNakamaye.Pending.CompanionBasic.companion_basic}
\leanok
\(\mathcal O(\underline{\mathcal I})\) is of maximal order along \(N\): at every \(x\in N\cap W\) the order of its ideal is at most its mark. Moreover
\(\operatorname{cosupp}\mathcal O(\underline{\mathcal I})=\operatorname{cosupp}(\mathcal N(\underline{\mathcal I}),o)\cap\operatorname{cosupp}\underline{\mathcal I}\cap W\).
\end{lemma}
\begin{proof}
\leanok
Write \(\nu=\operatorname{ord}_x\mathcal N(\underline{\mathcal I})\) and \(m=\operatorname{ord}_x\mathcal M(\underline{\mathcal I})\); by Lemma~\ref{hir:order-product}, \(\operatorname{ord}_x(\mathfrak a|_N)=m+\nu\).
If \(o<d\), the ideal of \(\mathcal O(\underline{\mathcal I})\) is \(\mathcal N^{d-o}+\mathcal M^{o}\) with mark \(o(d-o)\), and its order at \(x\) is
\(\min((d-o)\nu,\,om)\) (Lemma~\ref{hir:order-power} and the order of a sum). If \(x\in\operatorname{cosupp}\underline{\mathcal I}\), then \(\nu\le o\); otherwise
\(m+\nu<d\), so \(\nu\le o\) or \(m<d-o\); in each case the minimum is at most \(o(d-o)\). If \(o\ge d\), then \(\nu\le o\) on the cosupport, and off it
\(\nu\le m+\nu<d\le o\). The cosupport of the sum is the intersection of the cosupports (Lemma~\ref{hir:marked-ideal-operations}, along \(N\)
Lemma~\ref{hir:marked-operations-along}); at a point with \(\nu\ge o\) and \(\nu\le o\) one has \(m\ge d-o\) if and only if \(m+\nu\ge d\); and if \(o\ge d\),
\(\nu\ge o\) implies \(m+\nu\ge d\).
For \(N=M\) and an abstract factorization \(\mathfrak a=\mathfrak b\mathfrak c\) the maximal-order half is Lemma~\ref{hir:companion-maximal-order} (proved);
the factorization \(\mathfrak a|_N=\mathcal M\,\mathcal N\) along \(N\) is Proposition~\ref{hir:monomial-part-along}(3) (proved, with coherent ambient
representatives, (5)), and the maximal-order half along \(N\) for an abstract factorization is Lemma~\ref{hir:companion-maximal-order-along} (proved), its
cosupport half Lemma~\ref{hir:companion-reduction-along} with \(T\) the trivial tower (proved). On an open \(W\ne M\), with the bound
\(\operatorname{ord}\mathcal N\le o\) only on \(\operatorname{cosupp}\underline{\mathcal I}\cap W\), the statement is Lemma~\ref{bcalc:companion-open}(3) (proved), for the marked
ideals of the domain of Definition~\ref{hir:centre-function} and with the nonmonomial part of \(\underline{\mathcal I}|_W\), the restriction of
\(\mathcal N(\underline{\mathcal I})\) (Lemma~\ref{bcalc:monomial-part-etale}); this is the form the general value uses. The lemma as stated here, for every
\(\underline{\mathcal I}\) as in Lemma~\ref{hir:divisor-exponent} and with \(\mathcal N(\underline{\mathcal I})\) itself, follows without the domain
hypothesis: the companion ideal commutes with restriction to \(W\) (Lemma~\ref{bcalc:companion-rebase}(2)), and so do the nonmonomial part and its cosupports
(Lemma~\ref{bcalc:monomial-part-etale}); the restriction \(\underline{\mathcal I}|_W\) is well formed and nonzero along \(N\cap W\), and the bound holds on its
cosupport, so the maximal-order and cosupport statements of Lemma~\ref{bcalc:companion-open}(2) apply to \(\underline{\mathcal I}|_W\) (the latter with \(T\) the trivial tower).
This is the Lean form.
\end{proof}

\begin{lemma}[Reduction to maximal order along multiple test blow-ups]
\label{hir:companion-reduction}
\uses{hir:order-product,hir:marked-operations-transform,hir:marked-ideal-operations,hir:marked-sum-product,hir:geo-transform-ambient,hir:hadamard-bridge,hir:giraud-support,mc:inclusions,bcalc:controlled-transform}
\source{bierstone-milman08-functoriality,wlodarczyk05-hironaka}
\lean{TranscendentalNakamaye.Pending.CompanionReduction.mul_mem_of_isCoherent,TranscendentalNakamaye.Pending.CompanionReduction.pullback_le_controlledTransform,TranscendentalNakamaye.Pending.CompanionReduction.comp_mem_pullback',TranscendentalNakamaye.Pending.CompanionReduction.exists_mem_ne_zero_of_notMem_zeroLocus,TranscendentalNakamaye.Pending.CompanionReduction.notMem_cosupport_controlledTransform,TranscendentalNakamaye.Pending.CompanionReduction.notMem_cosupport_succ_controlledTransform,TranscendentalNakamaye.Pending.CompanionReduction.controlledTransform_le_controlledTransform,TranscendentalNakamaye.Pending.CompanionReduction.mem_cosupp_of_mem_cosupp_geoTransform,TranscendentalNakamaye.Pending.CompanionReduction.notMem_cosupport_succ_geoTransform,TranscendentalNakamaye.Pending.CompanionReduction.cosupp_add_eq_inter_cosupp_mul,TranscendentalNakamaye.Pending.CompanionReduction.reduction_add,TranscendentalNakamaye.Pending.CompanionReduction.reduction_le,TranscendentalNakamaye.Pending.CompanionReduction.bound_along,TranscendentalNakamaye.Pending.CompanionReduction.companion,TranscendentalNakamaye.Pending.CompanionReduction.mul_eq_of_le,TranscendentalNakamaye.Pending.CompanionReduction.companion_reduction,TranscendentalNakamaye.Pending.CompanionReduction.companion_resolution}
\leanok
Let \(N=M\), let \(\mathfrak b,\mathfrak c\) be coherent, \(\mathfrak a=\mathfrak b\mathfrak c\), \(d\ge1\), \(o\ge1\), \(E\) a common sequence, and suppose
\(\operatorname{ord}_x\mathfrak b\le o\) for every \(x\in\operatorname{cosupp}(\mathfrak a,d)\). Put \(\underline{\mathcal O}=(M,M,E,\mathfrak b,o)+(M,M,E,\mathfrak c,d-o)\) if \(o<d\),
and \(\underline{\mathcal O}=(M,M,E,\mathfrak b,o)\) if \(o\ge d\). Then every multiple test blow-up \(T\) of \(\underline{\mathcal O}\) is a multiple test blow-up of
\(\underline{\mathcal I}=(M,M,E,\mathfrak a,d)\) and of \((M,M,E,\mathfrak b,o)\); along it \(\operatorname{ord}_z\mathfrak b_T\le o\) for every
\(z\in\operatorname{cosupp}\underline{\mathcal I}_T\), and
\[
 \operatorname{cosupp}\underline{\mathcal O}_T=\operatorname{cosupp}(\mathfrak b_T,o)\cap\operatorname{cosupp}\underline{\mathcal I}_T .
\]
In particular, if \(T\) is a resolution of \(\underline{\mathcal O}\), then \(\operatorname{ord}_z\mathfrak b_T<o\) at every \(z\in\operatorname{cosupp}\underline{\mathcal I}_T\).
\end{lemma}
\begin{proof}
\leanok
Induction on the length of \(T\); write \(\mathfrak b_T,\mathfrak c_T,\mathfrak a_T\) for the transforms of \((\mathfrak b,o)\), \((\mathfrak c,d-o)\), \((\mathfrak a,d)\).
Case \(o<d\): by Lemma~\ref{hir:marked-operations-transform} the multiple test blow-ups of the sum are the simultaneous ones of its summands, these
are multiple test blow-ups of the product \((\mathfrak b,o)\cdot(\mathfrak c,d-o)=(\mathfrak a,d)\), and along them \(\mathfrak a_T=\mathfrak b_T\mathfrak c_T\) and
\(\operatorname{cosupp}\underline{\mathcal O}_T=\operatorname{cosupp}(\mathfrak b_T,o)\cap\operatorname{cosupp}(\mathfrak c_T,d-o)\). The latter is contained in
\(\operatorname{cosupp}(\mathfrak a_T,d)\) (Lemma~\ref{hir:order-basic}), and conversely a point of \(\operatorname{cosupp}(\mathfrak b_T,o)\cap\operatorname{cosupp}(\mathfrak a_T,d)\) with
\(\operatorname{ord}\mathfrak b_T\le o\) lies in \(\operatorname{cosupp}(\mathfrak c_T,d-o)\) by Lemma~\ref{hir:order-product}(2). Case \(o\ge d\): a centre \(C\) admissible for
\((\mathfrak b_T,o)\) lies in \(\operatorname{cosupp}(\mathfrak b_T,o)\subset\operatorname{cosupp}(\mathfrak a_T,d)\), because \(\mathfrak a_T\subset\mathfrak b_T\); and the inclusion persists:
\(\sigma^c(\mathfrak a_T,d)\subset\sigma^c(\mathfrak b_T,o)\), because a germ \(g\) with \(g\,\mathcal I_D^{d}\subset\sigma^*\mathfrak a_T\) satisfies
\(g\,\mathcal I_D^{o}\subset g\,\mathcal I_D^{d}\subset\sigma^*\mathfrak b_T\) (the controlled transform is the quotient of Definition~\ref{bcalc:controlled-transform-def}).
Persistence of the bound, in both cases: by Lemma~\ref{hir:giraud-support}, \(\operatorname{ord}_z\mathfrak b\le o\) if and only if \(1\in(\mathcal D^o\mathfrak b)_z\). If
\(z\in\operatorname{cosupp}(\mathfrak a_{T'},d)\) for \(T'=T\cdot\sigma\), then \(\sigma(z)\in\operatorname{cosupp}(\mathfrak a_T,d)\): otherwise
\(1\in(\mathcal D^{d-1}\mathfrak a_T)_{\sigma(z)}\), hence \(1\in\sigma^c(\mathcal D^{d-1}\mathfrak a_T,1)_z\subset(\mathcal D^{d-1}\sigma^c(\mathfrak a_T,d))_z\)
(Lemma~\ref{mc:inclusions}; the centre lies in \(\operatorname{cosupp}(\mathfrak a_T,d)\), so \(\mathfrak a_T\subset\mathcal I_C^d\) by Lemma~\ref{hir:hadamard-bridge}),
i.e.\ \(z\notin\operatorname{cosupp}(\mathfrak a_{T'},d)\). By induction \(1\in(\mathcal D^o\mathfrak b_T)_{\sigma(z)}\), so
\(1\in\sigma^c(\mathcal D^o\mathfrak b_T,0)_z\subset(\mathcal D^o\sigma^c(\mathfrak b_T,o))_z\) by Lemma~\ref{mc:inclusions}. The last assertion is the case
\(\operatorname{cosupp}\underline{\mathcal O}_T=\emptyset\).
\end{proof}
Along a submanifold \(N\ne M\) the same argument applies with the calculus of Lemma~\ref{hir:along-calculus}: Lemma~\ref{hir:companion-reduction-along}.

\begin{lemma}[The companion ideal of a factorization is of maximal order]
\label{hir:companion-maximal-order}
\uses{hir:companion-reduction,mc:maximal-order,hir:giraud-support,hir:order-product,hir:order-power,hir:order-basic,hir:marked-sum-product}
\source{wlodarczyk05-hironaka,bierstone-milman08-functoriality}
\lean{TranscendentalNakamaye.Pending.CompanionReduction.one_mem_of_apply_ne_zero,TranscendentalNakamaye.Pending.CompanionReduction.isMaximalOrder_of_forall_notMem,TranscendentalNakamaye.Pending.CompanionReduction.holomorphicAt_pow_of,TranscendentalNakamaye.Pending.CompanionReduction.pow_notMem_idealPow,TranscendentalNakamaye.Pending.CompanionReduction.exists_notMem_idealPow_of_notMem_cosupport,TranscendentalNakamaye.Pending.CompanionReduction.mem_idealPow_of_mem_mul,TranscendentalNakamaye.Pending.CompanionReduction.companion_isMaximalOrder}
\leanok
In the situation of Lemma~\ref{hir:companion-reduction} (\(N=M\), \(\mathfrak a=\mathfrak b\mathfrak c\) with \(\mathfrak b,\mathfrak c\) coherent, and
\(\operatorname{ord}_x\mathfrak b\le o\) for every \(x\in\operatorname{cosupp}(\mathfrak a,d)\)), the marked ideal \(\underline{\mathcal O}\) is of maximal order
(Definition~\ref{mc:maximal-order}): \(\operatorname{ord}_x\le\) mark at every \(x\in M\).
\end{lemma}
\begin{proof}
\leanok
By Lemma~\ref{hir:giraud-support}, \((\mathfrak e,\mu)\) is of maximal order as soon as no point lies in \(\operatorname{cosupp}(\mathfrak e,\mu+1)\): at such a point
\(\mathcal D^\mu\mathfrak e\) contains a germ not vanishing at \(x\), hence a unit. Let \(x\in M\) and write \(\nu=\operatorname{ord}_x\mathfrak b\).
Case \(o<d\): the ideal is \(\mathfrak b^{d-o}+\mathfrak c^{o}\) with mark \(o(d-o)\). If \(\nu\le o\), pick \(g\in\mathfrak b_x\) of order \(\le o\); then
\(g^{d-o}\) has order \(\le o(d-o)\) (Lemma~\ref{hir:order-product}, iterated). If \(\nu>o\), then \(x\notin\operatorname{cosupp}(\mathfrak a,d)\) by hypothesis, so some germ
of \(\mathfrak b\mathfrak c\) has order \(<d\); since every product \(ab\) with \(a\in\mathfrak b_x\) has order \(\ge o+1+\operatorname{ord}b\), some \(f\in\mathfrak c_x\) has
order \(<d-o\), and \(f^o\in(\mathfrak c^o)_x\) has order \(\le o(d-o-1)<o(d-o)+1\). Case \(o\ge d\): the ideal is \(\mathfrak b\) with mark \(o\); if
\(\nu\ge o+1\), then \(x\notin\operatorname{cosupp}(\mathfrak a,d)\) by hypothesis, but every germ of \(\mathfrak b\mathfrak c\) lies in
\(\mathfrak m_x^{o+1}\subset\mathfrak m_x^{d}\), a contradiction.
\end{proof}

\begin{lemma}[Reduction to maximal order inside a submanifold]
\label{hir:companion-reduction-along}
\uses{hir:companion-reduction,hir:marked-operations-along,hir:restricted-equivalence-basic,hir:marked-ideal-operations,hir:marked-sum-product,hir:along-transform,hir:giraud-along,hir:exponent-along,hir:local-retraction,hir:order-product,hir:order-basic,hir:geo-transform-coherent,hir:geo-transform-wellformed,hir:test-tower,bcalc:transform-off-centre}
\source{bierstone-milman08-functoriality,wlodarczyk05-hironaka}
\lean{TranscendentalNakamaye.Pending.CompanionReductionAlong.exists_not_vanishesToOrderOn_of_notMem,TranscendentalNakamaye.Pending.CompanionReductionAlong.mem_cosupport_of_mem_cosupport_mul_along,TranscendentalNakamaye.Pending.CompanionReductionAlong.mem_cosupp_of_mem_cosupp_geoTransform_along,TranscendentalNakamaye.Pending.CompanionReductionAlong.notMem_cosupport_succ_geoTransform_along,TranscendentalNakamaye.Pending.CompanionReductionAlong.geoTransform_ideal_le,TranscendentalNakamaye.Pending.CompanionReductionAlong.cosupp_add_eq_inter_cosupp_mul_along,TranscendentalNakamaye.Pending.CompanionReductionAlong.cosupp_subset_of_le,TranscendentalNakamaye.Pending.CompanionReductionAlong.bound_along_along,TranscendentalNakamaye.Pending.CompanionReductionAlong.reduction_add_along,TranscendentalNakamaye.Pending.CompanionReductionAlong.reduction_le_along,TranscendentalNakamaye.Pending.CompanionReductionAlong.companionAlong,TranscendentalNakamaye.Pending.CompanionReductionAlong.companionAlong_N,TranscendentalNakamaye.Pending.CompanionReductionAlong.companionAlong_E,TranscendentalNakamaye.Pending.CompanionReductionAlong.mul_eq_of_le_along,TranscendentalNakamaye.Pending.CompanionReductionAlong.companion_reduction_along,TranscendentalNakamaye.Pending.CompanionReductionAlong.companion_resolution_along}
\leanok
Let \(N\subset M\) be a closed submanifold, \(\mathfrak b,\mathfrak c\) coherent, \(o\ge1\), \(d\in\mathbb N\), \(E\) a common sequence, and suppose
\(\operatorname{ord}_x(\mathfrak b|_N)\le o\) for every \(x\in\operatorname{cosupp}((\mathfrak b\mathfrak c)|_N,d)\). Put
\(\underline{\mathcal O}=(M,N,E,\mathfrak b,o)+(M,N,E,\mathfrak c,d-o)\) if \(o<d\), and \(\underline{\mathcal O}=(M,N,E,\mathfrak b,o)\) if \(o\ge d\). Then every
multiple test blow-up \(T\) of \(\underline{\mathcal O}\) is a multiple test blow-up of \(\underline{\mathcal I}=(M,N,E,\mathfrak b\mathfrak c,d)\) and of
\((M,N,E,\mathfrak b,o)\); along it \(\operatorname{ord}_z(\mathfrak b_T|_{N_T})\le o\) for every \(z\in\operatorname{cosupp}\underline{\mathcal I}_T\), and
\[
 \operatorname{cosupp}\underline{\mathcal O}_T=\operatorname{cosupp}(\mathfrak b_T|_{N_T},o)\cap\operatorname{cosupp}\underline{\mathcal I}_T .
\]
In particular, if \(T\) is a resolution of \(\underline{\mathcal O}\), then \(\operatorname{ord}_z(\mathfrak b_T|_{N_T})<o\) at every \(z\in\operatorname{cosupp}\underline{\mathcal I}_T\).
For \(N=M\) this is Lemma~\ref{hir:companion-reduction}. A marked ideal \((M,N,E,\mathfrak a,d)\) with \(\mathfrak a|_N=(\mathfrak b\mathfrak c)|_N\) is restricted-equivalent to
\(\underline{\mathcal I}\), with the same multiple test blow-ups and transforms (Lemma~\ref{hir:restricted-equivalence-basic}); in the application
\(\mathfrak b=\mathcal N_N(\mathfrak a)\), \(\mathfrak c=\mathcal M_N(\mathfrak a)\) (Proposition~\ref{hir:monomial-part-along}).
\end{lemma}
\begin{proof}
\leanok
The proof of Lemma~\ref{hir:companion-reduction}, with orders along \(N\) and the calculus along \(N\).

Products along \(N\). Let \(x\in\operatorname{cosupp}((\mathfrak b\mathfrak c)|_N,d)\) with \(\operatorname{ord}_x(\mathfrak b|_N)\le o\le d\). Take a local holomorphic retraction
\(r\) onto \(N\) at \(x\) (Lemma~\ref{hir:local-retraction}); a germ \(g\) has order \(\ge q\) along \(N\) at \(x\) if and only if \(g\circ r\in\mathfrak m_x^q\)
(Lemma~\ref{hir:exponent-along}). Pick \(g\in\mathfrak b_x\) with \(g\circ r\notin\mathfrak m_x^{o+1}\). For \(f\in\mathfrak c_x\), \(gf\in(\mathfrak b\mathfrak c)_x\), so
\((f\circ r)(g\circ r)\in\mathfrak m_x^{d}=\mathfrak m_x^{(d-o)+o}\), and Lemma~\ref{hir:order-product}(2) gives \(f\circ r\in\mathfrak m_x^{d-o}\). Hence
\(x\in\operatorname{cosupp}(\mathfrak c|_N,d-o)\).

One blowup. Let \(\underline{\mathcal J}=(M,N,E,\mathfrak e,\mu)\) with \(N\) a closed submanifold and \(\mathfrak e\) coherent, \(C\) an admissible centre and
\(\sigma\) the blowup along \(C\). The cosupport of \(\underline{\mathcal J}'\) lies over \(\operatorname{cosupp}\underline{\mathcal J}\): over \(C\) because
\(C\subset\operatorname{cosupp}\underline{\mathcal J}\), off \(C\) by Lemma~\ref{bcalc:transform-off-centre}. If \(\operatorname{ord}_{\sigma(z)}(\mathfrak e|_N)\le\mu\) for
\(z\in N'\), then by Giraud along \(N\) (Lemma~\ref{hir:giraud-along}) some \(u\in\mathcal D^\mu_N(\mathfrak e)\) satisfies \(u(\sigma(z))\ne0\); \(u\circ\sigma\) lies in the transform
of \((M,N,E,\mathcal D_N^\mu\mathfrak e,0)\), hence, on \(N'\), in \(\mathcal D^\mu_{N'}(\mathfrak e')\) (Lemma~\ref{hir:along-transform}); so
\(\mathcal D^\mu_{N'}(\mathfrak e')\) does not vanish at \(z\), i.e.\ \(\operatorname{ord}_z(\mathfrak e'|_{N'})\le\mu\) (Lemma~\ref{hir:giraud-along} on \(N'\)). Finally the transform is
monotone: for marked ideals on the same \(N\) with \(\mathfrak a\subset\mathfrak b\) and marks \(d\le o\), a germ \(g\) with
\(g\,(\sigma^*\mathcal I_C)^d\subset\sigma^*\mathfrak a\) on \(N'\) satisfies \(g\,(\sigma^*\mathcal I_C)^o\subset\sigma^*\mathfrak b\) on \(N'\), so the transform of
\((\mathfrak a,d)\) is contained in that of \((\mathfrak b,o)\) (Definition~\ref{hir:test-tower}). The transforms keep closed submanifolds and coherent ideals
(Lemmas~\ref{hir:geo-transform-wellformed} and~\ref{hir:geo-transform-coherent}).

Persistence of the bound: if \(T\) is a multiple test blow-up of both \((M,N,E,\mathfrak b,o)\) and \(\underline{\mathcal I}\), induction on the length of \(T\) with the two
facts of the previous paragraph gives \(\operatorname{ord}_z(\mathfrak b_T|_{N_T})\le o\) on \(\operatorname{cosupp}\underline{\mathcal I}_T\).

Case \(o<d\). By Lemma~\ref{hir:marked-operations-along} the multiple test blow-ups of the sum are the simultaneous ones of its summands, these are multiple test
blow-ups of the product \((\mathfrak b,o)\cdot(\mathfrak c,d-o)=(\mathfrak b\mathfrak c,d)\), and along them the transforms of the sum and of the product are
restricted-equivalent to the sum and the product of the transforms, so they have the same cosupports (Lemma~\ref{hir:restricted-equivalence-basic}). The cosupport
of the sum of the transforms is \(\operatorname{cosupp}(\mathfrak b_T,o)\cap\operatorname{cosupp}(\mathfrak c_T,d-o)\) (Lemma~\ref{hir:marked-operations-along}); it lies in
the cosupport of the product (Lemma~\ref{hir:marked-ideal-operations}), and conversely, by the bound and the paragraph on products,
\(\operatorname{cosupp}(\mathfrak b_T,o)\cap\operatorname{cosupp}\underline{\mathcal I}_T\subset\operatorname{cosupp}(\mathfrak c_T,d-o)\).
Case \(o\ge d\). Since \(\mathfrak b\mathfrak c\subset\mathfrak b\) and \(d\le o\), \(\operatorname{cosupp}(\mathfrak b|_N,o)\subset\operatorname{cosupp}((\mathfrak b\mathfrak c)|_N,d)\), so a centre
admissible for \((M,N,E,\mathfrak b,o)\) is admissible for \(\underline{\mathcal I}\); by monotonicity of the transform the inclusion of the ideals, hence of the
cosupports, persists along \(T\) (induction). The last assertion is the case \(\operatorname{cosupp}\underline{\mathcal O}_T=\emptyset\).
\end{proof}

\begin{lemma}[The companion ideal of a factorization along a submanifold is of maximal order]
\label{hir:companion-maximal-order-along}
\uses{hir:companion-reduction-along,hir:companion-maximal-order,mc:maximal-order-along,hir:exponent-along,hir:local-retraction,hir:order-product,hir:order-power,hir:marked-ideal-operations,hir:marked-sum-product}
\source{wlodarczyk05-hironaka,bierstone-milman08-functoriality}
\lean{TranscendentalNakamaye.Pending.CompanionReductionAlong.companion_isMaximalOrderAlong}
\leanok
In the situation of Lemma~\ref{hir:companion-reduction-along} (only that \(N\) is a submanifold, i.e.\ regular at each of its points, is needed), the marked
ideal \(\underline{\mathcal O}\) is of maximal order along \(N\) (Definition~\ref{mc:maximal-order-along}): the order along \(N\) of its ideal is at most its mark at
every point of \(N\).
\end{lemma}
\begin{proof}
\leanok
Let \(x\in N\) and take a local retraction \(r\) onto \(N\) at \(x\); orders along \(N\) at \(x\) are the orders at \(x\) of the composites with \(r\)
(Lemma~\ref{hir:exponent-along}), and \((g^k)\circ r=(g\circ r)^k\). Case \(o<d\): the ideal is \(\mathfrak b^{d-o}+\mathfrak c^{o}\) with mark \(o(d-o)\).
If \(\operatorname{ord}_x(\mathfrak b|_N)\le o\), pick \(g\in\mathfrak b_x\) with \(g\circ r\notin\mathfrak m_x^{o+1}\); then \((g\circ r)^{d-o}\notin\mathfrak m_x^{o(d-o)+1}\)
(Lemma~\ref{hir:order-product}, iterated, as in Lemma~\ref{hir:companion-maximal-order}). Otherwise \(x\notin\operatorname{cosupp}((\mathfrak b\mathfrak c)|_N,d)\) by
hypothesis, and since \(\operatorname{cosupp}(\mathfrak b|_N,o+1)\cap\operatorname{cosupp}(\mathfrak c|_N,d-o)\subset\operatorname{cosupp}((\mathfrak b\mathfrak c)|_N,d+1)\)
(Lemma~\ref{hir:marked-ideal-operations}), \(x\notin\operatorname{cosupp}(\mathfrak c|_N,d-o)\): some \(f\in\mathfrak c_x\) has \(f\circ r\notin\mathfrak m_x^{d-o}\), and
\(f^o\in(\mathfrak c^o)_x\) (Lemma~\ref{hir:order-power}) has \((f\circ r)^o\notin\mathfrak m_x^{o(d-o-1)+1}\supset\mathfrak m_x^{o(d-o)+1}\). Case \(o\ge d\): the ideal is \(\mathfrak b\) with mark
\(o\); if \(\operatorname{ord}_x(\mathfrak b|_N)\ge o+1\), every germ of \(\mathfrak b\mathfrak c\subset\mathfrak b\) has order \(\ge o+1\ge d\) along \(N\) at \(x\), so
\(x\in\operatorname{cosupp}((\mathfrak b\mathfrak c)|_N,d)\), contradicting the hypothesis.
\end{proof}

\begin{lemma}[Uniqueness of the monomial factorization]
\label{hir:monomial-factorization-unique}
\uses{hir:monomial-part-ambient,hir:monomial-factorization,hir:exponent-locally-constant,hir:monomial-invariant-geometry,hir:order-product,hir:companion-maximal-order}
\source{bierstone-milman08-functoriality}
\lean{TranscendentalNakamaye.Pending.NonmonomialTransform.divExp_eq_of_mul_principal,TranscendentalNakamaye.Pending.NonmonomialTransform.monomialPart_mem_iff_of_mul_principal,TranscendentalNakamaye.Pending.NonmonomialTransform.nonmonomialPart_mem_iff_of_mul_principal}
\leanok
Let \(n\ge1\), \(E\) a simple normal crossings divisor on \(M\), \(z\in M\), and \(u\) coordinates centred at \(z\) in which the members \(H_i\) through \(z\) are the
hyperplanes \(\{u_{\tau(i)}=0\}\) (\(\tau\) injective); let \(m=\prod_{z\in H_i}u_{\tau(i)}^{b_i}\). Let \(\mathfrak n\) be coherent and suppose that near \(z\) the ideal
\(\mathfrak b\) (of holomorphic germs) is \(\mathfrak n\cdot(m)\), and that \(\mathfrak n_y\not\subset\mathcal I_{H_i,y}\) at the points \(y\in H_i\) near \(z\). Then
\(e_{H_i}(z)=b_i\) for the members through \(z\), \(\mathcal M(\mathfrak b)_z=(m)\) and \(\mathcal N(\mathfrak b)_z=\mathfrak n_z\).
\end{lemma}
\begin{proof}
\leanok
Every germ of \(\mathfrak b\) near \(z\) is a combination of products \(am\), whose order at points of \(H_i\) is at least \(b_i\); so \(e_{H_i}(z)\ge b_i\)
(Lemma~\ref{hir:exponent-locally-constant}). Conversely, the points \(y_0\in H_i\) near \(z\) off the other hyperplanes accumulate at \(z\); at such a point pick
\(a\in\mathfrak n_{y_0}\setminus\mathcal I_{H_i,y_0}\); \(a\) does not vanish identically on \(H_i\) near \(y_0\), so there are points \(y\in H_i\) near \(y_0\), still off the other
hyperplanes, with \(a(y)\ne0\) and \(a\in\mathfrak n_y\) (coherence). There \(am=(a\prod_{j\ne i}u_{\tau(j)}^{b_j})\,u_{\tau(i)}^{b_i}\) is a unit times a power of a
coordinate, of order exactly \(b_i\) (Lemma~\ref{hir:order-product}), so \(e_{H_i}(z)<b_i+1\). Hence \(\mathcal M(\mathfrak b)_z\), the germs divisible by
\(u_{\tau(i)}^{b_i}\) for every member through \(z\), is \((m)\) (Lemma~\ref{hir:monomial-factorization}(1)). Finally \(h\in\mathcal N(\mathfrak b)_z\) gives
\(hm\in\mathfrak b_z\), i.e.\ \((h-\sum c_ka_k)m=0\) near \(z\) with \(a_k\in\mathfrak n_z\); \(m\) vanishes identically near no point, so \(h=\sum c_ka_k\in\mathfrak n_z\); the
converse is immediate.
\end{proof}

\begin{lemma}[One blowup of a product with a monomial factor, \(N=M\)]
\label{hir:nonmonomial-step}
\uses{hir:monomial-step,hir:monomial-part-ambient,hir:monomial-factorization,hir:exponent-locally-constant,hir:monomial-invariant-geometry,hir:order-product,hir:hadamard-bridge,hir:controlled-operations,hir:geo-transform-ambient,bcalc:controlled-transform,bcalc:exceptional-divisor,bcalc:exceptional-pushforward,bcalc:controlled-not-exceptional,bcalc:monomial-chart,bcalc:monomial-transform-coords,bcalc:isomorphism-off-centre,bcalc:snc-transform,bcalc:strict-geo,bcalc:hyperplane-ideal}
\source{bierstone-milman08-functoriality,wlodarczyk05-hironaka}
\lean{TranscendentalNakamaye.Pending.NonmonomialTransformTower.mul_mem_iff_of_principal,TranscendentalNakamaye.Pending.NonmonomialTransformTower.isNonzeroAt_of_notMem,TranscendentalNakamaye.Pending.NonmonomialTransformTower.frequently_mem_notMem_of_hasNormalCrossings,TranscendentalNakamaye.Pending.NonmonomialTransformTower.le_sum_divExp_of_subset_cosupport,TranscendentalNakamaye.Pending.NonmonomialTransformTower.exceptionalDivisor_mem_of_vanishingIdeal_mem,TranscendentalNakamaye.Pending.NonmonomialTransformTower.isNonzeroAt_controlledTransform_of_notMem,TranscendentalNakamaye.Pending.NonmonomialTransformTower.idealPow_mem_iff_of_generator,TranscendentalNakamaye.Pending.NonmonomialTransformTower.step}
\leanok
Let \(n\ge1\), \(E=(H_1,\dots,H_s)\) a simple normal crossings divisor on \(M\), \(\mathfrak b\) and \(\mathcal M\) coherent ideals, \(\mathcal M\) principal monomial at
every point (Lemma~\ref{hir:monomial-step}), and suppose that \(\mathfrak b_y\) contains a germ not vanishing near \(y\) for every \(y\in M\) and that
\(\mathfrak b_y\not\subset\mathcal I_{H_i,y}\) for every \(i\) and \(y\in H_i\). Let \(o,d\in\mathbb N\) and let \(C\) be a closed submanifold with only normal crossings
with \(E\), \(C\subset\operatorname{cosupp}(\mathfrak b,o)\cap\operatorname{cosupp}(\mathfrak b\mathcal M,d)\), \(\operatorname{ord}_x\mathfrak b\le o\) for \(x\in C\). Let
\(\sigma\colon M'\to M\) be the blowup along \(C\), \(D=\sigma^{-1}(C)\) and \(E'=(H_1',\dots,H_s',D)\). Then
\[
 \sigma^c(\mathfrak b\mathcal M,d)=\sigma^c(\mathfrak b,o)\cdot\mathcal M',\qquad \mathcal M'=\mathcal O_{M'}(-(o\mathbin{\dot-}d)D)\,\sigma^c(\mathcal M,d\mathbin{\dot-}o)
\]
(\(\dot-\) the truncated difference), \(\mathcal M'\) is principal monomial for \(E'\) at every point of \(M'\), and \(\sigma^c(\mathfrak b,o)\) has the two properties of
\(\mathfrak b\) for \(E'\). Along the way:
\begin{enumerate}
\item at every point of \(D\), \(\mathcal I_D\subset\mathcal O_{M'}(-D)\) (so the two agree there);
\item if \(\mathcal M\) is monomial, \(\mathcal M_y\) contains a germ not vanishing near \(y\) for every \(y\), and \(C\subset\operatorname{cosupp}(\mathcal M,\mu)\), then at every
\(x\in C\), \(\mu\le\sum e_{H_i}(x)\), the sum over the members through \(x\) that contain \(C\) near \(x\).
\end{enumerate}
\end{lemma}
\begin{proof}
\leanok
(2) Take coordinates \(w\) centred at \(x\) in which \(C=\{w_t=0,\ t\in S_0\}\) and every member through \(x\) is \(\{w_{t_i}=0\}\). If \(t_i\in S_0\), then
\(C\subset H_i\) near \(x\); so every member through \(x\) not containing \(C\) near \(x\) has \(t_i\notin S_0\). The points with \(w_{t_i}\ne0\) for these \(i\) and all
other coordinates \(0\) lie on \(C\) and accumulate at \(x\). At such a point \(y\) close to \(x\) the members through \(y\) are among those through \(x\) containing
\(C\) near \(x\), with the exponents they have at \(x\) (Lemma~\ref{hir:exponent-locally-constant}); since \(y\in\operatorname{cosupp}(\mathcal M,\mu)\) and \(\mathcal M\) is
monomial, Lemma~\ref{hir:monomial-invariant-geometry}(4) gives \(\mu\le\sum_{y\in H_i}e_{H_i}(y)\le\sum_{C\subset H_i}e_{H_i}(x)\).

(1) Let \(z\in D\). \(D\) is a smooth hypersurface (Theorem~\ref{bcalc:snc-transform}), so \(\mathcal I_{D,z}=(u)\) for a coordinate \(u\) centred at \(z\)
(Lemma~\ref{bcalc:hyperplane-ideal}). In the monomial chart through \(z\) (Lemma~\ref{bcalc:monomial-chart}), \(\mathcal O(-D)=(W_j)\) near \(z\) with \(W_j(z)=0\) and \(dW_j(z)\ne0\);
\(W_j\) vanishes exactly on \(D\) near \(z\), so \(W_j=cu\) with \(c\) holomorphic. If \(c(z)=0\), then \(W_j\in\mathfrak m_z^2\), contradicting \(dW_j(z)\ne0\). Hence
\(u=c^{-1}W_j\in\mathcal O(-D)_z\), and \(\mathcal I_{D,z}\subset\mathcal O(-D)_z\).

The factorization. Since \(\operatorname{ord}\mathfrak b\le o\) on \(C\subset\operatorname{cosupp}(\mathfrak b\mathcal M,d)\), \(C\subset\operatorname{cosupp}(\mathcal M,d\mathbin{\dot-}o)\)
(Lemma~\ref{hir:order-product}(2); trivial if \(o\ge d\)). By the Hadamard bridge (Lemma~\ref{hir:hadamard-bridge}) \(\mathfrak b\subset\mathcal I_C^{o}\) and
\(\mathcal M\subset\mathcal I_C^{d\dot-o}\), so (Lemma~\ref{bcalc:controlled-transform}, products and powers as in Lemma~\ref{hir:controlled-operations})
\[
 \sigma^*(\mathfrak b\mathcal M)=\mathcal O(-oD)\sigma^c(\mathfrak b,o)\,\mathcal O(-(d\mathbin{\dot-}o)D)\sigma^c(\mathcal M,d\mathbin{\dot-}o)
 =\mathcal O(-dD)\cdot\sigma^c(\mathfrak b,o)\,\mathcal O(-(o\mathbin{\dot-}d)D)\sigma^c(\mathcal M,d\mathbin{\dot-}o),
\]
as \(o+(d\mathbin{\dot-}o)=d+(o\mathbin{\dot-}d)\); the controlled transform is the unique coherent quotient by \(\mathcal O(-dD)\).

\(\mathcal M'\) is principal monomial. If \(\sigma(z)\notin C\), \(\mathcal O(-D)\) is the unit ideal at \(z\), and \(\sigma^c(\mathcal M,d\mathbin{\dot-}o)_z=(\sigma^*\mathcal M)_z\) is
principal monomial with the exponents of \(\sigma(z)\), as in Lemma~\ref{hir:monomial-step}(2). If \(x=\sigma(z)\in C\), then \(\mathcal M_x=(\prod_{x\in H_i}u_{\tau(i)}^{e_{H_i}(x)})\)
(Lemma~\ref{hir:monomial-step}(1)) and \(d\mathbin{\dot-}o\le\sum_{C\subset H_i}e_{H_i}(x)\) by (2); Corollary~\ref{bcalc:monomial-transform-coords} gives coordinates \(\Theta\)
centred at \(z\), adapted to \(E'\), with \(\sigma^c(\mathcal M,d\mathbin{\dot-}o)_z=(\prod_{z\in E'_i}\Theta_{\tau'(i)}^{a'_i})\). As \(D=\{\Theta_{\tau'(D)}=0\}\) near \(z\),
\(\mathcal O(-D)_z=\mathcal I_{D,z}=(\Theta_{\tau'(D)})\) by (1) and the inclusion \(\mathcal O(-D)\subset\mathcal I_D\); so \(\mathcal M'_z\) is generated by the monomial with the
exponent on \(D\) raised by \(o\mathbin{\dot-}d\).

The properties of \(\sigma^c(\mathfrak b,o)\). At \(z\in D\): \(\operatorname{ord}_{\sigma(z)}\mathfrak b\le o\), so some germ of \(\sigma^c(\mathfrak b,o)_z\) is not in \(\mathcal O(-D)_z\)
(Corollary~\ref{bcalc:controlled-not-exceptional}, resting on Theorem~\ref{bcalc:exceptional-pushforward}); by (1) it is not in \(\mathcal I_{D,z}\), and it does not vanish
near \(z\). At \(z\in H_i'\): \(\mathfrak b\not\subset\mathcal I_{H_i}\) at the points of \(H_i\setminus C\) near \(\sigma(z)\), so the second clause of
Corollary~\ref{bcalc:controlled-not-exceptional} (with Lemma~\ref{bcalc:strict-geo}) gives a germ of \(\sigma^c(\mathfrak b,o)_z\) outside \(\mathcal I_{H_i',z}\). Off \(D\),
\(\sigma\) has a holomorphic section over \(M\setminus C\) (Lemma~\ref{bcalc:isomorphism-off-centre}); for \(a\in\mathfrak b_{\sigma(z)}\) not vanishing near \(\sigma(z)\), the germ
\(a\circ\sigma\in(\sigma^*\mathfrak b)_z\subset\sigma^c(\mathfrak b,o)_z\) does not vanish near \(z\).
\end{proof}

\begin{theorem}[Transforms of the monomial and nonmonomial parts, \(N=M\)]
\label{hir:nonmonomial-transform-ambient}
\uses{hir:nonmonomial-step,hir:monomial-factorization-unique,hir:monomial-part-ambient,hir:monomial-factorization,hir:monomial-step,hir:companion-reduction,hir:test-tower,hir:geo-transform-ambient,bcalc:transform-wellformed}
\source{bierstone-milman08-functoriality,wlodarczyk05-hironaka}
\lean{TranscendentalNakamaye.Pending.NonmonomialTransformTower.isNonzeroAt_mul_of_isPrincipalAt,TranscendentalNakamaye.Pending.NonmonomialTransformTower.parts_eq_of_factor,TranscendentalNakamaye.Pending.NonmonomialTransformTower.isPrincipalAt_monomialPart,TranscendentalNakamaye.Pending.NonmonomialTransformTower.factor_along,TranscendentalNakamaye.Pending.NonmonomialTransformTower.nonmonomialPart_geoTransformAlong,TranscendentalNakamaye.Pending.NonmonomialTransformTower.nonmonomialPart_geoTransformAlong_of_resolution}
\leanok
Let \(n\ge1\) and \(\underline{\mathcal I}=(M,M,E,\mathfrak a,d)\) be well formed with \(d\ge1\) and \(\mathfrak a_y\) containing a germ not vanishing near \(y\) for every \(y\). Let
\(\mathcal N=\mathcal N(\mathfrak a)\), \(\mathcal M=\mathcal M(\mathfrak a)\) (Proposition~\ref{hir:monomial-part-ambient}), \(o\ge1\) with \(\operatorname{ord}_x\mathcal N\le o\) for every
\(x\in\operatorname{cosupp}\underline{\mathcal I}\), and \(\underline{\mathcal O}=(M,M,E,\mathcal N,o)+(M,M,E,\mathcal M,d-o)\) if \(o<d\), \(\underline{\mathcal O}=(M,M,E,\mathcal N,o)\) if
\(o\ge d\). Let \(T\) be a multiple test blow-up of \(\underline{\mathcal O}\). Then \(T\) is a multiple test blow-up of \(\underline{\mathcal I}\); \(\underline{\mathcal I}_T\) is well
formed with \(N_T=M_T\) and \(\mathfrak a_T\) nowhere locally zero; \(\mathcal N(\mathfrak a_T)\) is the transform of \((M,M,E,\mathcal N,o)\) along \(T\); \(\mathcal M(\mathfrak a_T)\)
is principal monomial at every point; \(\mathfrak a_T=\mathcal N(\mathfrak a_T)\,\mathcal M(\mathfrak a_T)\); and \(\operatorname{ord}_z\mathcal N(\mathfrak a_T)\le o\) for every
\(z\in\operatorname{cosupp}\underline{\mathcal I}_T\). If \(T\) resolves \(\underline{\mathcal O}\), then \(\operatorname{ord}_z\mathcal N(\mathfrak a_T)<o\) for every
\(z\in\operatorname{cosupp}\underline{\mathcal I}_T\).
\end{theorem}
\begin{proof}
\leanok
\(\mathcal M\) is principal monomial at every point (simple normal crossings coordinates at \(x\) and Lemma~\ref{hir:monomial-factorization}(1), as in
Lemma~\ref{hir:monomial-step}(1)); \(\mathcal N\) is coherent, \(\mathcal N_y\not\subset\mathcal I_{H,y}\) at the points \(y\) of a member \(H\), \(\mathfrak a=\mathcal N\mathcal M\), and
\(\mathcal N_y\supset\mathfrak a_y\) contains a germ not vanishing near \(y\) (Proposition~\ref{hir:monomial-part-ambient}). Conversely, if \(\mathfrak a'=\mathfrak b\mathcal M'\) with
\(\mathcal M'\) principal monomial at every point and \(\mathfrak b\) coherent with \(\mathfrak b_y\not\subset\mathcal I_{H,y}\) at the points of every member \(H\), then
\(\mathcal N(\mathfrak a')=\mathfrak b\) and \(\mathcal M(\mathfrak a')=\mathcal M'\) (Lemma~\ref{hir:monomial-factorization-unique} at each point, the factorization near \(z\) being
\(\mathfrak b\cdot(m)\) for a generator \(m\) of \(\mathcal M'_z\)).
By Lemma~\ref{hir:companion-reduction} (\(\mathfrak b=\mathcal N\), \(\mathfrak c=\mathcal M\)), \(T\) is a multiple test blow-up of \(\underline{\mathcal I}\) and of
\((M,M,E,\mathcal N,o)\); hence each centre \(C_k\) of \(T\) lies in \(\operatorname{cosupp}(\mathfrak a_{T_k},d)\cap\operatorname{cosupp}(\mathcal N_{T_k},o)\) and has only normal
crossings with \(E_{T_k}\) (\(T_k\) the first \(k\) blowups), and \(\operatorname{ord}\mathcal N_{T_k}\le o\) on \(\operatorname{cosupp}(\mathfrak a_{T_k},d)\). By induction on \(k\),
\(\mathfrak a_{T_k}=\mathcal N_{T_k}\mathcal M_k\) with \(\mathcal M_k\) principal monomial at every point and \(\mathcal N_{T_k}\) (the transform of \((\mathcal N,o)\)) coherent,
nowhere locally zero and contained in the ideal of no member of \(E_{T_k}\) at its points: Lemma~\ref{hir:nonmonomial-step} carries this from \(k\) to \(k+1\), the
transforms being well formed (Lemma~\ref{bcalc:transform-wellformed}) with \(N=M\) (Lemma~\ref{hir:geo-transform-ambient}). The converse above identifies
\(\mathcal N(\mathfrak a_T)=\mathcal N_T\) and \(\mathcal M(\mathfrak a_T)=\mathcal M_T\); \(\mathfrak a_T\) is nowhere locally zero since a germ not vanishing near \(z\) times
a monomial is not (orders add, Lemma~\ref{hir:order-product}). The bound is that of Lemma~\ref{hir:companion-reduction}, and so is the last assertion
(\(\operatorname{ord}\mathcal N_T<o\) on \(\operatorname{cosupp}\underline{\mathcal I}_T\) after a resolution of \(\underline{\mathcal O}\)).
\end{proof}

\begin{lemma}[One blowup of a product with a monomial factor, inside a submanifold]
\label{hir:nonmonomial-step-along}
\uses{hir:nonmonomial-step,hir:companion-reduction-along,hir:monomial-part-along,hir:exponent-along,hir:monomial-invariant-along,hir:monomial-step-along,hir:marked-operations-along-step,hir:restricted-calculus,hir:restricted-equivalence-basic,hir:geo-transform-coherent,hir:geo-transform-wellformed,ana:restricted-calculus,bcalc:monomial-transform-along-coords,bcalc:monomial-transform-along-off-centre,bcalc:controlled-not-exceptional-along,bcalc:trace-strict-transform,bcalc:member-trace-ideal,bcalc:exceptional-trace-ideal,bcalc:isomorphism-off-centre,bcalc:transform-wellformed}
\source{bierstone-milman08-functoriality,wlodarczyk05-hironaka}
\lean{TranscendentalNakamaye.Pending.NonmonomialStepAlong.isInSpanOnAt_self_eq,TranscendentalNakamaye.Pending.NonmonomialStepAlong.isInSpanOnAt_iff_of_ambient,TranscendentalNakamaye.Pending.NonmonomialStepAlong.isInSpanOnAt_mul_iff_of_generators,TranscendentalNakamaye.Pending.NonmonomialStepAlong.isInSpanOnAt_of_eventually_eq_zero,TranscendentalNakamaye.Pending.NonmonomialStepAlong.isNonzeroAlongAt_of_isInSpanOnAt,TranscendentalNakamaye.Pending.NonmonomialStepAlong.le_sum_divExpAlong_of_subset_cosupport,TranscendentalNakamaye.Pending.NonmonomialStepAlong.monomialFactor,TranscendentalNakamaye.Pending.NonmonomialStepAlong.isCoherent_monomialFactor,TranscendentalNakamaye.Pending.NonmonomialStepAlong.restrictLe_geoTransform_mul,TranscendentalNakamaye.Pending.NonmonomialStepAlong.isPrincipalAlongAt_monomialFactor,TranscendentalNakamaye.Pending.NonmonomialStepAlong.exists_not_isInSpanOnAt_geoTransform,TranscendentalNakamaye.Pending.NonmonomialStepAlong.isNonzeroAlongAt_geoTransform,TranscendentalNakamaye.Pending.NonmonomialStepAlong.step_along}
\leanok
Let \(n\ge1\), \(N\subset M\) a closed submanifold, \(E\) a simple normal crossings divisor transverse to \(N\), \(\mathcal M\) a coherent ideal which is nonzero and
monomial along \(N\) (Definition~\ref{hir:monomial-part}), and \(\mathfrak b\) a coherent ideal with \(\mathfrak b|_N\) vanishing identically near no point of \(N\) and,
at every point \(y\in N\) of a member \(H\), \(\mathfrak b_y|_N\not\subset(\mathcal I_H)_y|_N\). Let \(o,d\in\mathbb N\) and let \(C\) be a closed submanifold with only normal
crossings with \(E\), \(C\subset\operatorname{cosupp}(\mathfrak b|_N,o)\cap\operatorname{cosupp}((\mathfrak b\mathcal M)|_N,d)\) and
\(\operatorname{ord}_x(\mathfrak b|_N)\le o\) for \(x\in C\). Let \(\sigma\colon M'\to M\) be the blowup along \(C\), \(N'\) the strict transform of \(N\),
\(D=\sigma^{-1}(C)\), \(E'=(H_1',\dots,H_s',D)\), and \(\mathfrak a'\), \(\mathfrak b'\), \(\mathcal M''\) the transforms of \((M,N,E,\mathfrak b\mathcal M,d)\),
\((M,N,E,\mathfrak b,o)\), \((M,N,E,\mathcal M,d\mathbin{\dot-}o)\). Then, with \(\mathcal M'=(\sigma^*\mathcal I_C)^{o\dot-d}\,\mathcal M''\):
\(\mathfrak a'|_{N'}=(\mathfrak b'\mathcal M')|_{N'}\); \(\mathcal M'\) is principal monomial along \(N'\) for \(E'\) at every point of \(N'\); and \(\mathfrak b'\) has the two
properties of \(\mathfrak b\) along \(N'\) for \(E'\). For \(N=M\) this is Lemma~\ref{hir:nonmonomial-step}.
\end{lemma}
\begin{proof}
\leanok
The proof of Lemma~\ref{hir:nonmonomial-step}, with the calculus along \(N\).

The generic-point bound. If \(C\subset\operatorname{cosupp}(\mathcal M|_N,\mu)\), then at \(x\in C\), \(\mu\le\sum e_H(x)\) over the members \(H\ni x\) containing \(C\)
near \(x\) (exponents along \(N\), Lemma~\ref{hir:exponent-along}): points of \(C\) on no member through \(x\) that does not contain \(C\) near \(x\) accumulate at \(x\) (as in
Lemma~\ref{hir:nonmonomial-step}(2)); near \(x\) the members through such a point carry the exponents they have at \(x\) (local constancy), and the cosupport of a
monomial ideal along \(N\) is where the exponents add up to the mark (Lemma~\ref{hir:monomial-invariant-along}).

The factorization. \(C\subset\operatorname{cosupp}(\mathcal M|_N,d\mathbin{\dot-}o)\) (Lemma~\ref{hir:companion-reduction-along}, the paragraph on products;
trivial if \(o\ge d\)), so \(C\) is admissible for the three marked ideals, and \(\sigma^*\mathfrak e|_{N'}=((\sigma^*\mathcal I_C)^{\mu}\mathfrak e')|_{N'}\) for each of them
(Lemma~\ref{hir:marked-operations-along-step}). With \(\sigma^*(\mathfrak b\mathcal M)=\sigma^*\mathfrak b\,\sigma^*\mathcal M\) and \(o+(d\mathbin{\dot-}o)=d+(o\mathbin{\dot-}d)\),
\(((\sigma^*\mathcal I_C)^d\mathfrak a')|_{N'}=((\sigma^*\mathcal I_C)^d\,\mathfrak b'\mathcal M')|_{N'}\) (restricted products, Lemma~\ref{hir:restricted-calculus}), and
\((\sigma^*\mathcal I_C)^d\) is invertible along \(N'\) (Lemma~\ref{hir:geo-transform-wellformed}), so it cancels (Lemma~\ref{ana:restricted-calculus}).

\(\mathcal M'\) is principal monomial along \(N'\). Let \(z\in N'\). Take coordinates at \(x=\sigma(z)\) adapted to \(N\) and \(E\) in which \(\mathcal M|_N\) is generated by the
monomial with the exponents along \(N\) (Lemma~\ref{hir:monomial-step-along}). If \(x\in C\), the generic-point bound and Corollary~\ref{bcalc:monomial-transform-along-coords}
give coordinates \(\Theta\) centred at \(z\), adapted to \(N'\) and \(E'\), in which \(\mathcal M''|_{N'}\) is generated by a monomial; at \(z\) the ideal
\(\sigma^*\mathcal I_C\) lies in \(\mathcal I_D\) and \(\mathcal I_D\subset\mathcal O(-D)=\sigma^*\mathcal I_C\) (Lemma~\ref{hir:nonmonomial-step}(1)), so it is generated by the coordinate
\(\Theta_D\) of \(D\), and restricted generators multiply: \(\mathcal M'|_{N'}\) is generated by the monomial with the exponent on \(D\) raised by \(o\mathbin{\dot-}d\). If \(x\notin C\),
Lemma~\ref{bcalc:monomial-transform-along-off-centre} gives the generator of \(\mathcal M''|_{N'}\), and \(\sigma^*\mathcal I_C\) is the unit ideal at \(z\).

The properties of \(\mathfrak b'\). At \(z\in H'\cap N'\): \(H'\cap N'=(H\cap N)'\) (Lemma~\ref{bcalc:trace-strict-transform}); at the points \(y\in H\cap N\),
\(\mathfrak b_y|_N\not\subset(\mathcal I_H)_y|_N\) means that some germ of \(\mathfrak b_y\) does not vanish on \(H\cap N\) near \(y\) (Lemma~\ref{bcalc:member-trace-ideal}), so clause 2 of
Corollary~\ref{bcalc:controlled-not-exceptional-along} with \(A=H\cap N\) gives a germ of \(\mathfrak b'_z\) not vanishing on \((H\cap N)'\) near \(z\), i.e.\ not in
\((\mathcal I_{H'})_z\) on \(N'\) (Lemma~\ref{bcalc:member-trace-ideal} for the well-formed transform, Lemma~\ref{bcalc:transform-wellformed}). At \(z\in D\cap N'\): since
\(\operatorname{ord}_{\sigma(z)}(\mathfrak b|_N)\le o\), clause 1 of Corollary~\ref{bcalc:controlled-not-exceptional-along} gives a germ of \(\mathfrak b'_z\) not in \(\sigma^*\mathcal I_C\) on
\(N'\), and \((\sigma^*\mathcal I_C)|_{N'}=(\mathcal I_{D\cap N'})|_{N'}\supset(\mathcal I_D)|_{N'}\) (Lemma~\ref{bcalc:exceptional-trace-ideal}). Nonvanishing: over \(C\) the germ just found
does not vanish identically on \(N'\) near \(z\) (otherwise it would lie in every ideal on \(N'\)); off \(C\), for \(a\in\mathfrak b_{\sigma(z)}\) not vanishing identically on \(N\)
near \(\sigma(z)\), \(a\circ\sigma\) lies on \(N'\) in the span of \(\mathfrak b'_z\) (Lemma~\ref{bcalc:monomial-transform-along-off-centre}) and does not vanish identically on \(N'\)
near \(z\), along the holomorphic section of \(\sigma\) over \(M\setminus C\) (Lemma~\ref{bcalc:isomorphism-off-centre}); a combination that does not vanish identically has a
term that does not.
\end{proof}

\begin{lemma}[The factorization along the companion tower, inside a submanifold]
\label{hir:nonmonomial-factor-along}
\uses{hir:nonmonomial-step-along,hir:companion-reduction-along,hir:restricted-equivalence-basic,hir:monomial-invariant-along,hir:test-tower,bcalc:transform-wellformed,hir:geo-transform-coherent}
\source{bierstone-milman08-functoriality,wlodarczyk05-hironaka}
\lean{TranscendentalNakamaye.Pending.NonmonomialTransformAlong.factor_along_along}
\leanok
In the situation of Lemma~\ref{hir:nonmonomial-step-along} (\(N\), \(E\), \(\mathfrak b\), \(\mathcal M\)), let \(o\ge1\), \(d\in\mathbb N\), and suppose
\(\operatorname{ord}_x(\mathfrak b|_N)\le o\) for every \(x\in\operatorname{cosupp}((\mathfrak b\mathcal M)|_N,d)\). Let \(T\) be a multiple test blow-up of the companion ideal
\(\underline{\mathcal O}\) of Lemma~\ref{hir:companion-reduction-along}. Then the transform \(\underline{\mathcal I}_T\) of \((M,N,E,\mathfrak b\mathcal M,d)\) is well formed; the
transform \(\mathfrak b_T\) of \((M,N,E,\mathfrak b,o)\) lives on the same \(N_T\) and \(E_T\), is coherent, does not vanish identically near any point of \(N_T\), and at the points of
\(N_T\) is contained, restricted to \(N_T\), in the ideal of no member of \(E_T\); and \(\mathfrak a_T|_{N_T}=(\mathfrak b_T\mathcal M_T)|_{N_T}\) with \(\mathcal M_T\) coherent, nonzero
and monomial along \(N_T\).
\end{lemma}
\begin{proof}
\leanok
Induction on the length of \(T\). By Lemma~\ref{hir:companion-reduction-along}, \(T\) is a multiple test blow-up of \(\underline{\mathcal I}\) and of \((M,N,E,\mathfrak b,o)\), so the
next centre \(C\) is admissible for both, and \(\operatorname{ord}(\mathfrak b_T|_{N_T})\le o\) on \(\operatorname{cosupp}\underline{\mathcal I}_T\supset C\). By induction
\(\underline{\mathcal I}_T\) is restricted-equivalent to \((M_T,N_T,E_T,\mathfrak b_T\mathcal M_T,d)\), so the two have the same cosupport and the same transform
(Lemma~\ref{hir:restricted-equivalence-basic}); Lemma~\ref{hir:nonmonomial-step-along} applies to \(\mathfrak b_T\), \(\mathcal M_T\) and \(C\), and its principal monomial factor is nonzero
and monomial along \(N_{T'}\) (Lemma~\ref{hir:monomial-invariant-along}). The transforms stay well formed (Lemma~\ref{bcalc:transform-wellformed}).
\end{proof}

\begin{lemma}[Uniqueness of the monomial factorization inside a submanifold]
\label{hir:monomial-factorization-unique-along}
\uses{hir:monomial-factorization-unique,hir:exponent-along,hir:monomial-part-along,hir:companion-reduction-along,hir:marked-ideal-operations,hir:restricted-calculus,bcalc:member-trace-ideal,ana:restricted-colon,ana:restricted-ideals}
\source{bierstone-milman08-functoriality}
\lean{TranscendentalNakamaye.Pending.NonmonomialTransformAlong.isDivNearAlong_eq_of_cosupport_eq,TranscendentalNakamaye.Pending.NonmonomialTransformAlong.divExpAlong_eq_of_isDivNearAlong_eq,TranscendentalNakamaye.Pending.NonmonomialTransformAlong.isDivNearAlong_mul_iff,TranscendentalNakamaye.Pending.NonmonomialTransformAlong.monomialPartAlong_eq_of_divExpAlong_eq,TranscendentalNakamaye.Pending.NonmonomialTransformAlong.monomialPartAlong_eq_of_factor,TranscendentalNakamaye.Pending.NonmonomialTransformAlong.exists_isInSpanOnAt_of_mul_generator,TranscendentalNakamaye.Pending.NonmonomialTransformAlong.le_nonmonomialPartAlong_of_factor,TranscendentalNakamaye.Pending.NonmonomialTransformAlong.restrictLe_nonmonomialPartAlong_of_factor}
\leanok
Let \(n\ge1\), \(N\subset M\) a closed submanifold, \(E\) a simple normal crossings divisor transverse to \(N\), \(\mathfrak a\) an ideal, \(\mathfrak b\) and \(\mathcal M\) coherent
ideals with \(\mathfrak a|_N=(\mathfrak b\mathcal M)|_N\), and suppose that at every point \(y\in N\) of a member \(H\), \(\mathfrak b_y|_N\not\subset(\mathcal I_H)_y|_N\). Then:
\begin{enumerate}
\item near every point of \(N\cap H\), the points of \(N\cap H\) lie in \(\operatorname{cosupp}(\mathfrak a|_N,e)\) if and only if they lie in \(\operatorname{cosupp}(\mathcal M|_N,e)\);
hence \(\mathfrak a\) and \(\mathcal M\) have the same exponents along \(N\) (Lemma~\ref{hir:exponent-along}) and the same monomial part along \(N\)
(Proposition~\ref{hir:monomial-part-along});
\item if moreover \(\mathcal M\) is nonzero and monomial along \(N\), then \(\mathfrak b\subset\mathcal N_N(\mathfrak a)\) and \(\mathcal N_N(\mathfrak a)|_N=\mathfrak b|_N\).
\end{enumerate}
For \(N=M\) this is Lemma~\ref{hir:monomial-factorization-unique}.
\end{lemma}
\begin{proof}
\leanok
(1) Exponents and monomial parts along \(N\) depend on \(\mathfrak a\) only through the cosupports \(\operatorname{cosupp}(\mathfrak a|_N,e)\), which are those of
\((\mathfrak b\mathcal M)|_N\). If the points of \(N\cap H\) near \(x\) lie in \(\operatorname{cosupp}(\mathcal M|_N,e)\), they lie in
\(\operatorname{cosupp}((\mathfrak b\mathcal M)|_N,e)\supset\operatorname{cosupp}(\mathfrak b|_N,0)\cap\operatorname{cosupp}(\mathcal M|_N,e)\) (Lemma~\ref{hir:marked-ideal-operations}).
Conversely, let \(x\in N\cap H\) and \(a\in\mathfrak b_x\) with \(a|_N\notin(\mathcal I_H)_x|_N\), i.e.\ \(a\) does not vanish on \(N\cap H\) near \(x\)
(Lemma~\ref{bcalc:member-trace-ideal}). At the points \(y\in N\cap H\) near \(x\) with \(a(y)\ne0\), \(\operatorname{ord}_y(\mathfrak b|_N)=0\), so near \(y\)
\(\operatorname{cosupp}((\mathfrak b\mathcal M)|_N,e)\subset\operatorname{cosupp}(\mathcal M|_N,e)\) (Lemma~\ref{hir:companion-reduction-along}, the paragraph on
products, with \(o=0\)); such points accumulate at \(x\), and the identity principle along \(N\cap H\) (Lemma~\ref{hir:exponent-along}) concludes.

(2) A germ \(g\) of \(\mathcal M_N(\mathfrak a)=\mathcal M_N(\mathcal M)\) restricts into \(\mathcal M|_N\) (\(\mathcal M\) is monomial along \(N\)), so for \(f\in\mathfrak b\),
\((gf)|_N\in(\mathfrak b\mathcal M)|_N=\mathfrak a|_N\): \(f\in\mathcal N_N(\mathfrak a)=(\mathfrak a|_N:\mathcal M_N(\mathfrak a))\). Conversely, at \(y\in N\), \(\mathcal M_N(\mathcal M)\) and
\(\mathcal M\) restrict to \((m)|_N\) for the local generator \(m\) of Proposition~\ref{hir:monomial-part-along}; for \(f\in\mathcal N_N(\mathfrak a)_y\),
\((mf)|_N\in\mathfrak a|_N=(\mathfrak b\,(m))|_N\), so \(mf=mh\) on \(N\) near \(y\) with \(h|_N\in\mathfrak b_y|_N\) (restricted spans,
Lemma~\ref{hir:restricted-calculus}); \(m\) vanishes identically on no nonempty open subset of \(N\), so \(f=h\) on \(N\) near \(y\)
(cancellation along \(N\), Lemma~\ref{ana:restricted-colon}).
\end{proof}

\begin{theorem}[Transforms of the monomial and nonmonomial parts, inside a submanifold]
\label{hir:nonmonomial-transform-along}
\uses{hir:nonmonomial-factor-along,hir:monomial-factorization-unique-along,hir:companion-reduction-along,hir:monomial-part-along,hir:monomial-invariant-along,hir:marked-equivalence-basic,hir:restricted-equivalence-basic,hir:wellformed-tower,hir:nonmonomial-transform-ambient}
\source{bierstone-milman08-functoriality,wlodarczyk05-hironaka}
\lean{TranscendentalNakamaye.Pending.NonmonomialTransformAlong.le_nonmonomialPartAlong_self,TranscendentalNakamaye.Pending.NonmonomialTransformAlong.isPrincipalMonomialAlongAt_monomialPartAlong,TranscendentalNakamaye.Pending.NonmonomialTransformAlong.parts_hypotheses,TranscendentalNakamaye.Pending.NonmonomialTransformAlong.nonmonomialPart_geoTransformAlong_along,TranscendentalNakamaye.Pending.NonmonomialTransformAlong.nonmonomialPart_geoTransformAlong_along_of_resolution}
\leanok
Let \(n\ge1\) and \(\underline{\mathcal I}=(M,N,E,\mathfrak a,d)\) be well formed with \(\mathfrak a|_N\) vanishing identically near no point of \(N\). Let
\(\mathcal N=\mathcal N_N(\mathfrak a)\), \(\mathcal M=\mathcal M_N(\mathfrak a)\) (Proposition~\ref{hir:monomial-part-along}), \(o\ge1\) with
\(\operatorname{ord}_x(\mathcal N|_N)\le o\) for every \(x\in\operatorname{cosupp}\underline{\mathcal I}\), and \(\underline{\mathcal O}\) the companion ideal of
Lemma~\ref{hir:companion-reduction-along} with \(\mathfrak b=\mathcal N\), \(\mathfrak c=\mathcal M\). Let \(T\) be a multiple test blow-up of \(\underline{\mathcal O}\). Then \(T\) is a
multiple test blow-up of \(\underline{\mathcal I}\); \(\underline{\mathcal I}_T\) is well formed; the transform \(\mathcal N_T\) of \((M,N,E,\mathcal N,o)\) lives on \(N_T\), is
contained in the nonmonomial part \(\mathcal N_{N_T}(\mathfrak a_T)\) and equal to it on \(N_T\); the monomial part \(\mathcal M_{N_T}(\mathfrak a_T)\) is that of a coherent ideal
which is nonzero and monomial along \(N_T\); and \(\operatorname{ord}_z(\mathcal N_{N_T}(\mathfrak a_T)|_{N_T})\le o\) for every \(z\in\operatorname{cosupp}\underline{\mathcal I}_T\).
If \(T\) resolves \(\underline{\mathcal O}\), then \(\operatorname{ord}_z(\mathcal N_{N_T}(\mathfrak a_T)|_{N_T})<o\) for every \(z\in\operatorname{cosupp}\underline{\mathcal I}_T\). For \(N=M\)
this is Theorem~\ref{hir:nonmonomial-transform-ambient}.
\end{theorem}
\begin{proof}
\leanok
\(\mathcal N\) and \(\mathcal M\) are coherent, \(\mathcal M\) is principal monomial at every point of \(N\) (Proposition~\ref{hir:monomial-part-along}), hence nonzero and monomial
along \(N\) (Lemma~\ref{hir:monomial-invariant-along}); \(\mathcal N\supset\mathfrak a\) is nonzero along \(N\) and in the ideal of no member on \(N\); and \(\underline{\mathcal I}\) is
restricted-equivalent to \((M,N,E,\mathcal N\mathcal M,d)\), so the two have the same multiple test blow-ups and, from the first blowup on, the same transforms
(Lemmas~\ref{hir:restricted-equivalence-basic} and~\ref{hir:marked-equivalence-basic}). Lemma~\ref{hir:nonmonomial-factor-along} gives
\(\mathfrak a_T|_{N_T}=(\mathcal N_T\mathcal M_T)|_{N_T}\) with \(\mathcal N_T\) in the ideal of no member of \(E_T\) on \(N_T\) and \(\mathcal M_T\) nonzero and monomial along \(N_T\), and
Lemma~\ref{hir:monomial-factorization-unique-along} identifies \(\mathcal M_{N_T}(\mathfrak a_T)=\mathcal M_{N_T}(\mathcal M_T)\) and \(\mathcal N_{N_T}(\mathfrak a_T)|_{N_T}=\mathcal N_T|_{N_T}\),
with \(\mathcal N_T\subset\mathcal N_{N_T}(\mathfrak a_T)\). The bound and its strict form after a resolution are those of Lemma~\ref{hir:companion-reduction-along}, read through
\(\mathcal N_T\subset\mathcal N_{N_T}(\mathfrak a_T)\) (so \(\operatorname{cosupp}(\mathcal N_{N_T}(\mathfrak a_T)|_{N_T},q)\subset\operatorname{cosupp}(\mathcal N_T|_{N_T},q)\)).
\end{proof}

\begin{lemma}[Transforms of the monomial and nonmonomial parts]
\label{hir:nonmonomial-transform}
\uses{hir:nonmonomial-transform-ambient,hir:nonmonomial-transform-along,hir:monomial-factorization-unique-along,hir:nonmonomial-step,hir:nonmonomial-step-along,hir:nonmonomial-factor-along,hir:monomial-factorization-unique,hir:companion-reduction,hir:companion-basic,hir:monomial-part,hir:divisor-exponent,bcalc:exceptional-pushforward,bcalc:controlled-not-exceptional,bcalc:monomial-lift,bcalc:exceptional-divisor,bcalc:strict-geo,hir:companion-ideal,bcalc:companion-open,bcalc:companion-rebase,bcalc:monomial-part-etale,hir:open-restriction-marked,hir:monomial-part-along,hir:exponent-along,ana:restricted-colon,oka:coherent-colon}
\source{bierstone-milman08-functoriality}
\lean{TranscendentalNakamaye.Pending.HironakaNonmonomialTransform.monomialPartAlong_le_colonAlong,TranscendentalNakamaye.Pending.HironakaNonmonomialTransform.monomialPartAlong_restrictLe_colonAlong,TranscendentalNakamaye.Pending.HironakaNonmonomialTransform.colonAlong_le_monomialPartAlong,TranscendentalNakamaye.Pending.HironakaNonmonomialTransform.colonAlong_restrictLe_monomialPartAlong,TranscendentalNakamaye.Pending.HironakaNonmonomialTransform.nonmonomial_transform_of_bound,TranscendentalNakamaye.Pending.HironakaNonmonomialTransform.nonmonomial_transform,TranscendentalNakamaye.Pending.HironakaNonmonomialTransform.nonmonomial_transform_of_resolution}
\leanok
Let \(T\) be a multiple test blow-up of the companion ideal \(\mathcal O(\underline{\mathcal I})\) on \(W\). Then \(T\) is a multiple test blow-up of
\(\underline{\mathcal I}|_W\), the nonmonomial part of \(\underline{\mathcal I}_T\) is the transform \(\mathcal N(\underline{\mathcal I})_T\) of
\((\mathcal N(\underline{\mathcal I}),o)\), and its monomial part is the quotient \((\mathfrak a_T|_{N_T}:\mathcal N(\underline{\mathcal I})_T)\), a product of powers of the
ideals of the members of \(E_T\). In particular, if \(T\) resolves \(\mathcal O(\underline{\mathcal I})\), then \(\operatorname{ord}\mathcal N(\underline{\mathcal I}_T)<o\) at
every point of \(\operatorname{cosupp}\underline{\mathcal I}_T\).
\end{lemma}
\begin{proof}
\leanok
\emph{Reduction to \(W=M\).} The companion ideal on \(W\) is the restriction to \(W\) of the companion ideal of \(\underline{\mathcal I}\), and it is the companion ideal of
\(\underline{\mathcal I}|_W\): the monomial and nonmonomial parts commute with restriction to open sets (Lemma~\ref{bcalc:monomial-part-etale},
Lemma~\ref{bcalc:companion-rebase}). \(\underline{\mathcal I}|_W\) is well formed and its ideal vanishes identically near no point of \(N\cap W\)
(Lemma~\ref{hir:open-restriction-marked}); its cosupport is \(\operatorname{cosupp}\underline{\mathcal I}\cap W\), where \(\operatorname{ord}\mathcal N(\underline{\mathcal I})\le o\)
(Lemma~\ref{bcalc:companion-open}(3)); and the transform of \((N,E,\mathcal N(\underline{\mathcal I}),o)|_W\) is that of \((N\cap W,E|_W,\mathcal N(\underline{\mathcal I}|_W),o)\).
So we may assume \(W=M\), with \(\operatorname{ord}_x\mathcal N\le o\) at every \(x\in\operatorname{cosupp}\underline{\mathcal I}\).

\emph{The parts.} One blowup \(\sigma\) along \(C\subset\operatorname{cosupp}(\mathcal N,o)\), \(\mathcal N=\mathcal N(\underline{\mathcal I})\); let \(D=\sigma^{-1}(C)\). By the Hadamard bridge
\(\sigma^*\mathcal N\subset\mathcal I_D^{o}\), and the transform is the controlled transform \(\mathcal N'=\sigma^c(\mathcal N,o)=\mathcal I_D^{-o}\sigma^*\mathcal N\).
It is contained in \(\mathcal I_D\) at no point \(z\in D\): otherwise \(\operatorname{ord}_{\sigma(z)}\mathcal N\ge o+1\) by Corollary~\ref{bcalc:controlled-not-exceptional}
(which rests on Theorem~\ref{bcalc:exceptional-pushforward}), while \(\operatorname{ord}\mathcal N\le o\) on \(C\) (Lemma~\ref{hir:companion-reduction}). It is
contained in \(\mathcal I_{H'}\) at no point \(z\) of the strict transform \(H'\) of a member \(H\): by Lemma~\ref{hir:divisor-exponent}, \(\mathcal N\) is contained in
\(\mathcal I_H\) at no point of \(H\) near \(\sigma(z)\), so the second clause of Corollary~\ref{bcalc:controlled-not-exceptional} gives a germ of \(\mathcal N'\) at \(z\) outside
\(\mathcal I_{H'}\) (the points of \(\sigma^{-1}(H\setminus C)\) accumulate at every point of \(H'\), Lemma~\ref{bcalc:strict-geo}). The transform of \(\mathfrak a\) is, on
\(N'\), \(\mathcal M'\mathcal N'\) with \(\mathcal M'\) the transform of the monomial factor, nonzero and monomial along \(N'\). By induction on the length of \(T\), with
Lemma~\ref{hir:companion-reduction} for the multiple test blow-ups and the bound, this is Theorem~\ref{hir:nonmonomial-transform-along}: \(T\) is a multiple test blow-up
of \(\underline{\mathcal I}\), the transform \(\mathcal N_T\) of \((\mathcal N,o)\) is contained in \(\mathcal N(\underline{\mathcal I}_T)\) and equal to it on \(N_T\), the monomial
part \(\mathcal M(\underline{\mathcal I}_T)\) is that of a monomial ideal \(\mathcal M_T\) (exponents are additive in the factorial local rings of \(N_T\), so the factorization of
Definition~\ref{hir:monomial-part} is unique, Lemma~\ref{hir:monomial-factorization-unique-along}), hence principal monomial at every point of \(N_T\): generated on \(N_T\)
by a product of powers of local equations of the members of \(E_T\) (Proposition~\ref{hir:monomial-part-along}); and, if \(T\) resolves the companion ideal,
\(\operatorname{ord}\mathcal N(\underline{\mathcal I}_T)<o\) on \(\operatorname{cosupp}\underline{\mathcal I}_T\).

\emph{The monomial part is the quotient.} Let \(\mathfrak c=(\mathfrak a_T|_{N_T}:\mathcal N_T)\) (Lemma~\ref{ana:restricted-colon}). If \(g\) is a germ of
\(\mathcal M(\underline{\mathcal I}_T)\) and \(h\) one of \(\mathcal N_T\subset\mathcal N(\underline{\mathcal I}_T)=(\mathfrak a_T|_{N_T}:\mathcal M(\underline{\mathcal I}_T))\), then
\(gh\in\mathfrak a_T|_{N_T}\); so \(\mathcal M(\underline{\mathcal I}_T)\subset\mathfrak c\). Conversely \(\mathfrak c\) is coherent, as the colon of two coherent ideals
(Lemma~\ref{oka:coherent-colon}), and \(\mathcal N_T\mathfrak c\subset\mathfrak a_T\) on \(N_T\). Let \(H\) be a member of \(E_T\) through \(x\in N_T\) and \(e=e_H(x)\) the
exponent of \(\mathfrak a_T\) along \(N_T\) (Lemma~\ref{hir:exponent-along}): near \(x\) the points of \(H\cap N_T\) lie in \(\operatorname{cosupp}(\mathfrak a_T|_{N_T},e)\), hence in
\(\operatorname{cosupp}((\mathcal N_T\mathfrak c)|_{N_T},e)\). As \(\mathcal N_T\) is contained in \(\mathcal I_H\) on \(N_T\) at no point of \(H\cap N_T\), at the points of
\(H\cap N_T\) where a germ of \(\mathcal N_T\) does not vanish the factor \(\mathcal N_T\) has order \(0\), so these points lie in \(\operatorname{cosupp}(\mathfrak c|_{N_T},e)\); they
accumulate at every point of \(H\cap N_T\) near \(x\), and by the identity principle along \(H\cap N_T\) (Lemma~\ref{hir:exponent-along}) every germ of \(\mathfrak c_x\) is, on
\(N_T\), a germ of \(\mathcal I_H^{e}\) (compose with a retraction onto \(N_T\) preserving the members). So \(\mathfrak c\subset\mathcal M(\underline{\mathcal I}_T)\), and the two
have the same restriction to \(N_T\): the monomial part is the quotient \((\mathfrak a_T|_{N_T}:\mathcal N(\underline{\mathcal I})_T)\).
\end{proof}

\begin{lemma}[Monomial marked ideals under the blowup of an intersection of members]
\label{hir:monomial-blowup}
\uses{bcalc:monomial-lift,bcalc:snc-chart,bcalc:snc-transform-nc,bcalc:member-pullback,hir:geo-transform,hir:monomial-part,bcalc:monomial-chart,bcalc:monomial-transform,hir:monomial-invariant-geometry,bcalc:monomial-transform-coords,hir:monomial-step,hir:monomial-part-along,hir:exponent-along,bcalc:monomial-transform-along-coords,bcalc:monomial-transform-along-off-centre,bcalc:member-intersection-along,hir:monomial-invariant-along,hir:monomial-step-along,bcalc:boundary-count-drop-along,hir:open-restriction-marked,hir:divisor-exponent,hir:monomial-invariance-along}
\source{bierstone-milman08-functoriality,wlodarczyk05-hironaka}
\lean{TranscendentalNakamaye.Pending.HironakaMonomialBlowup.filter_eventually_subset_inter_memberInter_eq,TranscendentalNakamaye.Pending.HironakaMonomialBlowup.eventually_mem_cosupp_of_le_sum,TranscendentalNakamaye.Pending.HironakaMonomialBlowup.monomial_blowup_of_isAdmissibleCentre,TranscendentalNakamaye.Pending.HironakaMonomialBlowup.divExpAlong_pullback_val_get,TranscendentalNakamaye.Pending.HironakaMonomialBlowup.monomial_blowup}
\leanok
Let \(\underline{\mathcal I}\) be well formed and monomial, \(x\in N\), and near \(x\) let \(\mathfrak a|_N=\prod_{H\in S}\mathcal I_{H\cap N}^{a_H}\), \(S\) the set of members
through \(x\). Let \(J\subset S\) with \(\sum_{H\in J}a_H\ge d\) and \(C=N\cap\bigcap_{H\in J}H\) near \(x\). Then \(C\) is an admissible centre near \(x\), and near
every point \(z\) over \(x\) of the blowup along \(C\) the transform is monomial:
\[
 \mathfrak a'|_{N'}=\mathcal I_{D\cap N'}^{\,a_D}\prod_{H\in S(z)}\mathcal I_{H'\cap N'}^{a_H},\qquad a_D=\sum_{H\in J}a_H-d,
\]
where \(D=\sigma^{-1}(C)\) and \(S(z)\subset S\) is the set of members whose strict transforms pass through \(z\); \(S(z)\) does not contain \(J\).
\end{lemma}
\begin{proof}
\leanok
In the coordinates of Lemma~\ref{hir:divisor-exponent} at \(x\), adapted to \(N\) and \(E\), the members through \(x\) are distinct coordinate hyperplanes
\(\{z_H=0\}\) with \(z_H\) a coordinate of \(N\), and \(C=N\cap\{z_j=0,\ j\in J\}\) is a coordinate subspace of \(N\): a closed submanifold with only normal crossings
with \(E\) (Lemma~\ref{bcalc:member-intersection-along}). The members containing \(C\) near \(x\) are exactly those of \(J\): moving along \(z_H\) for \(H\in S\setminus J\)
stays in \(N\) and on the members of \(J\) and leaves \(H\). Near \(x\) the exponents along the members of \(J\) are those at \(x\)
(Lemma~\ref{hir:exponent-along}), and a point \(y\) of \(N\) lies in the cosupport of the monomial marked ideal iff the exponents of the members through \(y\) add up
to at least \(d\) (Lemma~\ref{hir:monomial-invariant-along}); at \(y\in C\) near \(x\) the members of \(J\) pass through \(y\), so \(\sum_{H\in J}a_H\ge d\) puts \(y\) in the
cosupport. Hence \(C\) is an admissible centre of the restriction of \(\underline{\mathcal I}\) to a neighbourhood \(U\) of \(x\)
(Lemma~\ref{hir:open-restriction-marked}; the restriction is well formed, monomial, with the same exponents, Lemma~\ref{hir:monomial-invariance-along}).
Let \(z\in N'\) lie over \(x\). In the monomial chart in which \(z_{j_0}\) (\(j_0\in J\)) generates the exceptional ideal, \(z_j=z_{j_0}w_j\) for
\(j\in J\setminus\{j_0\}\), so \(\sigma^*\prod_Sz_H^{a_H}=z_{j_0}^{\sum_Ja}\prod_{J\setminus\{j_0\}}w_j^{a_j}\prod_{S\setminus J}z_H^{a_H}\); dividing by \(z_{j_0}^d\) gives the
formula, with \(D=\{z_{j_0}=0\}\), \(H'_j=\{w_j=0\}\) for \(j\in J\setminus\{j_0\}\), and \(H'_{j_0}\) not meeting the chart. In coordinates at \(z\) adapted to \(N'\) and
\(E'\) this is Corollary~\ref{bcalc:monomial-transform-along-coords}, fed with the monomial generator of \(\mathfrak a|_N\) at \(x\)
(Proposition~\ref{hir:monomial-part-along}(6)) and with \(\sum_{H\supset C}a_H=\sum_{H\in J}a_H\ge d\). Finally the strict transforms of the members of \(J\) do not
all pass through \(z\) (Theorem~\ref{bcalc:boundary-count-drop-along}), so \(S(z)\) does not contain \(J\).
\end{proof}

\begin{lemma}[The invariant of monomial exponent data]
\label{hir:monomial-invariant}
\uses{hir:invariant-order}
\source{wlodarczyk05-hironaka,bierstone-milman03-desingularization-algorithms,bierstone-milman08-functoriality}
\lean{TranscendentalNakamaye.Pending.MonomialInvariant.Val,TranscendentalNakamaye.Pending.MonomialInvariant.key,TranscendentalNakamaye.Pending.MonomialInvariant.feasible,TranscendentalNakamaye.Pending.MonomialInvariant.inv,TranscendentalNakamaye.Pending.MonomialInvariant.mem_feasible,TranscendentalNakamaye.Pending.MonomialInvariant.key_injective,TranscendentalNakamaye.Pending.MonomialInvariant.key_lt_key_iff,TranscendentalNakamaye.Pending.MonomialInvariant.feasible_nonempty_iff,TranscendentalNakamaye.Pending.MonomialInvariant.inv_eq_bot_iff,TranscendentalNakamaye.Pending.MonomialInvariant.exists_top,TranscendentalNakamaye.Pending.MonomialInvariant.eq_of_inv_eq_key,TranscendentalNakamaye.Pending.MonomialInvariant.card_le_of_top,TranscendentalNakamaye.Pending.MonomialInvariant.sum_erase_lt_of_top,TranscendentalNakamaye.Pending.MonomialInvariant.inv_mono,TranscendentalNakamaye.Pending.MonomialInvariant.inv_eq_iff_subset,TranscendentalNakamaye.Pending.MonomialInvariant.key_lt_of_blowup,TranscendentalNakamaye.Pending.MonomialInvariant.inv_lt_of_blowup,TranscendentalNakamaye.Pending.MonomialInvariant.wellFoundedLT_val,TranscendentalNakamaye.Pending.MonomialInvariant.exists_eq_bot_of_lt_succ}
\leanok
Let \(n,d\in\mathbb N\), \(S\subset\mathbb N\) finite with \(|S|\le n\) (indices of the members through a point) and \(a\colon\mathbb N\to\mathbb N\) (exponents). Call
\(K\subset S\) \emph{feasible} if \(\sum_Ka\ge d\), give it the key \(\kappa(K)=(n-|K|,\sum_Ka,K)\), ordered lexicographically with finite subsets of \(\mathbb N\) ordered
colexicographically, and let \(\gamma(S,a)\) be the largest key of a feasible set (\(\bot\) if there is none). Then:
\begin{enumerate}
\item \(\gamma(S,a)=\bot\) if and only if \(\sum_Sa<d\); otherwise \(\gamma(S,a)=\kappa(I)\) for a unique feasible \(I\) (the \emph{top set}), and \(I\) has minimal
cardinality among the feasible sets, so \(\sum_{I\setminus\{j\}}a<d\) for \(j\in I\);
\item \(S'\subset S\) implies \(\gamma(S',a)\le\gamma(S,a)\), with equality if and only if \(I\subset S'\);
\item if \(j_0\in I\), \(S'\subset(S\setminus\{j_0\})\cup\{m\}\), \(a'=a\) on \(S\) and \(a'_m+d=\sum_Ia\), then \(\gamma(S',a')<\gamma(S,a)\);
\item the values are well-founded, so a sequence of values that strictly decreases as long as it is \(\ne\bot\) reaches \(\bot\).
\end{enumerate}
\end{lemma}
\begin{proof}
\leanok
(1) \(S\) itself is feasible when \(\sum_Sa\ge d\); the key determines the set (third entry). If a feasible \(K\) had \(|K|<|I|\le n\), its first entry would be
larger. (2) Feasible sets of \(S'\) are feasible sets of \(S\); if \(\gamma(S',a)=\kappa(I)\), some feasible \(K\subset S'\) has \(\kappa(K)=\kappa(I)\), so \(K=I\).
(3) Let \(K'\) be feasible for \((S',a')\). If \(m\in K'\), then \(K=(K'\setminus\{m\})\cup\{j_0\}\subset S\) has \(|K|=|K'|\) and
\(\sum_Ka=a_{j_0}+\sum_{K'\setminus\{m\}}a>a'_m+\sum_{K'\setminus\{m\}}a'=\sum_{K'}a'\ge d\), because \(a'_m=\sum_Ia-d<a_{j_0}\) by (1); so \(K\) is feasible and
\(\kappa'(K')<\kappa(K)\le\kappa(I)\). If \(m\notin K'\), then \(K'\subset S\setminus\{j_0\}\) is feasible for \((S,a)\) with the same key, and \(K'\ne I\).
(4) Colex on finite subsets of \(\mathbb N\) is order-isomorphic to \(\mathbb N\) (binary expansion); lexicographic products and \(\{\bot\}\cup\) of well-founded
orders are well-founded; Lemma~\ref{hir:invariant-order}.
\end{proof}

\begin{lemma}[The invariant of a monomial marked ideal at one stage]
\label{hir:monomial-invariant-geometry}
\uses{hir:monomial-invariant,hir:exponent-locally-constant,hir:monomial-part-ambient,hir:monomial-factorization,hir:snc,hir:cosupport,hir:marked-ideal,hir:transform,hir:order-basic,hir:order-product,hir:order-taylor,hir:giraud-support,hir:giraud-along,hir:normal-power-pullback,hir:companion-maximal-order,hir:invariant-order,hir:resolution-driver,bcalc:hyperplane-ideal,bcalc:member-intersection,bex:hadamard,lag:regular-point-basics}
\source{wlodarczyk05-hironaka,bierstone-milman08-functoriality}
\lean{TranscendentalNakamaye.Pending.MonomialCase.feasible_congr,TranscendentalNakamaye.Pending.MonomialCase.inv_congr,TranscendentalNakamaye.Pending.MonomialCase.memberIdx,TranscendentalNakamaye.Pending.MonomialCase.mem_memberIdx,TranscendentalNakamaye.Pending.MonomialCase.card_memberIdx_le,TranscendentalNakamaye.Pending.MonomialCase.eventually_memberIdx_subset,TranscendentalNakamaye.Pending.MonomialCase.expData,TranscendentalNakamaye.Pending.MonomialCase.gamma,TranscendentalNakamaye.Pending.MonomialCase.eventually_expData_eq,TranscendentalNakamaye.Pending.MonomialCase.eventually_gamma_le,TranscendentalNakamaye.Pending.MonomialCase.eventually_gamma_eq_iff,TranscendentalNakamaye.Pending.MonomialCase.flatCoord_notMem_idealPow_two,TranscendentalNakamaye.Pending.MonomialCase.prod_flatCoord_pow_order,TranscendentalNakamaye.Pending.MonomialCase.exists_snc_chart,TranscendentalNakamaye.Pending.MonomialCase.mem_cosupport_iff_of_monomial,TranscendentalNakamaye.Pending.MonomialCase.gamma_eq_bot_iff_of_monomial,TranscendentalNakamaye.Pending.MonomialCase.topSet,TranscendentalNakamaye.Pending.MonomialCase.topSet_key,TranscendentalNakamaye.Pending.MonomialCase.topIdx,TranscendentalNakamaye.Pending.MonomialCase.top_of_gamma_eq,TranscendentalNakamaye.Pending.MonomialCase.mem_memberInter_topIdx_iff,TranscendentalNakamaye.Pending.MonomialCase.eventually_gamma_eq_iff_mem_memberInter,TranscendentalNakamaye.Pending.MonomialCase.mem_memberInter_of_gamma_eq,TranscendentalNakamaye.Pending.MonomialCase.isInvUSCOn_gamma,TranscendentalNakamaye.Pending.MonomialCase.isClosed_gamma_eq,TranscendentalNakamaye.Pending.MonomialCase.isCoordinateFamilyAt_of_inter_eq,TranscendentalNakamaye.Pending.MonomialCase.isAdmissibleCentre_gamma_eq,TranscendentalNakamaye.Pending.MonomialCase.interior_get_eq_empty,TranscendentalNakamaye.Pending.MonomialCase.isAdmissibleCentre_maxLocus_gamma,TranscendentalNakamaye.Pending.MonomialCase.holomorphicAt_flatMonomial,TranscendentalNakamaye.Pending.MonomialCase.eventually_isNonzeroAt_of_principal,TranscendentalNakamaye.Pending.MonomialCase.flatMonomial_mem_idealPow_vanishingIdeal,TranscendentalNakamaye.Pending.MonomialCase.divExp_eq_of_principal,TranscendentalNakamaye.Pending.MonomialCase.eventually_monomialPart_le_of_principal,TranscendentalNakamaye.Pending.MonomialCase.gamma_eq_of_principal,TranscendentalNakamaye.Pending.MonomialCase.eventually_isInSpanAt_iff_of_isCoherent}
\leanok
Let \(n\ge1\), \(E=(E_0,\dots,E_{r-1})\) a simple normal crossings divisor on \(M\), \(\mathfrak a\) a coherent ideal such that \(\mathfrak a_y\) contains a germ not
vanishing near \(y\) for every \(y\), and \(d\in\mathbb N\). For \(x\in M\) let \(S(x)=\{i:x\in E_i\}\), \(a_x(i)=e_{E_i}(x)\) the exponents of
Lemma~\ref{hir:exponent-locally-constant}, and \(\gamma(x)=\gamma(S(x),a_x)\) the invariant of Lemma~\ref{hir:monomial-invariant} (for the bound \(n\) and the mark \(d\)). Then:
\begin{enumerate}
\item \(|S(x)|\le n\), and \(\gamma(S,a)\) depends only on the values of \(a\) on \(S\);
\item (upper semicontinuity) near \(x\), \(S(y)\subset S(x)\) and \(a_y=a_x\) on \(S(y)\); hence \(\gamma(y)\le\gamma(x)\) near \(x\);
\item (the locus of a value near a point) if \(\gamma(x)=v\ne\bot\) with top set \(I\), then near \(x\): \(\gamma(y)=v\) if and only if \(y\in\bigcap_{i\in I}E_i\);
\item (the cosupport) if \(\mathfrak a\) is monomial (\(\mathcal M(\mathfrak a)\subset\mathfrak a\), Proposition~\ref{hir:monomial-part-ambient}; this item needs no
non-vanishing hypothesis), then \(x\in\operatorname{cosupp}(\mathfrak a,d)\) if and only if \(\sum_{i\in S(x)}a_x(i)\ge d\), i.e.\ \(\gamma(x)\ne\bot\);
\item (the maximum locus) if moreover \(\underline{\mathcal I}=(M,M,E,\mathfrak a,d)\), the set of points of \(\operatorname{cosupp}\underline{\mathcal I}\) at which \(\gamma\)
attains its maximum on \(\operatorname{cosupp}\underline{\mathcal I}\) is an admissible centre, with empty interior if \(d\ge1\); more generally, if \(\gamma\le v\) on \(M\), the
locus \(\{\gamma=v\}\) is an admissible centre;
\item (principal monomial ideals) if near \(x\) the ideal \(\mathfrak a\) is generated by \(\prod_{i\in S(x)}u_{\sigma(i)}^{b_i}\) in coordinates \(u\) centred at \(x\)
in which the members \(E_i\), \(i\in S(x)\), are the hyperplanes \(\{u_{\sigma(i)}=0\}\) (\(\sigma\) injective), then \(\mathfrak a_y\) contains a germ not vanishing near \(y\)
for \(y\) near \(x\), \(e_{E_i}(x)=b_i\) for \(i\in S(x)\), so \(\gamma(x)=\gamma(S(x),b)\), and \(\mathcal M(\mathfrak a)\subset\mathfrak a\) near \(x\) (for \(\mathfrak a\) coherent). For \(\mathfrak a\) coherent it suffices that \(\mathfrak a_x\) be generated by the
monomial: a stalk generator of a coherent ideal generates the nearby stalks.
\end{enumerate}
\end{lemma}
\begin{proof}
\leanok
(1) The members through \(x\) are distinct coordinate hyperplanes. (2) The members are closed, and \(e_{E_i}\) is locally constant along \(E_i\)
(Lemma~\ref{hir:exponent-locally-constant}(4)); then Lemma~\ref{hir:monomial-invariant}(2). (3) Lemma~\ref{hir:monomial-invariant}(2) with \(S(y)\subset S(x)\);
the third entry of \(v\) is \(I\). (4) Take simple normal crossings coordinates \(u\) at \(x\), the members through \(x\) being \(\{u_{\sigma(i)}=0\}\). The monomial
\(m=\prod_{i\in S(x)}u_{\sigma(i)}^{a_x(i)}\) lies in \(\mathcal M(\mathfrak a)_x\subset\mathfrak a_x\) and has order exactly \(\sum a_x(i)\) at \(x\) (a coordinate has order
one, Lemma~\ref{hir:companion-maximal-order}, and orders add, Lemma~\ref{hir:order-product}); so \(x\in\operatorname{cosupp}\) forces \(\sum a_x(i)\ge d\). Conversely
every germ of \(\mathfrak a_x\subset\mathcal M(\mathfrak a)_x\) is a multiple of \(m\) (Lemma~\ref{hir:monomial-factorization}(1)), so has order \(\ge\sum a_x(i)\).
(5) Let \(x_0\) be a point of the maximum locus and \(v=\gamma(x_0)\ne\bot\) (by (4)); then \(\gamma\le v\) on \(M\) (off the cosupport \(\gamma=\bot\)) and the
maximum locus is \(\{\gamma=v\}\). This set is closed by (2), lies in the cosupport by (4), and by (3) agrees near each of its points with \(\bigcap_{i\in I}E_i\),
a closed submanifold with only normal crossings with \(E\) (Lemma~\ref{bcalc:member-intersection}); both properties are local
(Lemma~\ref{lag:regular-point-basics} for regularity; the same shrinking of the coordinate chart for adapted coordinates). If \(d\ge1\), \(I\ne\emptyset\), so the
locus lies in a member, which has empty interior (a coordinate is not in \(\mathfrak m_x^2\)).
(6) Let \(m=\prod u_{\sigma(i)}^{b_i}\). The points off all the hyperplanes accumulate at every point near \(x\), and \(m\) does not vanish there. The
monomial lies in \((\mathcal I_{E_i}^{b_i})_x\), so \(\mathfrak a\) is divisible by \(\mathcal I_{E_i}^{b_i}\) near \(x\); at the points of \(E_i\) near \(x\) off the other
hyperplanes \(m\) has order exactly \(b_i\) (the other factors are units, \(u_{\sigma(i)}\notin\mathfrak m^2\), orders add), so \(\mathfrak a\) is not divisible by
\(\mathcal I_{E_i}^{b_i+1}\) near \(x\); hence \(e_{E_i}(x)=b_i\) (Lemma~\ref{hir:exponent-locally-constant}(2)). Near \(x\) the exponents are locally constant
(Lemma~\ref{hir:exponent-locally-constant}(4)); a germ of \(\mathcal M(\mathfrak a)_y\) is divisible by \(u_{\sigma(i)}^{b_i}\) for \(E_i\ni y\) and trivially for the other
\(i\in S(x)\) (\(u_{\sigma(i)}\) is a unit at \(y\)), hence by \(m\) (Lemma~\ref{hir:monomial-factorization}(1)), so it lies in \(\mathfrak a_y=(m)_y\).
For the last sentence: local generators \(s_k\) of \(\mathfrak a\) near \(x\) are multiples \(c_km\) at \(x\), hence near \(x\), and \(m\in\mathfrak a_y\) for \(y\) near \(x\)
(coherence).
\end{proof}

\begin{lemma}[Resolution of monomial marked ideals]
\label{hir:monomial-case}
\uses{hir:monomial-invariance,hir:monomial-blowup,hir:monomial-part,hir:invariant-order,hir:monomial-invariant,hir:monomial-invariant-geometry,bcalc:boundary-count-drop,hir:monomial-step,hir:monomial-resolution,hir:exponent-along,hir:monomial-part-along,bcalc:boundary-count-drop-along,bcalc:member-intersection-along,hir:monomial-invariant-along,hir:monomial-step-along,hir:monomial-resolution-along,hir:monomial-invariance-along}
\source{bierstone-milman08-functoriality,wlodarczyk05-hironaka,bierstone-milman03-desingularization-algorithms}
\notready
Let \(\underline{\mathcal I}\) be well formed and monomial, \(d\ge1\). For \(x\in\operatorname{cosupp}\underline{\mathcal I}\) with exponents \(a_H\), \(H\in S(x)\), put
\(\gamma_1(x)=\min\{|I|:I\subset S(x),\ \sum_Ia_H\ge d\}\), \(\gamma_2(x)=\max\{\sum_Ia_H:I\subset S(x),\ |I|=\gamma_1(x)\}\), let \(\gamma_3(x)\) be the largest
\(I\subset S(x)\) with \(|I|=\gamma_1(x)\) and \(\sum_Ia_H=\gamma_2(x)\) (sets of members compared colexicographically through their indices in the
sequence \(E\); the exceptional member of a blowup is appended, so it has the largest index), and \(\gamma(x)=(-\gamma_1(x),\gamma_2(x),\gamma_3(x))\), ordered
lexicographically. With \(n-\gamma_1\) in place of \(-\gamma_1\) this is the invariant \(\gamma(S(x),a)\) of Lemma~\ref{hir:monomial-invariant}. Then:
\begin{enumerate}
\item \(\gamma\) is upper semicontinuous on \(\operatorname{cosupp}\underline{\mathcal I}\), and on every open set on which \(\gamma\le\lambda\) the locus \(\{\gamma=\lambda\}\) is
\(N\cap\bigcap_{H\in\lambda_3}H\), an admissible centre;
\item after blowing up the maximum locus \(C\), \(\gamma(z)<\gamma(\sigma(z))\) at every point \(z\) of the new cosupport over \(C\), and
\(\gamma(z)=\gamma(\sigma(z))\) elsewhere;
\item if \(\operatorname{cosupp}\underline{\mathcal I}\) is compact, the iterated blowup of the maximum locus of \(\gamma\) reaches an empty cosupport after
finitely many steps;
\item \(\gamma\) is unchanged by restriction to open sets, by biholomorphisms and by steps with empty centre.
\end{enumerate}
\end{lemma}
\begin{proof}
(1) Near \(x\) the members through \(y\) are among those through \(x\), with the same exponents, so every subset available at \(y\) is available at \(x\);
hence \(\gamma(y)\le\gamma(x)\) (Lemma~\ref{hir:monomial-invariant}(2)). If \(\gamma\le\lambda=\gamma(x)\) near \(x\), a point \(y\) near \(x\) lying on all members of \(I=\lambda_3\) has \(\gamma(y)\ge\gamma(x)\),
so \(\gamma(y)=\lambda\); conversely \(\gamma(y)=\lambda\) gives \(\gamma_3(y)=I\), so \(y\) lies on the members of \(I\). The locus is admissible by
Lemma~\ref{hir:monomial-blowup}. For \(N=M\), (1) is proved in Lean: Lemma~\ref{hir:monomial-invariant-geometry}(2), (3) and (5) (the maximum locus on the
cosupport is an admissible centre with empty interior, in the form used by the driver, Proposition~\ref{hir:resolution-driver}); for \(N\ne M\), (1) is
proved in Lean: Lemma~\ref{hir:monomial-invariant-along}(1), (3).
(2) The combinatorial part is Lemma~\ref{hir:monomial-invariant}(3), with the exponent data of Lemma~\ref{hir:monomial-blowup}, read off its generator
by Lemma~\ref{hir:monomial-invariant-geometry}(6) (proved); a member \(H_{j_0}\), \(j_0\in I\), whose strict transform misses \(z\) exists by
Lemma~\ref{bcalc:boundary-count-drop} (TN3, proved; along \(N\): Theorem~\ref{bcalc:boundary-count-drop-along}, proved). For \(N=M\), (2) is proved in Lean: Lemma~\ref{hir:monomial-step}(2), (4), (5); for \(N\ne M\): Lemma~\ref{hir:monomial-step-along}(3)--(5). In words:
let \(z\) lie over \(x\in C\), \(I=\gamma_3(x)\), in the chart of \(j_0\in I\). By Lemma~\ref{hir:monomial-blowup} the exponent data at \(z\) arise from those at
\(x\) by deleting members, among them \(H_{j_0}\), and adding \(D\) with \(a_D=\sum_Ia-d<a_{j_0}\) (every proper subset of \(I\) has fewer than \(\gamma_1(x)\)
elements, hence sum \(<d\)). A subset at \(z\) containing \(D\) has sum smaller than the subset at \(x\) obtained by replacing \(D\) by \(H_{j_0}\); a subset not
containing \(D\) is a subset at \(x\) different from \(I\). Hence \(\gamma_1(z)\ge\gamma_1(x)\); if equal, \(\gamma_2(z)\le\gamma_2(x)\), with equality only for subsets
not containing \(D\) and different from \(I\), which are smaller than \(I\) in the order used for \(\gamma_3\). Off \(C\) the blowup is biholomorphic and the data agree.
(3) After the blowup of the maximum locus \(C=\{\gamma=\lambda\}\) of a compact cosupport, \(\gamma<\lambda\) at every point of the new cosupport: over
\(C\) by (2), elsewhere because \(\gamma(z)=\gamma(\sigma(z))<\lambda\). So the maximum strictly decreases, and the values are well-founded
(Lemma~\ref{hir:monomial-invariant}(4)); this is the formal route. For \(N=M\), (3) is proved in Lean: Theorem~\ref{hir:monomial-resolution} (over a
compact base; the monomial form and the non-vanishing persist along the canonical segments, Lemma~\ref{hir:monomial-step}(5)); for \(N\ne M\):
Theorem~\ref{hir:monomial-resolution-along}; (3) for a compact cosupport in a noncompact \(M\), as stated: Lemma~\ref{hir:monomial-case-local}. The original count (for the lexicographic order of the source, which is not well-founded
on its own): the values of \((-\gamma_1,\gamma_2)\) lie in a finite set, because exponents never increase (\(a_D<a_{j_0}\)) and at most \(n\) members pass through a
point; by (2) the maximum \(\lambda\) of \((-\gamma_1,\gamma_2)\) does not increase. While it stays equal to \(\lambda\), let \(\Sigma_\lambda\) be the finite set of
subsets \(I\) of members with \(|I|=-\lambda_1\), \(\sum_Ia=\lambda_2\) and \(N\cap\bigcap_IH\cap\operatorname{cosupp}\ne\emptyset\). Each step blows up
\(N\cap\bigcap_IH\) for the largest \(I\in\Sigma_\lambda\) and removes \(I\) from \(\Sigma_\lambda\) (the strict transforms of its members no longer meet over the
centre); by the computation in (2) no subset containing a new exceptional member enters \(\Sigma_\lambda\) (it would have sum smaller than a subset at level
\(>\lambda\)). So \(\lambda\) drops after at most \(|\Sigma_\lambda|\) steps, and \(\lambda\) can drop only finitely often.
(4) \(\gamma(x)\) depends only on the exponents of the members through \(x\) and on their relative order in \(E\); a step with empty centre appends a
member through no point. For \(N=M\), (4) is proved in Lean: Lemma~\ref{hir:monomial-invariance}; for \(N\ne M\): Lemma~\ref{hir:monomial-invariance-along}.
\end{proof}
The sources choose the centre as the largest minimal subset (Włodarczyk's \(\rho\), Bierstone--Milman's \(J\)) and prove termination through the decrease of
the exponent sum \(\mu\) over the centre (Bierstone--Milman 2003, Theorem~1.14 and its proof in \S3.3); the blueprint first minimizes the number of members
and maximizes their exponent sum, which gives the termination argument (3) directly.

\emph{Superseded.} As stated, the second half of (1) is false when it is read globally on the open set. Members may be disconnected, and their
exponents are only locally constant (Lemma~\ref{hir:divisor-exponent}). Take \(M=N=\mathbf C^2\), \(H_1=\{z_1=0\}\sqcup\{z_1=1\}\), \(H_2=\{z_2=0\}\),
\(\mathfrak a=(z_1^2(z_1-1)z_2^2)\) and \(d=3\). The cosupport is \(\{(0,0),(1,0)\}\) and \(\gamma\le\gamma(0,0)\) everywhere. But \((1,0)\in H_1\cap H_2\) has
\(\gamma_2(1,0)=3<4=\gamma_2(0,0)\), so the locus \(\{\gamma=\gamma(0,0)\}=\{(0,0)\}\) is not \(H_1\cap H_2\). The proof of (1) shows only the local statement.
Lemma~\ref{hir:monomial-case-local} is the corrected statement, and it is proved; the statement above is kept unchanged.

\begin{lemma}[Resolution of monomial marked ideals, corrected form]
\label{hir:monomial-case-local}
\uses{hir:monomial-invariant,hir:monomial-invariant-along,hir:monomial-step-along,hir:monomial-resolution-along,hir:monomial-invariance-along,hir:monomial-step,hir:monomial-resolution,hir:monomial-part-along,hir:exponent-along,hir:resolution-driver,hir:invariant-order,hir:open-restriction-marked,hir:giraud-along,bcalc:transform-off-centre,bcalc:member-intersection-along,hir:principalization-local}
\source{bierstone-milman08-functoriality,wlodarczyk05-hironaka,bierstone-milman03-desingularization-algorithms}
\lean{TranscendentalNakamaye.Pending.MonomialCaseAlong.isInvUSCOn_gammaAlong,TranscendentalNakamaye.Pending.MonomialCaseAlong.eventually_gammaAlong_eq_iff_mem_memberInter,TranscendentalNakamaye.Pending.MonomialCaseAlong.isAdmissibleCentre_gammaAlong_eq,TranscendentalNakamaye.Pending.MonomialCaseLocal.isMonomialAlong_pullback_val_of_isMonomialAlong,TranscendentalNakamaye.Pending.MonomialCaseLocal.isNonzeroAlongAt_pullback_val,TranscendentalNakamaye.Pending.MonomialCaseLocal.isAdmissibleCentre_gammaAlong_level,TranscendentalNakamaye.Pending.MonomialStepAlong.gammaAlong_geoTransform_lt_of_mem,TranscendentalNakamaye.Pending.MonomialStepAlong.gammaAlong_geoTransform_eq_of_notMem,TranscendentalNakamaye.Pending.MonomialStepAlong.monomial_step_along,TranscendentalNakamaye.Pending.MonomialCaseLocal.isCompact_cosupp_geoTransformAlong,TranscendentalNakamaye.Pending.MonomialCaseLocal.exists_isGeoResolution_of_wellFounded_invariant_of_isCompact,TranscendentalNakamaye.Pending.MonomialCaseLocal.exists_isGeoResolution_of_monomialAlong_of_isCompact,TranscendentalNakamaye.Pending.MonomialCaseLocal.exists_isGeoResolution_of_monomial_of_isCompact,TranscendentalNakamaye.Pending.MonomialInvariantFunctorialAlong.gammaAlong_pullback_val,TranscendentalNakamaye.Pending.MonomialInvariantFunctorialAlong.gammaAlong_pullback_diffeomorph,TranscendentalNakamaye.Pending.MonomialInvariantFunctorialAlong.gammaAlong_geoTransform_empty}
\leanok
Let \(n\ge1\) and \(\underline{\mathcal I}=(M,N,E,\mathfrak a,d)\) be well formed and monomial (Definition~\ref{hir:monomial-part}), with \(\mathfrak a_y\) containing, for every
\(y\in N\), a germ which does not vanish identically on \(N\) near \(y\), and \(d\ge1\). For \(x\in\operatorname{cosupp}\underline{\mathcal I}\) let \(\gamma(x)=(-\gamma_1(x),\gamma_2(x),\gamma_3(x))\)
be the invariant of the exponents of the members through \(x\) (with the order of Lemma~\ref{hir:monomial-invariant}: \(\gamma_1\) the least number of members whose exponents
add up to at least \(d\), \(\gamma_2\) the largest such sum, \(\gamma_3\) the largest such set of members). Then:
\begin{enumerate}
\item \(\gamma\) is upper semicontinuous on \(\operatorname{cosupp}\underline{\mathcal I}\). If \(W\subset M\) is open and \(\gamma\le\lambda\) at the points of
\(\operatorname{cosupp}\underline{\mathcal I}\cap W\), then \(\{x\in\operatorname{cosupp}\underline{\mathcal I}\cap W:\gamma(x)=\lambda\}\) is an admissible centre of
\(\underline{\mathcal I}|_W\). Near each of its points it is \(N\cap\bigcap_{H\in\lambda_3}H\).
\item After blowing up the maximum locus \(C\) of \(\gamma\) on the cosupport, \(\gamma(z)<\gamma(\sigma(z))\) at every point \(z\) of the new \(N\) over \(C\), and
\(\gamma(z)=\gamma(\sigma(z))\) at the other points of the new \(N\).
\item If \(\operatorname{cosupp}\underline{\mathcal I}\) is compact (\(M\) need not be), the iterated blowup of the maximum locus of \(\gamma\) reaches an empty cosupport after
finitely many steps.
\item \(\gamma\) is unchanged by restriction to open sets, by biholomorphisms and by steps with empty centre.
\end{enumerate}
\end{lemma}
\begin{proof}
\leanok
(1) Upper semicontinuity: Lemma~\ref{hir:monomial-invariant-along}. The restriction \(\underline{\mathcal I}|_W\) is again well formed, nonzero along \(N\cap W\) and monomial (the
monomial part commutes with restriction, Proposition~\ref{hir:monomial-part-along}; Lemma~\ref{hir:open-restriction-marked}), and \(\gamma\) of the restriction is the
restriction of \(\gamma\) (Lemma~\ref{hir:monomial-invariance-along}). Off the cosupport \(\gamma=\bot\), so \(\gamma\le\lambda\) holds on all of \(N\cap W\), and
Lemma~\ref{hir:monomial-invariant-along}(3) applied on \(W\) gives an admissible centre. Near a point \(x\) of value \(\lambda\), the points of \(N\) of value \(\lambda\) are those
on the members of \(\lambda_3\) (Lemma~\ref{hir:monomial-invariant-along}; Theorem~\ref{bcalc:member-intersection-along}).
(2) Lemma~\ref{hir:monomial-step-along}.
(3) The cosupports along a multiple test blow-up stay closed (Lemma~\ref{hir:giraud-along}) and lie over the original cosupport (Lemma~\ref{bcalc:transform-off-centre}). The
projection of a tower is proper (Theorem~\ref{hir:principalization-local}), so they stay compact, and the upper semicontinuous \(\gamma\) attains its maximum on them
(Lemma~\ref{hir:invariant-order}). The argument of Proposition~\ref{hir:resolution-driver} then runs with these compact cosupports in place of a compact base: the maximum locus
is an admissible centre with empty interior by (1), by (2) every value of the new cosupport is below the old maximum, and the values are well-ordered
(Lemma~\ref{hir:monomial-invariant}(4)). This is Theorem~\ref{hir:monomial-resolution-along} with compact cosupports; for \(N=M\), Theorem~\ref{hir:monomial-resolution}.
(4) Lemma~\ref{hir:monomial-invariance-along}.
\end{proof}

\begin{lemma}[One step of the monomial case, \(N=M\)]
\label{hir:monomial-step}
\uses{hir:monomial-invariant,hir:monomial-invariant-geometry,hir:monomial-part-ambient,hir:monomial-factorization,hir:exponent-locally-constant,hir:geo-transform,hir:geo-transform-ambient,hir:geo-transform-coherent,bcalc:monomial-transform-coords,bcalc:boundary-count-drop,bcalc:transform-off-centre,bcalc:isomorphism-off-centre,bcalc:transform-wellformed,bcalc:snc-transform,bcalc:hyperplane-ideal,bcalc:strict-geo}
\source{wlodarczyk05-hironaka,bierstone-milman08-functoriality}
\lean{TranscendentalNakamaye.Pending.MonomialResolution.IsPrincipalAt,TranscendentalNakamaye.Pending.MonomialResolution.extExp,TranscendentalNakamaye.Pending.MonomialResolution.IsPrincipalAt.eventually,TranscendentalNakamaye.Pending.MonomialResolution.IsPrincipalAt.isNonzeroAt,TranscendentalNakamaye.Pending.MonomialResolution.IsPrincipalAt.monomialPart_le,TranscendentalNakamaye.Pending.MonomialResolution.IsPrincipalAt.gamma_eq,TranscendentalNakamaye.Pending.MonomialResolution.isPrincipalAt_of_monomial,TranscendentalNakamaye.Pending.MonomialResolution.exists_chart_comp,TranscendentalNakamaye.Pending.MonomialResolution.geoTransform_mem_iff_pullback,TranscendentalNakamaye.Pending.MonomialResolution.isInSpanAt_eq_congr,TranscendentalNakamaye.Pending.MonomialResolution.mem_get_geoTransform_iff_of_notMem,TranscendentalNakamaye.Pending.MonomialResolution.castE,TranscendentalNakamaye.Pending.MonomialResolution.isPrincipalAt_geoTransform_of_notMem,TranscendentalNakamaye.Pending.MonomialResolution.isPrincipalAt_geoTransform_of_notMem_of_monomial,TranscendentalNakamaye.Pending.MonomialResolution.gamma_geoTransform_eq_of_notMem,TranscendentalNakamaye.Pending.MonomialResolution.maxLocusγ,TranscendentalNakamaye.Pending.MonomialResolution.mem_maxLocusγ_iff,TranscendentalNakamaye.Pending.MonomialResolution.eventually_mem_maxLocusγ_iff,TranscendentalNakamaye.Pending.MonomialResolution.filter_eventually_subset_eq_topIdx,TranscendentalNakamaye.Pending.MonomialResolution.sum_topIdx_eq,TranscendentalNakamaye.Pending.MonomialResolution.isPrincipalAt_geoTransform_of_mem,TranscendentalNakamaye.Pending.MonomialResolution.gamma_geoTransform_lt_of_mem,TranscendentalNakamaye.Pending.MonomialResolution.monomial_step}
\leanok
Let \(n\ge1\) and \(\underline{\mathcal I}=(M,M,E,\mathfrak a,d)\) be well formed, with \(\mathfrak a\) monomial (\(\mathcal M(\mathfrak a)\subset\mathfrak a\),
Proposition~\ref{hir:monomial-part-ambient}) and \(\mathfrak a_y\) containing a germ not vanishing near \(y\) for every \(y\). Let \(\gamma\) be the invariant of
Lemma~\ref{hir:monomial-invariant-geometry}, \(C\) its maximum locus on \(\operatorname{cosupp}\underline{\mathcal I}\), \(\sigma\colon M'\to M\) the blowup along \(C\),
\(\underline{\mathcal I}'=(M',N',E',\mathfrak a',d)\) the transform (Definition~\ref{hir:geo-transform}) and \(\gamma'\) its invariant. Say that an ideal \(\mathfrak b\) is
\emph{principal monomial at \(z\) with exponents \(b\)} if \(\mathfrak b_z=(\prod_{z\in H_i}u_{\tau(i)}^{b_i})\) in coordinates \(u\) centred at \(z\) in which every member
\(H_i\) through \(z\) is the hyperplane \(\{u_{\tau(i)}=0\}\) (\(\tau\) injective on these members).
\begin{enumerate}
\item If \(\mathfrak b\) is coherent and principal monomial at \(z\) with exponents \(b\), then \(\mathfrak b_z\) contains a germ not vanishing near \(z\),
\(\mathcal M(\mathfrak b)_z\subset\mathfrak b_z\), and \(\gamma(z)=\gamma(S(z),b)\). The ideal \(\mathfrak a\) is principal monomial at every point \(x\), with the exponents
\(e_{H_i}(x)\) of Lemma~\ref{hir:exponent-locally-constant}.
\item (off the centre) If \(\sigma(z)\notin C\), then \(\mathfrak a'_z=(\sigma^*\mathfrak a)_z\), \(\mathfrak a'\) is principal monomial at \(z\) with the exponents of \(\sigma(z)\) on
the strict transforms of the members through \(\sigma(z)\), and \(\gamma'(z)=\gamma(\sigma(z))\).
\item (the maximum locus) Let \(x\in C\), \(\gamma(x)=v\) with top set \(I\). Then \(y\in C\) if and only if \(\gamma(y)=v\); near \(x\), \(C=\bigcap_{i\in I}H_i\); and the
members \(H_i\) through \(x\) with \(C\subset H_i\) near \(x\) are exactly those with \(i\in I\).
\item (over the centre) If \(x=\sigma(z)\in C\), then \(\mathfrak a'\) is principal monomial at \(z\), with the exponents \(e_{H_i}(x)\) on the strict transforms \(H_i'\)
through \(z\) and \(a_D=\sum_{i\in I}e_{H_i}(x)-d\) on \(D=\sigma^{-1}(C)\); and \(\gamma'(z)<\gamma(x)\).
\item Hence \(\underline{\mathcal I}'\) is well formed with \(N'=M'\), \(\mathfrak a'\) is monomial and \(\mathfrak a'_z\) contains a germ not vanishing near \(z\) for every
\(z\); \(\gamma'(z)<\gamma(\sigma(z))\) over \(C\) and \(\gamma'(z)=\gamma(\sigma(z))\) elsewhere; and \(\gamma'(z)<\gamma(x_0)\) for every \(z\in\operatorname{cosupp}\underline{\mathcal I}'\)
and \(x_0\in C\).
\end{enumerate}
\end{lemma}
\begin{proof}
\leanok
(1) The first sentence is Lemma~\ref{hir:monomial-invariant-geometry}(6) (a stalk generator of a coherent ideal generates the nearby stalks). Take simple
normal crossings coordinates \(u\) at \(x\), the members through \(x\) being \(\{u_{\tau(i)}=0\}\), and \(m=\prod_{x\in H_i}u_{\tau(i)}^{e_{H_i}(x)}\). A germ of
\(\mathfrak a_x\subset\mathcal M(\mathfrak a)_x\) is divisible by every \(u_{\tau(i)}^{e_{H_i}(x)}\), hence by \(m\) (Lemma~\ref{hir:monomial-factorization}(1)); conversely
\(m\in\mathcal M(\mathfrak a)_x\subset\mathfrak a_x\). So \(\mathfrak a_x=(m)\).
(2) Off \(D\) the blowup is a local biholomorphism (Lemma~\ref{bcalc:isomorphism-off-centre}), so for a flattening \(\Phi\) at \(\sigma(z)\) centred there,
\(\Phi\circ\sigma\) (read in charts) is a coordinate system centred at \(z\). The ideal of \(C\) pulls back to the unit ideal at \(z\), so the colon defining
\(\mathfrak a'\) is \((\sigma^*\mathfrak a)_z\) (\(N'=M'\), Lemma~\ref{hir:geo-transform-ambient}), which is generated by \(m\circ\sigma\) for a generator \(m\) of
\(\mathfrak a_{\sigma(z)}\) (Lemma~\ref{bcalc:hyperplane-ideal}). Near \(z\) the strict transform \(H_i'\) is \(\sigma^{-1}(H_i)\) (Lemma~\ref{bcalc:transform-off-centre}) and
\(D\) does not pass through \(z\); so with the generator \(m\) of (1) at \(\sigma(z)\), \(\mathfrak a'\) is principal monomial at \(z\) in the coordinates
\(\Phi\circ\sigma\), with the same members and exponents, and \(\gamma'(z)=\gamma(\sigma(z))\) by (1).
(3) Off the cosupport \(\gamma=\bot\) (Lemma~\ref{hir:monomial-invariant-geometry}(4)), so \(\gamma\le v\) everywhere and \(C=\{\gamma=v\}\); near \(x\) this is
\(\bigcap_{i\in I}H_i\) (Lemma~\ref{hir:monomial-invariant-geometry}(3)). If \(x\in H_i\), \(i\notin I\), take coordinates \(u\) as in (1): the points with
\(u_{\tau(i)}\ne0\) and all other coordinates \(0\) accumulate at \(x\); they lie on every \(H_j\), \(j\in I\) (\(\tau\) is injective), hence in \(C\), but not on \(H_i\).
(4) By (1) at \(x\) and (3), \(\mathfrak a_x=(\prod u_{\tau(i)}^{e_{H_i}(x)})\) and \(e=\sum_{x\in H_i,\ C\subset H_i\text{ near }x}e_{H_i}(x)=\sum_{i\in I}e_{H_i}(x)\ge d\),
\(I\) being feasible. Corollary~\ref{bcalc:monomial-transform-coords} gives the principal monomial form at \(z\), with \(a_D=e-d\). A strict transform through
\(z\) comes from a member through \(x\) (strict transforms lie in preimages, Lemma~\ref{bcalc:strict-geo}), and the last member of \(E'\) is \(D\)
(Theorem~\ref{bcalc:snc-transform}); by Lemma~\ref{bcalc:boundary-count-drop} (near \(x\), \(\bigcap_{i\in I}H_i\subset C\) by (3)) some \(j_0\in I\) has
\(z\notin H_{j_0}'\). So the members through \(z\) are among \((S(x)\setminus\{j_0\})\cup\{D\}\), the exponents of the old members are unchanged and
\(a_D+d=\sum_Ia\); Lemma~\ref{hir:monomial-invariant}(3) gives \(\gamma'(z)<\gamma(x)\).
(5) \(C\) is an admissible centre (Lemma~\ref{hir:monomial-invariant-geometry}(5)), so \(\underline{\mathcal I}'\) is well formed (Lemma~\ref{bcalc:transform-wellformed}),
\(N'=M'\) (Lemma~\ref{hir:geo-transform-ambient}) and \(\mathfrak a'\) is coherent (Lemma~\ref{hir:geo-transform-coherent}). By (2) and (4) \(\mathfrak a'\) is principal monomial at
every point, so (1) gives the non-vanishing and monomiality. Let \(z\in\operatorname{cosupp}\underline{\mathcal I}'\) and \(x_0\in C\). If \(\sigma(z)\in C\), then
\(\gamma'(z)<\gamma(\sigma(z))=\gamma(x_0)\) by (4) and (3). Otherwise \(\sigma(z)\in\operatorname{cosupp}\underline{\mathcal I}\setminus C\)
(Lemma~\ref{bcalc:transform-off-centre}), so \(\gamma(\sigma(z))<\gamma(x_0)\), and \(\gamma'(z)=\gamma(\sigma(z))\) by (2).
\end{proof}

\begin{theorem}[Resolution of monomial marked ideals, \(N=M\)]
\label{hir:monomial-resolution}
\uses{hir:monomial-step,hir:monomial-invariant,hir:monomial-invariant-geometry,hir:resolution-driver,hir:test-tower,hir:geo-transform-ambient,hir:giraud-support}
\source{wlodarczyk05-hironaka,bierstone-milman08-functoriality}
\lean{TranscendentalNakamaye.Pending.MonomialResolution.towerInv,TranscendentalNakamaye.Pending.MonomialResolution.invariants_canonicalSegment,TranscendentalNakamaye.Pending.MonomialResolution.exists_isGeoResolution_of_monomial}
\leanok
Let \(M\) be compact, \(n\ge1\), and \(\underline{\mathcal I}=(M,M,E,\mathfrak a,d)\) well formed with \(d\ge1\), \(\mathfrak a\) monomial and \(\mathfrak a_y\) containing a germ not
vanishing near \(y\) for every \(y\). For a tower \(T\) let \(\gamma_T\) be the invariant of Lemma~\ref{hir:monomial-invariant-geometry} of the transform
\(\underline{\mathcal I}_T\). Along every canonical segment for \(\gamma\) (Theorem~\ref{hir:resolution-driver}) that is a multiple test blow-up, \(\underline{\mathcal I}_T\) is
well formed with \(N_T=M_T\), \(\mathfrak a_T\) is monomial and nowhere locally zero; and some canonical segment for \(\gamma\) is a resolution of
\(\underline{\mathcal I}\).
\end{theorem}
\begin{proof}
\leanok
The first sentence is induction on the tower with Lemma~\ref{hir:monomial-step}(5) (each centre of a canonical segment is the maximum locus). Then
Theorem~\ref{hir:resolution-driver}(2) applies: the cosupports are closed (coherent ideals, Lemma~\ref{hir:giraud-support}); \(\gamma_T\) is upper semicontinuous and
its maximum locus is an admissible centre with empty interior (Lemma~\ref{hir:monomial-invariant-geometry}(2), (5), \(d\ge1\)); after the blowup of the
maximum locus \(\gamma\) is below the old maximum on the new cosupport (Lemma~\ref{hir:monomial-step}(5)); the values are well-founded
(Lemma~\ref{hir:monomial-invariant}(4)).
\end{proof}

\begin{lemma}[The monomial invariant under restriction, biholomorphisms and empty steps, \(N=M\)]
\label{hir:monomial-invariance}
\uses{hir:monomial-invariant-geometry,hir:exponent-locally-constant,hir:monomial-step,hir:monomial-resolution,hir:open-restriction-marked,hir:open-restriction-tower,hir:biholomorphism-transform,hir:tower-transport,hir:biholomorphism-blowup}
\source{wlodarczyk05-hironaka,bierstone-milman08-functoriality}
\lean{TranscendentalNakamaye.Pending.MonomialInvariantFunctorial.isDivNear_pullback_iff,TranscendentalNakamaye.Pending.MonomialInvariantFunctorial.divExp_pullback,TranscendentalNakamaye.Pending.MonomialInvariantFunctorial.memberIdx_map_preimage,TranscendentalNakamaye.Pending.MonomialInvariantFunctorial.gamma_pullback,TranscendentalNakamaye.Pending.MonomialInvariantFunctorial.gamma_pullback_val,TranscendentalNakamaye.Pending.MonomialInvariantFunctorial.gamma_pullback_diffeomorph,TranscendentalNakamaye.Pending.MonomialInvariantFunctorial.gamma_geoTransform_empty,TranscendentalNakamaye.Pending.MonomialInvariantFunctorial.towerInv_restrict,TranscendentalNakamaye.Pending.MonomialInvariantFunctorial.towerInv_transport}
\leanok
Let \(\gamma=\gamma_{\mathfrak a,E,d}\) be the invariant of Lemma~\ref{hir:monomial-invariant-geometry}, read off the members of \(E\) through a point and the
exponents \(e_{H}\) of Lemma~\ref{hir:exponent-locally-constant}.
\begin{enumerate}
\item If \(\varphi\colon M'\to M\) satisfies \(\operatorname{cosupp}(\varphi^*\mathfrak a,e)=\varphi^{-1}\operatorname{cosupp}(\mathfrak a,e)\) for every \(e\) and maps the
neighbourhoods of \(x'\) onto neighbourhoods of \(\varphi(x')\), then \(e_{\varphi^{-1}H}(\varphi^*\mathfrak a)(x')=e_H(\mathfrak a)(\varphi(x'))\) for every \(H\), and
\(\gamma_{\varphi^*\mathfrak a,\varphi^{-1}E,d}(x')=\gamma_{\mathfrak a,E,d}(\varphi(x'))\).
\item In particular \(\gamma\) of the restriction \(\underline{\mathcal I}|_U\) to an open \(U\) at \(p\) is \(\gamma\) of \(\underline{\mathcal I}\) at \(p\), and \(\gamma\) of
\(\Phi^*\underline{\mathcal I}\) at \(p\) is \(\gamma\) of \(\underline{\mathcal I}\) at \(\Phi(p)\) for a biholomorphism \(\Phi\).
\item If \(\underline{\mathcal I}=(M,M,E,\mathfrak a,d)\) is monomial and \(\sigma\) is the blowup along \(\emptyset\), then \(\gamma'(z)=\gamma(\sigma(z))\).
\item Along the restriction to \(U\) of a multiple test blow-up \(T\) of \(\underline{\mathcal I}\) (Lemma~\ref{hir:open-restriction-tower}), the invariants of the
transforms of \(\underline{\mathcal I}|_U\) and of \(\underline{\mathcal I}\) agree at the points over \(U\); along the transport of \(T\) by a biholomorphism of the
base (Lemma~\ref{hir:tower-transport}), they correspond through the last-stage biholomorphism.
\end{enumerate}
\end{lemma}
\begin{proof}
\leanok
(1) The exponent \(e_H(\mathfrak a)(x)\) is the largest \(e\) such that \(\operatorname{ord}\mathfrak a\ge e\) at the points of \(H\) near \(x\), i.e.\ such that near \(x\),
\(H\subset\operatorname{cosupp}(\mathfrak a,e)\); under the hypotheses this condition at \(x'\) for \((\varphi^*\mathfrak a,\varphi^{-1}H)\) is the condition at
\(\varphi(x')\) for \((\mathfrak a,H)\). The members of \(\varphi^{-1}E\) through \(x'\) are the preimages of the members of \(E\) through \(\varphi(x')\), with the same
indices; so the data defining \(\gamma\) agree. (2) The hypotheses of (1) hold for the inclusion of an open set (Lemma~\ref{hir:open-restriction-marked}) and for a
biholomorphism (Lemma~\ref{hir:biholomorphism-transform}). (3) \(\sigma\) is a biholomorphism and \(\sigma(z)\notin\emptyset\), so this is the off-centre case of
Lemma~\ref{hir:monomial-step}(2); the appended member \(\sigma^{-1}(\emptyset)\) passes through no point. (4) The transforms along the restricted tower are the
restrictions of the transforms (Lemma~\ref{hir:open-restriction-tower}) and, from the first blowup on, the transforms along the transported tower are equal to
those along \(T\) (Lemma~\ref{hir:tower-transport}); apply (2).
\end{proof}

\begin{lemma}[The invariant of a monomial marked ideal along a submanifold at one stage]
\label{hir:monomial-invariant-along}
\uses{hir:monomial-invariant,hir:monomial-invariant-geometry,hir:exponent-along,hir:monomial-part-along,hir:monomial-factorization-along,hir:local-retraction,hir:point-power-pullback,bcalc:member-intersection-along,bcalc:snc-simultaneous,hir:snc,hir:cosupport,hir:marked-ideal,hir:order-product,hir:resolution-driver,hir:invariant-order,lag:regular-point-basics}
\source{wlodarczyk05-hironaka,bierstone-milman08-functoriality}
\lean{TranscendentalNakamaye.Pending.MonomialCaseAlong.expDataAlong,TranscendentalNakamaye.Pending.MonomialCaseAlong.gammaAlong,TranscendentalNakamaye.Pending.MonomialCaseAlong.gammaAlong_univ,TranscendentalNakamaye.Pending.MonomialCaseAlong.eventually_expDataAlong_eq,TranscendentalNakamaye.Pending.MonomialCaseAlong.eventually_gammaAlong_le,TranscendentalNakamaye.Pending.MonomialCaseAlong.eventually_gammaAlong_eq_iff,TranscendentalNakamaye.Pending.MonomialCaseAlong.mem_cosupport_iff_of_isMonomialAlong,TranscendentalNakamaye.Pending.MonomialCaseAlong.gammaAlong_eq_bot_iff_of_isMonomialAlong,TranscendentalNakamaye.Pending.MonomialCaseAlong.top_of_gammaAlong_eq,TranscendentalNakamaye.Pending.MonomialCaseAlong.eventually_gammaAlong_eq_iff_mem_memberInter,TranscendentalNakamaye.Pending.MonomialCaseAlong.mem_memberInter_of_gammaAlong_eq,TranscendentalNakamaye.Pending.MonomialCaseAlong.isInvUSCOn_gammaAlong,TranscendentalNakamaye.Pending.MonomialCaseAlong.isClosed_gammaAlong_eq,TranscendentalNakamaye.Pending.MonomialCaseAlong.isAdmissibleCentre_gammaAlong_eq,TranscendentalNakamaye.Pending.MonomialCaseAlong.isAdmissibleCentre_maxLocus_gammaAlong,TranscendentalNakamaye.Pending.MonomialCaseAlong.eventually_isInSpanOnAt_iff_of_isCoherent,TranscendentalNakamaye.Pending.MonomialCaseAlong.IsPrincipalMonomialAlongAt,TranscendentalNakamaye.Pending.MonomialCaseAlong.eventually_isNonzeroAlongAt_of_principal,TranscendentalNakamaye.Pending.MonomialCaseAlong.divExpAlong_eq_of_principal,TranscendentalNakamaye.Pending.MonomialCaseAlong.gammaAlong_eq_of_principal,TranscendentalNakamaye.Pending.MonomialCaseAlong.isNonzeroAlongAt_and_isMonomialAlong_of_principal}
\leanok
Let \(n\ge1\), \(N\subset M\) a closed submanifold, \(E=(E_0,\dots,E_{r-1})\) a simple normal crossings divisor transverse to \(N\), \(\mathfrak a\) a coherent
ideal with \(\mathfrak a_y\) containing a germ not vanishing identically on \(N\) near \(y\) for every \(y\in N\), and \(d\in\mathbb N\). For \(x\in N\) let
\(S(x)=\{i:x\in E_i\}\), \(a^N_x(i)=e^N_{E_i}(x)\) the exponents along \(N\) of Lemma~\ref{hir:exponent-along}, and \(\gamma_N(x)=\gamma(S(x),a^N_x)\) the invariant of
Lemma~\ref{hir:monomial-invariant}. For \(N=M\) this is the invariant of Lemma~\ref{hir:monomial-invariant-geometry}.
\begin{enumerate}
\item (upper semicontinuity on \(N\)) Near \(x\in N\), for \(y\in N\): \(S(y)\subset S(x)\) and \(a^N_y=a^N_x\) on \(S(y)\); hence \(\gamma_N(y)\le\gamma_N(x)\), and if
\(\gamma_N(x)=v\ne\bot\) with top set \(I\), then \(\gamma_N(y)=v\) if and only if \(y\in\bigcap_{i\in I}E_i\).
\item (the cosupport) If \(\mathfrak a\) is monomial along \(N\) (\(\operatorname{Res}_N\mathcal M_N(\mathfrak a)\subset\operatorname{Res}_N\mathfrak a\),
Proposition~\ref{hir:monomial-part-along}), then for \(x\in N\): \(x\in\operatorname{cosupp}(\mathfrak a|_N,d)\) if and only if \(\sum_{i\in S(x)}a^N_x(i)\ge d\), i.e.\
\(\gamma_N(x)\ne\bot\).
\item (the maximum locus) If moreover \(\underline{\mathcal I}=(M,N,E,\mathfrak a,d)\), and \(\gamma_N\le v\) on \(N\), the set \(\{y\in N:\gamma_N(y)=v\}\) is an admissible
centre; the points of \(\operatorname{cosupp}\underline{\mathcal I}\) at which \(\gamma_N\) attains its maximum on the cosupport form an admissible centre, with empty
interior if \(d\ge1\).
\item (principal monomial ideals along \(N\)) Suppose that in coordinates \(u\) at \(x\in N\) in which \(N=\mathbf C^k\times0\) and the members \(E_i\), \(i\in T\), are the
hyperplanes \(\{u_{\sigma(i)}=0\}\) (\(\sigma(i)<k\), \(\sigma\) injective on \(T\)), the restriction of \(\mathfrak a_y\) to \(N\) is generated by \(m=\prod_{i\in T}u_{\sigma(i)}^{b_i}\) for
\(y\) near \(x\) (for \(\mathfrak a\) coherent it suffices at \(y=x\)). Then \(\mathfrak a_y\) contains a germ not vanishing identically on \(N\) near \(y\) for \(y\in N\) near \(x\),
\(a^N_x(i)=b_i\) for \(i\in T\) with \(x\in E_i\), and \(\gamma_N(x)=\gamma(S(x),b)\) when \(T=S(x)\). If this holds at every point of \(N\) (with \(T=S(x)\)), \(\mathfrak a\) is
monomial along \(N\) and nonzero along \(N\) near every point. (Item (4) does not use the nonvanishing hypothesis above.)
\end{enumerate}
\end{lemma}
\begin{proof}
\leanok
(1) The members are closed and \(e^N_{E_i}\) is locally constant on \(N\cap E_i\) (Lemma~\ref{hir:exponent-along}(4), (5)); then
Lemma~\ref{hir:monomial-invariant}(2), and the third entry of \(v\) is \(I\).
(2) Take coordinates at \(x\) adapted to \(N\) and \(E\) (Lemma~\ref{bcalc:snc-simultaneous}) and the monomial \(m=\prod_{i\in S(x)}u_{\sigma(i)}^{a^N_x(i)}\) of
Proposition~\ref{hir:monomial-part-along}(2). Its order along \(N\) at \(x\) is its ambient order \(\sum a^N_x(i)\): the flat retraction \(r\) onto \(N\) keeps the coordinates
\(u_{\sigma(i)}\) (Lemma~\ref{hir:monomial-factorization-along}(1)), so \(m\circ r=m\) near \(x\), and orders along \(N\) are ambient orders of retracted germs
(Lemma~\ref{hir:exponent-along}(1)); a coordinate has order one and orders add (Lemma~\ref{hir:order-product}). If \(x\in\operatorname{cosupp}\), then \(m\in\mathcal M_N(\mathfrak a)_x\),
whose restriction lies in \(\operatorname{Res}_N\mathfrak a_x\) by monomiality, has order \(\ge d\) along \(N\), so \(\sum a^N_x(i)\ge d\). Conversely every germ of \(\mathfrak a_x\) is a
multiple of \(m\) on \(N\) (Proposition~\ref{hir:monomial-part-along}(1), (2)), hence of order \(\ge\sum a^N_x(i)\) along \(N\).
(3) As in Lemma~\ref{hir:monomial-invariant-geometry}(5): the set is closed by (1), lies in the cosupport by (2), and agrees near each of its points with
\(N\cap\bigcap_{i\in I}E_i\), a closed submanifold with only normal crossings with \(E\) (Lemma~\ref{bcalc:member-intersection-along}); both properties are local
(Lemma~\ref{lag:regular-point-basics}). Off the cosupport \(\gamma_N=\bot\) on \(N\) by (2), so the maximum locus is such a level set; if \(d\ge1\), \(I\ne\emptyset\) and the locus
lies in a member, which has empty interior.
(4) Divisibility: near \(x\), the restriction of every germ of \(\mathfrak a_z\), \(z\in N\cap E_i\), is a multiple of \(m\), and \(m\in\mathcal I_{E_i}^{b_i}\), so the points of
\(N\cap E_i\) near \(x\) lie in \(\operatorname{cosupp}(\mathfrak a|_N,b_i)\). Non-divisibility by \(b_i+1\): the points \(z\in N\cap E_i\) near \(x\) off the other hyperplanes
(moving the other coordinates of index \(<k\)) accumulate at \(x\); \(m|_N\) is a combination of restrictions of germs of \(\mathfrak a_z\), so if they all had order \(\ge b_i+1\)
along \(N\) at \(z\), so would \(m\), hence (retraction, as in (2)) ambiently, while \(m\) is a unit times \(u_{\sigma(i)}^{b_i}\) at \(z\). A germ of \(\mathfrak a_x\) not vanishing on
\(N\) exists because \(m|_N\) is a combination of such germs and \(m\) does not vanish identically on \(N\) near \(x\) (Proposition~\ref{hir:monomial-part-along}(2)). Restricted
generators propagate to nearby points by coherence. The last sentence: the exponents along \(N\) are the \(b_i\), so the generator of \(\mathcal M_N(\mathfrak a)\) of
Proposition~\ref{hir:monomial-part-along}(2) in the same coordinates is \(m\), whose restriction lies in \(\operatorname{Res}_N\mathfrak a\).
\end{proof}

\begin{lemma}[One step of the monomial case along a submanifold]
\label{hir:monomial-step-along}
\uses{hir:monomial-invariant-along,hir:monomial-invariant,hir:monomial-step,hir:monomial-part-along,hir:exponent-along,bcalc:snc-simultaneous,bcalc:monomial-transform-along-coords,bcalc:monomial-transform-along-off-centre,bcalc:boundary-count-drop-along,bcalc:member-intersection-along,bcalc:transform-wellformed,bcalc:transform-off-centre,bcalc:strict-geo,bcalc:snc-transform,hir:geo-transform}
\source{wlodarczyk05-hironaka,bierstone-milman08-functoriality}
\lean{TranscendentalNakamaye.Pending.MonomialStepAlong.IsPrincipalAlongAt,TranscendentalNakamaye.Pending.MonomialStepAlong.IsPrincipalAlongAt.isPrincipalMonomialAlongAt,TranscendentalNakamaye.Pending.MonomialStepAlong.IsPrincipalAlongAt.gammaAlong_eq,TranscendentalNakamaye.Pending.MonomialStepAlong.isPrincipalAlongAt_of_data,TranscendentalNakamaye.Pending.MonomialStepAlong.isPrincipalAlongAt_of_isMonomialAlong,TranscendentalNakamaye.Pending.MonomialStepAlong.maxLocusA,TranscendentalNakamaye.Pending.MonomialStepAlong.mem_maxLocusA_iff,TranscendentalNakamaye.Pending.MonomialStepAlong.eventually_mem_maxLocusA_iff,TranscendentalNakamaye.Pending.MonomialStepAlong.filter_eventually_subset_eq_topIdx_along,TranscendentalNakamaye.Pending.MonomialStepAlong.sum_topIdx_eq_along,TranscendentalNakamaye.Pending.MonomialStepAlong.isAdmissibleCentre_maxLocusA,TranscendentalNakamaye.Pending.MonomialStepAlong.isPrincipalAlongAt_geoTransform_of_mem,TranscendentalNakamaye.Pending.MonomialStepAlong.gammaAlong_geoTransform_lt_of_mem,TranscendentalNakamaye.Pending.MonomialStepAlong.isPrincipalAlongAt_geoTransform_of_notMem,TranscendentalNakamaye.Pending.MonomialStepAlong.gammaAlong_geoTransform_eq_of_notMem,TranscendentalNakamaye.Pending.MonomialStepAlong.monomial_step_along}
\leanok
Let \(n\ge1\) and \(\underline{\mathcal I}=(M,N,E,\mathfrak a,d)\) be well formed, with \(\mathfrak a\) nonzero and monomial along \(N\) (Lemma~\ref{hir:monomial-invariant-along}), \(\gamma_N\)
its invariant, \(C\) the maximum locus of \(\gamma_N\) on \(\operatorname{cosupp}\underline{\mathcal I}\), \(\sigma\colon M'\to M\) the blowup along \(C\), \(\underline{\mathcal I}'=(M',N',E',\mathfrak a',d)\)
the transform and \(\gamma'_N\) its invariant along \(N'\).
\begin{enumerate}
\item \(\mathfrak a\) is principal monomial along \(N\) at every \(x\in N\), in centred coordinates adapted to \(N\) and \(E\), with the exponents \(a^N_x\).
\item Near a point \(x\in C\) with top set \(I\), \(C=N\cap\bigcap_{i\in I}E_i\), and the members through \(x\) containing \(C\) near \(x\) are exactly those with \(i\in I\).
\item (over the centre) If \(z\in N'\) and \(\sigma(z)\in C\), then \(\mathfrak a'\) is principal monomial along \(N'\) at \(z\), with the exponents of \(\sigma(z)\) on the strict transforms and
\(\sum_{i\in I}a^N_{\sigma(z)}(i)-d\) on \(D=\sigma^{-1}(C)\); and \(\gamma'_N(z)<\gamma_N(\sigma(z))\).
\item (off the centre) If \(z\in N'\) and \(\sigma(z)\notin C\), then \(\mathfrak a'\) is principal monomial along \(N'\) at \(z\), with the exponents of \(\sigma(z)\), and
\(\gamma'_N(z)=\gamma_N(\sigma(z))\).
\item Hence \(\underline{\mathcal I}'\) is well formed, \(\mathfrak a'\) is nonzero and monomial along \(N'\), and \(\gamma'_N(z)<\gamma_N(x_0)\) for every
\(z\in\operatorname{cosupp}\underline{\mathcal I}'\) and \(x_0\in C\).
\end{enumerate}
\end{lemma}
\begin{proof}
\leanok
(1) Proposition~\ref{hir:monomial-part-along}(6) in the coordinates of Lemma~\ref{bcalc:snc-simultaneous}; they are centred by subtracting the image of \(x\), which lies in
\(\mathbf C^k\times0\) and has vanishing \(\sigma(i)\)-th coordinates for the members through \(x\) (this changes neither \(N\), nor the hyperplanes, nor the monomial).
(2) As in Lemma~\ref{hir:monomial-step}(3), from Lemma~\ref{hir:monomial-invariant-along}(1)--(3); the curve leaving the hyperplane of a member outside \(I\) moves a
coordinate of index \(<k\), so it stays in \(N\).
(3) By (1) at \(x=\sigma(z)\) and (2), \(\sum_{C\subset E_i\text{ near }x}a^N_x(i)=\sum_{i\in I}a^N_x(i)\ge d\); Corollary~\ref{bcalc:monomial-transform-along-coords} gives the principal
monomial form at \(z\) in coordinates adapted to \(N'\) and \(E'\). By Theorem~\ref{bcalc:boundary-count-drop-along} some \(j_0\in I\) has \(z\notin E'_{j_0}\); a strict transform
through \(z\) comes from a member through \(x\) (Lemma~\ref{bcalc:strict-geo}) and the last member of \(E'\) is \(D\) (Theorem~\ref{bcalc:snc-transform}), so
Lemma~\ref{hir:monomial-invariant-along}(4) and Lemma~\ref{hir:monomial-invariant}(3) give \(\gamma'_N(z)<\gamma_N(x)\).
(4) Lemma~\ref{bcalc:monomial-transform-along-off-centre} with the centred coordinates of (1); the members through \(z\) are the strict transforms of those through
\(\sigma(z)\) (Lemma~\ref{bcalc:transform-off-centre}), with the same exponents (Lemma~\ref{hir:monomial-invariant-along}(4)).
(5) \(C\) is admissible (Lemma~\ref{hir:monomial-invariant-along}(3)), so \(\underline{\mathcal I}'\) is well formed (Lemma~\ref{bcalc:transform-wellformed}); by (3), (4) \(\mathfrak a'\) is
principal monomial along \(N'\) at every point of \(N'\), so Lemma~\ref{hir:monomial-invariant-along}(4) gives nonvanishing and monomiality. For \(z\) in the new
cosupport: over \(C\) use (3) and \(\gamma_N(\sigma(z))=\gamma_N(x_0)\) (the maximum locus is a level set); off \(C\), \(\sigma(z)\) lies in the old cosupport outside \(C\)
(Lemma~\ref{bcalc:transform-off-centre}), so \(\gamma'_N(z)=\gamma_N(\sigma(z))<\gamma_N(x_0)\).
\end{proof}

\begin{theorem}[Resolution of monomial marked ideals along a submanifold]
\label{hir:monomial-resolution-along}
\uses{hir:monomial-step-along,hir:monomial-invariant-along,hir:monomial-invariant,hir:resolution-driver,hir:test-tower,hir:giraud-along}
\source{wlodarczyk05-hironaka,bierstone-milman08-functoriality}
\lean{TranscendentalNakamaye.Pending.MonomialStepAlong.towerInvAlong,TranscendentalNakamaye.Pending.MonomialStepAlong.invariants_canonicalSegment_along,TranscendentalNakamaye.Pending.MonomialStepAlong.exists_isGeoResolution_of_monomialAlong}
\leanok
Let \(M\) be compact, \(n\ge1\), and \(\underline{\mathcal I}=(M,N,E,\mathfrak a,d)\) well formed with \(d\ge1\), \(\mathfrak a\) nonzero and monomial along \(N\). For a tower \(T\) let
\(\gamma_{N,T}\) be the invariant of Lemma~\ref{hir:monomial-invariant-along} of the transform \(\underline{\mathcal I}_T\) (along \(N_T\)). Along every canonical segment for
\(\gamma_N\) (Theorem~\ref{hir:resolution-driver}) that is a multiple test blow-up, \(\underline{\mathcal I}_T\) is well formed and \(\mathfrak a_T\) is nonzero and monomial along \(N_T\);
and some canonical segment for \(\gamma_N\) is a resolution of \(\underline{\mathcal I}\).
\end{theorem}
\begin{proof}
\leanok
As Theorem~\ref{hir:monomial-resolution}, with Lemma~\ref{hir:monomial-step-along}(5) for the induction and the drop, Lemma~\ref{hir:monomial-invariant-along}(1), (3) for
upper semicontinuity and the admissible maximum loci with empty interior, closed cosupports along \(N_T\) (Lemma~\ref{hir:giraud-along}), and well-founded values
(Lemma~\ref{hir:monomial-invariant}(4)).
\end{proof}

\begin{lemma}[The invariant of the monomial case along a submanifold under open sets, biholomorphisms and empty steps]
\label{hir:monomial-invariance-along}
\uses{hir:monomial-invariant-along,hir:monomial-step-along,hir:monomial-resolution-along,hir:exponent-along,hir:monomial-invariance,hir:open-restriction-marked,hir:open-restriction-tower,hir:biholomorphism-transform,hir:tower-transport,hir:biholomorphism-blowup}
\source{wlodarczyk05-hironaka,bierstone-milman08-functoriality}
\lean{TranscendentalNakamaye.Pending.MonomialInvariantFunctorialAlong.isDivNearAlong_pullback_iff,TranscendentalNakamaye.Pending.MonomialInvariantFunctorialAlong.divExpAlong_pullback,TranscendentalNakamaye.Pending.MonomialInvariantFunctorialAlong.expDataAlong_pullback,TranscendentalNakamaye.Pending.MonomialInvariantFunctorialAlong.gammaAlong_pullback,TranscendentalNakamaye.Pending.MonomialInvariantFunctorialAlong.gammaAlong_pullback_val,TranscendentalNakamaye.Pending.MonomialInvariantFunctorialAlong.gammaAlong_pullback_diffeomorph,TranscendentalNakamaye.Pending.MonomialInvariantFunctorialAlong.gammaAlong_geoTransform_empty,TranscendentalNakamaye.Pending.MonomialInvariantFunctorialAlong.towerInvAlong_restrict,TranscendentalNakamaye.Pending.MonomialInvariantFunctorialAlong.towerInvAlong_transport}
\leanok
Let \(\gamma_N=\gamma_{N,\mathfrak a,E,d}\) be the invariant of Lemma~\ref{hir:monomial-invariant-along}, read off the members of \(E\) through a point of \(N\) and
the exponents along \(N\), \(e^N_H\), of Lemma~\ref{hir:exponent-along}.
\begin{enumerate}
\item If \(\varphi\colon M'\to M\) satisfies \(\operatorname{cosupp}((\varphi^*\mathfrak a)|_{\varphi^{-1}N},e)=\varphi^{-1}\operatorname{cosupp}(\mathfrak a|_N,e)\) for
every \(e\) and maps the neighbourhoods of \(x'\) onto neighbourhoods of \(\varphi(x')\), then
\(e^{\varphi^{-1}N}_{\varphi^{-1}H}(\varphi^*\mathfrak a)(x')=e^N_H(\mathfrak a)(\varphi(x'))\) for every \(H\), and
\(\gamma_{\varphi^{-1}N,\varphi^*\mathfrak a,\varphi^{-1}E,d}(x')=\gamma_{N,\mathfrak a,E,d}(\varphi(x'))\).
\item In particular \(\gamma_{N\cap U}\) of the restriction \(\underline{\mathcal I}|_U\) to an open \(U\) at \(p\) is \(\gamma_N\) of \(\underline{\mathcal I}\) at \(p\), and
\(\gamma_{\Phi^{-1}N}\) of \(\Phi^*\underline{\mathcal I}\) at \(p\) is \(\gamma_N\) of \(\underline{\mathcal I}\) at \(\Phi(p)\) for a biholomorphism \(\Phi\).
\item If \(n\ge1\), \(\underline{\mathcal I}=(M,N,E,\mathfrak a,d)\) is well formed, nonzero and monomial along \(N\), and \(\sigma\) is the blowup along \(\emptyset\), then
\(\gamma'_{N'}(z)=\gamma_N(\sigma(z))\) for \(z\in N'\).
\item Along the restriction to \(U\) of a multiple test blow-up \(T\) of a marked ideal \(\underline{\mathcal I}\) with \(N\) a closed submanifold and \(\mathfrak a\) coherent, the
invariants along \(N_T\) of the transforms of \(\underline{\mathcal I}|_U\) and of \(\underline{\mathcal I}\) agree at the points over \(U\); along the transport of a tower by a
biholomorphism of the base, they correspond through the last-stage biholomorphism.
\end{enumerate}
\end{lemma}
\begin{proof}
\leanok
As Lemma~\ref{hir:monomial-invariance}, along \(N\). (1) The exponent \(e^N_H(\mathfrak a)(x)\) is the largest \(e\) such that, near \(x\), the points of \(N\cap H\) lie in
\(\operatorname{cosupp}(\mathfrak a|_N,e)\); under the hypotheses this condition at \(x'\) for \((\varphi^{-1}N,\varphi^*\mathfrak a,\varphi^{-1}H)\) is the condition at
\(\varphi(x')\) for \((N,\mathfrak a,H)\). The members of \(\varphi^{-1}E\) through \(x'\) are the preimages of the members of \(E\) through \(\varphi(x')\), with the same
indices; so the data defining \(\gamma\) agree. (2) The cosupports of the marked ideals \((M,N,\emptyset,\mathfrak a,e)\) pull back along the inclusion of an open set
(Lemma~\ref{hir:open-restriction-marked}) and along a biholomorphism (Lemma~\ref{hir:biholomorphism-transform}); both maps are open embeddings. (3) The empty
centre is admissible, so the transform is well formed with coherent ideal, and \(\sigma(z)\notin\emptyset\): this is the off-centre case
Lemma~\ref{hir:monomial-step-along}(4). (4) The transforms along the restricted tower are the restrictions of the transforms
(Lemma~\ref{hir:open-restriction-tower}) and, from the first blowup on, the transforms along the transported tower are equal to those along \(T\)
(Lemma~\ref{hir:tower-transport}); apply (2).
\end{proof}


\subsection{The resolution algorithm: marked ideals of maximal order}

\begin{lemma}[Members of \(E\) away from the cosupport]
\label{hir:drop-members}
\uses{hir:geo-transform-basic,hir:test-tower,hir:snc,hir:open-restriction-marked,bcalc:transform-off-centre-tower,bcalc:sub-divisor}
\source{wlodarczyk05-hironaka}
\lean{TranscendentalNakamaye.Pending.DropMembers.hasNormalCrossings_of_forall_disjoint,TranscendentalNakamaye.Pending.DropMembers.geoTransformAlong_withE,TranscendentalNakamaye.Pending.DropMembers.cosupp_geoTransformAlong_withE,TranscendentalNakamaye.Pending.DropMembers.isAdmissibleCentre_geoTransformAlong_withE_iff,TranscendentalNakamaye.Pending.DropMembers.isGeoTestTower_withE_iff_of_disjoint,TranscendentalNakamaye.Pending.DropMembers.disjoint_cosupp_of_disjoint}
\leanok
Let \(\underline{\mathcal I}=(M,N,E,\mathfrak a,d)\) with \(N\) a closed submanifold and \(\mathfrak a\) coherent, and let \(E^\circ\) be a list of members of \(E\) such that
every member of \(E\) missing from \(E^\circ\) is closed and disjoint from \(\operatorname{cosupp}\underline{\mathcal I}\) (for instance \(E\) without the members
disjoint from the cosupport). Then \(\underline{\mathcal I}\) and \(\underline{\mathcal I}^\circ=(M,N,E^\circ,\mathfrak a,d)\) have the same multiple test blow-ups;
along every tower the two transforms have the same submanifold, ideal and cosupport; and along a multiple test blow-up the transforms of the deleted
members stay disjoint from the cosupports.
\end{lemma}
\begin{proof}
\leanok
The submanifold, the ideal and the mark of a transform do not depend on \(E\) (induction on the tower), hence neither does the cosupport; only
admissibility does, through the normal crossings of the centre with the divisor. Along any tower every member of the transformed \(E\) missing from the
transformed \(E^\circ\) is the strict transform of a deleted member \(H_0\), and lies over \(H_0\) (a strict transform
\(\overline{\pi^{-1}(K\setminus C)}\) lies in the closed set \(\pi^{-1}(K)\)). Along a multiple test blow-up the cosupports lie over
\(\operatorname{cosupp}\underline{\mathcal I}\) (Corollary~\ref{bcalc:transform-off-centre-tower}), which misses \(H_0\); so every centre (contained in the cosupport)
misses the transforms of the deleted members. Normal crossings of a centre \(C\) with a divisor only involve the members through points of \(C\): with fewer
members it persists (Lemma~\ref{bcalc:sub-divisor}), and members missing \(C\) may be added (at a point \(x\in C\) a member through \(x\) is a member of the
smaller list, with the same coordinate). By induction on the tower the multiple test blow-ups coincide.
\end{proof}

\begin{lemma}[Deleting the members that contain every centre]
\label{hir:drop-containing-members}
\uses{bcalc:nc-comparison,bcalc:sub-divisor,bcalc:nc-comparison-transform,hir:drop-members,hir:snc,mc:strict-along,hir:geo-transform-basic,hir:test-tower,hir:wellformed-tower}
\source{wlodarczyk05-hironaka}
\lean{TranscendentalNakamaye.Pending.DropMembers.isGeoTestTower_withE_iff_of_centresIn_of_snc,TranscendentalNakamaye.Pending.DropMembers.isGeoTestTower_withE_iff_of_centresIn}
\leanok
Let \(\underline{\mathcal I}=(M,N,E,\mathfrak a,d)\) be well formed, \(E'\) a list of members of \(E\), and \(S\) a set contained in every member of \(E\) missing from
\(E'\) (for instance \(S=N\cap\bigcap_{H\in A}H\) and \(E'=E_S\) the sequence \(E\) without the members of a set \(A\)). For towers whose centres lie in the
strict transforms \(S_T\) of \(S\): \(T\) is a multiple test blow-up of \(\underline{\mathcal I}\) if and only if it is one of
\(\underline{\mathcal I}'=(M,N,E',\mathfrak a,d)\), and along it the two transforms have the same submanifold, ideal and cosupport (Lemma~\ref{hir:drop-members});
their sequences differ by members containing \(S_T\).
\end{lemma}
\begin{proof}
\leanok
Induction on the tower. The members of the transformed \(E'\) are members of the transformed \(E\) (strict transforms of members of \(E'\), and the exceptional
divisor, which is in both). A member of the transformed \(E\) missing from the transformed \(E'\) is the strict transform of a member \(H\) missing from \(E'\),
and \(S_T\subset H_T\) by induction; strict transforms are monotone, so it contains \(S_{T'}\) after the next blowup. A centre \(C\subset S_T\) therefore
lies in every member missing from the transformed \(E'\), and, the transform of \(\underline{\mathcal I}\) being well formed along a multiple test blow-up
(Lemma~\ref{hir:wellformed-tower}), its divisor has simple normal crossings and the normal-crossings comparison (Theorem~\ref{bcalc:nc-comparison}) shows that \(C\)
has only normal crossings with one divisor if and only if with the other. Submanifold, ideal and cosupport agree by Lemma~\ref{hir:drop-members}, so \(C\) is
admissible for one transform if and only if for the other. Well-formedness is used only for the simple normal crossings of the transformed \(E\) along
the multiple test blow-ups of \(\underline{\mathcal I}\); the lemma holds whenever these hold.
\end{proof}

\begin{lemma}[Homogenized coefficient ideal]
\label{hir:homogenized-coefficient}
\uses{mc:homogenized-equivalent,mc:homogenized-maximal-order,mc:homogenized-tower,mc:homogenized-etale,mc:derivative-open-restriction,mc:pullback-finite-sum,mc:homogenized-ideal,mc:homogenized-basic,mc:coefficient-ideal,mc:coefficient-equivalent,mc:inclusions-tower,mc:maximal-order,mc:maximal-order-derivative,hir:marked-equivalence-basic,hir:marked-operations-transform,hir:controlled-operations,hir:derivative-invariance,hir:open-restriction-ideal,mc:homogenized-along,mc:homogenized-along-basic,mc:maximal-order-along,mc:homogenized-along-maximal-order,mc:homogenized-along-etale,mc:homogenized-tower-along,mc:homogenized-equivalent-along,mc:coefficient-equivalent-along}
\source{wlodarczyk05-hironaka}
\lean{TranscendentalNakamaye.Pending.HomogenizedEquivalence.hom_wEquiv,TranscendentalNakamaye.Pending.HomogenizedEquivalence.coeff_hom_wEquiv,TranscendentalNakamaye.Pending.HomogenizedEquivalence.isMaximalOrder_coeff_hom,TranscendentalNakamaye.Pending.HomogenizedEtale.pullback_homogenizedIdeal_of_inverse,TranscendentalNakamaye.Pending.HomogenizedEtale.pullback_coeffMarked_of_inverse,TranscendentalNakamaye.Pending.HomogenizedEtale.pullback_val_homogenizedIdeal,TranscendentalNakamaye.Pending.HomogenizedEtale.pullback_val_coeffMarked}
\leanok
Let \(\underline{\mathcal J}=(M,M,E,\mathfrak j,\mu)\) be of maximal order, \(\mathfrak j\) coherent, \(\mu\ge1\), and let
\(\overline{\underline{\mathcal J}}=\mathcal C(\mathcal H(\underline{\mathcal J}))\) be the coefficient ideal of the homogenized ideal. Then
\(\mathcal H(\underline{\mathcal J})\simeq\underline{\mathcal J}\), \(\overline{\underline{\mathcal J}}\simeq\underline{\mathcal J}\) (Definition~\ref{hir:marked-equivalence}),
\(\overline{\underline{\mathcal J}}\) is of maximal order, and \(\mathcal H\) and \(\mathcal C\) commute with restriction to open sets and with biholomorphisms between
open sets.
\end{lemma}
\begin{proof}
\leanok
Since \(\mathcal H(\mathfrak j)\supset\mathfrak j\), every multiple test blow-up of \(\mathcal H(\underline{\mathcal J})\) is one of \(\underline{\mathcal J}\), and along it the
transform of \(\mathcal H(\mathfrak j)\) contains that of \(\mathfrak j\). Conversely, by Lemma~\ref{mc:inclusions-tower} a multiple test blow-up of
\(\underline{\mathcal J}\) is one of every \((\mathcal D^i\mathfrak j,\mu-i)\), \(i<\mu\), and of \((T(\mathfrak j),1)\), \(T(\mathfrak j)=\mathcal D^{\mu-1}\mathfrak j\); by
Lemma~\ref{hir:marked-operations-transform} it is one of each product \((\mathcal D^i\mathfrak j\cdot T(\mathfrak j)^i,\mu)\), whose transform is the product of
the transforms. The ideal \(\mathcal H(\mathfrak j)\) is the sum of these products, all with mark \(\mu\); the order of a sum is the minimum of the orders, and
controlled transforms with a common mark commute with sums (Lemma~\ref{hir:controlled-operations}). So the tower is a multiple test blow-up of
\(\mathcal H(\underline{\mathcal J})\) and the cosupport of its transform is the intersection of the cosupports of the transformed products, which contains
\(\operatorname{cosupp}\underline{\mathcal J}_T\) (each factor's cosupport does, by Lemma~\ref{mc:inclusions-tower}); with the first inclusion the cosupports
agree.
Then \(\mathcal H(\underline{\mathcal J})\) is of maximal order (\(\mathcal D^\mu\mathcal H(\mathfrak j)\supset\mathcal D^\mu\mathfrak j\)), Lemma~\ref{mc:coefficient-equivalent}
gives \(\mathcal C(\mathcal H(\underline{\mathcal J}))\simeq\mathcal H(\underline{\mathcal J})\), and transitivity (Lemma~\ref{hir:marked-equivalence-basic}) gives the
second equivalence. A sum one of whose summands is of maximal order is of maximal order (the order of a sum is the minimum of the scaled orders). Both
constructions are built from \(\mathcal D\), products and sums, which commute with restriction to open sets and with biholomorphisms
(Lemma~\ref{hir:derivative-invariance}).
For \(N=M\) the two equivalences, the transform formula along multiple test blow-ups and
the maximal order are Lemmas~\ref{mc:homogenized-equivalent}, \ref{mc:homogenized-tower} and~\ref{mc:homogenized-maximal-order}, and the commutation of
\(\mathcal H\) and \(\mathcal C\) with restriction to open sets and with biholomorphisms (``le:etale2'', ``le:etale3'') is Lemma~\ref{mc:homogenized-etale}
(a biholomorphism between open sets is a restriction followed by a biholomorphism). Proposition~\ref{hir:maximal-order-step} uses the lemma along
\(N\ne M\) (Lemma~\ref{hir:along-calculus}(4)). Along \(N\), at one stage: \(\mathcal H_N\), \(\mathcal C_N\), their coherence, cosupports and
maximal order are Lemmas~\ref{mc:homogenized-along}, \ref{mc:homogenized-along-basic}, \ref{mc:maximal-order-along} and~\ref{mc:homogenized-along-maximal-order},
their invariance under open restriction and biholomorphisms preserving \(N\) is Lemma~\ref{mc:homogenized-along-etale}; along multiple test blow-ups,
the transform formula for \(\mathcal H_N\) and the equivalences \(\mathcal H_N(\underline{\mathcal J})\simeq\underline{\mathcal J}\),
\(\mathcal C_N(\underline{\mathcal J})\simeq\underline{\mathcal J}\) are Lemmas~\ref{mc:homogenized-tower-along}, \ref{mc:homogenized-equivalent-along} and
\ref{mc:coefficient-equivalent-along}.
\end{proof}

\begin{lemma}[Restriction of the coefficient ideal to a submanifold]
\label{hir:coefficient-restriction}
\uses{mc:coefficient-ideal,mc:coefficient-equivalent,mc:inclusions-tower,hir:geo-transform,hir:marked-equivalence,hir:snc,mc:strict-along,hir:restricted-equivalence-basic,mc:restrict-marked,mc:restrict-tower,mc:coefficient-wequiv,mc:coefficient-restriction-support,mc:coefficient-restriction-tower,mc:coefficient-restriction-full,hir:drop-containing-members,hir:wellformed-tower}
\source{wlodarczyk05-hironaka}
\lean{TranscendentalNakamaye.Pending.CoefficientRestrictionDrop.centresIn_of_isGeoTestTower_withN_withE,TranscendentalNakamaye.Pending.CoefficientRestrictionDrop.coeff_restrict_withE_iff}
\leanok
Let \(\underline{\mathcal J}=(M,M,E,\mathfrak j,\mu)\) be well formed with \(\mu\ge1\), let \(S\subset M\) be a closed submanifold, let \(E_S\) be a list of members of \(E\)
such that every member of \(E\) missing from \(E_S\) contains \(S\) (for instance the members of \(E\) not containing \(S\)), and let
\(\underline{\mathcal S}=(M,S,E_S,\mathcal C(\mathfrak j),\mu_{\mathcal C})\) with the ideal and mark of the coefficient ideal. Then
\(\operatorname{cosupp}\underline{\mathcal S}=\operatorname{cosupp}\underline{\mathcal J}\cap S\), and a tower is a multiple test blow-up of \(\underline{\mathcal S}\) if and only
if it is a multiple test blow-up of \(\underline{\mathcal J}\) whose centres lie in the strict transforms \(S_T\) of \(S\); along it
\(\operatorname{cosupp}\underline{\mathcal S}_T=\operatorname{cosupp}\underline{\mathcal J}_T\cap S_T\).
\end{lemma}
\begin{proof}
\leanok
With \(E\) in place of \(E_S\) this is Włodarczyk's Lemmas ``le:cf'', ``le:S'' and ``le:S2'' (Lemma~\ref{mc:coefficient-restriction-full}, with
Lemmas~\ref{mc:coefficient-restriction-support} and~\ref{mc:coefficient-restriction-tower}): the restriction of \(\mathcal C(\mathfrak j)\) to \(S\) is the marked
ideal with submanifold \(S\), and by Lemma~\ref{hir:restricted-equivalence-basic} its transform along a centre in \(S\cap\operatorname{cosupp}\) is the restriction to
\(S'\) of the controlled transform; the proof of the source (expansion of the germs of \(\mathfrak j\) in coordinates normal to \(S\), whose coefficients are, up to
factorials, derivatives lying in \(\mathcal D^{|\alpha|}\mathfrak j\), followed by the monomial charts of the blowup) applies verbatim to convergent expansions.
The centres of a multiple test blow-up of \(\underline{\mathcal S}\) lie in its submanifolds, the strict transforms \(S_T\) (the submanifold of a transform does
not depend on the divisor). Along towers with centres in \(S_T\), the multiple test blow-ups of \(\underline{\mathcal S}\) and of the marked ideal with \(E\) in place
of \(E_S\) coincide and their cosupports agree (Lemma~\ref{hir:drop-containing-members}); the simple normal crossings of the transformed divisor that this uses
hold because these towers are multiple test blow-ups of the well-formed \(\underline{\mathcal J}\) (Lemma~\ref{hir:wellformed-tower}). The case of the empty tower
is the support statement.
In Step 1a (Proposition~\ref{hir:maximal-order-step}) the lemma is applied to \(\overline{\underline{\mathcal J}}=\mathcal C(\mathcal H(\underline{\mathcal J}))\), of maximal
order, and to \(S\) an intersection of members of \(E\). Along a submanifold \(N\ne M\) (Step 1a inside \(N\)) the single-stage cosupport identity
\(\operatorname{cosupp}(\mathcal C_N(\underline{\mathcal J})|_S)=\operatorname{cosupp}\underline{\mathcal J}\cap S\) for \(S\subset N\) is Lemma~\ref{mc:coefficient-restriction-along} (TN23, 15bf903a); the tower form
along \(N\) is part of Lemma~\ref{hir:along-calculus}(4).
\end{proof}

\begin{lemma}[Reduction to a hypersurface of maximal contact]
\label{hir:contact-equivalence}
\uses{mc:contact-equivalence,mc:contact-equivalence-open,mc:tangent-direction-open,hir:coefficient-restriction,mc:giraud-tower,mc:tangent-direction,mc:tangent-direction-support,bcalc:hypersurface,hir:open-restriction-marked,hir:marked-equivalence,mc:coefficient-restriction-along,mc:glueing-lemma-along}
\source{wlodarczyk05-hironaka}
\lean{TranscendentalNakamaye.Pending.ContactEquivalence.contact_wEquiv,TranscendentalNakamaye.Pending.ContactEquivalenceOpen.contact_wEquiv_open}
\leanok
Let \(\underline{\mathcal J}=(M,M,\emptyset,\mathfrak j,\mu)\) be of maximal order, \(\mathfrak j\) coherent, \(\mu\ge1\), and let \(u\) be a tangent direction of
\((\mathfrak j,\mu)\) on an open set \(U\) such that \(\operatorname{cosupp}\underline{\mathcal J}\cap U\) contains no nonempty open subset of \(V(u)\). Then
\(\underline{\mathcal J}|_U\simeq(U,V(u)\cap U,\emptyset,\mathcal C(\mathfrak j),\mu_{\mathcal C})\).
\end{lemma}
\begin{proof}
\leanok
By Lemma~\ref{mc:giraud-tower} every admissible centre along every multiple test blow-up of \(\underline{\mathcal J}|_U\) lies in the strict transform of
\(V(u)\cap U\); every multiple test blow-up of the right-hand side has its centres in its submanifolds, which are these strict transforms. Lemma~\ref{hir:coefficient-restriction}
with \(S=V(u)\cap U\) (a closed submanifold of \(U\), Lemma~\ref{bcalc:hypersurface}) identifies the multiple test blow-ups and the cosupports, the
cosupports of \(\underline{\mathcal J}|_U\) lying in the strict transforms.
This is Lemma~\ref{mc:contact-equivalence-open} (with Lemma~\ref{mc:contact-equivalence} for \(U=M\)), for any
\(E\) and without the hypothesis on the open subsets of \(V(u)\), which the proof does not use. Proposition~\ref{hir:maximal-order-step} uses it along
\(N\ne M\) (Lemma~\ref{hir:along-calculus}(3)--(4)): the cosupport identity for \(S_w\) with submanifold
\(V(w)\cap N\cap U\) is Lemma~\ref{mc:coefficient-restriction-along}, its form along multiple test blow-ups Lemma~\ref{mc:coefficient-restriction-full-along}, and the
glueing of two tangent directions along \(N\) is Lemma~\ref{mc:glueing-lemma-along}, lifted along multiple test blow-ups by Lemma~\ref{mc:glueing-lift-along}.
\end{proof}
\subsection{The resolution algorithm: invariants, canonicity and termination}

\begin{lemma}[Upper semicontinuous invariants]
\label{hir:invariant-order}
\source{bierstone-milman08-functoriality,kollar07-resolution,encinas-hauser02-strong-resolution}
\lean{TranscendentalNakamaye.Pending.ResolutionInvariantOrder.IsInvUSCOn,TranscendentalNakamaye.Pending.ResolutionInvariantOrder.IsInvUSCOn.mono,TranscendentalNakamaye.Pending.ResolutionInvariantOrder.IsInvUSCOn.comp,TranscendentalNakamaye.Pending.ResolutionInvariantOrder.IsInvUSCOn.exists_isMaxOn,TranscendentalNakamaye.Pending.ResolutionInvariantOrder.IsInvUSCOn.isClosed_superlevel,TranscendentalNakamaye.Pending.ResolutionInvariantOrder.IsInvUSCOn.isClosed_maxLocus,TranscendentalNakamaye.Pending.ResolutionInvariantOrder.maxLocus_inter_eq,TranscendentalNakamaye.Pending.ResolutionInvariantOrder.exists_not_of_lt_succ,TranscendentalNakamaye.Pending.ResolutionInvariantOrder.wellFoundedLT_lex}
\leanok
Let \(\Lambda\) be a totally ordered set, \(s\) a subset of a topological space, and \(f\colon s\to\Lambda\) \emph{upper semicontinuous on \(s\)}: every \(x\in s\)
has a neighbourhood whose points in \(s\) have values at most \(f(x)\). This property passes to subsets and to pullbacks along continuous maps. If \(s\) is
compact and nonempty, \(f\) attains its maximum on \(s\). If \(s\) is closed, the sets \(\{x\in s:f(x)\ge\lambda\}\) are closed, in particular the maximum locus;
on a set on which \(f\le\lambda\), the points of value \(\lambda\) form the trace of \(\{f\ge\lambda\}\). If \(\Lambda\) is well ordered and \(P(0),P(1),\dots\) are
properties such that \(m_{i+1}<m_i\) whenever \(P(i)\) and \(P(i+1)\) hold, then some \(P(i)\) fails. Tuples of a fixed length of elements of a well-ordered set,
compared lexicographically, form a well-ordered set.
\end{lemma}
\begin{proof}
\leanok
Cover \(s\) by finitely many neighbourhoods \(U_i\) of points \(x_i\in s\) with \(f\le f(x_i)\) on \(U_i\cap s\); the largest \(f(x_i)\) is the maximum. If \(y\notin s\), a
neighbourhood of \(y\) misses the closed set \(s\); if \(y\in s\) and \(f(y)<\lambda\), then \(f<\lambda\) on a neighbourhood of \(y\) in \(s\); so \(\{f\ge\lambda\}\) is closed.
The maximum locus is the superlevel set of the maximum. Restriction and pullback are immediate from the definition. If every \(P(i)\) held, \((m_i)\) would be a
strictly decreasing sequence in a well-ordered set. Lexicographic products of finitely many well orders are well ordered.
\end{proof}
Encinas and Hauser take their invariant in \(\mathbf N^{4n}\) with the lexicographic order and conclude termination from its strict decrease under the
blowup of the top locus and the well ordering (``Decrease of invariant'', ``Resolution of mobiles''). The invariant below takes values in a totally ordered
set that need not be well ordered (the order on subsets of \(E\) is not), and termination is proved separately, as local finiteness
(Definition~\ref{hir:centre-function}(8)).

\begin{definition}[Isomorphism of towers]
\label{hir:tower-isomorphism}
\uses{hir:blowup-tower,hir:tower-transport,hir:biholomorphism-blowup}
\source{wlodarczyk05-hironaka}
\lean{TranscendentalNakamaye.Pending.BlowupTowerIsomorphism.IsTowerIso,TranscendentalNakamaye.Pending.BlowupTowerIsomorphism.transport_isTowerIso}
\leanok
Let \(T\colon X=X_0\xleftarrow{\sigma_1}\cdots\xleftarrow{\sigma_r}X_r\) and \(T'\colon X'=X'_0\xleftarrow{\sigma'_1}\cdots\xleftarrow{\sigma'_r}X'_r\) be towers of the same
length with centres \(C_i\), \(C'_i\). An \emph{isomorphism} \(\chi\colon T'\to T\) over a biholomorphism \(\Phi\colon X'\to X\) is a family of biholomorphisms
\(\chi_i\colon X'_i\to X_i\) with \(\chi_0=\Phi\), \(\sigma_i\circ\chi_i=\chi_{i-1}\circ\sigma'_i\) and \(\chi_{i-1}^{-1}(C_{i-1})=C'_{i-1}\) for \(i\ge1\). The transport
\(\Phi^*T\) (Definition~\ref{hir:tower-transport}) is isomorphic to \(T\) over \(\Phi\) with \(\chi_i=\mathrm{id}\) for \(i\ge1\); the stage-by-stage lifts of a
biholomorphism between open sets along a segment (Theorem~\ref{mc:glueing-lift}, Lemma~\ref{mc:contact-independence-along}) form an isomorphism between the
restricted segments.
\end{definition}

\begin{definition}[Resolution invariant; canonical segments]
\label{hir:centre-function}
\uses{hir:test-tower,hir:geo-transform,hir:tower-transport,hir:tower-isomorphism,hir:open-restriction-tower,hir:open-restriction-marked,hir:restricted-equivalence,hir:invariant-order,hir:resolution-driver}
\source{bierstone-milman08-functoriality,wlodarczyk05-hironaka}
\lean{TranscendentalNakamaye.Pending.ResolutionInvariantSpec.DimLE,TranscendentalNakamaye.Pending.ResolutionInvariantSpec.IsInvDomain,TranscendentalNakamaye.Pending.ResolutionInvariantSpec.IsCanonical,TranscendentalNakamaye.Pending.ResolutionInvariantSpec.countCentresMeeting,TranscendentalNakamaye.Pending.ResolutionInvariantExt.IsEmptyTower,TranscendentalNakamaye.Pending.ResolutionInvariantExt.ResolutionInvariant}
\leanok
Fix \(n\) and \(k\le n\) and a totally ordered set \(\Lambda\). A \emph{resolution invariant in dimension \(\le k\)} assigns to every complex manifold \(X\) of
dimension \(n\), every well-formed marked ideal \(\underline{\mathcal I}=(X,N,E,\mathfrak a,d)\) with \(\dim N\le k\), \(d\ge1\) and \(\mathfrak a|_N\) vanishing identically on no
nonempty open subset of \(N\), every tower \(T\) over \(X\) and every \(x\in\operatorname{cosupp}\underline{\mathcal I}_T\) a value
\(\operatorname{inv}_{\underline{\mathcal I}}(T)(x)\in\Lambda\). A tower \(T\) is a \emph{canonical segment} of \(\underline{\mathcal I}\) if each of its centres is empty or
equal to the set where \(\operatorname{inv}_{\underline{\mathcal I}}\) of the preceding part of \(T\) attains its maximum on the cosupport. The following are
required for the marked ideals of the domain, along every canonical segment \(T\) of \(\underline{\mathcal I}\) and at the points of
\(\operatorname{cosupp}\underline{\mathcal I}_T\).
\begin{enumerate}
\item (Admissibility) For a canonical segment \(T\) and an open set \(W\) of its last stage on which \(\operatorname{inv}_{\underline{\mathcal I}}(T)\le\lambda\), the set of
points of \(W\cap\operatorname{cosupp}\underline{\mathcal I}_T\) where the value is \(\lambda\) is an admissible centre for \(\underline{\mathcal I}_T|_W\).
\item (Semicontinuity) \(\operatorname{inv}_{\underline{\mathcal I}}(T)\) is upper semicontinuous on \(\operatorname{cosupp}\underline{\mathcal I}_T\)
(Lemma~\ref{hir:invariant-order}).
\item (Decrease) If \(T\) is a canonical segment and \(T'\) is \(T\) followed by the blowup \(\sigma\) along the maximum locus \(C\ne\emptyset\), then at every point \(z\)
of \(\operatorname{cosupp}\underline{\mathcal I}_{T'}\) the value is smaller than \(\operatorname{inv}_{\underline{\mathcal I}}(T)(\sigma(z))\) if \(\sigma(z)\in C\), and equal to
it otherwise.
\item (Empty steps) If \(T'\) is \(T\) followed by a blowup \(\pi\) along a centre \(C=\emptyset\), \(\operatorname{inv}_{\underline{\mathcal I}}(T')=\operatorname{inv}_{\underline{\mathcal I}}(T)\circ\pi\)
on \(\operatorname{cosupp}\underline{\mathcal I}_{T'}\).
\item[(4')] (Empty prefixes) If all centres of \(T_0\) are empty and \(T_0T_1\) is a canonical segment, then
\(\operatorname{inv}_{\underline{\mathcal I}}(T_0T_1)=\operatorname{inv}_{\underline{\mathcal I}_{T_0}}(T_1)\) on \(\operatorname{cosupp}\underline{\mathcal I}_{T_0T_1}\): the
invariant may be re-based across a part of the segment consisting of isomorphisms.
\item (Open sets) For \(U\subset X\) open, \(\operatorname{inv}_{\underline{\mathcal I}|_U}(T|_U)\) is the restriction of \(\operatorname{inv}_{\underline{\mathcal I}}(T)\) to the
preimage of \(U\) (Lemma~\ref{hir:open-restriction-tower}).
\item (Biholomorphisms) For a biholomorphism \(\Phi\colon X'\to X\) and an isomorphism \(\chi\colon T'\to T\) of towers over \(\Phi\)
(Definition~\ref{hir:tower-isomorphism}), \(\operatorname{inv}_{\Phi^*\underline{\mathcal I}}(T')=\operatorname{inv}_{\underline{\mathcal I}}(T)\circ\chi_r\); in particular
\(\operatorname{inv}_{\Phi^*\underline{\mathcal I}}(\Phi^*T)=\operatorname{inv}_{\underline{\mathcal I}}(T)\circ\Phi_T\) for the transported tower and its last-stage biholomorphism
(Definition~\ref{hir:tower-transport}).
\item (Representatives) Restricted-equivalent marked ideals (Definition~\ref{hir:restricted-equivalence}) have the same invariant.
\item (Local finiteness) For every compact \(K\subset X\) there is \(m\) such that every canonical segment of \(\underline{\mathcal I}\) has at most \(m\) nonempty centres
meeting the preimage of \(K\).
\end{enumerate}
\end{definition}
By (5), the restriction to an open set \(U\) of a canonical segment of \(\underline{\mathcal I}\) is a canonical segment of \(\underline{\mathcal I}|_U\): a restricted centre
is either empty, or the maximum of the restricted invariant is the global maximum and the restricted centre is its maximum locus. This is Włodarczyk's
notion of an \emph{extension} of a resolution, with steps of empty centre in place of isomorphisms (Lemma~\ref{hir:canonical-segments}). In the source,
extensions insert isomorphisms anywhere in a resolution, and the invariants of points transformed isomorphically do not change; (4') is the form of this
the inductive step needs: in the maximal-order case and the general step the lower-dimensional invariant at a point is computed along the part of the segment
since a stage depending on the point, and nearby points, whose stages differ, are compared by re-basing across steps that are empty near them.
This is the second Lean form of the definition (independent review, 3 October 2026): the first form (Definition~\ref{hir:centre-function-v1}) asked (2) and
(4)--(7) along every multiple test blow-up, which the inductive step cannot provide (along a non-canonical tower the restriction to an intersection of
members or to a hypersurface of maximal contact need not be a multiple test blow-up of the restricted marked ideal), and lacked (4'), which does not follow
from its eight properties. Property (8) in this uniform form is superseded by its per-sequence form (8\('\)), Definition~\ref{hir:centre-function-seq}
(human decision, 3 October 2026); this definition and its Lean form are kept, superseded on the route: the resolution invariants the route
constructs (Theorem~\ref{hir:canonical-centres}) and consumes (Theorems~\ref{hir:canonical-resolution}, \ref{hir:factorization-from-resolution},
\ref{hir:embedded-from-resolution}) are those of Definition~\ref{hir:centre-function-seq}, and the results stated for this form (Theorem~\ref{hir:invariant-driver},
the first typing of the Step 1 value) are carried over by Corollary~\ref{hir:step-one-seq}.

\begin{definition}[First Lean form of resolution invariants (superseded)]
\label{hir:centre-function-v1}
\uses{hir:centre-function,hir:canonical-centres-zero}
\source{wlodarczyk05-hironaka}
\lean{TranscendentalNakamaye.Pending.ResolutionInvariantSpec.ResolutionInvariant,TranscendentalNakamaye.Pending.ResolutionInvariantSpec.ResolutionInvariant.ofDimZero}
\leanok
Superseded by Definition~\ref{hir:centre-function} and kept: the same data and properties (1)--(8) of Definition~\ref{hir:centre-function} except
(4'), with (2) and (4)--(7) required along every multiple test blow-up of \(\underline{\mathcal I}\) instead of along canonical segments, and (4) stated for
the literal empty centre. In dimension \(\le0\) every constant function is one (as in Lemma~\ref{hir:canonical-centres-zero}). Not used by the route.
\end{definition}

\begin{lemma}[Results for the first Lean form (superseded)]
\label{hir:centre-function-v1-results}
\uses{hir:centre-function-v1,hir:invariant-driver,hir:canonical-segments,hir:step-one-b-independence,hir:resolution-driver,mc:glueing-lift-tail}
\source{wlodarczyk05-hironaka}
\lean{TranscendentalNakamaye.Pending.ResolutionInvariantDriver.isAdmissibleCentre_maxLocus,TranscendentalNakamaye.Pending.ResolutionInvariantDriver.isGeoTestTower_of_isCanonical,TranscendentalNakamaye.Pending.ResolutionInvariantDriver.isInvDomain_of_isCanonical,TranscendentalNakamaye.Pending.ResolutionInvariantDriver.exists_isGeoResolution_isCanonicalSegment,TranscendentalNakamaye.Pending.ResolutionInvariantDriver.exists_isGeoResolution_isCanonical,TranscendentalNakamaye.Pending.ResolutionInvariantSegments.isCanonical_restrict,TranscendentalNakamaye.Pending.ResolutionInvariantSegments.isCanonical_of_isTowerIso,TranscendentalNakamaye.Pending.MaximalOrderStepOneB.inv_contactMarked_eq}
\leanok
For an invariant in the sense of Definition~\ref{hir:centre-function-v1}, the conclusions of Theorem~\ref{hir:invariant-driver}(3)--(4), of
Lemma~\ref{hir:canonical-segments}(1)--(2) and of Lemma~\ref{hir:step-one-b-independence} (the latter with ``multiple test blow-up'' in place of ``canonical
segment'' in its hypothesis) hold. Superseded with Definition~\ref{hir:centre-function-v1}; kept.
\end{lemma}
\begin{proof}
\leanok
The proofs of those statements, which use (1), (2), (5), (6) and (8) only along canonical segments.
\end{proof}

\begin{definition}[Resolution invariant with local finiteness along growing sequences]
\label{hir:centre-function-seq}
\uses{hir:centre-function,hir:canonical-centres-zero}
\source{wlodarczyk05-hironaka,bierstone-milman08-functoriality}
\lean{TranscendentalNakamaye.Pending.ResolutionInvariantSeq.GrowingSeq,TranscendentalNakamaye.Pending.ResolutionInvariantSeq.IsInvValue,TranscendentalNakamaye.Pending.ResolutionInvariantSeq.SeqLocallyFinite,TranscendentalNakamaye.Pending.ResolutionInvariantSeq.ResolutionInvariant,TranscendentalNakamaye.Pending.ResolutionInvariantSeq.isInvValue_ext,TranscendentalNakamaye.Pending.ResolutionInvariantSeq.seqLocallyFinite_ext,TranscendentalNakamaye.Pending.ResolutionInvariantSeq.ResolutionInvariant.ofExt,TranscendentalNakamaye.Pending.ResolutionInvariantSeq.ResolutionInvariant.ofDimZero}
\leanok
A \emph{growing sequence} of towers over \(X\) is a sequence \(T_0,T_1,\dots\) with \(T_0\) the empty tower and \(T_{i+1}\) the tower \(T_i\) followed by one blowup. A
\emph{resolution invariant in dimension \(\le k\)} (the form used by the route) is a value function \(\operatorname{inv}\) as in Definition~\ref{hir:centre-function}
with properties (1)--(7) and (4') of Definition~\ref{hir:centre-function}, and, in place of (8):
\begin{enumerate}
\item[(8\('\))] (Local finiteness along growing sequences) For every marked ideal \(\underline{\mathcal I}\) of the domain, every compact \(K\subset X\) and every growing
sequence \((T_i)\) of canonical segments of \(\underline{\mathcal I}\), only finitely many nonempty centres of the sequence meet the preimage of \(K\): the numbers of
centres of \(T_i\) meeting the preimage of \(K\) are bounded.
\end{enumerate}
Property (8) of Definition~\ref{hir:centre-function} (one bound for all canonical segments) implies (8\('\)), so every resolution invariant in the sense of
Definition~\ref{hir:centre-function} is one in this sense; in particular every constant function is one in dimension \(\le0\)
(Lemma~\ref{hir:canonical-centres-zero}).
\end{definition}
The uniform form (8) cannot be obtained in the inductive step from the other properties: canonical segments branch at every step (a centre may be empty or the
maximum locus), an empty step inserted in the middle of a segment changes the values of the general step (Lemma~\ref{bcalc:monomial-invariant-empty-member}), and a
finitely branching tree may have finitely many marked nodes on every path and still unboundedly many on finite paths. What the inductive step proves is (8\('\)), and
(8\('\)) suffices to resolve (Theorem~\ref{hir:invariant-driver-seq}). Definition~\ref{hir:centre-function} and its Lean form are kept.

\begin{lemma}[Canonical segments on open sets and along isomorphisms of towers]
\label{hir:canonical-segments}
\uses{hir:centre-function,hir:invariant-driver,hir:open-restriction-tower,hir:open-restriction-marked,bcalc:tower-restrict,hir:tower-isomorphism,mc:tower-isomorphism-transform}
\source{wlodarczyk05-hironaka}
\lean{TranscendentalNakamaye.Pending.ResolutionInvariantSegments.maxLocus_preimage,TranscendentalNakamaye.Pending.ResolutionInvariantSegments.maxLocus_preimage_of_surjective,TranscendentalNakamaye.Pending.ResolutionInvariantSegments.cosupp_geoTransformAlong_restrict,TranscendentalNakamaye.Pending.ResolutionInvariantExt.isCanonical_restrict,TranscendentalNakamaye.Pending.ResolutionInvariantExt.isCanonical_of_isTowerIso,TranscendentalNakamaye.Pending.ResolutionInvariantExt.isCanonical_of_isEmptyTower,TranscendentalNakamaye.Pending.ResolutionInvariantExt.isGeoTestTower_of_isEmptyTower,TranscendentalNakamaye.Pending.ResolutionInvariantExt.maxLocus_congr,TranscendentalNakamaye.Pending.ResolutionInvariantExt.isCanonical_append_iff,TranscendentalNakamaye.Pending.ResolutionInvariantSegments.countCentresMeeting_empty,TranscendentalNakamaye.Pending.ResolutionInvariantSegments.countCentresMeeting_mono,TranscendentalNakamaye.Pending.ResolutionInvariantSegments.countCentresMeeting_union_le,TranscendentalNakamaye.Pending.ResolutionInvariantSegments.countCentresMeeting_biUnion_le,TranscendentalNakamaye.Pending.ResolutionInvariantSegments.countCentresMeeting_restrict,TranscendentalNakamaye.Pending.ResolutionInvariantSegments.countCentresMeeting_restrict_of_subset}
\leanok
Let \(\operatorname{inv}\) be a resolution invariant in dimension \(\le k\) and \(\underline{\mathcal I}\) a marked ideal of its domain.
\begin{enumerate}
\item For an open \(U\subset X\), the restriction \(T|_U\) of a canonical segment \(T\) of \(\underline{\mathcal I}\) is a canonical segment of \(\underline{\mathcal I}|_U\); along
\(T|_U\) the cosupport is the trace of that along \(T\).
\item If \(\chi\colon T'\to T\) is an isomorphism of towers over a biholomorphism \(\Phi\colon X'\to X\) (Definition~\ref{hir:tower-isomorphism}) and \(T\) is a canonical
segment of \(\underline{\mathcal I}\), then \(T'\) is a canonical segment of \(\Phi^*\underline{\mathcal I}\).
\item For a tower \(T\) over \(X\), the number \(c_K(T)\) of nonempty centres of \(T\) meeting the preimage of \(K\subset X\) is monotone in \(K\), satisfies
\(c_{K_1\cup\dots\cup K_m}(T)\le\sum_jc_{K_j}(T)\), and \(c_{K\cap U}(T)=c_{K\cap U}(T|_U)\) for an open \(U\subset X\) (the latter count taken in \(U\)).
\item A tower whose centres are all empty is a canonical segment and a multiple test blow-up; if \(T_0\) is such a tower, \(T_0T_1\) is a canonical segment of
\(\underline{\mathcal I}\) if and only if \(T_1\) is a canonical segment of \(\underline{\mathcal I}_{T_0}\).
\end{enumerate}
\end{lemma}
\begin{proof}
\leanok
A maximum locus pulls back along a map preserving the values on the set as soon as the pulled-back locus is nonempty (the maximum is then attained in the
image), and along any surjection. (1) Induction along \(T\); the segment is a multiple test blow-up (Theorem~\ref{hir:invariant-driver}), so the cosupport of
\((\underline{\mathcal I}|_U)_{T|_U}\) is the trace of that of \(\underline{\mathcal I}_T\) (Lemmas~\ref{hir:open-restriction-tower} and~\ref{hir:open-restriction-marked}). A
centre of \(T\) is empty or the maximum locus \(C\); its trace is empty, or nonempty and then, the values agreeing on the trace of the cosupport by
Definition~\ref{hir:centre-function}(5), the maximum locus of the restricted invariant. (2) Induction along the isomorphism: \(C'=\chi^{-1}(C)\), the
cosupports correspond (Lemma~\ref{mc:tower-isomorphism-transform}) and the values are composed with the biholomorphism \(\chi\) by
Definition~\ref{hir:centre-function}(6). (3) Induction along \(T\): a centre meeting the preimage of \(K_1\cup K_2\) meets that of \(K_1\) or of \(K_2\); and a centre
\(C_i\) of \(T|_U\) is the trace of the centre of \(T\) on the preimage of \(U\), whose points over \(K\) are its points over \(K\cap U\) (Lemma~\ref{bcalc:tower-restrict}).
(4) The empty set is an admissible centre. Induction along \(T_1\): when \(T_0T_1'\) (equivalently \(T_1'\)) is canonical, the two invariants agree on the common
cosupport \(\operatorname{cosupp}\underline{\mathcal I}_{T_0T_1'}=\operatorname{cosupp}(\underline{\mathcal I}_{T_0})_{T_1'}\) by Definition~\ref{hir:centre-function}(4'), so
their maximum loci agree.
\end{proof}

\begin{lemma}[Splitting a segment at a stage]
\label{hir:tower-split}
\uses{hir:blowup-tower,hir:geo-transform-basic,hir:test-tower,hir:centre-function,hir:canonical-segments}
\source{wlodarczyk05-hironaka}
\lean{TranscendentalNakamaye.Pending.BlowupTowerSplit.split,TranscendentalNakamaye.Pending.BlowupTowerSplit.head,TranscendentalNakamaye.Pending.BlowupTowerSplit.tail,TranscendentalNakamaye.Pending.BlowupTowerSplit.head_step_succ,TranscendentalNakamaye.Pending.BlowupTowerSplit.tail_step_succ,TranscendentalNakamaye.Pending.BlowupTowerSplit.head_append_tail,TranscendentalNakamaye.Pending.BlowupTowerSplit.length_tail,TranscendentalNakamaye.Pending.BlowupTowerSplit.length_head,TranscendentalNakamaye.Pending.BlowupTowerSplit.proj_eq_comp,TranscendentalNakamaye.Pending.BlowupTowerSplit.geoTransformAlong_head_tail,TranscendentalNakamaye.Pending.BlowupTowerSplit.isGeoTestTower_iff_head_tail,TranscendentalNakamaye.Pending.BlowupTowerSplit.isCanonical_of_append,TranscendentalNakamaye.Pending.BlowupTowerSplit.isCanonical_head,TranscendentalNakamaye.Pending.BlowupTowerSplit.isEmptyTower_append_iff,TranscendentalNakamaye.Pending.BlowupTowerSplit.inv_eq_of_isEmptyTower_head}
\leanok
Every tower \(T\colon X\leftarrow\cdots\leftarrow Y\) splits, for each \(j\ge0\), as \(T=T_0T_1\) with \(T_1\) its last \(\min(j,\operatorname{length}T)\) blowups. Then
\(\operatorname{proj}_T=\operatorname{proj}_{T_0}\circ\operatorname{proj}_{T_1}\), \(\underline{\mathcal I}_T=(\underline{\mathcal I}_{T_0})_{T_1}\), \(T\) is a multiple test
blow-up of \(\underline{\mathcal I}\) iff \(T_0\) is one of \(\underline{\mathcal I}\) and \(T_1\) one of \(\underline{\mathcal I}_{T_0}\), a prefix of a canonical segment is a
canonical segment, and a concatenation has only empty centres iff both parts have. If \(T\) is a canonical segment for a resolution invariant and \(T_0\) has
only empty centres, \(\operatorname{inv}_{\underline{\mathcal I}}(T)=\operatorname{inv}_{\underline{\mathcal I}_{T_0}}(T_1)\) on the cosupport.
\end{lemma}
\begin{proof}
\leanok
Induction on \(j\) and the tower: splitting before the last \(j+1\) blowups of a tower ending with a blowup \(\pi\) splits the tower before \(\pi\) before its last
\(j\) blowups and appends \(\pi\) to the tail. The other statements follow from the concatenation formulas (Lemmas~\ref{hir:blowup-tower}
and~\ref{hir:geo-transform-basic}), the definition of canonical segments by induction along the tower, and Definition~\ref{hir:centre-function}(4').
\end{proof}

\begin{lemma}[Splitting a segment along open sets and isomorphisms of towers]
\label{hir:tower-split-iso}
\uses{hir:tower-split,hir:tower-isomorphism,bcalc:tower-restrict,mc:tower-isomorphism-append,mc:tower-isomorphism-restrict-append,mc:tower-isomorphism-basic}
\source{wlodarczyk05-hironaka}
\lean{TranscendentalNakamaye.Pending.BlowupTowerSplitIso.IsTowerIso.congr_last,TranscendentalNakamaye.Pending.BlowupTowerSplitIso.IsTowerIso.congr_base,TranscendentalNakamaye.Pending.BlowupTowerSplitIso.isTowerIso_restrict_of_append_eq,TranscendentalNakamaye.Pending.BlowupTowerSplitIso.openPreimage_head_tail,TranscendentalNakamaye.Pending.BlowupTowerSplitIso.isTowerIso_restrict_head_tail,TranscendentalNakamaye.Pending.BlowupTowerSplitIso.exists_isTowerIso_split_restrict,TranscendentalNakamaye.Pending.BlowupTowerSplitIso.proj_tail_split_restrict,TranscendentalNakamaye.Pending.BlowupTowerSplitIso.length_tail_restrict,TranscendentalNakamaye.Pending.BlowupTowerSplitIso.IsTowerIso.exists_split}
\leanok
Let \(T\) be a tower over \(X\), \(j\ge0\), and \(T=T_0T_1\) its split before the last \(j\) blowups (Lemma~\ref{hir:tower-split}).
\begin{enumerate}
\item For an open \(U\subset X\), let \(T|_U=T'_0T'_1\) be the split of the restricted tower before its last \(j\) blowups. There is an isomorphism
\(\chi_0\colon T'_0\to T_0|_U\) of towers over \(\mathrm{id}_U\) (Definition~\ref{hir:tower-isomorphism}), and \(T'_1\) is isomorphic over \(\chi_0\) to
\(T_1|_{p_{T_0}^{-1}U}\), with last-stage map the identity of \(p_T^{-1}U\). In particular \(T|_U\) is isomorphic over \(\mathrm{id}_U\) to the concatenation
\(T_0|_U\,T_1|_{p_{T_0}^{-1}U}\), the tails of the two splits have the same length, and \(\chi_0\) maps the image of a point of \(p_T^{-1}U\) in the head stage of
\(T|_U\) to its image in the head stage of \(T\).
\item If \(\chi\colon T'\to T\) is an isomorphism of towers over a biholomorphism \(\Phi\), and \(T'=T'_0T'_1\) is the split of \(T'\) before its last \(j\) blowups,
there is an isomorphism \(\chi_0\colon T'_0\to T_0\) over \(\Phi\), and \(T'_1\) is isomorphic to \(T_1\) over \(\chi_0\) with last-stage map that of \(\chi\).
\end{enumerate}
An isomorphism of towers may be stated with any pointwise equal base or last-stage map.
\end{lemma}
\begin{proof}
\leanok
(1) Induction on \(T\) and \(j\). If \(j=0\) or \(T\) is empty, the head of the split of \(T|_U\) is \(T|_U=T_0|_U\) and both tails are empty; the identity is an
isomorphism (Lemma~\ref{mc:tower-isomorphism-append}(1)). If \(T\) ends with a blowup \(\sigma\) along \(C\) and \(j\ge1\), the head of the split of \(T|_U\) is that of the
restriction of the shorter tower before its last \(j-1\) blowups, and both tails end with the restriction of \(\sigma\) to the open preimage along the trace of \(C\)
(Lemma~\ref{bcalc:tower-restrict}); the inductive isomorphism \(\chi_0\) is kept, and the tail isomorphism is extended by the identity of the common open subset of the
new stage, which commutes with the two restrictions of \(\sigma\) and maps the trace of \(C\) onto itself. The concatenation form is
Lemma~\ref{mc:tower-isomorphism-restrict-append} applied to \(T=T_0T_1\); the lengths are \(\min(j,\operatorname{length}T)\) on both sides
(\(T|_U\) has the length of \(T\)), and the last statement is \(\Phi\circ p_{T'}=p_T\circ\chi\) for the tail isomorphism (Lemma~\ref{mc:tower-isomorphism-basic}) with
\(p_{T_1|_V}=p_{T_1}\) on points. (2) Induction on \(\chi\) and \(j\): for \(j=0\) the heads are the towers and \(\chi_0=\chi\); for empty towers \(\chi_0=\Phi\); at a step
with \(j\ge1\) the head isomorphism comes from the shorter towers and the tail isomorphism is extended by the step of \(\chi\). The last statement: equal maps.
\end{proof}

\begin{lemma}[The stage at which a point entered its regime]
\label{hir:regime-run}
\uses{hir:tower-split,hir:blowup-tower,bcalc:tower-restrict,hir:tower-isomorphism,mc:tower-isomorphism-basic}
\source{wlodarczyk05-hironaka}
\lean{TranscendentalNakamaye.Pending.StepOneRegime.Regime,TranscendentalNakamaye.Pending.StepOneRegime.regimeRun,TranscendentalNakamaye.Pending.StepOneRegime.regimeRun_refl,TranscendentalNakamaye.Pending.StepOneRegime.regimeRun_step,TranscendentalNakamaye.Pending.StepOneRegime.regimeRun_step_of,TranscendentalNakamaye.Pending.StepOneRegime.regimeRun_step_of_not,TranscendentalNakamaye.Pending.StepOneRegime.regimeRun_le_length,TranscendentalNakamaye.Pending.StepOneRegime.regime_head_of_le_regimeRun,TranscendentalNakamaye.Pending.StepOneRegime.not_regime_head_regimeRun,TranscendentalNakamaye.Pending.StepOneRegime.regimeRun_restrict,TranscendentalNakamaye.Pending.StepOneRegime.regimeRun_isTowerIso,TranscendentalNakamaye.Pending.StepOneRegime.eventually_le_regimeRun}
\leanok
A \emph{regime} \(P\) is a property of the points of the stages of the towers over \(X\). For a tower \(T\) and a point \(y\) of its last stage, the \emph{run}
\(r_P(T,y)\) is the number of trailing blowups of \(T\) along which the images of \(y\) before the last stage have property \(P\): \(r_P=0\) for the empty tower, and
after a blowup \(\sigma\), \(r_P(T\sigma,z)=r_P(T,\sigma(z))+1\) if \(\sigma(z)\) has \(P\) and \(0\) otherwise. Write \(T=T_0T_1\) for the split of \(T\) before its last
\(r_P(T,y)\) blowups (Lemma~\ref{hir:tower-split}).
\begin{enumerate}
\item \(r_P(T,y)\le\operatorname{length}T\); the image of \(y\) at every stage of \(T_1\) before the last one has \(P\), and if \(T_0\) is not empty, the image of \(y\) at
the stage before the last blowup of \(T_0\) does not.
\item If \(U\subset X\) is open and \(P'\) is a regime over \(U\) such that a point of \(p_T^{-1}U\) has \(P'\) for \(T|_U\) iff it has \(P\) for \(T\), then
\(r_{P'}(T|_U,y)=r_P(T,y)\) for \(y\in p_T^{-1}U\).
\item If \(\chi\colon T'\to T\) is an isomorphism of towers over \(\Phi\) and \(P'\), \(P\) correspond along every isomorphism of towers over \(\Phi\), then
\(r_{P'}(T',y')=r_P(T,\chi(y'))\).
\item Let \(G\) be a regime inherited by images under blowups such that near each point having \(P\) every point having \(G\) has \(P\). Then near \(y\) every point
having \(G\) has run at least \(r_P(T,y)\).
\end{enumerate}
\end{lemma}
\begin{proof}
\leanok
Induction on \(T\), along the recursion defining the run; the splits before the last \(i+1\) blowups of \(T\sigma\) and before the last \(i\) blowups of \(T\) have the same
head (Lemma~\ref{hir:tower-split}). (2) A blowup of \(T|_U\) is the restriction of the corresponding blowup of \(T\) (Lemma~\ref{bcalc:tower-restrict}), so the images
agree. (3) Induction along the isomorphism, using \(\sigma\circ\chi_i=\chi_{i-1}\circ\sigma'\) (Definition~\ref{hir:tower-isomorphism}). (4) Induction on \(T\): if
\(\sigma(z)\) has \(P\), by continuity of \(\sigma\) the images \(\sigma(z')\) of the points \(z'\) near \(z\) lie in the neighbourhood of \(\sigma(z)\) given by induction and by
the hypothesis; if \(z'\) has \(G\), so has \(\sigma(z')\), which then has \(P\), and \(r_P(T\sigma,z')=r_P(T,\sigma(z'))+1\ge r_P(T,\sigma(z))+1\).
\end{proof}

\begin{lemma}[Dropping the original members of \(E\)]
\label{hir:drop-original-members}
\uses{hir:geo-transform-basic,hir:test-tower,bcalc:sub-divisor,bcalc:nc-comparison-transform,hir:drop-members,mc:homogenized-along,hir:step-one-b-independence}
\source{wlodarczyk05-hironaka}
\lean{TranscendentalNakamaye.Pending.StepOneRegime.geoTransformAlong_dropE,TranscendentalNakamaye.Pending.StepOneRegime.isGeoTestTower_dropE,TranscendentalNakamaye.Pending.StepOneRegime.cosupp_geoTransformAlong_dropE,TranscendentalNakamaye.Pending.StepOneRegime.coeffHom_dropE,TranscendentalNakamaye.Pending.StepOneRegime.geoTransformAlong_coeffHom_dropE,TranscendentalNakamaye.Pending.StepOneRegime.isGeoTestTower_stepOneMarked_of_coeffHom,TranscendentalNakamaye.Pending.StepOneRegime.cosupp_geoTransformAlong_stepOneMarked}
\leanok
Let \(\underline{\mathcal I}=(X,N,E,\mathfrak a,d)\) and \(m\le\#E\), and write \(\underline{\mathcal I}^{(m)}\) for \(\underline{\mathcal I}\) with its first \(m\) members deleted.
Along every tower \(T\), \((\underline{\mathcal I}^{(m)})_T=(\underline{\mathcal I}_T)^{(m)}\); every multiple test blow-up of \(\underline{\mathcal I}\) is one of
\(\underline{\mathcal I}^{(m)}\), with the same cosupports. For \(\underline{\mathcal J}=(X,N,E,\mathfrak a,d)\) and
\(\overline{\underline{\mathcal J}}=\mathcal C_N(\mathcal H_N(\underline{\mathcal J}))\), deleting the \(\#E\) original members gives
\(\mathcal C_N(\mathcal H_N(\mathfrak a),d)\) with \(E=\emptyset\), the marked ideal of Lemma~\ref{hir:step-one-b-independence}: its transform along every tower is
\(\overline{\underline{\mathcal J}}_T\) without the strict transforms of the original members, and the multiple test blow-ups of \(\overline{\underline{\mathcal J}}\) are
multiple test blow-ups of it, with the same cosupports.
\end{lemma}
\begin{proof}
\leanok
Only the divisor of a transform depends on \(E\) (Lemma~\ref{hir:drop-members}); the divisor of the transform is the list of the strict transforms of the members
followed by the exceptional divisor (Lemma~\ref{bcalc:nc-comparison-transform}), and \(E_T\) has \(\#E+\operatorname{length}T\ge m\) members
(Lemma~\ref{hir:geo-transform-basic}), so deleting the first \(m\) members commutes with each step. Admissibility asks normal crossings with the members, which pass
to fewer members (Lemma~\ref{bcalc:sub-divisor}). The coefficient and homogenized ideals along \(N\) carry the divisor of \(\underline{\mathcal J}\) unchanged
(Lemma~\ref{mc:homogenized-along}).
\end{proof}
Step 1b of Proposition~\ref{hir:maximal-order-step} does \emph{not} delete the original members: re-basing a segment across steps with empty centres (Definition~\ref{hir:centre-function}(4')) replaces the base by one with \(\#E+\operatorname{length}P\) members, so the number of deleted members would change and the positions of the remaining members would shift. The original members are kept and masked (Definition~\ref{hir:step-one-value-masked}); the lemma serves to compare the multiple test blow-ups and the cosupports of the two marked ideals.

\begin{lemma}[Step 1b does not depend on the hypersurface of maximal contact]
\label{hir:step-one-b-independence}
\uses{hir:centre-function,hir:tower-isomorphism,mc:glueing-lift-tail,mc:contact-independence-along,mc:homogenized-along,mc:restrict-marked,mc:strict-along,mc:glueing-biholomorphism-along}
\source{wlodarczyk05-hironaka}
\lean{TranscendentalNakamaye.Pending.MaximalOrderStepOneB.stepOneMarked,TranscendentalNakamaye.Pending.MaximalOrderStepOneB.contactMarked,TranscendentalNakamaye.Pending.MaximalOrderStepOneB.inv_contactMarked_eq_ext}
\leanok
Let \(\operatorname{inv}\) be a resolution invariant in dimension \(\le k\), \(\mathfrak a\) coherent, \(N\) a closed submanifold, \(\mu\ge1\),
\(\overline{\underline{\mathcal J}}=\mathcal C_N(\mathcal H_N(\mathfrak a),\mu)\), and \(u\), \(v\) tangent directions of \((\mathfrak a|_N,\mu)\) at a point \(x\) of the cosupport.
For a tower \(T_0\) and an open \(W\) of its last stage write \(\underline{\mathcal S}_{u,T_0}|_W\) for the restriction to \(W\) of the coefficient marked ideal of
\(\overline{\underline{\mathcal J}}_{T_0}\) on the strict transform along \(T_0\) of \(N\cap V(u)\). There are open neighbourhoods \(U'\), \(U\) of \(x\) (the domain and image
of a glueing biholomorphism) such that for every segment \(T_0T_1\) which is a multiple test blow-up of \(\overline{\underline{\mathcal J}}\) and every point \(p\) of
\(\operatorname{cosupp}\overline{\underline{\mathcal J}}_{T_0T_1}\) over \(U'\): if \(\underline{\mathcal S}_{u,T_0}|_{T_0^{-1}U}\) lies in the domain of \(\operatorname{inv}\),
\(T_1|_{T_0^{-1}U}\) is one of its canonical segments and \(p\) lies in the cosupport along it, then
\[\operatorname{inv}_{\underline{\mathcal S}_{v,T_0}|_{T_0^{-1}U'}}(T_1|_{T_0^{-1}U'})(p)=\operatorname{inv}_{\underline{\mathcal S}_{u,T_0}|_{T_0^{-1}U}}(T_1|_{T_0^{-1}U})(p).\]
\end{lemma}
\begin{proof}
\leanok
By Lemma~\ref{mc:glueing-lift-tail} (with the glueing data of Lemma~\ref{mc:glueing-biholomorphism-along}) there are a lift \(\Phi_0\) of the glueing
biholomorphism along \(T_0\) with \(\Phi_0^*(\underline{\mathcal S}_{u,T_0}|_{T_0^{-1}U})=\underline{\mathcal S}_{v,T_0}|_{T_0^{-1}U'}\) and an isomorphism
\(\chi\colon T_1|_{T_0^{-1}U'}\to T_1|_{T_0^{-1}U}\) of towers over \(\Phi_0\) (Definition~\ref{hir:tower-isomorphism}) fixing
\(\operatorname{cosupp}\overline{\underline{\mathcal J}}_{T_0T_1}\) pointwise. Definition~\ref{hir:centre-function}(6) gives
\(\operatorname{inv}_{\Phi_0^*\underline{\mathcal S}_u}(T_1|_{T_0^{-1}U'})(p)=\operatorname{inv}_{\underline{\mathcal S}_u}(T_1|_{T_0^{-1}U})(\chi(p))\), and \(\chi(p)=p\).
\end{proof}

\begin{definition}[The value function of Step 1, first form (superseded in Step 1b)]
\label{hir:step-one-value}
\uses{hir:centre-function,hir:regime-run,hir:drop-original-members,hir:tower-split,bcalc:boundary-count,bcalc:member-intersection-tower,bcalc:member-locus-def,bcalc:codim-one-part-along-def,bcalc:member-restriction-domain,mc:tangent-direction-along,hir:step-one-b-independence,mc:homogenized-along}
\source{wlodarczyk05-hironaka}
\lean{TranscendentalNakamaye.Pending.StepOneValue.barMarked,TranscendentalNakamaye.Pending.StepOneValue.origMembers,TranscendentalNakamaye.Pending.StepOneValue.level,TranscendentalNakamaye.Pending.StepOneValue.IsSpecial,TranscendentalNakamaye.Pending.StepOneValue.stepRegime,TranscendentalNakamaye.Pending.StepOneValue.maskA,TranscendentalNakamaye.Pending.StepOneValue.tangentChoice,TranscendentalNakamaye.Pending.StepOneValue.maskB,TranscendentalNakamaye.Pending.StepOneValue.lowerMask,TranscendentalNakamaye.Pending.StepOneValue.lowerMarked,TranscendentalNakamaye.Pending.StepOneValue.lowerValue,TranscendentalNakamaye.Pending.StepOneValue.stepOneInv,TranscendentalNakamaye.Pending.StepOneValue.LowerData}
\leanok
Let \(\operatorname{inv}\) be a resolution invariant in dimension \(\le k\) with values in \(\Lambda\) (Definition~\ref{hir:centre-function}),
\(\underline{\mathcal J}=(X,N,E,\mathfrak a,d)\) a marked ideal and \(\overline{\underline{\mathcal J}}=\mathcal C_N(\mathcal H_N(\underline{\mathcal J}))\), with the divisor \(E\).
Let \(T\) be a tower over \(X\) and \(y\) a point of its last stage.
\begin{enumerate}
\item The \emph{original members} at the last stage are the strict transforms along \(T\) of the members of \(E\) (the first \(\#E\) members of the divisor of
\(\overline{\underline{\mathcal J}}_T\)); the \emph{level} \(s(y)\) is the number of them through \(y\).
\item The point \(y\) is \emph{special} if \(s(y)=0\) and \(y\) lies in the codimension-one part of \(\operatorname{cosupp}\overline{\underline{\mathcal J}}_T\) along
\(N_T\) (Step 1ba), or \(s(y)\ge1\) and \(y\) lies in the member locus of \(\operatorname{cosupp}\overline{\underline{\mathcal J}}_T\) along \(N_T\) and the original
members (near \(y\), \(N_T\) meets the original members through \(y\) inside the cosupport; Step 1aa). The \emph{regime} of level \(s\) is: level \(s\) and not
special. Let \(T=T_0T_1\) be the split of \(T\) at the start of the run of \(y\) in the regime of its level (Lemma~\ref{hir:regime-run}), \(Y_0\) the last stage
of \(T_0\) and \(y_0\) the image of \(y\) in \(Y_0\).
\item The \emph{mask} and the \emph{lower marked ideal} at level \(s\): for \(s\ge1\), the interior \(W\) of the set of points of \(Y_0\) on no original member missing
\(y_0\) and not special of the level of \(y_0\), and the restriction to \(W\) of \((Y_0,N_{T_0}\cap H_{S(y_0)},E_{S(y_0)},\mathcal C_{N_{T_0}}(\overline{\underline{\mathcal J}}_{T_0}),
\mu_{\mathcal C})\), \(S(y_0)\) the original members through \(y_0\) (Lemma~\ref{bcalc:member-restriction-domain}); for \(s=0\), with a tangent direction \(w\) of
\((\mathfrak a|_N,d)\) on an open \(U\ni p_T(y)\) chosen for the base point \(p_T(y)\) (if there is one), the interior \(W\) of the set of points of \(Y_0\) over \(U\), on
no original member and off the codimension-one part, and the restriction to \(W\) of
\(\underline{\mathcal S}_{w,T_0}=(Y_0,(N\cap V(w))_{T_0},E_{T_0}\setminus E,\mathcal C(\overline{\underline{\mathcal J}}_{T_0}),\mu_{\mathcal C})\), built from
\(\mathcal C_N(\mathcal H_N(\mathfrak a),d)\) with the original members deleted (Lemmas~\ref{hir:drop-original-members} and~\ref{hir:step-one-b-independence}).
\item The \emph{lower value} is \(\operatorname{inv}\) of the lower marked ideal along the restricted tail \(T_1|_W\) at \(y\) (\(\infty\) if \(y_0\notin W\)), and the
\emph{value} is \(\overline{\operatorname{inv}}(y)=(s(y),\infty)\) if \(y\) is special and \((s(y),\text{lower value})\) otherwise, in \(\mathbf N\times(\Lambda\cup\{\infty\})\)
ordered lexicographically.
\item The point has \emph{lower data} if \(y_0\in W\), the lower marked ideal lies in the domain of \(\operatorname{inv}\), \(T_1|_W\) is one of its canonical segments and
\(y\) lies in its cosupport along \(T_1|_W\).
\end{enumerate}
In Step 1b this first form deletes the original members; it is superseded there by the masked form, Definition~\ref{hir:step-one-value-masked}, and kept.
\end{definition}

\begin{lemma}[The value of Step 1 after a blowup]
\label{hir:step-one-value-off-centre}
\uses{hir:step-one-value,hir:centre-function,hir:regime-run,hir:tower-split,bcalc:boundary-count,bcalc:boundary-count-canonical,bcalc:codim-one-part-off-centre,bcalc:member-locus-canonical,bcalc:member-intersection-tower,mc:homogenized-along-basic}
\source{wlodarczyk05-hironaka}
\lean{TranscendentalNakamaye.Pending.StepOneValue.inv_step_of_notMem,TranscendentalNakamaye.Pending.StepOneValue.inv_step_lt_of_mem,TranscendentalNakamaye.Pending.StepOneValue.stepOneInv_of_isSpecial,TranscendentalNakamaye.Pending.StepOneValue.stepOneInv_of_not_isSpecial,TranscendentalNakamaye.Pending.StepOneValue.stepOneInv_fst,TranscendentalNakamaye.Pending.StepOneValue.stepOneInv_lt_of_level_lt,TranscendentalNakamaye.Pending.StepOneValue.stepOneInv_le_of_isSpecial,TranscendentalNakamaye.Pending.StepOneValue.lowerValue_eq_of_lowerData,TranscendentalNakamaye.Pending.StepOneValue.lowerValue_step_of_notMem,TranscendentalNakamaye.Pending.StepOneValue.lowerValue_step_lt_of_mem,TranscendentalNakamaye.Pending.StepOneValue.barMarked_E,TranscendentalNakamaye.Pending.StepOneValue.barMarked_N,TranscendentalNakamaye.Pending.StepOneValue.isCoherent_barMarked,TranscendentalNakamaye.Pending.StepOneValue.origMembers_eq,TranscendentalNakamaye.Pending.StepOneValue.level_step_of_notMem,TranscendentalNakamaye.Pending.StepOneValue.level_step_le,TranscendentalNakamaye.Pending.StepOneValue.isSpecial_step_iff_of_notMem,TranscendentalNakamaye.Pending.StepOneValue.stepRegime_step_iff_of_notMem,TranscendentalNakamaye.Pending.StepOneValue.stepOneInv_step_eq_of_regime,TranscendentalNakamaye.Pending.StepOneValue.stepOneInv_step_of_notMem,TranscendentalNakamaye.Pending.StepOneValue.stepOneInv_step_lt_of_level_lt,TranscendentalNakamaye.Pending.StepOneValue.stepOneInv_step_lt_of_mem}
\leanok
In the situation of Definition~\ref{hir:step-one-value}, let \(\mathfrak a\) be coherent, \(N\) a closed submanifold, the members of \(E\) closed, \(T\) a multiple test blow-up
of \(\overline{\underline{\mathcal J}}\) and \(\sigma\) a blowup along \(C\) of its last stage; write \(T\sigma\) for \(T\) followed by \(\sigma\).
\begin{enumerate}
\item For a canonical segment \(T'\sigma\) of a marked ideal of the domain of \(\operatorname{inv}\) and \(z\) in the cosupport along \(T'\sigma\),
\(\operatorname{inv}(T'\sigma)(z)=\operatorname{inv}(T')(\sigma(z))\) if \(\sigma(z)\notin C\), and \(\operatorname{inv}(T'\sigma)(z)<\operatorname{inv}(T')(\sigma(z))\) if \(\sigma(z)\in C\).
\item The level does not increase: \(s_{T\sigma}(z)\le s_T(\sigma(z))\); if \(\sigma(z)\notin C\), the level, the special points and the regimes at \(z\) and \(\sigma(z)\) agree.
\item The first component of \(\overline{\operatorname{inv}}\) is the level; a point of smaller level has a smaller value, and a special point has the largest value among
the points of its level.
\item If \(\sigma(z)\notin C\) and, when \(\sigma(z)\) is not special, \(z\) has lower data for the split of \(T\) at the start of the run of \(\sigma(z)\) followed by \(\sigma\),
then \(\overline{\operatorname{inv}}(T\sigma)(z)=\overline{\operatorname{inv}}(T)(\sigma(z))\) (Definition~\ref{hir:centre-function}(3), equal off the centre, and (4)).
\item If \(\sigma(z)\in C\), then \(\overline{\operatorname{inv}}(T\sigma)(z)<\overline{\operatorname{inv}}(T)(\sigma(z))\) if \(s_{T\sigma}(z)<s_T(\sigma(z))\), and also if the levels agree,
neither \(z\) nor \(\sigma(z)\) is special and \(z\) has lower data as in (4).
\end{enumerate}
\end{lemma}
\begin{proof}
\leanok
(1) \(T'\sigma\) canonical means \(T'\) canonical and \(C\) empty or the maximum locus; Definition~\ref{hir:centre-function}(3) and (4) (an empty maximum locus is an empty
centre). (2) The strict transform of a closed member lies in its preimage, and off the centre it agrees with the preimage near the point
(Lemmas~\ref{bcalc:boundary-count} and~\ref{bcalc:boundary-count-canonical}); the codimension-one part and the member locus are local and correspond off the centre along
multiple test blow-ups (Lemmas~\ref{bcalc:codim-one-part-off-centre} and~\ref{bcalc:member-locus-canonical}); the coefficient ideal of the homogenized ideal is coherent
(Lemma~\ref{mc:homogenized-along-basic}). (3) Lexicographic order, \(\infty\) the largest element. (4) By (2), if \(\sigma(z)\) is special so is \(z\), with the same level and
value \((s,\infty)\). Otherwise \(\sigma(z)\) is in the regime of \(z\), so the run of \(z\) is that of \(\sigma(z)\) plus one (Lemma~\ref{hir:regime-run}): the split of \(T\sigma\)
has the same head as that of \(T\), with the tail followed by \(\sigma\) (Lemma~\ref{hir:tower-split}), the same image \(y_0\), mask and lower marked ideal; the restricted
tail is followed by the restriction of \(\sigma\), whose centre is the trace of \(C\), not containing the image of \(z\), and (1) applies to the canonical segment given by
the lower data. (5) The first case is (3); in the second, as in (4) the lower value at \(z\) is computed along the restricted tail followed by the restriction of
\(\sigma\), whose centre contains the image of \(z\), and (1) gives the strict decrease.
\end{proof}

\begin{lemma}[The masks and the domain of the lower marked ideals of Step 1]
\label{hir:step-one-lower-domain}
\uses{hir:step-one-value,hir:wellformed-tower,hir:along-transform,mc:homogenized-along-maximal-order,mc:homogenized-along-basic,bcalc:transform-off-centre,bcalc:codim-one-local-along,bcalc:codim-one-part-along,bcalc:member-locus-centre,bcalc:member-restriction-domain,bcalc:inv-domain-stable,mc:contact-stage-domain-masked,hir:drop-original-members}
\source{wlodarczyk05-hironaka}
\lean{TranscendentalNakamaye.Pending.StepOneLowerDomain.isWellFormed_barMarked,TranscendentalNakamaye.Pending.StepOneLowerDomain.isMaximalOrderAlong_barMarked,TranscendentalNakamaye.Pending.StepOneLowerDomain.cosupp_barMarked,TranscendentalNakamaye.Pending.StepOneLowerDomain.isClosed_of_mem_origMembers,TranscendentalNakamaye.Pending.StepOneLowerDomain.isMaximalOrderAlong_geoTransformAlong,TranscendentalNakamaye.Pending.StepOneLowerDomain.proj_mem_cosupp_of_mem_cosupp,TranscendentalNakamaye.Pending.StepOneLowerDomain.tangentChoice_spec,TranscendentalNakamaye.Pending.StepOneLowerDomain.exists_tangent_of_mem_cosupp,TranscendentalNakamaye.Pending.StepOneLowerDomain.mem_maskB,TranscendentalNakamaye.Pending.StepOneLowerDomain.isInvDomain_lowerMarked_zero,TranscendentalNakamaye.Pending.StepOneLowerDomain.mem_maskA,TranscendentalNakamaye.Pending.StepOneLowerDomain.isInvDomain_lowerMarked_succ,TranscendentalNakamaye.Pending.StepOneLowerDomain.mem_lowerMask_and_isInvDomain}
\leanok
In the situation of Definition~\ref{hir:step-one-value}, let \(\underline{\mathcal J}\) be well formed with \(d\ge1\), of maximal order along \(N\), \(\dim N\le k\),
\(n\ge1\), and \(T_0\) a multiple test blow-up of \(\overline{\underline{\mathcal J}}\). Then \(\overline{\underline{\mathcal J}}\) is well formed and of maximal order along \(N\)
with the cosupport of \(\underline{\mathcal J}\), maximal order is kept along \(T_0\), the images of points of the cosupport along \(T_0\) lie in the cosupport, and the
original members are closed. Let \(y_0\) be a point of \(\operatorname{cosupp}\overline{\underline{\mathcal J}}_{T_0}\) which is not special, of level \(s\). Then \(y_0\) lies
in its mask, and the lower marked ideal at level \(s\) lies in the domain of resolution invariants in dimension \(\le k-1\) (Definition~\ref{hir:centre-function}).
\end{lemma}
\begin{proof}
\leanok
\(\overline{\underline{\mathcal J}}\) has the submanifold and divisor of \(\underline{\mathcal J}\) and a coherent ideal (Lemma~\ref{mc:homogenized-along-basic}), the
cosupport of \(\underline{\mathcal J}\), and maximal order (Lemma~\ref{mc:homogenized-along-maximal-order}); maximal order passes to transforms by admissible blowups
(Lemma~\ref{hir:along-transform}) and so do well formedness and \(\dim N\le k\) (Lemmas~\ref{hir:wellformed-tower} and~\ref{bcalc:inv-domain-stable}). A point of the
cosupport over the centre of a blowup maps into the centre, which lies in the cosupport; off it, Lemma~\ref{bcalc:transform-off-centre}.

\(s=0\). The base point \(p_{T_0}(y_0)\) lies in \(\operatorname{cosupp}\underline{\mathcal J}\), so a tangent direction \(w\) on an open \(U\ni p_{T_0}(y_0)\) exists
(Lemma~\ref{bcalc:codim-one-local-along}). The set of points over \(U\), on no original member and off the codimension-one part is open (the members are closed, the
codimension-one part is closed by maximal order, Theorem~\ref{bcalc:codim-one-part-along}) and contains \(y_0\) (level \(0\), not special); so \(y_0\) lies in the mask
\(W\subset p_{T_0}^{-1}U\). The lower marked ideal is \(\underline{\mathcal S}_{w,T_0}|_W\), formed from \(\mathcal C_N(\mathcal H_N(\mathfrak a),d)\), along whose
multiple test blow-ups \(T_0\) is one with the same cosupports and codimension-one parts (Lemma~\ref{hir:drop-original-members}); \(W\) misses the codimension-one part, and
Lemma~\ref{mc:contact-stage-domain-masked} applies.

\(s\ge1\). On the open set \(V\) of points on no original member missing \(y_0\) (Lemma~\ref{bcalc:member-restriction-domain}) the level is at most \(s\), so for the
restriction of \(\overline{\underline{\mathcal J}}_{T_0}\) to \(V\) the points of the member locus of level \(s\) form an admissible centre
(Lemma~\ref{bcalc:member-locus-centre}), in particular a closed subset of \(V\), not containing \(y_0\); its complement in \(V\) is an open neighbourhood of \(y_0\) inside
the defining set of the mask. At a point \(z\) of the mask on \(N_{T_0}\cap H_{S(y_0)}\), \(z\) lies on exactly the original members through \(y_0\), so it has level \(s\)
and is not special: it is off the member locus, and the restricted coefficient ideal is nonzero along \(N_{T_0}\cap H_{S(y_0)}\) near \(z\); with well formedness and
\(\dim(N_{T_0}\cap H_{S(y_0)})\le k-1\) (Lemma~\ref{bcalc:member-restriction-domain}) the restriction to the mask lies in the domain in dimension \(\le k-1\).
\end{proof}

\begin{lemma}[Points over a mask keep the regime]
\label{hir:step-one-mask-tower}
\uses{hir:step-one-value,hir:step-one-value-off-centre,hir:step-one-lower-domain,hir:regime-run,bcalc:member-locus-level,bcalc:codim-one-part-tower,bcalc:member-intersection-tower}
\source{wlodarczyk05-hironaka}
\lean{TranscendentalNakamaye.Pending.StepOneMaskTower.le_regimeRun_append,TranscendentalNakamaye.Pending.StepOneMaskTower.origMembers_append,TranscendentalNakamaye.Pending.StepOneMaskTower.level_append_le,TranscendentalNakamaye.Pending.StepOneMaskTower.level_le_of_forall_mem_imp,TranscendentalNakamaye.Pending.StepOneMaskTower.level_le_of_mem_maskA,TranscendentalNakamaye.Pending.StepOneMaskTower.level_eq_zero_of_mem_maskB,TranscendentalNakamaye.Pending.StepOneMaskTower.not_isSpecial_of_mem_maskA,TranscendentalNakamaye.Pending.StepOneMaskTower.not_isSpecial_of_mem_maskB,TranscendentalNakamaye.Pending.StepOneMaskTower.stepRegime_and_le_regimeRun_of_mem_lowerMask}
\leanok
In the situation of Lemma~\ref{hir:step-one-lower-domain}, let \(T_0T_1\) be a multiple test blow-up of \(\overline{\underline{\mathcal J}}\), \(y_0\) a point of the last
stage of \(T_0\) of level \(s\), \(W\) its mask at level \(s\) (Definition~\ref{hir:step-one-value}), and \(z\) a point of the last stage over \(W\). Then the level of \(z\) is
at most \(s\), and \(0\) if \(s=0\). If the level of \(z\) is \(s\), then \(z\) is not special; more precisely \(z\) is in the regime of level \(s\) and its run in it is at least
the length of \(T_1\): \(z\) entered the regime at the last stage of \(T_0\) or before. In general, if a regime \(G\) of the towers over the last stage of \(T_0\) is inherited
by images and implies a regime \(P\) after \(T_0\), the run in \(P\) of a point having \(G\) after \(T_1\) is at least the length of \(T_1\).
\end{lemma}
\begin{proof}
\leanok
The original members along \(T_0T_1\) are the strict transforms along \(T_1\) of those at \(T_0\), and the level does not increase along \(T_1\)
(Lemma~\ref{hir:step-one-value-off-centre}(2)); a point of \(W\) lies on no original member missing \(y_0\) (\(s\ge1\)), resp. on none (\(s=0\)). If \(s\ge1\) and \(z\) of
level \(s\) were special, it would lie over the closure of the special points of level \(s\) at \(T_0\) (Lemma~\ref{bcalc:member-locus-level}); but \(W\) is an open
neighbourhood of the image of \(z\) missing them. If \(s=0\), the codimension-one part along \(T_1\) lies over that at \(T_0\) (Lemma~\ref{bcalc:codim-one-part-tower},
maximal order kept, Lemma~\ref{hir:step-one-lower-domain}), which \(W\) misses. Both statements hold at every stage of \(T_1\) for the images of \(z\) (their level is at
least that of \(z\), at most \(s\)), so the regime holds along the whole tail and the run covers it (induction along \(T_1\), Lemma~\ref{hir:regime-run}).
\end{proof}

\begin{lemma}[Comparing nearby points in Step 1]
\label{hir:step-one-nearby}
\uses{hir:step-one-value,hir:step-one-value-off-centre,hir:step-one-lower-domain,hir:step-one-mask-tower,hir:regime-run,hir:tower-split,hir:centre-function,hir:canonical-segments,bcalc:restrict-nested,bcalc:restrict-empty-tower,bcalc:member-restriction-domain,mc:tower-isomorphism-transform,hir:open-restriction-tower}
\source{wlodarczyk05-hironaka}
\lean{TranscendentalNakamaye.Pending.StepOneNearby.exists_split_le,TranscendentalNakamaye.Pending.StepOneNearby.mem_cosupp_restrict_nested_iff,TranscendentalNakamaye.Pending.StepOneNearby.inv_eq_of_two_opens,TranscendentalNakamaye.Pending.StepOneNearby.mem_origMembers_iff_of_mem_maskA,TranscendentalNakamaye.Pending.StepOneNearby.membersThrough_eq_of_mem_maskA,TranscendentalNakamaye.Pending.StepOneNearby.HasLowerData,TranscendentalNakamaye.Pending.StepOneNearby.stepOneInv_lt_top_of_hasLowerData,TranscendentalNakamaye.Pending.StepOneNearby.toLex_top_le_stepOneInv,TranscendentalNakamaye.Pending.StepOneNearby.stepOneInv_lt_of_mem_lowerMask,TranscendentalNakamaye.Pending.StepOneNearby.not_stepRegime_proj,TranscendentalNakamaye.Pending.StepOneNearby.disjoint_lowerMask_of_isCanonical,TranscendentalNakamaye.Pending.StepOneNearby.centresAvoid_lowerMask}
\leanok
In the situation of Lemma~\ref{hir:step-one-lower-domain}, say that a tower \(T\) \emph{has lower data} if every point of
\(\operatorname{cosupp}\overline{\underline{\mathcal J}}_T\) which is not special has lower data (Definition~\ref{hir:step-one-value}(5)).
\begin{enumerate}
\item For \(a\le b\), the split \(T=T_0^bT_1^b\) before the last \(b\) blowups and the split \(T=T_0^aT_1^a\) before the last \(a\) blowups are related by a tower
\(M\) with \(T_0^bM=T_0^a\) and \(MT_1^a=T_1^b\).
\item Let \(\operatorname{inv}\) be a resolution invariant, \(\underline{\mathcal I}\) a marked ideal with a coherent ideal and \(N\) a closed submanifold, \(T\) a tower over its
base and \(A\), \(B\) open sets of the base such that \(\underline{\mathcal I}|_A\), \(\underline{\mathcal I}|_B\) lie in the domain and \(T|_A\), \(T|_B\) are canonical segments
of them. At a point over \(A\cap B\) of \(\operatorname{cosupp}(\underline{\mathcal I}|_A)_{T|_A}\) the values of \(\underline{\mathcal I}|_A\) along \(T|_A\) and of
\(\underline{\mathcal I}|_B\) along \(T|_B\) agree, and the point lies in \(\operatorname{cosupp}(\underline{\mathcal I}|_B)_{T|_B}\).
\item A point of the mask of \(y_0\) at level \(s\ge1\) (at the same stage) of level \(s\) lies on exactly the original members through \(y_0\).
\item If \(T\) has lower data, the value at a non-special point \(y\) of the cosupport is \((s(y),\lambda)\) with \(\lambda\ne\infty\); a point \(w\) of level at least \(s\)
which is not in the regime of level \(s\) has value at least \((s,\infty)\).
\item Let \(T_0M\) be a multiple test blow-up of \(\overline{\underline{\mathcal J}}\) with lower data, \(p_1\) a point of the last stage of \(T_0\) of level \(s\) and \(w\) a point
of the last stage of \(T_0M\) of level at least \(s\), not in the regime of level \(s\). Every point of the cosupport over the mask of \(p_1\) has a smaller value than \(w\).
\item If, after a further blowup \(\sigma\), \(w\) is a point of level at least \(s\) of the cosupport not in the regime of level \(s\), then so is \(\sigma(w)\).
\item Let \(T_0M\sigma\) be a multiple test blow-up of \(\overline{\underline{\mathcal J}}\) which is a canonical segment of \(\overline{\operatorname{inv}}\), such that every
prefix \(T_0M'\) with \(M'\) at most as long as \(M\) which is a multiple test blow-up and a canonical segment has lower data. Let \(p_1\) be a point of the last stage of
\(T_0\) of level \(s\) and \(w\) a point of the cosupport at the last stage of \(T_0M\) of level at least \(s\), not in the regime of level \(s\). Then the centres of
\(M\sigma\) miss the points over the mask of \(p_1\).
\end{enumerate}
\end{lemma}
\begin{proof}
\leanok
(1) Induction along \(T\): splitting a blowup off the end of \(T\) shifts both indices by one, and for \(a=0<b\) the middle part is the tail at \(b\).
(2) Apply Lemma~\ref{bcalc:restrict-nested}(2) to \(A\cap B\subset A\) and \(A\cap B\subset B\); the cosupports along \(T|_{A\cap B}\) and \(T|_A\) correspond under the
isomorphism of Lemma~\ref{bcalc:restrict-nested}(1) (Lemma~\ref{mc:tower-isomorphism-transform}) and the trace of the cosupport along a restricted multiple test blow-up
is the cosupport of the restriction (Lemma~\ref{hir:canonical-segments}), and canonical segments are multiple test blow-ups (Theorem~\ref{hir:invariant-driver}).
(3) A point \(w\) of the mask lies on no original member missing \(y_0\); the original members through \(w\) are among those through \(y_0\), and there are as many,
so they are the same.
(4) Lower data give the lower value as a value of \(\operatorname{inv}\), an element of \(\Lambda\). If the level of \(w\) is larger than \(s\), the first component is
larger; if it is \(s\), \(w\) is special, with value \((s,\infty)\).
(5) Over the mask the level is at most \(s\) (Lemma~\ref{hir:step-one-mask-tower}). A point of smaller level has a smaller value; a point of level \(s\) is in
the regime (Lemma~\ref{hir:step-one-mask-tower}), so not special, and by lower data and (4) its value is below \((s,\infty)\), which is at most the value of \(w\) by (4).
(6) The image of a cosupport point lies in the cosupport (Lemma~\ref{hir:step-one-lower-domain}) and the level does not increase
(Lemma~\ref{hir:step-one-value-off-centre}(2)). If \(\sigma(w)\) were in the regime of level \(s\), it would lie in its own mask
(Lemma~\ref{hir:step-one-lower-domain}), and \(w\), over that mask, of level between \(s\) and the level \(s\) of \(\sigma(w)\), would be in the regime as well
(Lemma~\ref{hir:step-one-mask-tower}).
(7) Induction along \(M\). The centre of \(\sigma\) is empty or the maximum locus of \(\overline{\operatorname{inv}}\) on the cosupport at the last stage of \(T_0M\); a point
of it over the mask of \(p_1\) would have value at least that of \(w\), against (5). For the earlier centres apply the induction hypothesis to the image of \(w\), which
by (6) satisfies the same hypotheses. In the terms of Lemma~\ref{bcalc:restrict-empty-tower}, the steps of \(M\sigma\) restricted to the mask of \(p_1\) have only empty
centres, so a value computed from the stage of \(T_0\) can be re-based to the stage after \(M\sigma\) (Definition~\ref{hir:centre-function}(4')).
\end{proof}

\begin{definition}[The value function of Step 1, original members masked]
\label{hir:step-one-value-masked}
\uses{hir:step-one-value,hir:centre-function,hir:regime-run,hir:tower-split,mc:contact-stage-domain-masked,mc:homogenized-along,mc:strict-along,mc:tangent-direction-along}
\source{wlodarczyk05-hironaka}
\lean{TranscendentalNakamaye.Pending.StepOneValueMasked.lowerMarkedB,TranscendentalNakamaye.Pending.StepOneValueMasked.lowerMarkedM,TranscendentalNakamaye.Pending.StepOneValueMasked.LowerFamily,TranscendentalNakamaye.Pending.StepOneValueMasked.lowerValueL,TranscendentalNakamaye.Pending.StepOneValueMasked.stepOneInvL,TranscendentalNakamaye.Pending.StepOneValueMasked.LowerDataL,TranscendentalNakamaye.Pending.StepOneValueMasked.stepOneInvM,TranscendentalNakamaye.Pending.StepOneValueMasked.lowerValue_eq_lowerValueL,TranscendentalNakamaye.Pending.StepOneValueMasked.stepOneInv_eq_stepOneInvL,TranscendentalNakamaye.Pending.StepOneValueMasked.lowerData_iff_lowerDataL,TranscendentalNakamaye.Pending.StepOneValueMasked.lowerMarkedM_succ,TranscendentalNakamaye.Pending.StepOneValueMasked.lowerMarkedM_zero}
\leanok
In the situation of Definition~\ref{hir:step-one-value}, a \emph{family of lower marked ideals} assigns to every level \(s\), tower \(T_0\) and point \(y_0\) of its last
stage a marked ideal \(L_s(T_0,y_0)\) on the mask of \(y_0\) at level \(s\). For such a family, the \emph{lower value}, the \emph{value} \(\overline{\operatorname{inv}}_L\)
and \emph{lower data} are defined as in Definition~\ref{hir:step-one-value}(4)--(5) with \(L_s(T_0,y_0)\) in place of the lower marked ideal there; Definition~\ref{hir:step-one-value}
is the case of its lower marked ideals.
The \emph{masked} family is: for \(s\ge1\) the lower marked ideal of Definition~\ref{hir:step-one-value}(3); for \(s=0\), with \(w\) the tangent direction chosen at the base
point and \(W\) the mask, the restriction to \(W\) of
\[\underline{\mathcal S}^{E}_{w,T_0}=\bigl(Y_0,(N\cap V(w))_{T_0},E_{T_0},\mathcal C_{N_{T_0}}(\overline{\underline{\mathcal J}}_{T_0}),\mu_{\mathcal C}\bigr),\]
formed from \(\overline{\underline{\mathcal J}}=\mathcal C_N(\mathcal H_N(\underline{\mathcal J}))\) with the divisor \(E\) of \(\underline{\mathcal J}\) \emph{kept}: the original members
are masked (the mask misses their strict transforms, so their traces on \(W\) are empty) instead of deleted. The \emph{masked value function} is
\(\overline{\operatorname{inv}}_L\) for the masked family.
\end{definition}
This is Włodarczyk's \((E_p\setminus E)_{|V(u)_p}\), which removes the original members as a set: the divisor keeps the positions of all members, so that
re-basing across steps with empty centres (Definition~\ref{hir:centre-function}(4')), which replaces \(\underline{\mathcal J}\) by \(\underline{\mathcal J}_P\) with
\(\#E+\operatorname{length}P\) members, does not shift the positions of the members of the lower marked ideal; a resolution invariant may depend on these positions (the
monomial invariant of Lemma~\ref{hir:monomial-invariant} does). Deleting the \(\#E\) original members, as in Definition~\ref{hir:step-one-value}, would shift them.

\begin{lemma}[The masked value function: domain, finiteness, one blowup, nearby points]
\label{hir:step-one-value-masked-basic}
\uses{hir:step-one-value-masked,hir:step-one-value-off-centre,hir:step-one-lower-domain,hir:step-one-mask-tower,hir:step-one-nearby,mc:contact-stage-domain-masked,hir:regime-run,hir:tower-split,bcalc:restrict-empty-tower}
\source{wlodarczyk05-hironaka}
\lean{TranscendentalNakamaye.Pending.StepOneValueMasked.mem_lowerMask_and_isInvDomain_M,TranscendentalNakamaye.Pending.StepOneValueMasked.stepRegime_and_mem_cosupp_head,TranscendentalNakamaye.Pending.StepOneValueMasked.mem_lowerMask_of_stepRegime,TranscendentalNakamaye.Pending.StepOneValueMasked.lowerValueL_ne_top,TranscendentalNakamaye.Pending.StepOneValueMasked.stepOneInvL_of_isSpecial,TranscendentalNakamaye.Pending.StepOneValueMasked.stepOneInvL_of_not_isSpecial,TranscendentalNakamaye.Pending.StepOneValueMasked.stepOneInvL_fst,TranscendentalNakamaye.Pending.StepOneValueMasked.stepOneInvL_lt_of_level_lt,TranscendentalNakamaye.Pending.StepOneValueMasked.stepOneInvL_le_of_isSpecial,TranscendentalNakamaye.Pending.StepOneValueMasked.lowerValueL_eq_of_mem,TranscendentalNakamaye.Pending.StepOneValueMasked.lowerValueL_step_of_notMem,TranscendentalNakamaye.Pending.StepOneValueMasked.lowerValueL_step_lt_of_mem,TranscendentalNakamaye.Pending.StepOneValueMasked.stepOneInvL_step_eq_of_regime,TranscendentalNakamaye.Pending.StepOneValueMasked.stepOneInvL_step_of_notMem,TranscendentalNakamaye.Pending.StepOneValueMasked.stepOneInvL_step_lt_of_level_lt,TranscendentalNakamaye.Pending.StepOneValueMasked.stepOneInvL_step_lt_of_mem,TranscendentalNakamaye.Pending.StepOneValueMasked.HasLowerDataL,TranscendentalNakamaye.Pending.StepOneValueMasked.stepOneInvL_lt_top_of_hasLowerData,TranscendentalNakamaye.Pending.StepOneValueMasked.toLex_top_le_stepOneInvL,TranscendentalNakamaye.Pending.StepOneValueMasked.stepOneInvL_lt_of_mem_lowerMask,TranscendentalNakamaye.Pending.StepOneValueMasked.disjoint_lowerMask_of_isCanonicalL,TranscendentalNakamaye.Pending.StepOneValueMasked.centresAvoid_lowerMaskL}
\leanok
In the situation of Lemma~\ref{hir:step-one-lower-domain}:
\begin{enumerate}
\item at a point \(y_0\) of the cosupport along \(T_0\) in the regime of level \(s\), \(y_0\) lies in its mask and the masked lower marked ideal lies in the domain in
dimension \(\le k-1\);
\item along a multiple test blow-up \(T\) of \(\overline{\underline{\mathcal J}}\), the image of a non-special point \(y\) of the cosupport at the start of its run is in the
regime of the level of \(y\) and in the cosupport, so it lies in its mask, and the lower value at \(y\) is not \(\infty\) (for every family);
\item for every family, Lemma~\ref{hir:step-one-value-off-centre}(3)--(5) and Lemma~\ref{hir:step-one-nearby}(4)--(7) hold with \(\overline{\operatorname{inv}}_L\) and lower
data for \(L\) in place of \(\overline{\operatorname{inv}}\) and lower data.
\end{enumerate}
\end{lemma}
\begin{proof}
\leanok
(1) For \(s\ge1\) this is Lemma~\ref{hir:step-one-lower-domain}. For \(s=0\) the mask lies over the open set of \(w\), misses the codimension-one part and the original
members, and Lemma~\ref{mc:contact-stage-domain-masked} applies to \(\underline{\mathcal S}^E_{w,T_0}\) restricted to it. (2) If the run is positive, the image at its start has
the regime (Lemma~\ref{hir:regime-run}); if it is \(0\), the image is \(y\), of its own level and not special. Images of points of the cosupport lie in the cosupport
(Lemma~\ref{hir:step-one-lower-domain}), and (1) gives the mask; the lower value is then a value of \(\operatorname{inv}\). (3) The proofs of the two lemmas use the lower
marked ideals only through lower data, verbatim.
\end{proof}

\begin{lemma}[The value of a nearby point through the reference point's lower marked ideal]
\label{hir:step-one-compare}
\uses{hir:step-one-value-masked,hir:step-one-value-masked-basic,hir:step-one-nearby,hir:step-one-mask-tower,bcalc:member-restriction-rebase,bcalc:inv-rebase,mc:contact-stage-rebase,mc:contact-stage-tail,hir:coefficient-restriction,bcalc:member-intersection-tower,hir:tower-isomorphism,hir:centre-function,mc:tangent-direction-along}
\source{wlodarczyk05-hironaka}
\lean{TranscendentalNakamaye.Pending.StepOneCompare.inv_restrict_append,TranscendentalNakamaye.Pending.StepOneCompare.memberRestrict_membersThrough_congr,TranscendentalNakamaye.Pending.StepOneCompare.lowerValue_eq_succ,TranscendentalNakamaye.Pending.StepOneCompare.contactKept,TranscendentalNakamaye.Pending.StepOneCompare.lowerMarkedB_eq_contactKept,TranscendentalNakamaye.Pending.StepOneCompare.tangentChoice_spec_of_mem,TranscendentalNakamaye.Pending.StepOneCompare.contactKept_rebase,TranscendentalNakamaye.Pending.StepOneCompare.lowerValue_eq_zero_rebase,TranscendentalNakamaye.Pending.StepOneCompare.mem_cosupp_lowerMarked_succ_iff,TranscendentalNakamaye.Pending.StepOneCompare.level_eq_iff_of_mem_maskA,TranscendentalNakamaye.Pending.StepOneCompare.isGeoTestTower_cosupp_lowerMarkedB}
\leanok
In the situation of Lemma~\ref{hir:step-one-lower-domain}, with the masked lower marked ideals (Definition~\ref{hir:step-one-value-masked}), let
\(T_0'MT_1\) be a multiple test blow-up of \(\overline{\underline{\mathcal J}}\), \(T_0=T_0'M\), \(y_0\) a point of the last stage of \(T_0\) of level \(s\) with mask \(W\)
and lower marked ideal \(\underline{\mathcal S}(y_0)\), and \(p\) a point of the last stage over \(W\) with lower data at the split \(T_0'(MT_1)\), such that the
centres of \(M\) miss its mask \(W'\) at \(T_0'\). Write \(\underline{\mathcal S}^E_{w,T_0}|_V\) for the restriction to an open \(V\) of the marked ideal of
Definition~\ref{hir:step-one-value-masked} formed with a function \(w\).
\begin{enumerate}
\item For a marked ideal on \(W'\) in the domain, \((MT_1)|_{W'}\) is a canonical segment of it iff \(M|_{W'}\) followed by \(T_1|_{M^{-1}W'}\) is, with the same
cosupports and values.
\item (\(s\ge1\).) If \(\underline{\mathcal S}(y_0)\) lies in the domain, \(T_1|_W\) is a canonical segment of it, the image of \(p\) at \(T_0\) has the level of \(y_0\)
and \(p\) lies in the cosupport of \(\underline{\mathcal S}(y_0)\) along \(T_1|_W\), then the lower value of \(p\) is
\(\operatorname{inv}_{\underline{\mathcal S}(y_0)}(T_1|_W)(p)\).
\item (\(s=0\).) \(p\) lies over \(M^{-1}W'\), \(\underline{\mathcal S}^E_{v,T_0}|_{M^{-1}W'}\) lies in the domain with \(T_1|_{M^{-1}W'}\) a canonical segment of it, \(v\) the
tangent direction chosen at the base point of \(p\), and the lower value of \(p\) is \(\operatorname{inv}\) of it along \(T_1|_{M^{-1}W'}\) at \(p\).
\item (Cosupports.) If \(s\ge1\), \(y_0\) lies in its mask and \(T_1|_W\) is a multiple test blow-up of \(\underline{\mathcal S}(y_0)\), a point over \(W\) lies in its cosupport
along \(T_1|_W\) iff it lies in \(\operatorname{cosupp}\overline{\underline{\mathcal J}}_{T_0T_1}\) and on the strict transforms of the original members through \(y_0\); over
\(W\) the latter holds iff the point has the level of \(y_0\). If \(s=0\) and \(y_0\) lies in the cosupport, \(T_1|_W\) is a multiple test blow-up of
\(\underline{\mathcal S}(y_0)\) and its cosupport along it is the trace of \(\operatorname{cosupp}\overline{\underline{\mathcal J}}_{T_0T_1}\).
\end{enumerate}
\end{lemma}
\begin{proof}
\leanok
(1) The two towers are isomorphic over the identity of \(W'\), the identity on points at the last stage (Lemma~\ref{hir:tower-isomorphism} with the restriction of a
concatenation); apply Definition~\ref{hir:centre-function}(6) and Lemma~\ref{hir:canonical-segments}(2). (2) By (1) and Lemma~\ref{bcalc:member-restriction-rebase}
(re-basing along \(M\), empty over \(W'\)), the lower value of \(p\) is the value of the marked ideal formed at \(T_0\) at the image \(p_0\) of \(p\), restricted to
\(M^{-1}W'\), along \(T_1|_{M^{-1}W'}\). Its original members are the strict transforms along \(M\) of those at \(T_0'\) (Lemma~\ref{hir:step-one-mask-tower}); \(p_0\) lies
in the mask of \(y_0\) with the same level, so on the same original members (Lemma~\ref{hir:step-one-nearby}(3)), and the marked ideal is the one formed at \(y_0\). The
values on the two open sets \(M^{-1}W'\) and \(W\) agree at \(p\) (Lemma~\ref{hir:step-one-nearby}(2)). (3) By (1) and Lemma~\ref{mc:contact-stage-rebase} in its form for
any divisor, with Lemma~\ref{bcalc:inv-rebase}. (4) For \(s\ge1\): move the restriction to \(W\) inside the restricted marked ideal; Lemma~\ref{hir:coefficient-restriction}
on \(W\) gives the cosupport \(\operatorname{cosupp}\cap\) the strict transform of \(N_{T_0}\cap H_S\), and the strict transform of \(N_{T_0}\cap H_S\) is \(N_{T_0T_1}\)
intersected with the strict transforms of the members of \(S\) (Lemma~\ref{bcalc:member-intersection-tower}). Over the mask a point lies on no original member missing
\(y_0\); strict transforms of closed members lie over them, so the members through the point are among the strict transforms of those through \(y_0\), and the counts agree
iff it lies on all of them. For \(s=0\): Lemma~\ref{mc:contact-stage-tail} in its form for any divisor, the mask lying over the open set of the tangent direction.
\end{proof}

\begin{theorem}[The canonical-tail theorem of Step 1]
\label{hir:step-one-canonical-tail}
\uses{hir:step-one-value-masked,hir:step-one-value-masked-basic,hir:step-one-compare,hir:step-one-nearby,hir:step-one-mask-tower,hir:regime-run,hir:tower-split,bcalc:max-locus-local,mc:tangent-change-opens,hir:centre-function,hir:invariant-driver}
\source{wlodarczyk05-hironaka}
\lean{TranscendentalNakamaye.Pending.StepOneCanonicalTail.centresAvoid_of_length_eq_zero,TranscendentalNakamaye.Pending.StepOneCanonicalTail.exists_eq_step_refl,TranscendentalNakamaye.Pending.StepOneCanonicalTail.mem_cosupp_head,TranscendentalNakamaye.Pending.StepOneCanonicalTail.level_le_head,TranscendentalNakamaye.Pending.StepOneCanonicalTail.mem_cosupp_lowerMarkedM_iff,TranscendentalNakamaye.Pending.StepOneCanonicalTail.lowerDataL_refl,TranscendentalNakamaye.Pending.StepOneCanonicalTail.lowerValueL_eq_succ_of_eq,TranscendentalNakamaye.Pending.StepOneCanonicalTail.lowerValueL_eq_zero_of_eq,TranscendentalNakamaye.Pending.StepOneCanonicalTail.exists_middle,TranscendentalNakamaye.Pending.StepOneCanonicalTail.level_le_of_mem_lowerMask,TranscendentalNakamaye.Pending.StepOneCanonicalTail.stepOneInvM_eq_lower,TranscendentalNakamaye.Pending.StepOneCanonicalTail.lowerDataL_step,TranscendentalNakamaye.Pending.StepOneCanonicalTail.hasLowerDataL_of_isCanonical,TranscendentalNakamaye.Pending.StepOneCanonicalTail.stepOneInvM_step_of_notMem_of_isCanonical,TranscendentalNakamaye.Pending.StepOneCanonicalTail.stepOneInvM_step_lt_of_mem_of_isCanonical}
\leanok
In the situation of Lemma~\ref{hir:step-one-lower-domain}, let \(\operatorname{inv}\) be a resolution invariant in dimension \(\le k-1\) and \(\overline{\operatorname{inv}}\)
the masked value function of Step 1 (Definition~\ref{hir:step-one-value-masked}). Along every multiple test blow-up \(T\) of \(\overline{\underline{\mathcal J}}\) which is a
canonical segment of \(\overline{\operatorname{inv}}\), every point of \(\operatorname{cosupp}\overline{\underline{\mathcal J}}_T\) which is not special has lower data at the split of
\(T\) at the start of its run: its lower value is the value of \(\operatorname{inv}\) along a canonical segment of its lower marked ideal, at a point of its cosupport.
Consequently, along such segments followed by a blowup \(\sigma\) with centre \(C\), \(\overline{\operatorname{inv}}(T\sigma)(z)=\overline{\operatorname{inv}}(T)(\sigma(z))\) at the
points \(z\) of the cosupport with \(\sigma(z)\notin C\), and \(\overline{\operatorname{inv}}(T\sigma)(z)<\overline{\operatorname{inv}}(T)(\sigma(z))\) if \(\sigma(z)\in C\) and \(z\),
\(\sigma(z)\) are not special and of the same level.
\end{theorem}
\begin{proof}
\leanok
Induction on the length of \(T\); the empty tower and a point which has just entered its regime have lower data along the empty tail
(Lemma~\ref{hir:step-one-value-masked-basic}(1) and Lemma~\ref{hir:step-one-compare}(4)). Let \(T=T'\sigma\) and \(y\) a non-special point of the cosupport whose image
\(z=\sigma(y)\) keeps the regime of the level \(s\) of \(y\); by induction \(z\) has lower data at the split \(T'=T_0T_1\) of its run, with mask \(W\) and lower marked ideal
\(\underline{\mathcal S}\), and the split of \(y\) is \(T_0(T_1\sigma)\). It remains to see that the trace of \(C\) on \(T_1|_W\) is empty or the maximum locus of
\(\operatorname{inv}_{\underline{\mathcal S}}(T_1|_W)\) on its cosupport, and that \(y\) lies in the cosupport along \((T_1\sigma)|_W\), which then follows from
Lemma~\ref{hir:step-one-compare}(4) since \((T_1\sigma)|_W\) is a canonical segment, hence a multiple test blow-up (Theorem~\ref{hir:invariant-driver}).
\(C\) is empty or the maximum locus of \(\overline{\operatorname{inv}}(T')\) on the cosupport. Let \(q\) be a point of the cosupport of \(\underline{\mathcal S}\) along \(T_1|_W\).
By Lemma~\ref{hir:step-one-compare}(4), \(q\) lies in the cosupport of \(\overline{\underline{\mathcal J}}_{T'}\) and has level \(s\); by Lemma~\ref{hir:step-one-mask-tower} it is in
the regime and its run is at least the length of \(T_1\), so its split is \(T_0'(MT_1)\) with \(T_0'M=T_0\) (Lemma~\ref{hir:step-one-nearby}(1)), and by induction \(q\) has
lower data there. If its run is longer than that of \(z\), the image of \(z\) one stage before \(T_0\) is out of the regime (Lemma~\ref{hir:regime-run}), so the centres
of \(M\) miss the mask of \(q\) (Lemma~\ref{hir:step-one-nearby}(7), Lemma~\ref{hir:step-one-value-masked-basic}(3)). Lemma~\ref{hir:step-one-compare}(2), resp.\ (3)
followed by the change of tangent direction on the masks (Lemma~\ref{mc:tangent-change-opens}), gives
\(\overline{\operatorname{inv}}(T')(q)=(s,\operatorname{inv}_{\underline{\mathcal S}}(T_1|_W)(q))\). A point of the cosupport over \(W\) which is not in the cosupport of
\(\underline{\mathcal S}\) has level \(<s\) (Lemma~\ref{hir:step-one-compare}(4) and the level bound over the mask), hence a smaller value than \(z\), which is in it. The
trace of the maximum locus is then empty or the maximum locus of \(\operatorname{inv}_{\underline{\mathcal S}}(T_1|_W)\) (Lemma~\ref{bcalc:max-locus-local}, with the strictly
monotone map \(\lambda\mapsto(s,\lambda)\)). The consequences are Lemma~\ref{hir:step-one-value-masked-basic}(3) with these lower data.
\end{proof}

\begin{theorem}[Step 1: the value drops over the centre]
\label{hir:step-one-special-drop}
\uses{hir:step-one-value-masked,hir:step-one-value-masked-basic,hir:step-one-canonical-tail,hir:step-one-mask-tower,hir:step-one-value-off-centre,hir:step-one-lower-domain,bcalc:boundary-count-drop-along,bcalc:sub-divisor,bcalc:codim-one-along,bcalc:codim-one-part-along,bcalc:analytic-fill,hir:centre-function}
\source{wlodarczyk05-hironaka}
\lean{TranscendentalNakamaye.Pending.StepOneSpecialDrop.level_le_of_mem_maxLocus,TranscendentalNakamaye.Pending.StepOneSpecialDrop.mem_maxLocus_of_isSpecial,TranscendentalNakamaye.Pending.StepOneSpecialDrop.notMem_cosupp_step_of_isSpecial_zero,TranscendentalNakamaye.Pending.StepOneSpecialDrop.level_step_lt_of_isSpecial,TranscendentalNakamaye.Pending.StepOneSpecialDrop.stepOneInvL_step_lt_of_isSpecial,TranscendentalNakamaye.Pending.StepOneSpecialDrop.stepOneInvM_decrease,TranscendentalNakamaye.Pending.StepOneSpecialDrop.stepOneInvM_emptyStep}
\leanok
In the situation of Lemma~\ref{hir:step-one-lower-domain} (with \(d\ge1\)), let \(\overline{\operatorname{inv}}\) be the Step 1 value function for any family of lower
marked ideals (Definition~\ref{hir:step-one-value-masked}), \(T\sigma\) a multiple test blow-up of \(\overline{\underline{\mathcal J}}\) whose last centre \(C\) is the maximum
locus of \(\overline{\operatorname{inv}}(T)\) on the cosupport, and \(x\in C\) a special point. Then \(\overline{\operatorname{inv}}(T\sigma)(z)<\overline{\operatorname{inv}}(T)(x)\)
at every point \(z\) of the cosupport over \(x\): if \(x\) has level \(0\) there is no such point, and if it has level \(s\ge1\) the level of \(z\) is \(<s\).
Consequently, for the masked value function and a resolution invariant in dimension \(\le k-1\), property (3) of Definition~\ref{hir:centre-function} holds along
every canonical segment \(T\) of \(\overline{\operatorname{inv}}\) such that \(T\sigma\) is a multiple test blow-up, \(\sigma\) the blowup of the maximum locus: the value
drops over \(C\) and is unchanged off it; and property (4) holds along every canonical segment which is a multiple test blow-up: a blowup along the
empty centre does not change the values.
\end{theorem}
\begin{proof}
\leanok
A special point \(x\) of level \(s\) has the value \((s,\infty)\), the maximum, so every point of the cosupport has level \(\le s\), and the special points of level
\(s\) lie in \(C\). If \(s=0\) (Włodarczyk's Step~1ba), \(x\) lies in the codimension-one part \(Z_1\) of the cosupport, which near \(x\) is the whole cosupport
(Lemma~\ref{bcalc:codim-one-part-along}); its points there are special of level \(0\), so near \(x\) the centre is the cosupport, hence regular of codimension one in
\(N\) (Lemma~\ref{bcalc:analytic-fill}), and by maximal order (Lemma~\ref{hir:step-one-lower-domain}) nothing of the cosupport lies over \(x\)
(Lemma~\ref{bcalc:codim-one-along}). If \(s\ge1\) (Step~1aa), \(x\) lies in the member locus of the original members \(E\) (a sub-divisor,
Lemma~\ref{bcalc:sub-divisor}): near \(x\), the points of \(N\) on the members of \(E\) through \(x\) are in the cosupport, of level \(\ge s\), hence of level \(s\) and special,
so they lie in \(C\); Lemma~\ref{bcalc:boundary-count-drop-along} then makes the number of original members through \(z\) smaller than \(s\). In both cases the first
coordinate drops (Lemma~\ref{hir:step-one-value-masked-basic}). For the consequence, the points with \(\sigma(z)\notin C\) and the non-special points of \(C\) of the same
level are Theorem~\ref{hir:step-one-canonical-tail}: such a \(\sigma(z)\) is in the regime of its level, so its mask contains it and \(z\), of the same level, is not
special (Lemma~\ref{hir:step-one-mask-tower}); if the level drops (it never grows, Lemma~\ref{hir:step-one-value-off-centre}), so does the value. Property (4) is the case \(C=\emptyset\) of
the points off the centre (the empty centre is admissible).
\end{proof}

\begin{lemma}[Step 1: admissibility at the top values]
\label{hir:step-one-admissible-top}
\uses{hir:step-one-value-masked,hir:step-one-value-masked-basic,hir:step-one-lower-domain,hir:wellformed-tower,hir:open-restriction-marked,hir:snc-invariance,bcalc:member-locus-centre,bcalc:member-locus-canonical,bcalc:boundary-count-canonical,bcalc:codim-one-part-along,bcalc:codim-one-admissible-along,mc:contact-stage-domain-masked,hir:centre-function}
\source{wlodarczyk05-hironaka}
\lean{TranscendentalNakamaye.Pending.StepOneAdmissibleTop.one_le_boundaryCount_of_mem,TranscendentalNakamaye.Pending.StepOneAdmissibleTop.level_le_of_stepOneInvL_le,TranscendentalNakamaye.Pending.StepOneAdmissibleTop.stepOneInvL_eq_top_iff,TranscendentalNakamaye.Pending.StepOneAdmissibleTop.isAdmissibleCentre_top_of_one_le,TranscendentalNakamaye.Pending.StepOneAdmissibleTop.isAdmissibleCentre_top_zero,TranscendentalNakamaye.Pending.StepOneAdmissibleTop.stepOneInvL_admissible_top}
\leanok
In the situation of Lemma~\ref{hir:step-one-lower-domain} (with \(d\ge1\)), let \(\overline{\operatorname{inv}}\) be the Step 1 value function for any family of lower
marked ideals (Definition~\ref{hir:step-one-value-masked}) and \(T\) a multiple test blow-up of \(\overline{\underline{\mathcal J}}\). A point of the cosupport has the value
\((s,\infty)\) iff it is special of level \(s\). If \(W\) is an open set of the last stage of \(T\) on which \(\overline{\operatorname{inv}}(T)\le(s,\infty)\) on the cosupport,
the points of \(W\) in the cosupport with value \((s,\infty)\) form an admissible centre of the restriction of \(\overline{\underline{\mathcal J}}_T\) to \(W\); that is,
property (1) of Definition~\ref{hir:centre-function} holds at the values \((s,\infty)\).
\end{lemma}
\begin{proof}
\leanok
A lower value is never \(\infty\) (Lemma~\ref{hir:step-one-value-masked-basic}), so the value \((s,\infty)\) is taken exactly at the special points of level \(s\), and
on \(W\) every point of the cosupport has level \(\le s\). The restriction to \(W\) is well formed (Lemmas~\ref{hir:wellformed-tower} and~\ref{hir:open-restriction-marked})
and its member locus and boundary count along the restricted original members are the traces of those of \(\overline{\underline{\mathcal J}}_T\)
(Lemmas~\ref{bcalc:member-locus-canonical} and~\ref{bcalc:boundary-count-canonical}). If \(s\ge1\) (Włodarczyk's Step~1aa), the special points of level \(s\) in \(W\)
are the level-\(s\) part of the member locus of the restriction, with \(s\) the maximal level there, an admissible centre by Lemma~\ref{bcalc:member-locus-centre}.
If \(s=0\) (Step~1ba), they are the trace on \(W\) of the codimension-one part \(Z_1\) of the cosupport, a closed submanifold by maximal order
(Lemmas~\ref{hir:step-one-lower-domain} and~\ref{bcalc:codim-one-part-along}) transverse to the exceptional members (Lemma~\ref{bcalc:codim-one-admissible-along}),
and over \(W\) the cosupport meets no original member, so the trace is transverse to the whole restricted divisor (Lemma~\ref{hir:snc-invariance} and
Lemma~\ref{mc:contact-stage-domain-masked}).
\end{proof}

\begin{lemma}[Step 1: upper semicontinuity of the value]
\label{hir:step-one-usc}
\uses{hir:step-one-value-masked,hir:step-one-value-masked-basic,hir:step-one-canonical-tail,hir:step-one-compare,bcalc:boundary-count,hir:centre-function,hir:invariant-driver}
\source{wlodarczyk05-hironaka}
\lean{TranscendentalNakamaye.Pending.StepOneUSC.eventually_forall_mem_of_eventually,TranscendentalNakamaye.Pending.StepOneUSC.eventually_level_le,TranscendentalNakamaye.Pending.StepOneUSC.stepOneInvM_eventually_le_of_not_isSpecial,TranscendentalNakamaye.Pending.StepOneUSC.stepOneInvM_usc}
\leanok
In the situation of Theorem~\ref{hir:step-one-canonical-tail}, along every multiple test blow-up \(T\) of \(\overline{\underline{\mathcal J}}\) which is a canonical
segment of the masked value function \(\overline{\operatorname{inv}}\), the function \(\overline{\operatorname{inv}}(T)\) is upper semicontinuous on
\(\operatorname{cosupp}\overline{\underline{\mathcal J}}_T\): every point \(x\) of the cosupport has a neighbourhood whose points in the cosupport have values
\(\le\overline{\operatorname{inv}}(T)(x)\). That is, property (2) of Definition~\ref{hir:centre-function} holds along these segments.
\end{lemma}
\begin{proof}
\leanok
The level \(s(y)\), the number of original members through \(y\), is upper semicontinuous, the members being closed (Lemma~\ref{bcalc:boundary-count}). At a
special point \(x\) of level \(s\) the value is \((s,\infty)\), the largest value of level \(s\) (Lemma~\ref{hir:step-one-value-masked-basic}), and near \(x\) every
level is \(\le s\). Let \(x\) be non-special of level \(s\). By Theorem~\ref{hir:step-one-canonical-tail} it has lower data at the split \(T=T_0T_1\) of its run: with mask
\(W\) and lower marked ideal \(\underline{\mathcal S}\), the restricted tail \(T_1|_W\) is a canonical segment of \(\underline{\mathcal S}\), hence a multiple test blow-up
(Theorem~\ref{hir:invariant-driver}), and \(x\) lies in the cosupport of \(\underline{\mathcal S}\) along it. Over the open preimage of \(W\) the points of the cosupport have
level \(\le s\), and those of level \(s\) lie in the cosupport of \(\underline{\mathcal S}\) along \(T_1|_W\) (Lemma~\ref{hir:step-one-compare}(4)), where
\(\overline{\operatorname{inv}}(T)=(s,\operatorname{inv}_{\underline{\mathcal S}}(T_1|_W))\) (the value at a point of the cosupport of the lower marked ideal, in the proof of
Theorem~\ref{hir:step-one-canonical-tail}); \(x\) is one of them. Property (2) of \(\operatorname{inv}\) for \(\underline{\mathcal S}\) along \(T_1|_W\) gives a neighbourhood
of \(x\) on which \(\operatorname{inv}_{\underline{\mathcal S}}(T_1|_W)\le\operatorname{inv}_{\underline{\mathcal S}}(T_1|_W)(x)\) on that cosupport; inside the open preimage
of \(W\), the points of level \(<s\) have smaller values and those of level \(s\) have values \(\le\overline{\operatorname{inv}}(T)(x)\).
\end{proof}

\begin{lemma}[Step 1: admissibility of the level sets]
\label{hir:step-one-admissible}
\uses{hir:step-one-value-masked,hir:step-one-value-masked-basic,hir:step-one-canonical-tail,hir:step-one-compare,hir:step-one-admissible-top,hir:step-one-usc,bcalc:member-restrict-admissible-local,mc:contact-admissible,bcalc:centre-local,hir:biholomorphism-invariance,hir:wellformed-tower,hir:giraud-along,hir:centre-function,hir:invariant-driver}
\source{wlodarczyk05-hironaka}
\lean{TranscendentalNakamaye.Pending.StepOneAdmissible.AdmissibleTransfer,TranscendentalNakamaye.Pending.StepOneAdmissible.HasAdmissibleTransfer,TranscendentalNakamaye.Pending.StepOneAdmissible.hasAdmissibleTransfer_succ,TranscendentalNakamaye.Pending.StepOneAdmissible.hasAdmissibleTransfer_zero,TranscendentalNakamaye.Pending.StepOneAdmissible.hasAdmissibleTransfer,TranscendentalNakamaye.Pending.StepOneAdmissible.exists_opens_preimage_eq_empty,TranscendentalNakamaye.Pending.StepOneAdmissible.exists_opens_isAdmissibleCentre_of_eq,TranscendentalNakamaye.Pending.StepOneAdmissible.stepOneInvM_admissible_coe_of_hasAdmissibleTransfer,TranscendentalNakamaye.Pending.StepOneAdmissible.stepOneInvM_admissible_succ_coe,TranscendentalNakamaye.Pending.StepOneAdmissible.stepOneInvM_admissible_coe,TranscendentalNakamaye.Pending.StepOneAdmissible.stepOneInvM_admissible}
\leanok
In the situation of Theorem~\ref{hir:step-one-canonical-tail}, let \(T\) be a multiple test blow-up of \(\overline{\underline{\mathcal J}}\) which is a canonical segment
of the masked value function \(\overline{\operatorname{inv}}\), \(W\) an open set of its last stage and \(l\) a value with \(\overline{\operatorname{inv}}(T)\le l\) on
\(W\cap\operatorname{cosupp}\overline{\underline{\mathcal J}}_T\). Then the points of \(W\) in the cosupport with value \(l\) form an admissible centre of the restriction of
\(\overline{\underline{\mathcal J}}_T\) to \(W\); that is, property (1) of Definition~\ref{hir:centre-function} holds along these segments.
\end{lemma}
\begin{proof}
\leanok
For \(l=(s,\infty)\) this is Lemma~\ref{hir:step-one-admissible-top}. Let \(l=(s,\lambda)\) with \(\lambda\) a value of \(\operatorname{inv}\). Admissible centres are local
relative to \(W\) (Lemma~\ref{bcalc:centre-local}): it suffices that every \(x\in W\) has an open neighbourhood \(V\subset W\) on which the trace of the level set is an
admissible centre of the restriction to \(V\). If \(x\) is not in the cosupport, which is closed (Lemmas~\ref{hir:wellformed-tower} and~\ref{hir:giraud-along}), or has a value \(<l\), a
neighbourhood misses the level set (Lemma~\ref{hir:step-one-usc}), and the empty set is admissible. Let \(x\) have the value \(l\). It is not special (its value is not
\((s,\infty)\)), of level \(s\), and by Theorem~\ref{hir:step-one-canonical-tail} it has lower data at the split \(T=T_0T_1\) of its run: a mask \(W_x\), a lower marked
ideal \(\underline{\mathcal S}\) on it in the domain of \(\operatorname{inv}\), and \(T_1|_{W_x}\) a canonical segment of it. Put \(V=W\cap T_1^{-1}W_x\). Over \(W_x\) the points
of the cosupport have level \(\le s\), those of level \(s\) are the points of the cosupport of \(\underline{\mathcal S}\) along \(T_1|_{W_x}\)
(Lemma~\ref{hir:step-one-compare}(4)), and there \(\overline{\operatorname{inv}}(T)=(s,\operatorname{inv}_{\underline{\mathcal S}}(T_1|_{W_x}))\) (proof of
Theorem~\ref{hir:step-one-canonical-tail}). Hence on \(V\) the level set is \(\{\operatorname{inv}_{\underline{\mathcal S}}(T_1|_{W_x})=\lambda\}\) on the cosupport of
\(\underline{\mathcal S}\), on which \(\operatorname{inv}_{\underline{\mathcal S}}(T_1|_{W_x})\le\lambda\); by property (1) of \(\operatorname{inv}\) it is an admissible centre of
the restriction to \(V\) of the transform of \(\underline{\mathcal S}\) along \(T_1|_{W_x}\). It is then an admissible centre of the restriction to \(V\) of
\(\overline{\underline{\mathcal J}}_T\): for \(s\ge1\) by Włodarczyk's le:S for one centre (Lemma~\ref{bcalc:member-restrict-admissible-local}), for \(s=0\) because the
divisor is kept and the cosupport of the marked ideal of Step 1b lies in that of \(\overline{\underline{\mathcal J}}_T\) (Lemma~\ref{mc:contact-admissible}); the open set
\(V\) seen inside \(T_1^{-1}W_x\) and \(V\) itself are identified by the identity on points (Lemma~\ref{bcalc:centre-local}, Lemma~\ref{hir:biholomorphism-invariance}).
\end{proof}

\begin{corollary}[Step 1: canonical segments are multiple test blow-ups]
\label{hir:step-one-canonical-test}
\uses{hir:step-one-admissible,hir:step-one-canonical-tail,hir:step-one-usc,hir:step-one-special-drop,hir:centre-function,hir:invariant-driver}
\source{wlodarczyk05-hironaka}
\lean{TranscendentalNakamaye.Pending.StepOneAdmissible.isAdmissibleCentre_maxLocus_stepOneInvM,TranscendentalNakamaye.Pending.StepOneAdmissible.isGeoTestTower_of_isCanonical_stepOneInvM,TranscendentalNakamaye.Pending.StepOneAdmissible.hasLowerDataL_stepOneInvM_of_isCanonical,TranscendentalNakamaye.Pending.StepOneAdmissible.stepOneInvM_admissible_of_isCanonical,TranscendentalNakamaye.Pending.StepOneAdmissible.stepOneInvM_usc_of_isCanonical,TranscendentalNakamaye.Pending.StepOneAdmissible.stepOneInvM_decrease_of_isCanonical,TranscendentalNakamaye.Pending.StepOneAdmissible.stepOneInvM_emptyStep_of_isCanonical}
\leanok
In the situation of Theorem~\ref{hir:step-one-canonical-tail}, every canonical segment of the masked value function \(\overline{\operatorname{inv}}\) is a multiple test
blow-up of \(\overline{\underline{\mathcal J}}\). Consequently Theorem~\ref{hir:step-one-canonical-tail}, Lemma~\ref{hir:step-one-admissible},
Lemma~\ref{hir:step-one-usc} and the properties (3), (4) of Theorem~\ref{hir:step-one-special-drop} hold along every canonical segment of \(\overline{\operatorname{inv}}\),
without assuming that it is a multiple test blow-up: properties (1)--(4) of Definition~\ref{hir:centre-function} hold for \(\overline{\operatorname{inv}}\).
\end{corollary}
\begin{proof}
\leanok
Induction on the length. If \(T\) is a canonical segment and a multiple test blow-up, the maximum locus of \(\overline{\operatorname{inv}}(T)\) on the cosupport is empty or
the level set of its maximum value \(l\), on which \(\overline{\operatorname{inv}}(T)\le l\) everywhere; by Lemma~\ref{hir:step-one-admissible} with \(W\) the whole last
stage it is an admissible centre (Theorem~\ref{hir:invariant-driver}: admissible for the restriction to the whole space means admissible). So the next centre, empty or
the maximum locus, is admissible, and \(T\sigma\) is a multiple test blow-up. The consequences are the cited statements with this hypothesis supplied.
\end{proof}

\begin{theorem}[Step 1: the value re-based across an empty prefix]
\label{hir:step-one-empty-prefix}
\uses{hir:step-one-value-masked,hir:step-one-value-masked-basic,hir:step-one-canonical-tail,hir:step-one-canonical-test,hir:step-one-compare,hir:regime-run,hir:tower-split,mc:contact-empty-prefix,mc:contact-lower-value,mc:coefficient-representatives-along,mc:contact-representative,bcalc:member-restriction-rebase,bcalc:inv-rebase,bcalc:member-locus-def,bcalc:boundary-count,hir:restricted-equivalence-basic,hir:marked-equivalence-basic,hir:centre-function}
\source{wlodarczyk05-hironaka}
\lean{TranscendentalNakamaye.Pending.StepOneEmptyPrefix.boundaryCount_append_of_forall_notMem,TranscendentalNakamaye.Pending.StepOneEmptyPrefix.memberInterLocus_append_of_forall_eq_empty,TranscendentalNakamaye.Pending.StepOneEmptyPrefix.drop_E_geoTransformAlong_of_isEmptyTower,TranscendentalNakamaye.Pending.StepOneEmptyPrefix.strictAlong_empty,TranscendentalNakamaye.Pending.StepOneEmptyPrefix.restrictEquiv_barMarked_append,TranscendentalNakamaye.Pending.StepOneEmptyPrefix.origMembers_emptyPrefix,TranscendentalNakamaye.Pending.StepOneEmptyPrefix.level_emptyPrefix,TranscendentalNakamaye.Pending.StepOneEmptyPrefix.isSpecial_emptyPrefix_iff,TranscendentalNakamaye.Pending.StepOneEmptyPrefix.stepRegime_emptyPrefix_iff,TranscendentalNakamaye.Pending.StepOneEmptyPrefix.maskA_emptyPrefix,TranscendentalNakamaye.Pending.StepOneEmptyPrefix.membersThrough_append_of_forall_eq_empty,TranscendentalNakamaye.Pending.StepOneEmptyPrefix.restrictEquiv_memberRestrict_membersThrough,TranscendentalNakamaye.Pending.StepOneEmptyPrefix.lowerData_congr,TranscendentalNakamaye.Pending.StepOneEmptyPrefix.lowerDataL_succ_emptyPrefix,TranscendentalNakamaye.Pending.StepOneEmptyPrefix.lowerValueL_succ_emptyPrefix_refl,TranscendentalNakamaye.Pending.StepOneEmptyPrefix.regimeRun_append_of_lt,TranscendentalNakamaye.Pending.StepOneEmptyPrefix.regimeRun_append_of_eq,TranscendentalNakamaye.Pending.StepOneEmptyPrefix.regime_base_of_regimeRun_eq,TranscendentalNakamaye.Pending.StepOneEmptyPrefix.regimeRun_eq_length_of_isEmptyTower,TranscendentalNakamaye.Pending.StepOneEmptyPrefix.split_zero,TranscendentalNakamaye.Pending.StepOneEmptyPrefix.split_append_of_le,TranscendentalNakamaye.Pending.StepOneEmptyPrefix.split_of_length_le,TranscendentalNakamaye.Pending.StepOneEmptyPrefix.hypotheses_emptyPrefix,TranscendentalNakamaye.Pending.StepOneEmptyPrefix.stepOneInvM_emptyPrefix_eq,TranscendentalNakamaye.Pending.StepOneEmptyPrefix.stepOneInvM_emptyPrefix}
\leanok
In the situation of Theorem~\ref{hir:step-one-canonical-tail}, let \(P\) be a sequence of blowups with only empty centres and
\(\underline{\mathcal J}_P\) the transform of \(\underline{\mathcal J}\) along it. If \(PT_1\) is a canonical segment of the masked value function
\(\overline{\operatorname{inv}}\) of \(\underline{\mathcal J}\), then \(T_1\) is a canonical segment of the masked value function of \(\underline{\mathcal J}_P\), and at every point of
the cosupport the value of \(\underline{\mathcal J}\) along \(PT_1\) is that of \(\underline{\mathcal J}_P\) along \(T_1\); that is, property (4') of
Definition~\ref{hir:centre-function} holds.
\end{theorem}
\begin{proof}
\leanok
Induction on \(T_1\); along a canonical segment of both, the maximum loci agree once the values and the cosupports agree, so it suffices to compare the values along
\(PT_1\) and \(T_1\) when both are canonical segments (multiple test blow-ups by Corollary~\ref{hir:step-one-canonical-test}). The marked ideal
\(\overline{\underline{\mathcal J}_P}\) of Step 1 of \(\underline{\mathcal J}_P\) is restricted-equivalent to the transform \(\overline{\underline{\mathcal J}}_P\)
(Lemma~\ref{mc:contact-empty-prefix}); restricted-equivalent marked ideals have the same transforms and cosupports
(Lemmas~\ref{hir:restricted-equivalence-basic} and~\ref{hir:marked-equivalence-basic}), so along \(T_1\) and \(PT_1\) the two have the same \(N\), divisor and
cosupport. The original members of \(\underline{\mathcal J}_P\) are those of \(\underline{\mathcal J}\) followed by the exceptional members of \(P\), which are empty
(the blowups are along \(\emptyset\)); empty members pass through no point and do not change the member locus (Lemmas~\ref{bcalc:boundary-count}
and~\ref{bcalc:member-locus-def}). Hence levels, special points, regimes and the masks of Step 1ab agree, and so do the values at special points. At a non-special point
\(z\) of level \(s\) let \(r\) be its run along \(T_1\) for \(\underline{\mathcal J}_P\) (Lemma~\ref{hir:regime-run}). If \(r<\operatorname{length}T_1\), its run along
\(PT_1\) is \(r\) too, and the splits are \(P(H_0)\cdot T_1'\) and \(H_0\cdot T_1'\) (Lemma~\ref{hir:tower-split}). If \(r=\operatorname{length}T_1\), the image of \(z\) at the
start of \(T_1\) is in the regime, empty steps keep the regime, and the run along \(PT_1\) is the whole sequence: the splits start at the two bases. For \(s=0\) the lower
values agree by Lemma~\ref{mc:contact-lower-value}. For \(s\ge1\), Step 1ab's restricted marked ideals of restricted-equivalent marked ideals at the same original
members are restricted-equivalent (Lemma~\ref{mc:coefficient-representatives-along}), the members of the empty prefix through the point being none; on the common mask the
two lower marked ideals have the same domain, canonical segments and values by property (7) of the lower invariant (Lemmas~\ref{mc:contact-representative}
and~\ref{bcalc:inv-rebase}). When the run covers \(T_1\), the lower value of \(\underline{\mathcal J}\) formed at the base of \(PT_1\) is first re-based across \(P\),
whose centres are empty (Lemma~\ref{hir:step-one-compare}(2), Lemma~\ref{bcalc:member-restriction-rebase}).
\end{proof}

\begin{theorem}[Step 1: independence of the representative]
\label{hir:step-one-representative}
\uses{hir:step-one-value-masked,hir:step-one-value-masked-basic,hir:step-one-canonical-tail,hir:step-one-canonical-test,hir:regime-run,mc:contact-representative,mc:contact-lower-value,bcalc:member-restrict-restrict-repr,hir:marked-equivalence-basic,hir:restricted-equivalence-basic,hir:canonical-segments,hir:centre-function}
\source{wlodarczyk05-hironaka}
\lean{TranscendentalNakamaye.Pending.StepOneRepresentative.regimeRun_congr,TranscendentalNakamaye.Pending.StepOneRepresentative.isSpecial_iff_of_eq,TranscendentalNakamaye.Pending.StepOneRepresentative.restrictEquiv_facts,TranscendentalNakamaye.Pending.StepOneRepresentative.stepOneInvM_representative_eq,TranscendentalNakamaye.Pending.StepOneRepresentative.stepOneInvM_representative}
\leanok
In the situation of Theorem~\ref{hir:step-one-canonical-tail}, let \(\underline{\mathcal J}'\) be restricted-equivalent to \(\underline{\mathcal J}\) and satisfy the same
hypotheses. Every canonical segment \(T\) of the masked value function of \(\underline{\mathcal J}\) is one of that of \(\underline{\mathcal J}'\), and the two values agree
at every point of the cosupport of \(\overline{\underline{\mathcal J}}_T\); that is, property (7) of Definition~\ref{hir:centre-function} holds.
\end{theorem}
\begin{proof}
\leanok
The marked ideals of Step 1 of \(\underline{\mathcal J}\) and \(\underline{\mathcal J}'\) are restricted-equivalent (Lemma~\ref{mc:contact-representative}), so along every
sequence of blowups their transforms have the same \(N\), the same cosupport and the same multiple test blow-ups (Lemmas~\ref{hir:restricted-equivalence-basic}
and~\ref{hir:marked-equivalence-basic}); the divisors are equal, so the original members, the levels, the special points, the regimes and the runs
(Lemma~\ref{hir:regime-run}) are the same. At a non-special point both have lower data at the same split (Theorem~\ref{hir:step-one-canonical-tail} on both sides),
and the lower values agree: at level \(0\) by Lemma~\ref{mc:contact-lower-value}, at levels \(\ge1\) by Lemma~\ref{bcalc:member-restrict-restrict-repr}. By induction
on the length, the maximum loci of the two values on the common cosupport agree, so the canonical segments agree.
\end{proof}

\begin{theorem}[Step 1: restriction to open sets]
\label{hir:step-one-restrict-open}
\uses{hir:step-one-value-masked,hir:step-one-value-masked-basic,hir:step-one-canonical-tail,hir:step-one-canonical-test,mc:contact-stage-kept,hir:open-restriction-tower,hir:open-restriction-marked,bcalc:codim-one-part-canonical,bcalc:member-locus-canonical,bcalc:nonmonomial-order-regime,hir:tower-split-iso,mc:contact-lower-value,bcalc:member-restrict-restrict-repr,hir:canonical-segments,hir:centre-function}
\source{wlodarczyk05-hironaka}
\lean{TranscendentalNakamaye.Pending.StepOneRestrictOpen.geoTransformAlong_barMarked_restrict,TranscendentalNakamaye.Pending.StepOneRestrictOpen.restrict_facts,TranscendentalNakamaye.Pending.StepOneRestrictOpen.hypotheses_restrict,TranscendentalNakamaye.Pending.StepOneRestrictOpen.stepOneInvM_restrict_eq,TranscendentalNakamaye.Pending.StepOneRestrictOpen.stepOneInvM_restrict}
\leanok
In the situation of Theorem~\ref{hir:step-one-canonical-tail}, let \(U\subset X\) be open. For every canonical segment \(T\) of the masked value function of
\(\underline{\mathcal J}\), the restriction \(T|_U\) is a canonical segment of the masked value function of \(\underline{\mathcal J}|_U\), and at the points of the
cosupport over \(U\) the two values agree; that is, property (5) of Definition~\ref{hir:centre-function} holds.
\end{theorem}
\begin{proof}
\leanok
The marked ideal of Step 1 commutes with restriction (Lemma~\ref{mc:contact-stage-kept}), and so do transforms along multiple test blow-ups
(Lemma~\ref{hir:open-restriction-tower}); the restricted marked ideal satisfies the hypotheses of Step 1 (Lemma~\ref{hir:open-restriction-marked}). Hence the original
members, the levels and the cosupport along \(T|_U\) are the traces of those along \(T\), and so are the codimension-one part and the member locus
(Lemmas~\ref{bcalc:codim-one-part-canonical} and~\ref{bcalc:member-locus-canonical}): special points and regimes correspond, and so do the runs
(Lemma~\ref{bcalc:nonmonomial-order-regime}). The split of \(T|_U\) is isomorphic over \(\operatorname{id}_U\) to the restriction of the split of \(T\)
(Lemma~\ref{hir:tower-split-iso}); with lower data on both sides (Theorem~\ref{hir:step-one-canonical-tail}), the lower values agree at level \(0\) by
Lemma~\ref{mc:contact-lower-value} and at levels \(\ge1\) by Lemma~\ref{bcalc:member-restrict-restrict-repr}. By induction on the length, the trace of each centre is
empty or the maximum locus of the restricted value (Lemma~\ref{hir:canonical-segments}), so \(T|_U\) is a canonical segment.
\end{proof}

\begin{theorem}[Step 1: isomorphisms of towers]
\label{hir:step-one-tower-iso}
\uses{hir:step-one-value-masked,hir:step-one-value-masked-basic,hir:step-one-canonical-tail,hir:step-one-canonical-test,mc:contact-stage-kept,mc:tower-isomorphism-transform,hir:biholomorphism-invariance,bcalc:codim-one-part-canonical,bcalc:member-locus-canonical,bcalc:nonmonomial-order-regime,hir:tower-split-iso,mc:contact-lower-value,bcalc:member-restrict-tower-iso,hir:canonical-segments,hir:tower-isomorphism,hir:centre-function}
\source{wlodarczyk05-hironaka}
\lean{TranscendentalNakamaye.Pending.StepOneTowerIso.geoTransformAlong_barMarked_isTowerIso,TranscendentalNakamaye.Pending.StepOneTowerIso.isTowerIso_facts,TranscendentalNakamaye.Pending.StepOneTowerIso.hypotheses_pullback,TranscendentalNakamaye.Pending.StepOneTowerIso.stepOneInvM_isTowerIso_eq,TranscendentalNakamaye.Pending.StepOneTowerIso.stepOneInvM_isTowerIso}
\leanok
In the situation of Theorem~\ref{hir:step-one-canonical-tail}, let \(\Phi\colon X'\to X\) be a biholomorphism and \(\chi\colon T'\to T\) an isomorphism of sequences of
blowups over \(\Phi\) (Definition~\ref{hir:tower-isomorphism}). If \(T\) is a canonical segment of the masked value function of \(\underline{\mathcal J}\), then \(T'\) is
one of that of \(\Phi^*\underline{\mathcal J}\), and the value of \(\Phi^*\underline{\mathcal J}\) along \(T'\) at \(p\) is that of \(\underline{\mathcal J}\) along \(T\) at
\(\chi(p)\), at every \(p\) with \(\chi(p)\) in the cosupport; that is, property (6) of Definition~\ref{hir:centre-function} holds.
\end{theorem}
\begin{proof}
\leanok
The marked ideal of Step 1 commutes with biholomorphisms (Lemma~\ref{mc:contact-stage-kept}) and transforms correspond along isomorphisms of towers
(Lemma~\ref{mc:tower-isomorphism-transform}); \(\Phi^*\underline{\mathcal J}\) satisfies the hypotheses of Step 1 (Lemma~\ref{hir:biholomorphism-invariance}). Hence the
original members, the levels and the cosupport along \(T'\) are the \(\chi\)-preimages of those along \(T\), and so are the codimension-one part and the member locus
(Lemmas~\ref{bcalc:codim-one-part-canonical} and~\ref{bcalc:member-locus-canonical}): special points and regimes correspond, and so do the runs
(Lemma~\ref{bcalc:nonmonomial-order-regime}). The splits correspond along isomorphisms (Lemma~\ref{hir:tower-split-iso}); with lower data on both sides
(Theorem~\ref{hir:step-one-canonical-tail}), the lower values agree at level \(0\) by Lemma~\ref{mc:contact-lower-value} and at levels \(\ge1\) by
Lemma~\ref{bcalc:member-restrict-tower-iso}. By induction on the isomorphism, the centres of \(T'\), the \(\chi\)-preimages of those of \(T\), are empty or the maximum loci
(Lemma~\ref{hir:canonical-segments}, \(\chi\) being surjective).
\end{proof}

\begin{corollary}[Step 1: the properties for the marked ideal itself]
\label{hir:step-one-bridge}
\uses{hir:step-one-admissible,hir:step-one-canonical-test,hir:step-one-usc,hir:step-one-special-drop,hir:step-one-empty-prefix,hir:step-one-restrict-open,hir:step-one-tower-iso,hir:step-one-representative,hir:along-calculus,hir:marked-equivalence-basic,hir:open-restriction-marked,hir:centre-function}
\source{wlodarczyk05-hironaka}
\lean{TranscendentalNakamaye.Pending.StepOneBridge.isAdmissibleCentre_pullback_val_iff_of_eq,TranscendentalNakamaye.Pending.StepOneBridge.barMarked_wEquiv,TranscendentalNakamaye.Pending.StepOneBridge.isCanonical_iff,TranscendentalNakamaye.Pending.StepOneBridge.isGeoTestTower_of_isCanonical,TranscendentalNakamaye.Pending.StepOneBridge.admissible_field,TranscendentalNakamaye.Pending.StepOneBridge.usc_field,TranscendentalNakamaye.Pending.StepOneBridge.decrease_field,TranscendentalNakamaye.Pending.StepOneBridge.emptyStep_field,TranscendentalNakamaye.Pending.StepOneBridge.emptyPrefix_field,TranscendentalNakamaye.Pending.StepOneBridge.restrictOpen_field,TranscendentalNakamaye.Pending.StepOneBridge.towerIso_field,TranscendentalNakamaye.Pending.StepOneBridge.representative_field}
\leanok
Let \(\underline{\mathcal J}\) lie in the domain of Definition~\ref{hir:centre-function} in dimension \(\le k\), be of maximal order along \(N\), and let
\(\operatorname{inv}\) be a resolution invariant in dimension \(\le k-1\). The canonical segments of the masked value function \(\overline{\operatorname{inv}}\) of Step 1
defined with the maximum loci on \(\operatorname{cosupp}\underline{\mathcal J}_T\) are those defined with \(\operatorname{cosupp}\overline{\underline{\mathcal J}}_T\), the two
cosupports agree along them, and properties (1)--(7) and (4') of Definition~\ref{hir:centre-function} hold for \(\overline{\operatorname{inv}}\) as stated there, for
the marked ideals of maximal order in the domain.
\end{corollary}
\begin{proof}
\leanok
\(\overline{\underline{\mathcal J}}=\mathcal C_N(\mathcal H_N(\underline{\mathcal J}))\) is equivalent to \(\underline{\mathcal J}\) (Lemma~\ref{hir:along-calculus}(4)): the same divisor, the same
multiple test blow-ups and, along them, the same cosupports. By induction on the length the two notions of canonical segment agree, the maximum loci being taken on equal
sets along segments which are multiple test blow-ups (Corollary~\ref{hir:step-one-canonical-test}). Admissible centres on an open set depend only on the divisor and the
cosupport (Lemma~\ref{hir:marked-equivalence-basic}, the cosupport of a restriction being the trace). The properties are then
Lemmas~\ref{hir:step-one-admissible} and~\ref{hir:step-one-usc}, Theorems~\ref{hir:step-one-special-drop}, \ref{hir:step-one-empty-prefix},
\ref{hir:step-one-restrict-open}, \ref{hir:step-one-tower-iso} and~\ref{hir:step-one-representative}, with the cosupports rewritten.
\end{proof}

\begin{theorem}[The maximum drops along canonical segments; well-founded values resolve]
\label{hir:step-one-well-founded}
\uses{hir:step-one-bridge,hir:centre-function,hir:resolution-driver,hir:invariant-driver,hir:canonical-segments,bcalc:transform-off-centre,bcalc:inv-domain-stable}
\source{wlodarczyk05-hironaka,bierstone-milman08-functoriality,kollar07-resolution}
\lean{TranscendentalNakamaye.Pending.StepOneWellFounded.StepDrop,TranscendentalNakamaye.Pending.StepOneWellFounded.lt_of_mem_maxLocus,TranscendentalNakamaye.Pending.StepOneWellFounded.lt_of_isCanonical_append,TranscendentalNakamaye.Pending.StepOneWellFounded.lt_of_isCanonical_step,TranscendentalNakamaye.Pending.StepOneWellFounded.proj_mem_cosupp_step,TranscendentalNakamaye.Pending.StepOneWellFounded.stepDrop_of_resolutionInvariant,TranscendentalNakamaye.Pending.StepOneWellFounded.exists_isGeoResolution_of_wellFoundedLT,TranscendentalNakamaye.Pending.StepOneWellFounded.stepDrop_stepOneInvM,TranscendentalNakamaye.Pending.StepOneWellFounded.stepOneInvM_lt_of_isCanonical_append,TranscendentalNakamaye.Pending.StepOneWellFounded.isAdmissibleCentre_maxLocus_stepOneInvM,TranscendentalNakamaye.Pending.StepOneWellFounded.exists_isGeoResolution_stepOneInvM}
\leanok
Let \(\underline{\mathcal I}\) be a marked ideal on \(X\) and \(\operatorname{inv}\) a function on the towers over \(X\) such that after every step of a canonical segment
(the centre \(C\) empty or the maximum locus) every point \(z\) of the new cosupport lies over the cosupport, with \(\operatorname{inv}(z)<\operatorname{inv}(\pi z)\) if
\(\pi z\in C\) and \(\operatorname{inv}(z)=\operatorname{inv}(\pi z)\) otherwise.
\begin{enumerate}
\item If \(C\neq\emptyset\) is the maximum locus after a canonical segment \(T\), with value \(\lambda\), then along every canonical segment extending \(T\) by the blowup of
\(C\) every value on the cosupport is \(<\lambda\).
\item A resolution invariant (Definition~\ref{hir:centre-function}) has this property along its canonical segments; if its values are well-ordered, it resolves every
marked ideal of its domain over a compact base by a canonical segment without empty centres, without using property (8).
\item The masked Step 1 value of a marked ideal of maximal order in the domain has this property (with the maximum loci on \(\operatorname{cosupp}\underline{\mathcal J}_T\)),
its maximum locus is an admissible centre, and if the values \(\Lambda\) of the lower invariant are well-ordered, it resolves \(\underline{\mathcal J}\) over a compact base
by a canonical segment without empty centres.
\end{enumerate}
\end{theorem}
\begin{proof}
\leanok
(1) By induction along the extension. After the blowup of \(C\): a point over \(C\) has a value \(<\operatorname{inv}(\pi z)=\lambda\); a point off \(C\) keeps the value of
\(\pi z\), a point of the cosupport outside the maximum locus, so \(<\lambda\). At a later step every value is \(<\) or \(=\) the value of its image, which is \(<\lambda\) by
induction. The points of the cosupport lie over the cosupport: over the centre since the centre is admissible, off it since the blowup is an isomorphism there
(Lemma~\ref{bcalc:transform-off-centre}).
(2) Properties (3) and (4) of Definition~\ref{hir:centre-function} give the hypothesis, canonical segments being multiple test blow-ups
(Theorem~\ref{hir:invariant-driver}). With (1) for one step, the hypotheses of the well-founded form of Theorem~\ref{hir:resolution-driver} hold: closed cosupports
(Lemma~\ref{bcalc:inv-domain-stable}), (2), admissible maximum loci with empty interior (Theorem~\ref{hir:invariant-driver}), and the drop.
(3) The same with Corollary~\ref{hir:step-one-bridge} in place of the fields; \(\mathbf N\times(\Lambda\cup\{\infty\})\), ordered lexicographically, is well-ordered
when \(\Lambda\) is.
\end{proof}

\begin{theorem}[The values of the nonempty centres; the phases of the Step 1 value]
\label{hir:step-one-phase}
\uses{hir:step-one-well-founded,hir:centre-function,hir:canonical-segments,hir:step-one-value-masked,hir:step-one-value-masked-basic}
\source{wlodarczyk05-hironaka,bierstone-milman08-functoriality}
\lean{TranscendentalNakamaye.Pending.StepOnePhase.centreValues,TranscendentalNakamaye.Pending.StepOnePhase.centreValues_refl,TranscendentalNakamaye.Pending.StepOnePhase.centreValues_step_of_nonempty,TranscendentalNakamaye.Pending.StepOnePhase.centreValues_step_of_eq_empty,TranscendentalNakamaye.Pending.StepOnePhase.centreValues_append,TranscendentalNakamaye.Pending.StepOnePhase.centreValues_isPrefix_append,TranscendentalNakamaye.Pending.StepOnePhase.exists_eq_of_mem_centreValues,TranscendentalNakamaye.Pending.StepOnePhase.length_centreValues,TranscendentalNakamaye.Pending.StepOnePhase.countCentresMeeting_le_countCentresMeeting_univ,TranscendentalNakamaye.Pending.StepOnePhase.countCentresMeeting_le_length_centreValues,TranscendentalNakamaye.Pending.StepOnePhase.pairwise_and_lt_centreValues,TranscendentalNakamaye.Pending.StepOnePhase.pairwise_centreValues,TranscendentalNakamaye.Pending.StepOnePhase.lt_of_mem_centreValues,TranscendentalNakamaye.Pending.StepOnePhase.nodup_centreValues,TranscendentalNakamaye.Pending.StepOnePhase.pairwise_fst_of_pairwise,TranscendentalNakamaye.Pending.StepOnePhase.pairwise_snd_filter_fst_eq,TranscendentalNakamaye.Pending.StepOnePhase.length_eq_sum_filter_fst,TranscendentalNakamaye.Pending.StepOnePhase.level_le_length_E,TranscendentalNakamaye.Pending.StepOnePhase.fst_stepOneInvM_le,TranscendentalNakamaye.Pending.StepOnePhase.fst_le_of_mem_centreValues_stepOneInvM,TranscendentalNakamaye.Pending.StepOnePhase.pairwise_centreValues_stepOneInvM,TranscendentalNakamaye.Pending.StepOnePhase.pairwise_fst_centreValues_stepOneInvM,TranscendentalNakamaye.Pending.StepOnePhase.length_filter_special_centreValues_le,TranscendentalNakamaye.Pending.StepOnePhase.length_centreValues_eq_sum_phases,TranscendentalNakamaye.Pending.StepOnePhase.pairwise_snd_phase_stepOneInvM}
\leanok
Let \(\underline{\mathcal I}\) and \(\operatorname{inv}\) be as in Theorem~\ref{hir:step-one-well-founded}: after every step of a canonical segment every point of the new
cosupport lies over the cosupport, with a smaller value over the centre and the same value off it. For a tower \(T\) let \(v(T)\) be the sequence, in the order of
the steps, of the values of \(\operatorname{inv}\) of the preceding part at a point of each nonempty centre of \(T\).
\begin{enumerate}
\item For every tower, \(v\) of a prefix is a prefix of \(v(T)\), the length of \(v(T)\) is the number of nonempty centres of \(T\), and at most that many centres
meet the preimage of a given set.
\item If \(T\) is a canonical segment, \(v(T)\) is strictly decreasing, and every value of \(\operatorname{inv}\) on \(\operatorname{cosupp}\underline{\mathcal I}_T\) is smaller
than every entry of \(v(T)\).
\item For the masked Step 1 value of a marked ideal \(\underline{\mathcal J}\) of maximal order in the domain (Theorem~\ref{hir:step-one-well-founded}(3)) and a canonical
segment \(T\): the first entries (levels \(s\)) of the entries of \(v(T)\) do not increase and are at most \(\#E\); at most \(\#E+1\) entries have the second entry
\(\infty\), since each special value \((s,\infty)\) occurs at most once; the length of \(v(T)\) is the sum over the phases \(s=0,\dots,\#E\) of the numbers of entries of
level \(s\); and inside the phase \(s\) the second entries strictly decrease.
\end{enumerate}
So the number of nonempty centres of a canonical segment of the Step 1 value is at most \(\#E+1\) plus, for each phase \(s\), the length of a strictly decreasing
sequence of lower values occurring as maxima; bounding the latter is where property (8) of the lower-dimensional invariant enters.
\end{theorem}
\begin{proof}
\leanok
(1) By induction along the tower: a step appends one entry if its centre is nonempty and none otherwise.
(2) By induction along the segment. After a step along \(C\), a point \(z\) of the new cosupport lies over a point \(\pi z\) of the old cosupport and
\(\operatorname{inv}(z)\le\operatorname{inv}(\pi z)\), which is smaller than every earlier entry by induction. If \(C\ne\emptyset\), it is the maximum locus (a canonical
segment), so the new entry is the value \(\lambda\) at a point of \(C\subset\operatorname{cosupp}\), smaller than every earlier entry by induction, and every value after
the step is \(<\lambda\) (Theorem~\ref{hir:step-one-well-founded}(1) for one step).
(3) The level is the number of original members through the point, at most \(\#E\) (Definition~\ref{hir:step-one-value-masked}). In a strictly decreasing sequence of
lexicographic pairs the first entries do not increase, and among the entries with a given first entry the second entries strictly decrease. Two entries with the
same level and second entry \(\infty\) are equal, and the entries are pairwise distinct by (2); so the map to the level is injective on these entries, with values in
\(\{0,\dots,\#E\}\). Partitioning the entries by their level gives the sum.
\end{proof}

\begin{definition}[Value functions for the marked ideals of maximal order]
\label{hir:max-order-value}
\uses{hir:centre-function,hir:canonical-segments,hir:restricted-equivalence}
\source{wlodarczyk05-hironaka,bierstone-milman08-functoriality}
\lean{TranscendentalNakamaye.Pending.MaxOrderValue.IsMaxOrderValue}
\leanok
A \emph{value function for the marked ideals of maximal order in dimension \(\le k\)} assigns to every marked ideal \(\underline{\mathcal J}\), every tower \(T\)
over its base and every point \(y\) of the last stage a value \(M_{\underline{\mathcal J}}(T)(y)\) in a totally ordered set, such that for every marked ideal of the
domain of Definition~\ref{hir:centre-function} in dimension \(\le k\) which is of maximal order along its submanifold, properties (1)--(7) and (4') of
Definition~\ref{hir:centre-function} hold along its canonical segments (in (7) both representatives of maximal order). Property (8) is not required.
\end{definition}

\begin{lemma}[Value functions of maximal order: the Step 1 value, and canonical segments]
\label{hir:max-order-value-basic}
\uses{hir:max-order-value,hir:step-one-bridge,hir:step-one-well-founded,hir:invariant-driver,hir:canonical-segments,hir:restricted-equivalence-basic,bcalc:inv-rebase,bcalc:transform-off-centre}
\source{wlodarczyk05-hironaka,bierstone-milman08-functoriality}
\lean{TranscendentalNakamaye.Pending.MaxOrderValue.isMaxOrderValue_stepOneInvM,TranscendentalNakamaye.Pending.MaxOrderValue.IsMaxOrderValue.step_of_notMem,TranscendentalNakamaye.Pending.MaxOrderValue.IsMaxOrderValue.step_lt_of_mem,TranscendentalNakamaye.Pending.MaxOrderValue.IsMaxOrderValue.isAdmissibleCentre_maxLocus,TranscendentalNakamaye.Pending.MaxOrderValue.IsMaxOrderValue.isGeoTestTower_of_isCanonical,TranscendentalNakamaye.Pending.MaxOrderValue.IsMaxOrderValue.stepDrop,TranscendentalNakamaye.Pending.MaxOrderValue.IsMaxOrderValue.isCanonical_restrict,TranscendentalNakamaye.Pending.MaxOrderValue.IsMaxOrderValue.isCanonical_of_isTowerIso,TranscendentalNakamaye.Pending.MaxOrderValue.IsMaxOrderValue.isCanonical_append_iff,TranscendentalNakamaye.Pending.MaxOrderValue.IsMaxOrderValue.isCanonical_of_restrictEquiv,TranscendentalNakamaye.Pending.MaxOrderValue.IsMaxOrderValue.value_rebase}
\leanok
\begin{enumerate}
\item Given a resolution invariant in dimension \(\le k-1\), the masked Step 1 value (Definition~\ref{hir:step-one-value-masked}) is a value function for the marked
ideals of maximal order in dimension \(\le k\) (for \(n\ge1\)).
\end{enumerate}
Let \(M\) be a value function for the marked ideals of maximal order in dimension \(\le k\) and \(\underline{\mathcal J}\) a marked ideal of the domain, of maximal order.
\begin{enumerate}\setcounter{enumi}{1}
\item After a step of a canonical segment, a point of the new cosupport has a smaller value than its image over the centre and the same value off it; the points of
the cosupport lie over the cosupport (the hypothesis of Theorem~\ref{hir:step-one-well-founded}).
\item The maximum locus along a canonical segment which is a multiple test blow-up is an admissible centre; every canonical segment is a multiple test blow-up.
\item The restriction of a canonical segment to an open set \(U\) is a canonical segment of \(\underline{\mathcal J}|_U\); canonical segments correspond along isomorphisms
of towers; after a prefix \(T_0\) with only empty centres, \(T_0T_1\) is a canonical segment of \(\underline{\mathcal J}\) if and only if \(T_1\) is one of
\(\underline{\mathcal J}_{T_0}\); and canonical segments of restricted-equivalent marked ideals of the domain, both of maximal order, coincide.
\item (Re-basing the values.) If \(T_0\) has only empty centres, \(\underline{\mathcal J}_{T_0}\) is of maximal order, \(\underline{\mathcal N}\) (coherent ideal, of maximal
order) is restricted-equivalent to \(\underline{\mathcal J}_{T_0}\) and \(T_0T_1\) is a canonical segment of \(\underline{\mathcal J}\), then \(\underline{\mathcal N}\) lies in the
domain, \(T_1\) is a canonical segment of it, and \(M_{\underline{\mathcal J}}(T_0T_1)=M_{\underline{\mathcal N}}(T_1)\) on the cosupport.
\end{enumerate}
\end{lemma}
\begin{proof}
\leanok
(1) is Corollary~\ref{hir:step-one-bridge}. (2): by (3) and (4) of Definition~\ref{hir:centre-function} (a nonempty centre is the maximum locus), and since a test
blowup is an isomorphism off its centre (Lemma~\ref{bcalc:transform-off-centre}) while its centre lies in the cosupport. (3)--(5) are the proofs of
Theorem~\ref{hir:invariant-driver}(3), Lemma~\ref{hir:canonical-segments} and Lemma~\ref{bcalc:inv-rebase}, which use only properties (1) and (4)--(7) and (4') of the
invariant, applied to \(M\) on marked ideals of maximal order: the maximum locus of the level of the maximum is admissible on the whole last stage by (1); the maximum
loci correspond under restriction, isomorphisms, empty prefixes and restricted equivalence by (5), (6), (4') and (7), the cosupports corresponding
(Lemma~\ref{hir:restricted-equivalence-basic}); the domain is invariant under restricted equivalence.
\end{proof}

\begin{corollary}[The Step 1 value for a lower invariant with local finiteness along growing sequences]
\label{hir:step-one-seq}
\uses{hir:max-order-value-basic,hir:step-one-bridge,hir:step-one-well-founded,hir:step-one-phase,hir:centre-function-seq,hir:step-one-value-masked}
\source{wlodarczyk05-hironaka,bierstone-milman08-functoriality}
\lean{TranscendentalNakamaye.Pending.StepOneSeq.stepOneInvM_ofExt,TranscendentalNakamaye.Pending.StepOneSeq.isMaxOrderValue_stepOneInvM,TranscendentalNakamaye.Pending.StepOneSeq.stepDrop_stepOneInvM,TranscendentalNakamaye.Pending.StepOneSeq.isAdmissibleCentre_maxLocus_stepOneInvM,TranscendentalNakamaye.Pending.StepOneSeq.exists_isGeoResolution_stepOneInvM}
\leanok
Let \(n\ge1\) and let \(\operatorname{inv}\) be a resolution invariant in dimension \(\le k-1\) in the sense of Definition~\ref{hir:centre-function-seq}. Then the masked
Step 1 value defined with \(\operatorname{inv}\) (Definition~\ref{hir:step-one-value-masked}) is a value function for the marked ideals of maximal order in dimension
\(\le k\) (Definition~\ref{hir:max-order-value}); along its canonical segments it drops as in Theorem~\ref{hir:step-one-well-founded}(1), its maximum loci are admissible
centres, and if the values of \(\operatorname{inv}\) are well-ordered it resolves every marked ideal of maximal order in the domain over a compact base. The same holds for
the statements of Theorems~\ref{hir:step-one-well-founded} and~\ref{hir:step-one-phase} and Corollary~\ref{hir:step-one-bridge}. For a resolution invariant in the sense of
Definition~\ref{hir:centre-function}, read in the sense of Definition~\ref{hir:centre-function-seq}, the value is the same.
\end{corollary}
\begin{proof}
\leanok
The proofs of Lemma~\ref{hir:max-order-value-basic}(1), Corollary~\ref{hir:step-one-bridge} and Theorems~\ref{hir:step-one-well-founded} and~\ref{hir:step-one-phase}, with
everything they rest on, use the lower invariant only through its values and its properties (1)--(7) and (4'), never through property (8). Since
Definitions~\ref{hir:centre-function} and~\ref{hir:centre-function-seq} share these properties word for word, the same proofs apply.
\end{proof}

\begin{proposition}[Marked ideals of maximal order]
\label{hir:maximal-order-step}
\uses{bcalc:member-intersection,bcalc:boundary-count,bcalc:boundary-count-drop,hir:drop-containing-members,hir:centre-function,hir:homogenized-coefficient,hir:coefficient-restriction,hir:contact-equivalence,bcalc:codim-one-part,bcalc:codim-one-step,bcalc:codim-one-step-rest,hir:drop-members,hir:along-calculus,mc:glueing-lemma,mc:glueing-trivial,mc:glueing-lift,mc:glueing-tower-pullback,mc:lift-biholomorphism,mc:giraud-tower,hir:marked-equivalence-basic,hir:tower-transport-open,hir:open-biholomorphism,bcalc:snc-transform-nc,bcalc:transform-wellformed,bcalc:member-intersection-along,bcalc:boundary-count-drop-along,bcalc:member-intersection-drop-along,bcalc:codim-one-along,bcalc:codim-one-local-along,bcalc:codim-one-part-along-def,bcalc:codim-one-part-along,bcalc:codim-one-step-along,hir:giraud-tower-along,hir:tower-isomorphism,mc:contact-independence-along,mc:glueing-lift-tail,hir:canonical-segments,hir:step-one-b-independence,hir:tower-split,hir:tower-split-iso,hir:regime-run,hir:drop-original-members,hir:step-one-value,hir:step-one-value-off-centre,hir:step-one-lower-domain,hir:step-one-mask-tower,hir:step-one-nearby,hir:step-one-value-masked,hir:step-one-value-masked-basic,hir:step-one-compare,hir:step-one-canonical-tail,hir:step-one-special-drop,hir:step-one-admissible-top,hir:step-one-usc,hir:step-one-admissible,hir:step-one-canonical-test,hir:step-one-empty-prefix,hir:step-one-representative,hir:step-one-restrict-open,hir:step-one-tower-iso,hir:step-one-bridge,hir:step-one-well-founded,hir:step-one-phase,hir:centre-function-seq,hir:step-one-seq,hir:invariant-well-founded,mc:step-one-phase-bound,hir:max-order-value,hir:max-order-value-basic,bcalc:growing-seq-tail}
\source{wlodarczyk05-hironaka,bierstone-milman08-functoriality}
\lean{TranscendentalNakamaye.Pending.MaxOrderSeqFinite.seqLocallyFiniteMaxOrder_stepOneInvM,TranscendentalNakamaye.Pending.MaxOrderSeqFinite.mem_centreValues_step,TranscendentalNakamaye.Pending.MaxOrderSeqFinite.eq_of_mem_maxLocus,TranscendentalNakamaye.Pending.MaxOrderSeqFinite.maximal_order_step,TranscendentalNakamaye.Pending.MaxOrderSeqFinite.maximal_order_step',TranscendentalNakamaye.Pending.MaxOrderSeqFinite.isInvDomain_zero_of_eq_zero,TranscendentalNakamaye.Pending.MaxOrderSeqFinite.cosupp_eq_empty_of_eq_zero,TranscendentalNakamaye.Pending.MaxOrderSeqFinite.countCentresMeeting_eq_zero_of_eq_zero,TranscendentalNakamaye.Pending.MaxOrderSeqFinite.isMaxOrderValue_of_eq_zero,TranscendentalNakamaye.Pending.MaxOrderValue.isMaxOrderValue_stepOneInvM}
\leanok
Assume a resolution invariant in dimension \(\le k-1\). Then there is an invariant with the properties of Definition~\ref{hir:centre-function} for the marked
ideals of maximal order along \(N\) with \(\dim N=k\) (with values in \(\mathbf N\times(\Lambda\cup\{\infty\})\), ordered lexicographically), property (8) in the
per-sequence form (8\('\)) of Definition~\ref{hir:centre-function-seq}.
\end{proposition}
\begin{proof}
\leanok
Let \(\underline{\mathcal J}\) be of maximal order along \(N\), \(\dim N=k\). Replace it by \(\overline{\underline{\mathcal J}}=\mathcal C(\mathcal H(\underline{\mathcal J}))\), which is
equivalent to it and of maximal order (Lemma~\ref{hir:homogenized-coefficient}, along \(N\) by Lemma~\ref{hir:along-calculus}); the multiple test blow-ups and
resolutions of the two coincide (Lemma~\ref{hir:marked-equivalence-basic}). For \(x\) in a cosupport along a canonical segment write \(s(x)\) for the number of
members of the original \(E\) through \(x\) (strict transforms).

\emph{Step 1a (separating the cosupport from \(E\)).} Let \(s=\max s\) on the cosupport. Points of an intersection \(H^s_\alpha\cap N\) of \(s\) members contained in the
cosupport get the value \((s,\infty)\); these intersections are admissible centres (Lemma~\ref{bcalc:member-intersection}), and blowing them up separates the
members (\(s\ge2\), Lemma~\ref{bcalc:snc-transform-nc}) or, for \(s=1\), is the blowup of a codimension-one component of the cosupport
(Theorem~\ref{bcalc:codim-one-components}); in both cases the new points have \(s(z)<s\) (Theorem~\ref{bcalc:boundary-count-drop}, in its local form: near each of its points the centre
is the intersection of the members through that point, which is how Włodarczyk's intersections of irreducible components read for possibly disconnected
members; \(s\) is upper semicontinuous and does not increase, Lemma~\ref{bcalc:boundary-count}). At the other points \(x\) with \(s(x)=s\),
lying on exactly one such intersection \(S=H^s_\alpha\cap N\), the value is \((s,\operatorname{inv}_{\underline{\mathcal S}}(x))\), where
\(\underline{\mathcal S}\) is the restriction of \(\overline{\underline{\mathcal J}}\) to \(S\) (Lemma~\ref{hir:coefficient-restriction}, along \(N\)), of dimension \(k-s\le k-1\),
after deleting from \(E\) the \(s\) members containing \(S\) (Lemma~\ref{hir:drop-containing-members}; Lemma~\ref{hir:coefficient-restriction} is then applied
with \(E_S\) on both sides).
By maximality of \(s\) the intersections meeting the cosupport have disjoint traces on it, the maximum locus is a union of maximum loci of these
lower-dimensional invariants (admissible by the inductive hypothesis (1) and Lemma~\ref{hir:coefficient-restriction}), and the centres lie in or are
disjoint from each \(S\), so Lemma~\ref{hir:coefficient-restriction} applies along the whole segment. By the inductive local finiteness (8) on the finitely
many intersections meeting the compact sets considered, \(s\) drops after finitely many steps; repeat until \(s=0\) on the cosupport. Then the members of
the original \(E\) miss the cosupport and are dropped (Lemma~\ref{hir:drop-members}).

\emph{Step 1b (no boundary).} Points of the codimension-one part \(Z_1\) of the cosupport, a closed submanifold open in the cosupport
(Theorem~\ref{bcalc:codim-one-part}), get the value \((0,\infty)\); its blowup removes them from the cosupport, and the rest of the cosupport has codimension
at least two (Corollaries~\ref{bcalc:codim-one-step} and~\ref{bcalc:codim-one-step-rest}). At the remaining points, choose a tangent direction \(w\) of the current transform on an open set \(U_w\) (it exists by
Lemma~\ref{mc:giraud-tower}(1) from a tangent direction of \(\overline{\underline{\mathcal J}}\), and \(V(w)\) is the strict transform of the original hypersurface of
maximal contact); by Lemma~\ref{hir:contact-equivalence} the transform on \(U_w\) is equivalent to the marked ideal \(\underline{\mathcal S}_w\) with submanifold
\(V(w)\cap N\cap U_w\), of dimension \(k-1\), and the value at \(x\) is \((0,\operatorname{inv}_{\underline{\mathcal S}_w}(x))\) along the part of the segment after the
start of Step 1b. It does not depend on \(w\): for two tangent directions \(w,w'\) at \(x\), the glueing lemma (Lemma~\ref{mc:glueing-lemma}, applied to the
homogenized ideal at the start of Step 1, and lifted along the segment by Theorem~\ref{mc:glueing-lift}) gives a
biholomorphism \(\varphi\) between open neighbourhoods \(U'\), \(U\) of \(x\) whose stage-by-stage lifts form an isomorphism of the restricted segments over
\(\varphi\) (Definition~\ref{hir:tower-isomorphism}) with \(\varphi_T^*\underline{\mathcal S}_{w,T}=\underline{\mathcal S}_{w',T}\) at every stage
(Lemma~\ref{mc:contact-independence-along}) and which fix the cosupport pointwise; since \(\underline{\mathcal S}_w\) is formed at the stage \(T_0\) where
Step 1b starts and its invariant is computed along the tail \(T_1\) of the segment, the form used is the isomorphism of the restricted tails over the lift
\(\varphi_{T_0}\) (Lemma~\ref{mc:glueing-lift-tail}); properties (5) and (6) of the inductive invariant give (Lemma~\ref{hir:step-one-b-independence}, proved)
\(\operatorname{inv}_{\underline{\mathcal S}_{w'}}(x)=\operatorname{inv}_{\underline{\mathcal S}_w}(\varphi(x))=\operatorname{inv}_{\underline{\mathcal S}_w}(x)\).

\emph{The values as functions of the segment.} For a point \(x\) of the cosupport at the last stage of a canonical segment \(T\), let \(T=T_0T_1\) be the split
(Lemma~\ref{hir:tower-split}) at
the last stage at which the images of \(x\) changed their number \(s\) of original members (Step 1a), resp.\ left \(Z_1\) or reached \(s=0\) (Step 1b). The
lower-dimensional invariant is that of the marked ideal formed at the last stage of \(T_0\) (\(\underline{\mathcal S}\), resp.\ \(\underline{\mathcal S}_w\)), restricted to a
neighbourhood \(W\) of the image of \(x\) meeting no other original member, along \(T_1|_W\), which is a canonical segment of it: near \(x\) the centres of \(T_1\) are
empty or the maximum loci of the second component. For a nearby point whose split is earlier, the steps between the two splits are empty near it, and
Definition~\ref{hir:centre-function}(4') of the inductive invariant re-bases its value to the later split; this is how (2), (5) and (6) compare nearby points
(the source re-bases in the same way across the isomorphisms of an extension, in its lemma on étale morphisms). The stage of the split is the start of the
run of the point in its regime (its number \(s\) of original members, and whether it lies in the part blown up first), Lemma~\ref{hir:regime-run}. In Step 1a the
original members not through the point are kept in \(E\), masked by the restriction to \(W\) (their traces on \(W\) are empty); in Step 1b all original members
are kept and masked the same way, \(W\) missing them (Definition~\ref{hir:step-one-value-masked}), which is the source's set-level \((E_p\setminus E)_{|V(u)_p}\). In both cases
the divisor of the lower marked ideal is the divisor of the transform minus the members through the point, restricted to \(W\), so re-basing across steps with empty
centres appends the same empty exceptional members on both sides and does not shift the positions of the members; deleting the \(\#E\) original members in Step 1b
would shift them, since \(\#E\) itself grows along a re-based prefix.

\emph{Adaptations of the source.} The lower marked ideals are restrictions of the coefficient ideal \(\mathcal C_{N_{T_0}}(\overline{\underline{\mathcal J}}_{T_0})\) to
\(N_{T_0}\cap H_S\), resp.\ \((N\cap V(w))_{T_0}\), where the source restricts \(\overline{\underline{\mathcal J}}_{p}\) itself; the two have the same multiple test
blow-ups and cosupports by Włodarczyk's Lemma ``le:S'' (Lemma~\ref{hir:coefficient-restriction}), and the restriction of the coefficient ideal is compatible with
re-basing. The tangent directions are those of \((\mathfrak a|_N,d)\) along \(N\) (Definition~\ref{mc:tangent-direction-along}), chosen at the base point, where the source
takes \(u\in\mathcal D^{\bar\mu-1}(\overline{\mathcal J})\) at the current stage; the strict transform of \(V(w)\) is the source's hypersurface along the segment
(Lemma~\ref{mc:giraud-tower}). Values are assigned at every stage of a segment, not only ``at the moment the points are getting blown up'', as in Bierstone--Milman's
functorial form; and the intersections of \(s\) members contained in the cosupport are replaced by the member locus, defined near each point.

\emph{Properties.} (1)--(3) hold for the values \((s,\cdot)\) and \((0,\cdot)\) by the inductive hypothesis on the open sets \(U_w\) (resp.\ near the intersections
\(S\)), by Theorem~\ref{bcalc:codim-one-part} for \(Z_1\), and because \(s\) is upper semicontinuous and does not increase along the segment. (4)--(7) hold since
every ingredient (\(\mathcal H\), \(\mathcal C\), \(s\), tangent directions, restriction to \(S\) and to \(V(w)\)) commutes with restriction to open sets and with
biholomorphisms, and is unchanged by empty steps. (4') needs, besides the masked divisor (Definition~\ref{hir:step-one-value-masked}), three comparisons across an
empty prefix \(P\): the run of a point along \(PT_1\) may exceed the length of \(T_1\), so on one side the head of the split is a prefix of \(P\) and on the other the
base of \(\underline{\mathcal J}_P\), and the two lower values agree by Definition~\ref{hir:centre-function}(4') of the inductive invariant; \(\mathcal C_N(\mathcal H_N(\underline{\mathcal J}_P))\)
and \(\overline{\underline{\mathcal J}}_P\) agree up to restricted equivalence (Definition~\ref{hir:restricted-equivalence}), absorbed by (7) of the inductive invariant;
and the tangent directions chosen at the base points in \(X\) and in the last stage of \(P\) differ, absorbed by Lemma~\ref{hir:step-one-b-independence} on the masks. (8): cover a compact set by finitely many relatively compact \(K_j\subset U_{w_j}\) (or neighbourhoods of the
intersections \(S\)); the restricted segment over \(U_{w_j}\) is a canonical segment of \(\underline{\mathcal S}_{w_j}\), which has at most \(m_j\) nonempty centres meeting
\(K_j\) (the counts add up over the \(K_j\) and do not change under restriction, Lemma~\ref{hir:canonical-segments}(3)), and Steps 1aa and 1ba contribute one step per value of \(s\).
Inputs inside \(N\) (proved): for Step 1a, the admissibility of \(N\cap H_S\) (Lemma~\ref{bcalc:member-intersection-along}), the drop of \(s\) at the points of
\(N'\) over the centre (Theorem~\ref{bcalc:boundary-count-drop-along}), and the transversality of \(E_S\) to \(N\cap H_S\) with the well formedness of the
restriction (Lemma~\ref{bcalc:member-intersection-drop-along}); for Step 1b, the codimension-one part of the cosupport along \(N\) and its blowup
(Lemmas~\ref{bcalc:codim-one-along} and~\ref{bcalc:codim-one-local-along}, Theorem~\ref{bcalc:codim-one-part-along}, Definition~\ref{bcalc:codim-one-part-along-def},
Corollary~\ref{bcalc:codim-one-step-along}; the theorems cited above are their forms for \(N=M\)), and the tower half of Lemma~\ref{hir:along-calculus}(3),
Theorem~\ref{hir:giraud-tower-along}.
Proved form: for a lower invariant in the sense of Definition~\ref{hir:centre-function-seq}, properties (1)--(7) and (4') are Corollary~\ref{hir:step-one-seq};
when the values of the lower invariant are well-ordered, so are those of \(\overline{\operatorname{inv}}\), and (8\('\)) is Lemma~\ref{hir:invariant-well-founded}.
For a lower invariant with values in an arbitrary totally ordered set \(\Lambda\), as stated, the properties (1)--(7) and (4') are those of
Lemma~\ref{hir:max-order-value-basic}(1), and (8\('\)) is assembled over the phases (the Lean form). Along a growing sequence \((T_i)\) of canonical segments the values
\(v(T_i)\) of the nonempty centres form one strictly decreasing sequence growing by prefixes (Theorem~\ref{hir:step-one-phase}). If their lengths are bounded, so are the
numbers of centres meeting a compact set. Otherwise let \(s_0\) be the least level of an entry. The levels do not increase, so from the first entry of level \(s_0\) on
every entry has level \(s_0\), and at most one of them is the special value \((s_0,\infty)\), so some entry of level \(s_0\) has a finite second entry. Let \(a\) be the stage at
which the first such centre is blown up. The points of a centre all have the value of its entry, and the value \((s,\infty)\) is taken exactly at the special points of level
\(s\). So every centre before \(a\) has level \(>s_0\) or is special of level \(s_0\), and every centre after \(a\) has level \(\ge s_0\). At \(a\) every point of the cosupport
has a value at most the maximum \((s_0,\lambda)\), \(\lambda<\infty\), so its points of level \(\ge s_0\) are in the regime of level \(s_0\). Thus a phase of level \(s_0\)
starts at \(a\) and never ends, and Theorem~\ref{mc:step-one-phase-bound} bounds the counts by property (8) of the lower invariant (Lemma~\ref{bcalc:growing-seq-tail}).
For \(n=0\) every cosupport of the domain is empty, every centre of a canonical segment is empty, and every value function has the properties.
\end{proof}

\begin{definition}[The value function of the general step]
\label{hir:general-value}
\uses{hir:max-order-value,bcalc:nonmonomial-order-local,bcalc:nonmonomial-order-regime,bcalc:nonmonomial-order-tail,bcalc:companion-rebase,hir:monomial-invariant-along,hir:tower-split,hir:regime-run,hir:centre-function}
\source{wlodarczyk05-hironaka,bierstone-milman08-functoriality}
\lean{TranscendentalNakamaye.Pending.GeneralValue.compAt,TranscendentalNakamaye.Pending.GeneralValue.genLowerValue,TranscendentalNakamaye.Pending.GeneralValue.genSecond,TranscendentalNakamaye.Pending.GeneralValue.genInv,TranscendentalNakamaye.Pending.GeneralValue.GenLowerData}
\leanok
Let \(M\) be a value function (Definition~\ref{hir:max-order-value}, with values in a totally ordered set \(\Lambda\)), \(\underline{\mathcal I}=(X,N,E,\mathfrak a,d)\)
a marked ideal, \(T\) a tower over \(X\) and \(y\) a point of its last stage. Let \(o=\nu(y)\) be the order at \(y\) of the nonmonomial part of the transform
\(\underline{\mathcal I}_T\) along \(N_T\) (Lemma~\ref{bcalc:nonmonomial-order-local}).
\begin{itemize}
\item If \(o=0\), the value of the general step is \(\operatorname{inv}(y)=(0,\gamma(y))\), with \(\gamma\) the invariant of the monomial case of \(\underline{\mathcal I}_T\)
at \(y\) (Lemma~\ref{hir:monomial-invariant-along}).
\item If \(o\ge1\), let \(r\) be the run of \(y\) in the regime ``in the cosupport, of order \(o\)'' (Lemmas~\ref{hir:regime-run} and~\ref{bcalc:nonmonomial-order-regime}) and
\(T=T_0T_1\) the split with \(T_1\) of length \(r\) (Lemma~\ref{hir:tower-split}): the last stage of \(T_0\) is where the images of \(y\) reached the order \(o\). Let
\(W=\{\nu\le o\}\), an open set of the last stage of \(T_0\) (Lemma~\ref{bcalc:nonmonomial-order-tail}), and \(\mathcal O=\mathcal O(\underline{\mathcal I}_{T_0}|_W,o)\) the
companion ideal (Lemma~\ref{bcalc:companion-rebase}). Then \(\operatorname{inv}(y)=(o,M_{\mathcal O}(T_1|_W)(y))\) if the image of \(y\) at \(T_0\) lies in \(W\), and \((o,\infty)\)
otherwise.
\end{itemize}
The values are ordered lexicographically; the second entries lie in the ordered sum of the values of \(\gamma\) followed by \(\Lambda\cup\{\infty\}\) (the first entry
decides between the two cases). The data at \(y\) are \emph{lower data} if the image of \(y\) lies in \(W\), \(\mathcal O\) lies in the domain in dimension \(\le k\) and is
of maximal order along its submanifold, \(T_1|_W\) is a canonical segment of \(\mathcal O\) for \(M\), and \(y\) lies in the cosupport of \(\mathcal O_{T_1|_W}\).
(The source's first entry is \(o/d\); \(d\) is constant along the transforms and under restriction, isomorphisms and restricted equivalence, so the order is the same.)
\end{definition}

\begin{lemma}[The general value: first entry, finiteness, one blowup]
\label{hir:general-value-basic}
\uses{hir:general-value,hir:max-order-value-basic,bcalc:nonmonomial-order-canonical,bcalc:nonmonomial-order-tail,bcalc:nonmonomial-order-regime,bcalc:nonmonomial-order-local,hir:regime-run}
\source{wlodarczyk05-hironaka,bierstone-milman08-functoriality}
\lean{TranscendentalNakamaye.Pending.GeneralValue.genInv_fst,TranscendentalNakamaye.Pending.GeneralValue.ordMonotone_genInv,TranscendentalNakamaye.Pending.GeneralValue.genInv_lt_of_nmOrd_lt,TranscendentalNakamaye.Pending.GeneralValue.genInv_of_nmOrd_eq_zero,TranscendentalNakamaye.Pending.GeneralValue.genInv_of_nmOrd_ne_zero,TranscendentalNakamaye.Pending.GeneralValue.genLowerValue_eq_of_mem,TranscendentalNakamaye.Pending.GeneralValue.genLowerValue_ne_top,TranscendentalNakamaye.Pending.GeneralValue.genInv_lt_top_of_genLowerData,TranscendentalNakamaye.Pending.GeneralValue.genLowerValue_step_of_notMem,TranscendentalNakamaye.Pending.GeneralValue.genLowerValue_step_lt_of_mem,TranscendentalNakamaye.Pending.GeneralValue.genInv_step_eq_of_regime,TranscendentalNakamaye.Pending.GeneralValue.genInv_step_of_notMem,TranscendentalNakamaye.Pending.GeneralValue.genInv_step_lt_of_mem}
\leanok
Let \(M\) be a value function for the marked ideals of maximal order in dimension \(\le k\) and \(\operatorname{inv}\) the general value (Definition~\ref{hir:general-value}).
\begin{enumerate}
\item The first entry of \(\operatorname{inv}\) is \(\nu\); so \(\operatorname{inv}(y)\le\operatorname{inv}(y')\) implies \(\nu(y)\le\nu(y')\), and a point of smaller order has a
smaller value.
\item Along a multiple test blow-up of a marked ideal of the domain (\(n\ge1\)), at a point of the cosupport the companion value is finite: the image of the point at the
start of its run lies in \(W\). Under lower data the value at a point of order \(o\ge1\) is below \((o,\infty)\).
\item Let \(T\) be followed by a blowup \(\pi\) along \(C\) and \(z\) a point of the new cosupport such that \(\pi z\) has order \(o\ge1\) and \(z\) has lower data at the split of
\(T\) at the start of the run of \(\pi z\), followed by \(\pi\). If \(\pi z\notin C\) (\(T\pi\) a multiple test blow-up of a marked ideal of the domain, \(n\ge1\)), then
\(\operatorname{inv}(z)=\operatorname{inv}(\pi z)\). If \(\pi z\in C\), \(z\) has order \(o\) and \(\pi z\) is in the regime of order \(o\), then
\(\operatorname{inv}(z)<\operatorname{inv}(\pi z)\).
\end{enumerate}
\end{lemma}
\begin{proof}
\leanok
(1) By the definition of the lexicographic order. (2) The image of a point of the cosupport at the start of its run lies in the regime and in \(W\)
(Lemma~\ref{bcalc:nonmonomial-order-tail}). (3) Off an admissible centre the order of the nonmonomial part and the regime do not change
(Lemmas~\ref{bcalc:nonmonomial-order-local} and~\ref{bcalc:nonmonomial-order-regime}), so the run of \(z\) is that of \(\pi z\) plus one and the split of \(T\pi\) at the start of
the run of \(z\) is the split of \(T\) followed by \(\pi\) (Lemma~\ref{hir:regime-run}); the companion value is computed from the same companion ideal along the restricted
tail followed by the restriction of \(\pi\), whose centre is the trace of \(C\). Under lower data, Lemma~\ref{hir:max-order-value-basic}(2) for \(M\) gives equality off the
centre and a drop over it.
\end{proof}

\begin{lemma}[The general value of a nearby point through the reference point's companion ideal]
\label{hir:general-compare}
\uses{hir:general-value,hir:general-value-basic,hir:max-order-value,hir:max-order-value-basic,bcalc:companion-rebase,bcalc:companion-open,bcalc:companion-domain,bcalc:restrict-nested,mc:tower-isomorphism-restrict-append,bcalc:restrict-empty-tower,bcalc:nonmonomial-order-tail,bcalc:nonmonomial-order-local,hir:step-one-nearby}
\source{wlodarczyk05-hironaka,bierstone-milman08-functoriality}
\lean{TranscendentalNakamaye.Pending.GeneralCompare.value_restrict_nestedOpens,TranscendentalNakamaye.Pending.GeneralCompare.value_eq_of_two_opens,TranscendentalNakamaye.Pending.GeneralCompare.value_restrict_append,TranscendentalNakamaye.Pending.GeneralCompare.compAt_eq_pullback,TranscendentalNakamaye.Pending.GeneralCompare.isInvDomain_compAt,TranscendentalNakamaye.Pending.GeneralCompare.isMaximalOrderAlong_compAt,TranscendentalNakamaye.Pending.GeneralCompare.genLowerValue_eq_rebase}
\leanok
Let \(M\) be a value function for the marked ideals of maximal order in dimension \(\le k\) (Definition~\ref{hir:max-order-value}).
\begin{enumerate}
\item (Nested and two open sets.) Let \(\underline{\mathcal N}\) have a coherent ideal and a closed submanifold, and \(A\), \(B\) open sets of the base of a tower \(T\)
such that \(\underline{\mathcal N}|_A\), \(\underline{\mathcal N}|_B\) lie in the domain, are of maximal order, and \(T|_A\), \(T|_B\) are canonical segments of them. At a point
over \(A\cap B\) in the cosupport along \(T|_A\) the two values of \(M\) agree, and the point lies in the cosupport along \(T|_B\). For \(V\subset W\) the data on \(W\) restrict
to \(V\), with the same values.
\item (Restricted concatenation.) If \((M'T_1)|_W\) is a canonical segment of a marked ideal on \(W\) of the domain, of maximal order, then so is \(M'|_W\) followed by
\(T_1|_{M'^{-1}W}\), with the same cosupport and values.
\item The companion ideal formed at a stage \(T_0\) on its mask \(W=\{\nu\le o\}\) (\(o\ge1\)) is the restriction to \(W\) of the companion ideal of
\(\underline{\mathcal I}_{T_0}\); it lies in the domain and is of maximal order.
\item (Comparison.) Let \(T_0'M'\) be a multiple test blow-up of \(\underline{\mathcal I}\) (of the domain, \(n\ge1\)), \(o\ge1\), the centres of \(M'\) missing the mask
\(W'\) of \(T_0'\), and \(p\) a point of the last stage of \(T_0'M'T_1\) with lower data at the split \(T_0'(M'T_1)\) (Definition~\ref{hir:general-value}). If the image
of \(p\) at \(T_0=T_0'M'\) lies in the mask \(W\) of \(T_0\) and \(T_1|_W\) is a canonical segment of the companion ideal \(\mathcal O_{T_0}\) formed at \(T_0\), then the
companion value of \(p\) is \(M_{\mathcal O_{T_0}}(T_1|_W)(p)\), and \(p\) lies in the cosupport of \(\mathcal O_{T_0}\) along \(T_1|_W\).
\end{enumerate}
\end{lemma}
\begin{proof}
\leanok
(1) and (2) are the proofs of the corresponding statements for resolution invariants (Lemma~\ref{hir:step-one-nearby} and the comparison of
Lemma~\ref{hir:step-one-compare}): restriction to \(W\cap V\subset W\) and the isomorphism of towers onto the restriction to \(V\)
(Lemma~\ref{bcalc:restrict-nested}), resp.\ the isomorphism of Lemma~\ref{mc:tower-isomorphism-restrict-append}, with properties (5) and (6) of \(M\) and the
canonical segments of Lemma~\ref{hir:max-order-value-basic}; maximal order passes to restrictions and along biholomorphisms.
(3) The companion ideal commutes with restriction to open sets (Lemma~\ref{bcalc:companion-rebase}); it lies in the domain (Lemma~\ref{bcalc:companion-domain}) and is of
maximal order where \(\nu\le o\) on the cosupport (Lemma~\ref{bcalc:companion-open}), which is the mask (Lemma~\ref{bcalc:nonmonomial-order-tail}).
(4) By (2), the companion value of \(p\) is the value along \(M'|_{W'}\) followed by \(T_1|_{M'^{-1}W'}\). The first part has only empty centres
(Lemma~\ref{bcalc:restrict-empty-tower}), and the companion ideal formed at \(T_0'\) on \(W'\), transformed along it, is restricted-equivalent to the companion ideal of
\(\underline{\mathcal I}_{T_0}\) restricted to \(M'^{-1}W'\) (Lemma~\ref{bcalc:companion-rebase}); the latter is of maximal order since the order is \(\le o\) over
\(W'\) and does not change over \(W'\) (Lemma~\ref{bcalc:nonmonomial-order-local}). Re-basing the values (Lemma~\ref{hir:max-order-value-basic}(5)) gives the value of
\(M\) for that companion ideal along \(T_1|_{M'^{-1}W'}\), and (1) for the two open sets \(M'^{-1}W'\) and \(W\) of the last stage of \(T_0\), with (3), moves it to
\(\mathcal O_{T_0}\) along \(T_1|_W\).
\end{proof}

\begin{theorem}[The canonical-tail theorem of the general step; the drop of the general value]
\label{hir:general-canonical-tail}
\uses{hir:general-value,hir:general-value-basic,hir:general-compare,hir:max-order-value-basic,bcalc:nonmonomial-order-tail,bcalc:nonmonomial-order-regime,bcalc:nonmonomial-order-canonical,bcalc:nonmonomial-order-local,bcalc:nonmonomial-order-monomial-canonical,bcalc:companion-rebase,hir:tower-split,hir:regime-run,hir:step-one-nearby,hir:step-one-canonical-tail}
\source{wlodarczyk05-hironaka,bierstone-milman08-functoriality}
\lean{TranscendentalNakamaye.Pending.GeneralCanonicalTail.mem_cosupp_compAt_iff,TranscendentalNakamaye.Pending.GeneralCanonicalTail.genLowerData_refl,TranscendentalNakamaye.Pending.GeneralCanonicalTail.exists_middle_ord,TranscendentalNakamaye.Pending.GeneralCanonicalTail.genLowerValue_eq_rebase_of_eq,TranscendentalNakamaye.Pending.GeneralCanonicalTail.HasGenLowerData,TranscendentalNakamaye.Pending.GeneralCanonicalTail.genInv_eq_of_mem_cosupp_compAt,TranscendentalNakamaye.Pending.GeneralCanonicalTail.genLowerData_step,TranscendentalNakamaye.Pending.GeneralCanonicalTail.hasGenLowerData_of_isCanonical,TranscendentalNakamaye.Pending.GeneralCanonicalTail.ordZeroValue_genInv,TranscendentalNakamaye.Pending.GeneralCanonicalTail.genInv_drop}
\leanok
Let \(\underline{\mathcal I}\) lie in the domain in dimension \(\le k\) (\(n\ge1\)), \(M\) a value function for the marked ideals of maximal order in dimension
\(\le k\) and \(\operatorname{inv}\) the general value (Definition~\ref{hir:general-value}).
\begin{enumerate}
\item Let \(T_0\), \(T_1\) be multiple test blow-ups of \(\underline{\mathcal I}\) and of \(\underline{\mathcal I}_{T_0}\), \(o\ge1\), and \(T_1|_W\) one of the companion ideal formed
at \(T_0\) on its mask \(W\). A point over \(W\) lies in the cosupport of the companion's transform if and only if it lies in \(\operatorname{cosupp}\underline{\mathcal I}_{T_0T_1}\)
and has order \(o\).
\item (Canonical tails.) Along every multiple test blow-up \(T\) of \(\underline{\mathcal I}\) which is a canonical segment of \(\operatorname{inv}\), every point of the cosupport of
positive order has lower data at the split of \(T\) at the start of its run: the restricted tail is a canonical segment of the companion ideal formed there.
\item (Drop.) If moreover \(T\) followed by a blowup \(\pi\) along \(C\) is a multiple test blow-up and a canonical segment, every point \(z\) of the new cosupport lies over
the cosupport, with \(\operatorname{inv}(z)<\operatorname{inv}(\pi z)\) if \(\pi z\in C\) and \(\operatorname{inv}(z)=\operatorname{inv}(\pi z)\) otherwise: properties (3) and (4)
of Definition~\ref{hir:centre-function} for the general value.
\end{enumerate}
\end{theorem}
\begin{proof}
\leanok
(1) The companion's cosupport along a multiple test blow-up of it is the set of points of order \(\ge o\) in the cosupport of the transform of the restricted marked ideal
(Lemma~\ref{bcalc:nonmonomial-order-tail}), and over \(W\) the order stays \(\le o\) along such a tower (Lemma~\ref{bcalc:nonmonomial-order-regime}); restriction to open sets
commutes with the transforms, the cosupports and the order.
(2) Induction on the tower. For the empty tower, and for a point whose image leaves its regime, the run is empty and the companion ideal formed at the last stage gives
lower data along the empty tail: the point lies in \(W\) (its order is \(o\)), the companion ideal lies in the domain and is of maximal order
(Lemma~\ref{hir:general-compare}(3)), and the point lies in its cosupport by (1). Let \(y\) be a point after a further blowup \(\pi\) along \(C\) whose image \(\pi y\) keeps its
order \(o\ge1\); the split of \(y\) is that of \(\pi y\), with the tail extended by \(\pi\) (Lemma~\ref{hir:regime-run}), and \(\pi y\) has lower data by induction. The extended
restricted tail is a multiple test blow-up of the companion ideal (Lemma~\ref{bcalc:nonmonomial-order-tail}), so \(y\) lies in the companion's cosupport by (1). If \(C\ne\emptyset\),
it is the maximum locus of \(\operatorname{inv}\); its trace on the restricted tail is empty or the maximum locus of \(M\) for the companion ideal: on the companion's cosupport
the value of \(\operatorname{inv}\) is \((o,M(q))\) and the other points over \(W\) have order \(<o\), hence smaller value than \(\pi y\). For the value at a point \(q\) of the
companion's cosupport: \(q\) has order \(o\), its run started no later than that of \(\pi y\) (Lemma~\ref{bcalc:nonmonomial-order-regime}), the steps between the two starts
have centres of order \(\ge o+1\) and miss the mask of the earlier start (Lemma~\ref{bcalc:nonmonomial-order-canonical}), and Lemma~\ref{hir:general-compare}(4) moves the
value of \(q\), computed from its lower data (by induction), to the companion ideal formed at the start of \(\pi y\).
(3) The order does not go up along a canonical segment (Lemma~\ref{bcalc:nonmonomial-order-canonical}). At a point whose image has order \(0\) this is
Lemma~\ref{bcalc:nonmonomial-order-monomial-canonical}. Otherwise, if the order drops, so does the value; if not, the image is in the regime of its order, and with the lower data
of (2) at the extended split, Lemma~\ref{hir:general-value-basic}(3) gives the drop over the centre and the equality off it (where the order does not change,
Lemma~\ref{bcalc:nonmonomial-order-local}).
\end{proof}

\begin{lemma}[Upper semicontinuity of the general value]
\label{hir:general-usc}
\uses{hir:general-canonical-tail,hir:general-value,hir:max-order-value,bcalc:nonmonomial-order-canonical,bcalc:nonmonomial-order-tail,bcalc:nonmonomial-order-monomial-canonical,hir:step-one-usc}
\source{wlodarczyk05-hironaka,bierstone-milman08-functoriality}
\lean{TranscendentalNakamaye.Pending.GeneralUSC.genInv_eventually_le_of_pos,TranscendentalNakamaye.Pending.GeneralUSC.genInv_usc}
\leanok
Let \(\underline{\mathcal I}\) lie in the domain in dimension \(\le k\) (\(n\ge1\)), \(M\) a value function for the marked ideals of maximal order in dimension \(\le k\), and
\(T\) a canonical segment of the general value \(\operatorname{inv}\) which is a multiple test blow-up. Near a point \(y\) of the cosupport of order \(o\ge1\), the points of the
cosupport of order \(o\) have values \(\le\operatorname{inv}(y)\); and \(\operatorname{inv}\) is upper semicontinuous on \(\operatorname{cosupp}\underline{\mathcal I}_T\) (property (2) of
Definition~\ref{hir:centre-function} for the general value).
\end{lemma}
\begin{proof}
\leanok
By Theorem~\ref{hir:general-canonical-tail}(2), \(y\) has lower data at the start \(T_0\) of its run: the restricted tail \(T_1|_W\) is a canonical segment of the companion ideal
\(\mathcal O\) formed at \(T_0\) on its mask \(W\). The points of order \(o\) of the cosupport over \(W\) (an open neighbourhood of \(y\)) are the points of the cosupport of
\(\mathcal O\) along \(T_1|_W\) (Theorem~\ref{hir:general-canonical-tail}(1)), where \(\operatorname{inv}=(o,M_{\mathcal O}(T_1|_W))\) (the comparison in the proof of
Theorem~\ref{hir:general-canonical-tail}(2)); \(M_{\mathcal O}\) is upper semicontinuous along its canonical segment (Definition~\ref{hir:max-order-value}, property (2)). At points of
order \(0\) the corresponding statement is Lemma~\ref{bcalc:nonmonomial-order-monomial-canonical}. Nearby points have order at most \(\nu(y)\), and a smaller order gives a smaller
value (Lemma~\ref{bcalc:nonmonomial-order-canonical}), so it suffices to compare points of the same order.
\end{proof}

\begin{theorem}[Admissibility for the general value; its canonical segments are multiple test blow-ups]
\label{hir:general-admissible}
\uses{hir:general-canonical-tail,hir:general-usc,hir:general-value,hir:max-order-value,bcalc:nonmonomial-order-phase,bcalc:nonmonomial-order-monomial-canonical,hir:invariant-driver,hir:step-one-admissible,hir:step-one-well-founded}
\source{wlodarczyk05-hironaka,bierstone-milman08-functoriality}
\lean{TranscendentalNakamaye.Pending.GeneralAdmissible.exists_opens_preimage_eq_empty,TranscendentalNakamaye.Pending.GeneralAdmissible.exists_opens_isAdmissibleCentre_of_pos,TranscendentalNakamaye.Pending.GeneralAdmissible.genInv_admissible_pos,TranscendentalNakamaye.Pending.GeneralAdmissible.genInv_admissible,TranscendentalNakamaye.Pending.GeneralAdmissible.isAdmissibleCentre_maxLocus_genInv,TranscendentalNakamaye.Pending.GeneralAdmissible.isGeoTestTower_of_isCanonical_genInv,TranscendentalNakamaye.Pending.GeneralAdmissible.stepDrop_genInv}
\leanok
Let \(\underline{\mathcal I}\) lie in the domain in dimension \(\le k\) (\(n\ge1\)), \(M\) a value function for the marked ideals of maximal order in dimension \(\le k\), and
\(\operatorname{inv}\) the general value.
\begin{enumerate}
\item Along a canonical segment \(T\) of \(\operatorname{inv}\) which is a multiple test blow-up, on an open set \(W\) of the last stage on which \(\operatorname{inv}\le\lambda\) on the
cosupport, the points of the cosupport over \(W\) with value \(\lambda\) form an admissible centre of \(\underline{\mathcal I}_T|_W\) (property (1) of
Definition~\ref{hir:centre-function} for the general value); in particular the maximum locus is an admissible centre.
\item Every canonical segment of \(\operatorname{inv}\) is a multiple test blow-up, so Theorem~\ref{hir:general-canonical-tail} and Lemma~\ref{hir:general-usc} hold along every
canonical segment, and \(\operatorname{inv}\) has the drop property of Theorem~\ref{hir:step-one-well-founded} along its canonical segments.
\end{enumerate}
\end{theorem}
\begin{proof}
\leanok
(1) Admissible centres are local (Theorem~\ref{hir:invariant-driver}(1)). Near a point off the level set, a neighbourhood misses it: upper semicontinuity
(Lemma~\ref{hir:general-usc}) and the closed cosupport. At a value \((0,\gamma)\) this is Lemma~\ref{bcalc:nonmonomial-order-monomial-canonical}. At a value
\((o,\lambda)\) with \(o\ge1\) and \(\lambda\) finite, near a point \(x\) of the level set take \(W\) intersected with the open set over the mask of the start of \(x\): there the
level set is the level set \(\{M=\lambda\}\) of the companion ideal formed at that start, on the cosupport of its transform (the comparison in the proof of
Theorem~\ref{hir:general-canonical-tail}), admissible for the restriction of the companion's transform by property (1) of \(M\) (with \(M\le\lambda\) there), hence admissible for the
restriction of \(\underline{\mathcal I}_T\), the companion's transform having the same divisor and a cosupport inside it (Lemma~\ref{bcalc:nonmonomial-order-phase}). The other
values do not occur: the companion value is finite on the domain (Lemma~\ref{hir:general-value-basic}(2)).
(2) By induction along the segment, each centre being empty or the maximum locus, admissible by (1); then Theorem~\ref{hir:general-canonical-tail}(3) applies at every step.
\end{proof}

\begin{lemma}[The general value does not depend on the representative]
\label{hir:general-representative}
\uses{hir:general-value,hir:general-value-basic,bcalc:nonmonomial-order-transport,bcalc:companion-rebase,bcalc:general-value-order-zero,hir:marked-equivalence-basic,mc:contact-representative}
\source{wlodarczyk05-hironaka,bierstone-milman08-functoriality}
\lean{TranscendentalNakamaye.Pending.GeneralRepresentative.genLowerValue_eq_refl,TranscendentalNakamaye.Pending.GeneralRepresentative.genLowerValue_congr_of_eq,TranscendentalNakamaye.Pending.GeneralRepresentative.genLowerValue_refl_congr,TranscendentalNakamaye.Pending.GeneralRepresentative.genLowerValue_congr,TranscendentalNakamaye.Pending.GeneralRepresentative.genInv_congr_of_pos,TranscendentalNakamaye.Pending.GeneralRepresentative.genInv_congr,TranscendentalNakamaye.Pending.GeneralRepresentative.representative_field}
\leanok
Let \(M\) be any value function for the marked ideals of maximal order and \(\operatorname{inv}\) the general value (Definition~\ref{hir:general-value}).
\begin{enumerate}
\item The companion value at a split \(T_0T_1\) only depends on the transform \(\underline{\mathcal I}_{T_0}\): two splits with equal transforms at their heads have the same
companion value.
\item If \(\underline{\mathcal I}\simeq\underline{\mathcal J}\) are restricted-equivalent (closed \(N\), coherent ideals), then
\(\operatorname{inv}_{\underline{\mathcal I}}(T)(y)=\operatorname{inv}_{\underline{\mathcal J}}(T)(y)\) along every tower \(T\) at every point \(y\) of its last stage. In particular
property (7) of Definition~\ref{hir:centre-function} holds for the general value, along its canonical segments.
\end{enumerate}
\end{lemma}
\begin{proof}
\leanok
(1) The companion value is defined from the transform at the head, its mask and the tail only.
(2) The orders of the nonmonomial parts agree along every tower, hence so do the regimes, their runs and the masks (Lemma~\ref{bcalc:nonmonomial-order-transport}). At a point
of order \(0\) both values are \((0,\gamma)\) of transforms with the same \(\gamma\) (Lemma~\ref{bcalc:general-value-order-zero}). At a point of order \(o\ge1\) both values are
\((o,\cdot)\), the companion values being computed at the same split \(T_0T_1\) (the same run). If \(T_0\) has at least one step, the transforms at \(T_0\) are equal
(Lemma~\ref{hir:marked-equivalence-basic}) and (1) applies. If \(T_0\) is empty, the restrictions of \(\underline{\mathcal I}\) and \(\underline{\mathcal J}\) to the common mask are
restricted-equivalent (Theorem~\ref{mc:contact-representative}), so their companion ideals are equal (Lemma~\ref{bcalc:companion-rebase}), and the companion values agree. No
property of \(M\) and no canonicity is used.
\end{proof}

\begin{lemma}[The general value along isomorphisms of towers]
\label{hir:general-tower-iso}
\uses{hir:general-value,hir:general-value-basic,hir:general-canonical-tail,hir:general-admissible,hir:max-order-value,bcalc:companion-tower-iso,bcalc:nonmonomial-order-transport,bcalc:nonmonomial-order-regime,bcalc:nonmonomial-order-local,bcalc:general-value-order-zero,hir:tower-split-iso,hir:tower-isomorphism}
\source{wlodarczyk05-hironaka,bierstone-milman08-functoriality}
\lean{TranscendentalNakamaye.Pending.GeneralTowerIso.value_companionOf_isTowerIso,TranscendentalNakamaye.Pending.GeneralTowerIso.genInv_isTowerIso_of_pos,TranscendentalNakamaye.Pending.GeneralTowerIso.towerIso_field}
\leanok
Let \(\underline{\mathcal I}\) lie in the domain in dimension \(\le k\) (\(n\ge1\)), \(M\) a value function for the marked ideals of maximal order in dimension \(\le k\) and
\(\operatorname{inv}\) the general value. Let \(\chi\colon T'\to T\) be an isomorphism of towers over a biholomorphism \(\Phi\colon X'\to X\)
(Definition~\ref{hir:tower-isomorphism}) and \(T\) a canonical segment of \(\operatorname{inv}\). Then
\(\operatorname{inv}_{\Phi^*\underline{\mathcal I}}(T')(p)=\operatorname{inv}_{\underline{\mathcal I}}(T)(\chi(p))\) at every \(p\) with \(\chi(p)\in\operatorname{cosupp}\underline{\mathcal I}_T\):
property (6) of Definition~\ref{hir:centre-function} for the general value.
\end{lemma}
\begin{proof}
\leanok
The segment is a multiple test blow-up (Theorem~\ref{hir:general-admissible}). At order \(0\) this is Lemma~\ref{bcalc:general-value-order-zero}. Let \(\chi(p)\) have order \(o\ge1\)
and run \(r\). The orders and the runs at \(p\) and \(\chi(p)\) agree (Lemmas~\ref{bcalc:nonmonomial-order-local}, \ref{bcalc:nonmonomial-order-regime}); the splits of \(T'\) and \(T\)
at \(r\) correspond, the heads being isomorphic over \(\Phi\) by some \(\chi_0\) and the tails over \(\chi_0\) (Lemma~\ref{hir:tower-split-iso}); the masks correspond along
\(\chi_0\) (Lemma~\ref{bcalc:nonmonomial-order-transport}), and the companion ideal formed at the head of \(T'\) on its mask is the pullback, along the restriction of \(\chi_0\),
of the one formed at the head of \(T\) (Lemma~\ref{bcalc:companion-tower-iso}). The restricted tail of \(T\) is a canonical segment of the latter, with \(\chi(p)\) in its
cosupport (Theorem~\ref{hir:general-canonical-tail}), and the restricted tails are isomorphic; property (6) of \(M\) (Definition~\ref{hir:max-order-value}) gives the equality
of the companion values, hence of the values \((o,\cdot)\).
\end{proof}

\begin{lemma}[The general value under restriction to open sets]
\label{hir:general-restrict-open}
\uses{hir:general-tower-iso,hir:general-value,hir:general-value-basic,hir:general-canonical-tail,hir:general-compare,hir:general-admissible,hir:max-order-value,bcalc:restrict-nested,bcalc:nonmonomial-order-transport,bcalc:nonmonomial-order-regime,bcalc:nonmonomial-order-local,bcalc:companion-rebase,bcalc:general-value-order-zero,hir:tower-split-iso,hir:open-restriction-tower,hir:open-restriction-marked,hir:step-one-nearby}
\source{wlodarczyk05-hironaka,bierstone-milman08-functoriality}
\lean{TranscendentalNakamaye.Pending.GeneralRestrictOpen.value_restrict_restrict,TranscendentalNakamaye.Pending.GeneralRestrictOpen.genInv_restrict_of_pos,TranscendentalNakamaye.Pending.GeneralRestrictOpen.restrictOpen_field}
\leanok
Let \(\underline{\mathcal I}\) lie in the domain in dimension \(\le k\) (\(n\ge1\)), \(M\) a value function for the marked ideals of maximal order in dimension \(\le k\) and
\(\operatorname{inv}\) the general value.
\begin{enumerate}
\item For a marked ideal \(\mathcal C\) (closed \(N\), coherent ideal) and opens \(V\), \(W\): if \(\mathcal C|_W\) lies in the domain, is of maximal order and \(T_1|_W\) is a
canonical segment of it for \(M\), then so for \((\mathcal C|_V)|_{W\cap V}\) and \((T_1|_V)|_{W\cap V}\), with the same cosupport and the same values of \(M\) at its points.
\item For \(U\subset X\) open and a canonical segment \(T\) of \(\operatorname{inv}\),
\(\operatorname{inv}_{\underline{\mathcal I}|_U}(T|_U)\) is the restriction of \(\operatorname{inv}_{\underline{\mathcal I}}(T)\) to the cosupport over \(U\): property (5) of
Definition~\ref{hir:centre-function} for the general value.
\end{enumerate}
\end{lemma}
\begin{proof}
\leanok
(1) Lemma~\ref{bcalc:restrict-nested} moves the data from \(W\) to \(W\cap V\); restricting to \(V\) and then to \(W\cap V\) seen in \(V\) is isomorphic over the identity on
points to restricting to \(W\cap V\), and property (6) of \(M\) transports the data along this isomorphism; the cosupports correspond as in Lemma~\ref{hir:step-one-nearby}.
(2) At order \(0\) this is Lemma~\ref{bcalc:general-value-order-zero}. Let \(p\) have order \(o\ge1\) and run \(r\), \(T=T_0T_1\) the split at \(r\); \(p\) has lower data there
(Theorem~\ref{hir:general-canonical-tail}): with the mask \(W\) of \(\underline{\mathcal I}_{T_0}\) and \(\mathcal O\) the companion ideal of \(\underline{\mathcal I}_{T_0}\), \(T_1|_W\) is a
canonical segment of \(\mathcal O|_W\). The orders and runs for \(\underline{\mathcal I}|_U\) along \(T|_U\) are those for \(\underline{\mathcal I}\) along \(T\)
(Lemmas~\ref{bcalc:nonmonomial-order-local}, \ref{bcalc:nonmonomial-order-regime}), and the split of \(T|_U\) at \(r\) is isomorphic over the identity of \(U\) to the restriction of
the split of \(T\) (Lemma~\ref{hir:tower-split-iso}). By Lemma~\ref{hir:general-tower-iso} (its first step, with \(\Phi=\mathrm{id}_U\)) the companion value of \(p\) for
\(\underline{\mathcal I}|_U\) is the value of \(M\) for the companion ideal of \((\underline{\mathcal I}_{T_0})|_V\), \(V=T_0^{-1}(U)\), on its mask, along \((T_1|_V)\) restricted there; this
companion ideal is \(\mathcal O|_V\) (Lemma~\ref{bcalc:companion-rebase}) and its mask is the trace of \(W\) (Lemma~\ref{bcalc:nonmonomial-order-transport}). By (1) this value is that
of \(M\) for \(\mathcal O|_W\) along \(T_1|_W\) at \(p\), the companion value of \(p\) for \(\underline{\mathcal I}\).
\end{proof}

\begin{lemma}[Re-basing the general value after an empty prefix]
\label{hir:general-empty-prefix}
\uses{hir:general-value,hir:general-value-basic,hir:general-canonical-tail,hir:general-compare,hir:general-representative,hir:general-admissible,hir:max-order-value,hir:max-order-value-basic,bcalc:companion-rebase,bcalc:nonmonomial-order-transport,bcalc:restrict-empty-tower,bcalc:general-value-order-zero,hir:step-one-empty-prefix,hir:tower-split}
\source{wlodarczyk05-hironaka,bierstone-milman08-functoriality}
\lean{TranscendentalNakamaye.Pending.GeneralEmptyPrefix.genLowerValue_refl_append_of_isEmptyTower,TranscendentalNakamaye.Pending.GeneralEmptyPrefix.genInv_emptyPrefix_of_pos,TranscendentalNakamaye.Pending.GeneralEmptyPrefix.emptyPrefix_field}
\leanok
Let \(\underline{\mathcal I}\) lie in the domain in dimension \(\le k\) (\(n\ge1\)), \(M\) a value function for the marked ideals of maximal order in dimension \(\le k\) and
\(\operatorname{inv}\) the general value. If all centres of \(P\) are empty and \(PT_1\) is a canonical segment of \(\operatorname{inv}\), then
\(\operatorname{inv}_{\underline{\mathcal I}}(PT_1)=\operatorname{inv}_{\underline{\mathcal I}_P}(T_1)\) on \(\operatorname{cosupp}\underline{\mathcal I}_{PT_1}\): property (4') of
Definition~\ref{hir:centre-function} for the general value.
\end{lemma}
\begin{proof}
\leanok
At order \(0\) this is Lemma~\ref{bcalc:general-value-order-zero}. Let \(z\) have order \(o\ge1\). The orders agree, and the regimes along \(PH\) and \(H\) correspond
(Lemma~\ref{bcalc:nonmonomial-order-transport}). If the run of \(z\) along \(T_1\) is shorter than \(T_1\), it is its run along \(PT_1\), and the split of \(PT_1\) at it is
\(P\) followed by the split of \(T_1\) (Lemmas~\ref{hir:step-one-empty-prefix}, \ref{hir:tower-split}); both companion values are computed from the same transform
(Lemma~\ref{hir:general-representative}(1)). Otherwise the run covers \(T_1\) and then \(P\) (whose steps are empty), and the companion ideals are formed at the base, for
\(\underline{\mathcal I}\) along \(PT_1\) and for \(\underline{\mathcal I}_P\) along \(T_1\). The point has lower data at the first split (Theorem~\ref{hir:general-canonical-tail}); the
restriction of \(PT_1\) to the mask \(W\) is \(P|_W\) followed by the restriction of \(T_1\) (Lemma~\ref{hir:general-compare}), \(P|_W\) has only empty centres
(Lemma~\ref{bcalc:restrict-empty-tower}), the companion ideal of \(\underline{\mathcal I}_P\) on \(P^{-1}W\) is restricted-equivalent to the transform of the companion of
\(\underline{\mathcal I}\) along \(P|_W\) (Lemma~\ref{bcalc:companion-rebase}), and re-basing \(M\) across \(P|_W\) (Lemma~\ref{hir:max-order-value-basic}) gives the equality;
\(P^{-1}W\) is exactly the mask of \(\underline{\mathcal I}_P\) (Lemma~\ref{bcalc:nonmonomial-order-transport}), so no canonicity at the later stage is needed.
\end{proof}

\begin{theorem}[The general value has the properties of a resolution invariant except local finiteness]
\label{hir:general-step-value}
\uses{hir:centre-function-seq,hir:general-admissible,hir:general-usc,hir:general-canonical-tail,hir:general-empty-prefix,hir:general-restrict-open,hir:general-tower-iso,hir:general-representative,hir:max-order-value}
\source{wlodarczyk05-hironaka,bierstone-milman08-functoriality}
\lean{TranscendentalNakamaye.Pending.GeneralStepValue.isInvValue_genInv}
\leanok
Let \(n\ge1\) and \(M\) be a value function for the marked ideals of maximal order in dimension \(\le k\) (Definition~\ref{hir:max-order-value}). Then the general value
\(\operatorname{inv}\) has properties (1)--(7) and (4') of Definition~\ref{hir:centre-function} for the marked ideals of the domain in dimension \(\le k\), along their canonical
segments; with (8\('\)) it is a resolution invariant in the sense of Definition~\ref{hir:centre-function-seq}.
\end{theorem}
\begin{proof}
\leanok
Canonical segments are multiple test blow-ups (Theorem~\ref{hir:general-admissible}(2)). (1) Theorem~\ref{hir:general-admissible}(1); (2) Lemma~\ref{hir:general-usc}; (3) and
(4) Theorem~\ref{hir:general-canonical-tail}(3), a step along an empty centre or the maximum locus extending the canonical segment; (4')
Lemma~\ref{hir:general-empty-prefix}; (5) Lemma~\ref{hir:general-restrict-open}; (6) Lemma~\ref{hir:general-tower-iso}; (7) Lemma~\ref{hir:general-representative}.
\end{proof}

\begin{theorem}[Local finiteness of the general value; the general value is a resolution invariant]
\label{hir:general-seq-finite}
\uses{hir:general-step-value,hir:centre-function-seq,hir:max-order-value,bcalc:general-companion-phase,bcalc:general-phase-count,bcalc:nonmonomial-order-global-phase}
\source{wlodarczyk05-hironaka,bierstone-milman08-functoriality}
\lean{TranscendentalNakamaye.Pending.GeneralSeqFinite.SeqLocallyFiniteMaxOrder,TranscendentalNakamaye.Pending.GeneralSeqFinite.seqLocallyFinite_genInv,TranscendentalNakamaye.Pending.GeneralSeqFinite.ofMaxOrder,TranscendentalNakamaye.Pending.GeneralSeqFinite.ofMaxOrder_inv}
\leanok
Let \(n\ge1\) and let \(M\) be a value function for the marked ideals of maximal order in dimension \(\le k\) (Definition~\ref{hir:max-order-value}) which has property
(8\('\)) of Definition~\ref{hir:centre-function-seq} for the marked ideals of the domain in dimension \(\le k\) of maximal order along their submanifold. Then the general
value \(\operatorname{inv}\) has property (8\('\)) for every marked ideal of the domain in dimension \(\le k\); with Theorem~\ref{hir:general-step-value} it is a resolution
invariant in dimension \(\le k\) in the sense of Definition~\ref{hir:centre-function-seq}.
\end{theorem}
\begin{proof}
\leanok
Let \((T_i)\) be a growing sequence of canonical segments of the general value of \(\underline{\mathcal I}\) and \(K\) compact. The points of the nonempty centres meeting the
preimage of \(K\) have order \(\le q\) for one \(q\) (Lemma~\ref{bcalc:nonmonomial-order-global-phase}), so the number of centres of \(T_i\) meeting the preimage of \(K\) is
at most the sum over \(o=0,\dots,q\) of the numbers \(c_o(K,T_i)\) of centres having a point over \(K\) of order \(o\) (Lemma~\ref{bcalc:general-companion-phase}). For
\(o=0\), \(c_0(K,T_i)\) is at most the number of values of nonempty centres of order \(0\), which is bounded along the sequence (Lemma~\ref{bcalc:general-phase-count}: the
monomial phase is finite). For \(o\ge1\), \(c_o(K,T_i)\) is bounded along the sequence by (8\('\)) of \(M\) for one companion ideal of the phase
(Lemma~\ref{bcalc:general-companion-phase}). The sum of the finitely many bounds bounds the numbers of centres meeting \(K\).
\end{proof}

\begin{lemma}[Local finiteness from well-ordered values]
\label{hir:invariant-well-founded}
\uses{hir:centre-function-seq,hir:step-one-well-founded,hir:step-one-phase,bcalc:general-phase-count,hir:max-order-value,hir:max-order-value-basic,hir:general-seq-finite}
\source{wlodarczyk05-hironaka,bierstone-milman08-functoriality}
\lean{TranscendentalNakamaye.Pending.CanonicalCentres.stepDrop_of_isInvValue,TranscendentalNakamaye.Pending.CanonicalCentres.seqLocallyFinite_of_wellFoundedLT,TranscendentalNakamaye.Pending.CanonicalCentres.ofWellFounded,TranscendentalNakamaye.Pending.CanonicalCentres.seqLocallyFiniteMaxOrder_of_wellFoundedLT}
\leanok
Let \(\operatorname{inv}\) be a value function with properties (1)--(7) and (4') of Definition~\ref{hir:centre-function} for the marked ideals of the domain in dimension
\(\le k\), with values in a well-ordered set. Then \(\operatorname{inv}\) has property (8\('\)) of Definition~\ref{hir:centre-function-seq}, so it is a resolution invariant
in dimension \(\le k\) in that sense. Likewise a value function for the marked ideals of maximal order (Definition~\ref{hir:max-order-value}) with well-ordered values has
(8\('\)) for the marked ideals of the domain of maximal order, the hypothesis of Theorem~\ref{hir:general-seq-finite}.
\end{lemma}
\begin{proof}
\leanok
By (3) and (4), after every step of a canonical segment a point of the new cosupport lies over the cosupport, with a smaller value over the centre and the same value off
it (the hypothesis of Theorem~\ref{hir:step-one-well-founded}; for the marked ideals of maximal order this is Lemma~\ref{hir:max-order-value-basic}(2)). Let \((T_i)\) be a
growing sequence of canonical segments and \(v(T_i)\) the sequences of values of their nonempty centres (Theorem~\ref{hir:step-one-phase}). Each \(v(T_i)\) is strictly
decreasing (Theorem~\ref{hir:step-one-phase}(2)) and is a prefix of \(v(T_{i+1})\) (Theorem~\ref{hir:step-one-phase}(1)), so the \(v(T_i)\) are the initial segments of
one strictly decreasing sequence; since the values are well-ordered this sequence is finite (Lemma~\ref{bcalc:general-phase-count}), so the lengths of the \(v(T_i)\) are
bounded. The number of centres of \(T_i\) meeting the preimage of any set is at most the length of \(v(T_i)\) (Theorem~\ref{hir:step-one-phase}(1)).
\end{proof}
This is how the source terminates the procedure (the values are discrete, Włodarczyk, end of the proof of Proposition ``pr:resol''). Without well-ordered
values, Theorem~\ref{hir:general-seq-finite} proves (8\('\)) of the general value from (8\('\)) of the maximal-order value by the phases of the general value; for the
Step 1 value itself the phases give a bound within one lasting phase (Lemma~\ref{mc:step-one-phase-bound}), and the assembly over all phases is carried out in the
proof of Proposition~\ref{hir:maximal-order-step}. Neither is needed for Theorem~\ref{hir:canonical-centres}.

\begin{theorem}[The general step for well-ordered values]
\label{hir:general-step-wf}
\uses{hir:step-one-seq,hir:general-step-value,hir:invariant-well-founded,hir:centre-function-seq,bcalc:nonmonomial-order-global-phase,hir:general-value}
\source{wlodarczyk05-hironaka,bierstone-milman08-functoriality}
\lean{TranscendentalNakamaye.Pending.CanonicalCentres.generalStep,TranscendentalNakamaye.Pending.CanonicalCentres.generalStep_inv}
\leanok
Let \(n\ge1\) and let \(\operatorname{inv}\) be a resolution invariant in dimension \(\le k\) (Definition~\ref{hir:centre-function-seq}) with values in a well-ordered set
\(\Lambda\). Then the general value (Definition~\ref{hir:general-value}) of the masked Step 1 value defined with \(\operatorname{inv}\) is a resolution invariant in dimension
\(\le k+1\), with values in the well-ordered set \(\mathbf N\times(\Gamma_{\bot}\sqcup(\mathbf N\times(\Lambda\cup\{\infty\}))\cup\{\infty\})\), ordered lexicographically,
\(\Gamma\) the values of the monomial invariant.
\end{theorem}
\begin{proof}
\leanok
The masked Step 1 value is a value function for the marked ideals of maximal order in dimension \(\le k+1\) (Corollary~\ref{hir:step-one-seq}), so the general value has
properties (1)--(7) and (4') (Theorem~\ref{hir:general-step-value}). Its values are well-ordered when \(\Lambda\) is (Lemma~\ref{bcalc:nonmonomial-order-global-phase}), so
it has (8\('\)) by Lemma~\ref{hir:invariant-well-founded}.
\end{proof}

\begin{proposition}[General marked ideals]
\label{hir:general-step}
\uses{hir:maximal-order-step,hir:companion-ideal,hir:companion-basic,hir:companion-reduction,hir:centre-function,hir:invariant-order,hir:max-order-value,hir:max-order-value-basic,hir:general-value,hir:general-value-basic,hir:general-compare,hir:general-canonical-tail,hir:general-usc,hir:general-admissible,hir:general-representative,hir:general-tower-iso,hir:general-restrict-open,hir:general-empty-prefix,hir:general-step-value,hir:centre-function-seq,hir:general-seq-finite,hir:invariant-well-founded,hir:general-step-wf,hir:nonmonomial-transform-ambient,hir:nonmonomial-transform-along,hir:monomial-case-local,hir:canonical-centres-zero}
\source{bierstone-milman08-functoriality,wlodarczyk05-hironaka}
\lean{TranscendentalNakamaye.Pending.MaxOrderSeqFinite.generalStepOfExt,TranscendentalNakamaye.Pending.MaxOrderSeqFinite.general_step,TranscendentalNakamaye.Pending.MaxOrderSeqFinite.general_step',TranscendentalNakamaye.Pending.MaxOrderSeqFinite.resolutionInvariantOfEqZero,TranscendentalNakamaye.Pending.GeneralSeqFinite.ofMaxOrder}
\leanok
Assume a resolution invariant in dimension \(\le k-1\). Then there is a resolution invariant in dimension \(\le k\).
\end{proposition}
\begin{proof}
\leanok
Let \(\underline{\mathcal I}\) have \(\dim N=k\). If \(k=0\), the cosupport is empty (\(d\ge1\), \(\mathfrak a|_N\ne0\) near each point) and there is nothing to define.
At a point \(x\) of the cosupport along a canonical segment, let \(o=\operatorname{ord}_x\mathcal N(\underline{\mathcal I}_T)\). If \(o\ge1\), take an open \(W\ni x\) on which
\(\operatorname{ord}\mathcal N\le o\) on the cosupport (upper semicontinuity), the companion ideal \(\mathcal O\) on \(W\) (Definition~\ref{hir:companion-ideal}), which is of
maximal order (Lemma~\ref{hir:companion-basic}), and set \(\operatorname{inv}(x)=(o/d,\overline{\operatorname{inv}}_{\mathcal O}(x))\) with the invariant of
Proposition~\ref{hir:maximal-order-step}, computed along the part of the segment since the value \(o\) was first reached near \(x\) (before that, all centres near
\(x\) were empty, since they lay over the closed set where \(\operatorname{ord}\mathcal N>o\)); values computed from different starting stages are compared by
re-basing across those empty steps (Definition~\ref{hir:centre-function}(4')). If \(o=0\) the marked ideal is monomial near \(x\) and
\(\operatorname{inv}(x)=(0,\gamma(x))\) (Lemma~\ref{hir:monomial-case}). The companion ideals for different \(W\) agree on overlaps, so the value is well defined.
By Lemmas~\ref{hir:companion-reduction} and~\ref{hir:nonmonomial-transform}, a canonical segment of \(\mathcal O\) is a multiple test blow-up of
\(\underline{\mathcal I}|_W\), whose cosupport points with \(\operatorname{ord}\mathcal N=o\) are exactly the cosupport of the transform of \(\mathcal O\); when \(\mathcal O\) is
resolved, \(o\) has dropped at every point. Properties (1)--(7) and (4') follow from those of Proposition~\ref{hir:maximal-order-step} and Lemma~\ref{hir:monomial-case} (proved for the general value in
Theorem~\ref{hir:general-step-value}), and
from the upper semicontinuity of \(x\mapsto\operatorname{ord}_x\mathcal N\). (8\('\)) of Definition~\ref{hir:centre-function-seq}: along a growing sequence of
canonical segments, over a compact set the values of \(o\) are bounded, each value is left after finitely many steps by Proposition~\ref{hir:maximal-order-step}(8\('\))
along the restricted sequences of the companion ideals, and the monomial phase is finite by Lemma~\ref{hir:monomial-case}(3).
With well-ordered values, (8\('\)) is immediate (Lemma~\ref{hir:invariant-well-founded}); this is the form proved, Theorem~\ref{hir:general-step-wf}.
For a lower invariant with values in an arbitrary totally ordered set, as stated, the step uses (8\('\)) of the maximal-order value of
Proposition~\ref{hir:maximal-order-step} (proved by the assembly over the phases); Theorem~\ref{hir:general-seq-finite} reduces (8\('\)) of the general value to it,
and the resolution invariant obtained in dimension \(\le k\) is one in the sense of Definition~\ref{hir:centre-function-seq}, the property (8\('\)) proved above. This is
the Lean form. For \(n=0\) every value function is a resolution invariant, all cosupports of the domain being empty (Lemma~\ref{hir:canonical-centres-zero}).
\end{proof}

\begin{lemma}[Resolution invariants in dimension zero]
\label{hir:canonical-centres-zero}
\uses{hir:centre-function,hir:wellformed-tower,bcalc:transform-off-centre,hir:test-tower}
\source{bierstone-milman08-functoriality}
\lean{TranscendentalNakamaye.Pending.ResolutionInvariantSpec.cosupp_eq_empty_of_dimLE_zero,TranscendentalNakamaye.Pending.ResolutionInvariantSpec.cosupp_geoTransformAlong_eq_empty_of_dimLE_zero,TranscendentalNakamaye.Pending.ResolutionInvariantSpec.isAdmissibleCentre_empty,TranscendentalNakamaye.Pending.ResolutionInvariantSpec.countCentresMeeting_eq_zero_of_dimLE_zero,TranscendentalNakamaye.Pending.ResolutionInvariantExt.ResolutionInvariant.ofDimZero}
\leanok
If \(\dim N\le0\), \(d\ge1\) and \(\mathfrak a|_N\) vanishes identically near no point of \(N\), then \(\operatorname{cosupp}\underline{\mathcal I}_T=\emptyset\) along every multiple test
blow-up \(T\) of \(\underline{\mathcal I}\), and every canonical segment has only empty centres. Hence every constant function is a resolution invariant in
dimension \(\le0\) (Definition~\ref{hir:centre-function}).
\end{lemma}
\begin{proof}
\leanok
A point \(x\in N\) is isolated in \(N\), so a germ of \(\mathfrak a_x\) not vanishing identically on \(N\) near \(x\) does not vanish at \(x\); its order along \(N\) is \(0<d\), and
\(x\notin\operatorname{cosupp}\underline{\mathcal I}\). Along a multiple test blow-up every centre lies in the (empty) cosupport, so it is empty, and a blowup along \(\emptyset\)
creates no cosupport (Lemma~\ref{bcalc:transform-off-centre}; the stages stay well formed, Lemma~\ref{hir:wellformed-tower}). So properties (2)--(7) hold vacuously, the level
sets of (1) are empty (the empty set is an admissible centre), and a canonical segment, whose centres are empty or maximum loci on the empty cosupport, has
no nonempty centre (8).
\end{proof}

\begin{theorem}[Canonical resolution invariant]
\label{hir:canonical-centres}
\uses{hir:general-step-wf,hir:invariant-well-founded,hir:centre-function,hir:centre-function-seq,hir:canonical-centres-zero,bcalc:nonmonomial-order-global-phase}
\source{bierstone-milman08-functoriality,wlodarczyk05-hironaka}
\lean{TranscendentalNakamaye.Pending.CanonicalCentres.ValTy,TranscendentalNakamaye.Pending.CanonicalCentres.valTyInst,TranscendentalNakamaye.Pending.CanonicalCentres.linearOrder_valTy_succ,TranscendentalNakamaye.Pending.CanonicalCentres.invFamily,TranscendentalNakamaye.Pending.CanonicalCentres.invFamily_succ,TranscendentalNakamaye.Pending.CanonicalCentres.canonicalInvariant,TranscendentalNakamaye.Pending.CanonicalCentres.exists_resolutionInvariant}
\leanok
For every \(n\) and every \(k\le n\) there is a resolution invariant in dimension \(\le k\) on complex manifolds of dimension \(n\), in the sense of
Definition~\ref{hir:centre-function-seq}, with values in a well-ordered set. Explicitly, for \(n\ge1\) and every \(k\), with \(\Lambda_0=\mathbf N\) and
\(\Lambda_{k+1}=\mathbf N\times(\Gamma_{\bot}\sqcup(\mathbf N\times(\Lambda_k\cup\{\infty\}))\cup\{\infty\})\) ordered lexicographically (\(\Gamma\) the values of the
monomial invariant), there is a resolution invariant \(\operatorname{inv}_k\) in dimension \(\le k\) with values in \(\Lambda_k\): \(\operatorname{inv}_0\) a constant and
\(\operatorname{inv}_{k+1}\) the general value of the masked Step 1 value defined with \(\operatorname{inv}_k\).
\end{theorem}
\begin{proof}
\leanok
Induction on \(k\). For \(k=0\) every cosupport is empty and every constant function is a resolution invariant (Lemma~\ref{hir:canonical-centres-zero}). If
\(\operatorname{inv}_k\) is a resolution invariant in dimension \(\le k\) with values in the well-ordered set \(\Lambda_k\), Theorem~\ref{hir:general-step-wf} gives one in
dimension \(\le k+1\) with values in \(\Lambda_{k+1}\), which is well-ordered (Lemma~\ref{bcalc:nonmonomial-order-global-phase}(5)). For \(n=0\) only \(k=0\) occurs.
\end{proof}

\begin{lemma}[The canonical resolution of a smooth centre]
\label{hir:smooth-centre}
\uses{hir:centre-function,bcalc:smooth-centre-locus,hir:biholomorphism-invariance,hir:open-restriction-marked,hir:canonical-segments,hir:tower-isomorphism,mc:smooth-centre-translation,mc:smooth-centre-invariant}
\source{wlodarczyk05-hironaka}
\lean{TranscendentalNakamaye.Pending.SmoothCentreInvariant.inv_refl_locallyConstant,TranscendentalNakamaye.Pending.SmoothCentreInvariant.inv_refl_eq_of_isPreconnected,TranscendentalNakamaye.Pending.SmoothCentreInvariant.mem_cosupp_geoTransformAlong_of_isEmptyTower,TranscendentalNakamaye.Pending.SmoothCentreInvariant.inv_eq_inv_refl_of_isEmptyTower,TranscendentalNakamaye.Pending.SmoothCentreInvariant.inv_eq_of_isEmptyTower_of_isPreconnected,TranscendentalNakamaye.Pending.SmoothCentreInvariant.maxLocus_eq_cosupp_of_isEmptyTower,TranscendentalNakamaye.Pending.SmoothCentreInvariant.inv_refl_eq_of_regularAtDim,TranscendentalNakamaye.Pending.SmoothCentreInvariant.inv_refl_eq_of_pureDim,TranscendentalNakamaye.Pending.SmoothCentreInvariant.inv_eq_of_isEmptyTower_of_pureDim,TranscendentalNakamaye.Pending.SmoothCentreInvariant.maxLocus_eq_cosupp_of_isEmptyTower_of_pureDim}
\leanok
Let \(\operatorname{inv}\) be a resolution invariant in dimension \(\le n\) (Definition~\ref{hir:centre-function}), \(Z\subset U\) a closed connected submanifold of pure
codimension \(c\ge1\) of a complex manifold \(U\) of dimension \(n\), and \(\underline{\mathcal I}=(U,U,\emptyset,\mathcal I_Z,1)\). Then
\(\operatorname{inv}_{\underline{\mathcal I}}\) of the empty tower is constant on \(\operatorname{cosupp}\underline{\mathcal I}=Z\). Along a canonical segment \(T\) of
\(\underline{\mathcal I}\) all of whose centres are empty, \(\operatorname{inv}_{\underline{\mathcal I}}(T)\) is constant on \(\operatorname{cosupp}\underline{\mathcal I}_T\);
so a nonempty maximum locus after such a segment is the whole cosupport. Connectedness is not needed: the same holds for a closed submanifold
\(Z\) of pure dimension, and \(\operatorname{inv}_{\underline{\mathcal I}}\) of the empty tower takes the same value at any two points where \(Z\) has the same dimension.
\end{lemma}
\begin{proof}
\leanok
\(\underline{\mathcal I}\) lies in the domain and its cosupport is \(Z\) (Proposition~\ref{bcalc:smooth-centre-locus}). Let \(x,x'\in Z\) be close. In a chart in which
\(Z\) is a coordinate subspace, a translation along that subspace is a biholomorphism \(\varphi\colon U_1\to U_2\) between neighbourhoods of \(x\) and \(x'\) with
\(\varphi(x)=x'\) and \(\varphi^{-1}(Z\cap U_2)=Z\cap U_1\), so \(\varphi^*(\underline{\mathcal I}|_{U_2})=\underline{\mathcal I}|_{U_1}\)
(Lemma~\ref{mc:smooth-centre-translation}; Lemma~\ref{hir:biholomorphism-invariance}: the vanishing ideal of \(Z\) pulls back to that of its preimage). By Definition~\ref{hir:centre-function}(5) and (6) for the empty
towers, \(\operatorname{inv}_{\underline{\mathcal I}}(x)=\operatorname{inv}_{\underline{\mathcal I}|_{U_1}}(x)=\operatorname{inv}_{\underline{\mathcal I}|_{U_2}}(x')=
\operatorname{inv}_{\underline{\mathcal I}}(x')\). So the value is locally constant on \(Z\), hence constant, \(Z\) being connected. Along \(T\), by
Definition~\ref{hir:centre-function}(4) and induction (each prefix is again a canonical segment with empty centres), \(\operatorname{inv}_{\underline{\mathcal I}}(T)\) is the
value of the empty tower composed with the projection, and \(\operatorname{cosupp}\underline{\mathcal I}_T\) is the preimage of \(Z\). (An explicit computation through
the companion ideal and maximal contact gives the value \((1,0,\infty)\) if \(c=1\) and \((1,0,\operatorname{inv}')\) with \(\operatorname{inv}'\) the invariant of
\((U,V(u),\emptyset,\mathcal I_Z,1)\) otherwise; it is not needed.) The two claims are Lemma~\ref{mc:smooth-centre-invariant}, which takes \(Z\) preconnected and the
domain hypothesis as given by Proposition~\ref{bcalc:smooth-centre-locus}. Without connectedness: any two points where \(Z\) has the same dimension have
neighbourhoods matched by a biholomorphism preserving \(Z\) (the composite of two flat charts on a common ball, Lemma~\ref{mc:smooth-centre-translation}),
so the same argument compares the values at any two such points directly (Lemma~\ref{mc:smooth-centre-invariant}).
\end{proof}

\subsection{The resolution theorem}



\begin{theorem}[The driver: iterating the blowup of the maximum locus]
\label{hir:resolution-driver}
\uses{hir:geo-transform,hir:geo-transform-basic,hir:test-tower,bex:blowup-exists,ana:blowup-compact,hir:invariant-order}
\source{bierstone-milman08-functoriality,wlodarczyk05-hironaka,kollar07-resolution}
\lean{TranscendentalNakamaye.Pending.ResolutionDriver.compactSpace_of_tower,TranscendentalNakamaye.Pending.ResolutionDriver.exists_blowup,TranscendentalNakamaye.Pending.ResolutionDriver.exists_isGeoResolution_of_bounded,TranscendentalNakamaye.Pending.ResolutionDriver.maxLocus,TranscendentalNakamaye.Pending.ResolutionDriver.IsCanonicalSegment,TranscendentalNakamaye.Pending.ResolutionDriver.isCanonicalSegment_refl,TranscendentalNakamaye.Pending.ResolutionDriver.isCanonicalSegment_step_iff,TranscendentalNakamaye.Pending.ResolutionDriver.maxLocus_nonempty,TranscendentalNakamaye.Pending.ResolutionDriver.exists_isGeoResolution_of_bounded_invariant,TranscendentalNakamaye.Pending.ResolutionDriver.exists_isGeoResolution_of_wellFounded_invariant}
\leanok
Let \(X\) be a compact complex manifold, \(\underline{\mathcal I}\) a marked ideal on \(X\), \(\Lambda\) a totally ordered set, and \(\operatorname{inv}\) a function assigning to
every tower \(T\) over \(X\) a function \(\operatorname{inv}(T)\colon X_T\to\Lambda\) on its last stage. Call \(T\) a \emph{canonical segment} if each of its centres is
nonempty and is the maximum locus of \(\operatorname{inv}\) of the preceding part on the cosupport of the transform. Assume that for every canonical segment \(T\)
which is a multiple test blow-up of \(\underline{\mathcal I}\): the cosupport of \(\underline{\mathcal I}_T\) is closed, \(\operatorname{inv}(T)\) is upper semicontinuous on it
(Definition~\ref{hir:invariant-order}), and, if the cosupport is nonempty, its maximum locus is an admissible centre with empty interior. Assume moreover
either
\begin{enumerate}
\item that such canonical segments have length at most some \(k\), or
\item that \(\Lambda\) is well-founded and, after blowing up the maximum locus, \(\operatorname{inv}\) at every point of the new cosupport is smaller than the old
maximum.
\end{enumerate}
Then some canonical segment is a resolution of \(\underline{\mathcal I}\). More abstractly, for any class of towers containing the empty tower in which every
multiple test blow-up with nonempty cosupport can be extended by the blowup of some admissible centre with empty interior, a bound on the lengths gives a
resolution in the class.
\end{theorem}
\begin{proof}
\leanok
Every stage is compact (Lemma~\ref{ana:blowup-compact}, induction on the tower), so a nonempty closed cosupport is compact and the upper semicontinuous
invariant attains its maximum on it (Lemma~\ref{hir:invariant-order}): the maximum locus is nonempty. An admissible centre with empty interior can be blown up
(Lemma~\ref{bex:blowup-exists}), and the result is again a canonical segment and a multiple test blow-up. Under (1), induction on \(k\) minus the length: a
canonical segment of length \(k\) with nonempty cosupport would extend to one of length \(k+1\). Under (2), well-founded induction on an upper bound \(l\) for the
values on the cosupport: after one blowup the new values are below the old maximum, which is \(\le l\).
\end{proof}

\begin{theorem}[A resolution invariant resolves its domain]
\label{hir:invariant-driver}
\uses{hir:centre-function,hir:resolution-driver,hir:invariant-order,hir:open-restriction-marked,hir:test-tower,hir:giraud-along,bcalc:inv-domain-stable,hir:cosupport}
\source{wlodarczyk05-hironaka}
\lean{TranscendentalNakamaye.Pending.ResolutionInvariantDriver.isAdmissibleCentre_of_forall_exists_opens,TranscendentalNakamaye.Pending.ResolutionInvariantDriver.isAdmissibleCentre_iff_forall_exists_opens,TranscendentalNakamaye.Pending.ResolutionInvariantDriver.isAdmissibleCentre_of_pullback_val_top,TranscendentalNakamaye.Pending.ResolutionInvariantDriver.interior_cosupp_eq_empty,TranscendentalNakamaye.Pending.ResolutionInvariantDriver.maxLocus_eq_level,TranscendentalNakamaye.Pending.ResolutionInvariantDriver.maxLocus_subset,TranscendentalNakamaye.Pending.ResolutionInvariantExt.isAdmissibleCentre_maxLocus,TranscendentalNakamaye.Pending.ResolutionInvariantExt.isGeoTestTower_of_isCanonical,TranscendentalNakamaye.Pending.ResolutionInvariantExt.isInvDomain_of_isCanonical,TranscendentalNakamaye.Pending.ResolutionInvariantDriver.interior_maxLocus_eq_empty,TranscendentalNakamaye.Pending.ResolutionInvariantDriver.isCanonical_of_isCanonicalSegment,TranscendentalNakamaye.Pending.ResolutionInvariantDriver.countCentresMeeting_univ_eq_length,TranscendentalNakamaye.Pending.ResolutionInvariantExt.exists_isGeoResolution_isCanonicalSegment,TranscendentalNakamaye.Pending.ResolutionInvariantExt.exists_isGeoResolution_isCanonical}
\leanok
Let \(\operatorname{inv}\) be a resolution invariant in dimension \(\le k\) (Definition~\ref{hir:centre-function}) and \(\underline{\mathcal I}\) a marked ideal of its
domain.
\begin{enumerate}
\item (Admissible centres are local.) A subset \(C\) of a complex manifold \(M\) is an admissible centre of a marked ideal \(\underline{\mathcal I}\) on \(M\) (ideal of
holomorphic germs) if and only if every point of \(M\) has an open neighbourhood \(W\) such that \(C\cap W\) is an admissible centre of \(\underline{\mathcal I}|_W\).
\item The cosupport of a marked ideal of the domain has empty interior.
\item Along a canonical segment \(T\) of \(\underline{\mathcal I}\) which is a multiple test blow-up, the maximum locus of \(\operatorname{inv}_{\underline{\mathcal I}}(T)\) on
\(\operatorname{cosupp}\underline{\mathcal I}_T\) is an admissible centre with empty interior. Every canonical segment of \(\underline{\mathcal I}\) is a multiple test
blow-up, and its transform lies in the domain.
\item If \(X\) is compact, some canonical segment of \(\underline{\mathcal I}\) without empty centres is a resolution of \(\underline{\mathcal I}\).
\end{enumerate}
\end{theorem}
\begin{proof}
\leanok
(1) Restriction preserves admissibility (Lemma~\ref{hir:open-restriction-marked}). Conversely, closedness is local; at \(x\in C\) the regularity of \(C\) and the adapted
coordinates for \(C\) and the members of \(E\) through \(x\) are read off in \(W\) (the charts of \(W\) are those of \(M\)), and
\(\operatorname{cosupp}(\underline{\mathcal I}|_W)=\operatorname{cosupp}\underline{\mathcal I}\cap W\) (Lemma~\ref{hir:open-restriction-marked}).
(2) Let \(y\) be an interior point of the cosupport. Near \(y\) the ideal is generated by finitely many sections \(s_j\) holomorphic on an open \(U\ni y\). At every point
\(z\) of \(U\) in the interior of the cosupport, \(s_j\in\mathfrak a_z\) vanishes to order \(d\ge1\) along \(N\ni z\), hence \(s_j(z)=0\). So every germ of \(\mathfrak a_y\), a
holomorphic combination of the \(s_j\), vanishes on a neighbourhood of \(y\), against \(\mathfrak a\) being nonzero along \(N\) at \(y\in N\).
(3) If the maximum locus is nonempty, it is the level set of its value \(\lambda\) on the whole last stage \(W\), on which \(\operatorname{inv}\le\lambda\); by
Definition~\ref{hir:centre-function}(1) it is an admissible centre of the restriction to \(W\), hence of \(\underline{\mathcal I}_T\) by (1) (the ideal of
\(\underline{\mathcal I}_T\) is coherent, Lemma~\ref{bcalc:inv-domain-stable}); the empty set is admissible. By induction along the segment, every centre (empty or the
maximum locus) is admissible, so the segment is a multiple test blow-up (Definition~\ref{hir:test-tower}), and \(\underline{\mathcal I}_T\) lies in the domain
(Proposition~\ref{bcalc:inv-domain-stable}); its cosupport has empty interior by (2), hence so has the maximum locus.
(4) Apply Theorem~\ref{hir:resolution-driver}(1) to the canonical segments without empty centres: the cosupports are closed (Lemma~\ref{hir:giraud-along}), the invariant is
upper semicontinuous on them (Definition~\ref{hir:centre-function}(2)), the maximum loci are admissible with empty interior (3), and a canonical segment without empty
centres has as many centres meeting \(X\) as its length, at most \(m\) for \(K=X\) by Definition~\ref{hir:centre-function}(8). This is the end of the proof of
Proposition~\texttt{pr:resol} in the source: the centres are the maximum loci of the invariant, blown up until the cosupport is empty.
\end{proof}
Superseded on the route by Theorem~\ref{hir:invariant-driver-seq} (local finiteness along growing sequences); kept.

\begin{theorem}[A resolution invariant with local finiteness along growing sequences resolves its domain]
\label{hir:invariant-driver-seq}
\uses{hir:centre-function-seq,hir:invariant-driver,hir:resolution-driver,hir:giraud-along,bex:blowup-exists}
\source{wlodarczyk05-hironaka}
\lean{TranscendentalNakamaye.Pending.ResolutionInvariantSeq.IsInvValue.isAdmissibleCentre_maxLocus,TranscendentalNakamaye.Pending.ResolutionInvariantSeq.IsInvValue.isGeoTestTower_of_isCanonical,TranscendentalNakamaye.Pending.ResolutionInvariantSeq.GrowingSeq.ofIterate,TranscendentalNakamaye.Pending.ResolutionInvariantSeq.exists_isGeoResolution_of_seqFinite,TranscendentalNakamaye.Pending.ResolutionInvariantSeq.exists_isGeoResolution_isCanonicalSegment,TranscendentalNakamaye.Pending.ResolutionInvariantSeq.exists_isGeoResolution_isCanonical}
\leanok
Let \(X\) be compact.
\begin{enumerate}
\item Let \(P\) be a class of towers over \(X\) containing the empty tower such that every multiple test blow-up \(T\in P\) of \(\underline{\mathcal I}\) with nonempty
cosupport has a nonempty admissible centre with empty interior all of whose blowups extend \(T\) inside \(P\). If along every growing sequence of multiple test
blow-ups in \(P\) the numbers of centres are bounded, some tower of \(P\) is a resolution of \(\underline{\mathcal I}\).
\item For a resolution invariant \(\operatorname{inv}\) in dimension \(\le k\) in the sense of Definition~\ref{hir:centre-function-seq}, along a canonical segment the
maximum locus is an admissible centre and canonical segments are multiple test blow-ups; every marked ideal of the domain is resolved by a canonical segment
without empty centres.
\end{enumerate}
\end{theorem}
\begin{proof}
\leanok
(1) Otherwise every tower of \(P\) which is a multiple test blow-up has nonempty cosupport. Choose for each such tower a centre as in the hypothesis and a blowup
along it (Theorem~\ref{bex:blowup-exists}); iterating from the empty tower (dependent choice) gives a growing sequence in \(P\) of multiple test blow-ups whose
\(i\)-th tower has \(i\) nonempty centres, against the bound.
(2) The first two statements are proved as in Theorem~\ref{hir:invariant-driver}(3) from properties (1) and (3) only. Take for \(P\) the canonical segments without
empty centres: the cosupports are closed (Lemma~\ref{hir:giraud-along}), the invariant is upper semicontinuous on them, so on the compact stage the maximum locus of a
nonempty cosupport is nonempty, admissible and with empty interior (Theorem~\ref{hir:invariant-driver}(2)); a growing sequence of such segments is a growing
sequence of canonical segments, along which (8\('\)) with \(K=X\) bounds the numbers of nonempty centres. Apply (1).
\end{proof}

\begin{theorem}[Resolution of marked ideals]
\label{hir:canonical-resolution}
\uses{hir:resolution-driver,hir:invariant-driver,hir:invariant-driver-seq,hir:centre-function-seq,hir:canonical-segments,hir:test-tower,hir:giraud-support,hir:hadamard-bridge,hir:tower-transport-resolution,hir:open-restriction-tower,hir:canonical-centres,hir:centre-function,hir:invariant-order,bex:blowup-exists,hir:tower-split,hir:tower-isomorphism,mc:lift-biholomorphism}
\source{wlodarczyk05-hironaka,bierstone-milman08-functoriality,kollar07-resolution}
\lean{TranscendentalNakamaye.Pending.CanonicalResolution.isCanonicalSegment_of_append,TranscendentalNakamaye.Pending.CanonicalResolution.length_eq_zero_of_cosupp_eq_empty,TranscendentalNakamaye.Pending.CanonicalResolution.exists_isTowerIso_of_length_eq,TranscendentalNakamaye.Pending.CanonicalResolution.exists_isTowerIso_of_isGeoResolution,TranscendentalNakamaye.Pending.CanonicalResolution.canonicalResolution,TranscendentalNakamaye.Pending.CanonicalResolution.canonicalResolution_canonicalInvariant}
\leanok
Let \(M\) be a compact complex manifold and \(\underline{\mathcal I}=(M,N,E,\mathfrak a,d)\) a well-formed marked ideal with \(d\ge1\) such that
\(\mathfrak a|_N\) vanishes identically on no nonempty open subset of \(N\). There is a resolution of \(\underline{\mathcal I}\). Moreover a resolution \(\mathcal R(\underline{\mathcal I})\) can be chosen for every such
\(\underline{\mathcal I}\) so that
\begin{enumerate}
\item (centres are maximum loci) every centre of \(\mathcal R(\underline{\mathcal I})\) is the locus where an upper semicontinuous invariant
of the current transform, with values in a totally ordered set, attains its maximum;
\item (open embeddings) for an open set \(U\subset M\), the restriction of \(\mathcal R(\underline{\mathcal I})\) to \(U\) (a tower over \(U\) whose
centres are the intersections with the preimages of \(U\), some of them empty) is a canonical segment of \(\underline{\mathcal I}|_U\) for the same invariant;
\item (biholomorphisms) if \(\varphi\colon M'\to M\) is a biholomorphism and \(\underline{\mathcal I}'=\varphi^*\underline{\mathcal I}\), then \(\varphi\) lifts to
biholomorphisms of the stages of \(\mathcal R(\underline{\mathcal I}')\) and \(\mathcal R(\underline{\mathcal I})\) commuting with the blowups.
\end{enumerate}
\end{theorem}
\begin{proof}
\leanok
The hypotheses say that \(\underline{\mathcal I}\) lies in the domain of Definition~\ref{hir:centre-function} in dimension \(\le\dim M\). Let \(\operatorname{inv}\) be the
resolution invariant of Theorem~\ref{hir:canonical-centres} in dimension \(\le\dim M\) (Definition~\ref{hir:centre-function-seq}); call a \emph{canonical resolution} a
canonical segment without empty centres which is a resolution. One exists by Theorem~\ref{hir:invariant-driver-seq}(2): starting from the empty tower, as long as the
cosupport of the current transform is nonempty, blow up the locus where \(\operatorname{inv}\) attains its maximum on it (the cosupport is closed in the compact last
stage, \(\operatorname{inv}\) is upper semicontinuous on it, Lemma~\ref{hir:invariant-order}; the maximum locus is an admissible centre with empty interior, and its blowup
exists, Lemma~\ref{bex:blowup-exists}); by (8\('\)) of Definition~\ref{hir:centre-function-seq} along the growing sequence so obtained, with \(K=M\), the process stops,
with empty cosupport. Let \(\mathcal R(\underline{\mathcal I})\) be any canonical resolution. (1) holds by construction and Definition~\ref{hir:centre-function}(2).
(2): by Definition~\ref{hir:centre-function}(5) (Lemma~\ref{hir:canonical-segments}(1)), the restriction of \(\mathcal R(\underline{\mathcal I})\) to \(U\) is a canonical
segment of \(\underline{\mathcal I}|_U\), each restricted centre being empty or the maximum locus of the restricted invariant. (No comparison with a resolution of
\(\underline{\mathcal I}|_U\) is claimed: \(U\) need not be compact, and the restricted segment has empty steps in the middle.) (3): every canonical resolution of
\(\underline{\mathcal I}'\) is isomorphic to \(\mathcal R(\underline{\mathcal I})\) over \(\varphi\) (Definition~\ref{hir:tower-isomorphism}). First, two canonical segments
without empty centres of \(\underline{\mathcal I}'\) and of \(\underline{\mathcal I}\) with the same number of steps are isomorphic over \(\varphi\), by induction: if the
parts before the last step are isomorphic over \(\varphi\) with last-stage biholomorphism \(\chi\), then by Definition~\ref{hir:centre-function}(6) the invariants
correspond under \(\chi\), so do the cosupports and their maximum loci, i.e.\ the last centres; blowups along corresponding centres are isomorphic over \(\chi\) by the
universal property (Lemma~\ref{mc:lift-biholomorphism}). Second, two canonical resolutions have the same number of steps: otherwise the prefix of the longer one with
the length of the shorter one (Lemma~\ref{hir:tower-split}) is isomorphic to the shorter one, so its cosupport is empty, and the next centre, the maximum locus on the
empty set, is empty, which is excluded. So \(\varphi\) lifts to biholomorphisms of the stages of \(\mathcal R(\underline{\mathcal I}')\) and \(\mathcal R(\underline{\mathcal I})\)
commuting with the blowups, whichever canonical resolutions are chosen.
\end{proof}

\subsection{The two resolution statements}

\begin{lemma}[The controlled transform with mark one (design)]
\label{hir:h5a-controlled-transform}
\uses{mc:factorization-ambient,hir:transform,hir:test-tower,hir:test-tower-basic,hir:restricted-equivalence,ana:blowups,hir:geo-transform-ambient,hir:cosupport-basic,hir:hadamard-bridge,hir:bridge-consumers,bcalc:tower-exceptional,bcalc:strict-ideal,hir:open-restriction-tower,hir:giraud-support}
\source{bierstone-milman08-functoriality,wlodarczyk05-hironaka}
\lean{TranscendentalNakamaye.Pending.HironakaControlledTransform.controlledTransformAlong,TranscendentalNakamaye.Pending.HironakaControlledTransform.controlledTransformAlong_refl,TranscendentalNakamaye.Pending.HironakaControlledTransform.controlledTransformAlong_step,TranscendentalNakamaye.Pending.HironakaControlledTransform.controlledTransformAlong_eq,TranscendentalNakamaye.Pending.HironakaControlledTransform.pullback_controlledTransformAlong_step,TranscendentalNakamaye.Pending.HironakaControlledTransform.pullback_controlledTransformAlong_step_of_isGeoTestTower,TranscendentalNakamaye.Pending.HironakaControlledTransform.pullback_eq_mul_controlledTransformAlong,TranscendentalNakamaye.Pending.HironakaControlledTransform.pullback_vanishingIdeal_N_eq_mul,TranscendentalNakamaye.Pending.HironakaControlledTransform.mem_one_of_notMem_zeroLocus,TranscendentalNakamaye.Pending.HironakaControlledTransform.controlled_transform}
\leanok
Let \(\underline{\mathcal I}=(M,N,E,\mathfrak a,1)\) be a well-formed marked ideal with mark \(1\) and \(T\colon M=M_0\leftarrow\cdots\leftarrow M_r\) a multiple test blow-up of it
with centres \(C_i\subsetneq N_{i-1}\) and exceptional divisors \(D_i\). Define the \emph{controlled transforms} \(\mathfrak a^c_0=\mathfrak a\),
\(\mathfrak a^c_i=(\sigma_i^*\mathfrak a^c_{i-1}:\mathcal I_{D_i})\), ideals on the whole of \(M_i\), and let \(H_r\) be the divisor \(\sum_i(\sigma_{i+1}\circ\cdots\circ\sigma_r)^*D_i\).
Then:
\begin{enumerate}
\item \(\sigma_i^*\mathfrak a^c_{i-1}=\mathcal I_{D_i}\cdot\mathfrak a^c_i\) for every \(i\), hence \(f^*\mathfrak a=\mathcal I_{H_r}\cdot\mathfrak a^c_r\) for the composite \(f\colon M_r\to M\); and
\(f^*\mathcal I_N=\mathcal I_{H_r}\cdot\mathcal I_{N_r}\);
\item \(\mathfrak a^c_r\) and the ideal \(\mathfrak a_r\) of the transform \(\underline{\mathcal I}_T\) have the same restriction to \(N_r\);
\item if \(\operatorname{cosupp}\underline{\mathcal I}_T=\emptyset\), then \(\mathfrak a^c_r\) is the unit ideal on a neighbourhood of \(N_r\).
\end{enumerate}
\end{lemma}
\begin{proof}
\leanok
Let \(\underline{\mathcal I}'=(M,M,E,\mathfrak a,1)\). By Lemma~\ref{mc:factorization-ambient}, \(T\) is also a multiple test blow-up of \(\underline{\mathcal I}'\), the
ideal \(\mathfrak a_r\) of \(\underline{\mathcal I}_T\) has the same restriction to \(N_r\) as the ideal \(\mathfrak a'_r\) of \(\underline{\mathcal I}'_T\), and
\(\operatorname{cosupp}\underline{\mathcal I}_T=N_r\cap Z(\mathfrak a'_r)\). The same holds for every prefix \(M_0\leftarrow\cdots\leftarrow M_i\) of \(T\), which is a multiple test
blow-up of \(\underline{\mathcal I}\). First, \(\mathfrak a^c_i=\mathfrak a'_i\) for every \(i\), by induction on \(i\): with \(N=M\), mark one and a coherent ideal, the ideal of
the transform by \(\sigma_i\) is the colon ideal \((\sigma_i^*\mathfrak a'_{i-1}:\mathcal I_{D_i})\) and is coherent (Lemma~\ref{hir:geo-transform-ambient}). This gives (2).
(1) The centre \(C_i\) is an admissible centre of \(\underline{\mathcal I}'_{i-1}\), so it lies in its cosupport, the zero locus of \(\mathfrak a'_{i-1}=\mathfrak a^c_{i-1}\)
(Lemma~\ref{hir:cosupport-basic}); every germ of \(\mathfrak a^c_{i-1}\) vanishes on \(C_i\), its pull-back is divisible by a local equation of \(D_i\)
(Lemma~\ref{hir:hadamard-bridge}), and \(\sigma_i^*\mathfrak a^c_{i-1}=\mathcal I_{D_i}\cdot(\sigma_i^*\mathfrak a^c_{i-1}:\mathcal I_{D_i})=\mathcal I_{D_i}\cdot\mathfrak a^c_i\)
(Corollary~\ref{hir:bridge-consumers}). Pulling back along the later blowups and multiplying gives \(f^*\mathfrak a=\mathcal I_{H_r}\cdot\mathfrak a^c_r\), where \(H_r\) is the
exceptional divisor of the tower (Definition~\ref{bcalc:tower-exceptional}, Lemma~\ref{hir:geo-transform-ambient}). For \(\mathcal I_N\): \(N_{i-1}\) is a closed submanifold
containing \(C_i\) (Lemma~\ref{hir:open-restriction-tower}), so \(\sigma_i^*\mathcal I_{N_{i-1}}=\mathcal I_{D_i}\cdot\mathcal I_{N_i}\) with \(N_i\) the strict transform of
\(N_{i-1}\) (Theorem~\ref{bcalc:strict-ideal}); by induction on \(i\), using that the total transform of a product of coherent ideals is the product of the total transforms,
\(f^*\mathcal I_N=\mathcal I_{H_r}\cdot\mathcal I_{N_r}\).
(3) If \(\operatorname{cosupp}\underline{\mathcal I}_T=N_r\cap Z(\mathfrak a^c_r)\) is empty, then \(N_r\) lies in \(U=M_r\setminus Z(\mathfrak a^c_r)\), which is open as
\(\mathfrak a^c_r\) is coherent (Lemma~\ref{hir:giraud-support}). At \(x\in U\) some germ \(b\) of \(\mathfrak a^c_r\) has \(b(x)\ne0\), so \(1=b^{-1}\cdot b\) lies in
\(\mathfrak a^c_r\) at \(x\): \(\mathfrak a^c_r\) is the unit ideal on the open neighbourhood \(U\) of \(N_r\).
\end{proof}

\begin{lemma}[Relative principalization, first stage (design)]
\label{hir:h5a-first-stage}
\uses{mc:factorization-tower,mc:factorization-thin-centre,hir:h5a-controlled-transform,mc:factorization-domain,hir:invariant-driver-seq,hir:centre-function-seq,ana:blowups,ana:strict-transform-sequence,hir:geo-transform-ambient,hir:resolution-driver,hir:canonical-centres}
\source{transcendental-nakamaye-copy4-20260911:lem:relative-principalization,bierstone-milman08-functoriality}
\lean{TranscendentalNakamaye.Pending.HironakaControlledTransform.first_stage,TranscendentalNakamaye.Pending.HironakaControlledTransform.exists_first_stage}
\leanok
Let \(Y\) be a compact connected complex manifold of dimension \(n\), \(W\subsetneq Y\) a closed connected submanifold of dimension \(k\) and \(\mathfrak a\) a coherent
ideal with \(\mathfrak a\not\subset\mathcal I_W\). Let \(T_A\) be a multiple test blow-up of \(\underline{\mathcal I}_A=(Y,W,\emptyset,\mathfrak a,1)\) with
\(\operatorname{cosupp}(\underline{\mathcal I}_A)_{T_A}=\emptyset\) (it exists if there is a resolution invariant in dimension \(\le k\): \(\underline{\mathcal I}_A\) lies in the domain by
Lemma~\ref{mc:factorization-domain}; Theorem~\ref{hir:invariant-driver-seq}). Let \(f_A\colon Y_t\to Y\) be its composite. Then \(f_A\) is a sequence of blowups adapted to \(W\)
(every centre lies in the cosupport, a proper analytic subset of the current strict transform of \(W\)); the strict transform \(W_t\) is a closed connected
submanifold of dimension \(k\); \(Y_t\) is compact and connected; with \(H_A\) and \(\mathfrak a_t=\mathfrak a^c\) as in Lemma~\ref{hir:h5a-controlled-transform},
\(f_A^*\mathcal I_W=\mathcal I_{H_A}\cdot\mathcal I_{W_t}\), \(f_A^*\mathfrak a=\mathcal I_{H_A}\cdot\mathfrak a_t\), the support of \(H_A\) lies in the exceptional locus, and \(\mathfrak a_t\) is
a coherent ideal, the unit ideal near \(W_t\), in particular not the zero ideal.
\end{lemma}
\begin{proof}
\leanok
Here \(k<n\): a closed submanifold of dimension \(n\) is open, and \(W\ne Y\) with \(Y\) connected. \(\underline{\mathcal I}_A\) lies in the domain in dimension \(\le k\)
(Lemma~\ref{mc:factorization-domain}); in particular it is well formed. Theorem~\ref{mc:factorization-tower} gives that \(f_A\) is a sequence of blowups adapted to \(W\)
(every centre lies in the cosupport, a proper analytic subset of the current strict transform of \(W\)), that \(W_t=N_{T_A}\) is the strict transform of \(W\)
(Lemma~\ref{ana:strict-transform-sequence}) and a closed connected submanifold of dimension \(k\), that \(f_A^*\mathcal I_W=\mathcal I_{H_A}\cdot\mathcal I_{W_t}\), and that the
support of \(H_A\) lies in the exceptional locus. Lemma~\ref{hir:h5a-controlled-transform} gives \(f_A^*\mathfrak a=\mathcal I_{H_A}\cdot\mathfrak a_t\); \(\mathfrak a_t\) is the
ideal of the transform of \((Y,Y,\emptyset,\mathfrak a,1)\), hence coherent (Lemma~\ref{hir:geo-transform-ambient}); and, the final cosupport being empty, \(\mathfrak a_t\) is the
unit ideal on an open neighbourhood of \(W_t\). As \(W_t\) is nonempty, \(\mathfrak a_t\) contains the germ \(1\) at some point, so it is not the zero ideal. \(Y_t\) is compact,
the blowups along closed centres being proper (Theorem~\ref{hir:resolution-driver}), and connected: a blowup of a connected manifold along a nowhere dense closed centre is
connected, since the complement of the exceptional divisor is dense in it and biholomorphic to the complement of the centre, which is connected. Finally such a \(T_A\) exists
without hypothesis: the canonical resolution invariant (Theorem~\ref{hir:canonical-centres}) and the driver (Theorem~\ref{hir:invariant-driver-seq}) resolve
\(\underline{\mathcal I}_A\), which lies in the domain in dimension \(\le k\).
\end{proof}

\begin{lemma}[Relative principalization, second stage (design)]
\label{hir:h5a-second-stage}
\uses{mc:factorization-second,hir:h5a-controlled-transform,hir:h5a-first-stage,mc:factorization-domain,hir:invariant-driver-seq,hir:centre-function-seq,ana:blowups,hir:geo-transform-ambient,hir:cosupport-basic,hir:giraud-support,bcalc:restrict-empty-tower,hir:principalization,hir:principalization-local}
\source{transcendental-nakamaye-copy4-20260911:lem:relative-principalization,bierstone-milman08-functoriality}
\lean{TranscendentalNakamaye.Pending.HironakaControlledTransform.second_stage,TranscendentalNakamaye.Pending.HironakaControlledTransform.second_stage_of_first_stage}
\leanok
Let \(Y_t\), \(W_t\) and \(\mathfrak a_t\) be as in Lemma~\ref{hir:h5a-first-stage}, and \(T_B\) a multiple test blow-up of \(\underline{\mathcal I}_B=(Y_t,Y_t,\emptyset,\mathfrak a_t,1)\)
with empty final cosupport (it exists if there is a resolution invariant in dimension \(\le n\): \(\underline{\mathcal I}_B\) lies in the domain by
Lemma~\ref{mc:factorization-domain} with \(N=Y_t\)). Let \(g\colon Y'\to Y_t\) be its composite and \(H'\) the divisor of Lemma~\ref{hir:h5a-controlled-transform}. Then every
centre of \(T_B\) is disjoint from the strict transform of \(W_t\), \(g\) is an isomorphism near it, \(g^*\mathfrak a_t=\mathcal I_{H'}\), \(H'\) is disjoint from the strict
transform \(W'\) of \(W_t\), and \(g^*\mathcal I_{W_t}=\mathcal I_{W'}\).
\end{lemma}
\begin{proof}
\leanok
The argument works for any coherent ideal \(\mathfrak b\) on a complex manifold \(X\) and any closed connected submanifold \(A\) of dimension \(k\) disjoint from \(Z(\mathfrak b)\),
with \(T_B\) resolving \((X,X,\emptyset,\mathfrak b,1)\); here \(\mathfrak b=\mathfrak a_t\), \(A=W_t\). By Lemma~\ref{hir:h5a-first-stage}, \(\mathfrak a_t\) is coherent and the unit
ideal on an open neighbourhood of \(W_t\), so \(W_t\) misses \(Z(\mathfrak a_t)\), which is closed (Lemma~\ref{hir:giraud-support}); \(U=Y_t\setminus Z(\mathfrak a_t)\) is an open
neighbourhood of \(W_t\). With \(N=Y_t\) and mark one, the ideal of each transform of \(\underline{\mathcal I}_B\) divides the total transform of \(\mathfrak a_t\)
(Lemma~\ref{hir:geo-transform-ambient}), so the cosupport, its zero locus (Lemma~\ref{hir:cosupport-basic}), lies over \(Z(\mathfrak a_t)\): every centre of \(T_B\) lies over
\(Z(\mathfrak a_t)\) and misses the preimage of \(U\) (Lemma~\ref{bcalc:restrict-empty-tower}), which contains the strict transform of \(W_t\) at its stage.
Lemma~\ref{mc:factorization-second} gives: \(W'=g^{-1}(W_t)\) is the strict transform of \(W_t\) and a closed connected submanifold of dimension \(k\), \(H'\) is disjoint from it,
and \(g^*\mathcal I_{W_t}=\mathcal I_{W'}\). The final transform has empty cosupport with \(N=Y'\), so its ideal is the unit ideal, and
\(g^*\mathfrak a_t=\mathcal I_{H'}\) (Lemma~\ref{hir:geo-transform-ambient}). Hence the support of \(H'\) is \(g^{-1}(Z(\mathfrak a_t))\); the exceptional locus of \(g\) lies in it
(Theorem~\ref{hir:principalization}), so \(g\) is an isomorphism over \(U\) (Theorem~\ref{hir:principalization-local}), in particular near \(W_t\).
\end{proof}

\begin{lemma}[Relative principalization from two resolutions (design)]
\label{hir:h5a-assembly}
\uses{mc:factorization-assembly,hir:h5a-first-stage,hir:h5a-second-stage,hir:h5a-controlled-transform,mc:factorization-tower,mc:factorization-second,ana:relative-principalization-from-core,ana:blowups,ana:strict-transform-composition}
\source{transcendental-nakamaye-copy4-20260911:lem:relative-principalization}
\lean{TranscendentalNakamaye.Pending.HironakaControlledTransform.assembly}
\leanok
Let \(Y\), \(W\), \(k\), \(\mathfrak a\) be as in Lemma~\ref{hir:h5a-first-stage}. If the marked ideals \((Y,W,\emptyset,\mathfrak a,1)\) and \((Y_t,Y_t,\emptyset,\mathfrak a_t,1)\)
have multiple test blow-ups with empty final cosupport, then the factorization data of Lemma~\ref{ana:relative-principalization-from-core} exist for \(Y\), \(W\), \(k\),
\(\mathfrak a\): \(f=f_A\circ g\) is adapted to \(W\), its strict transform \(W'\) is connected and regular of dimension \(k\), and with \(H=g^*H_A\) and \(H'\) as in
Lemma~\ref{hir:h5a-second-stage}, \(f^*\mathcal I_W=\mathcal I_H\cdot\mathcal I_{W'}\) and \(f^*\mathfrak a=\mathcal I_H\cdot\mathcal I_{H'}\).
\end{lemma}
\begin{proof}
\leanok
By Lemma~\ref{hir:h5a-controlled-transform}, \(\mathfrak a_t\) is the ideal of the transform of \((Y,Y,\emptyset,\mathfrak a,1)\) along \(T_A\), so this is
Theorem~\ref{mc:factorization-assembly} with \(T_1=T_A\) and \(T_2=T_B\); compactness of \(Y\) is not used. In words: by Theorem~\ref{mc:factorization-tower} and
Lemma~\ref{mc:factorization-second} (the parts of Lemmas~\ref{hir:h5a-first-stage} and~\ref{hir:h5a-second-stage} that do not use compactness), the centres of \(g\) are
disjoint from the strict transform \(W_t\), so \(f=f_A\circ g\) is adapted to \(W\) and its strict transform is \(W'=g^{-1}(W_t)\) (Lemma~\ref{ana:strict-transform-composition}),
connected and regular of dimension \(k\). Pull back the identities of the first stage by \(g\) and use the second: \(f^*\mathfrak a=g^*\mathcal I_{H_A}\cdot g^*\mathfrak a_t=
\mathcal I_H\cdot\mathcal I_{H'}\) and \(f^*\mathcal I_W=\mathcal I_H\cdot g^*\mathcal I_{W_t}=\mathcal I_H\cdot\mathcal I_{W'}\); \(\operatorname{supp}H=g^{-1}(\operatorname{supp}H_A)\)
lies in \(g^{-1}(\operatorname{Exc}f_A)\subset\operatorname{Exc}f\); and \(H'\) is disjoint from \(W'\).
\end{proof}

\begin{lemma}[Relative principalization from the resolution of marked ideals]
\label{hir:factorization-from-resolution}
\uses{hir:canonical-resolution,hir:test-tower-basic,hir:canonical-centres,hir:invariant-driver-seq,hir:centre-function-seq,mc:hironaka-from-seq-invariant,mc:factorization-assembly,mc:factorization-ambient,mc:factorization-domain}
\source{transcendental-nakamaye-copy4-20260911:lem:relative-principalization,bierstone-milman08-functoriality}
\lean{TranscendentalNakamaye.Pending.RelativePrincipalization.exists_factorizationData}
\leanok
The geometric core of Lemma~\ref{ana:relative-principalization} holds: for \(Y\), \(W\), \(k\) and \(\mathfrak a\) as there, there is a sequence of
blowups adapted to \(W\) with the factorization data of Lemma~\ref{ana:relative-principalization-from-core}.
\end{lemma}
\begin{proof}
\leanok
By Theorem~\ref{hir:canonical-centres} there is a resolution invariant in dimension \(\le n\), \(n=\dim Y\), in the sense of Definition~\ref{hir:centre-function-seq}, and
Theorem~\ref{mc:hironaka-from-seq-invariant}(2) with \(j=n\) gives the factorization data: the driver (Theorem~\ref{hir:invariant-driver-seq}) resolves the two marked
ideals below, which lie in the domain in dimension \(\le n\) (Lemma~\ref{mc:factorization-domain}), and Theorem~\ref{mc:factorization-assembly} turns the two resolutions into the
factorization data. The second marked ideal uses the ambient controlled transform \(\mathfrak a_t\) on the whole of \(Y_t\), not the ideal of the transform of the
marked ideal along \(W\), which is only determined along \(W_t\) (Lemma~\ref{mc:factorization-ambient}). The design Lemmas~\ref{hir:h5a-controlled-transform},
\ref{hir:h5a-first-stage}, \ref{hir:h5a-second-stage} and~\ref{hir:h5a-assembly} state the same argument with the controlled transforms written as colon ideals. In words:
Two applications of Theorem~\ref{hir:canonical-resolution} with \(d=1\) and \(E=\emptyset\): first to \((Y,W,\emptyset,\mathfrak a,1)\), whose centres
lie in the cosupport, a proper analytic subset of the current strict transform of \(W\), so that the sequence is adapted to \(W\); then to
\((Y_t,Y_t,\emptyset,\mathfrak a_t,1)\) for the transform \(\mathfrak a_t\), which is the unit ideal near the strict transform \(W_t\), so that its centres are
disjoint from \(W_t\). The divisor \(H\) is the total transform of the exceptional divisors of the first stage and \(H'\) is the divisor of the
second stage; the two equalities of ideals follow from the definition of the transform and the emptiness of the final cosupports.
\end{proof}

\begin{lemma}[Embedded resolution from the canonical resolution of marked ideals]
\label{hir:embedded-from-resolution}
\uses{hir:canonical-resolution,hir:canonical-centres,hir:test-tower-basic,hir:centre-function,hir:canonical-segments,hir:smooth-centre,bcalc:smooth-centre-locus,mc:smooth-part-invariant,bcalc:restrict-empty-tower,ana:cartan-coherence,mc:embedded-from-invariant,mc:embedded-from-first-centre,mc:first-centre-exists,mc:first-centre-smooth-part,mc:subvariety-thin,mc:subvariety-full-dimension,hir:centre-function-seq,mc:hironaka-from-seq-invariant}
\source{transcendental-nakamaye-copy4-20260911:lem:resolution,bierstone-milman08-functoriality,wlodarczyk05-hironaka}
\lean{TranscendentalNakamaye.Pending.EmbeddedResolution.exists_embeddedResolution_smoothData}
\leanok
The resolution input of Lemma~\ref{ana:embedded-resolution} holds: for a closed irreducible analytic subvariety \(V\) of a compact complex
manifold \(X\) there is a finite composition of blowups with smooth centres whose exceptional image lies in \(V_{\mathrm{sing}}\) and whose
strict transform of \(V\) is regular of dimension \(\dim V\) at each of its points.
\end{lemma}
\begin{proof}
\leanok
Let \(\operatorname{inv}\) be the resolution invariant of Theorem~\ref{hir:canonical-centres} in dimension \(\le n\) (Definition~\ref{hir:centre-function-seq}). The ideal
\(\mathcal I_V\) is coherent by Cartan's theorem (Theorem~\ref{ana:cartan-coherence}), so Theorem~\ref{mc:hironaka-from-seq-invariant}(3) applies (the argument of
Theorem~\ref{mc:embedded-from-invariant} for this form of resolution invariant); its proof runs as follows.
If \(\dim V=n\), then \(V\) is a connected component of \(X\), hence smooth, and the identity works (Lemma~\ref{mc:subvariety-full-dimension}). Otherwise
\(\mathcal I_V\) vanishes identically on no nonempty open set and \((X,X,\emptyset,\mathcal I_V,1)\) lies in the domain (Lemma~\ref{mc:subvariety-thin}); take its canonical
resolution (Theorem~\ref{hir:canonical-resolution}). Over \(U=X\setminus V_{\mathrm{sing}}\) the marked ideal is \((U,U,\emptyset,\mathcal I_{V_{\mathrm{reg}}},1)\) with
\(V_{\mathrm{reg}}\) a closed connected submanifold of \(U\) of pure codimension \(\operatorname{codim}V\) (\(V\) is irreducible). Some centre meets the preimage of \(U\), since
\(V_{\mathrm{reg}}\ne\emptyset\) lies in the cosupport until it is blown up; let \(C_r\) be the first one (Lemma~\ref{mc:first-centre-exists}). The centres before \(C_r\) lie
over \(V_{\mathrm{sing}}\), so the restriction to \(U\) of the segment before \(C_r\) is a canonical segment of \((U,U,\emptyset,\mathcal I_{V_{\mathrm{reg}}},1)\)
(Lemma~\ref{hir:canonical-segments}(1)) all of whose centres are empty (Lemma~\ref{bcalc:restrict-empty-tower}). By Definition~\ref{hir:centre-function}(5) and
Lemma~\ref{hir:smooth-centre}, the invariant at stage \(r\) is constant, say \(\lambda\), on the cosupport over \(U\), which is the preimage of \(V_{\mathrm{reg}}\). As \(C_r\)
is the maximum locus and meets this set, \(\lambda\) is the maximum, so \(C_r\) contains the whole cosupport over \(U\) (these two steps are
Lemma~\ref{mc:smooth-part-invariant}). Hence \(C_r\) contains the strict transform \(V_r\) of \(V\) as a union of components of the smooth \(C_r\) (identity principle:
\(V_r\) is irreducible of dimension \(\dim V\) and contains the points over \(V_{\mathrm{reg}}\); Lemma~\ref{mc:first-centre-smooth-part}), so \(V_r\) is regular of dimension
\(\dim V\), and the exceptional image of the first \(r\) blowups lies in \(V_{\mathrm{sing}}\). Stop before the \(r\)-th blowup (Theorem~\ref{mc:embedded-from-first-centre}).
(This route uses only the canonical segments of the one resolution of \((X,X,\emptyset,\mathcal I_V,1)\), not a resolution of the restriction to the non-compact \(U\).)
So the lemma rests only on the existence of a resolution invariant in dimension \(\le n\) (Theorem~\ref{hir:canonical-centres}).
\end{proof}

\subsection{Principalization of a coherent ideal}

The two resolution statements above are the forms the development needs. The canonical resolution gives more; the first of these
stronger forms is the principalization of an arbitrary coherent ideal on a compact manifold, with a normal-crossings exceptional
divisor (Włodarczyk's ``(1) \(\Rightarrow\) (2)'' of his \S2.3; Bierstone--Milman's principalization theorem).

\begin{theorem}[Principalization of a coherent ideal]
\label{hir:principalization}
\uses{hir:canonical-centres,hir:centre-function-seq,hir:invariant-driver-seq,hir:geo-transform,hir:geo-transform-ambient,hir:marked-ideal,hir:snc,hir:test-tower,hir:wellformed-tower,bcalc:inv-domain-stable,bcalc:tower-exceptional,bcalc:exceptional-divisor,bcalc:isomorphism-off-centre,ana:blowup-sequence-fibres,ana:blowups,mc:factorization-domain}
\source{wlodarczyk05-hironaka,bierstone-milman08-functoriality}
\lean{TranscendentalNakamaye.Pending.HironakaPrincipalization.isClosed_of_mem_E_geoTransformAlong,TranscendentalNakamaye.Pending.HironakaPrincipalization.zeroLocus_left_subset_mul,TranscendentalNakamaye.Pending.HironakaPrincipalization.support_add,TranscendentalNakamaye.Pending.HironakaPrincipalization.support_exceptionalDivisor_eq_biUnion,TranscendentalNakamaye.Pending.HironakaPrincipalization.exceptionalLocus_subset_support,TranscendentalNakamaye.Pending.HironakaPrincipalization.isInvDomain_univ,TranscendentalNakamaye.Pending.HironakaPrincipalization.exists_tower_principalization,TranscendentalNakamaye.Pending.HironakaPrincipalization.exists_principalization,TranscendentalNakamaye.Pending.HironakaPrincipalization.exists_principalization_of_connected}
\leanok
Let \(X\) be a compact complex manifold of dimension \(n\) and \(\mathfrak a\subset\mathcal O_X\) a coherent ideal such that at every point \(x\) some germ of
\(\mathfrak a_x\) does not vanish identically near \(x\) (on a connected \(X\): \(\mathfrak a\ne0\)). There are a finite composition of blowups with smooth
centres \(f\colon X'\to X\), an effective divisor \(D\) on \(X'\) and a divisor \(E=(E_1,\dots,E_r)\) on \(X'\) with simple normal crossings such that
\[
 f^*\mathfrak a=\mathcal O_{X'}(-D),\qquad \operatorname{Supp}D=E_1\cup\dots\cup E_r=f^{-1}(Z(\mathfrak a)),
\]
and \(f\) is an isomorphism over \(X\setminus Z(\mathfrak a)\). More precisely, \(f\) can be taken to be a multiple test blow-up \(T\) of the marked ideal
\((X,X,\emptyset,\mathfrak a,1)\) with empty final cosupport, \(D\) its exceptional divisor \(E_T\) (Definition~\ref{bcalc:tower-exceptional}) and \(E\) the divisor of
the transform \((X,X,\emptyset,\mathfrak a,1)_T\).
\end{theorem}
\begin{proof}
\leanok
The marked ideal \(\underline{\mathcal I}=(X,X,\emptyset,\mathfrak a,1)\) lies in the domain in dimension \(\le n\) (Definition~\ref{hir:centre-function-seq}): \(X\) is a
closed submanifold of itself of dimension \(n\), the empty divisor has simple normal crossings and is transverse to it (Definition~\ref{hir:snc}), \(\mathfrak a\) is
coherent and nonzero at every point, and the mark is \(1\) (on a connected \(X\), \(\mathfrak a\ne0\) gives the nonvanishing, Lemma~\ref{mc:factorization-domain}).
The resolution invariant of Theorem~\ref{hir:canonical-centres} in dimension \(\le n\) and the driver of Theorem~\ref{hir:invariant-driver-seq} give a canonical
segment \(T\) of \(\underline{\mathcal I}\) which is a multiple test blow-up with \(\operatorname{cosupp}\underline{\mathcal I}_T=\emptyset\). With \(N=X\) and mark one,
\(T^*\mathfrak a=\mathcal O(-E_T)\cdot\mathfrak a_T\) with \(\mathfrak a_T\) the ideal of \(\underline{\mathcal I}_T\), whose zero locus is the empty cosupport, so
\(\mathfrak a_T\) is the unit ideal and \(T^*\mathfrak a=\mathcal O(-E_T)\) (Lemma~\ref{hir:geo-transform-ambient}). The transforms along \(T\) stay in the domain
(Proposition~\ref{bcalc:inv-domain-stable}, Lemma~\ref{hir:wellformed-tower}), so the divisor \(E\) of \(\underline{\mathcal I}_T\) has simple normal crossings.
Its members are the strict transforms \(\overline{\pi^{-1}(H\setminus C)}\) of the earlier members and the preimages \(\pi^{-1}(C)\) of the centres
(Definition~\ref{hir:geo-transform}); by induction on \(T\), \(\operatorname{Supp}E_T=\bigcup E\): after a blowup \(\pi\) along \(C\),
\(\operatorname{Supp}(\pi^*E_T+D_\pi)=\pi^{-1}(\operatorname{Supp}E_T)\cup\pi^{-1}(C)\) (Lemma~\ref{bcalc:exceptional-divisor}), a point of \(\pi^{-1}(H)\) lies in
\(\pi^{-1}(C)\) or in \(\pi^{-1}(H\setminus C)\), and \(\overline{\pi^{-1}(H\setminus C)}\subset\pi^{-1}(H)\) since \(H\) is closed. The support of \(\mathcal O(-E_T)=T^*\mathfrak a\) is
\(Z(T^*\mathfrak a)=T^{-1}(Z(\mathfrak a))\). Finally the exceptional locus of a blowup lies over its centre (Lemma~\ref{bcalc:isomorphism-off-centre}) and
\(\operatorname{Exc}(g\circ\pi)=\operatorname{Exc}(\pi)\cup\pi^{-1}\operatorname{Exc}(g)\), so by induction \(\operatorname{Exc}(T)\subset\operatorname{Supp}E_T=T^{-1}(Z(\mathfrak a))\);
a blowup sequence of a compact manifold whose exceptional locus maps into \(Z(\mathfrak a)\) is an isomorphism over its complement
(Lemma~\ref{ana:blowup-sequence-fibres}).
\end{proof}

\subsection{Embedded resolution with a normal-crossings exceptional divisor}

\begin{theorem}[Embedded resolution with a normal-crossings exceptional divisor]
\label{hir:embedded-resolution-snc}
\uses{hir:embedded-from-resolution,hir:principalization,hir:canonical-centres,hir:centre-function-seq,hir:invariant-driver-seq,hir:transform,hir:geo-transform,hir:snc,hir:wellformed-tower,bcalc:inv-domain-stable,bcalc:exceptional-divisor,bcalc:normal-crossings-local,ana:regular-component,ana:cartan-coherence,ana:strict-transform-sequence,ana:blowup-sequence-fibres,ana:blowups,mc:hironaka-from-seq-invariant,mc:embedding-properties-smooth-centre,mc:embedded-from-first-centre,mc:first-centre-exists,mc:first-centre-smooth-part,mc:subvariety-thin,mc:subvariety-full-dimension}
\source{wlodarczyk05-hironaka,bierstone-milman08-functoriality}
\lean{TranscendentalNakamaye.Pending.HironakaEmbeddedResolution.eventually_mem_connectedComponentIn_iff,TranscendentalNakamaye.Pending.HironakaEmbeddedResolution.hasNormalCrossings_connectedComponentIn,TranscendentalNakamaye.Pending.HironakaEmbeddedResolution.support_exceptionalDivisor_subset_of_centresAvoid,TranscendentalNakamaye.Pending.HironakaEmbeddedResolution.exists_tower_embeddedResolution,TranscendentalNakamaye.Pending.HironakaEmbeddedResolution.exists_embeddedResolution}
\leanok
Let \(X\) be a compact complex manifold and \(V\subset X\) a closed irreducible analytic subvariety. There are a finite composition of blowups with smooth
centres \(f\colon X'\to X\) and a divisor \(E=(E_1,\dots,E_r)\) on \(X'\) with simple normal crossings such that
\(\operatorname{Exc}(f)\subset E_1\cup\dots\cup E_r\subset f^{-1}(V_{\mathrm{sing}})\), \(f\) is an isomorphism over \(X\setminus V_{\mathrm{sing}}\), the strict transform
\(V'\) of \(V\) is a closed connected submanifold of dimension \(\dim V\), and \(V'\) has only normal crossings with \(E\).
\end{theorem}
\begin{proof}
\leanok
If \(\dim V=n\), \(V\) is a connected component of \(X\) (Lemma~\ref{mc:subvariety-full-dimension}); take \(f=\mathrm{id}\) and \(E=\emptyset\). Otherwise
\(\mathcal I_V\) is coherent (Theorem~\ref{ana:cartan-coherence}) and \(\underline{\mathcal I}=(X,X,\emptyset,\mathcal I_V,1)\) lies in the domain in dimension
\(\le n\) (Lemma~\ref{mc:subvariety-thin}). Let \(T\) be a canonical segment of the resolution invariant of Theorem~\ref{hir:canonical-centres} with empty final
cosupport (Theorem~\ref{hir:invariant-driver-seq}), \(C_r\) its first centre meeting the preimage of \(U=X\setminus V_{\mathrm{sing}}\)
(Lemma~\ref{mc:first-centre-exists}) and \(T'\) the part of \(T\) before \(C_r\). As in the proof of Lemma~\ref{hir:embedded-from-resolution}
(Theorem~\ref{mc:embedded-from-first-centre}), the strict transform \(V'\) of \(V\) by \(f=T'\) is the connected component of \(C_r\) through a point over \(U\),
regular of dimension \(\dim V\), and \(f(\operatorname{Exc}f)\subset V_{\mathrm{sing}}\); it is connected and closed (Lemma~\ref{ana:strict-transform-sequence}), and \(f\) is an
isomorphism over \(U\) (Lemma~\ref{ana:blowup-sequence-fibres}). Take for \(E\) the divisor of the transform \(\underline{\mathcal I}_{T'}\); it has simple normal crossings,
the transforms staying in the domain (Proposition~\ref{bcalc:inv-domain-stable}). As in the proof of Theorem~\ref{hir:principalization}, \(\bigcup E\) is the support of the
exceptional divisor \(E_{T'}\) (Definition~\ref{bcalc:tower-exceptional}) and contains \(\operatorname{Exc}(f)\); this support is the union of the total transforms of the
preimages of the centres of \(T'\) (Lemma~\ref{bcalc:exceptional-divisor}), and every centre of \(T'\) misses the preimage of \(U\), so \(\bigcup E\subset f^{-1}(V_{\mathrm{sing}})\).
Finally \(C_r\) is an admissible centre of \(\underline{\mathcal I}_{T'}\) (Corollary~\ref{mc:embedding-properties-smooth-centre}; Definition~\ref{hir:transform}): a closed
submanifold with only normal crossings with \(E\). Every point of \(C_r\) is regular, so its connected components are the components of its regular locus, open in
\(C_r\) (Proposition~\ref{ana:regular-component}); \(V'\) agrees with \(C_r\) near each of its points, and normal crossings are a local condition
(Lemma~\ref{bcalc:normal-crossings-local}). So \(V'\) has only normal crossings with \(E\).
\end{proof}
The divisor \(E\) is the support of the exceptional divisor of the tower and lies over \(V_{\mathrm{sing}}\); in fact
\(\operatorname{Exc}(f)=E_1\cup\dots\cup E_r=f^{-1}(V_{\mathrm{sing}})\) (Theorem~\ref{hir:embedded-resolution-exceptional}).

\subsection{Local principalization on a non-compact manifold}

The canonical process needs a maximum of the invariant on the cosupport, which a non-compact base need not provide. Over a neighbourhood of a compact set
it is run on shrinking open sets (Bierstone--Milman, \S13: the algorithm over relatively compact open subsets).

\begin{theorem}[Local principalization]
\label{hir:principalization-local}
\uses{hir:principalization,hir:canonical-centres,hir:centre-function-seq,hir:invariant-driver-seq,hir:invariant-driver,hir:resolution-driver,hir:invariant-order,hir:canonical-segments,hir:open-restriction-tower,hir:giraud-along,hir:tower-isomorphism,mc:tower-isomorphism-transform,mc:hironaka-from-seq-invariant,mc:embedding-properties-smooth-centre,bcalc:restrict-nested,bcalc:inv-domain-stable,bcalc:transform-off-centre,bcalc:exceptional-divisor,bcalc:isomorphism-off-centre,ana:blowup-compact,ana:blowup-surjective,ana:blowup-sequence-fibres,ana:blowups}
\source{bierstone-milman97-invent-math:page-0010,bierstone-milman97-invent-math:page-0009,bierstone-milman08-functoriality,wlodarczyk05-hironaka}
\lean{TranscendentalNakamaye.Pending.HironakaLocalPrincipalization.exists_mem_nhds_isCompact_preimage,TranscendentalNakamaye.Pending.HironakaLocalPrincipalization.isProperMap_of_isBlowupAlong,TranscendentalNakamaye.Pending.HironakaLocalPrincipalization.isCompact_preimage_of_isBlowupAlong,TranscendentalNakamaye.Pending.HironakaLocalPrincipalization.surjective_of_isBlowupAlong,TranscendentalNakamaye.Pending.HironakaLocalPrincipalization.IsBlowupSequence.surjective,TranscendentalNakamaye.Pending.HironakaLocalPrincipalization.isIsomorphismOver_compl_of_image_exceptionalLocus_subset,TranscendentalNakamaye.Pending.HironakaLocalPrincipalization.isProperMap_proj,TranscendentalNakamaye.Pending.HironakaLocalPrincipalization.isCanonical_restrict,TranscendentalNakamaye.Pending.HironakaLocalPrincipalization.pullback_symm_nestedDiffeo,TranscendentalNakamaye.Pending.HironakaLocalPrincipalization.exists_local_resolution_of_disjoint,TranscendentalNakamaye.Pending.HironakaLocalPrincipalization.exists_shrink_step,TranscendentalNakamaye.Pending.HironakaLocalPrincipalization.exists_resolution_of_nested,TranscendentalNakamaye.Pending.HironakaLocalPrincipalization.exists_local_resolution_of_lt,TranscendentalNakamaye.Pending.HironakaLocalPrincipalization.exists_local_resolution,TranscendentalNakamaye.Pending.HironakaLocalPrincipalization.exists_tower_principalization_local,TranscendentalNakamaye.Pending.HironakaLocalPrincipalization.exists_principalization_local}
\leanok
Let \(X\) be a complex manifold of dimension \(n\), not necessarily compact, \(K\subset X\) compact, and \(\mathfrak a\subset\mathcal O_X\) a coherent ideal such that at
every point \(x\) some germ of \(\mathfrak a_x\) does not vanish identically near \(x\). There are an open neighbourhood \(U\) of \(K\), a finite composition of blowups
with smooth centres \(f\colon U'\to U\), an effective divisor \(D\) on \(U'\) and a divisor \(E=(E_1,\dots,E_r)\) on \(U'\) with simple normal crossings such that
\[
 f^*(\mathfrak a|_U)=\mathcal O_{U'}(-D),\qquad \operatorname{Supp}D=E_1\cup\dots\cup E_r=f^{-1}(Z(\mathfrak a)\cap U),
\]
and \(f\) is an isomorphism over \(U\setminus Z(\mathfrak a)\). More generally, for a resolution invariant in dimension \(\le k\) with well-ordered values
(Definition~\ref{hir:centre-function-seq}) and a marked ideal \(\underline{\mathcal I}\) of its domain on any complex manifold \(X\), there are an open neighbourhood
\(U\) of \(K\) and a multiple test blow-up of \(\underline{\mathcal I}|_U\) with empty final cosupport.
\end{theorem}
\begin{proof}
\leanok
\emph{Blowups are proper and surjective.} Every point \(y\) of the base of a blowup \(\pi\) has a neighbourhood with compact preimage (off the centre by the section over
the complement of the centre, on the centre by a flattening chart and the monomial lifts; the local half of Lemma~\ref{ana:blowup-compact}). An ultrafilter whose image
converges to \(y\) then lives on a compact set, so it converges, to a point over \(y\): \(\pi\) is proper, hence closed. Its image is closed and contains the dense
complement of the centre (Lemma~\ref{bcalc:isomorphism-off-centre}), so \(\pi\) is surjective; so is every blowup sequence, and the projection of a tower is proper. With
surjectivity in place of compactness, a blowup sequence whose exceptional locus maps into a set \(A\) is an isomorphism over \(X\setminus A\)
(as in Lemma~\ref{ana:blowup-sequence-fibres}).

\emph{The local driver.} Let \(T\) be a canonical segment of \(\underline{\mathcal I}\) all of whose values over \(K\) are \(<\mu\) (\(\mu\) in the values with a top added). We
induct on \(\mu\), the values being well-ordered. If no point of \(\operatorname{cosupp}\underline{\mathcal I}_T\) lies over \(K\), the cosupport has closed image (the
projection is closed, the cosupport closed, Lemma~\ref{hir:giraud-along}); over its complement \(U\) the restricted tower \(T|_U\) (Lemma~\ref{hir:open-restriction-tower})
has empty cosupport (Lemma~\ref{hir:canonical-segments}). Otherwise the part of the cosupport over \(K\) is compact, and the invariant, upper semicontinuous on the
cosupport, takes a maximum \(\mu_0<\mu\) there (Lemma~\ref{hir:invariant-order}). The points off the cosupport or with value \(\le\mu_0\) form an open set \(W\), and since the
projection is closed, \(O=X\setminus T(W^c)\) is an open neighbourhood of \(K\) whose preimage lies in \(W\). The restriction \(T|_O\) is a canonical segment of
\(\underline{\mathcal I}|_O\) (Lemma~\ref{hir:canonical-segments}(1), from properties (5) and (1)) with the same values (property (5)); on its cosupport the values are \(\le\mu_0\),
with \(\mu_0\) attained, so its maximum locus \(C\) is the level set \(\mu_0\), an admissible centre with empty interior (Corollary~\ref{mc:embedding-properties-smooth-centre},
Theorem~\ref{hir:invariant-driver}). Blow it up. The new cosupport lies over the old one (Lemma~\ref{bcalc:transform-off-centre}); by property (3) a point over \(C\) has a
value \(<\mu_0\), and a point off \(C\) keeps its value, which is \(\le\mu_0\) and \(\ne\mu_0\). So every value over \(K\) is now \(<\mu_0\), and the induction hypothesis,
applied over \(O\) to \(\underline{\mathcal I}|_O\) and the compact \(K\subset O\), gives an open \(U_1\subset O\) containing \(K\) and a resolution of
\((\underline{\mathcal I}|_O)|_{U_1}\). Seen as an open set \(U\) of \(X\), \(U_1\) is biholomorphic to \(U\) by the identity on points (Lemma~\ref{bcalc:restrict-nested}), the
transported tower (Lemma~\ref{hir:tower-isomorphism}) resolves the transported marked ideal \(\underline{\mathcal I}|_U\) (Lemma~\ref{mc:tower-isomorphism-transform}), and the
transports compose because the ideals are coherent. The domain is stable under restriction (Proposition~\ref{bcalc:inv-domain-stable}); the resolution invariant used is that of
Theorem~\ref{hir:canonical-centres}, through Lemma~\ref{mc:hironaka-from-seq-invariant}. Start with the empty tower and \(\mu=\top\).

\emph{Principalization.} Apply the local driver to \((X,X,\emptyset,\mathfrak a,1)\), whose restriction to \(U\) is \((U,U,\emptyset,\mathfrak a|_U,1)\), and conclude as in the
proof of Theorem~\ref{hir:principalization}, with the isomorphism off \(Z(\mathfrak a)\) from the first paragraph in place of compactness.
\end{proof}
The neighbourhood \(U\) depends on \(K\); for a relatively compact open \(X_0\subset X\) take \(K=\overline{X_0}\). A locally finite principalization of all of \(X\)
(Bierstone--Milman, Theorem~13.4) is not claimed.

\subsection{Local embedded resolution on a non-compact manifold}

\begin{theorem}[Local embedded resolution]
\label{hir:embedded-resolution-local}
\uses{hir:embedded-resolution-snc,hir:principalization-local,hir:embedded-from-resolution,hir:canonical-centres,hir:centre-function-seq,hir:invariant-driver-seq,hir:canonical-segments,hir:tower-isomorphism,hir:open-restriction-tower,hir:snc,hir:wellformed-tower,mc:tower-isomorphism-transform,mc:hironaka-from-seq-invariant,mc:embedding-properties-smooth-centre,mc:first-centre-exists,mc:first-centre-smooth-part,mc:smooth-part-invariant,mc:subvariety-thin,mc:subvariety-full-dimension,bcalc:restrict-nested,bcalc:restrict-empty-tower,bcalc:inv-domain-stable,bcalc:exceptional-divisor,bcalc:tower-exceptional,bcalc:normal-crossings-local,ana:blowup-exceptional-image,ana:blowup-surjective,ana:closure-difference-analytic,ana:analytic-preimage,ana:strict-transform-composition,ana:strict-transform-blowup,ana:strict-transform-sequence,ana:regular-transport,ana:regular-component,ana:cartan-coherence,ana:blowup-sequence-fibres,ana:blowups}
\source{bierstone-milman97-invent-math:page-0009,bierstone-milman08-functoriality,wlodarczyk05-hironaka}
\lean{TranscendentalNakamaye.Pending.HironakaEmbeddedResolutionLocal.isClosed_image_exceptionalLocus,TranscendentalNakamaye.Pending.HironakaEmbeddedResolutionLocal.isComplexAnalyticSubset_image_exceptionalLocus,TranscendentalNakamaye.Pending.HironakaEmbeddedResolutionLocal.isComplexAnalyticSubset_strictTransform_of_isBlowupAlong,TranscendentalNakamaye.Pending.HironakaEmbeddedResolutionLocal.isComplexAnalyticSubset_strictTransform_proj,TranscendentalNakamaye.Pending.HironakaEmbeddedResolutionLocal.exists_strictTransform_subvariety_of_isBlowupAlong,TranscendentalNakamaye.Pending.HironakaEmbeddedResolutionLocal.exists_strictTransform_subvariety_of_isBlowupSequence,TranscendentalNakamaye.Pending.HironakaEmbeddedResolutionLocal.isConnected_strictTransform_of_isBlowupSequence,TranscendentalNakamaye.Pending.HironakaEmbeddedResolutionLocal.strictTransform_eq_connectedComponentIn,TranscendentalNakamaye.Pending.HironakaEmbeddedResolutionLocal.strictTransform_eq_connectedComponentIn_maxLocus,TranscendentalNakamaye.Pending.HironakaEmbeddedResolutionLocal.not_centresAvoid_of_cosupp_eq_empty,TranscendentalNakamaye.Pending.HironakaEmbeddedResolutionLocal.exists_embedded_of_isCanonical,TranscendentalNakamaye.Pending.HironakaEmbeddedResolutionLocal.isCanonical_of_isTowerIso,TranscendentalNakamaye.Pending.HironakaEmbeddedResolutionLocal.exists_local_resolution_isCanonical_of_disjoint,TranscendentalNakamaye.Pending.HironakaEmbeddedResolutionLocal.exists_canonical_resolution_of_nested,TranscendentalNakamaye.Pending.HironakaEmbeddedResolutionLocal.exists_local_resolution_isCanonical_of_lt,TranscendentalNakamaye.Pending.HironakaEmbeddedResolutionLocal.exists_local_resolution_isCanonical,TranscendentalNakamaye.Pending.HironakaEmbeddedResolutionLocal.subset_closure_inter_regularOpen,TranscendentalNakamaye.Pending.HironakaEmbeddedResolutionLocal.strictTransform_of_centresAvoid,TranscendentalNakamaye.Pending.HironakaEmbeddedResolutionLocal.mem_maxLocus_iff_of_centresAvoid_of_pureDim,TranscendentalNakamaye.Pending.HironakaEmbeddedResolutionLocal.exceptionalDivisor_facts_of_centresAvoid,TranscendentalNakamaye.Pending.HironakaEmbeddedResolutionLocal.exists_tower_embeddedResolution_local,TranscendentalNakamaye.Pending.HironakaEmbeddedResolutionLocal.exists_embeddedResolution_local}
\leanok
Let \(X\) be a complex manifold of dimension \(n\), not necessarily compact, \(V\subset X\) a closed irreducible analytic subvariety and \(K\subset X\) compact.
There are an open neighbourhood \(U\) of \(K\), a finite composition of blowups with smooth centres \(f\colon U'\to U\) and a divisor \(E=(E_1,\dots,E_r)\) on
\(U'\) with simple normal crossings such that \(\operatorname{Exc}(f)\subset E_1\cup\dots\cup E_r\subset f^{-1}(V_{\mathrm{sing}})\), \(f\) is an isomorphism over
\(U\setminus V_{\mathrm{sing}}\), the strict transform \(V'\) of \(V\cap U\) is a closed submanifold of \(U'\), of dimension \(\dim V\) at each of its points, and \(V'\)
has only normal crossings with \(E\). Moreover, for a resolution invariant in dimension \(\le k\) with well-ordered values
(Definition~\ref{hir:centre-function-seq}), a marked ideal \(\underline{\mathcal I}\) of its domain on any complex manifold \(X\) and a compact \(K\subset X\), the
neighbourhood \(U\) of Theorem~\ref{hir:principalization-local} can be chosen together with a \emph{canonical segment} of \(\underline{\mathcal I}|_U\) with empty final
cosupport.
\end{theorem}
\begin{proof}
\leanok
\emph{Canonical output of the local driver.} Run the local driver of Theorem~\ref{hir:principalization-local}, keeping track of canonicity: the restriction of a canonical
segment to an open set is canonical (Lemma~\ref{hir:canonical-segments}(1)), the blowup of the maximum locus extends a canonical segment, and canonical segments correspond
along isomorphisms of towers over a biholomorphism (Lemma~\ref{hir:canonical-segments}; for any value function with properties (1), (4), (5), (6) of
Definition~\ref{hir:centre-function-seq}, the centres of the transported tower being the preimages of the old ones and the values being transported by property (6),
Definition~\ref{hir:tower-isomorphism}). The last step of the driver brings a resolution over an open set of an open set back to an open set of \(X\) by such an isomorphism
(Lemma~\ref{bcalc:restrict-nested}), so the resolution it produces is a canonical segment.

\emph{Compactness is not needed for the basic facts.} A blowup is proper and surjective (proof of Theorem~\ref{hir:principalization-local}), so the image
\(\pi(\operatorname{Exc}\pi)\) of its exceptional locus is closed, and the proof of Lemma~\ref{ana:blowup-exceptional-image} shows that it is analytic. Hence the strict
transform \(\overline{\pi^{-1}(A)\setminus\pi^{-1}\pi(\operatorname{Exc}\pi)}\) of an analytic set \(A\) is analytic (Proposition~\ref{ana:closure-difference-analytic}), and by
composition (Lemma~\ref{ana:strict-transform-composition}) so is its strict transform along any blowup sequence. The proofs of
Lemmas~\ref{ana:strict-transform-blowup} and~\ref{ana:strict-transform-sequence} and of the identity-principle and stopping steps of
Lemma~\ref{hir:embedded-from-resolution} (Lemmas~\ref{mc:first-centre-smooth-part} and~\ref{mc:first-centre-exists}) then hold on any complex manifold.

\emph{The local statement.} If \(\dim V=n\), \(V\) is a connected component of \(X\) (Lemma~\ref{mc:subvariety-full-dimension}); take \(U=X\), \(f=\mathrm{id}\) and
\(E=\emptyset\). Otherwise \(\mathcal I_V\) is coherent (Theorem~\ref{ana:cartan-coherence}) and \((X,X,\emptyset,\mathcal I_V,1)\) lies in the domain in dimension
\(\le n\) (Lemma~\ref{mc:subvariety-thin}). The first part, with the resolution invariant of Theorem~\ref{hir:canonical-centres}, gives an open \(U\supset K\) and a canonical
segment \(T\) of \((U,U,\emptyset,\mathcal I_{W},1)\), \(W=V\cap U\), with empty final cosupport. Let \(O=U\setminus V_{\mathrm{sing}}\); \(W\cap O\) is a closed submanifold
of \(O\) of pure dimension \(\dim V\), and \(W\subset\overline{W\cap O}\) since the regular points of dimension \(\dim V\) are dense in the irreducible \(V\)
(Proposition~\ref{ana:regular-component}) and \(U\) is open. If every centre of \(T\) misses the preimage of \(O\), the cosupport over \(O\) is the preimage of \(W\cap O\)
(Lemma~\ref{bcalc:restrict-empty-tower}), which is therefore empty, the projection being onto; then \(W=\emptyset\) and \(f=T\) works with an empty strict transform.
Otherwise let \(C\) be the first centre meeting the preimage of \(O\) and \(T'\) the segment before it (Lemma~\ref{mc:first-centre-exists}). As in
Lemma~\ref{mc:smooth-part-invariant}, \(C\) contains, over \(O\), exactly the preimage of \(W\), and it is an admissible centre (Corollary~\ref{mc:embedding-properties-smooth-centre}):
a closed submanifold with only normal crossings with the transformed divisor \(E\).

The strict transform \(S\) of \(W\) by \(T'\) need not be irreducible, so the identity-principle step is replaced by the following argument. Over \(O\) the projection has no
exceptional point, so \(S\) is the preimage of \(W\) there, which is \(C\) there. Every point of \(T'^{-1}(W)\setminus\operatorname{Exc}\) is a limit of points over \(W\cap O\),
the projection being open off the exceptional locus; hence \(S=\overline{C\cap T'^{-1}(O)}\subset C\). Let \(x\in S\) and \(C_x\) the connected component of the closed
submanifold \(C\) through \(x\), a neighbourhood of \(x\) in \(C\) (Proposition~\ref{ana:regular-component}). It meets \(T'^{-1}(O)\), at a point near \(x\); there \(C_x\subset S\)
and \(S\) is analytic, so the identity theorem along \(C_x\) gives \(C_x\subset S\). Thus \(S\) agrees with \(C\) near each of its points, and its dimension at \(x\) is that of
\(C_x\), read off at a point over \(O\), where \(S\) is regular of dimension \(\dim V\) (Lemma~\ref{ana:regular-transport}). Normal crossings are local
(Lemma~\ref{bcalc:normal-crossings-local}), so \(S\) has only normal crossings with \(E\). As in Theorem~\ref{hir:embedded-resolution-snc}, \(E\) has simple normal crossings
(Proposition~\ref{bcalc:inv-domain-stable}), its union is the support of the exceptional divisor of \(T'\) (Definition~\ref{bcalc:tower-exceptional},
Lemma~\ref{bcalc:exceptional-divisor}), which contains the exceptional locus and lies over \(V_{\mathrm{sing}}\) since the centres of \(T'\) miss the preimage of \(O\); and
\(f=T'\) is an isomorphism over \(U\setminus V_{\mathrm{sing}}\) (proof of Theorem~\ref{hir:principalization-local}).
\end{proof}

\subsection{Embedded resolution: the exceptional divisor is the preimage of the singular locus}

\begin{theorem}[Embedded resolution: \(\operatorname{Exc}(f)=\bigcup E=f^{-1}(V_{\mathrm{sing}})\)]
\label{hir:embedded-resolution-exceptional}
\uses{hir:embedded-resolution-snc,hir:embedded-resolution-local,ana:regular-transport,ana:strict-transform-composition,ana:blowups}
\source{bierstone-milman97-invent-math:page-0009,wlodarczyk05-hironaka}
\lean{TranscendentalNakamaye.Pending.HironakaEmbeddedExceptional.LocalHolomorphicEquivAt.symmAt,TranscendentalNakamaye.Pending.HironakaEmbeddedExceptional.regularAtDim_of_regularAtDim_preimage,TranscendentalNakamaye.Pending.HironakaEmbeddedExceptional.preimage_singularLocus_subset_exceptionalLocus,TranscendentalNakamaye.Pending.HironakaEmbeddedExceptional.exceptionalLocus_eq_preimage_singularLocus,TranscendentalNakamaye.Pending.HironakaEmbeddedExceptional.exists_embeddedResolution_exceptional,TranscendentalNakamaye.Pending.HironakaEmbeddedExceptional.exists_embeddedResolution_local_exceptional}
\leanok
Let \(f\colon X'\to X\) be a holomorphic map between complex manifolds of the same dimension and \(V\subset X\) a closed irreducible analytic subvariety whose strict
transform \(V'=\overline{f^{-1}(V)\setminus\operatorname{Exc}(f)}\) is regular of dimension \(\dim V\) at each of its points. Then
\(f^{-1}(V_{\mathrm{sing}})\subset\operatorname{Exc}(f)\), with equality when \(f(\operatorname{Exc}f)\subset V_{\mathrm{sing}}\). In particular the resolutions of
Theorems~\ref{hir:embedded-resolution-snc} and~\ref{hir:embedded-resolution-local} satisfy
\[
 \operatorname{Exc}(f)=E_1\cup\dots\cup E_r=f^{-1}(V_{\mathrm{sing}})\qquad(\text{resp. }f^{-1}(V_{\mathrm{sing}}\cap U)),
\]
with all their other properties.
\end{theorem}
\begin{proof}
\leanok
Let \(y\in f^{-1}(V_{\mathrm{sing}})\) and suppose \(y\notin\operatorname{Exc}(f)\). The exceptional locus is closed, and on its complement \(V'\) agrees with
\(f^{-1}(V)\); so \(y\in V'\), and \(V'\) is regular of dimension \(\dim V\) at \(y\). The differential of \(f\) at \(y\) is bijective, so in charts \(f\) has a local
holomorphic inverse near \(f(y)\) (inverse function theorem). Composing a flattening chart of \(V'\) at \(y\) with this local inverse gives a flattening chart of \(V\)
at \(f(y)\) of the same dimension: regular points push forward along a local biholomorphism, the converse of Lemma~\ref{ana:regular-transport}. Hence \(f(y)\) is a
regular point of \(V\) of dimension \(\dim V\), i.e.\ \(f(y)\notin V_{\mathrm{sing}}\), a contradiction. If moreover \(f(\operatorname{Exc}f)\subset V_{\mathrm{sing}}\),
then \(\operatorname{Exc}(f)\subset f^{-1}(V_{\mathrm{sing}})\) as well. The resolutions of Theorems~\ref{hir:embedded-resolution-snc}
and~\ref{hir:embedded-resolution-local} (the latter over \(U\), for \(V\cap U\), whose singular points are those of \(V\) in \(U\)) have
\(\operatorname{Exc}(f)\subset\bigcup E\subset f^{-1}(V_{\mathrm{sing}})\) and a strict transform regular of dimension \(\dim V\) at each of its points; with
\(f^{-1}(V_{\mathrm{sing}})\subset\operatorname{Exc}(f)\) the three sets are equal.
\end{proof}
This is the form of Bierstone--Milman, Theorem~1.6: the exceptional divisor is the critical locus of \(f\) and equals \(f^{-1}(V_{\mathrm{sing}})\).

\subsection{Admissible sequences; principalization and embedded resolution with compact support}

\begin{definition}[Admissible blowup sequence]
\label{hir:admissible-sequence}
\uses{hir:snc,hir:test-tower,hir:geo-transform,hir:marked-ideal,ana:blowups,bcalc:transform-off-centre}
\source{bierstone-milman97-invent-math:page-0009,bierstone-milman97-invent-math:page-0010,wlodarczyk05-hironaka}
\lean{TranscendentalNakamaye.Pending.HironakaAdmissibleSequence.IsAdmissibleBlowupSequence,TranscendentalNakamaye.Pending.HironakaAdmissibleSequence.IsAdmissibleBlowupSequence.isBlowupSequence,TranscendentalNakamaye.Pending.HironakaAdmissibleSequence.IsAdmissibleBlowupSequence.mono,TranscendentalNakamaye.Pending.HironakaAdmissibleSequence.isAdmissibleBlowupSequence_proj,TranscendentalNakamaye.Pending.HironakaAdmissibleSequence.centresOver_compl_of_centresAvoid,TranscendentalNakamaye.Pending.HironakaAdmissibleSequence.centresOver_cosupp}
\leanok
Let \(X\) be a complex manifold, \(E_0\) a finite sequence of closed subsets of \(X\) and \(F\subset X\). An \emph{admissible blowup sequence} of \((X,E_0)\) with
centres over \(F\) is a finite composition \(f\colon X_r\xrightarrow{\pi_r}\cdots\xrightarrow{\pi_1}X_0=X\) of blowups \(\pi_{j+1}\) along closed submanifolds
\(C_j\subset X_j\) with empty interior, together with divisors \(E_j\) on \(X_j\), such that \(C_j\) has only normal crossings with \(E_j\)
(Definition~\ref{hir:snc}) and \(C_j\subset(\pi_1\circ\dots\circ\pi_j)^{-1}(F)\), where \(E_{j+1}\) consists of the strict transforms
\(\overline{\pi_{j+1}^{-1}(H\setminus C_j)}\), \(H\in E_j\), followed by \(\pi_{j+1}^{-1}(C_j)\). Its divisor is \(E=E_r\). A multiple test blow-up of a marked ideal
\((X,N,E_0,\mathfrak a,d)\) (Definition~\ref{hir:test-tower}) whose centres lie over \(F\) is an admissible blowup sequence of \((X,E_0)\) with centres over \(F\),
its divisor being that of the transform (Definition~\ref{hir:geo-transform}); the centres of a multiple test blow-up lie over the cosupport of the marked ideal.
\end{definition}

\begin{theorem}[Principalization with a compact zero locus]
\label{hir:principalization-compact-support}
\uses{hir:admissible-sequence,hir:principalization,hir:principalization-local,hir:monomial-case-local,hir:canonical-centres,hir:centre-function-seq,hir:invariant-driver-seq,hir:invariant-order,hir:geo-transform-ambient,hir:giraud-along,mc:hironaka-from-seq-invariant,mc:embedding-properties-smooth-centre,mc:factorization-domain,bcalc:transform-off-centre,bcalc:inv-domain-stable,bcalc:tower-exceptional,bcalc:exceptional-divisor,bcalc:isomorphism-off-centre,ana:blowup-sequence-fibres,ana:blowups}
\source{bierstone-milman97-invent-math:page-0010,bierstone-milman97-invent-math:page-0009,wlodarczyk05-hironaka}
\lean{TranscendentalNakamaye.Pending.HironakaAdmissibleSequence.exists_isGeoResolution_isCanonical_of_isCompact,TranscendentalNakamaye.Pending.HironakaAdmissibleSequence.isNonzeroAlongAt_univ_iff_exists,TranscendentalNakamaye.Pending.HironakaAdmissibleSequence.cosupp_univ_one,TranscendentalNakamaye.Pending.HironakaAdmissibleSequence.exists_principalization_of_isCompact,TranscendentalNakamaye.Pending.HironakaAdmissibleSequence.exists_principalization_admissible}
\leanok
Let \(X\) be a complex manifold, not necessarily compact, and \(\mathfrak a\subset\mathcal O_X\) a coherent ideal, nonzero at every point (some germ of \(\mathfrak a_x\)
does not vanish identically near \(x\)), whose zero locus \(Z(\mathfrak a)\) is compact. There are an admissible blowup sequence \(f\colon X'\to X\) of \((X,\emptyset)\) with
centres over \(Z(\mathfrak a)\) (Definition~\ref{hir:admissible-sequence}), with divisor \(E\), and an effective divisor \(D\) on \(X'\), such that \(E\) has simple normal
crossings, \(f^*\mathfrak a=\mathcal O_{X'}(-D)\), \(\operatorname{Supp}D=\bigcup E=f^{-1}(Z(\mathfrak a))\), and \(f\) is an isomorphism over \(X\setminus Z(\mathfrak a)\).
In particular this holds on every compact \(X\) (Theorem~\ref{hir:principalization}, with the admissibility of the sequence and the position of its centres).
More generally, a resolution invariant in dimension \(\le k\) with well-ordered values (Definition~\ref{hir:centre-function-seq}) resolves every marked ideal of its domain with
compact cosupport, on any complex manifold, by a canonical segment.
\end{theorem}
\begin{proof}
\leanok
The canonical process of Theorem~\ref{hir:invariant-driver-seq} uses the compactness of the base only to find a maximum of the invariant on the cosupport. Along a
multiple test blow-up the cosupport is closed and lies over the original cosupport (Lemmas~\ref{hir:giraud-along} and~\ref{bcalc:transform-off-centre}), and the
projection is proper (proof of Theorem~\ref{hir:principalization-local}); so the cosupports stay compact. The invariant is upper semicontinuous on them, so its maximum
is attained; its maximum locus is an admissible centre with empty interior (Corollary~\ref{mc:embedding-properties-smooth-centre}); after its blowup every value on the new
cosupport is smaller than the old maximum (property (3), Definition~\ref{hir:centre-function-seq}: a point over the centre drops, a point off it keeps a value that is at most
the maximum and different from it). The values being well-ordered, the process stops: this is the compact-cosupport driver of Lemma~\ref{hir:monomial-case-local}, applied to the
resolution invariant of Theorem~\ref{hir:canonical-centres}. The marked ideal \((X,X,\emptyset,\mathfrak a,1)\) lies in the domain (Lemma~\ref{mc:factorization-domain}) and its
cosupport is \(Z(\mathfrak a)\). Conclude as in the proof of Theorem~\ref{hir:principalization}, with the isomorphism off \(Z(\mathfrak a)\) from the proof of
Theorem~\ref{hir:principalization-local} (blowups are proper and surjective). The resolving canonical segment is a multiple test blow-up whose centres lie in the cosupport,
hence over \(Z(\mathfrak a)\): an admissible blowup sequence (Definition~\ref{hir:admissible-sequence}).
\end{proof}

\begin{theorem}[Embedded resolution of a compact subvariety]
\label{hir:embedded-resolution-compact-support}
\uses{hir:admissible-sequence,hir:principalization-compact-support,hir:embedded-resolution-snc,hir:embedded-resolution-local,hir:embedded-resolution-exceptional,hir:embedded-from-resolution,hir:canonical-centres,hir:centre-function-seq,hir:canonical-segments,mc:first-centre-exists,mc:first-centre-smooth-part,mc:smooth-part-invariant,mc:subvariety-thin,mc:subvariety-full-dimension,mc:embedding-properties-smooth-centre,bcalc:restrict-empty-tower,bcalc:inv-domain-stable,bcalc:exceptional-divisor,bcalc:tower-exceptional,bcalc:normal-crossings-local,ana:strict-transform-sequence,ana:regular-component,ana:cartan-coherence,ana:blowup-sequence-fibres,ana:blowups}
\source{bierstone-milman97-invent-math:page-0009,bierstone-milman08-functoriality,wlodarczyk05-hironaka}
\lean{TranscendentalNakamaye.Pending.HironakaAdmissibleSequence.exists_tower_embedded_of_isCanonical,TranscendentalNakamaye.Pending.HironakaAdmissibleSequence.exists_tower_embeddedResolution_of_isCompact,TranscendentalNakamaye.Pending.HironakaAdmissibleSequence.exists_embeddedResolution_of_isCompact,TranscendentalNakamaye.Pending.HironakaAdmissibleSequence.exists_embeddedResolution_admissible,TranscendentalNakamaye.Pending.HironakaAdmissibleSequence.exceptionalLocus_nonempty_of_singular,TranscendentalNakamaye.Pending.HironakaAdmissibleSequence.ne_nil_of_singular}
\leanok
Let \(X\) be a complex manifold, not necessarily compact, and \(V\subset X\) a compact irreducible analytic subvariety. There is an admissible blowup sequence
\(f\colon X'\to X\) of \((X,\emptyset)\) with centres over \(V_{\mathrm{sing}}\) (Definition~\ref{hir:admissible-sequence}), with divisor \(E\), such that \(E\) has simple
normal crossings, \(\operatorname{Exc}(f)=\bigcup E=f^{-1}(V_{\mathrm{sing}})\), \(f\) is an isomorphism over \(X\setminus V_{\mathrm{sing}}\), the strict transform \(V'\) of \(V\)
is a closed connected submanifold of dimension \(\dim V\), and \(V'\) has only normal crossings with \(E\). In particular this holds on every compact \(X\)
(Theorems~\ref{hir:embedded-resolution-snc} and~\ref{hir:embedded-resolution-exceptional}, with the admissibility of the sequence and the position of its centres). If \(V\) is
singular, \(E\neq\emptyset\) and \(f\) is not biholomorphic.
\end{theorem}
\begin{proof}
\leanok
If \(\dim V=n\), \(V\) is a connected component of \(X\) (Lemma~\ref{mc:subvariety-full-dimension}) and smooth; take \(f=\mathrm{id}\), \(E=\emptyset\). Otherwise
\((X,X,\emptyset,\mathcal I_V,1)\) lies in the domain (Theorem~\ref{ana:cartan-coherence}, Lemma~\ref{mc:subvariety-thin}) and its cosupport is the compact \(V\), so
Theorem~\ref{hir:principalization-compact-support} gives a canonical segment with empty final cosupport. Its part before the first centre meeting the preimage of
\(X\setminus V_{\mathrm{sing}}\) is, as in the proof of Theorem~\ref{hir:embedded-resolution-local} (which needs neither compactness nor irreducibility of the strict
transform), a multiple test blow-up whose centres lie over \(V_{\mathrm{sing}}\) and whose strict transform \(V'\) is a closed submanifold, regular of dimension
\(\dim V\), with only normal crossings with the transformed divisor \(E\); \(E\) has simple normal crossings, and its union is the support of the exceptional divisor, which
contains \(\operatorname{Exc}(f)\) and lies over \(V_{\mathrm{sing}}\). So \(f(\operatorname{Exc}f)\subset V_{\mathrm{sing}}\), \(V'\) is connected (the proof of
Lemma~\ref{ana:strict-transform-sequence}, which needs only surjectivity of blowups, not compactness), \(f\) is an isomorphism over \(X\setminus V_{\mathrm{sing}}\) (proof of
Theorem~\ref{hir:principalization-local}), and \(\operatorname{Exc}(f)=\bigcup E=f^{-1}(V_{\mathrm{sing}})\) by Theorem~\ref{hir:embedded-resolution-exceptional}. The segment
is an admissible blowup sequence with centres over \(V_{\mathrm{sing}}\) (Definition~\ref{hir:admissible-sequence}). If \(V_{\mathrm{sing}}\neq\emptyset\), a point of it has a
preimage (blowup sequences are onto), which is exceptional by Theorem~\ref{hir:embedded-resolution-exceptional}; so \(E\ne\emptyset\).
\end{proof}

\begin{theorem}[Principalization and embedded resolution over relatively compact open sets]
\label{hir:resolution-relatively-compact}
\uses{hir:admissible-sequence,hir:principalization-compact-support,hir:embedded-resolution-compact-support,hir:embedded-resolution-local,hir:embedded-resolution-exceptional,hir:principalization-local,hir:canonical-segments,hir:open-restriction-tower,hir:tower-isomorphism,mc:tower-isomorphism-transform,bcalc:restrict-nested,bcalc:inv-domain-stable,ana:regular-transport,ana:blowups}
\source{bierstone-milman97-invent-math:page-0009,bierstone-milman08-functoriality,wlodarczyk05-hironaka}
\lean{TranscendentalNakamaye.Pending.HironakaRelativelyCompact.exists_isCanonical_resolution_restrict,TranscendentalNakamaye.Pending.HironakaRelativelyCompact.exists_isCanonical_resolution_of_subset_isCompact,TranscendentalNakamaye.Pending.HironakaRelativelyCompact.preimage_val_singularLocus_subset_exceptionalLocus,TranscendentalNakamaye.Pending.HironakaRelativelyCompact.exists_principalization_of_subset_isCompact,TranscendentalNakamaye.Pending.HironakaRelativelyCompact.exists_principalization_local_admissible,TranscendentalNakamaye.Pending.HironakaRelativelyCompact.exists_embeddedResolution_of_subset_isCompact,TranscendentalNakamaye.Pending.HironakaRelativelyCompact.exists_embeddedResolution_local_admissible}
\leanok
Let \(X\) be a complex manifold, not necessarily compact, and \(X_0\subset X\) an open set contained in a compact set (for instance a relatively compact open set).
\begin{enumerate}
\item For a coherent ideal \(\mathfrak a\subset\mathcal O_X\), nonzero at every point, there are an admissible blowup sequence \(f\colon X_0'\to X_0\) of
\((X_0,\emptyset)\) with centres over \(Z(\mathfrak a)\) (Definition~\ref{hir:admissible-sequence}), with divisor \(E\), and an effective divisor \(D\) on \(X_0'\), such that
\(E\) has simple normal crossings, \(f^*(\mathfrak a|_{X_0})=\mathcal O_{X_0'}(-D)\), \(\operatorname{Supp}D=\bigcup E=f^{-1}(Z(\mathfrak a)\cap X_0)\), and \(f\) is an
isomorphism over \(X_0\setminus Z(\mathfrak a)\).
\item For a closed irreducible analytic subvariety \(V\subset X\) there is an admissible blowup sequence \(f\colon X_0'\to X_0\) of \((X_0,\emptyset)\) with centres over
\(V_{\mathrm{sing}}\), with divisor \(E\), such that \(E\) has simple normal crossings, \(\operatorname{Exc}(f)=\bigcup E=f^{-1}(V_{\mathrm{sing}}\cap X_0)\), \(f\) is an
isomorphism over \(X_0\setminus V_{\mathrm{sing}}\), and the strict transform of \(V\cap X_0\) is a closed submanifold of \(X_0'\), of dimension \(\dim V\) at each of its
points, with only normal crossings with \(E\).
\end{enumerate}
In particular the local forms of Theorems~\ref{hir:principalization-local} and~\ref{hir:embedded-resolution-local} hold with a relatively compact neighbourhood \(U\) of
\(K\) and with the sequence admissible.
\end{theorem}
\begin{proof}
\leanok
Let \(K\supset X_0\) be compact. The local driver with canonical output (Theorem~\ref{hir:embedded-resolution-local}) gives an open \(U\supset K\) and a canonical segment
\(T\) resolving the restriction to \(U\) of \((X,X,\emptyset,\mathfrak a,1)\), resp.\ of \((X,X,\emptyset,\mathcal I_V,1)\). Restrict \(T\) to the open set \(X_0\) of \(U\): the
restriction of a canonical segment is canonical (Lemma~\ref{hir:canonical-segments}(1)), its cosupport is the trace of the old one, hence empty
(Lemma~\ref{hir:open-restriction-tower}), and the open set \(X_0\) of \(U\) is biholomorphic to the open set \(X_0\) of \(X\) by the identity on points
(Lemma~\ref{bcalc:restrict-nested}); transporting along this biholomorphism keeps the segment canonical and a resolution (Definition~\ref{hir:tower-isomorphism},
Lemma~\ref{mc:tower-isomorphism-transform}). This canonical resolution over \(X_0\) replaces the compact-support driver in the proofs of
Theorems~\ref{hir:principalization-compact-support} and~\ref{hir:embedded-resolution-compact-support}, in the second case with the first-centre argument of
Theorem~\ref{hir:embedded-resolution-local} for \(V\cap X_0\) (which need not be irreducible, so the strict transform is not claimed connected) and
\(f^{-1}(V_{\mathrm{sing}}\cap X_0)\subset\operatorname{Exc}(f)\) from Theorem~\ref{hir:embedded-resolution-exceptional}, applied over \(X_0\), where the regular points of
\(V\cap X_0\) of dimension \(\dim V\) are those of \(V\). If \(\dim V=n\), \(V\cap X_0\) is open and closed in \(X_0\) and the identity works. For the local forms take for
\(X_0\) an open neighbourhood of \(K\) with compact closure (manifolds are locally compact).
\end{proof}

\begin{theorem}[Principalization relative to a simple normal crossings boundary]
\label{hir:principalization-boundary}
\uses{hir:admissible-sequence,hir:principalization-compact-support,hir:principalization,hir:snc,hir:geo-transform,hir:geo-transform-ambient,hir:wellformed-tower,bcalc:inv-domain-stable,bcalc:tower-exceptional,bcalc:exceptional-divisor,mc:factorization-domain,ana:blowups}
\source{kollar07-resolution,bierstone-milman97-invent-math:page-0010,wlodarczyk05-hironaka}
\lean{TranscendentalNakamaye.Pending.HironakaPrincipalizationBoundary.isInvDomain_univ_of_isSNCDivisor,TranscendentalNakamaye.Pending.HironakaPrincipalizationBoundary.support_exceptionalDivisor_eq_biUnion_drop,TranscendentalNakamaye.Pending.HironakaPrincipalizationBoundary.exists_principalization_boundary,TranscendentalNakamaye.Pending.HironakaPrincipalizationBoundary.exists_principalization_boundary_of_compactSpace}
\leanok
Let \(X\) be a complex manifold, not necessarily compact, \(E_0=(H_1,\dots,H_s)\) a divisor on \(X\) with simple normal crossings, and \(\mathfrak a\subset\mathcal O_X\) a
coherent ideal, nonzero at every point, whose zero locus \(Z(\mathfrak a)\) is compact (for instance \(X\) compact). There are an admissible blowup sequence
\(f\colon X'\to X\) of \((X,E_0)\) with centres over \(Z(\mathfrak a)\) (Definition~\ref{hir:admissible-sequence}), with divisor \(E=(H_1',\dots,H_s',E_1,\dots,E_r)\), and an
effective divisor \(D\) on \(X'\), such that \(E\) has simple normal crossings (\(H_i'\) the strict transform of \(H_i\)), \(f^*\mathfrak a=\mathcal O_{X'}(-D)\),
\(\operatorname{Supp}D=E_1\cup\dots\cup E_r=f^{-1}(Z(\mathfrak a))\), and \(f\) is an isomorphism over \(X\setminus Z(\mathfrak a)\).
\end{theorem}
\begin{proof}
\leanok
A divisor with simple normal crossings is transverse to \(X\) itself, so \((X,X,E_0,\mathfrak a,1)\) is a well-formed marked ideal of the domain in dimension \(\le n\)
(Definition~\ref{hir:snc}; the nonvanishing as in Lemma~\ref{mc:factorization-domain}); its cosupport is \(Z(\mathfrak a)\). As in Theorem~\ref{hir:principalization-compact-support}
its canonical segment resolves it; the divisors of the transforms have simple normal crossings (Proposition~\ref{bcalc:inv-domain-stable},
Lemma~\ref{hir:wellformed-tower}), and their first \(s\) members are the strict transforms of the \(H_i\) (Definition~\ref{hir:geo-transform}). With \(N=X\) and mark one,
\(T^*\mathfrak a=\mathcal O(-E_T)\) whatever the initial divisor (Lemma~\ref{hir:geo-transform-ambient}). The support of \(E_T\) is the union of the members created by the
blowups: after a blowup \(\pi\) along \(C\) it is \(\pi^{-1}(\operatorname{Supp}E_T)\cup\pi^{-1}(C)\) (Lemma~\ref{bcalc:exceptional-divisor}), and as in the proof of
Theorem~\ref{hir:principalization} this is the union of the strict transforms of the earlier created members and of \(\pi^{-1}(C)\); the strict transforms of the \(H_i\) do not
enter. The rest is as in Theorem~\ref{hir:principalization-compact-support}.
\end{proof}
This is an analytic form of Koll\'ar's ``Principalization, III'' (stated there for smooth algebraic varieties, where no compactness is needed), without its functoriality.

\begin{theorem}[Embedded resolution of an analytic set of pure dimension]
\label{hir:embedded-resolution-pure-dimensional}
\uses{hir:admissible-sequence,hir:embedded-resolution-compact-support,hir:resolution-relatively-compact,hir:principalization-compact-support,hir:embedded-resolution-local,hir:embedded-resolution-exceptional,mc:subvariety-thin,ana:cartan-coherence,ana:regular-transport,fpgeom:complex-analytic-subvariety,ana:blowups}
\source{bierstone-milman97-invent-math:page-0009,wlodarczyk05-hironaka}
\lean{TranscendentalNakamaye.Pending.HironakaEmbeddedReducible.isClosed_diff_regularLocus,TranscendentalNakamaye.Pending.HironakaEmbeddedReducible.diff_regularLocus_eq_of_pure,TranscendentalNakamaye.Pending.HironakaEmbeddedReducible.interior_eq_empty_of_pure,TranscendentalNakamaye.Pending.HironakaEmbeddedReducible.isNonzeroAlongAt_vanishingIdeal_of_interior_eq_empty,TranscendentalNakamaye.Pending.HironakaEmbeddedReducible.isNonzeroAlongAt_vanishingIdeal_of_pure,TranscendentalNakamaye.Pending.HironakaEmbeddedReducible.isInvDomain_vanishingIdeal_of_interior_eq_empty,TranscendentalNakamaye.Pending.HironakaEmbeddedReducible.isInvDomain_vanishingIdeal_of_pure,TranscendentalNakamaye.Pending.HironakaEmbeddedReducible.preimage_diff_regularLocus_subset_exceptionalLocus,TranscendentalNakamaye.Pending.HironakaEmbeddedReducible.exists_embeddedResolution_pure_of_isCanonical,TranscendentalNakamaye.Pending.HironakaEmbeddedReducible.exists_embeddedResolution_pure_of_isCompact,TranscendentalNakamaye.Pending.HironakaEmbeddedReducible.exists_embeddedResolution_pure,TranscendentalNakamaye.Pending.HironakaEmbeddedReducible.preimage_val_diff_regularLocus,TranscendentalNakamaye.Pending.HironakaEmbeddedReducible.exists_embeddedResolution_pure_of_subset_isCompact,TranscendentalNakamaye.Pending.HironakaEmbeddedReducible.subset_closure_regularLocus,TranscendentalNakamaye.Pending.HironakaEmbeddedReducible.subset_closure_regularLocus_union}
\leanok
Let \(X\) be a complex manifold of dimension \(n\), not necessarily compact, and \(W\subset X\) an analytic set of pure dimension \(d<n\): the points at which \(W\) is a
submanifold of dimension \(d\) are dense in \(W\). \(W\) may be reducible. Put \(W_{\mathrm{sing}}=W\setminus\operatorname{Reg}_d(W)\); it is closed, and it is the set of
points of \(W\) at which \(W\) is not a submanifold of any dimension.
\begin{enumerate}
\item If \(W\) is compact (for instance if \(X\) is compact), there is an admissible blowup sequence \(f\colon X'\to X\) of \((X,\emptyset)\) with centres over
\(W_{\mathrm{sing}}\) (Definition~\ref{hir:admissible-sequence}), with divisor \(E\), such that \(E\) has simple normal crossings,
\(\operatorname{Exc}(f)=\bigcup E=f^{-1}(W_{\mathrm{sing}})\), \(f\) is an isomorphism over \(X\setminus W_{\mathrm{sing}}\), and the strict transform \(W'\) of \(W\) is a
closed submanifold of \(X'\), of dimension \(d\) at each of its points (not necessarily connected), with only normal crossings with \(E\).
\item Without compactness of \(W\), the same holds for \(W\cap X_0\) over every open set \(X_0\subset X\) contained in a compact set, with \(X_0\) in place of \(X\) and centres
over \(W_{\mathrm{sing}}\cap X_0\).
\end{enumerate}
An irreducible subvariety of dimension \(<n\) is of pure dimension, and so is a disjoint union of two closed sets of pure dimension \(d\).
\end{theorem}
\begin{proof}
\leanok
\(W_{\mathrm{sing}}\) is closed because regularity of dimension \(d\) at a point of \(W\) persists at nearby points of \(W\) (Definition~\ref{fpgeom:complex-analytic-subvariety}). If
\(W\) is a submanifold of dimension \(k\) near \(x\in W\), the nearby points of \(W\) are regular of dimension \(k\); one of them is regular of dimension \(d\) by density, so
\(k=d\). An interior point of \(W\) would have nearby regular points of dimension \(d\) which are interior points, hence regular of dimension \(n>d\); so \(W\) has empty
interior. Hence \(\mathcal I_W\) is nonzero at every point (as in Lemma~\ref{mc:subvariety-thin}: off \(W\) it contains \(1\); at a point of \(W\) some local equation of \(W\) does
not vanish identically near it, since otherwise \(W\) would contain a neighbourhood), it is coherent by Cartan's theorem (Theorem~\ref{ana:cartan-coherence}, valid for every
analytic set), and \((X,X,\emptyset,\mathcal I_W,1)\) lies in the domain. Its cosupport is \(W\).

(1) Theorem~\ref{hir:principalization-compact-support} gives a canonical segment with empty final cosupport. The first-centre argument of
Theorem~\ref{hir:embedded-resolution-local} uses only that \(W\) is analytic and regular of dimension \(d\) at every point of the dense open subset \(W\setminus W_{\mathrm{sing}}\)
of \(W\); irreducibility is never used. With the open set \(O=X\setminus W_{\mathrm{sing}}\) it gives the part of the segment before the first centre meeting \(f^{-1}(O)\): a
multiple test blow-up whose centres lie over \(W_{\mathrm{sing}}\), whose strict transform \(W'\) is a closed submanifold regular of dimension \(d\), with only normal crossings
with the transformed divisor \(E\), which has simple normal crossings and whose union is the support of the exceptional divisor, containing \(\operatorname{Exc}(f)\) and lying
over \(W_{\mathrm{sing}}\). A point \(y\notin\operatorname{Exc}(f)\) over \(W\) lies on \(W'\), and near \(y\) the set \(W'\) is the preimage of \(W\) under a local biholomorphism,
so \(f(y)\) is a regular point of \(W\) of dimension \(d\) (Lemma~\ref{ana:regular-transport}); hence \(f^{-1}(W_{\mathrm{sing}})\subset\operatorname{Exc}(f)\), and
\(\operatorname{Exc}(f)=\bigcup E=f^{-1}(W_{\mathrm{sing}})\) as in Theorem~\ref{hir:embedded-resolution-exceptional}. The isomorphism over \(X\setminus W_{\mathrm{sing}}\) and the
admissibility are as in Theorem~\ref{hir:embedded-resolution-compact-support}. (2) Replace the compact-support driver by the canonical resolution over \(X_0\) of
Theorem~\ref{hir:resolution-relatively-compact}; \(W\cap X_0\) is of pure dimension \(d\) in \(X_0\), with the same regular points.

For the last sentence: an irreducible subvariety is the closure of its regular points of dimension \(\dim V\) (Theorem~\ref{hir:embedded-resolution-compact-support}, through
Lemma~\ref{mc:subvariety-thin}); near a point of one of two disjoint closed sets the union agrees with that set.
\end{proof}
This is Bierstone--Milman's Theorem~1.6 for a closed analytic subset of pure dimension, which may be reducible; when the components have different dimensions the canonical
resolution of \((\mathcal I_W,1)\) is not enough (it may blow up a smooth component before the others are resolved), and that case is not covered here.
That case is covered, by one admissible sequence, in Theorem~\ref{hir:embedded-resolution-analytic}.

\begin{theorem}[Embedded resolution relative to a simple normal crossings boundary]
\label{hir:embedded-resolution-boundary}
\uses{hir:admissible-sequence,hir:embedded-resolution-pure-dimensional,hir:principalization-boundary,hir:principalization-compact-support,hir:resolution-relatively-compact,hir:embedded-resolution-local,hir:embedded-resolution-compact-support,hir:embedded-from-resolution,hir:smooth-centre,mc:smooth-part-invariant,mc:first-centre-exists,hir:snc,bcalc:inv-domain-stable,bcalc:exceptional-divisor,bcalc:tower-exceptional,bcalc:restrict-empty-tower,ana:cartan-coherence,ana:blowups}
\source{bierstone-milman97-invent-math:page-0009,bierstone-milman97-invent-math:page-0010,kollar07-resolution,wlodarczyk05-hironaka}
\lean{TranscendentalNakamaye.Pending.HironakaEmbeddedBoundary.map_preimage_of_forall_eq_empty,TranscendentalNakamaye.Pending.HironakaEmbeddedBoundary.cosupp_mk_vanishingIdeal,TranscendentalNakamaye.Pending.HironakaEmbeddedBoundary.pullback_val_mk_vanishingIdeal,TranscendentalNakamaye.Pending.HironakaEmbeddedBoundary.pullback_pullback_val_mk_of_forall_eq_empty,TranscendentalNakamaye.Pending.HironakaEmbeddedBoundary.inv_refl_eq_of_regularAtDim_of_forall_eq_empty,TranscendentalNakamaye.Pending.HironakaEmbeddedBoundary.inv_eq_of_isEmptyTower_of_pureDim_of_forall_eq_empty,TranscendentalNakamaye.Pending.HironakaEmbeddedBoundary.inv_eq_over_open_of_isEmptyTower_of_pureDim_of_disjoint,TranscendentalNakamaye.Pending.HironakaEmbeddedBoundary.mem_cosupp_iff_of_centresAvoid,TranscendentalNakamaye.Pending.HironakaEmbeddedBoundary.mem_maxLocus_iff_of_centresAvoid_of_pureDim,TranscendentalNakamaye.Pending.HironakaEmbeddedBoundary.exists_tower_embedded_of_isCanonical,TranscendentalNakamaye.Pending.HironakaEmbeddedBoundary.centresOver_of_inter_subset,TranscendentalNakamaye.Pending.HironakaEmbeddedBoundary.support_exceptionalDivisor_subset_of_centresOver,TranscendentalNakamaye.Pending.HironakaEmbeddedBoundary.isInvDomain_boundary,TranscendentalNakamaye.Pending.HironakaEmbeddedBoundary.subset_closure_inter_compl,TranscendentalNakamaye.Pending.HironakaEmbeddedBoundary.isClosed_biUnion_of_forall_isClosed,TranscendentalNakamaye.Pending.HironakaEmbeddedBoundary.exists_embeddedResolution_boundary_of_isCanonical,TranscendentalNakamaye.Pending.HironakaEmbeddedBoundary.exists_embeddedResolution_boundary_of_isCompact,TranscendentalNakamaye.Pending.HironakaEmbeddedBoundary.exists_embeddedResolution_boundary,TranscendentalNakamaye.Pending.HironakaEmbeddedBoundary.subset_closure_diff_biUnion_of_forall_disjoint,TranscendentalNakamaye.Pending.HironakaEmbeddedBoundary.preimage_biUnion_list,TranscendentalNakamaye.Pending.HironakaEmbeddedBoundary.exists_embeddedResolution_boundary_of_subset_isCompact}
\leanok
Let \(X\) be a complex manifold of dimension \(n\), not necessarily compact, \(E_0=(H_1,\dots,H_s)\) a divisor on \(X\) with simple normal crossings, with support
\(|E_0|=H_1\cup\dots\cup H_s\), and \(W\subset X\) an analytic set of pure dimension \(d<n\) (Theorem~\ref{hir:embedded-resolution-pure-dimensional}; \(W\) may be reducible)
such that no member of \(E_0\) contains a nonempty relatively open subset of \(W\), that is, \(W\subset\overline{W\setminus|E_0|}\). Put
\(W_{\mathrm{sing}}=W\setminus\operatorname{Reg}_d(W)\) and \(F=W_{\mathrm{sing}}\cup(W\cap|E_0|)\).
\begin{enumerate}
\item If \(W\) is compact (for instance if \(X\) is compact), there is an admissible blowup sequence \(f\colon X'\to X\) of \((X,E_0)\) with centres over \(F\)
(Definition~\ref{hir:admissible-sequence}), with divisor \(E=(H_1',\dots,H_s',E_1,\dots,E_r)\) (\(H_i'\) the strict transform of \(H_i\)), such that \(E\) has simple normal
crossings, \(\operatorname{Exc}(f)\subset E_1\cup\dots\cup E_r\subset f^{-1}(F)\), \(f^{-1}(W_{\mathrm{sing}})\subset\operatorname{Exc}(f)\), \(f\) is an isomorphism over
\(X\setminus F\), and the strict transform \(W'\) of \(W\) is a closed submanifold of \(X'\), of dimension \(d\) at each of its points, with only normal crossings with \(E\).
\item Without compactness of \(W\), the same holds for \(W\cap X_0\) over every open set \(X_0\subset X\) contained in a compact set, with \(X_0\) in place of \(X\), the
restriction of \(E_0\) to \(X_0\) in place of \(E_0\), and centres over \(F\cap X_0\).
\end{enumerate}
If no member of \(E_0\) meets \(W\) (for instance if \(E_0=\emptyset\)), the hypothesis on \(E_0\) holds and \(F=W_{\mathrm{sing}}\).
\end{theorem}
\begin{proof}
\leanok
As in the proof of Theorem~\ref{hir:embedded-resolution-pure-dimensional}, \(W\) has empty interior, \(\mathcal I_W\) is coherent (Theorem~\ref{ana:cartan-coherence}) and
nonzero at every point. A divisor with simple normal crossings is transverse to \(X\) (Definition~\ref{hir:snc}), so \(\mathfrak m=(X,X,E_0,\mathcal I_W,1)\) is a well-formed
marked ideal of the domain in dimension \(\le n\); its cosupport is \(W\) (the cosupport does not involve the divisor).

(1) Theorem~\ref{hir:principalization-compact-support} gives a canonical segment \(T_0\) of \(\mathfrak m\) with empty final cosupport. Let
\(O=X\setminus(W_{\mathrm{sing}}\cup|E_0|)\); it is open, since \(W_{\mathrm{sing}}\) and the members of \(E_0\) are closed. \(W\) is regular of dimension \(d\) at every
point of \(W\cap O\), and \(W\cap O\) is dense in \(W\): a neighbourhood \(U\) of \(x\in W\) contains a point \(y\in W\setminus|E_0|\); then \(U\setminus|E_0|\) is a
neighbourhood of \(y\), so it contains a regular point of \(W\) of dimension \(d\), which lies in \(W\cap O\).

The new point is the behaviour of the invariant over \(O\). The restriction of \(\mathfrak m\) to \(O\) is \((O,O,(\emptyset,\dots,\emptyset),\mathcal I_{W\cap O},1)\): every
member of \(E_0\) restricts to the empty set. Let \(z,z'\) be points of the closed submanifold \(Z=W\cap O\) of \(O\) at which it has the same dimension. The biholomorphism
\(\Phi\colon U_1\to U_2\) between neighbourhoods of \(z\) and \(z'\) of Lemma~\ref{hir:smooth-centre} (homogeneity) with \(\Phi(z)=z'\) pulls the restriction of
\((O,O,\emptyset,\mathcal I_Z,1)\) to \(U_2\) back to its restriction to \(U_1\); since \(\Phi^{-1}(\emptyset)=\emptyset\), it also pulls the restriction of
\((O,O,(\emptyset,\dots,\emptyset),\mathcal I_Z,1)\) to \(U_2\) back to its restriction to \(U_1\). So the argument of Lemma~\ref{mc:smooth-part-invariant} (restriction to open sets,
invariance under isomorphisms of towers, and empty steps) applies verbatim to \(\mathfrak m\) over \(O\): along a canonical segment of \(\mathfrak m\) whose centres miss the preimage
of \(O\), the value of the invariant is constant on the part of the cosupport over \(O\), which is the preimage of \(W\cap O\) (Lemma~\ref{bcalc:restrict-empty-tower}).

Now argue as in Lemma~\ref{hir:embedded-from-resolution} and the first-centre argument of Theorem~\ref{hir:embedded-resolution-local}. If no centre of \(T_0\) meets the
preimage of \(O\), the final cosupport contains the preimage of \(W\cap O\), which is therefore empty (blowup sequences are onto), and \(W=\emptyset\) by density; take
\(f=T_0\). Otherwise let \(T\) be the part of \(T_0\) before the first centre \(C\) meeting the preimage of \(O\) (Lemma~\ref{mc:first-centre-exists}). \(C\) is the maximum
locus of the invariant on the cosupport, an admissible centre; since it meets the preimage of \(O\), the maximum is the constant value there, and over \(O\) the centre \(C\) is
exactly the preimage of \(W\). The identity-principle step of Theorem~\ref{hir:embedded-resolution-local}, which involves no marked ideal, shows that the strict transform
\(W'\) of \(W\) lies in \(C\), agrees with \(C\) near each of its points and is regular of dimension \(d\) at each of its points; as \(C\) has only normal crossings with the
transformed divisor \(E\), so has \(W'\).

\(E\) has simple normal crossings (Proposition~\ref{bcalc:inv-domain-stable}) and its first \(s\) members are the strict transforms of the \(H_i\). The centres of \(T\) lie in
the cosupports, hence over \(W\), and they miss the preimages of \(O\), hence lie over \(W\setminus O=F\): \(T\) is an admissible blowup sequence of \((X,E_0)\) with centres over
\(F\) (Definition~\ref{hir:admissible-sequence}). The support of the exceptional divisor of \(T\) (Definition~\ref{bcalc:tower-exceptional}) contains \(\operatorname{Exc}(f)\), it is
\(E_1\cup\dots\cup E_r\) (proof of Theorem~\ref{hir:principalization-boundary}), and it lies over \(F\): after a blowup \(\pi\) along a centre \(C'\) it is
\(\pi^{-1}(\operatorname{Supp}E_T)\cup\pi^{-1}(C')\) (Lemma~\ref{bcalc:exceptional-divisor}). Hence \(f(\operatorname{Exc}f)\subset F\), and \(f\) is an isomorphism over
\(X\setminus F\) as in Theorem~\ref{hir:embedded-resolution-compact-support}. Finally \(f^{-1}(W_{\mathrm{sing}})\subset\operatorname{Exc}(f)\) as in the proof of
Theorem~\ref{hir:embedded-resolution-pure-dimensional}, which uses only that \(W'\) is regular of dimension \(d\).

(2) Replace the compact-support driver by the canonical resolution over \(X_0\) of Theorem~\ref{hir:resolution-relatively-compact}, applied to \(\mathfrak m\). Its restriction to
\(X_0\) is \((X_0,X_0,E_0|_{X_0},\mathcal I_{W\cap X_0},1)\); \(E_0|_{X_0}\) has simple normal crossings, \(W\cap X_0\) is of pure dimension \(d\) with the same regular points, and
\(W\cap X_0\subset\overline{(W\cap X_0)\setminus|E_0|}\) because closures commute with restriction to an open set. Then argue as in (1) on \(X_0\).

If no member of \(E_0\) meets \(W\), then \(W\setminus|E_0|=W\) and \(W\cap|E_0|=\emptyset\).
\end{proof}
This is an analytic form of Bierstone--Milman's Theorem~1.6 with an initial divisor (Koll\'ar's ``Principalization, III'' setting), for a closed analytic set of pure dimension.
As in Theorem~1.6~(1) of Bierstone--Milman, where a centre may lie in \(X_j\cap E_j\) once \(X_j\) is smooth, \(f\) may blow up smooth points of \(W\) lying on \(|E_0|\). The
sharper form in which \(f\) is an isomorphism over every point at which \(W\) is smooth and has normal crossings with \(E_0\) would need the invariant of \(\mathfrak m\) to be
constant along the smooth part of \(W\) across the strata of \(E_0\); it is not covered here. With \(E_0=\emptyset\), Theorem~\ref{hir:embedded-resolution-pure-dimensional} gives
the sharper conclusion \(\operatorname{Exc}(f)=\bigcup E=f^{-1}(W_{\mathrm{sing}})\).

\begin{theorem}[Embedded resolution of the dimension parts of an analytic set]
\label{hir:embedded-resolution-dimension-parts}
\uses{hir:embedded-resolution-pure-dimensional,hir:embedded-resolution-boundary,hir:admissible-sequence,ana:dimension-strata,lag:regular-locus-dense,fpgeom:complex-analytic-subvariety}
\source{bierstone-milman97-invent-math:page-0009,wlodarczyk05-hironaka,demailly-cadg:II.5.3}
\lean{TranscendentalNakamaye.Pending.HironakaEmbeddedDimensionParts.dimensionPart,TranscendentalNakamaye.Pending.HironakaEmbeddedDimensionParts.regularStratum_eq_regularLocus,TranscendentalNakamaye.Pending.HironakaEmbeddedDimensionParts.isComplexAnalyticSubset_dimensionPart,TranscendentalNakamaye.Pending.HironakaEmbeddedDimensionParts.dimensionPart_subset,TranscendentalNakamaye.Pending.HironakaEmbeddedDimensionParts.regularAtDim_dimensionPart_of_regularAtDim,TranscendentalNakamaye.Pending.HironakaEmbeddedDimensionParts.subset_closure_regularLocus_dimensionPart,TranscendentalNakamaye.Pending.HironakaEmbeddedDimensionParts.eq_iUnion_dimensionPart,TranscendentalNakamaye.Pending.HironakaEmbeddedDimensionParts.dimensionPart_eq_self_of_pure,TranscendentalNakamaye.Pending.HironakaEmbeddedDimensionParts.diff_regularLocus_dimensionPart_subset,TranscendentalNakamaye.Pending.HironakaEmbeddedDimensionParts.exists_embeddedResolution_dimensionPart_of_isCompact,TranscendentalNakamaye.Pending.HironakaEmbeddedDimensionParts.exists_embeddedResolution_dimensionPart,TranscendentalNakamaye.Pending.HironakaEmbeddedDimensionParts.exists_embeddedResolution_dimensionPart_boundary_of_isCompact}
\leanok
Let \(X\) be a complex manifold of dimension \(n\) and \(W\subset X\) a closed analytic set, whose irreducible components may have different dimensions. For
\(k\in\mathbb N\) let \(W^{(k)}=\overline{\operatorname{Reg}_k(W)}\) be the \emph{\(k\)-dimensional part} of \(W\), the closure of the set of points at which \(W\) is a
submanifold of dimension \(k\), and let \(W^{(k)}_{\mathrm{sing}}=W^{(k)}\setminus\operatorname{Reg}_k(W^{(k)})\).
\begin{enumerate}
\item \(W^{(k)}\subset W\) is analytic and of pure dimension \(k\); every regular point of \(W\) of dimension \(k\) is a regular point of \(W^{(k)}\) of dimension \(k\), so
\(W^{(k)}_{\mathrm{sing}}\subset W^{(k)}\setminus\operatorname{Reg}_k(W)\); \(W=W^{(0)}\cup\dots\cup W^{(n)}\); and \(W^{(d)}=W\) if \(W\) is of pure dimension \(d\).
\item Let \(k<n\). If \(W\) is compact (for instance if \(X\) is compact), there is an admissible blowup sequence \(f\colon X'\to X\) of \((X,\emptyset)\) with centres
over \(W^{(k)}_{\mathrm{sing}}\) (Definition~\ref{hir:admissible-sequence}), with divisor \(E\), such that \(E\) has simple normal crossings,
\(\operatorname{Exc}(f)=\bigcup E=f^{-1}(W^{(k)}_{\mathrm{sing}})\), \(f\) is an isomorphism over \(X\setminus W^{(k)}_{\mathrm{sing}}\), and the strict transform of
\(W^{(k)}\) is a closed submanifold of \(X'\), of dimension \(k\) at each of its points, with only normal crossings with \(E\). Relative to a divisor \(E_0\) with simple
normal crossings such that \(W^{(k)}\subset\overline{W^{(k)}\setminus|E_0|}\), the conclusions of Theorem~\ref{hir:embedded-resolution-boundary}(1) hold for \(W^{(k)}\).
\end{enumerate}
\end{theorem}
\begin{proof}
\leanok
(1) \(W^{(k)}\) is the closure of the dimension stratum \(R_k\), which is analytic (Proposition~\ref{ana:dimension-strata}); it lies in the closed set \(W\). If \(W\) is a
submanifold of dimension \(k\) near \(x\), then on a neighbourhood \(U\) of \(x\) every point of \(W\) is regular of dimension \(k\)
(Definition~\ref{fpgeom:complex-analytic-subvariety}), so \(W\cap U\subset R_k\subset W^{(k)}\subset W\) and \(W^{(k)}\cap U=W\cap U\): \(x\) is a regular point of \(W^{(k)}\) of
dimension \(k\). Hence \(R_k\subset\operatorname{Reg}_k(W^{(k)})\), \(W^{(k)}=\overline{R_k}\subset\overline{\operatorname{Reg}_k(W^{(k)})}\) (pure dimension \(k\)), and a point
of \(W^{(k)}\) not regular of dimension \(k\) for \(W^{(k)}\) is not regular of dimension \(k\) for \(W\). Every neighbourhood of a point of \(W\) contains a regular point of
\(W\) (Theorem~\ref{lag:regular-locus-dense}), of some dimension \(\le n\); so \(W\subset\overline{R_0\cup\dots\cup R_n}=W^{(0)}\cup\dots\cup W^{(n)}\) (a finite union).
If \(W\) is of pure dimension \(d\), \(W\subset\overline{R_d}=W^{(d)}\subset W\).

(2) \(W^{(k)}\) is closed in the compact set \(W\), hence compact, analytic and of pure dimension \(k<n\); apply
Theorem~\ref{hir:embedded-resolution-pure-dimensional}(1), and Theorem~\ref{hir:embedded-resolution-boundary}(1) for the boundary form.
\end{proof}
Each dimension part is resolved by its own blowup sequence. One sequence resolving all the dimension parts at once, with smooth and pairwise disjoint strict transforms
(Bierstone--Milman's Theorem~1.6 for a non-equidimensional closed subspace), is not covered: the canonical resolution of \((\mathcal I_W,1)\) can be cut only before the
first centre meeting the smooth part of the parts of largest invariant, and that centre contains their strict transform while the other parts may still be singular;
resolving the parts one after the other would need their strict transforms to be separated first (for instance by principalizing \(\mathcal I_A+\mathcal I_B\)), since a
later centre may meet an earlier, already smooth, strict transform.
Theorem~\ref{hir:embedded-resolution-analytic} carries out exactly this: it separates the parts by principalizing \(\mathcal I_A+\mathcal I_B\) (Lemma~\ref{hir:separation}) and then resolves them one after the other, each resolution taking place away from the other parts.

\begin{theorem}[Resolution relative to a boundary over relatively compact open sets]
\label{hir:boundary-relatively-compact}
\uses{hir:principalization-boundary,hir:embedded-resolution-boundary,hir:resolution-relatively-compact,hir:principalization-compact-support,hir:admissible-sequence,hir:geo-transform-ambient,bcalc:inv-domain-stable,bcalc:exceptional-divisor,hir:principalization-local}
\source{bierstone-milman97-invent-math:page-0010,kollar07-resolution,wlodarczyk05-hironaka}
\lean{TranscendentalNakamaye.Pending.HironakaBoundaryRelativelyCompact.exists_principalization_boundary_of_subset_isCompact,TranscendentalNakamaye.Pending.HironakaBoundaryRelativelyCompact.exists_principalization_boundary_local,TranscendentalNakamaye.Pending.HironakaBoundaryRelativelyCompact.exists_embeddedResolution_boundary_local,TranscendentalNakamaye.Pending.HironakaAPI2.exists_principalization_boundary_of_subset_isCompact,TranscendentalNakamaye.Pending.HironakaAPI2.exists_principalization_boundary_local,TranscendentalNakamaye.Pending.HironakaAPI2.exists_embeddedResolution_boundary_local}
\leanok
Let \(X\) be a complex manifold, not necessarily compact, \(E_0\) a divisor on \(X\) with simple normal crossings, and \(X_0\subset X\) an open set contained in a compact
set.
\begin{enumerate}
\item For a coherent ideal \(\mathfrak a\subset\mathcal O_X\), nonzero at every point (no compactness of \(Z(\mathfrak a)\)), there are an admissible blowup sequence
\(f\colon X_0'\to X_0\) of \((X_0,E_0|_{X_0})\) with centres over \(Z(\mathfrak a)\cap X_0\) (Definition~\ref{hir:admissible-sequence}), with divisor
\(E=(H_1',\dots,H_s',E_1,\dots,E_r)\), and an effective divisor \(D\) on \(X_0'\) such that \(E\) has simple normal crossings,
\(f^*(\mathfrak a|_{X_0})=\mathcal O_{X_0'}(-D)\), \(\operatorname{Supp}D=E_1\cup\dots\cup E_r=f^{-1}(Z(\mathfrak a)\cap X_0)\), and \(f\) is an isomorphism over
\(X_0\setminus Z(\mathfrak a)\).
\item For every compact \(K\subset X\) there is a relatively compact open neighbourhood \(U\) of \(K\) over which (1) holds, and over which the conclusions of
Theorem~\ref{hir:embedded-resolution-boundary}(2) hold for an analytic set \(W\) of pure dimension \(d<n\) with \(W\subset\overline{W\setminus|E_0|}\).
\end{enumerate}
\end{theorem}
\begin{proof}
\leanok
(1) As in the proof of Theorem~\ref{hir:principalization-boundary}, \((X,X,E_0,\mathfrak a,1)\) lies in the domain. Theorem~\ref{hir:resolution-relatively-compact} gives a
canonical resolution over \(X_0\) of its restriction \((X_0,X_0,E_0|_{X_0},\mathfrak a|_{X_0},1)\), which lies in the domain (Proposition~\ref{bcalc:inv-domain-stable});
its cosupport is \(Z(\mathfrak a)\cap X_0\). With \(N=X_0\) and mark one, the pullback of \(\mathfrak a|_{X_0}\) is the ideal of the exceptional divisor \(E_T\) of the
resolving tower (Lemma~\ref{hir:geo-transform-ambient}); its support is \(f^{-1}(Z(\mathfrak a)\cap X_0)\) and, as in the proof of Theorem~\ref{hir:principalization-boundary}
(Lemma~\ref{bcalc:exceptional-divisor}), the union of the members created by the blowups. The centres lie in the cosupports, hence over \(Z(\mathfrak a)\cap X_0\), and \(f\) is
an isomorphism over \(X_0\setminus Z(\mathfrak a)\) as in Theorem~\ref{hir:principalization-local}. The divisor has simple normal crossings and the sequence is admissible as in
Theorem~\ref{hir:principalization-compact-support}. (2) Manifolds are locally compact: take for \(U\) an open neighbourhood of \(K\) with compact closure and apply (1), resp.
Theorem~\ref{hir:embedded-resolution-boundary}(2), with \(X_0=U\).
\end{proof}

\begin{lemma}[Calculus of admissible sequences]
\label{hir:admissible-calculus}
\uses{hir:admissible-sequence,ana:blowups,ana:strict-transform-composition,bcalc:isomorphism-off-centre,hir:principalization-local,hir:embedded-resolution-local,ana:regular-transport,hir:snc-invariance,hir:snc}
\source{bierstone-milman97-invent-math:page-0009,kollar07-resolution}
\lean{TranscendentalNakamaye.Pending.HironakaAdmissibleCalculus.isProperMap_of_isBlowupSequence,TranscendentalNakamaye.Pending.HironakaAdmissibleCalculus.isComplexAnalyticSubset_strictTransform,TranscendentalNakamaye.Pending.HironakaAdmissibleCalculus.strictTransform_subset_preimage,TranscendentalNakamaye.Pending.HironakaAdmissibleCalculus.isCompact_strictTransform,TranscendentalNakamaye.Pending.HironakaAdmissibleCalculus.strictTransform_union,TranscendentalNakamaye.Pending.HironakaAdmissibleCalculus.strictTransform_empty,TranscendentalNakamaye.Pending.HironakaAdmissibleCalculus.strictTransform_mono,TranscendentalNakamaye.Pending.HironakaAdmissibleCalculus.mem_strictTransform_iff_of_notMem,TranscendentalNakamaye.Pending.HironakaAdmissibleCalculus.diff_regularLocus_strictTransform_subset,TranscendentalNakamaye.Pending.HironakaAdmissibleCalculus.strictTransform_subset_closure_diff,TranscendentalNakamaye.Pending.HironakaAdmissibleCalculus.strictTransform_eq_preimage_of_disjoint,TranscendentalNakamaye.Pending.HironakaAdmissibleCalculus.regularAtDim_preimage_of_disjoint,TranscendentalNakamaye.Pending.HironakaAdmissibleCalculus.IsAdmissibleBlowupSequence.comp,TranscendentalNakamaye.Pending.HironakaAdmissibleCalculus.IsAdmissibleBlowupSequence.exists_decomposition,TranscendentalNakamaye.Pending.HironakaAdmissibleCalculus.IsAdmissibleBlowupSequence.exceptionalLocus_subset,TranscendentalNakamaye.Pending.HironakaAdmissibleCalculus.IsAdmissibleBlowupSequence.isIsomorphismOver_compl,TranscendentalNakamaye.Pending.HironakaAdmissibleCalculus.IsAdmissibleBlowupSequence.exceptionalLocus_subset_drop,TranscendentalNakamaye.Pending.HironakaAdmissibleCalculus.IsAdmissibleBlowupSequence.biUnion_subset,TranscendentalNakamaye.Pending.HironakaAdmissibleCalculus.IsAdmissibleBlowupSequence.isClosed_mem,TranscendentalNakamaye.Pending.HironakaAdmissibleCalculus.isCoordinateFamilyAt_of_forall_mem,TranscendentalNakamaye.Pending.HironakaAdmissibleCalculus.IsAdmissibleBlowupSequence.hasNormalCrossings_preimage_of_disjoint}
\leanok
Let \(f\colon Y\to X\) be an admissible blowup sequence of \((X,E_0)\) with centres over a set \(F\subset X\) (Definition~\ref{hir:admissible-sequence}), with divisor
\(E\), the members of \(E_0=(H_1,\dots,H_s)\) being closed.
\begin{enumerate}
\item (Composition.) If \(g\colon Z\to Y\) is an admissible blowup sequence of \((Y,E)\) with centres over \(f^{-1}(F)\), with divisor \(E'\), then \(f\circ g\) is an
admissible blowup sequence of \((X,E_0)\) with centres over \(F\), with divisor \(E'\).
\item (The divisor.) \(E=(T(H_1),\dots,T(H_s),N_1,\dots,N_r)\), where for every closed \(H\subset X\) the set \(T(H)\) is closed,
\(T(H)\subset f^{-1}(H)\subset T(H)\cup N_1\cup\dots\cup N_r\), and \(z\in T(H)\) if and only if \(f(z)\in H\), for every \(z\) with \(f(z)\notin F\); the members
\(N_j\) are closed and lie over \(F\), and \(\operatorname{Exc}(f)\subset N_1\cup\dots\cup N_r\). In particular \(f\) is an isomorphism over \(X\setminus F\), and
\(\bigcup E\subset f^{-1}(|E_0|\cup F)\).
\item (Away from the centres.) Let \(F\) be closed and \(S\subset X\) closed and disjoint from \(F\). Then the strict transform of \(S\) is \(f^{-1}(S)\); a regular point
of \(S\) of dimension \(d\) pulls back to regular points of \(f^{-1}(S)\) of dimension \(d\); and if \(S\) has only normal crossings with \(E_0\), then \(f^{-1}(S)\) has only
normal crossings with \(E\).
\item (Strict transforms.) Along any blowup sequence \(f\), the strict transform of an analytic set \(A\) is analytic, and compact if \(A\) is; for closed \(A\) it lies in
\(f^{-1}(A)\); strict transforms commute with finite unions. If \(\operatorname{Exc}(f)\subset f^{-1}(F)\) and the closed set \(A\) is regular of dimension \(d\) at each of
its points off \(F\), its strict transform is regular of dimension \(d\) at each of its points not over \(F\); if \(A\setminus F\) is dense in \(A\), the points of the strict
transform not over \(F\) are dense in it.
\end{enumerate}
\end{lemma}
\begin{proof}
\leanok
(1) By induction on \(g\): a step of \(g\) along a centre \(C\) having only normal crossings with the current divisor, with \(C\subset g_j^{-1}(f^{-1}(F))=(f\circ
g_j)^{-1}(F)\), is a step of \(f\circ g\).

(2) By induction on \(f\). For the identity take \(T(H)=H\) and no new member. After a blowup \(\pi\) along a centre \(C\subset f^{-1}(F)\) the divisor consists of the
\(\overline{\pi^{-1}(T(H_i)\setminus C)}\), the \(\overline{\pi^{-1}(N_j\setminus C)}\) and \(\pi^{-1}(C)\); put \(T'(H)=\overline{\pi^{-1}(T(H)\setminus C)}\), and take
the last \(r+1\) of these sets as the new members. \(T'(H)\) lies in the closed set \(\pi^{-1}(T(H))\subset(f\pi)^{-1}(H)\). A point \(z\) over \(H\) either lies over \(C\),
hence in \(\pi^{-1}(C)\), or \(\pi(z)\notin C\) lies in \(T(H)\) or in some \(N_j\), and then \(z\in T'(H)\) or \(z\in\overline{\pi^{-1}(N_j\setminus C)}\). If
\(f\pi(z)\notin F\), then \(\pi(z)\notin C\) and \(z\in T'(H)\iff\pi(z)\in T(H)\iff f\pi(z)\in H\). The new members lie over the \(N_j\) or over \(C\), hence over \(F\).
Finally \(\operatorname{Exc}(f\pi)=\operatorname{Exc}(\pi)\cup\pi^{-1}(\operatorname{Exc}f)\) by the chain rule, \(\operatorname{Exc}(\pi)\subset\pi^{-1}(C)\) because a blowup
is an isomorphism off its centre (Lemma~\ref{bcalc:isomorphism-off-centre}), and \(\pi^{-1}(N_1\cup\dots\cup N_r)\) lies in the new members as above. A blowup sequence
\(f\) with \(f(\operatorname{Exc}f)\subset F\) is an isomorphism over \(X\setminus F\) (Theorem~\ref{hir:principalization-local}).

(3) No point of \(f^{-1}(S)\) lies over \(F\), so none is exceptional by (2); hence \(f^{-1}(S)\setminus\operatorname{Exc}(f)=f^{-1}(S)\) is closed and equals the strict
transform. At such a point \(z\) the map \(f\) is a local biholomorphism, and regular points pull back (Lemma~\ref{ana:regular-transport}). On the open set
\(O=f^{-1}(X\setminus F)\ni z\) each member \(T(H_i)\) agrees with \(f^{-1}(H_i)\) and no new member meets \(O\), by (2). Adapted coordinates only constrain the members
through the point, so adapted coordinates for \(S\) and the members of \(E_0\) through \(f(z)\) pull back to adapted coordinates for \(f^{-1}(S)\) and the members of \(E\)
through \(z\) (Lemma~\ref{hir:snc-invariance}).

(4) A blowup is proper (Theorem~\ref{hir:principalization-local}), hence so is a blowup sequence; strict transforms compose (Lemma~\ref{ana:strict-transform-composition})
and the strict transform of an analytic set along one blowup is analytic (Theorem~\ref{hir:embedded-resolution-local}), so by induction the strict transform along a
sequence is analytic; for closed \(A\) it lies in the closed set \(f^{-1}(A)\), compact if \(A\) is. The closure of a finite union is the union of the closures. Off
\(\operatorname{Exc}(f)\) the strict transform of a closed set agrees with its preimage, and at a point \(z\) of the strict transform with \(f(z)\notin F\) the point \(z\) is
not exceptional and \(f(z)\) is a regular point of \(A\) of dimension \(d\), which pulls back (Lemma~\ref{ana:regular-transport}). For the density, \(f\) is open near a point
\(y\in f^{-1}(A)\setminus\operatorname{Exc}(f)\) (Lemma~\ref{ana:strict-transform-composition}); the image of a small neighbourhood of \(y\) is a neighbourhood of
\(f(y)\in A\), which contains a point of \(A\setminus F\); so \(y\), and hence the whole strict transform, lies in the closure of the non-exceptional points over
\(A\setminus F\).
\end{proof}

\begin{lemma}[Separation by principalizing \(\mathcal I_A+\mathcal I_B\)]
\label{hir:separation}
\uses{hir:principalization-boundary,hir:admissible-sequence,ana:ideal-sheaves,ana:cartan-coherence,ana:invertible-colon,hir:marked-ideal-operations,hir:embedded-resolution-pure-dimensional,ana:strict-transform-composition,ana:blowups}
\source{bierstone-milman97-invent-math:page-0009,bierstone-milman97-invent-math:page-0010,kollar07-resolution}
\lean{TranscendentalNakamaye.Pending.HironakaSeparation.zeroLocus_sumIdeal_vanishingIdeal,TranscendentalNakamaye.Pending.HironakaSeparation.exists_unit_of_isInSpanAt,TranscendentalNakamaye.Pending.HironakaSeparation.notMem_strictTransform_of_eventually,TranscendentalNakamaye.Pending.HironakaSeparation.disjoint_strictTransform_of_pullback_sumIdeal,TranscendentalNakamaye.Pending.HironakaSeparation.exists_separation}
\leanok
Let \(X\) be a complex manifold of dimension \(n\) and \(A,B\subset X\) closed sets with \(A\subset\overline{A\setminus B}\) and \(B\subset\overline{B\setminus A}\).
\begin{enumerate}
\item The zero locus of \(\mathcal I_A+\mathcal I_B\) is \(A\cap B\).
\item Let \(f\colon Y\to X\) be a holomorphic map of complex manifolds of dimension \(n\) and \(D\) an effective divisor on \(Y\) with
\(f^*(\mathcal I_A+\mathcal I_B)=\mathcal O_Y(-D)\) and \(\operatorname{Supp}D\subset f^{-1}(A\cap B)\). Then the strict transforms of \(A\) and \(B\) by \(f\) are
disjoint.
\item Let moreover \(A,B\) be analytic, \(A\) with empty interior, \(A\cap B\) compact, and \(E_0\) a divisor with simple normal crossings. There is an admissible blowup
sequence of \((X,E_0)\) with centres over \(A\cap B\) (Definition~\ref{hir:admissible-sequence}), with a divisor with simple normal crossings, after which the strict
transforms of \(A\) and \(B\) are disjoint.
\end{enumerate}
\end{lemma}
\begin{proof}
\leanok
(1) Off \(A\) the constant \(1\) lies in \(\mathcal I_A\), and off \(B\) in \(\mathcal I_B\); at a point of \(A\cap B\) every germ of \(\mathcal I_A\) and of
\(\mathcal I_B\), hence every combination of them, vanishes.

(2) Suppose that \(z\) lies in both strict transforms, and let \(h\) be a local generator of \(\mathcal O_Y(-D)\) on an open neighbourhood \(U\) of \(z\); it vanishes
identically on no nonempty open subset of \(U\). Since \(h\) lies in \(f^*(\mathcal I_A+\mathcal I_B)_z\), it is a holomorphic combination
\(h=\sum_ia_i\,(\tau_i\circ f)\) near \(z\) of germs \(\tau_i\) of \(\mathcal I_{A,f(z)}\) or of \(\mathcal I_{B,f(z)}\). Each \(\tau_i\circ f\) lies in
\(\mathcal O_Y(-D)_z\), so \(\tau_i\circ f=u_ih\) near \(z\) with \(u_i\) holomorphic. Hence \(h\cdot1=h\cdot\Phi\) with \(\Phi=\sum_ia_iu_i\), and cancelling \(h\)
(Lemma~\ref{ana:invertible-colon}) gives \(\Phi=1\) near \(z\); so \(u_i(z)\neq0\) for some \(i\). Say \(\tau_i\in\mathcal I_{A,f(z)}\); the other case is symmetric. Near
\(z\) we have \(u_i\neq0\), and every \(y\) with \(f(y)\in A\) satisfies \(\tau_i(f(y))=0\), hence \(h(y)=0\), so \(y\in\operatorname{Supp}D\) and \(f(y)\in B\). But \(z\) lies
in the closure of \(f^{-1}(A)\setminus\operatorname{Exc}(f)\), so such a point \(y\) exists arbitrarily near \(z\); \(f\) is open near \(y\)
(Lemma~\ref{ana:strict-transform-composition}), and since \(A\setminus B\) is dense in \(A\), the image of a small neighbourhood of \(y\) contains a point of \(A\setminus B\):
it is the image of a point near \(z\) over \(A\) and not over \(B\), a contradiction.

(3) \(\mathfrak a=\mathcal I_A+\mathcal I_B\) is coherent (Theorem~\ref{ana:cartan-coherence}; a sum of coherent ideals is coherent, generated by the union of local
generators), and nonzero at every point since \(\mathcal I_A\) is (\(A\) has empty interior; proof of Theorem~\ref{hir:embedded-resolution-pure-dimensional}); its zero locus
\(A\cap B\) is compact by (1). Principalization relative to \(E_0\) (Theorem~\ref{hir:principalization-boundary}) gives an admissible sequence \(f\) of \((X,E_0)\) with centres
over \(A\cap B\) and \(D\) with \(f^*\mathfrak a=\mathcal O(-D)\) and \(\operatorname{Supp}D=f^{-1}(A\cap B)\); apply (2).
\end{proof}

\begin{theorem}[Embedded resolution of an analytic set by one blowup sequence]
\label{hir:embedded-resolution-analytic}
\uses{hir:separation,hir:admissible-calculus,hir:embedded-resolution-boundary,hir:embedded-resolution-dimension-parts,hir:admissible-sequence,hir:embedded-resolution-pure-dimensional,hir:embedded-resolution-exceptional,bcalc:normal-crossings-local,lag:regular-point-basics,mc:subvariety-thin,ana:strict-transform-composition,ana:blowups,hir:snc}
\source{bierstone-milman97-invent-math:page-0009,bierstone-milman97-invent-math:page-0010,wlodarczyk05-hironaka,kollar07-resolution}
\lean{TranscendentalNakamaye.Pending.HironakaEmbeddedAnalytic.IsGoodPart,TranscendentalNakamaye.Pending.HironakaEmbeddedAnalytic.IsGoodPart.isClosed,TranscendentalNakamaye.Pending.HironakaEmbeddedAnalytic.IsGoodPart.pure,TranscendentalNakamaye.Pending.HironakaEmbeddedAnalytic.IsGoodPart.interior_eq_empty,TranscendentalNakamaye.Pending.HironakaEmbeddedAnalytic.IsGoodPart.subset_closure_diff_biUnion,TranscendentalNakamaye.Pending.HironakaEmbeddedAnalytic.IsGoodPart.subset_closure_diff_of_inter_subset,TranscendentalNakamaye.Pending.HironakaEmbeddedAnalytic.IsGoodPart.inter,TranscendentalNakamaye.Pending.HironakaEmbeddedAnalytic.IsGoodPart.strictTransform,TranscendentalNakamaye.Pending.HironakaEmbeddedAnalytic.exists_sepOne,TranscendentalNakamaye.Pending.HironakaEmbeddedAnalytic.exists_resolution_list,TranscendentalNakamaye.Pending.HironakaEmbeddedAnalytic.singularSet,TranscendentalNakamaye.Pending.HironakaEmbeddedAnalytic.mem_singularSet_iff,TranscendentalNakamaye.Pending.HironakaEmbeddedAnalytic.singularSet_subset,TranscendentalNakamaye.Pending.HironakaEmbeddedAnalytic.isClosed_singularSet,TranscendentalNakamaye.Pending.HironakaEmbeddedAnalytic.notMem_dimensionPart_of_regularAtDim,TranscendentalNakamaye.Pending.HironakaEmbeddedAnalytic.inter_dimensionPart_subset_singularSet,TranscendentalNakamaye.Pending.HironakaEmbeddedAnalytic.diff_regularLocus_dimensionPart_subset_singularSet,TranscendentalNakamaye.Pending.HironakaEmbeddedAnalytic.dimensionPart_eq_empty_of_le,TranscendentalNakamaye.Pending.HironakaEmbeddedAnalytic.eq_iUnion_dimensionPart_fin,TranscendentalNakamaye.Pending.HironakaEmbeddedAnalytic.dimensionPart_subset_closure_diff,TranscendentalNakamaye.Pending.HironakaEmbeddedAnalytic.preimage_singularSet_subset_exceptionalLocus,TranscendentalNakamaye.Pending.HironakaEmbeddedAnalytic.exists_embeddedResolution_analytic_boundary_of_isCompact,TranscendentalNakamaye.Pending.HironakaEmbeddedAnalytic.exists_embeddedResolution_analytic_of_isCompact,TranscendentalNakamaye.Pending.HironakaEmbeddedAnalytic.exists_embeddedResolution_analytic,TranscendentalNakamaye.Pending.HironakaEmbeddedAnalytic.exists_embeddedResolution_analytic_boundary}
\leanok
Let \(X\) be a complex manifold of dimension \(n\) and \(W\subset X\) a compact analytic set with empty interior, whose irreducible components may have different
dimensions (for instance any analytic set with empty interior in a compact \(X\)). Let \(W^{(k)}=\overline{\operatorname{Reg}_k(W)}\) be its dimension parts
(Theorem~\ref{hir:embedded-resolution-dimension-parts}) and \(W_{\mathrm{sing}}\) the set of points of \(W\) at which \(W\) is not a submanifold of any dimension.
\begin{enumerate}
\item There are an admissible blowup sequence \(f\colon X'\to X\) of \((X,\emptyset)\) with centres over \(W_{\mathrm{sing}}\) (Definition~\ref{hir:admissible-sequence})
and its divisor \(E\) such that \(E\) has simple normal crossings, \(\operatorname{Exc}(f)=\bigcup E=f^{-1}(W_{\mathrm{sing}})\), \(f\) is an isomorphism over
\(X\setminus W_{\mathrm{sing}}\), and the strict transform \(W'\) of \(W\) is a closed submanifold of \(X'\) with only normal crossings with \(E\); \(W'\) is the union of the
strict transforms \(W'^{(k)}\) of the \(W^{(k)}\), \(k<n\), which are pairwise disjoint, and \(W'\) has dimension \(k\) at every point of \(W'^{(k)}\).
\item Let \(E_0=(H_1,\dots,H_s)\) be a divisor with simple normal crossings such that \(W\subset\overline{W\setminus|E_0|}\), and \(F=W_{\mathrm{sing}}\cup(W\cap|E_0|)\).
There are an admissible blowup sequence \(f\) of \((X,E_0)\) with centres over \(F\) and its divisor \(E=(H_1',\dots,H_s',E_1,\dots,E_r)\) such that \(E\) has simple normal
crossings, \(\operatorname{Exc}(f)\subset E_1\cup\dots\cup E_r\subset f^{-1}(F)\), \(f^{-1}(W_{\mathrm{sing}})\subset\operatorname{Exc}(f)\), \(f\) is an isomorphism over
\(X\setminus F\), and the strict transform \(W'\) has the properties listed in (1).
\end{enumerate}
\end{theorem}
\begin{proof}
\leanok
Call \((A,d)\) a \emph{good part} relative to a divisor \(E\) with simple normal crossings and a closed set \(G\) if \(A\) is a compact analytic set, \(d<n\), the points at
which \(A\) is not a submanifold of dimension \(d\) and the points of \(A\cap|E|\) lie in \(G\), and \(A\setminus G\) is dense in \(A\). A good part is of pure dimension
\(d\) (its points off \(G\) are regular of dimension \(d\)), hence has empty interior, and no member of \(E\) contains a nonempty relatively open subset of it. By
Lemma~\ref{hir:admissible-calculus}(2),(4), if \(g\) is an admissible blowup sequence of \((X,E)\) with centres over \(G\), with divisor \(E'\), the strict transform of a good
part \((A,d)\) is a good part relative to \(E'\) and \(g^{-1}(G)\); and a good part relative to \(E,G\) is one relative to \(E,G\cap K\) for every closed \(K\supset A\).

\emph{Claim.} For good parts \((A_1,d_1),\dots,(A_m,d_m)\) relative to \(E,G\) with \(A_i\cap A_j\subset G\) for \(i\neq j\), there is an admissible blowup sequence \(f\) of
\((X,E)\) with centres over \(G\), with a divisor \(E'\) with simple normal crossings, such that each strict transform \(A_i'\) is regular of dimension \(d_i\) at each of its
points and has only normal crossings with \(E'\), and the \(A_i'\) are pairwise disjoint.

We argue by induction on \(m\); for \(m=0\) take the identity. Let \(m\ge1\).

\emph{Separation.} Separate \(A_1\) from \(A_2\), then (the strict transform of) \(A_1\) from (that of) \(A_3\), and so on. At each step the two current strict transforms are
good parts meeting in the current bad set, so each is the closure of its part off the other, and Lemma~\ref{hir:separation}(3) gives an admissible sequence with centres
over their intersection, which lies in the current bad set, after which their strict transforms are disjoint. A disjointness, once obtained, persists, because strict
transforms compose (Lemma~\ref{ana:strict-transform-composition}) and lie in preimages. The composite \(f_1\) is admissible with centres over \(G\)
(Lemma~\ref{hir:admissible-calculus}(1)), and the strict transform \(A_1^{(1)}\) of \(A_1\) misses the strict transforms \(A_j^{(1)}\), \(j\ge2\).

\emph{Resolution of the first part.} \(A_1^{(1)}\) is a good part relative to the divisor \(E^{(1)}\) of \(f_1\) and \(G^{(1)}=f_1^{-1}(G)\).
Theorem~\ref{hir:embedded-resolution-boundary}(1) gives an admissible sequence \(f_2\) of \((X^{(1)},E^{(1)})\) with centres over
\((A_1^{(1)})_{\mathrm{sing}}\cup(A_1^{(1)}\cap|E^{(1)}|)\subset A_1^{(1)}\cap G^{(1)}\), with divisor \(E^{(2)}\), after which the strict transform \(A_1^{(2)}\) of
\(A_1^{(1)}\) is regular of dimension \(d_1\) at each of its points and has only normal crossings with \(E^{(2)}\).

\emph{The other parts.} Let \(A_j^{(2)}\) (\(j\ge2\)) be the strict transforms of the \(A_j^{(1)}\) by \(f_2\), let \(K=A_2^{(2)}\cup\dots\cup A_m^{(2)}\), a closed set, and
\(G^{(2)}=f_2^{-1}(G^{(1)})\cap K\). The \((A_j^{(2)},d_j)\), \(j\ge2\), are good parts relative to \(E^{(2)},G^{(2)}\) meeting pairwise in \(G^{(2)}\). By induction there is
an admissible sequence \(f_3\) of \((X^{(2)},E^{(2)})\) with centres over \(G^{(2)}\) resolving them and making them pairwise disjoint. Since
\(A_1^{(2)}\subset f_2^{-1}(A_1^{(1)})\) and \(A_j^{(2)}\subset f_2^{-1}(A_j^{(1)})\), the set \(A_1^{(2)}\) misses \(K\supset G^{(2)}\). By
Lemma~\ref{hir:admissible-calculus}(3) the strict transform of \(A_1^{(2)}\) by \(f_3\) is \(f_3^{-1}(A_1^{(2)})\), regular of dimension \(d_1\) at each of its points, with only
normal crossings with the divisor of \(f_3\).

Take \(f=f_1\circ f_2\circ f_3\), admissible for \((X,E)\) with centres over \(G\) by Lemma~\ref{hir:admissible-calculus}(1) (the centres of \(f_2\) and \(f_3\) lie over
\(G^{(1)}\), resp. \(f_2^{-1}(G^{(1)})\)). Strict transforms compose; the strict transform of \(A_1\) lies over \(A_1^{(1)}\) and those of the \(A_j\) over the \(A_j^{(1)}\),
which proves the claim.

\emph{The theorem.} Take \(G=W_{\mathrm{sing}}\cup(W\cap|E_0|)\), with \(E_0=\emptyset\) in (1); it is closed, as \(W_{\mathrm{sing}}\) is (near a regular point of \(W\) all
points of \(W\) are regular of the same dimension). The parts are \((W^{(k)},k)\), \(k<n\): there is no part of dimension \(n\), since a regular point of dimension \(n\) is an
interior point (Lemma~\ref{mc:subvariety-thin}). A regular point of \(W\) of dimension \(k\) lies in no \(W^{(j)}\), \(j\neq k\): near it all points of \(W\) are regular of
dimension \(k\), and the dimension at a regular point is unique (Lemma~\ref{lag:regular-point-basics}). Hence \(W^{(k)}\cap W^{(j)}\subset W_{\mathrm{sing}}\) for \(k\neq j\),
and the points of \(W^{(k)}\) at which it is not a submanifold of dimension \(k\) lie in \(W_{\mathrm{sing}}\) (Theorem~\ref{hir:embedded-resolution-dimension-parts}(1)). The
regular points of \(W\) of dimension \(k\) off \(|E_0|\) are dense in \(W^{(k)}\): near a point of \(\operatorname{Reg}_k(W)\) all points of \(W\) are regular of dimension
\(k\), and \(W\setminus|E_0|\) is dense in \(W\). So the \((W^{(k)},k)\), \(k<n\), are good parts relative to \(E_0\) and \(G\), meeting pairwise in \(G\), and the claim
applies.

As \(W=\bigcup_{k<n}W^{(k)}\) and strict transforms commute with finite unions, \(W'=\bigcup_{k<n}W'^{(k)}\). Near a point of \(W'^{(k)}\) the other \(W'^{(j)}\), closed and
disjoint from it, are absent, so \(W'\) agrees there with \(W'^{(k)}\): it is regular of dimension \(k\) and has only normal crossings with \(E\)
(Lemma~\ref{bcalc:normal-crossings-local}). The statements on \(\operatorname{Exc}(f)\), the new members and the isomorphism are Lemma~\ref{hir:admissible-calculus}(2).
Finally \(f^{-1}(W_{\mathrm{sing}})\subset\operatorname{Exc}(f)\): at a non-exceptional point \(z\) over \(W\) the map \(f\) is a local biholomorphism and \(W'\) agrees with
\(f^{-1}(W)\) near \(z\), which is therefore regular at \(z\) of some dimension, since \(W'\) is a submanifold; so \(W\) is regular at \(f(z)\)
(Theorem~\ref{hir:embedded-resolution-exceptional}). With \(E_0=\emptyset\) this gives \(\operatorname{Exc}(f)\subset\bigcup E\subset f^{-1}(W_{\mathrm{sing}})\subset
\operatorname{Exc}(f)\).
\end{proof}
This is Bierstone--Milman's Theorem~1.6, with its consequence \(E'=\sigma^{-1}(\operatorname{Sing}X)\), for a closed analytic subset with compact support whose components may
have different dimensions, in the weak form given by the route through principalization: the centres lie over \(W_{\mathrm{sing}}\) (resp. \(F\)) but need not lie in the
singular locus of the current strict transform (condition (1) of Theorem~1.6 is not claimed), and with a boundary \(f\) may blow up smooth points of \(W\) on \(|E_0|\) at which
\(W\) already has normal crossings with \(E_0\).

\begin{theorem}[The public statements, in elementary terms]
\label{hir:public-statements}
\uses{hir:principalization-compact-support,hir:resolution-relatively-compact,hir:principalization-boundary,hir:embedded-resolution-compact-support,hir:embedded-resolution-pure-dimensional,hir:admissible-sequence}
\source{bierstone-milman97-invent-math:page-0009,bierstone-milman97-invent-math:page-0010,kollar07-resolution,wlodarczyk05-hironaka}
\lean{TranscendentalNakamaye.Pending.HironakaAPI.forall_isNonzeroAlongAt_univ_iff,TranscendentalNakamaye.Pending.HironakaAPI.exists_principalization,TranscendentalNakamaye.Pending.HironakaAPI.exists_principalization_of_isCompact,TranscendentalNakamaye.Pending.HironakaAPI.exists_principalization_of_subset_isCompact,TranscendentalNakamaye.Pending.HironakaAPI.exists_principalization_local,TranscendentalNakamaye.Pending.HironakaAPI.exists_principalization_boundary,TranscendentalNakamaye.Pending.HironakaAPI.exists_embeddedResolution,TranscendentalNakamaye.Pending.HironakaAPI.exists_embeddedResolution_of_isCompact,TranscendentalNakamaye.Pending.HironakaAPI.exists_embeddedResolution_of_subset_isCompact,TranscendentalNakamaye.Pending.HironakaAPI.exists_embeddedResolution_pure,TranscendentalNakamaye.Pending.HironakaAPI.exists_embeddedResolution_pure_of_isCompact}
\leanok
Let \(X\) be a complex manifold of dimension \(n\).
\begin{enumerate}
\item (Principalization.) Let \(\mathfrak a\subset\mathcal O_X\) be a coherent ideal such that at every point \(x\) some germ of \(\mathfrak a_x\) does not vanish
identically near \(x\). If \(X\) is compact, or more generally if \(Z(\mathfrak a)\) is compact, there are an admissible blowup sequence \(f\colon X'\to X\) of
\((X,\emptyset)\) with centres over \(Z(\mathfrak a)\) (Definition~\ref{hir:admissible-sequence}), with divisor \(E\), and an effective divisor \(D\) on \(X'\) such that
\(f^*\mathfrak a=\mathcal O_{X'}(-D)\), \(E\) has simple normal crossings, \(\operatorname{Supp}D=\bigcup E=f^{-1}(Z(\mathfrak a))\) and \(f\) is an isomorphism over
\(X\setminus Z(\mathfrak a)\). Without compactness the same holds over every open set contained in a compact set, in particular over a relatively compact neighbourhood
of any compact set; and, for \(Z(\mathfrak a)\) compact, relative to any divisor \(E_0\) with simple normal crossings (the sequence admissible for \((X,E_0)\), and
\(\operatorname{Supp}D\) the union of the members created by the blowups).
\item (Embedded resolution.) For a closed irreducible analytic subvariety \(V\) of a compact \(X\), or a compact one of any \(X\), there is an admissible blowup sequence
\(f\colon X'\to X\) of \((X,\emptyset)\) with centres over \(V_{\mathrm{sing}}\), with divisor \(E\) with simple normal crossings, such that
\(\operatorname{Exc}(f)=\bigcup E=f^{-1}(V_{\mathrm{sing}})\), \(f\) is an isomorphism over \(X\setminus V_{\mathrm{sing}}\), and the strict transform of \(V\) is a closed connected
submanifold of dimension \(\dim V\) with only normal crossings with \(E\). Without compactness the same holds, the strict transform being a closed submanifold of dimension
\(\dim V\) at each of its points, over every open set contained in a compact set. For a compact analytic set \(W\) of pure dimension \(d<n\), possibly reducible, the same holds
with \(W_{\mathrm{sing}}=W\setminus\operatorname{Reg}_d(W)\), the strict transform being a closed submanifold of dimension \(d\) at each of its points.
\end{enumerate}
\end{theorem}
\begin{proof}
\leanok
The hypothesis on \(\mathfrak a\) is the non-vanishing condition of Theorem~\ref{hir:principalization-compact-support} written out. (1) is
Theorem~\ref{hir:principalization-compact-support}, Theorem~\ref{hir:resolution-relatively-compact}(1) and Theorem~\ref{hir:principalization-boundary}; (2) is
Theorem~\ref{hir:embedded-resolution-compact-support}, Theorem~\ref{hir:resolution-relatively-compact}(2) and Theorem~\ref{hir:embedded-resolution-pure-dimensional}.
\end{proof}
These are the statements intended for a standalone library: none of them mentions marked ideals, towers, invariants or the domain of the invariants.

\begin{theorem}[The public statements, completed]
\label{hir:public-statements-complete}
\uses{hir:public-statements,hir:principalization,mc:factorization-domain,hir:principalization-boundary,hir:resolution-relatively-compact,hir:embedded-resolution-pure-dimensional,hir:embedded-resolution-boundary,hir:embedded-resolution-dimension-parts,hir:canonical-resolution,hir:canonical-centres}
\source{bierstone-milman97-invent-math:page-0009,bierstone-milman97-invent-math:page-0010,kollar07-resolution,wlodarczyk05-hironaka}
\lean{TranscendentalNakamaye.Pending.HironakaAPI2.not_le_vanishingIdeal_univ,TranscendentalNakamaye.Pending.HironakaAPI2.exists_principalization_of_connected,TranscendentalNakamaye.Pending.HironakaAPI2.exists_principalization_boundary_of_compactSpace,TranscendentalNakamaye.Pending.HironakaAPI2.exists_embeddedResolution_local,TranscendentalNakamaye.Pending.HironakaAPI2.exists_embeddedResolution_pure_of_subset_isCompact,TranscendentalNakamaye.Pending.HironakaAPI2.exists_embeddedResolution_boundary_of_isCompact,TranscendentalNakamaye.Pending.HironakaAPI2.exists_embeddedResolution_boundary,TranscendentalNakamaye.Pending.HironakaAPI2.exists_embeddedResolution_boundary_of_subset_isCompact,TranscendentalNakamaye.Pending.HironakaAPI2.eq_iUnion_dimensionPart,TranscendentalNakamaye.Pending.HironakaAPI2.isComplexAnalyticSubset_dimensionPart,TranscendentalNakamaye.Pending.HironakaAPI2.exists_embeddedResolution_dimensionPart_of_isCompact,TranscendentalNakamaye.Pending.HironakaAPI2.exists_embeddedResolution_dimensionPart,TranscendentalNakamaye.Pending.HironakaAPI2.exists_embeddedResolution_dimensionPart_boundary_of_isCompact,TranscendentalNakamaye.Pending.HironakaAPI2.canonicalResolution}
\leanok
Let \(X\) be a complex manifold of dimension \(n\). Besides the statements of Theorem~\ref{hir:public-statements}:
\begin{enumerate}
\item (Principalization.) If \(X\) is compact and connected and \(\mathfrak a\subset\mathcal O_X\) is a coherent ideal some germ of which does not vanish identically near
its point, the conclusion of Theorem~\ref{hir:public-statements}(1) holds. If \(X\) is compact, principalization relative to a divisor \(E_0\) with simple normal crossings
holds as in Theorem~\ref{hir:public-statements}(1).
\item (Embedded resolution.) For a closed irreducible analytic subvariety \(V\) of any \(X\) and a compact \(K\subset X\), the local form of
Theorem~\ref{hir:resolution-relatively-compact} holds over a relatively compact open neighbourhood of \(K\). For an analytic set of pure dimension \(d<n\), the
form of Theorem~\ref{hir:embedded-resolution-pure-dimensional}(2) holds over every open set contained in a compact set. Relative to a divisor \(E_0\) with simple
normal crossings, Theorem~\ref{hir:embedded-resolution-boundary} holds (\(W\) compact, \(X\) compact, or over an open set contained in a compact set).
\item (Analytic sets whose components have different dimensions.) An analytic set is the union of its dimension parts \(W^{(k)}=\overline{\operatorname{Reg}_k(W)}\),
\(k\le n\), which are analytic and of pure dimension \(k\), and each \(W^{(k)}\), \(k<n\), has an embedded resolution as in
Theorem~\ref{hir:embedded-resolution-dimension-parts} (\(W\) compact or \(X\) compact; also relative to a simple normal crossings boundary).
\item (Canonical resolution.) Over a compact \(X\), every marked ideal of the domain has the canonical resolution of Theorem~\ref{hir:canonical-resolution} for the
canonical invariant of Theorem~\ref{hir:canonical-centres}.
\end{enumerate}
\end{theorem}
\begin{proof}
\leanok
(1) If some germ of \(\mathfrak a\) does not vanish identically near its point, \(\mathfrak a\) is not contained in the ideal of \(X\); on a connected \(X\) a coherent
ideal which is not zero is nonzero at every point (Lemma~\ref{mc:factorization-domain}), and Theorem~\ref{hir:principalization} applies, in its admissible form
of Theorem~\ref{hir:public-statements}(1). The boundary form is Theorem~\ref{hir:principalization-boundary}. (2) is Theorem~\ref{hir:resolution-relatively-compact}(2) with
its last sentence, Theorem~\ref{hir:embedded-resolution-pure-dimensional}(2) and Theorem~\ref{hir:embedded-resolution-boundary}. (3) is
Theorem~\ref{hir:embedded-resolution-dimension-parts}. (4) is Theorem~\ref{hir:canonical-resolution} for the invariant of Theorem~\ref{hir:canonical-centres}.
\end{proof}
Together with Theorem~\ref{hir:public-statements} these are all the statements intended for the standalone library; only (4), the marked-ideal form of the canonical
resolution, mentions marked ideals and towers.

\begin{theorem}[The public statements, III: analytic sets whose components have different dimensions]
\label{hir:public-statements-analytic}
\uses{hir:public-statements,hir:public-statements-complete,hir:embedded-resolution-analytic,hir:separation,hir:admissible-calculus,hir:embedded-resolution-dimension-parts}
\source{bierstone-milman97-invent-math:page-0009,bierstone-milman97-invent-math:page-0010,kollar07-resolution,wlodarczyk05-hironaka}
\lean{TranscendentalNakamaye.Pending.HironakaAPI3.exists_embeddedResolution_analytic_of_isCompact,TranscendentalNakamaye.Pending.HironakaAPI3.exists_embeddedResolution_analytic,TranscendentalNakamaye.Pending.HironakaAPI3.exists_embeddedResolution_analytic_boundary_of_isCompact,TranscendentalNakamaye.Pending.HironakaAPI3.exists_embeddedResolution_analytic_boundary,TranscendentalNakamaye.Pending.HironakaAPI3.exists_separation,TranscendentalNakamaye.Pending.HironakaAPI3.isAdmissibleBlowupSequence_comp,TranscendentalNakamaye.Pending.HironakaAPI3.isIsomorphismOver_compl_of_isAdmissibleBlowupSequence,TranscendentalNakamaye.Pending.HironakaAPI3.exceptionalLocus_subset_of_isAdmissibleBlowupSequence}
\leanok
Let \(X\) be a complex manifold of dimension \(n\). Besides the statements of Theorems~\ref{hir:public-statements} and~\ref{hir:public-statements-complete}:
\begin{enumerate}
\item (Embedded resolution of an analytic set.) For a compact analytic set \(W\subset X\) with empty interior (for instance any analytic set with empty interior in a
compact \(X\)), whose components may have different dimensions, with \(W_{\mathrm{sing}}=W\setminus\bigcup_k\operatorname{Reg}_k(W)\) and dimension parts
\(W^{(k)}=\overline{\operatorname{Reg}_k(W)}\), the conclusions of Theorem~\ref{hir:embedded-resolution-analytic}(1) hold; relative to a divisor \(E_0\) with simple normal
crossings with \(W\subset\overline{W\setminus|E_0|}\), those of Theorem~\ref{hir:embedded-resolution-analytic}(2).
\item (Separation.) For analytic \(A,B\subset X\), \(A\) with empty interior, \(A\cap B\) compact, each the closure of its part off the other, and a divisor \(E_0\) with
simple normal crossings, the conclusion of Lemma~\ref{hir:separation}(3) holds.
\item (Admissible sequences.) Admissible blowup sequences compose as in Lemma~\ref{hir:admissible-calculus}(1); an admissible blowup sequence of \((X,E_0)\) with centres over
\(F\) is an isomorphism over \(X\setminus F\), and its exceptional locus lies in the members it creates, which lie over \(F\) (Lemma~\ref{hir:admissible-calculus}(2)).
\end{enumerate}
\end{theorem}
\begin{proof}
\leanok
(1) is Theorem~\ref{hir:embedded-resolution-analytic}, with \(W_{\mathrm{sing}}\) written out; (2) is Lemma~\ref{hir:separation}(3); (3) is
Lemma~\ref{hir:admissible-calculus}(1),(2).
\end{proof}
With Theorems~\ref{hir:public-statements} and~\ref{hir:public-statements-complete} these are the statements intended for the standalone library.


\begin{lemma}[Base change of admissible sequences along open embeddings]
\label{hir:base-change-open}
\uses{hir:admissible-sequence,ana:blowups,bcalc:restrict-open,hir:biholomorphism-blowup,hir:biholomorphism-invariance,hir:snc-invariance,ana:regular-transport,ana:strict-transform-composition,hir:snc}
\source{bierstone-milman08-functoriality,wlodarczyk05-hironaka}
\lean{TranscendentalNakamaye.Pending.HironakaBaseChange.mfderiv_val_bijective,TranscendentalNakamaye.Pending.HironakaBaseChange.exceptionalLocus_val,TranscendentalNakamaye.Pending.HironakaBaseChange.IsOpenHoloEmbedding,TranscendentalNakamaye.Pending.HironakaBaseChange.isOpenHoloEmbedding_val,TranscendentalNakamaye.Pending.HironakaBaseChange.IsOpenHoloEmbedding.contMDiff,TranscendentalNakamaye.Pending.HironakaBaseChange.IsOpenHoloEmbedding.exceptionalLocus_eq_empty,TranscendentalNakamaye.Pending.HironakaBaseChange.IsOpenHoloEmbedding.notMem_exceptionalLocus,TranscendentalNakamaye.Pending.HironakaBaseChange.IsOpenHoloEmbedding.isOpenEmbedding,TranscendentalNakamaye.Pending.HironakaBaseChange.IsOpenHoloEmbedding.continuous,TranscendentalNakamaye.Pending.HironakaBaseChange.IsOpenHoloEmbedding.isOpenMap,TranscendentalNakamaye.Pending.HironakaBaseChange.IsOpenHoloEmbedding.injective,TranscendentalNakamaye.Pending.HironakaBaseChange.IsOpenHoloEmbedding.isOpen_range,TranscendentalNakamaye.Pending.HironakaBaseChange.IsOpenHoloEmbedding.preimage_closure,TranscendentalNakamaye.Pending.HironakaBaseChange.IsOpenHoloEmbedding.pullsBackCoordinates,TranscendentalNakamaye.Pending.HironakaBaseChange.IsOpenHoloEmbedding.regularAtDim_preimage,TranscendentalNakamaye.Pending.HironakaBaseChange.IsOpenHoloEmbedding.hasNormalCrossings_preimage,TranscendentalNakamaye.Pending.HironakaBaseChange.IsOpenHoloEmbedding.isSNCDivisor_map_preimage,TranscendentalNakamaye.Pending.HironakaBaseChange.IsOpenHoloEmbedding.codRestrict,TranscendentalNakamaye.Pending.HironakaBaseChange.IsAdmissibleBlowupSequence.exists_baseChange,TranscendentalNakamaye.Pending.HironakaBaseChange.exceptionalLocus_eq_preimage,TranscendentalNakamaye.Pending.HironakaBaseChange.strictTransform_preimage_eq,TranscendentalNakamaye.Pending.HironakaBaseChange.IsAdmissibleBlowupSequence.exists_restrict_opens}
\leanok
Call a holomorphic map \(\iota\colon B\to A\) of complex manifolds of dimension \(n\) an \emph{open holomorphic embedding} if it is a biholomorphism of \(B\) onto an open
subset \(O\subset A\) followed by the inclusion; for instance the inclusion of an open subset, and its co-restriction to any larger open subset.
\begin{enumerate}
\item The inclusion of an open subset has bijective differential at every point; hence an open holomorphic embedding has no exceptional points, is an open embedding,
pulls back adapted coordinates, regular points of every dimension, normal crossings and divisors with simple normal crossings, and closures (\(\iota^{-1}\overline S=
\overline{\iota^{-1}S}\)).
\item (Base change.) Let \(f\colon Y\to A\) be an admissible blowup sequence of \((A,E_0)\) with centres over \(F\) (Definition~\ref{hir:admissible-sequence}), with divisor
\(E\), and \(\iota\colon B\to A\) an open holomorphic embedding. There are an admissible blowup sequence \(f'\colon Y'\to B\) of \((B,\iota^{-1}E_0)\) with centres over
\(\iota^{-1}(F)\), with divisor \(j^{-1}E\), and an open holomorphic embedding \(j\colon Y'\to Y\) with \(f\circ j=\iota\circ f'\) and \(j(Y')=f^{-1}(\iota(B))\).
\item (Transport.) In (2), \(\operatorname{Exc}(f')=j^{-1}\operatorname{Exc}(f)\), and for every \(S\subset A\) the strict transform of \(\iota^{-1}(S)\) by \(f'\) is
\(j^{-1}\) of the strict transform of \(S\) by \(f\).
\item (Restriction.) In particular, for an open \(O\subset A\) the sequence \(f\) restricts to an admissible blowup sequence of \((O,E_0|_O)\) with centres over
\(F\cap O\), whose total space is embedded by \(j\) onto \(f^{-1}(O)\), with \(\operatorname{Exc}\) and strict transforms as in (3).
\end{enumerate}
\end{lemma}
\begin{proof}
\leanok
(1) In the preferred charts the inclusion of an open subset is the identity near every point, so its differential is the identity; a biholomorphism has no exceptional
points (Lemma~\ref{hir:biholomorphism-invariance}), and \(\operatorname{Exc}(g\circ h)=\operatorname{Exc}(h)\cup h^{-1}\operatorname{Exc}(g)\) by the chain rule. Maps
without exceptional points pull back adapted coordinates and regular points (Lemmas~\ref{hir:snc-invariance} and~\ref{ana:regular-transport}). An open embedding is an
open map, so preimages commute with closures.

(2) By induction on \(f\). For the identity take \(Y'=B\), \(f'=\mathrm{id}\), \(j=\iota\). Let \(f\pi\) be obtained from \(f\) by a blowup \(\pi\colon Z\to Y\) along a
centre \(C\) having only normal crossings with \(E\) and lying over \(F\), and let \((Y',f',j)\) be given for \(f\), with \(j\) a biholomorphism \(\Phi\colon Y'\to O\)
onto the open set \(O=j(Y')\) followed by the inclusion. The restriction \(\pi^{-1}(O)\to O\) of \(\pi\) is the blowup of \(O\) along \(C\cap O\)
(Lemma~\ref{bcalc:restrict-open}), and composing it with \(\Phi^{-1}\) gives a blowup \(\sigma\colon\pi^{-1}(O)\to Y'\) along \(\Phi^{-1}(C\cap O)=j^{-1}(C)\)
(Lemma~\ref{hir:biholomorphism-blowup}). The centre \(j^{-1}(C)\) has only normal crossings with \(j^{-1}E\) by (1) and lies over \(\iota^{-1}(F)\) since
\(f\circ j=\iota\circ f'\). So \(f'\sigma\) is admissible; take for the new \(j\) the inclusion of \(\pi^{-1}(O)\) in \(Z\). Since \(j\sigma\) is \(\pi\) on
\(\pi^{-1}(O)\), the new divisor \(\big(\overline{\sigma^{-1}(j^{-1}H\setminus j^{-1}C)}\big)_{H\in E}\cup(\sigma^{-1}j^{-1}C)\) is the trace on \(\pi^{-1}(O)\) of
\(\big(\overline{\pi^{-1}(H\setminus C)}\big)_{H\in E}\cup(\pi^{-1}C)\) (closures commute with restriction to open sets), \(f\pi\) agrees with \(\iota f'\sigma\) on
\(\pi^{-1}(O)\), and \(\pi^{-1}(O)=\pi^{-1}f^{-1}(\iota(B))\).

(3) By the chain rule, \(j^{-1}\operatorname{Exc}(f)=\operatorname{Exc}(f\circ j)=\operatorname{Exc}(\iota\circ f')=\operatorname{Exc}(f')\), as \(\iota\) and \(j\) have no
exceptional points. Then \(f'^{-1}\iota^{-1}(S)\setminus\operatorname{Exc}(f')=j^{-1}(f^{-1}(S)\setminus\operatorname{Exc}(f))\), and closures commute with \(j^{-1}\) by
(1). (4) is (2) and (3) for the inclusion of \(O\).
\end{proof}
This is the functoriality of resolution sequences for open embeddings (Bierstone--Milman 2008, the remark after Theorem~1.3; Włodarczyk 2005, Proposition~2.4), for the
admissible sequences of Definition~\ref{hir:admissible-sequence}; no trivial blowup has to be inserted, because the restricted sequence starts from \(B\) itself.

\begin{theorem}[Embedded resolution of an analytic set over relatively compact open sets]
\label{hir:embedded-resolution-analytic-local}
\uses{hir:embedded-resolution-analytic,hir:base-change-open,hir:separation,hir:admissible-calculus,hir:boundary-relatively-compact,hir:embedded-resolution-boundary,hir:embedded-resolution-dimension-parts,hir:admissible-sequence,hir:principalization-local,ana:cartan-coherence,bcalc:normal-crossings-local,lag:regular-point-basics,mc:subvariety-thin,ana:strict-transform-composition,hir:snc}
\source{bierstone-milman97-invent-math:page-0009,bierstone-milman97-invent-math:page-0010,bierstone-milman08-functoriality,wlodarczyk05-hironaka}
\lean{TranscendentalNakamaye.Pending.HironakaEmbeddedAnalyticLocal.presentationUnion,TranscendentalNakamaye.Pending.HironakaEmbeddedAnalyticLocal.isComplexAnalyticSubset_union,TranscendentalNakamaye.Pending.HironakaEmbeddedAnalyticLocal.isComplexAnalyticSubset_biUnion_list,TranscendentalNakamaye.Pending.HironakaEmbeddedAnalyticLocal.IsLocGoodPart,TranscendentalNakamaye.Pending.HironakaEmbeddedAnalyticLocal.IsLocGoodPart.isClosed,TranscendentalNakamaye.Pending.HironakaEmbeddedAnalyticLocal.IsLocGoodPart.pure,TranscendentalNakamaye.Pending.HironakaEmbeddedAnalyticLocal.IsLocGoodPart.interior_eq_empty,TranscendentalNakamaye.Pending.HironakaEmbeddedAnalyticLocal.IsLocGoodPart.subset_closure_diff_biUnion,TranscendentalNakamaye.Pending.HironakaEmbeddedAnalyticLocal.IsLocGoodPart.subset_closure_diff_of_inter_subset,TranscendentalNakamaye.Pending.HironakaEmbeddedAnalyticLocal.IsLocGoodPart.inter,TranscendentalNakamaye.Pending.HironakaEmbeddedAnalyticLocal.IsLocGoodPart.strictTransform,TranscendentalNakamaye.Pending.HironakaEmbeddedAnalyticLocal.IsLocGoodPart.preimage,TranscendentalNakamaye.Pending.HironakaEmbeddedAnalyticLocal.strictTransform_comp_of_isOpenHoloEmbedding,TranscendentalNakamaye.Pending.HironakaEmbeddedAnalyticLocal.exists_separation_of_subset_isCompact,TranscendentalNakamaye.Pending.HironakaEmbeddedAnalyticLocal.exists_opens_of_range_subset,TranscendentalNakamaye.Pending.HironakaEmbeddedAnalyticLocal.exists_baseChange_of_subset,TranscendentalNakamaye.Pending.HironakaEmbeddedAnalyticLocal.exists_resolution_list_local,TranscendentalNakamaye.Pending.HironakaEmbeddedAnalyticLocal.preimage_val_singularSet_subset,TranscendentalNakamaye.Pending.HironakaEmbeddedAnalyticLocal.exists_embeddedResolution_analytic_boundary_of_subset_isCompact,TranscendentalNakamaye.Pending.HironakaEmbeddedAnalyticLocal.exists_embeddedResolution_analytic_of_subset_isCompact,TranscendentalNakamaye.Pending.HironakaEmbeddedAnalyticLocal.exists_embeddedResolution_analytic_local,TranscendentalNakamaye.Pending.HironakaEmbeddedAnalyticLocal.exists_embeddedResolution_analytic_boundary_local}
\leanok
Let \(X\) be a complex manifold of dimension \(n\), \(W\subset X\) an analytic set with empty interior, whose components may have different dimensions and which need
not be compact, \(W^{(k)}=\overline{\operatorname{Reg}_k(W)}\) its dimension parts and \(W_{\mathrm{sing}}\) the set of points of \(W\) at which \(W\) is not a submanifold of
any dimension. Let \(X_0\subset X\) be an open set contained in a compact set.
\begin{enumerate}
\item (Separation.) For analytic \(A,B\subset X\), \(A\) with empty interior, each the closure of its part off the other, and a divisor \(E_0\) with simple normal crossings,
there is an admissible blowup sequence of \((X_0,E_0|_{X_0})\) with centres over \(A\cap B\cap X_0\), with a divisor with simple normal crossings, after which the strict
transforms of \(A\cap X_0\) and \(B\cap X_0\) are disjoint (no compactness of \(A\cap B\) is needed).
\item There are an admissible blowup sequence \(f\colon X_0'\to X_0\) of \((X_0,\emptyset)\) with centres over \(W_{\mathrm{sing}}\cap X_0\) and its divisor \(E\) such that
\(E\) has simple normal crossings, \(\operatorname{Exc}(f)=\bigcup E=f^{-1}(W_{\mathrm{sing}}\cap X_0)\), \(f\) is an isomorphism over \(X_0\setminus W_{\mathrm{sing}}\), and the
strict transform \(W'\) of \(W\cap X_0\) is a closed submanifold of \(X_0'\) with only normal crossings with \(E\); \(W'\) is the union of the strict transforms
\(W'^{(k)}\) of the \(W^{(k)}\cap X_0\), \(k<n\), which are pairwise disjoint, and \(W'\) has dimension \(k\) at every point of \(W'^{(k)}\).
\item Let \(E_0\) be a divisor with simple normal crossings with \(W\subset\overline{W\setminus|E_0|}\), and \(F=W_{\mathrm{sing}}\cup(W\cap|E_0|)\). There are an
admissible blowup sequence \(f\) of \((X_0,E_0|_{X_0})\) with centres over \(F\cap X_0\) and its divisor \(E\) such that \(E\) has simple normal crossings,
\(\operatorname{Exc}(f)\) lies in the union of the members created by the blowups, which lies over \(F\cap X_0\), \(f^{-1}(W_{\mathrm{sing}}\cap X_0)\subset
\operatorname{Exc}(f)\), \(f\) is an isomorphism over \(X_0\setminus F\), and the strict transform \(W'\) of \(W\cap X_0\) has the properties listed in (2).
\item (Local form.) For every compact \(K\subset X\) there is an open neighbourhood \(U\) of \(K\) with compact closure over which (2) and (3) hold.
\end{enumerate}
\end{theorem}
\begin{proof}
\leanok
(1) The ideal \(\mathfrak a=\mathcal I_A+\mathcal I_B\) is coherent and nonzero at every point, with zero locus \(A\cap B\) (Lemma~\ref{hir:separation}(1) and the proof of
(3) there). Principalization relative to a boundary over the open set \(X_0\) contained in a compact set (Theorem~\ref{hir:boundary-relatively-compact}(1)) gives an
admissible sequence \(f\) of \((X_0,E_0|_{X_0})\) with centres over \(A\cap B\cap X_0\) and \(D\) with \(f^*(\mathfrak a|_{X_0})=\mathcal O(-D)\), \(\operatorname{Supp}D=
f^{-1}(A\cap B\cap X_0)\). With \(\iota\colon X_0\to X\) the inclusion, \((\iota f)^*\mathfrak a=f^*\iota^*\mathfrak a\), so Lemma~\ref{hir:separation}(2) applies to
\(\iota f\); and the strict transform of \(A\) by \(\iota f\) is that of \(A\cap X_0\) by \(f\) (strict transforms compose, Lemma~\ref{ana:strict-transform-composition},
and \(\iota\) has no exceptional points, Lemma~\ref{hir:base-change-open}(1)).

Call \((A,d)\) a \emph{locally good part} relative to \(E\) and a closed set \(G\) if it satisfies the conditions of a good part in the proof of
Theorem~\ref{hir:embedded-resolution-analytic} except compactness: \(A\) analytic, \(d<n\), the points where \(A\) is not a submanifold of dimension \(d\) and the points of
\(A\cap|E|\) in \(G\), and \(A\setminus G\) dense in \(A\). As there, locally good parts are of pure dimension \(d\), have empty interior, are carried along admissible
sequences with centres over \(G\) (Lemma~\ref{hir:admissible-calculus}(2),(4)), and the bad set may be cut down to any closed set containing the part; moreover they pull
back along open holomorphic embeddings (Lemma~\ref{hir:base-change-open}(1)). A finite union of analytic sets is analytic: near a point, the products
\(a_ib_j\) of local equations \(a_i\) of \(A\) and \(b_j\) of \(B\) cut out \(A\cup B\).

\emph{Claim.} For locally good parts \((A_1,d_1),\dots,(A_m,d_m)\) relative to \(E,G\) meeting pairwise in \(G\), and an open holomorphic embedding \(\iota\colon B\to X\)
whose image lies in a compact set, there is an admissible blowup sequence \(f\) of \((B,\iota^{-1}E)\) with centres over \(\iota^{-1}(G)\), with a divisor \(E'\) with
simple normal crossings, such that the strict transforms of the \(\iota^{-1}(A_i)\) are regular of dimension \(d_i\) at each of their points, have only normal crossings
with \(E'\), and are pairwise disjoint.

By induction on \(m\); for \(m=0\) take the identity of \(B\) (Lemma~\ref{hir:base-change-open}(1) for the divisor). Let \(m\ge1\). The closure \(\overline{\iota(B)}\) is
compact; let \(U_1\) be an open neighbourhood of it with compact closure. Let \(A'=A_2\cup\dots\cup A_m\), analytic; \(A_1\cap A'\subset G\), so \(A_1\) is the closure of
its part off \(A'\), and \(A'\) of its part off \(A_1\). (1) over \(U_1\) gives an admissible sequence \(g_1\) of \((U_1,E|_{U_1})\) with centres over \(A_1\cap A'\cap U_1
\subset G\), divisor \(E_1\), after which the strict transform \(A_1^{(1)}\) of \(A_1\cap U_1\) misses that of \(A'\cap U_1\), hence those \(A_j^{(1)}\) of the \(A_j\cap
U_1\), \(j\ge2\) (strict transforms are monotone). Base change along the co-restriction \(B\to U_1\) of \(\iota\) (Lemma~\ref{hir:base-change-open}(2)) gives an admissible
sequence \(g_1'\) of \((B,\iota^{-1}E)\) with centres over \(\iota^{-1}(G)\), divisor \(k_1^{-1}E_1\), and an open holomorphic embedding \(k_1\) into the total space of
\(g_1\), whose image \(g_1^{-1}(\iota(B))\) lies in the compact set \(g_1^{-1}(\overline{\iota(B)})\) (blowup sequences are proper,
Lemma~\ref{hir:admissible-calculus}(4)); the strict transform of \(\iota^{-1}(S)\) by \(g_1'\) is \(k_1^{-1}\) of that of \(S\cap U_1\) by \(g_1\)
(Lemma~\ref{hir:base-change-open}(3)). The \(A_j^{(1)}\) are locally good parts relative to \(E_1\) and \(G_1=g_1^{-1}(G\cap U_1)\).

Next \(A_1^{(1)}\) is resolved over an open neighbourhood \(U_2\) with compact closure of \(\overline{k_1(\cdot)}\), relative to \(E_1\)
(Theorem~\ref{hir:embedded-resolution-boundary}(2), over open sets contained in a compact set), by a sequence \(g_2\)
with centres over \((A_1^{(1)})_{\mathrm{sing}}\cup(A_1^{(1)}\cap|E_1|)\subset G_1\), after which the strict transform \(A_1^{(2)}\) of \(A_1^{(1)}\cap U_2\) is regular of
dimension \(d_1\) and has only normal crossings with the divisor \(E_2\). Base change along the co-restriction of \(k_1\) gives \(g_2'\) and \(k_2\) as before. Let
\(A_j^{(2)}\) be the strict transforms of the \(A_j^{(1)}\cap U_2\), \(j\ge2\), \(K_2\) their union and \(G_2=g_2^{-1}(G_1\cap U_2)\cap K_2\); they are locally good parts
relative to \(E_2,G_2\), meeting pairwise in \(G_2\), and \(A_1^{(2)}\) misses \(K_2\) as in the proof of Theorem~\ref{hir:embedded-resolution-analytic}. The induction
hypothesis along \(k_2\) gives \(h\) resolving the \(k_2^{-1}(A_j^{(2)})\) and making them disjoint, with centres over \(k_2^{-1}(G_2)\). Then \(f=g_1'g_2'h\) is admissible
for \((B,\iota^{-1}E)\) with centres over \(\iota^{-1}(G)\) (Lemma~\ref{hir:admissible-calculus}(1); the centres of \(g_2'\) and \(h\) lie over \(\iota^{-1}(G)\) since
\(g_1k_1=\iota g_1'\) and \(g_2k_2=k_1g_2'\)). By Lemma~\ref{hir:base-change-open}(3) twice, the strict transform of \(\iota^{-1}(A_j)\) by \(f\) is the strict transform by
\(h\) of \(k_2^{-1}(A_j^{(2)})\). For \(j=1\) the set \(k_2^{-1}(A_1^{(2)})\) misses \(k_2^{-1}(G_2)\), so by Lemma~\ref{hir:admissible-calculus}(3) its strict transform
is its preimage, regular of dimension \(d_1\) with only normal crossings with \(E'\) (regular points and normal crossings pull back along \(k_2\),
Lemma~\ref{hir:base-change-open}(1)); and it is disjoint from the others, which lie over the \(A_j^{(1)}\), \(j\ge2\). This proves the claim.

(2), (3) As in the proof of Theorem~\ref{hir:embedded-resolution-analytic}, with the claim for the inclusion \(\iota\colon X_0\to X\): the parts \((W^{(k)},k)\), \(k<n\), are
locally good parts relative to \(E_0\) and \(G=W_{\mathrm{sing}}\cup(W\cap|E_0|)\) meeting pairwise in \(G\); the strict transform of \(W\cap X_0\) is the union of those of the
\(W^{(k)}\cap X_0\); near a point of one of them it agrees with it; and \(f^{-1}(W_{\mathrm{sing}}\cap X_0)\subset\operatorname{Exc}(f)\) because a singular point of \(W\) in
\(X_0\) is a singular point of \(W\cap X_0\) (regular points are local) and the strict transform is a submanifold. With \(E_0=\emptyset\) this gives
\(\operatorname{Exc}(f)=\bigcup E=f^{-1}(W_{\mathrm{sing}}\cap X_0)\). (4) Manifolds are locally compact: take \(X_0=U\).
\end{proof}
This is Bierstone--Milman's Theorem~1.6 over relatively compact open sets (``finite over every relatively compact open subset''), for an analytic set whose components
may have different dimensions, in the weak form of Theorem~\ref{hir:embedded-resolution-analytic}. A finite sequence over all of a non-compact \(X\) is not claimed.

\begin{theorem}[The public statements, IV: analytic sets over relatively compact open sets; restriction to open sets]
\label{hir:public-statements-analytic-local}
\uses{hir:public-statements-analytic,hir:embedded-resolution-analytic-local,hir:base-change-open}
\source{bierstone-milman97-invent-math:page-0009,bierstone-milman97-invent-math:page-0010,bierstone-milman08-functoriality,wlodarczyk05-hironaka}
\lean{TranscendentalNakamaye.Pending.HironakaAPI4.exists_embeddedResolution_analytic_of_subset_isCompact,TranscendentalNakamaye.Pending.HironakaAPI4.exists_embeddedResolution_analytic_local,TranscendentalNakamaye.Pending.HironakaAPI4.exists_embeddedResolution_analytic_boundary_of_subset_isCompact,TranscendentalNakamaye.Pending.HironakaAPI4.exists_embeddedResolution_analytic_boundary_local,TranscendentalNakamaye.Pending.HironakaAPI4.exists_separation_of_subset_isCompact,TranscendentalNakamaye.Pending.HironakaAPI4.exists_restrict_opens}
\leanok
Let \(X\) be a complex manifold of dimension \(n\). Besides the statements of Theorems~\ref{hir:public-statements}, \ref{hir:public-statements-complete}
and~\ref{hir:public-statements-analytic}:
\begin{enumerate}
\item (Embedded resolution of an analytic set over relatively compact open sets.) For an analytic set \(W\subset X\) with empty interior, not necessarily compact, whose
components may have different dimensions, with \(W_{\mathrm{sing}}=W\setminus\bigcup_k\operatorname{Reg}_k(W)\), and an open set \(X_0\) contained in a compact set, the
conclusions of Theorem~\ref{hir:embedded-resolution-analytic-local}(2) hold; relative to a divisor \(E_0\) with simple normal crossings with
\(W\subset\overline{W\setminus|E_0|}\), those of Theorem~\ref{hir:embedded-resolution-analytic-local}(3); and both hold over a relatively compact open neighbourhood of any
compact set.
\item (Separation.) The conclusion of Theorem~\ref{hir:embedded-resolution-analytic-local}(1) holds.
\item (Restriction.) An admissible blowup sequence restricts to every open set as in Lemma~\ref{hir:base-change-open}(4).
\end{enumerate}
\end{theorem}
\begin{proof}
\leanok
(1) is Theorem~\ref{hir:embedded-resolution-analytic-local}(2)--(4), with \(W_{\mathrm{sing}}\) written out; (2) is Theorem~\ref{hir:embedded-resolution-analytic-local}(1);
(3) is Lemma~\ref{hir:base-change-open}(4).
\end{proof}
With Theorems~\ref{hir:public-statements}, \ref{hir:public-statements-complete} and~\ref{hir:public-statements-analytic} these are the statements intended for the
standalone library.
